% !TeX program = xelatex
% 《分析讲义》—— LECTURE NOTES IN ANALYSIS (2011) 中译本主文件
% Sergiu Klainerman, Princeton University
% 排版: 数译(mathtranslations.org) mathtranslation.cls v3.1
% 编译: XeLaTeX 连续编译两次(生成目录、交叉引用、术语索引页码).
% 依赖: mathtranslation.cls, logo.pdf
%
% 相对 mathtranslation.cls v3.1 默认行为的补充见 preamble.tex:
%   - 定理类环境共享同一计数器, 编号"节号.序号"(与原书一致);
%   - \mtsubsec / \mtsubsecn 节内带编号/独立编号子标题;
%   - \thmref 等按类型引用命令; 未编号 theorem* 等环境;
%   - mybibliography 兼容环境.

\documentclass[
  BookTitleCN={分析讲义},
  BookTitleEN={LECTURE NOTES IN ANALYSIS (2011)},
  OriginalAuthor={Sergiu Klainerman},
  OriginalEdition={2011 年讲义(普林斯顿大学)},
  OriginalPublisher={Princeton University},
  OriginalYear={2011},
  Translator={Fingolfin},
  ModelUsed={DeepSeek-V4-Flash+混元3},
  TranslationDate={28 Aug 2026}
]{mathtranslation}

% ================================================================
% 《分析讲义》中译本 —— 类文件之上的兼容层（mathtranslation.cls v3.1）
% ----------------------------------------------------------------
% mathtranslation.cls v3.1 基于 ctexbook，已原生提供 \chapter/\part/\appendix、
% 元信息命令 \BookTitleCN 等、\newterm/\addterm 术语系统、\printterminology、
% \headfont 标题字体、\makecover 封面、\longproof、\exercises/\answers、
% mtfang/mtsong 定理样式与 envenddiamond 结束标记。
%
% 因此本文件只保留 v3.1 未提供、且本书确需的部分；不改动正文内容：
%   1. 定理类环境共享同一计数器，编号“节号.序号”（与原书一致）。
%   2. \mtsubsec{标题} / \mtsubsecn{数字}{标题}：节内带编号/独立编号子标题。
%   3. 未编号环境 theorem* / definition* / ...（原书未编号对象用）。
%   4. \thmref 等按类型引用命令。
%   5. mybibliography 兼容环境（v3.0 起类文件只提供 \makebibliography）。
% ================================================================
% ---------- 1. 元信息命令 ----------
% v3.1 已把 \BookTitleCN / \OriginalAuthor / \Translator 等全部暴露为公开命令
% （见类文件第 7 节），旧版为 v1.2 写的 \MT@ 别名层在此整体删除。

% ---------- 2. 共享定理计数器 ----------
% v3.1 把定理类环境各自注册为 [section] 计数，各环境互相独立。本书要求它们
% 共用同一计数器，编号显示为“节号.序号”，且序号全书连续（与原书一致）。
% 故在此解除原定义并重建。
%
% 只 \let 环境命令，不释放 \c@theorem 等计数器本身：v3.1 已把这些计数器登记
% 进 \cl@section 的重置链，销毁计数器会让后续 \section 报
% “No counter 'definition' defined.”（报错位置会被误指到 \section 一行）。
\makeatletter
\let\theorem\relax\let\endtheorem\relax
\let\definition\relax\let\enddefinition\relax
\let\lemma\relax\let\endlemma\relax
\let\corollary\relax\let\endcorollary\relax
\let\proposition\relax\let\endproposition\relax
\let\remark\relax\let\endremark\relax
\let\example\relax\let\endexample\relax
\let\problem\relax\let\endproblem\relax
\makeatother

% 共享计数器：全书连续（不随 chapter/section 重置），显示“节号.序号”。
\newcounter{mtthm}
\renewcommand{\themtthm}{\arabic{section}.\arabic{mtthm}}

% 定理样式直接沿用类文件的 mtfang / mtsong（例、问题用宋体），不再另行定义。
\theoremstyle{mtfang}
\newtheorem{theorem}[mtthm]{定理}
\newtheorem{definition}[mtthm]{定义}
\newtheorem{lemma}[mtthm]{引理}
\newtheorem{corollary}[mtthm]{推论}
\newtheorem{proposition}[mtthm]{命题}
\newtheorem{remark}[mtthm]{注}
\theoremstyle{mtsong}
\newtheorem{example}[mtthm]{例}
\newtheorem{problem}[mtthm]{问题}

% 未编号版本（原书未编号的定义、定理、注等）。声明顺序照旧，因此它们沿用
% 上面生效的 mtsong 样式（宋体），与原书观感一致。
\newtheorem*{theorem*}{定理}
\newtheorem*{definition*}{定义}
\newtheorem*{lemma*}{引理}
\newtheorem*{corollary*}{推论}
\newtheorem*{proposition*}{命题}
\newtheorem*{remark*}{注}
\newtheorem*{example*}{例}
\newtheorem*{problem*}{问题}
\newtheorem*{conjecture*}{猜想}

% 类文件已用 \AtEndEnvironment 给 example/remark 挂了菱形结束标记（钩子存于
% \@end@<env>@hook，由通用 \end 调用，故不必也不可重复挂载，否则会出现双菱形）。
% 下面只为新定义的带 * 未编号环境补挂，保持与编号版本外观一致。
\AtEndEnvironment{example*}{\envenddiamond}
\AtEndEnvironment{remark*}{\envenddiamond}

% 节标题显示为“节号.”（原书样式，不带章号）。
\renewcommand{\thesection}{\arabic{section}.}

% ---------- 3. 节内子标题 ----------
% 带编号子标题: 推进共享计数器并显示 "节号.序号 标题".
\newcommand{\mtsubsec}[1]{%
  \par\medskip\noindent\refstepcounter{mtthm}%
  {\headfont\arabic{section}.\arabic{mtthm}\quad #1}\par\smallskip
}
% 独立编号子标题: 只显示 "数字. 标题", 不推进计数器.
% 用法: \mtsubsecn{数字}{标题}(数字为 0 或空时只显示标题).
\newcommand{\mtsubsecn}[2]{%
  \par\medskip\noindent%
  {\headfont\ifx\relax#1\relax\else #1.\quad\fi#2}\par\smallskip
}

% ---------- 4. 按类型区分的引用命令 ----------
% 所有定理类环境共享同一计数器, 故 \autoref 无法区分类型(一律显示"定理").
% 引用时用下列命令: \thmref{...} → "定理 N", \lemref{...} → "引理 N" 等.
\newcommand{\thmref}[1]{定理~\ref{#1}}
\newcommand{\lemref}[1]{引理~\ref{#1}}
\newcommand{\defref}[1]{定义~\ref{#1}}
\newcommand{\propref}[1]{命题~\ref{#1}}
\newcommand{\corref}[1]{推论~\ref{#1}}
\newcommand{\remref}[1]{注~\ref{#1}}
\newcommand{\exref}[1]{例~\ref{#1}}
\newcommand{\probref}[1]{问题~\ref{#1}}

% ---------- 5. 术语系统 ----------
% v3.1 的 \newterm{键}{中文}{英文} 与旧版行为完全一致（首次出现自动注册、
% 自动取页码、自动去重，并输出到 \MTTermList），故删除此处为 v1.2 写的
% 重新实现，直接使用类文件版本。

% ---------- 6. 公式/图表编号 ----------
% ctexbook（book 基类）默认把 equation / figure / table 按章编号（\theequation =
% 章号.序号）。原书与旧版 ctexart 类均为全书连续编号（公式 (1)--(446)，见"译者
% 说明"），故解除按章重置、恢复连续编号，保持与原文及既有成果一致。
\makeatletter
\@removefromreset{equation}{chapter}
\renewcommand{\theequation}{\arabic{equation}}
\@removefromreset{figure}{chapter}
\renewcommand{\thefigure}{\arabic{figure}}
\@removefromreset{table}{chapter}
\renewcommand{\thetable}{\arabic{table}}
\makeatother

% ---------- 6b. mybibliography 兼容环境 ----------
% v1.0 提供 mybibliography（参考文献列表，标题由正文手动 \section* 给出）。
% v1.2 已删除该环境；此处按 v1.2 的 thebibliography 的列表部分重建（不含标题），
% 供 partapx.tex 的"参考文献"使用。
\makeatletter
\newenvironment{mybibliography}[1]{%
  \begingroup
  \list{\@biblabel{\@arabic\c@enumiv}}%
       {\settowidth\labelwidth{\@biblabel{#1}}%
        \leftmargin\labelwidth
        \advance\leftmargin\labelsep
        \@openbib@code
        \usecounter{enumiv}%
        \let\p@enumiv\@empty
        \renewcommand\theenumiv{\@arabic\c@enumiv}}%
  \sloppy
  \clubpenalty4000
  \@clubpenalty \clubpenalty
  \widowpenalty4000%
  \sfcode`\.\@m
}{%
  \def\@noitemerr{\@latex@warning{Empty bibliography environment}}%
  \endlist
  \endgroup
}
\makeatother

% ---------- 7. 封面 ----------
% v3.1 的 \makecover 已按制作规范印出“译者 xx 翻译及重排”，并逐项拼接
% 版次·出版社·年份（空字段自动跳过），旧版为 v1.2 重写的封面宏在此整体删除。


\begin{document}

% ================================================================
%                             封面
% ================================================================
\makecover

% ================================================================
%                       译者说明与目录
% ================================================================
\frontmatter
\pagenumbering{roman}
\setcounter{page}{1}
\pagestyle{plain}

\chapter*{译者说明}
\addcontentsline{toc}{chapter}{译者说明}

本译本依据 Sergiu Klainerman 的讲义《LECTURE NOTES IN ANALYSIS (2011)》(普林斯顿大学数学系)翻译并重排. 原书分为三个部分: 第一部分``偏微分方程导论''(含导言``PDE 的世界''与第 1--4 章), 第二部分``分析的基本工具''(第 5--6 章), 第三部分``Minkowski 空间中波动方程的分析''(第 7--8 章), 其后为第 9 章``双线性估计''、附录 A``Riemann 与 Lorentz 几何的基本概念''与参考文献. 原书末尾附有作者论文《几何与 Fourier 方法在非线性波动方程中的应用》(GEOMETRIC AND FOURIER METHODS FOR NONLINEAR WAVE EQUATIONS), 因其为论文(article)而非专著(book), 本译本已将其单独抽取, 以简洁的 article 模板另行排版(见同目录 ipam/ 下独立文档).

\textbf{翻译原则.} 本译本遵循数译(mathtranslations.org)的《数学译本制作规范》: 忠于原文、数学正确优先于语言流畅、术语统一、保留原书数学结构、实现可点击的交叉引用. 正文中文使用宋体, 首次出现的专用术语以楷体标注并附英文原词, 定理、定义、命题、推论、引理、注使用仿宋, 英文与数学符号使用衬线 Computer Modern 字体. 除句子末尾使用英文句点外, 其余标点均使用中文标点.

\textbf{编号说明.} 为与原书一致, 定理类对象的编号按节进行, 显示为``节号.序号''(如``定理 3.4''); 无编号节的章(第 7、8 章)与各章引言部分显示为``0.序号''. 原书各节内带编号的小标题与定理共享同一计数序列(如第 1 章第 3 节的``3.5''、``3.6''), 本译本亦照此处理. 公式编号全书连续((1)--(446)), 与原文一致; 原书所附论文(已单独排版)中的公式另行编号. 原书中的个别编号跳号或重复(如第 4 章引言从命题 0.3 开始、第 6 章两处``注 2.3'')均如实保留.

\textbf{模型使用说明.} 本译本由数译团队使用大语言模型辅助翻译并逐章人工校对, 数学公式与符号依照原书排版还原. 原书 PDF 后半部分(所附论文)存在字体编码损坏, 其内容依据 OCR 识别结果重新排版, 个别无法还原的公式已在正文中以占位符标注.

\vspace{0.5cm}
\hfill \Translator

\hfill \TranslationDate

\newpage
\begingroup
\pagestyle{empty}
\thispagestyle{empty}
\tableofcontents
\clearpage
\endgroup

% ================================================================
%                         正文
% ================================================================
\mainmatter
\pagenumbering{arabic}
\setcounter{page}{1}
\pagestyle{fancy}

\part{偏微分方程导论}

% ==== part00 第一部分导言: PDE 的世界 ====
% 源: MinerU_latex_lecture nots.tex 93-867 (对应 pdf_text.txt 1-587, PDF 第 4-14 页)
% 术语: 偏微分方程=pde, 常微分方程=ode, 积分方程=integral-equation, 函数方程=functional-equation,
%       Laplace 方程=laplace-equation, 热方程=heat-equation, 波动方程=wave-equation, 演化方程=evolution-equation,
%       Schrödinger 方程=schrodinger-equation, Klein-Gordon 方程=klein-gordon-equation, 叠加原理=superposition,
%       线性=linear, 非线性=nonlinear, 极小曲面=minimal-surface, 微分算子=differential-operator, 局部性=locality,
%       不变群=invariance-group, 李群=lie-group, Dirichlet 问题=dirichlet-problem, 边界条件=boundary-condition,
%       初值条件=initial-condition, 初值问题=initial-value-problem, 边值问题=boundary-value-problem,
%       Cauchy-Riemann 方程组=cauchy-riemann, 复解析函数=complex-analytic, 确定方程组=determined-system,
%       超定(方程组)=overdetermined, 欠定(方程组)=underdetermined, 恰当形式=exact-form, 外微分=exterior-derivative,
%       de Rham 复形=de-rham-complex, D'Alembert 算子=dalembertian, Minkowski 内积=minkowski-product,
%       Lorentz 变换=lorentz-transformation, Galilei 不变=galilean-invariance, Dirac 算子=dirac-operator,
%       黎曼面=riemann-surface, Gauss 曲率=gauss-curvature, Gauss-Bonnet 公式=gauss-bonnet-formula,
%       欧拉示性数=euler-characteristic, 单值化定理=uniformization-theorem, 共形变换=conformal-transformation,
%       Laplace-Beltrami 算子=laplace-beltrami-operator, Ricci 流=ricci-flow, 调和坐标=harmonic-coordinates,
%       Ricci 曲率=ricci-curvature, 截面曲率=sectional-curvature, Einstein 真空方程=einstein-vacuum-equations,
%       Ricci 平坦=ricci-flat, Lorentz 流形=lorentzian-manifold, Lorentz 度量=lorentz-metric, 波坐标=wave-coordinates,
%       规范场论=gauge-field-theory, 规范变换=gauge-transformation, 控制论=control-theory,
%       随机微分方程=stochastic-differential-equation, Brown 运动=brownian-motion, Black-Scholes 方程=black-scholes-equation
% REFS:
%   - "我们将在第 2 节讨论"(section 2): 指向本部分(导言)下一节 "2. ...", 节标签由后续 part 文件给出, 待解析.
%   - 原书未解析交叉引用 "see (??)"(Ricci 流调和坐标段落, Δ_g 定义附近): 按原文保留 "(??)", 疑指 (21) 或 (8), 待人工确认.
%   - 原书笔误, 已按原文保留: (7) 右端 Σ e^{ij} X_i X_j 应为 X_i Y_j; (9) 末项 X_4Y_4 应为 X_3Y_3;
%     γ 矩阵列出 γ^0, γ^1, γ^3, γ^4(缺 γ^2); "Schrödinger 方程 (278)" 疑为 (4) 之误.
%   - 原书笔误, 已修正: 调和坐标 x=(x^1,x^2,x^2) 修正为 (x^1,x^2,x^3); ∂_i^2 = ∂²/∂x_i² 修正为 ∂²/(∂x^i)²;
%     E=m c^2 与 |x|^2=a^2 中 OCR 多余 dot 已去除.
%   - "见纲要文章"(compendium article)两处: 非本书章节引用, 保留文字即可.
%   - 脚注 1-4 已从 pdf_text.txt 页尾找回并翻译(可逆变换/Dirac 算子 Euclidean 版本/Ricci 流体积/波坐标).
% 编号说明: 本部分含 1 个编号定理 (Theorem 1.1, 标签 thm:c0-1.1) 与 28 个编号公式 (1)-(28);
%           正文 "Example."(Dirichlet 问题) 原书无编号, 译为 "例."; 9 个行间公式原书无编号, 用 \[ \] 不编号.

\section{PDE 的世界}

与常微分方程或积分方程一样, \newterm{pde}{偏微分方程}{partial differential equation}也是\newterm{functional-equation}{函数方程}{functional equation}.
也就是说, 我们试图确定的未知量(一个或多个)都是函数.
就偏微分方程(PDE)而言, 这些函数要从一些方程中确定, 这些方程除了通常的加法与乘法运算之外, 还涉及函数的偏导数.
下面是最基本的几个例子:

\begin{itemize}
\item (\newterm{laplace-equation}{Laplace 方程}{Laplace equation})
\begin{equation}\label{eq:c0-1}
\Delta u = 0
\end{equation}
其中 $\Delta u = \frac{\partial^{2}}{\partial x^{2}} u + \frac{\partial^{2}}{\partial y^{2}} u + \frac{\partial^{2}}{\partial z^{2}} u$.
在基础数学定义一节中描述的其他两个例子是
\item (\newterm{heat-equation}{热方程}{heat equation})
\begin{equation}\label{eq:c0-2}
-\partial_{t} u + k \Delta u = 0,
\end{equation}
\item (\newterm{wave-equation}{波动方程}{wave equation})
\begin{equation}\label{eq:c0-3}
-\partial_{t}^{2} u + c^{2} \Delta u = 0.
\end{equation}
\end{itemize}

在这两种情形下, 都要寻找一个依赖于变量 $t, x, y, z$ 的函数 $u$, 使之满足相应的方程.
注意 \eqref{eq:c0-2} 与 \eqref{eq:c0-3} 都含有与第一个方程中含义相同的符号 $\Delta$, 即
$\Delta u = \left(\frac{\partial^{2}}{\partial x^{2}} + \frac{\partial^{2}}{\partial y^{2}} + \frac{\partial^{2}}{\partial z^{2}}\right) u = \frac{\partial^{2} u}{\partial x^{2}} + \frac{\partial^{2} u}{\partial y^{2}} + \frac{\partial^{2} u}{\partial z^{2}}$.
在 \eqref{eq:c0-2} 与 \eqref{eq:c0-3} 中, $k>0$ 与 $c$ 都是固定常数(前者表示扩散速率, 后者表示光速).
只需研究特殊情形 $k = 1$ 与 $c = 1$ 下的方程求解即可.
事实上, 例如若 $u(t, x, y, z)$ 是 \eqref{eq:c0-3} 的解, 则 $v(t, x, y, z) = u(t, x/c, y/c, z/c)$ 满足 $c = 1$ 情形的同一方程.
这两个方程都称为\newterm{evolution-equation}{演化方程}{evolution equation}, 原因很简单:它们被用来描述某个特定物理对象相对于时间参数 $t$ 的变化.
注意 \eqref{eq:c0-1} 可以解释为 \eqref{eq:c0-3} 与 \eqref{eq:c0-2} 的特殊情形.
事实上, \eqref{eq:c0-3} 或 \eqref{eq:c0-2} 的与 $t$ 无关的解 $u = u(t, x, y, z)$, 即 $\partial_{t} u = 0$, 都满足 \eqref{eq:c0-1}.

下面是一些重要的 PDE 的其他例子:

\begin{itemize}
\item (\newterm{schrodinger-equation}{Schrödinger 方程}{Schr\"odinger equation})
\begin{equation}\label{eq:c0-4}
i \partial_{t} u + k \Delta u = 0
\end{equation}
其中 $u : \mathbb{R} \times \mathbb{R}^{3} \to \mathbb{C}$.
该方程描述一个有质量粒子的量子演化, $k = \frac{\hbar}{2 m}$, 其中 $\hbar > 0$ 是 Planck 常数, $m$ 是粒子的质量.
与热方程一样, 可以通过简单的变量替换把 $k$ 归一化为 1.
尽管该方程在形式上与热方程非常相似, 它的定性行为却大不相同.
在研究 PDE 时务必牢记:方程形式的微小改变可能导致解的性质发生很大变化.
\item (\newterm{klein-gordon-equation}{Klein-Gordon 方程}{Klein-Gordon equation})
\begin{equation}\label{eq:c0-5}
-\partial_{t}^{2} u + c^{2} \Delta u - \left(\frac{m c^{2}}{\hbar}\right)^{2} u = 0
\end{equation}
它是 Schrödinger 方程的相对论对应物;参数 $m$ 具有质量的物理含义, 而 $m c^{2}$ 具有静能量的物理含义(这反映了 Einstein 的著名方程 $E = m c^{2}$).
通过对时间与空间同时作适当的变量替换, 可以把常数 $c$ 与 $\frac{m c^{2}}{\hbar}$ 归一化, 使二者都等于 1.
\item (\newterm{minimal-surface}{极小曲面}{minimal surface})
\begin{equation}\label{eq:c0-6}
\partial_{x}\left(\frac{\partial_{x} u}{(1 + |\partial_{x} u|^{2} + |\partial_{y} u|^{2})^{\frac{1}{2}}}\right) + \partial_{y}\left(\frac{\partial_{y} u}{(1 + |\partial_{x} u|^{2} + |\partial_{y} u|^{2})^{\frac{1}{2}}}\right) = 0.
\end{equation}
这里 $\partial_{x}$ 与 $\partial_{y}$ 是偏导数 $\frac{\partial}{\partial x}$ 与 $\frac{\partial}{\partial y}$ 的简写记号.
\end{itemize}

注意上面提到的三个 PDE 都满足下列称为\newterm{superposition}{叠加原理}{principle of superposition}的简单性质:若 $u_{1}, u_{2}$ 是某个方程的解, 则它们的任意线性组合 $\lambda_{1} u_{1} + \lambda_{2} u_{2}$(其中 $\lambda_{1}, \lambda_{2}$ 是任意实数)也是该方程的解.
这样的方程称为\newterm{linear}{线性}{linear}方程.
下面这个关于未知函数 $u = u(x, y)$ 的方程显然不是线性的:

到目前为止我们遇到的方程都可以写成 $\mathcal{P}[u] = 0$ 的形式, 其中 $\mathcal{P}$ 是作用在 $u$ 上的\newterm{differential-operator}{微分算子}{differential operator}.
微分算子不过是这样一条规则:它把定义在 $\mathbb{R}^{n}$ 或它的某个开子集上的函数 $u$ 映为函数 $\mathcal{P}[u]$, 所执行的运算如下:

\begin{itemize}
\item 我们可以取 $u$ 关于 $\mathbb{R}^{n}$ 的变量 $x = (x^{1}, x^{2}, \ldots, x^{n})$ 的偏导数 $\partial_{i} u = \frac{\partial u}{\partial x^{i}}$.
还允许 $u$ 的高阶偏导数, 如混合二阶偏导数 $\partial_{i} \partial_{j} u = \frac{\partial^{2} u}{\partial x^{i} \partial x^{j}}$, 或 $\partial_{i}^{2} = \frac{\partial^{2}}{(\partial x^{i})^{2}}$.

方程 \eqref{eq:c0-2} 的伴随微分算子为 $\mathcal{P} = -\partial_{t} + \Delta$, 而方程 \eqref{eq:c0-3} 的为 $-\partial_{t}^{2} + \Delta$.
\item 可以对 $u$ 及其偏导数彼此相加、相乘, 也可以与变量 $x$ 的给定函数相加、相乘.
还可能出现与给定函数的复合.
\end{itemize}

对于方程 \eqref{eq:c0-1}, 伴随微分算子为 $\mathcal{P} = \Delta = \partial_{1}^{2} + \partial_{2}^{2} + \partial_{3}^{2} = \sum_{i, j = 1}^{3} e^{i j} \partial_{i} \partial_{j}$, 其中 $e^{i j}$ 是以 (1, 1, 1) 为元素的对角 $3 \times 3$ 矩阵, 对应于 $\mathbb{R}^{3}$ 中向量 $X, Y$ 的 Euclidean 内积

\begin{equation}\label{eq:c0-7}
\langle X, Y \rangle = X_{1} Y_{1} + X_{2} Y_{2} + X_{3} Y_{3} = \sum_{i, j = 1}^{3} e^{i j} X_{i} X_{j}.
\end{equation}

方程 \eqref{eq:c0-2}、\eqref{eq:c0-3} 与 \eqref{eq:c0-4} 的伴随微分算子分别是 $\mathcal{P} = -\partial_{t} + \Delta$、$\mathcal{P} = -\partial_{t}^{2} + \Delta$ 与 $\mathcal{P} = i \partial_{t} + \Delta$, 其变量为 $t, x^{1}, x^{2}, x^{3} \in \mathbb{R}^{1+3}$.
对于波动方程 \eqref{eq:c0-3} 这一特殊情形, 把变量 $t$ 记为 $x^{0}$ 是有益的.
波算子于是可以写成如下形式:

\begin{equation}\label{eq:c0-8}
\Box = -\partial_{0}^{2} + \partial_{1}^{2} + \partial_{2}^{2} + \partial_{3}^{2} = \sum_{\alpha, \beta = 0}^{3} m^{\alpha \beta} \partial_{\alpha} \partial_{\beta}
\end{equation}

其中 $m^{\alpha \beta}$ 是以 $(-1, 1, 1, 1)$ 为元素的对角 $4 \times 4$ 矩阵, 对应于 $\mathbb{R}^{1+3}$ 中的\newterm{minkowski-product}{Minkowski 内积}{Minkowski scalar product}.
这后一种内积对 $4$ 向量 $\boldsymbol{X} = (X_{0}, X_{1}, X_{2}, X_{3})$ 与 $Y = (Y_{0}, Y_{1}, Y_{2}, Y_{3})$ 定义为

\begin{equation}\label{eq:c0-9}
m(X, Y) = \sum_{\alpha, \beta = 0}^{3} m^{\alpha \beta} X_{\alpha} Y_{\beta} = -X_{0} Y_{0} + X_{1} Y_{1} + X_{2} Y_{2} + X_{4} Y_{4}
\end{equation}

微分算子 $\Box$ 以法国数学家的名字命名为\newterm{dalembertian}{D'Alembert 算子}{D'Alembertian}, 他是第一位在振动弦方程中引入该算子的数学家.

注意, 方程 \eqref{eq:c0-1}--\eqref{eq:c0-4} 的伴随微分算子都是线性的, 即

\[
\mathcal{P}[\lambda u + \mu v] = \lambda \mathcal{P}[u] + \mu \mathcal{P}[v],
\]

对任意函数 $u, v$ 与任意实数或复数 $\lambda, \mu$ 成立.
下面是线性微分算子的另一个简单例子

\begin{equation}\label{eq:c0-10}
\mathcal{P}[u] = a_{1}(x) \partial_{1} u + a_{2}(x) \partial_{2} u
\end{equation}

其中 $\boldsymbol{x} = (x_{1}, x_{2})$, $a_{1}, a_{2}$ 是 $x$ 的给定函数, 称为线性算子的系数.
形如

\begin{equation}\label{eq:c0-11}
\mathcal{P}[u] = f,
\end{equation}

对应线性微分算子 $\mathcal{P}$ 与给定函数 $f = f(x)$ 的方程称为线性非齐次方程.
这种方程的任意解 $u$ 都可以表示为 $\boldsymbol{u} = \boldsymbol{u}_{0} + \boldsymbol{v}$ 的形式, 其中 $u_{0}$ 是 \eqref{eq:c0-11} 的特解, 而 $v$ 是齐次方程

\begin{equation}\label{eq:c0-12}
\mathcal{P}[v] = 0.
\end{equation}

的解.

对于方程 \eqref{eq:c0-6}, 微分算子 $\mathcal{P}$ 相对于变量 $x^{1}$ 与 $x^{2}$ 可以写成如下形式:

\[
\mathcal{P}[u] = \sum_{i = 1}^{2} \partial_{i}\left(\frac{1}{(1 + |\partial u|^{2})^{\frac{1}{2}}} \partial_{i} u\right),
\]

其中 $|\partial u|^{2} = (\partial_{1} u)^{2} + (\partial_{2} u)^{2}$.
显然在这种情况下 $\mathcal{P}[u]$ 不是线性的.
我们称它为\newterm{nonlinear}{非线性}{nonlinear}算子, 相应的方程 \eqref{eq:c0-6} 称为非线性方程.
线性与非线性微分算子都有一个重要性质:\newterm{locality}{局部性}{locality}.
这意味着当我们把 $\mathcal{P}$ 作用于某个在开集 $D$ 中为零的函数 $u$ 时, 所得函数 $\mathcal{P}[u]$ 也在 $D$ 中为零.

还要注意, 我们的方程 \eqref{eq:c0-1}--\eqref{eq:c0-4} 还是平移不变的.
以 \eqref{eq:c0-1} 为例, 这意味着只要函数 $u = u(x)$ 是解, 函数 $u_{c}(x) := u(T_{c} x)$ 也是解, 其中 $T_{c}$ 是平移 $T_{c}(x) = x + c$.
另一方面, 由 \eqref{eq:c0-10} 定义的算子 $\mathcal{P}$ 所对应的方程 $\mathcal{P}[u] = 0$ 则不具有这种性质, 除非系数 $a_{1}, a_{2}$ 是常数.
显然, 把 $\mathcal{P}[u] = 0$ 的任意解 $u = u(x)$ 映为另一个解 $u_{T}(x) := u(T x)$ 的可逆变换\footnote{这些变换通常是线性映射.}$T : \mathbb{R}^{n} \to \mathbb{R}^{n}$ 的全体构成一个群, 称为方程的\newterm{invariance-group}{不变群}{invariance group}.
两个对称的复合仍然是对称, 对称的逆也是对称, 因此自然地把对称的集合看作构成一个群(通常是有限维或无限维\newterm{lie-group}{李群}{Lie group}).

Laplace 方程 \eqref{eq:c0-1} 不仅关于平移不变, 还关于旋转不变, 即保持 Euclidean 内积 \eqref{eq:c0-7} 的线性变换 $O : \mathbb{R}^{3} \to \mathbb{R}^{3}$, 也就是说, 对所有向量 $X, Y \in \mathbb{R}^{3}$ 有 $\langle O X, O Y \rangle = \langle X, Y \rangle$.
类似地, 波动方程 \eqref{eq:c0-3} 与 Klein-Gordon 方程 \eqref{eq:c0-5} 在\newterm{lorentz-transformation}{Lorentz 变换}{Lorentz transformation}下不变, 即保持 Minkowski 内积 \eqref{eq:c0-9} 的线性变换 $L : \mathbb{R}^{1+3} \to \mathbb{R}^{1+3}$, 也就是说, $m(L X, L Y) = m(X, Y)$.
我们的其他演化方程 \eqref{eq:c0-2} 与 \eqref{eq:c0-4} 显然关于空间变量 $x = (x^{1}, x^{2}, x^{3}) \in \mathbb{R}^{3}$ 的旋转不变, 其中 $t$ 保持固定.
它们还是\newterm{galilean-invariance}{Galilei 不变}{Galilean invariant}的, 这意味着, 就 Schrödinger \eqref{eq:c6-278} 这一特殊情形而言, 只要 $u = u(t, x)$ 是解, 则对任意向量 $v \in \mathbb{R}^{3}$, $u_{v}(t, x) = e^{i (x \cdot v)} e^{i t |v|^{2}} u(t, x - v t)$ 也是解.

到目前为止, 我们一直默认方程在整个空间中成立:Laplace 方程在 $\mathbb{R}^{3}$ 中, 热方程、波动方程与 Schrödinger 方程在 $\mathbb{R}^{4}$ 中, 极小曲面方程在 $\mathbb{R}^{2}$ 中.
实际上, 人们常常被限制在相应空间的一个区域上.
例如, 方程 \eqref{eq:c0-1} 通常在有界开区域上研究, 并附加以指定的\newterm{boundary-condition}{边界条件}{boundary condition}.
下面是一个典型例子.

\textbf{例.} 在 $\mathbb{R}^{3}$ 的开区域 $D$ 上的\newterm{dirichlet-problem}{Dirichlet 问题}{Dirichlet problem}是:求定义在 $D$ 的闭包 $\bar{D}$ 上、在 $D$ 内二次连续可微的连续函数 $u$, 使得在 $D$ 中 $\Delta u = 0$, 并且 $u$ 在 $D$ 的边界 $\partial D$ 上的限制被指定为某个连续函数 $u_{0}$.
更精确地说, 我们要求

\begin{equation}\label{eq:c0-13}
u|_{\partial D} = u_{0}
\end{equation}

对于 \eqref{eq:c0-6} 的解也可以施加同样的边界条件, 此时 $D$ 是 $\mathbb{R}^{2}$ 中的有界开区域.
在 $D$ 中满足边界条件 \eqref{eq:c0-13} 的 \eqref{eq:c0-6} 的解 $u = u(x, y)$ 解决了 Plateau 问题, 即在 $\mathbb{R}^{3}$ 中寻找经过给定曲线的极小曲面.
可以证明, 由图像 $\Gamma_{u} = \{(x, y, u(x, y)) \mid (x, y) \in D \subset \mathbb{R}^{2}\}$ 给出的曲面在所有满足相同边界条件 $v|_{\partial D} = u_{0}$ 的图像曲面 $\Gamma_{v}$ 中具有最小面积.

对演化方程 \eqref{eq:c0-2}--\eqref{eq:c0-4} 也可以施加自然的边界条件.
最简单的是指定 $u$ 在超平面 $t = 0$ 上的值.
对于热方程与 Schrödinger 方程, 我们令

\[
u|_{t = 0} = u_{0}
\]

而对于涉及 $t$ 的二阶导数的波动方程, 则施加两个条件

\begin{equation}\label{eq:c0-14}
u|_{t = 0} = u_{0} \quad \text{且} \quad \partial_{t} u|_{t = 0} = u_{1}
\end{equation}

其中 $u_{0}, u_{1}$ 是坐标 $(x, y, z)$ 的函数, 称为\newterm{initial-condition}{初值条件}{initial condition}.
在两种情形下解\newterm{initial-value-problem}{初值问题}{initial value problem}, 都意味着寻找方程在 $t > 0$ 时、且在 $t = 0$ 处满足相应初值条件的解.
此外, 还可以把变量 $(x, y, z)$ 限制在 $\mathbb{R}^{3}$ 的开区域 $D$ 上.
更进一步, 可以尝试在区域 $[0, \infty) \times D$ 上解\newterm{boundary-value-problem}{边值问题}{boundary value problem}:在 $[0, \infty) \times \partial D$ 上施加诸如 \eqref{eq:c0-13} 那样的边界条件, 并在 $t = 0$ 处施加初值条件.

对给定的 PDE, 边界条件与初值条件的选择非常重要.
找出哪些是恰当的边界条件与初值条件, 是 PDE 一般理论的一个重要方面, 我们将在第 2 节讨论.
对物理上有意义的方程, 这些条件会从其推导背景中自然地出现.
例如, 振动弦由一维波动方程 $\partial_{t}^{2} u - \partial_{x}^{2} u = 0$ 在区域 $(a, b) \times \mathbb{R}$ 中的解描述, 初值条件 $u = u_{0}$, $\partial_{t} u = u_{1}$(在 $t = t_{0}$ 处)相当于指定弦的初始位置与速度.
另一方面, 边界条件 $u(a) = u(b) = 0$ 只表示弦的两端被固定.

到目前为止我们只考虑了含一个未知函数的方程.
实际上, 许多重要的方程都以偏微分方程组的形式出现.
下面这个重要例子含有两个未知函数 $u_{1} = u_{1}(x^{1}, x^{2})$, $u_{2} = u_{2}(x^{1}, x^{2})$, 它们满足

\begin{itemize}
\item (\newterm{cauchy-riemann}{Cauchy-Riemann 方程组}{Cauchy-Riemann equations})
\begin{equation}\label{eq:c0-15}
\partial_{1} u_{2} - \partial_{2} u_{1} = 0, \qquad \partial_{1} u_{1} + \partial_{2} u_{2} = 0.
\end{equation}
\end{itemize}

Cauchy 首先观察到, $u = u_{1} + i u_{2}$ 作为 $z = x^{1} + i x^{2}$ 的函数是\newterm{complex-analytic}{复解析函数}{complex analytic function}, 当且仅当 \eqref{eq:c0-15} 成立.
再令 $\bar{\partial} = \partial_{1} + i \partial_{2}$, 注意 \eqref{eq:c0-15} 等价于

\begin{equation}\label{eq:c0-16}
\bar{\partial} u = 0.
\end{equation}

方程 \eqref{eq:c0-15} 也可以写成 $\mathcal{P}[u] = 0$ 的形式:引入列向量 $u = (u_{1}, u_{2})$ 以及微分算子 $\mathcal{P}[u]$,

\[
\mathcal{P}[u] = \begin{pmatrix} -\partial_{2} & \partial_{1} \\ \partial_{1} & \partial_{2} \end{pmatrix} \cdot \binom{u_{1}}{u_{2}}
\]

方程组 \eqref{eq:c0-15} 含有两个方程与两个未知量, 这是\newterm{determined-system}{确定方程组}{determined system}的标准情形.
若方程组所含方程多于未知量, 则称为\newterm{overdetermined}{超定方程组}{over-determined system};若所含方程少于未知量, 则称为\newterm{underdetermined}{欠定方程组}{underdetermined system}.
例如, 含两个方程与一个未知量的方程组 $\partial_{x} u(x, y) = f$, $\partial_{y} u(x, y) = g$ 显然是超定的.
解存在的必要条件是 $\partial_{y} f = \partial_{x} g$, 这个条件可以解释为要求一次形式 $w = f(x, y) d x + g(x, y) d y$ 是\newterm{exact-form}{恰当形式}{exact form}, 即它的\newterm{exterior-derivative}{外微分}{exterior derivative} $d \omega$ 恒等于零.
诸如\newterm{de-rham-complex}{de Rham 复形}{de Rham complex}之类的超定方程组在几何中起着非常重要的作用.

注意 \eqref{eq:c0-15} 是线性方程组.
还要注意, 算子 $\mathcal{P}$ 具有如下引人注目的性质:

\[
\mathcal{P}^{2}[u] = \mathcal{P}[\mathcal{P}[u]] = \binom{\Delta u_{1}}{\Delta u_{2}}
\]

换言之, $\mathcal{P}^{2} = \Delta \cdot I$, 其中 $I$ 是恒等算子 $I[u] = u$, 因此可以把 $\mathcal{P}$ 看作 $\Delta$ 的平方根.
对 D'Alembert 算子 $\Box$ 也可以定义类似类型的平方根.
为此我们需要满足下述性质的 $4 \times 4$ 复矩阵 $\gamma^{0}, \gamma^{1}, \gamma^{3}, \gamma^{4}$

\begin{equation}\label{eq:c0-17}
\gamma^{\alpha} \gamma^{\beta} + \gamma^{\beta} \gamma^{\alpha} = -2 m^{\alpha \beta} I
\end{equation}

其中 $I$ 是 $4 \times 4$ 单位矩阵, $m^{\alpha \beta}$ 如 \eqref{eq:c0-8} 所示.
利用 $\gamma$ 矩阵, 我们可以引入\newterm{dirac-operator}{Dirac 算子}{Dirac operator}, 它作用于定义在 $\mathbb{R}^{1+3}$ 上、取值于 $\mathbb{C}^{4}$ 的 $u = (u_{1}, u_{2}, u_{3}, u_{4})$, 定义为

\begin{equation}\label{eq:c0-18}
\mathcal{D} u = i \gamma^{\alpha} \partial_{\alpha} u
\end{equation}

利用 \eqref{eq:c0-17} 我们容易验证 $\mathcal{D}^{2} u = \Box u$.
因此 Dirac 算子 $\mathcal{D}$ 可以看作 D'Alembert 算子 $\Box$ 的平方根.
它导出了如下基本方程, 这是 Dirac 为自由、有质量、相对论性粒子(如电子)引入的方程:

\begin{itemize}
\item (Dirac 方程)
\begin{equation}\label{eq:c0-19}
\mathcal{D} u = k u
\end{equation}
\end{itemize}

偏微分方程在数学与科学中无处不在.
它们为一些最重要的物理理论提供了基本数学框架, 例如弹性力学、流体力学、电磁学、广义相对论与非相对论量子力学.
更现代的相对论量子场论原则上导出了含无穷多个未知量的方程, 这超出了偏微分方程的范围.
然而, 即使在这种情况下, 基本方程仍保留着 PDE 的局部性.
而且, 量子场论的出发点总是一个由 PDE 方程组描述的经典场论.
例如, 基于 Yang-Mills-Higgs 场论的弱相互作用与强相互作用标准模型就是如此.
如果再加上经典力学的常微分方程(它可以看作一维 PDE), 我们就会看到, 本质上全部物理学都是由微分方程描述的.
支撑我们一些最基本物理理论的偏微分方程的其他例子, 是 Maxwell、Einstein、Euler 与 Navier-Stokes 方程组.
注意, 每一个方程都处于整个物理学领域的核心, 即电动力学、广义相对论与流体力学.

物理学中出现的主要 PDE 的一个重要特征是它们表面上的普遍性.
例如, 波动方程最初由 D'Alembert 引入以描述振动弦的运动, 后来被发现与声波和电磁波的传播有关.
热方程最初由 Fourier 引入以描述热的传播, 也出现在许多其他耗散效应起重要作用的场合.
对于 Laplace、Schrödinger 方程以及许多其他基本方程, 情形也同样如此.

更令人惊讶的是, 最初为描述具体物理现象而引入的方程, 在那些被视为纯粹数学的领域, 如复分析、微分几何、拓扑与代数几何中, 也发挥着基本作用.
例如, 研究全纯函数性质的复分析, 可以被看作研究 Cauchy-Riemann 方程组 \eqref{eq:c0-15} 在 $\mathbb{R}^{2}$ 的某个区域中的解.
在拓扑与代数几何中起基本作用的 Hodge 理论, 其基础是研究流形上一类推广 Cauchy-Riemann 方程组的线性 PDE 方程组的解空间.
Atiyah-Singer 指标定理则是以流形上一类特殊的线性 PDE 来表述的, 它们与 Dirac 算子 \eqref{eq:c0-18} 的 Euclidean 版本有关\footnote{这是把 \eqref{eq:c0-17} 中的 Minkowski 度量 $m$ 换成 Euclidean 度量 $e$ 后得到的算子.}.

几何中的重要问题往往可以归结为求特定偏微分方程(通常是非线性的)的解.
在求经过给定曲线的、具有最小总面积的曲面的 Plateau 问题中, 我们已经看到了这样的例子.
著名的\newterm{uniformization-theorem}{单值化定理}{uniformization theorem}提供了另一个极好的例子.

为了陈述单值化定理, 我们需要回忆紧\newterm{riemann-surface}{黎曼面}{Riemann surface} $S$ 的定义.
它是一个配备光滑正定度量 $g$ 的二维紧流形.
\newterm{gauss-curvature}{Gauss 曲率}{Gauss curvature} $K = K(g)$ 是曲面的一个重要不变量, 它可以在每一点 $p \in S$ 处, 由 $p$ 附近局部坐标系 $x = (x^{1}, x^{2})$ 下的分量 $g_{a b}$ 显式算出.
计算涉及分量 $g_{a b}$ 关于 $x^{1}$ 与 $x^{2}$ 的一阶与二阶偏导数.
一个显著的事实是, $K$ 的最终值不依赖于进行计算的特定坐标系.
此外, 在 $S$ 是 $\mathbb{R}^{3}$ 中由方程 $|x|^{2} = a^{2}$ 给出的标准球面的特殊情形下, Gauss 曲率等于与我们对曲率的直觉相符的期望值, 即 $K = a^{-2}$.
Gauss 曲率的另一个显著性质是, 它沿 $S$ 的全积分不依赖于度量 $g$, 而只依赖于 $S$ 的拓扑性质.
更精确地说, 根据\newterm{gauss-bonnet-formula}{Gauss-Bonnet 公式}{Gauss-Bonnet formula}, 我们有

\[
\chi(S) = (2 \pi)^{-1} \int_{S} K \, d a_{g},
\]

其中 $d a_{g}$ 表示度量 $g$ 的面积元.
在坐标 $x^{1}, x^{2}$ 下有 $d a_{g} = \sqrt{|g|}\, d x^{1} d x^{2}$, 其中 $|g|$ 是矩阵 $(g_{a b})_{a, b = 1, 2}$ 的行列式.
数 $\chi(S)$ 是整数 $2, 0, -2, \ldots, -2 k, \ldots$ 之一, 称为 $S$ 的\newterm{euler-characteristic}{欧拉示性数}{Euler characteristic}, 并且有简单的拓扑解释.
因此, 任何可以连续变形为标准球面的曲面都有 $\chi(S) = 2$, 而任何可以连续变形为环面的曲面都有 $\chi(S) = 0$.
现在我们可以陈述单值化定理:

\begin{theorem}\label{thm:c0-1.1}
设 $S$ 是带度量 $g$、Gauss 曲率 $K = K(g)$ 与 欧拉示性数 $\chi(S)$ 的二维紧 黎曼面.
则存在度量 $g$ 的一个\newterm{conformal-transformation}{共形变换}{conformal transformation}, 即 $\tilde{g} = \Omega^{2} g$, 其中 $\Omega$ 是某个光滑的非零函数, 使得新度量 $\tilde{g}$ 的 Gauss 曲率 $\tilde{K}$ 恒等于 $1$、$0$ 或 $-1$, 视 $\chi(S) > 0$、$\chi(S) = 0$ 或 $\chi(S) < 0$ 而定.
\end{theorem}

为了证明这个非常重要的几何结果——它按 欧拉示性数给出了所有紧曲面的完整分类——我们被引向 $S$ 上的一个非线性偏微分方程.
事实上, 假设 $\chi(S) = 2$, 于是我们希望度量 $\tilde{g} = e^{2 u} g$ 的 Gauss 曲率 $\tilde{K}$ 恰好为 1.
通过简单计算, 很容易把 $\tilde{K}$ 与原度量 $g$ 的 Gauss 曲率 $K$ 联系起来.
这就导出了关于 $u$ 的如下方程:

\begin{equation}\label{eq:c0-20}
\Delta_{S} u + e^{2 u} = K
\end{equation}

其中 $\Delta_{S}$, 称为 $S$ 的\newterm{laplace-beltrami-operator}{Laplace-Beltrami 算子}{Laplace-Beltrami operator}, 是 Laplace 算子(见 \eqref{eq:c0-1})到曲面 $S$ 上的直接推广.
因此, 单值化定理的证明归结为解方程 \eqref{eq:c0-20}, 即对具有 Gauss 曲率 $K$ 的给定曲面 $S$, 求一个满足 \eqref{eq:c0-20} 的实值函数 $u$.

下面给出算子 $\Delta_{S}$ 相对于开坐标卡 $D \subset S$ 上的局部坐标系 $x = (x^{1}, x^{2})$ 的精确定义.
记 $G(x) = (g_{a b}(x))_{a, b = 1, 2}$ 为以我们的 Riemann 度量在 $D$ 上的分量为元素的 $2 \times 2$ 矩阵.
设 $G^{-1}(x)$ 表示 $G(x)$ 的逆矩阵, 并记其分量为 $(g^{a b}(x))_{a, b = 1, 2}$.
于是, 对所有 $x \in D$,

\[
\sum_{c} g_{a c}(x) g^{c b}(x) = \delta_{a b}
\]

其中 $\delta_{a b}$ 是通常的 Kronecker 记号.
如前, 我们还令 $|g(x)| = \det(G(x))$, 并定义

\[
\Delta_{S} u(x) = \frac{1}{\sqrt{|g(x)|}} \sum_{a, b = 1, 2} \partial_{b}\left(\sqrt{|g(x)|}\, g^{a b}(x)\, \partial_{a} u(x)\right)
\]

通常我们在上述公式中略去对 $x$ 的显式依赖.
采用 Einstein 对重复指标求和的约定也非常方便, 于是可以写成

\begin{equation}\label{eq:c0-21}
\Delta_{S} u = \frac{1}{\sqrt{|g|}} \partial_{b}\left(\sqrt{|g|}\, g^{a b}\, \partial_{a} u\right)
\end{equation}

作为第三个例子, 我们考虑紧 $n$ 维流形 $M$ 上的\newterm{ricci-flow}{Ricci 流}{Ricci flow}方程, 它在《纲要》(Compendium)的某篇文章中有描述.
在三维这一特殊情形下, 该方程最近被决定性地用于给出 Thurston 几何化猜想(包括著名的 Poincaré 猜想)的第一个证明.
几何化猜想在《纲要》的拓扑部分有描述, 它是上述二维单值化定理在三维空间中的精确类比.
Ricci 流在 $M$ 上任意局部坐标 $x = (x^{1}, x^{2}, x^{3})$ 下由如下方程定义:

\begin{itemize}
\item (Ricci 流)
\begin{equation}\label{eq:c0-22}
\partial_{t} g_{i j} = R_{i j}(g)
\end{equation}
\end{itemize}

这里 $g_{i j} = g_{i j}(t)$ 是光滑依赖于参数 $t$ 的 Riemann 度量族, $R_{i j}(g)$ 表示度量 $g_{i j}$ 的\newterm{ricci-curvature}{Ricci 曲率}{Ricci curvature}.
它不过是在单值化定理中遇到的 Gauss 曲率的三维推广.
在给定坐标系中, $R_{i j}(g)$ 可以用度量系数 $g_{i j}$ 及其一阶与二阶偏导数算出.
由于 $g_{i j}$ 与 $R_{i j}$ 都关于 $i, j = 1, 2, 3$ 对称, 我们可以把 \eqref{eq:c0-22} 解释为由六个方程与六个未知量组成的非线性方程组.
仔细观察会发现, \eqref{eq:c0-22} 与热方程 \eqref{eq:c0-2} 有关.
事实上, 相对于称为\newterm{harmonic-coordinates}{调和坐标}{harmonic coordinates}的特定坐标系 $x = (x^{1}, x^{2}, x^{3})$ 作直接计算, 可以证明 Ricci 流 \eqref{eq:c0-22} 具有如下形式

\begin{equation}\label{eq:c0-23}
\partial_{t} g_{i j} - \Delta_{g} g_{i j} = N_{i j}(g, \partial g)
\end{equation}

其中每个 $N_{i j}$, $i, j = 1, 2, 3$, 都是分量 $g_{i j}$ 及其关于坐标 $x$ 的一阶偏导数的函数, 而 $\Delta_{g}$ 又是与 $\mathbb{R}^{3}$ 中的 Laplacian $\Delta$ 非常相似的微分算子(见 (??)).
更精确地说, 若 $G^{-1} = (g^{a b})_{a, b = 1, 2, 3}$ 表示 $G = (g_{a b})_{a, b = 1, 2, 3}$ 的逆矩阵, 则利用求和约定可以写成

\[
\Delta_{g} = g^{a b} \partial_{a} \partial_{b} = \sum_{a, b = 1}^{3} g^{a b} \partial_{a} \partial_{b}.
\]

在点 $p \in M$ 的一个小邻域内, 我们可以选取调和坐标 $x^{a}$, 使得 $g^{a b}(p) = \delta^{a b}$, 其中 $\delta^{a b}$ 表示通常的 Kronecker 记号.
于是, 在 $p$ 附近, $\Delta_{g}$ 看起来确实像 $\Delta = \delta^{a b} \partial_{a} \partial_{b}$.

Ricci 流\footnote{实际上需要把方程 \eqref{eq:c0-22} 稍作修改, 以确保 $M$ 关于度量 $g(t)$ 计算的总体积保持不变.}允许我们把 $M$ 上的任意 Riemann 度量变形为一个具有常\newterm{sectional-curvature}{截面曲率}{sectional curvature}的简单度量.
其思想是从一个度量 $g$ 出发, 寻找满足初值条件 $g(0) = g$ 的 \eqref{eq:c0-22} 的解 $g(t)$.
人们希望当 $t \to \infty$ 时, 度量 $g(t)$ 收敛到一个常曲率度量.
直观上可以这样理解:这与热方程 \eqref{eq:c0-2} 中随着 $t$ 增大热量在空间中均匀分布的方式相同.
事实上, 由于 \eqref{eq:c0-22} 与 \eqref{eq:c0-2} 相似, 我们预期随着度量按 \eqref{eq:c0-22} 演化, $g(t)$ 的曲率变化会越来越小.
当 $t \to \infty$ 时在极限中得到的度量类型将使我们能够确定 $M$ 的拓扑特征.
然而, 在实现这一目标之前, 流可能会产生奇点.
为了克服这一重大技术困难, 需要对 \eqref{eq:c0-22} 的解的行为作详细的定性分析, 这一任务几乎需要几何 PDE 在过去一百年中的所有进展.

正如上面所见, 调和坐标的选取使我们能够把 Ricci 流写成非线性热方程组 \eqref{eq:c0-23}.
这一事实对几何方程相当典型.
此时讨论另一个重要例子——真空中的 Einstein 方程——是有益的.
这个方程的简介及其在广义相对论中重要性的简短讨论, 可以在《纲要》的文章中找到(见纲要文章).
\newterm{einstein-vacuum-equations}{Einstein 真空方程}{Einstein vacuum equations}的解由\newterm{ricci-flat}{Ricci 平坦}{Ricci flat}时空给出, 即\newterm{lorentzian-manifold}{Lorentz 流形}{Lorentzian manifold} $(M, g)$, 其中 $M$ 是四维流形, $g$ 是它上面的\newterm{lorentz-metric}{Lorentz 度量}{Lorentz metric}, 并且相应的 Ricci 曲率恒等于零.

\begin{itemize}
\item (Einstein 真空)
\begin{equation}\label{eq:c0-24}
\operatorname{Ric}(g) = 0.
\end{equation}
\end{itemize}

Lorentz 度量的 Ricci 曲率 $\operatorname{Ric}(g)$ 可以完全按照 Riemann 情形的方式定义.
因此, 相对于坐标系 $x^{\alpha}$, 其中 $\alpha = 0, 1, 2, 3$, Ricci 曲率(记作 $R_{\alpha \beta}$)可以用度量系数 $g_{\alpha \beta}$ 的一阶与二阶偏导数表示.
如前, 我们用 $g^{\alpha \beta}$ 表示逆度量的分量.
此外, 选取一个称为\newterm{wave-coordinates}{波坐标}{wave coordinates}\footnote{它们是上面讨论的调和坐标的精确类比.}的特定坐标系, 我们可以把 Einstein 真空方程 \eqref{eq:c0-24} 写成与波动方程 \eqref{eq:c0-3} 有关的方程组形式, 正如 Ricci 流方程组 \eqref{eq:c0-23} 与热方程 \eqref{eq:c0-2} 有关一样.
更精确地说,

\begin{equation}\label{eq:c0-25}
\Box_{g} g_{\alpha \beta} = N_{\alpha \beta}(g, \partial g)
\end{equation}

其中, 与 Ricci 流的情形一样, 项 $N_{\alpha \beta}(g, \partial g)$ 是可以显式算出的表达式, 它们依赖于度量 $g_{\alpha \beta}$、其逆 $g^{\alpha \beta}$ 以及 $g_{\alpha \beta}$ 关于坐标 $x^{\alpha}$ 的一阶导数.
这是关于度量 $(g_{\alpha \beta})_{\alpha, \beta = 0, 1, 2, 3}$ 的十个未知分量的十个方程的方程组.
出现在左端的微分算子

\[
\Box_{g} = \sum_{\mu, \nu} g^{\mu \nu} \partial_{\mu} \partial_{\nu}
\]

与我们之前在 \eqref{eq:c0-8} 中遇到的波算子 $\Box = m^{\mu \nu} \partial_{\mu} \partial_{\nu} = -\partial_{0}^{2} + \Delta$ 非常相似.
事实上, 在点 $p \in M$ 的一个邻域内, 我们可以选取波坐标 $x^{\alpha}$, 使得 $g^{\mu \nu}(p) = m^{\mu \nu}$.
于是, 局部地, $\Box_{g}$ 看起来就像 $\Box = \Box_{m}$, 从而我们可以把 \eqref{eq:c0-25} 解释为非线性波动方程组.

最后两个例子说明了, 对由几何量(如 Ricci 曲率)定义的方程, 选取好的坐标很重要.
要解这样的方程并寻找解的有趣性质, 选取一个良好适应的坐标系往往非常重要.
在诸如 Yang-Mills 方程之类的\newterm{gauge-field-theory}{规范场论}{gauge field theory}中, 坐标的角色被\newterm{gauge-transformation}{规范变换}{gauge transformation}所取代.

最后我们需要指出, PDE 不仅出现在物理学与几何学中, 也出现在应用科学的许多领域.
例如在工程中, 人们常常希望对 PDE 的解施加辅助条件, 它们对应于物理系统中我们能直接影响的部分, 比如小提琴琴弦与琴弓直接接触的部分, 以便控制它们的行为, 即获得优美的声音.
处理这一问题的数学理论称为\newterm{control-theory}{控制论}{Control Theory}.

在处理复杂物理系统时, 当我们不可能在任何给定时刻掌握系统的完整状态信息时, 人们常常对影响它的各种因素作各种随机性假设.
这就导出一类非常重要的方程, 称为\newterm{stochastic-differential-equation}{随机微分方程}{stochastic differential equation}.
举一个简单例子, 考虑 $N \times N$ 的常微分方程组

\begin{equation}\label{eq:c0-26}
\frac{d x}{d t} = f(x(t))
\end{equation}

这里 $f$ 是给定函数 $f : \mathbb{R}^{N} \to \mathbb{R}^{N}$.
解 $x(t)$ 是向量值函数 $x : [0, \infty) \to \mathbb{R}^{N}$.
给定初值 $x(0) = x_{0}$, 我们可以精确确定解在任意给定时刻 $t$ 的位置 $x(t)$ 与速度 $\frac{d x}{d t}$.
在应用情形中, 由于各种难以计入的因素, 解的状态可能无法如此精确地确定.
因此, 合理地修改方程以计入影响系统的随机效应.
于是人们考虑如下形式的方程

\begin{equation}\label{eq:c0-27}
\frac{d x}{d t} = f(x(t)) + B(x(t)) \frac{d W}{d t}(t)
\end{equation}

其中 $B(x)$ 是 $N \times M$ 维矩阵, $W(t)$ 表示 $\mathbb{R}^{M}$ 中的\newterm{brownian-motion}{Brown 运动}{Brownian motion}.
对偏微分方程也可以作类似的计入随机性的修改.
一个特别有趣的 PDE 例子是著名的\newterm{black-scholes-equation}{Black-Scholes 方程}{Black-Scholes equation}, 它由与金融中股票期权价格有关的随机过程导出.
股票期权的真实价格 $u(s, t)$(在时刻 $t$、价值为 $s$)满足如下 PDE:

\begin{equation}\label{eq:c0-28}
\partial_{t} u + r s \, \partial_{s} u + \frac{\sigma^{2}}{2} s^{2} \, \partial_{s}^{2} u - r u = 0, \qquad s > 0, \quad t \in [0, T],
\end{equation}

并受到期时刻 $T$ 的终值条件 $u = \max(0, (s - p))$ 与边界条件 $u(0, t) = 0$, $t \in [0, T]$ 约束.
这里 $p$ 是期权的执行价.
注意, 这个方程实际上是热方程 \eqref{eq:c0-2} 的一个(时间反演的)变体, 从而说明了上面提出的观点:同一类数学方程可以出现在几个完全不同的应用中(在这里是热力学与数理金融).

% ==== part01 第 1 章:由变分原理导出的方程 ====
% 源: MinerU_latex_lecture nots.tex 868-2576(pdf_text.txt 588-1962,PDF 第 15-38 页)
% 术语: 新术语(中文=英文):
%   时空=space-time; Lorentz 度量=Lorentz metric; 惯性坐标系=inertial coordinate system;
%   Lorentz 变换=Lorentz transformation; 体积元=volume element; 紧变分=compact variation;
%   驻定=stationary; 作用量原理=action principle; Christoffel 符号=Christoffel symbols;
%   水平集=level set; 第二基本形式=second fundamental form; 恰当 2-形式=exact two form;
%   规范自由度=gauge freedom; 电磁场=electromagnetic field; 外微分=exterior differentiation;
%   内部对称=internal symmetry; 李群=Lie group; 李代数=Lie algebra; Lie 括号=Lie bracket;
%   Killing 标量积=Killing scalar product; 正交群=orthogonal group; 特殊正交群=special orthogonal
%   group; Lorentz 群=Lorentz group; 酉群=unitary group; Hermite 双线性型=hermitian bilinear form;
%   Killing 型=Killing form; 规范条件=gauge condition; Coulomb 规范=Coulomb gauge;
%   Lorentz 规范=Lorentz gauge; 时间规范=temporal gauge; 物质场=matter field;
%   引力场=gravitational field; Einstein 张量=Einstein tensor;
%   两次缩并的 Bianchi 恒等式=twice contracted Bianchi identity; 微分同胚=diffeomorphism;
%   正能量条件=positive energy condition; 零锥=null cone; 零对=null pair;
%   零分解=null decomposition; 共形变换=conformal transformation; 变形张量=deformation tensor;
%   后向光锥=backward light cone; 能量密度=energy density; 能量通量密度=energy flux density;
%   有限传播速度=finite propagation speed; 影响域=domain of influence;
%   全局存在性=global existence; 能量色散=energy dispersion; 等距=isometry
% REFS: 跨章引用: 第 1 章第 3 节 "We know that (see appendix 4.2, up to a Lorentz transformation
%   the only Killing, future oriented timelike vectorfield is a constant multiple of ∂_t)"
%   → 目标: 附录 4.2(正文写"见附录 4.2",待统一解析). 
%   公式修正(与原书/OCR 差异,已按数学含义处理,均同时对照 pdf_text.txt):
%   - 源 1627 行 J_i 式补回被 OCR 丢失的 F_{i0}(两源一致缺失,按数学结构还原);
%   - 源 1216 行 Maxwell 变分积分末项 ∂_ν(√−gF^{μν}) → ∂_μ(√−gF^{μν})(与下文"D_μF^{μν}"一致);
%   - 源 2503 行 (77) 第一行 T_{+-} 项补回缺失的系数 1/4(由下划线注 2(t²+r²) 与次行 ½(t²+r²) 确定);
%   - 源 2117 行 T(E_(4),E_(4)) 首行去掉 ^*F_{4A}·^*F_{4A} 前多余的 1/4(否则等式链不成立);
%   - 源 2569 行 (80) 修正为 Q(L̲,L̲) = L̲(φ)²(原书误印 L(φ)²);
%   - 源 2497 行 K_0 表达式去掉结尾多余右括号与句点(任务指定);
%   - 源 1867 行 S_M[ψ,g;M] → S[ψ,g;M](微分同胚不变性应为 S[ψ,g;M]);
%   - 源 2456 行 Ė_c(0) → E_c(0)(点号为 OCR 噪音);
%   - 源 1472 行 su(p,q) 定义: A∈U(p,q) → A∈u(p,q),tr_M → tr_Q(由下一句 tr_Q A = Q^{ij}A_{ij} 确定);
%   - 源 1430 行 "dim δR O" → dim_R O; 源 2493、2557 行 ∂_η → ∂_r; 源 2528/2542 行 \boxed{d} → \Box;
%   书内疑似笔误(照排未改,供校对参考): O(p,q) 第二行 "LML^T = M"(M 未定义);
%   Lemma 3.7 中 "Ω = 4(n+1)t"(数学上应为 Ω = 4t,或应写 trπ = 4(n+1)t);
%   (56) 右端 "−ΔA_0 + ∂^i[A_0,A_i]" 为多余重复项; J_0 展开式(源 1615 行)同此. 
% 编号说明: 公式 (29)-(80) 顺次自动编号(共 52 个;(43)(44) 为两个独立公式,源文件中的空
%   enumerate 已处理为 (44));第 3 节 "3.5""3.6" 为 \mtsubsec 子标题,自然推进共享计数器,
%   故 Lemma 3.7 编号为 3.7;节内 5 个独立子标题(1. 标量场方程 ~ 5. Einstein 场方程)用
%   \mtsubsecn 不推进计数器;两处未编号 Definition、Remark、Example 用带 * 环境;
%   脚注 1-11 全部按 pdf_text.txt 对应页补全并翻译. 
% ====

\chapter{由变分原理导出的方程}

\section{基本概念}

本节我们将讨论一些线性与非线性方程的基本例子, 它们由相对论 Lagrange 量变分地导出. 相对论场论的基本对象是：

\begin{itemize}
\item \newterm{spacetime}{时空}{space-time} $(\mathbf{M}, \mathbf{g})$, 由 $n+1$ 维流形 $M$ 与 Lorentz 度量 $g$ 构成, 即定义在 $M$ 每点切空间上的、符号为 $(-1, 1, \ldots, 1)$ 的非退化二次型. 我们用 $x^\alpha, \ \alpha = 0, 1, \ldots, n$ 记 $M$ 中点的坐标. 

在本章的大部分内容中, 时空实际上将是最简单的可能例子——即 Minkowski 时空, 其中流形为 $\mathbb{R}^{n+1}$, 度量由下式给出：
\begin{equation}
ds^2 = \mathbf{m}_{\alpha\beta}\, dx^\alpha dx^\beta = -dt^2 + (dx^1)^2 + \dots + (dx^n)^2\label{eq:c1-29}
\end{equation}
其中 $t = x^0$, $m_{\alpha\beta} = \mathrm{diag}(-1, 1, \ldots, 1)$. 回忆：任何使度量具有形式 \eqref{eq:c1-29} 的坐标系称为 \newterm{inertial-frame}{惯性坐标系}{inertial coordinate system}. 任意两个惯性坐标系由 Lorentz 变换 相联系. 

\item 场的集合 $\boldsymbol{\psi} = \boldsymbol{\psi}^{(1)}, \boldsymbol{\psi}^{(2)}, \dots, \boldsymbol{\psi}^{(p)}$, 它们可以是标量、张量, 或其他一些定义在 $M$ 上的几何对象\footnote{为简单起见, 我们仅限于协变张量. }, 如旋量. 

\item \newterm{lagrangian-density}{Lagrange 密度}{Lagrangian density} L, 它是 $M$ 上的标量函数, 只依赖于张量场 $\psi$ 与度规\footnote{以及其逆 $g^{-1}$. } g. 
\end{itemize}

我们随后将相应的作用量 $S$ 定义为积分：
\[
\mathcal{S} = \mathcal{S}[\psi, \mathbf{g}:\mathcal{U}] = \int_{\mathcal{U}}L[\psi]\, dv_{\mathbf{g}}
\]
其中 $\mathcal{U}$ 是 $M$ 的任意相对紧集. 这里 $dv_{\mathbf{g}}$ 表示由度规 $\mathbf{g}$ 生成的 \newterm{volume-element}{体积元}{volume element}. 更精确地, 相对于局部坐标系 $x^\alpha$, 有
\[
dv_{\mathbf{g}} = \sqrt{-\mathbf{g}}\, dx^0dx^1\cdots dx^n = \sqrt{-\mathbf{g}}\, dx
\]
其中 $g$ 是矩阵 $(\mathbf{g}_{\alpha\beta})$ 的行列式. 

所谓场 $\psi$ 的 \newterm{compact-variation}{紧变分}{compact variation}, 是指一族定义于 $s \in (-\epsilon, \epsilon)$ 的光滑单参数场 $\psi_{(s)}$, 使得：
\begin{enumerate}
\item 在 $s = 0$ 处, $\psi_{(0)} = \psi$. 
\item 在 $\mathbf{M} \setminus \mathcal{U}$ 的所有点 $p$ 处, $\psi_{(s)} = \psi$. 
\end{enumerate}

对这样的变分, 记 $\delta\psi := \dot{\psi} := \frac{d\psi_{(s)}}{ds}\Big|_{s=0}$. 于是, 对小的 $s$, 
\[
\psi_{(s)} = \psi + s\dot{\psi} + O(s^2)
\]
若对 $\psi$ 的任意紧变分 $(\psi_{(s)}, \mathcal{U})$, 均有
\[
\left.\frac{d}{ds}\mathbf{S}(s)\right|_{s=0} = 0
\]
其中
\[
\mathbf{S}(s) = \mathbf{S}[\psi_{(s)}, \mathbf{g};\mathcal{U}]
\]
则称场 $\psi$ 关于 $S$ 是 \newterm{stationary}{驻定}{stationary} 的. 我们用简记记号将其写成
\[
\frac{\delta\mathbf{S}}{\delta\psi} = 0
\]

\newterm{action-principle}{作用量原理}{action principle}（亦称变分原理）断言：物理系统的一个可接受解必须关于给定的 Lagrange 密度（称为该系统的 Lagrange 量）是驻定的. 作用量原理使我们能够导出场 $\psi$ 所满足的偏微分方程, 称为 Euler-Lagrange 方程. 下面是一些简单的例子：

\mtsubsecn{1}{标量场方程}

从 Lagrange 密度出发：
\[
L[\phi] = -\frac{1}{2}\mathbf{g}^{\mu\nu}\partial_\mu\phi\, \partial_\nu\phi - V(\phi)
\]
其中 $\phi$ 是定义在 $(\mathbf{M}, \mathbf{g})$ 上的复值标量函数, $V(\phi)$ 是 $\phi$ 的给定实值函数. 

给定 $\phi$ 的紧变分 $(\phi_{(s)}, \mathcal{U})$, 我们设 $S(s) = S[\phi_{(s)}, \mathbf{g};\mathcal{U}]$. 分部积分给出：
\[
\begin{aligned}
\left.\frac{d}{ds}\mathcal{S}(s)\right|_{s=0} &= \int_{\mathcal{U}}\left[-\mathbf{g}^{\mu\nu}\partial_\mu\dot{\phi}\, \partial_\nu\phi - V'(\phi)\dot{\phi}\right]\sqrt{-\mathbf{g}}\, dx\\
&= \int_{\mathcal{U}}\dot{\phi}\left[\Box_{\mathbf{g}}\phi - V'(\phi)\right]dv_{\mathbf{g}}
\end{aligned}
\]
其中 $\Box_{\mathbf{g}}$ 是 D'Alembert 算子（波算子）：
\[
\Box_{\mathbf{g}}\phi = \frac{1}{\sqrt{-\mathbf{g}}}\partial_\mu\left(\mathbf{g}^{\mu\nu}\sqrt{-\mathbf{g}}\, \partial_\nu\phi\right)
\]

根据作用量原理以及 $\dot{\phi}$ 的任意性, 我们推断 $\phi$ 必满足如下的 Euler-Lagrange 方程：
\begin{equation}
\Box_{\mathbf{g}}\phi - V'(\phi) = 0\label{eq:c1-30}
\end{equation}

方程 \eqref{eq:c1-30} 称为带势 $V(\phi)$ 的标量波动方程. 

\mtsubsecn{2}{波映射}

波映射方程将在如下的框架中定义：时空 $(\mathbf{M}, \mathbf{g})$、带度量 $h$ 的 Riemann 流形 N, 以及映射
\[
\phi:\mathbf{M} \longrightarrow N.
\]
我们回忆：若 $X$ 是 $M$ 上的向量场, 则 $\phi_*X$ 是 $N$ 上由 $\phi_*X(f) = X(f\circ\phi)$ 定义的向量场. 若 $\omega$ 是 $N$ 上的 1-形式, 则其回拉 $\phi^*\omega$ 是 $M$ 上由 $\phi^*\omega(X) = \omega(\phi_*X)$ 定义的 1-形式, 其中 $X$ 是 $M$ 上的任意向量场. 类似地, 度量 $h$ 的回拉是 $M$ 上由公式 $(\phi^*h)(X, Y) = h(\phi_*X, \phi_*Y)$ 定义的对称 2-协变张量. 在 $M$ 上的局部坐标 $x^\alpha$ 与 $N$ 上的局部坐标 $y^a$ 中, 若 $\phi^a$ 表示 $\phi$ 关于 $y^a$ 的分量, 则有
\[
(\phi^*h)_{\alpha\beta}(p) = \frac{\partial\phi^a}{\partial x^\alpha}\frac{\partial\phi^b}{\partial x^\beta}h_{ab}(\phi(p)) = \left\langle\frac{\partial\phi}{\partial x^\alpha}, \frac{\partial\phi}{\partial x^\beta}\right\rangle
\]
其中 $\langle\cdot, \cdot\rangle$ 表示 $N$ 上的 Riemann 内积. 

考虑如下涉及映射 $\phi$ 的 Lagrange 密度：
\[
L = -\frac{1}{2}\mathrm{Tr}_{\mathbf{g}}(\phi^*h)
\]
其中 $\mathrm{Tr}_{\mathbf{g}}(\phi^*h)$ 表示 $\phi^*h$ 关于 $g$ 的迹. 在局部坐标中, 
\[
L[\phi] = -\frac{1}{2}\mathbf{g}^{\mu\nu}h_{ab}(\phi)\frac{\partial\phi^a}{\partial x^\mu}\frac{\partial\phi^b}{\partial x^\nu}.
\]

按定义, 波映射是相应作用量的驻点. 于是通过一个直接的计算：
\begin{equation}
\begin{aligned}
0 &= \left.\frac{d}{ds}\mathcal{S}(s)\right|_{s=0} = I_1 + I_2\\
I_1 &= -\frac{1}{2}\int_{\mathcal{U}}\mathbf{g}^{\mu\nu}\frac{\partial h_{ab}(\phi)}{\partial\phi^c}\dot{\phi}^c\, \partial_\mu\phi^a\, \partial_\nu\phi^b\sqrt{-\mathbf{g}}\, dx\\
I_2 &= -\int_{\mathcal{U}}\mathbf{g}^{\mu\nu}h_{ab}(\phi)\partial_\mu\dot{\phi}^a\, \partial_\nu\phi^b\sqrt{-\mathbf{g}}\, dx
\end{aligned}\label{eq:c1-31}
\end{equation}

经分部积分、重命名指标并利用关于 $b, c$ 的对称性, 我们可将 $I_2$ 改写为
\begin{equation}
\begin{aligned}
I_2 &= \int_{\mathcal{U}}\dot{\phi}^a\left(h_{ab}(\phi)\Box_{\mathbf{g}}\phi^b + \mathbf{g}^{\mu\nu}\frac{\partial h_{ab}}{\partial\phi^c}\partial_\mu\phi^c\, \partial_\nu\phi^b\right)dv_{\mathbf{g}}\\
&= \int_{\mathcal{U}}\dot{\phi}^a\left(h_{ab}(\phi)\Box_{\mathbf{g}}\phi^b + \frac{1}{2}\mathbf{g}^{\mu\nu}\left(\frac{\partial h_{ab}}{\partial\phi^c} + \frac{\partial h_{ac}}{\partial\phi^b}\right)\partial_\mu\phi^b\, \partial_\nu\phi^c\right)dv_{\mathbf{g}}
\end{aligned}\label{eq:c1-32}
\end{equation}
另外, 重命名指标：
\[
I_1 = -\frac{1}{2}\int_{\mathcal{U}}\mathbf{g}^{\mu\nu}\frac{\partial h_{bc}}{\partial\phi^a}\dot{\phi}^a\, \partial_\mu\phi^b\, \partial_\nu\phi^c\, dv_{\mathbf{g}}.
\]
因此, 
\[
\begin{aligned}
0 &= I_1 + I_2\\
&= \int_{\mathcal{U}}\dot{\phi}^a\left(h_{ab}\Box_{\mathbf{g}}\phi^b + \partial_\mu\phi^b\, \partial_\nu\phi^c\, \mathbf{g}^{\mu\nu}\frac{1}{2}\left(\frac{\partial h_{ab}}{\partial\phi^c} + \frac{\partial h_{ac}}{\partial\phi^b} - \frac{\partial h_{bc}}{\partial\phi^a}\right)\right)dv_{\mathbf{g}}\\
&= \int_{\mathcal{U}}\dot{\phi}^a\left(h_{ad}\Box_{\mathbf{g}}\phi^d + \partial_\mu\phi^b\, \partial_\nu\phi^c\, \mathbf{g}^{\mu\nu}\frac{1}{2}h^{ds}h_{ad}\cdot\left(\frac{\partial h_{sb}}{\partial\phi^c} + \frac{\partial h_{sc}}{\partial\phi^b} - \frac{\partial h_{bc}}{\partial\phi^s}\right)\right)dv_{\mathbf{g}}\\
&= \int_{\mathcal{U}}\dot{\phi}^a h_{ad}\left(\Box_{\mathbf{g}}\phi^d + \partial_\mu\phi^b\, \partial_\nu\phi^c\, \mathbf{g}^{\mu\nu}\Gamma^d_{bc}\right)dv_{\mathbf{g}}
\end{aligned}
\]
其中 $\Gamma^d_{bc} = \frac{1}{2}h^{ds}\left(\frac{\partial h_{sb}}{\partial\phi^c} + \frac{\partial h_{sc}}{\partial\phi^b} - \frac{\partial h_{bc}}{\partial\phi^s}\right)$ 是对应于 Riemann 度量 $h$ 的 \newterm{christoffel-symbols}{Christoffel 符号}{Christoffel symbols}. 由 $\dot{\phi}$ 的任意性, 得到波映射的如下方程：
\begin{equation}
\Box_{\mathbf{g}}\phi^a + \Gamma^a_{bc}\, \mathbf{g}^{\mu\nu}\partial_\mu\phi^b\, \partial_\nu\phi^c = 0\label{eq:c1-33}
\end{equation}

\begin{example*}
设 $N$ 是二维 Riemann 流形, 配有如下形式的度量 $h$：
\[
ds^2 = dr^2 + f(r)^2d\theta^2
\]
设 $\phi$ 是从 $M$ 到 $N$ 的波映射, 关于 $r, \theta$ 坐标的分量为 $\phi^1, \phi^2$. 则 $\Gamma^1_{11} = \Gamma^2_{11} = \Gamma^1_{12} = \Gamma^2_{22} = 0$, 且 $\Gamma^1_{22} = -f'(r)f(r)$, $\Gamma^2_{12} = \frac{f'(r)}{f(r)}$. 因此, 
\[
\begin{aligned}
\Box_{\mathbf{g}}\phi^1 &= f'(r)f(r)\, \mathbf{g}^{\mu\nu}\partial_\mu\phi^2\, \partial_\nu\phi^2\\
\Box_{\mathbf{g}}\phi^2 &= -\frac{f'(r)}{f(r)}\, \mathbf{g}^{\mu\nu}\partial_\mu\phi^1\, \partial_\nu\phi^2
\end{aligned}
\]
\end{example*}

当 $N$ 是 Euclidean 空间 $\mathbb{R}^m$ 的子流形时, 波映射方程可以有更简单的表述. 此时, 度量 $h$ 是 Euclidean 度量\footnote{使用外围 Euclidean 空间的标准坐标. }, 故 \eqref{eq:c1-31} 中的第一项消失：
\[
\begin{aligned}
\left.\frac{d}{ds}\mathcal{S}(s)\right|_{s=0} &= -\int_{\mathcal{U}}\mathbf{g}^{\alpha\beta}\left\langle\frac{\partial\phi}{\partial x^\alpha}, \frac{\partial\dot{\phi}}{\partial x^\beta}\right\rangle dv_{\mathbf{g}}\\
&= \int_{\mathcal{U}}\langle\Box\phi, \dot{\phi}\rangle dv_{\mathbf{g}}
\end{aligned}
\]
其中 $\langle\cdot, \cdot\rangle$ 是 Euclidean 内积, $\Box$ 是 $M$ 上的 D'Alembert 算子. 于是 Euler-Lagrange 方程取如下形式：
\begin{equation}
(\Box\phi(p))^T = 0\label{eq:c1-34}
\end{equation}
其中 $T$ 表示到 $N$ 在 $\phi(p)$ 处切空间上的投影. 

在 $N\subset \mathbb{R}^m$ 是超曲面的特殊情形, 我们可将 \eqref{eq:c1-34} 改写为更具体的形式. 设 $\nu$ 是 $N$ 上的单位法向量, k 是 \newterm{second-fundamental-form}{第二基本形式}{second fundamental form} $k(X, Y) = \langle\mathbf{D}_X\boldsymbol{\nu}, Y\rangle$, 其中 $\mathbf{D}_X$ 是 Euclidean 空间的标准协变导数. 超曲面 $N$（局部地）定义为某个实值函数 $f$ 的 \newterm{level-set}{水平集}{level set}. 对方程 $f(\phi(x)) = 0$ 关于 $M$ 上的局部坐标 $x^\mu$ 求导, 沿 $M$ 得 $0 = \langle\nu(\phi), \partial_\mu\phi\rangle$. 于是, 
\[
\begin{aligned}
0 &= \partial^\mu\langle\nu(\phi), \partial_\mu\phi\rangle = \langle\Box\phi, \nu\rangle + \mathbf{g}^{\mu\nu}\langle\partial_\nu\nu(\phi), \partial_\mu\phi\rangle\\
&= \langle\Box\phi, \nu\rangle + \mathbf{g}^{\mu\nu}\langle\nabla_{\phi_*(E_\nu)}\nu, \phi_*(E_\mu)\rangle
\end{aligned}
\]
其中 $\phi_*(E_\mu) = \frac{\partial\phi^i}{\partial x^\mu}\frac{\partial}{\partial y^i}$ 是 $E_\mu = \frac{\partial}{\partial x^\mu}$ 的 \newterm{pushforward}{前推}{pushforward}. 特别地, $\phi_*(E_\mu)$ 与 $N$ 相切. 因此, 
\begin{equation}
\langle\Box\phi, \nu\rangle = -k(\phi_*(E^\alpha), \phi_*(E_\alpha))\label{eq:c1-35}
\end{equation}
于是波映射方程变为：
\[
\Box\phi = -k(\phi_*(E^\alpha), \phi_*(E_\alpha))\, \nu
\]
当 $N$ 是标准球面 $\mathbb{S}^{m-1} \subset \mathbb{R}^m$ 时, $k(X, Y) = -\langle X, Y\rangle$；在 $\mathbb{R}^m$ 的标准坐标 $x^\alpha$ 下, $\nu^a(\phi) = \phi^a$, 波映射方程变为：
\[
\Box\phi^a = -\phi^a\, \mathbf{g}^{\alpha\beta}\left\langle\frac{\partial\phi}{\partial x^\alpha}, \frac{\partial\phi}{\partial x^\beta}\right\rangle
\]

\mtsubsecn{3}{Maxwell 方程组}

电磁场 $F$ 是四维流形 $M$ 上的 \newterm{exact-two-form}{恰当 2-形式}{exact two form}. 即 $F$ 是秩为二的反对称张量, 使得
\begin{equation}
F = dA\label{eq:c1-36}
\end{equation}
其中 $A$ 是 $M$ 上的 1-形式, 称为规范势或联络 1-形式. 注意 $A$ 不是唯一确定的——事实上, 若 $\chi$ 是任意标量函数, 则变换
\begin{equation}
A \longrightarrow \tilde{A} = A + d\chi\label{eq:c1-37}
\end{equation}
给出 $F$ 的另一个规范势 $\tilde{A}$. 这种任意性程度称为 \newterm{gauge-freedom}{规范自由度}{gauge freedom}, 变换 \eqref{eq:c1-37} 称为规范变换. 

\newterm{electromagnetic-field}{电磁场}{electromagnetic field}的 Lagrange 密度为
\[
L[F] = -\frac{1}{4}F_{\mu\nu}F^{\mu\nu}.
\]

F 的任意紧变分 $(F_{(s)}, \mathcal{U})$ 都可以用规范势 $A$ 的紧变分 $(A_{(s)}, \mathcal{U})$ 表示, 使得 $F_{(s)} = dA_{(s)}$. 记
\[
\dot{F} = \left.\frac{d}{ds}F_{(s)}\right|_{s=0}, \qquad \dot{A} = \left.\frac{d}{ds}A_{(s)}\right|_{s=0}
\]
则相对于坐标系 $x^\alpha$, 有 $F_{\mu\nu} = \partial_\mu A_\nu - \partial_\nu A_\mu$, 从而 $\dot{F}_{\mu\nu} = \partial_\mu\dot{A}_\nu - \partial_\nu\dot{A}_\mu$. 作用量原理给出：
\[
\begin{aligned}
0 &= \left.\frac{d}{ds}\mathcal{S}(s)\right|_{s=0} = -\frac{1}{2}\int_{\mathbf{M}}\dot{F}_{\mu\nu}F^{\mu\nu}dv_{\mathbf{g}}\\
&= -\frac{1}{2}\int_{\mathcal{U}}(\partial_\mu\dot{A}_\nu - \partial_\nu\dot{A}_\mu)F^{\mu\nu}dv_{\mathbf{g}}\\
&= -\int_{\mathcal{U}}\partial_\mu\dot{A}_\nu\, F^{\mu\nu}dv_{\mathbf{g}} = \int_{\mathcal{U}}\dot{A}_\nu\left(\frac{1}{\sqrt{-\mathbf{g}}}\partial_\mu\left(\sqrt{-\mathbf{g}}F^{\mu\nu}\right)\right)dv_{\mathbf{g}}
\end{aligned}
\]
注意被积函数中的第二个因子正是 $\mathbf{D}_\mu F^{\mu\nu}$, 其中 $D$ 是 $M$ 上对应于 $g$ 的协变导数. 因此 Euler-Lagrange 方程取如下形式：
\begin{equation}
\mathbf{D}_\nu F^{\mu\nu} = 0.\label{eq:c1-38}
\end{equation}
\eqref{eq:c1-36} 与 \eqref{eq:c1-38} 合在一起构成 Maxwell 方程组. 

\textbf{习题.} 给定 $M$ 上的向量场 $X^\alpha$, 证明
\[
\mathbf{D}_\alpha X^\alpha = \frac{1}{\sqrt{-\mathbf{g}}}\partial_\alpha\left(\sqrt{-\mathbf{g}}X^\alpha\right)
\]

利用 $F$ 的 Hodge 对偶, 我们可以把 Maxwell 方程组写成更对称的形式：
\[
{}^\star F_{\mu\nu} = \frac{1}{2}\epsilon_{\mu\nu\alpha\beta}F^{\alpha\beta}
\]
并注意到 \eqref{eq:c1-38} 等价于 $d\, {}^\star F = 0$. Maxwell 方程组于是取如下形式：
\begin{equation}
dF = 0, \qquad d\, {}^\star F = 0\label{eq:c1-39}
\end{equation}
或者等价地, 
\begin{equation}
\mathbf{D}_\nu F^{\mu\nu} = 0, \qquad \mathbf{D}_\nu\, {}^\star F^{\mu\nu} = 0\label{eq:c1-40}
\end{equation}
注意, 由于 Lorentz 变换与 Hodge 对偶和\newterm{exterior-differentiation}{外微分}{exterior differentiation}都可交换, Maxwell 方程组的 Lorentz 不变性在 \eqref{eq:c1-39} 中是显式的. 

\begin{definition*}
给定任意向量场 X, 我们可以定义缩并
\[
\begin{aligned}
E_\alpha = (i_XF)_\alpha &= X^\mu F_{\alpha\mu}\\
H_\alpha = (i_X\, {}^\star F)_\alpha &= X^\mu\, {}^\star F_{\alpha\mu}
\end{aligned}
\]
分别称为 $F$ 的电场分量与磁场分量. 注意这两个 1-形式都垂直于 X. 
\end{definition*}

我们特化到 $M$ 是 Minkowski 空间且 $X = \frac{d}{dx^0} = \frac{d}{dt}$ 的情形. 如前所述, E, H 垂直于 $\frac{d}{dt}$, 故 $E_0 = H_0 = 0$. 空间分量按定义为
\[
\begin{aligned}
E_i &= F_{0i}\\
H_i &= {}^\star F_{0i} = \frac{1}{2}\epsilon_{0ijk}F^{jk} = \frac{1}{2}\epsilon_{ijk}F^{jk}
\end{aligned}
\]

我们现在利用 \eqref{eq:c1-39} 由上推导 $E$ 与 $H$ 的方程, 它们分别蕴含
\begin{equation}
\mathbf{D}_\nu\, {}^\star F^{\mu\nu} = 0\label{eq:c1-41}
\end{equation}
与 \eqref{eq:c1-38}. 在 \eqref{eq:c1-40} 的两个方程中令 $\mu = 0$, 我们得到：
\begin{equation}
\partial^iE_i = 0, \quad \partial^iH_i = 0\label{eq:c1-42}
\end{equation}
令 $\mu = i$ 并注意到 $F_{ij} = \epsilon_{ijk}H^k$, ${}^\star F_{ij} = -\epsilon_{ijk}E^k$, 我们写出：
\[
\begin{aligned}
0 &= -\partial^0E_i + \partial^jF_{ij} = \partial_0E_i + \epsilon_{ijk}\partial^jH^k = \partial_tE_i + (\nabla\times H)_i\\
0 &= \partial_tH_i - \epsilon_{ijk}\partial_jE_k = \partial_tH_i - (\nabla\times E)_i
\end{aligned}
\]
因此, 
\begin{equation}
\partial_tE + \nabla\times H = 0\label{eq:c1-43}
\end{equation}
\begin{equation}
\partial_tH - \nabla\times E = 0\label{eq:c1-44}
\end{equation}

与 \eqref{eq:c1-43} 和 \eqref{eq:c1-44} 相伴, 我们可以在时刻 $t = 0$ 指定数据：
\[
E_i(0, x) = E_i^{(0)}, \qquad H_i(0, x) = H_i^{(0)}
\]

\textbf{习题.} 证明方程组 \eqref{eq:c1-42} 在系统 \eqref{eq:c1-43}–\eqref{eq:c1-44} 的时间演化下保持不变. 换言之, 若 $E^{(0)}, H^{(0)}$ 满足 \eqref{eq:c1-42}, 则 E, H 对所有时刻 $t \in \mathbb{R}$ 都满足 \eqref{eq:c1-42}. 

\mtsubsecn{4}{Yang-Mills 方程组}

所有经典场论的 Lagrange 量都表现出时空的对称性. 除了这些时空对称性, Lagrange 量还可能具有称为场的\newterm{internal-symmetry}{内部对称}{internal symmetry}的对称性. 一个简单的例子是复标量 Lagrange 量：
\[
\mathcal{L} = -\frac{1}{2}\mathbf{m}^{\alpha\beta}\partial_\alpha\phi\, \overline{\partial_\beta\phi} - V(|\phi|)
\]
其中 $\phi$ 是定义在 Minkowski 时空 $\mathbb{R}^{n+1}$ 上的复值标量, $\bar{\phi}$ 是其复共轭. 我们注意到 $L$ 在变换 $\phi \to e^{i\theta}\phi$ 下不变, 其中 $\theta$ 是固定的实数. 自然会问：Lagrange 量能否修改, 以允许形如 $\phi(x) \to e^{i\theta(x)}\phi(x)$ 的更一般的局部位相变换. 容易看到, 在此类变换下, 由于项 $\mathbf{m}^{\alpha\beta}\partial_\alpha\phi\, \overline{\partial_\beta\phi}$ 的存在, Lagrange 量不再不变. 为得到不变的 Lagrange 量, 需将导数 $\partial_\alpha\phi$ 替换为依赖于规范势 $A_\alpha$ 的协变导数 $\mathcal{D}_\alpha^{(A)}\phi \equiv \phi_{, \alpha} + iA_\alpha\phi$. 我们现在可以容易地验证, 新的 Lagrange 量
\[
\mathcal{L} = -\frac{1}{2}\mathbf{m}^{\alpha\beta}D_\alpha^{(A)}\phi\, \bar{D}_\beta^{(A)}\phi - V(|\phi|)
\]
相对于局部变换
\[
\phi(x^\alpha) \to e^{i\theta(x)}\phi(x^\alpha), \qquad A_\alpha \to A_\alpha - \theta_{, \alpha}
\]
不变, 后者称为规范变换. 

注意, 上面引入的规范变换与电磁场 $F$ 的定义配合得很好. 事实上, 令 $F = dA$, 我们注意到 $F$ 是不变的. 这使我们能够考虑一个包含 $F$ 的更一般的 Lagrange 量：
\[
\mathcal{L} = -\frac{1}{4}F_{\alpha\beta}F^{\alpha\beta} - \frac{1}{2}m^{\alpha\beta}\phi_{, \alpha}\bar{\phi}_{, \beta} - V(|\phi|)
\]
称为 Maxwell-Klein-Gordon Lagrange 量. 

Yang-Mills Lagrange 量是 Maxwell-Klein-Gordon Lagrange 量在如下情形的自然推广：与复标量 $\phi$ 的位相变换对应的群 $SU(1)$ 被替换为更一般的 李群 G. 此时, 规范势或联络 1-形式的作用由定义在 $M$ 上的 $\mathcal{G}$ 值 1-形式 $A = A_\mu dx^\mu$ 承担. 这里 $\mathcal{G}$ 是 李群 $G$ 的 \newterm{lie-algebra}{李代数}{Lie algebra}. 设 $[\cdot, \cdot]$ 为其 \newterm{lie-bracket}{Lie 括号}{Lie bracket}, $\langle\cdot, \cdot\rangle$ 为其 \newterm{killing-scalar-product}{Killing 标量积}{Killing scalar product}. 通常, 李群 $G$ 是经典矩阵群之一, 即 $\mathrm{Mat}(n, \mathbb{R})$ 或 $\mathrm{Mat}(n, \mathbb{C})$ 的子群. 我们稍作停顿, 回忆一些关于相关 李群及其 李代数的知识. 

\begin{enumerate}
\item \newterm{orthogonal-group}{正交群}{orthogonal group} $\mathbf{O}(p, q)$. 它们是 $\mathbb{R}^n$ 的线性变换群, 保持一个给定的、符号为 $p, q$（$p+q = n$）的非退化对称双线性型. 我们用 $\mathbf{R}^n_{p, q}$ 记相应的空间. $p = 0$ 的情形即 Euclidean 情形, 此时群简记为 $\mathbf{O}(n)$. $p = 1, q = n$ 的情形即 Minkowski 时空 $\mathbb{R}^{n+1}$, 群 $\mathbf{O}(1, n)$ 是 \newterm{lorentz-group}{Lorentz 群}{Lorentz group}. 一般地, 设 $Q$ 是对角矩阵, 其前 $p$ 个对角元为 $-1$, 其余为 $+1$. 则
\end{enumerate}
\[
\begin{aligned}
\mathbf{O}(p, q) &= \{L \in \mathrm{Mat}(n, \mathbb{R})\, |\, L^TQL = Q\}\\
&= \{L \in \mathrm{Mat}(n, \mathbb{R})\, |\, LML^T = M\}
\end{aligned}
\]
注意对 $L \in \mathbf{O}(p, q)$, $\det(L) = \pm 1$. 

回忆 \newterm{special-orthogonal-group}{特殊正交群}{special orthogonal group} $\mathbf{SO}(p, q)$ 定义为
\[
\mathbf{SO}(p, q) = \{L \in \mathbf{O}(p, q)\, |\, \det L = 1\}.
\]
它们对应于 $\mathbf{R}^n_{p, q}$ 的所有保定向等距. $\mathbf{O}(p, q)$ 与 $\mathbf{SO}(p, q)$ 都以\footnote{回忆 李群 $G$ 的 李代数就是 $G$ 在原点的切空间. }
\[
\mathbf{so}(p, q) = \{A \in \mathrm{Mat}(n, \mathbb{R})\, |\, AQ + QA^T = 0\}
\]
为 李代数, 且 $\dim_{\mathbf{R}}\mathbf{O}(p, q) = \dim_{\mathbf{R}}\mathbf{SO}(p, q) = n(n-1)/2$. so$(p, q)$ 的 Lie 括号是矩阵的通常 Lie 括号, 即 $[A, B] = AB - BA$, 并且我们有 Jacobi 恒等式：
\begin{equation}
[A, [B, C]] + [C, [A, B]] + [B, [C, A]] = 0\label{eq:c1-45}
\end{equation}
且其 Killing 标量积 $\langle A, B\rangle = -\mathrm{Tr}(AB^T)$（其中 Tr 是矩阵的通常迹）满足相容性条件：
\begin{equation}
\langle A, [B, C]\rangle = -\langle[A, B], C\rangle\label{eq:c1-46}
\end{equation}

\begin{enumerate}
\setcounter{enumi}{1}
\item \newterm{unitary-group}{酉群}{unitary group} $\mathbf{U}(p, q)$. 它们是正交群的复类比. 它们是 $\mathbb{C}^n$ 的线性变换群, 保持一个给定的非退化 \newterm{hermitian-bilinear-form}{Hermite 双线性型}{hermitian bilinear form}. 用 $\mathbb{C}^n_{p, q}$ 记相应的空间. 于是, 用上面的矩阵 $Q$：
\end{enumerate}
\[
\mathbf{U}(p, q) = \{U \in \mathrm{Mat}(n, \mathbb{C})\, |\, U^*QU = Q\}
\]
且
\[
\mathbf{SU}(p, q) = \{U \in \mathbf{U}(p, q)\, |\, \det U = 1\},
\]
相应的 李代数为：
\[
\begin{aligned}
\mathbf{u}(p, q) &= \{A \in \mathrm{Mat}(n, \mathbb{C})\, |\, AQ + QA^* = 0\}, \\
\mathbf{su}(p, q) &= \{A \in \mathbf{u}(p, q)\, |\, \mathrm{tr}_Q A = 0\},
\end{aligned}
\]
其中迹 $\mathrm{tr}_Q A = Q^{ij}A_{ij}$. Lie 括号仍是矩阵的通常括号. Killing 标量积由 $\langle A, B\rangle = -\mathrm{Tr}(AB^*)$ 给出. 另请注意 $\dim_{\mathbf{R}}\mathbf{U}(p, q) = n^2$, $\dim_{\mathbf{R}}\mathbf{SU}(p, q) = n^2 - 1$. 

在 Yang-Mills 理论中, 人们关心具有正定 \newterm{killing-form}{Killing 型}{Killing form}的紧 李群. $O(n), SO(n), U(n), SU(n)$ 即属此类. 

在给定的坐标系中, 联络 1-形式 $A$ 具有形式 $A_\mu dx^\mu$, 我们定义 $\mathcal{G}$ 值张量 $\psi$ 的（规范）协变导数为：
\begin{equation}
\mathbf{D}_\mu^{(A)}\psi = \mathbf{D}_\mu\psi + [A_\mu, \psi]\label{eq:c1-47}
\end{equation}
其中 $D$ 是 $M$ 上的协变导数. 注意 \eqref{eq:c1-47} 在如下的规范变换下不变：对给定的 $\mathcal{G}$ 值规范势 $A$ 与 $\mathcal{G}$ 值张量 $\psi$, 
\begin{equation}
\tilde{\psi} = U^{-1}\psi U, \qquad \tilde{A}_\alpha = U^{-1}A_\alpha U + (\mathbf{D}_\alpha U^{-1})U\label{eq:c1-48}
\end{equation}
其中 $U \in G$. 

\begin{proposition}\label{prop:c1-1.1}
\[
\begin{aligned}
\mathbf{D}_\mu^{(\tilde{A})}\tilde{\psi} &= U^{-1}\left(\mathbf{D}_\mu^{(A)}\psi\right)U\\
&= \widetilde{\mathbf{D}^A\psi}
\end{aligned}
\]
\end{proposition}

\begin{proof}
这只需要一点耐心. 首先我们将证明
\[
\mathbf{D}_\alpha\left(U^{-1}\psi U\right) = U^{-1}\left(\mathbf{D}_\alpha\psi + [\psi, U(\mathbf{D}_\alpha U^{-1})]\right)U
\]
事实上
\[
\begin{aligned}
\mathbf{D}_\alpha(U^{-1}\psi U) &= (\mathbf{D}_\alpha U^{-1})\psi U + U^{-1}(\mathbf{D}_\alpha\psi)U + U^{-1}\psi(\mathbf{D}_\alpha U)\\
&= U^{-1}\left(-(\mathbf{D}_\alpha U)U^{-1}\psi + \mathbf{D}_\alpha\psi + \psi(\mathbf{D}_\alpha U)U^{-1}\right)U\\
&= U^{-1}\left(\mathbf{D}_\alpha\psi + [\psi, (\mathbf{D}_\alpha U)U^{-1}]\right)U
\end{aligned}
\]
如所期望. 因此
\[
\begin{aligned}
\mathbf{D}_\alpha^{(\tilde{A})}\tilde{\psi} &= \mathbf{D}_\alpha\tilde{\psi} + [\tilde{A}_\alpha, \tilde{\psi}]\\
&= U^{-1}\left(\mathbf{D}_\alpha\psi + [\psi, U(\mathbf{D}_\alpha U^{-1})]\right) + [U^{-1}A_\alpha U + (\mathbf{D}_\alpha U^{-1})U, \, U^{-1}\psi U]\\
&= U^{-1}\left(\mathbf{D}_\alpha\psi + [\psi, (\mathbf{D}_\alpha U)U^{-1}] + [A_\alpha, \psi] + [U(\mathbf{D}_\alpha U^{-1}), \psi]\right)U\\
&= U^{-1}\left(\mathbf{D}_\alpha\psi + [A_\alpha, \psi]\right)U = \widetilde{\mathbf{D}_\alpha^{(A)}\psi}
\end{aligned}
\]
\end{proof}

如同在 Riemann 几何中, 交换两个（规范）协变导数产生一个称为\newterm{curvature}{曲率}{curvature}的基本对象, 这里记为 $F$：
\begin{equation}
\mathbf{D}_\alpha\mathbf{D}_\beta\psi - \mathbf{D}_\beta\mathbf{D}_\alpha\psi = [F_{\alpha\beta}, \psi]\label{eq:c1-49}
\end{equation}
其中曲率的分量 $F_{\alpha\beta}$ 可由如下直接的计算推出：
\[
\begin{aligned}
\mathbf{D}_\alpha\mathbf{D}_\beta\psi &= \mathbf{D}_\alpha(\mathbf{D}_\beta\psi) + [A_\alpha, \mathbf{D}_\beta\psi]\\
&= \mathbf{D}_\alpha(\mathbf{D}_\beta\psi + [A_\beta, \psi]) + [A_\alpha, \mathbf{D}_\beta\psi + [A_\beta, \psi]]\\
&= \mathbf{D}_\alpha\mathbf{D}_\beta\psi + [\mathbf{D}_\alpha A_\beta, \psi] + [A_\beta, \mathbf{D}_\alpha\psi] + [A_\alpha, \mathbf{D}_\beta\psi] + [A_\alpha, [A_\beta, \psi]]
\end{aligned}
\]
从而
\[
\begin{aligned}
(\mathbf{D}_\alpha\mathbf{D}_\beta - \mathbf{D}_\beta\mathbf{D}_\alpha)\psi &= [\mathbf{D}_\alpha A_\beta - \mathbf{D}_\beta A_\alpha, \psi]\\
&\quad + \underbrace{[A_\alpha, [A_\beta, \psi]] - [A_\beta, [A_\alpha, \psi]]}_{[[A_\alpha, A_\beta], \psi]}
\end{aligned}
\]
因此, 
\begin{equation}
F_{\alpha\beta} = \mathbf{D}_\alpha A_\beta - \mathbf{D}_\beta A_\alpha + [A_\alpha, A_\beta]\label{eq:c1-50}
\end{equation}

我们把如下事实留给读者证明：曲率张量 $F$ 在规范变换下不变. 即
\[
\widetilde{F^{(\tilde{A})}}\left(\equiv U^{-1}F^{(\tilde{A})}U\right) = F^{(A)}
\]
且 $F$ 满足 Bianchi 恒等式：
\begin{equation}
\mathbf{D}_\alpha F_{\beta\gamma} + \mathbf{D}_\gamma F_{\alpha\beta} + \mathbf{D}_\beta F_{\gamma\alpha} = 0\label{eq:c1-51}
\end{equation}

我们终于可以给出 Yang-Mills Lagrange 量所提供的 Maxwell 理论的推广：
\begin{equation}
L[A] = -\frac{1}{4}\left\langle F_{\alpha\beta}^{(A)}, F^{(A)\, \alpha\beta}\right\rangle_{\mathcal{G}}\label{eq:c1-52}
\end{equation}
我们像在 Maxwell 理论中那样推导 Euler-Lagrange 方程：
\[
\begin{aligned}
0 &= \left.\frac{d}{ds}\mathbf{S}(s)\right|_{s=0} = -\frac{1}{2}\int_{\mathcal{U}}\left\langle\dot{F}_{\alpha\beta}, F^{\alpha\beta}\right\rangle_{\mathcal{G}}dv_{\mathbf{g}}\\
&= -\frac{1}{2}\int_{\mathcal{U}}\left\langle\mathbf{D}_\alpha\dot{A}_\beta - \mathbf{D}_\beta\dot{A}_\alpha + [\dot{A}_\alpha, A_\beta] + [A_\alpha, \dot{A}_\beta], F^{\alpha\beta}\right\rangle_{\mathcal{G}}dv_{\mathbf{g}}\\
&= -\int_{\mathcal{U}}\left(\left\langle\mathbf{D}_\alpha\dot{A}_\beta, F^{\alpha\beta}\right\rangle + \left\langle[A_\alpha, \dot{A}_\beta], F^{\alpha\beta}\right\rangle_{\mathcal{G}}\right)dv_{\mathbf{g}}\\
&= \int_{\mathcal{U}}\left(\left\langle\dot{A}_\beta, \mathbf{D}_\alpha F^{\alpha\beta}\right\rangle_{\mathcal{G}} + \left\langle\dot{A}_\beta, [A_\alpha, F^{\alpha\beta}]\right\rangle_{\mathcal{G}}\right)dv_{\mathbf{g}}
\end{aligned}
\]
这蕴含
\begin{equation}
\mathbf{D}_\nu F^{\mu\nu} = 0\label{eq:c1-53}
\end{equation}
\eqref{eq:c1-51} 与 \eqref{eq:c1-53} 合在一起构成 Yang-Mills 方程组. 

注意这些方程在规范变换群下不变. 于是, Yang-Mills 方程组的一个解是规范等价势 $A_\alpha$ 的一个等价类, 其曲率 $F$ 满足 \eqref{eq:c1-53}. 

在我们以后对 Yang-Mills 的处理中, 几乎总是通过对 $A$ 施加附加的约束——称为 \newterm{gauge-condition}{规范条件}{gauge condition}——来指定解的等价类的一个代表元. 有三种标准做法, 每种都给出 Yang-Mills 方程组的一种形式, 各有其优缺点：

\begin{itemize}
\item \newterm{coulomb-gauge}{Coulomb 规范}{Coulomb gauge}定义为：
\begin{equation}
\nabla^iA_i(t, x) = 0, \qquad (t, x) \in \mathbb{R}^{n+1}\label{eq:c1-54}
\end{equation}
为简化记号, 先把 \eqref{eq:c1-53} 用流 $J$ 写出：
\begin{equation}
\mathbf{D}^\beta F_{\alpha\beta} = J_\alpha = -\left[A^\beta, F_{\alpha\beta}\right]\label{eq:c1-55}
\end{equation}
当 $\alpha = 0$ 时, \eqref{eq:c1-54} 使我们能把 \eqref{eq:c1-55} 写成：
\[
J_0 = \partial^iF_{0i} = \partial^i(\partial_0A_i - \partial_iA_0 + [A_0, A_i]) - \Delta A_0 + \partial^i[A_0, A_i]
\]
从而给出 $A$ 的时间分量：
\begin{equation}
\Delta A_0 = 2[\partial_iA_0, A_i] + [A_0, \partial_tA_i] + [A_i, [A_0, A_j]]\label{eq:c1-56}
\end{equation}
当 $\alpha = i$ 时, \eqref{eq:c1-55} 读作：
\[
J_i = -\partial_tF_{i0} + \partial^jF_{ij} = -\partial_t(\partial_iA_0 + [A_i, A_0]) + \partial^j(\partial_iA_j - \partial_jA_i + [A_i, A_j])
\]
化简后：
\begin{equation}
\begin{aligned}
\Box A_i &= -\partial_t\partial_iA_0 - 2[A_j, \partial_jA_i] + [A_j, \partial_iA_j] + [\partial_tA_i, A_j]\\
&\quad + 2[A_0, \partial_tA_i] - [A_0, \partial_iA_0] - [A_j, [A_j, A_i]] + [A_0, [A_0, A_i]]
\end{aligned}\label{eq:c1-57}
\end{equation}

\item \newterm{lorentz-gauge}{Lorentz 规范}{Lorentz gauge}由下式指定：
\begin{equation}
\partial^\mu A_\mu(t, x) = 0, \qquad (t, x) \in \mathbb{R}^{3+1}\label{eq:c1-58}
\end{equation}
Lorentz 规范对 $A$ 的时间分量与空间分量的对称处理颇具吸引力, 它还使 \eqref{eq:c1-53} 能写成一组波动方程：
\[
\begin{aligned}
\mathbf{D}^\beta F_{\alpha\beta} &= \mathbf{D}^\beta(\partial_\alpha A_\beta - \partial_\beta A_\alpha + [A_\alpha, A_\beta])\\
&= -\Box A_\alpha + \partial^\beta[A_\alpha, A_\beta] + [A_\beta, \partial_\alpha A_\beta] - [A^\beta, \partial_\beta A_\alpha] + [A_\beta, [A_\alpha, A^\beta]]
\end{aligned}
\]
该方程组可以示意性地写成：
\[
\Box\Phi = \Phi\cdot\partial\Phi + \Phi^3
\]
同样, 能否把任意解变换到 Lorentz 规范完全不清楚. 此外, 要为此纯双曲型的非线性波动方程组找到好的估计也很困难. 

\item \newterm{temporal-gauge}{时间规范}{temporal gauge}由条件 $A_0 = 0$ 指定. 
\end{itemize}

\mtsubsecn{5}{Einstein 场方程}

根据广义相对论变分原理, 时空度规 $g$ 自身相对于某个作用量是驻定的：
\[
\mathcal{S} = \int_{\mathcal{U}}L\, dv_{\mathbf{g}}.
\]
这里 $U$ 是 $(\mathbf{M}, \mathbf{g})$ 的相对紧区域, L（Lagrange 量）假定是 $M$ 上的标量函数, 其对度规的依赖不超过两阶导数\footnote{事实上我们只要求相应的 Euler-Lagrange 方程对度规的导数不超过二阶. }. 还假定它依赖于我们时空中存在的 \newterm{matter-field}{物质场}{matter field} $\psi = \psi^{(1)}, \psi^{(2)}, \dots, \psi^{(p)}$. 

事实上我们写
\[
\mathcal{S} = \mathcal{S}_G + \mathcal{S}_M
\]
其中
\[
\begin{aligned}
\mathcal{S}_G &= \int_{\mathcal{U}}L_G\, dv_{\mathbf{g}}\\
\mathcal{S}_M &= \int_{\mathcal{U}}L_M\, dv_{\mathbf{g}}
\end{aligned}
\]
分别表示\newterm{gravitational-field}{引力场}{gravitational field}与物质的（作用量）. 物质 Lagrange 量 $L_M$ 只依赖于物质场 $\psi$（假定为协变张量场）以及时空度规的逆 $\mathbf{g}^{\alpha\beta}$, 后者出现在张量场 $\psi$ 的缩并中, 以产生标量 $L_M$. 它还可能依赖于附加的正定度量, 这些度量不参与变分\footnote{波映射情形下的度量 $h$ 或 Yang-Mills 方程组情形下的 Killing 标量积即属此类. }. 

引力 Lagrange 量 $L_G$ 应当是其对应 Euler-Lagrange 方程至多含度规两阶导数的度规标量不变量, 而唯一的可能候选由\newterm{scalar-curvature}{标量曲率}{scalar curvature} $R$ 给出\footnote{至多相差一个加性常数. }. 因此我们设
\[
L_G = \mathbf{R}.
\]

现在考虑度规 $g$ 的紧变分 $(\mathbf{g}_{(s)}, \mathcal{U})$. 设 $\dot{\mathbf{g}}_{\mu\nu} = \left.\frac{d}{ds}\mathbf{g}_{\mu\nu}\right|_{s=0}$. 于是对小的 s, $\mathbf{g}_{\mu\nu}(s) = \mathbf{g}_{\mu\nu} + s\dot{\mathbf{g}}_{\mu\nu} + O(s^2)$. 又有 $\mathbf{g}^{\mu\nu}(s) = \mathbf{g}^{\mu\nu} - s\dot{\mathbf{g}}^{\mu\nu} + O(s^2)$, 其中 $\dot{\mathbf{g}}^{\mu\nu} = \mathbf{g}^{\alpha\mu}\mathbf{g}^{\beta\nu}\dot{\mathbf{g}}_{\alpha\beta}$. 则
\[
\left.\frac{d}{ds}\mathbf{S}_G(s)\right|_{s=0} = \int_{\mathcal{U}}\dot{\mathbf{R}}\, dv_{\mathbf{g}} + \int_{\mathcal{U}}\mathbf{R}\, d\dot{v}_{\mathbf{g}}
\]
现在
\[
\dot{dv_{\mathbf{g}}} = \frac{1}{2}\mathbf{g}^{\mu\nu}\dot{\mathbf{g}}_{\mu\nu}\, dv_{\mathbf{g}}
\]
事实上, 相对于一个坐标系, $dv_{\mathbf{g}} = \sqrt{-\mathbf{g}}\, dx^0dx^1\ldots dx^n$. 于是, 上述等式由
\[
\dot{\mathbf{g}} = \mathbf{g}\, \mathbf{g}^{\alpha\beta}\dot{\mathbf{g}}_{\alpha\beta}
\]
推出, 其中 $g$ 是 $\mathbf{g}_{\alpha\beta}$ 的行列式. 另一方面, 写 $\mathbf{R} = \mathbf{g}^{\mu\nu}\mathbf{R}_{\mu\nu}$ 并利用公式 $\left.\frac{d}{ds}\mathbf{g}_{(s)}^{\mu\nu}\right|_{s=0} = -\dot{\mathbf{g}}^{\mu\nu}$, 我们计算得 $\dot{\mathbf{R}} = -\dot{\mathbf{g}}^{\mu\nu}\mathbf{R}_{\mu\nu} + \mathbf{g}^{\mu\nu}\dot{\mathbf{R}}_{\mu\nu}$. 因此, 
\begin{equation}
\left.\frac{d}{ds}\mathbf{S}_G(s)\right|_{s=0} = -\int_{\mathcal{U}}\left(\mathbf{R}^{\mu\nu} - \frac{1}{2}\mathbf{g}^{\mu\nu}\mathbf{R}\right)\dot{\mathbf{g}}_{\mu\nu}\, dv_{\mathbf{g}} + \int_{\mathcal{U}}\mathbf{g}^{\mu\nu}\dot{\mathbf{R}}_{\mu\nu}\, dv_{\mathbf{g}}\label{eq:c1-59}
\end{equation}

为计算 $\dot{\mathbf{R}}_{\mu\nu}$, 我们利用下面的引理：

\begin{lemma}\label{lem:c1-1.2}
设 $\mathbf{g}_{\mu\nu}(s)$ 是一族时空度规, 满足 $\mathbf{g}(0) = \mathbf{g}$ 且 $\left.\frac{d}{ds}\mathbf{g}(0)\right. = \dot{\mathbf{g}}$. 又设 $\left.\frac{d}{ds}\mathbf{R}_{\alpha\beta}(s)\right|_{s=0} = \dot{\mathbf{R}}_{\alpha\beta}$. 则
\[
\dot{\mathbf{R}}_{\mu\nu} = \mathbf{D}_\alpha\dot{\Gamma}^\alpha_{\mu\nu} - \mathbf{D}_\mu\dot{\Gamma}^\alpha_{\alpha\nu}
\]
其中 $\dot{\Gamma}$ 是张量：
\[
\dot{\Gamma}^\alpha_{\beta\gamma} = \frac{1}{2}\mathbf{g}^{\alpha\lambda}\left(\mathbf{D}_\beta\dot{\mathbf{g}}_{\gamma\lambda} + \mathbf{D}_\gamma\dot{\mathbf{g}}_{\beta\lambda} - \mathbf{D}_\lambda\dot{\mathbf{g}}_{\beta\gamma}\right)
\]
\end{lemma}

\begin{proof}
由于恒等式两边都是张量, 只需在一点 $p$ 处、相对于一个使 Christoffel 符号 $\Gamma$ 在 $p$ 处为零的特殊坐标系来证明公式. 相对于这样的坐标系, Ricci 张量具有形式 $\mathbf{R}_{\mu\nu} = \partial_\alpha\Gamma^\alpha_{\mu\nu} - \partial_\mu\Gamma^\alpha_{\alpha\nu}$. 
\end{proof}

回到 \eqref{eq:c1-59}, 我们发现, 由于 $\mathbf{g}^{\mu\nu}\dot{\mathbf{R}}_{\mu\nu}$ 可以写成支集紧含于 $U$ 的张量的时空散度, 相应的积分恒为零. 因此我们推断：
\begin{equation}
\left.\frac{d}{ds}\mathbf{S}_G(s)\right|_{s=0} = -\int_{U}\mathbf{E}^{\mu\nu}\dot{\mathbf{g}}_{\mu\nu}\, dv_{\mathbf{g}}\label{eq:c1-60}
\end{equation}
其中 $\mathbf{E}^{\mu\nu} = \mathbf{R}^{\mu\nu} - \textstyle\frac{1}{2}\mathbf{g}^{\mu\nu}\mathbf{R}$. 现在我们考虑作用量积分 $\mathbf{S}_M$ 关于度规的变分. 如前所述, $L_M$ 通过度规的逆 $\mathbf{g}^{\mu\nu}$ 依赖于度规 g. 因此若记 $\mathbf{S}_M(s) = \mathbf{S}_M[\psi, \mathbf{g}_{(s)};\mathcal{U}]$, 并写 $d\dot{v}_{\mathbf{g}} = \frac{1}{2}\mathbf{g}_{\mu\nu}\dot{\mathbf{g}}^{\mu\nu}\, dv_{\mathbf{g}}$, 则有
\[
\begin{aligned}
\left.\frac{d}{ds}\mathbf{S}_M(s)\right|_{s=0} &= -\int_{\mathcal{U}}\frac{\partial L_M}{\partial\mathbf{g}^{\mu\nu}}\dot{\mathbf{g}}^{\mu\nu}\, dv_{\mathbf{g}} + \int_{\mathcal{U}}L_M\, d\dot{v}_{\mathbf{g}}\\
&= -\int_{\mathcal{U}}\left(\frac{\partial L_M}{\partial\mathbf{g}^{\mu\nu}} - \frac{1}{2}\mathbf{g}_{\mu\nu}L_M\right)\dot{\mathbf{g}}^{\mu\nu}\, dv_{\mathbf{g}}
\end{aligned}
\]

\begin{definition*}
对称张量
\[
\mathbf{T}_{\mu\nu} = -\left(\frac{\partial L_M}{\partial\mathbf{g}^{\mu\nu}} - \frac{1}{2}\mathbf{g}_{\mu\nu}L_M\right)
\]
称为作用量 $\mathbf{S}_M$ 的能量-动量张量. 
\end{definition*}

有了这个定义, 我们写
\begin{equation}
\left.\frac{d}{ds}\mathbf{S}_M(s)\right|_{s=0} = \int_{\mathcal{U}}\mathbf{T}^{\mu\nu}\dot{\mathbf{g}}_{\mu\nu}\, dv_{\mathbf{g}}\label{eq:c1-61}
\end{equation}
最后, 将 \eqref{eq:c1-60} 与 \eqref{eq:c1-61} 结合, 对总作用量 $\mathbf{S}$ 我们得到：
\[
\left.\frac{d}{ds}\mathbf{S}(s)\right|_{s=0} = -\int_{\mathcal{U}}\left(\mathbf{E}^{\mu\nu} - \mathbf{T}^{\mu\nu}\right)\dot{\mathbf{g}}_{\mu\nu}\, dv_{\mathbf{g}}
\]
由于 $\dot{\mathbf{g}}_{\mu\nu}$ 是支集紧含于 $U$ 的任意对称 2-张量, 我们得到 Einstein 场方程：
\[
\mathbf{E}^{\mu\nu} = \mathbf{T}^{\mu\nu}
\]
回忆 \newterm{einstein-tensor}{Einstein 张量}{Einstein tensor} $E$ 满足 \newterm{twice-contracted-bianchi}{两次缩并的 Bianchi 恒等式}{twice contracted Bianchi identity}：
\[
\mathbf{D}^\nu\mathbf{E}_{\mu\nu} = 0
\]
这意味着能量-动量张量 $T$ 也无散度：
\begin{equation}
\mathbf{D}_\nu\mathbf{T}^{\mu\nu} = 0\label{eq:c1-62}
\end{equation}
这是物质场能量-动量守恒律的简洁的时空表述. 

\section{能量-动量张量}

守恒律 \eqref{eq:c1-62} 是物质场的基本性质. 我们现在转向一个更直接的推导. 

我们考虑具有驻定解 $\psi$ 的任意 Lagrange 场论. 设 $\Phi_s$ 是由给定向量场 $X$ 生成的局部\newterm{diffeomorphism}{微分同胚}{diffeomorphism}单参数群. 我们将用流 $\Phi$ 按如下方式变分场 $\psi$：
\[
\begin{aligned}
\mathbf{g}_s &= (\Phi_s)_*\mathbf{g}\\
\psi_s &= (\Phi_s)_*\psi.
\end{aligned}
\]
由作用量积分在微分同胚下的不变性, 
\[
\mathbf{S}(s) = \mathbf{S}[\psi_s, \mathbf{g}_s;\mathbf{M}] = \mathbf{S}[\psi, \mathbf{g};\mathbf{M}].
\]
于是
\begin{equation}
0 = \left.\frac{d}{ds}\mathbf{S}(s)\right|_{s=0} = \int_{\mathbf{M}}\frac{\delta\mathbf{S}}{\delta\psi}\, dv_{\mathbf{g}} + \int_{\mathbf{M}}\mathbf{T}^{\mu\nu}\dot{\mathbf{g}}_{\mu\nu}\, dv_{\mathbf{g}}\label{eq:c1-63}
\end{equation}
第一项显然为零, 因为 $\psi$ 是驻定解. 在第二项中, 它表示关于度规的变分, 我们有
\[
\dot{\mathbf{g}}_{\mu\nu} = \left.\frac{d}{ds}(\mathbf{g}_s)_{\mu\nu}\right|_{s=0} = \mathcal{L}_X\mathbf{g}_{\mu\nu} = \mathbf{D}_\mu X_\nu + \mathbf{D}_\nu X_\mu
\]
因此
\[
0 = \int_{\mathbf{M}}\mathbf{T}^{\mu\nu}\mathcal{L}_X\mathbf{g}_{\mu\nu}\, dv_{\mathbf{g}} = 2\int_{\mathbf{M}}\mathbf{T}^{\mu\nu}\mathbf{D}_\nu X_\mu\, dv_{\mathbf{g}} = -2\int_{\mathbf{M}}\mathbf{D}_\nu\mathbf{T}^{\mu\nu}X_\mu\, dv_{\mathbf{g}}
\]
由于 $X$ 是任意的, 我们断定
\begin{equation}
\mathbf{D}_\nu\mathbf{T}^{\mu\nu} = 0.\label{eq:c1-64}
\end{equation}
这再次是能量-动量守恒律. 

下面列出前面讨论过的场论的能量-动量张量. 我们把利用定义进行计算的工作留给读者. 

\begin{enumerate}
\item 标量场方程的能量-动量张量为：
\[
\mathbf{T}_{\alpha\beta} = \frac{1}{2}\left(\phi_{, \alpha}\phi_{, \beta} - \frac{1}{2}\mathbf{g}_{\alpha\beta}\left(\mathbf{g}^{\mu\nu}\phi_{, \mu}\phi_{, \nu} + 2V(\phi)\right)\right)
\]
\item 波映射的能量-动量张量由下式给出：
\[
\mathbf{T}_{\alpha\beta} = \frac{1}{2}\left(\langle\phi_{, \alpha}, \phi_{, \beta}\rangle - \frac{1}{2}\mathbf{g}_{\alpha\beta}\left(\mathbf{g}^{\mu\nu}\langle\phi_{, \mu}, \phi_{, \nu}\rangle\right)\right)
\]
其中 $\langle\cdot, \cdot\rangle$ 表示目标流形上的 Riemann 内积. 
\item Maxwell 方程组的能量-动量张量为：
\[
\mathbf{T}_{\alpha\beta} = F_\alpha{}^\mu F_{\beta\mu} - \frac{1}{4}\mathbf{g}_{\alpha\beta}\left(F_{\mu\nu}F^{\mu\nu}\right)
\]
\item Yang-Mills 方程组的能量-动量张量为：
\[
\mathbf{T}_{\alpha\beta} = \left\langle F_\alpha{}^\mu, F_{\beta\mu}\right\rangle - \frac{1}{4}\mathbf{g}_{\alpha\beta}\left(\left\langle F_{\mu\nu}, F^{\mu\nu}\right\rangle\right)
\]
\end{enumerate}

能量-动量张量 $T$ 的一个可接受概念除守恒律 \eqref{eq:c1-64} 外还必须满足如下性质：
\begin{enumerate}
\item T 是对称的；
\item T 满足 \newterm{positive-energy-condition}{正能量条件}{positive energy condition}, 即对任意未来指向类时向量 X, Y, 有 $\mathbf{T}(X, Y) \geq 0$. 
\end{enumerate}

对称性在我们的构造中是自动成立的. 下面的命题断言, 上述场论的能量-动量张量满足正能量条件. 

\begin{proposition}\label{prop:c1-2.1}
若 $V$ 为正, 则标量波动方程的能量-动量张量满足正能量条件. 波映射、Maxwell 方程组与 Yang-Mills 的能量-动量张量满足正能量条件. 
\end{proposition}

\begin{proof}
为证明正性条件, 考虑某点 $p \in \mathbf{M}$ 处的两个向量 X, Y, 它们都是因果的且未来指向. 由 X, Y 张成的平面与 $p$ 处的\newterm{null-cone}{零锥}{null cone}沿两个零方向相交\footnote{若 X, Y 线性相关, 则任何过其公共方向的平面都可以. }. 设 L, $\underline{L}$ 是对应于两个互补零方向的未来指向零向量, 并按条件
\[
\langle L, \underline{L}\rangle = -2
\]
归一化, 即它们构成一个\newterm{null-pair}{零对}{null pair}. 由于向量场 X, Y 是 L, $\underline{L}$ 的系数为正的线性组合, 命题将由如下事实推出：$\mathbf{T}(L, L) \geq 0$、$\mathbf{T}(\underline{L}, \underline{L}) \geq 0$ 且 $\mathbf{T}(L, \underline{L}) \geq 0$. 为此, 我们考虑 $p$ 处由向量场 $E_{(n+1)} = L$、$E_{(n)} = \underline{L}$ 及 $E_{(1)}, \ldots, E_{(n-1)}$ 构成的标架, 它具有性质：
\[
\langle E_{(i)}, E_{(n)}\rangle = \langle E_{(i)}, E_{(n+1)}\rangle = 0
\]
且
\[
\langle E_{(i)}, E_{(j)}\rangle = \delta_{ij}
\]
对一切 $i, j = 1, \ldots, n-1$. 具有这些性质的标架称为零标架. 

\begin{enumerate}
\item 现在我们计算波动方程的情形：
\[
\begin{aligned}
\mathbf{T}(L, L) &= \frac{1}{2}E(\phi)^2\\
\mathbf{T}(\underline{L}, \underline{L}) &= \frac{1}{2}\underline{L}(\phi)^2.
\end{aligned}
\]
它们显然非负. 现在
\[
\mathbf{T}(L, \underline{L}) = \frac{1}{2}\left[L(\phi)\underline{L}(\phi) + \left(g^{\mu\nu}\phi_{, \mu}\phi_{, \nu} + 2V(\phi)\right)\right]
\]
我们想把 $g^{\mu\nu}\phi_{, \mu}\phi_{, \nu}$ 相对于零标架表示出来. 为此, 注意相对于零标架, 度量 $g_{\alpha\beta}$ 的仅有的非零分量为：
\[
g_{n(n+1)} = -2, \qquad g_{ii} = 1, \qquad i = 1, \dots, n-1
\]
而逆度量 $g^{\alpha\beta}$ 的非零分量为：
\[
g^{n(n+1)} = -\frac{1}{2}, \qquad g^{ii} = 1, \qquad i = 1, \ldots, n-1
\]
因此, 
\[
g^{\mu\nu}\phi_{, \mu}\phi_{, \nu} = -L(\phi)\underline{L}(\phi) + |\vec{\nabla}\phi|^2
\]
其中
\[
|\nabla\phi|^2 = (E_{(1)}(\phi))^2 + (E_{(2)}(\phi))^2 + \dots + E_{(n-1)}(\phi)^2.
\]
因此, 
\[
\mathbf{T}(L, \underline{L}) = \frac{1}{2}|\vec{\nabla}\phi|^2 + V(\phi).
\]
\item 对波映射, 按照同样的计算, 我们有
\[
\begin{aligned}
T(E, E) &= \frac{1}{2}\langle E(\phi), E(\phi)\rangle\\
T(\underline{E}, \underline{E}) &= \frac{1}{2}\langle\underline{E}(\phi), \underline{E}(\phi)\rangle\\
T(E, \underline{E}) &= \frac{1}{2}\sum_{i=1}^{n-1}\langle E_{(i)}(\phi), E_{(i)}(\phi)\rangle.
\end{aligned}
\]
T 的正性于是由目标流形 $N$ 上的 Riemann 度量 $h$ 推出. 
\item 为证明 3+1 维 Maxwell 方程组能量-动量张量的正性, 我们先把张量写成更对称的形式：
\begin{equation}
T_{\alpha\beta} = \frac{1}{2}\left(F_\alpha{}^\mu F_{\beta\mu} + {}^\star F_\alpha{}^\mu\, {}^\star F_{\beta\mu}\right)\label{eq:c1-65}
\end{equation}
其中 ${}^\star F$ 是 $F$ 的 Hodge 对偶, 即 ${}^\star F_{\alpha\beta} = \textstyle\frac{1}{2}\epsilon_{\alpha\beta\mu\nu}F^{\mu\nu}$. 

\textbf{习题.} 验证公式 \eqref{eq:c1-65}. 

我们在每点 $p \in \mathbf{M}$ 处引入 $F$ 的如下\newterm{null-decomposition}{零分解}{null decomposition}：
\[
\begin{aligned}
\alpha_A &= F_{A4}, \qquad \underline{\alpha}_A = F_{A3}\\
\rho &= \frac{1}{2}F_{34}, \qquad \sigma = \frac{1}{2}\, {}^\star F_{34}.
\end{aligned}
\]
它完全确定了张量 F. 这里指标 $A = 1, 2$ 对应于与球面相切的 $E_1, E_2$ 方向, 而指标 3, 4 对应于 $E_3 = \underline{L}$ 与 $E_4 = L$. 然后我们对 ${}^\star F$ 计算得：
\[
\begin{aligned}
{}^\star F_{A4} &= -{}^\star\alpha_A, \qquad {}^\star F_{A3} = {}^\star\underline{\alpha}_A\\
{}^\star F_{34} &= 2\sigma, \qquad {}^{\star\star}F_{34} = -2\rho
\end{aligned}
\]
其中 ${}^\star\alpha_A = \epsilon_{AB}\alpha_B$. 这里 $\epsilon_{AB}$ 是单位球面上的体积形式, 故 $\epsilon_{AB} = \textstyle\frac{1}{2}\epsilon_{AB34}$, 即 $\epsilon_{11} = \epsilon_{22} = 0$, $\epsilon_{12} = -\epsilon_{21} = 1$. 用这一记号, 我们计算得：
\[
\begin{aligned}
T(E_{(4)}, E_{(4)}) &= \frac{1}{2}\sum_{A=1}^2\left(F_{4A}\cdot F_{4A} + {}^\star F_{4A}\cdot{}^\star F_{4A}\right)\\
&= \frac{1}{2}\sum_{A=1}^2\left(\alpha_A\cdot\alpha_A + {}^\star\alpha_A\cdot{}^\star\alpha_A\right)\\
&= \sum_{A=1}^2\alpha_A\cdot\alpha_A = |\alpha|^2 \geq 0.
\end{aligned}
\]
类似地, 
\[
T(E_{(3)}, E_{(3)}) = \sum_{A=1}^2\underline{\alpha}_A\cdot\underline{\alpha}_A = |\underline{\alpha}|^2 \geq 0
\]
用同样的方法我们得到
\[
T(E, \underline{E}) = \rho^2 + \sigma^2 \geq 0
\]
这就证明了我们的断言. 
\item Yang-Mills 方程组能量-动量张量的正性用与 Maxwell 方程组完全相同的方式证明, 利用 Killing 标量积 $\langle\cdot, \cdot\rangle_{\mathcal{G}}$ 的正性. 
\end{enumerate}
\end{proof}

场论的能量-动量张量可能满足的另一个重要性质是无迹条件, 即
\[
\mathbf{g}_{\alpha\beta}\mathbf{T}^{\alpha\beta} = 0.
\]
事实证明, 这一条件被所有共形不变的场论满足. 

\begin{definition*}
若相应的作用量积分在度规的共形变换
\[
\mathbf{g}_{\alpha\beta} \longrightarrow \tilde{\mathbf{g}}_{\alpha\beta} = \Omega\mathbf{g}_{\alpha\beta}
\]
下不变, 则称场论是共形不变的, 其中 $\Omega$ 是时空上的正光滑函数. 
\end{definition*}

\begin{proposition}\label{prop:c1-2.2}
共形不变场论的能量-动量张量 $\mathbf{T}$ 是无迹的. 
\end{proposition}

\begin{proof}
考虑支集紧含于 $\mathcal{U} \subset \mathcal{M}$ 的任意光滑函数 f. 考虑给定度规 $g$ 的如下变分：
\[
\mathbf{g}_{\mu\nu}(s) = e^{sf}\mathbf{g}_{\mu\nu}.
\]
设 $S(s) = S_{\mathcal{U}}[\psi, \mathbf{g}(s)]$. 由 $S$ 的协变性, $S(s) = S(0)$. 因此
\[
0 = \left.\frac{d}{ds}\mathcal{S}(s)\right|_{s=0} = \int_{\mathcal{U}}T^{\mu\nu}\dot{\mathbf{g}}_{\mu\nu}\, dv_{\mathbf{g}}
\]
其中
\[
\dot{\mathbf{g}}_{\mu\nu} = \left.\frac{d}{ds}\mathbf{g}_{\mu\nu}(s)\right|_{s=0} = f\mathbf{g}_{\mu\nu}.
\]
于是 $\int_{\mathcal{U}}\left(T^{\mu\nu}\mathbf{g}_{\mu\nu}\right)f\, dv_{\mathbf{g}} = 0$, 又因 $f$ 任意, 我们推断：
\[
tr\, T = g^{\mu\nu}T_{\mu\nu} \equiv 0.
\]
\end{proof}

我们容易验证, Maxwell 方程组与 Yang-Mills 方程组在 $3\times 1$ 维中是共形不变的. 波映射场论在 $1+1$ 维中是共形不变的, 即当时空 $M$ 是二维时\footnote{线性标量波动方程亦然. }. 

\begin{remark*}
Maxwell 方程组的作用量积分 $\mathbf{S} = \int_{\mathcal{U}}F_{\alpha\beta}F^{\alpha\beta}dv_{\mathbf{g}}$ 在任意维数都是共形不变的, 前提是我们同时按比例缩放电磁场 F. 事实上, 若 $\tilde{\mathbf{g}}_{\alpha\beta} = \Omega^2\mathbf{g}_{\alpha\beta}$, 则 $dv_{\tilde{\mathbf{g}}} = \Omega^{n+1}dv_{\mathbf{g}}$；若再设 $\tilde{F}_{\alpha\beta} = \Omega^{-\frac{n-3}{2}}F_{\alpha\beta}$, 则得到：
\[
\begin{aligned}
\tilde{\mathbf{S}}[\tilde{F}, \tilde{\mathbf{g}}] &= \int \tilde{F}_{\alpha\beta}\tilde{F}_{\gamma\delta}\tilde{\mathbf{g}}^{\alpha\gamma}\tilde{\mathbf{g}}^{\beta\delta}dv_{\tilde{\mathbf{g}}}\\
&= \int F_{\alpha\beta}F_{\gamma\delta}\mathbf{g}^{\alpha\gamma}\mathbf{g}^{\beta\delta}dv_{\mathbf{g}}\\
&= \mathbf{S}[F, \mathbf{g}].
\end{aligned}
\]
\end{remark*}

我们以一个关于 $1+1$ 维共形场论的简单观察结束本节. 事实上我们特化到 Minkowski 空间 $\mathbb{R}^{1+1}$, 并考虑局部守恒律 $\partial^\mu\mathbf{T}_{\nu\mu} = 0$. 令 $\nu = 0, 1$, 我们得到：
\begin{equation}
\partial^0\mathbf{T}_{00} + \partial^1\mathbf{T}_{01} = 0, \qquad \partial^0\mathbf{T}_{01} + \partial^1\mathbf{T}_{11} = 0\label{eq:c1-66}
\end{equation}
由于能量-动量张量无迹, 我们有 $\mathbf{T}_{00} = \mathbf{T}_{11} := A$. 设 $\mathbf{T}_{01} = \mathbf{T}_{10} = B$. 因此 \eqref{eq:c1-66} 蕴含 $A$ 与 $B$ 都满足线性齐次波动方程：
\begin{equation}
\Box A = 0 = \Box B.\label{eq:c1-67}
\end{equation}
利用这一观察, 容易证明光滑初值对一切时间保持光滑. 

例如, 波映射在 $1+1$ 维中是共形不变的. 此时
\[
A = \mathbf{T}_{00} = \frac{1}{2}\left(\langle\partial_t\phi, \partial_t\phi\rangle + \langle\partial_x\phi, \partial_x\phi\rangle\right),
\]
给定 $C_0^\infty(\mathbb{R})$ 中的数据, \eqref{eq:c1-67} 蕴含 $\phi$ 的导数对一切正时间保持光滑. 这就证明了\newterm{global-existence}{全局存在性}{global existence}. 

\section{守恒律}

场论的能量-动量张量与守恒律密切相关. 这一联系通过 Noether 原理得以显现：

\textbf{Noether 原理:} 对任一保持作用的单参数变换群, 对应一个守恒律. 

我们如下阐明这一基本原理：设 $\mathbf{S} = \mathbf{S}[\psi, \mathbf{g}]$ 是场 $\psi$ 的作用量积分. 设 $\chi_t$ 是 $M$ 的\newterm{isometry}{等距}{isometry}单参数群, 即 $(\chi_t)_*\mathbf{g} = \mathbf{g}$. 则
\[
\begin{aligned}
\mathbf{S}[(\chi_t)_*\psi, \mathbf{g}] &= \mathbf{S}[(\chi_t)_*\psi, (\chi_t)_*\mathbf{g}]\\
&= \mathbf{S}[\psi, \mathbf{g}].
\end{aligned}
\]
于是作用在 $\psi \to (\chi_t)_*\psi$ 下保持不变. 根据 Noether 原理, 我们应该能为相应的 Euler-Lagrange 方程找到守恒律\footnote{同样的论证也适用于作用在共形不变场论上的共形等距. 因此在这种设定下我们也预期有守恒律. }. 我们用生成 $\chi_t$ 的 Killing 向量场 $X$ 来推导这些守恒律. 

我们从一个涉及 $\psi$ 的能量-动量张量 $T$ 与任意向量场 $X$ 的一般计算开始. 设 $P$ 是通过将 $T$ 与 $X$ 缩并得到的 1-形式：
\[
P_\alpha = \mathbf{T}_{\alpha\beta}X^\beta
\]
由于 $T$ 是对称且无散度的, 
\[
\mathbf{D}^\alpha P_\alpha = \left(\mathbf{D}^\alpha\mathbf{T}_{\alpha\beta}\right)X^\beta + \mathbf{T}_{\alpha\beta}\left(\mathbf{D}^\alpha X^\beta\right) = \frac{1}{2}\mathbf{T}^{\alpha\beta\, (X)}\pi_{\alpha\beta}
\]
其中 ${}^{(X)}\pi_{\alpha\beta}$ 是 $X$ 的\newterm{deformation-tensor}{变形张量}{deformation tensor}：
\[
{}^{(X)}\pi_{\alpha\beta} = (\mathcal{L}_X\mathbf{g})_{\alpha\beta} = \mathbf{D}_\alpha X_\beta + \mathbf{D}_\beta X_\alpha
\]

\textbf{记号.} 我们用
\[
\mathcal{N}^-(\bar{t}, \bar{x}) = \{(t, x)\, |\, 0 \leq t \leq \bar{t};\ |x-\bar{x}| = \bar{t} - t\}
\]
记顶点为 $p = (\bar{t}, \bar{x}) \in \mathbb{R}^{n+1}$ 的\newterm{backward-light-cone}{后向光锥}{backward light cone}. 该集合在某个时间区间 $[t_1, t_2]$（$t_1 \leq t_2 \leq \bar{t}$）上的限制记为 $\mathcal{N}^-_{[t_1, t_2]}(\bar{t}, \bar{x})$. 这些零超曲面是如下集合的零边界：
\[
\begin{aligned}
\mathcal{J}^-(\bar{t}, \bar{x}) &= \{(t, x)\, |\, 0 \leq t \leq \bar{t};\ |x-\bar{x}| \leq \bar{t} - t\}\\
\mathcal{J}^-_{[t_2, t_1]}(\bar{t}, \bar{x}) &= \{(t, x)\, |\, t_2 \leq t \leq t_1;\ |x-\bar{x}| \leq \bar{t} - t\}
\end{aligned}
\]
我们用 $S_t = S_t(\bar{t}, \bar{x})$ 与 $B_t = B_t(\bar{t}, \bar{x})$ 分别记时间切片 $\Sigma_t$ 与 $\mathcal{N}^-$ 及 $\mathcal{I}^-$ 的交. 

在沿 $\mathcal{N}^-(p)$ 的每一点 $q = (t, x)$ 处, 我们定义未来指向零向量构成的零对 $(E_+, E_-)$：
\[
\underline{L} = E_+ = \partial_t + \frac{x^i - \bar{x}^i}{|x-\bar{x}|}\partial_i, \qquad L = E_- = \partial_t - \frac{x^i - \bar{x}^i}{|x-\bar{x}|}\partial_i
\]
注意 $L$ 与 $\underline{L}$ 都是零的, 且 $\langle L, \underline{L}\rangle = -2$. 

以下是 Stokes 定理的简单推论, 采用如下形式. 

\begin{proposition}\label{prop:c1-3.1}
设 $P_\mu$ 是满足 $\partial^\mu P_\mu = F$ 的 1-形式. 则\footnote{(68) 中的括号 $\langle\cdot, \cdot\rangle$ 表示关于 Minkowski 度量的内积. }, 对一切 $t_1 \leq t_2 \leq \bar{t}$, 
\begin{equation}
\int_{B_{t_2}}\langle P, \partial_t\rangle + \int_{\mathcal{N}^-_{[t_1, t_2]}(p)}\langle P, E_-\rangle = \int_{B_{t_1}}\langle P, \partial_t\rangle - \int_{\mathcal{J}^-_{[t_1, t_2]}(p)}F\, dt\, dx\label{eq:c1-68}
\end{equation}
其中
\[
\int_{\mathcal{N}^-_{[t_1, t_2]}(p)}\langle P, E_-\rangle = \int_{t_1}^{t_2}dt\int_{S_t}\langle P, E_-\rangle\, da_t.
\]
\end{proposition}

将 Stokes 定理应用于 \eqref{eq:c1-68}, 即得：

\begin{theorem}\label{thm:c1-3.2}
设 $T$ 是与某个场论相关的能量-动量张量, X 是任意向量场. 则
\begin{equation}
\begin{aligned}
\int_{B_{t_2}}\mathbf{T}(\partial_t, X) + \int_{\mathcal{N}^-_{[t_1, t_2]}(p)}\mathbf{T}(E_-, X) &= \int_{B_{t_1}}\mathbf{T}(\partial_t, X)\\
&\quad - \int_{\mathcal{J}^-_{[t_1, t_2]}(p)}\mathbf{T}^{\alpha\beta\, (X)}\pi_{\alpha\beta}\, dt\, dx
\end{aligned}\label{eq:c1-69}
\end{equation}
\end{theorem}

在 $X$ 是 Killing 的特殊情形, 其变形张量 $\pi$ 恒为零. 因此：

\begin{corollary}\label{cor:c1-3.3}
若 $X$ 是 Killing 向量场, 则
\begin{equation}
\int_{B_{t_2}}\mathbf{T}(\partial_t, X) + \int_{\mathcal{N}^-_{[t_1, t_2]}(p)}\mathbf{T}(L, X) = \int_{B_{t_1}}\mathbf{T}(\partial_t, X)\label{eq:c1-70}
\end{equation}
\end{corollary}

此外, 若 $T$ 无迹且 $X$ 是共形 Killing 的, 则 \eqref{eq:c1-70} 仍然成立. 

恒等式 \eqref{eq:c1-70} 通常应用于类时未来指向的 Killing 向量场 X, 此时 $T$ 的正能量条件保证 \eqref{eq:c1-70} 中所有被积函数为正. 我们知道（见附录 4.2, 除 Lorentz 变换外, 唯一的 Killing 未来指向类时向量场是 $\partial_t$ 的常数倍）. 取 $X = \partial_t$, \eqref{eq:c1-70} 变为：
\begin{equation}
\int_{B_{t_2}}\mathbf{T}(\partial_t, \partial_t) + \int_{\mathcal{N}^-_{[t_1, t_2]}(p)}\mathbf{T}(E_-, \partial_t) = \int_{B_{t_1}}\mathbf{T}(\partial_t, \partial_t)\label{eq:c1-71}
\end{equation}
在共形场论的情形, 我们可以取 $X$ 为未来类时的共形 Killing 向量场 $X = K_0 = (t^2+|x|^2)\partial_t + 2tx^i\partial_i$. 于是：
\begin{equation}
\int_{B_{t_2}}\mathbf{T}(\partial_t, K_0) + \int_{\mathcal{N}^-_{[t_1, t_2]}(p)}\mathbf{T}(L, K_0) = \int_{B_{t_1}}\mathbf{T}(\partial_t, K_0)\label{eq:c1-72}
\end{equation}

在 \eqref{eq:c1-71} 中, 项 $\mathbf{T}(\partial_t, \partial_t)$ 称为\newterm{energy-density}{能量密度}{energy density}, 而 $\mathbf{T}(E_-, \partial_t)$ 称为\newterm{energy-flux-density}{能量通量密度}{energy flux density}. 相应的积分分别称为 $B_{t_1}$ 与 $B_{t_2}$ 中所含的能量, 以及穿过 $\mathcal{N}^-$ 的能量通量. \eqref{eq:c1-72} 中的相应项称为共形能量密度、共形通量等. 

方程 \eqref{eq:c1-71} 可用于推导相对论场论的如下基本性质：
\begin{enumerate}
\item \newterm{finite-propagation-speed}{有限传播速度}{finite propagation speed}
\item Cauchy 问题唯一性
\end{enumerate}

\begin{proof}
第一个性质由如下事实推出：若 $\int_{B_{t_1}}\mathbf{T}(\partial_t, \partial_t)$ 在时刻 $t = t_1$ 为零, 则积分 $\int_{B_{t_2}}\mathbf{T}(\partial_t, \partial_t)$ 与 $\int_{\mathcal{N}^-_{[t_1, t_2]}}\mathbf{T}(E_-, \partial_t)$ 也都必须为零. 由 $T$ 的正性性质可知, 相应的被积函数也必为零. 考虑具体理论中 $T$ 的具体形式, 即可证明场在球 $B_{t_1}$ 的\newterm{domain-of-influence}{影响域}{domain of influence}内也恒为零. 反之, 若场的初值在 $B_{t_1}$ 的补集中为零, 则场在 $B_{t_1}$ 影响域的补集中恒为零. 

第二个性质的证明对线性场论立即由第一个性质推出. 对非线性理论则需要多费一点功夫. 
\end{proof}

\textbf{习题 1.} 为我们迄今遇到的每个场论——标量波动方程 (SWE)、波映射 (WM)、Maxwell 方程组 (ME) 与 Yang-Mills (YM)——建立初值问题. 对光滑解证明初值问题解的唯一性. 

以下是 \eqref{eq:c1-71} 与 \eqref{eq:c1-72} 的另一个重要推论. 为陈述结果, 我们引入如下量：
\begin{equation}
\mathcal{E}(t) = \int_{\mathbb{R}^n}\mathbf{T}(\partial_t, \partial_t)(t, x)\, dx\label{eq:c1-73}
\end{equation}
\begin{equation}
\mathcal{E}_c(t) = \int_{\mathbb{R}^n}\mathbf{T}(K_0, \partial_t)(t, x)\, dx\label{eq:c1-74}
\end{equation}

\begin{theorem}[全局能量]\label{thm:c1-3.4}
对任意场论, 若 $\mathcal{E}(0) < \infty$, 则
\begin{equation}
\mathcal{E}(t) = \mathcal{E}(0)\label{eq:c1-75}
\end{equation}
此外, 对共形场论, 若 $\mathcal{E}_c(0) < \infty$, 则
\begin{equation}
\mathcal{E}_c(t) = \mathcal{E}_c(0)\label{eq:c1-76}
\end{equation}
\end{theorem}

\begin{proof}
将 \eqref{eq:c1-71} 与 \eqref{eq:c1-72} 应用于过去因果域 $\mathcal{I}^-(p)$（其中 $p = (\bar{t}, 0)$, 在 $t_1 = 0$ 与 $t_2 = t$ 之间）, 并令 $\bar{t} \to +\infty$, 即可轻易得到. 
\end{proof}

\textbf{习题 2.} 考虑 Lagrange 量
\[
\mathcal{L} = -\frac{1}{2}\mathbf{m}^{\alpha\beta}\partial_\alpha\phi\, \overline{\partial_\beta\phi} - V(|\phi|)
\]
其中 $\phi$ 是定义在 Minkowski 时空 $\mathbb{R}^{n+1}$ 上的复值标量, $\bar{\phi}$ 是其复共轭. 如前所述, L 在连续变换群 $\phi \to e^{i\theta}\phi$（$\theta \in \mathbb{R}$）下不变. 根据 Noether 定理, 相应的 Euler-Lagrange 方程应该有一个守恒律. 你能推导出它吗？

\mtsubsec{能量在共形不变情形下的色散}
在本节中, 我们将利用全局共形能量恒等式 \eqref{eq:c1-76} 来说明 Minkowski 空间中场论的能量如何耗散. 考虑定义在整个 $\mathbb{R}^{n+1}$ 上的共形场论. 在 $\mathbb{R}^{n+1}$ 的每一点（其中 $t \geq 0$）, 定义标准零标架：
\[
\begin{aligned}
L = E_+ &= \partial_t + \partial_r\\
\underline{L} = E_- &= \partial_t - \partial_r.
\end{aligned}
\]
注意共形 Killing 向量场 $K_0 = (t^2+r^2)\partial_t + 2rt\partial_r$ 可以表示为：
\[
K_0 = \frac{1}{2}\left[(t+r)^2E_+ + (t-r)^2E_-\right]
\]
于是, 
\begin{equation}
\begin{aligned}
\mathcal{E}_c(t) &= \int_{\mathbb{R}^n}\left[\frac{1}{4}(t+r)^2\mathbf{T}_{++} + \frac{1}{4}(t-r)^2\mathbf{T}_{--} + \frac{1}{4}\underbrace{\left((t+r)^2+(t-r)^2\right)}_{2(t^2+r^2)}\mathbf{T}_{+-}\right]dx\\
&= \int_{\mathbb{R}^n}\left[\frac{1}{4}(t+r)^2\mathbf{T}_{++} + \frac{1}{2}(t^2+r^2)\mathbf{T}_{+-} + \frac{1}{4}(t-r)^2\mathbf{T}_{--}\right]dx\\
\mathcal{E}_c(0) &= \int_{\mathbb{R}^n}\mathbf{T}(\partial_t, K_0)(0, x)\, dx = \int_{\mathbb{R}^n}|x|^2\mathbf{T}(\partial_t, \partial_t)\, dx
\end{aligned}\label{eq:c1-77}
\end{equation}
根据 \eqref{eq:c1-76}, 我们有 $\mathcal{E}_c(t) = \mathcal{E}_c(0)$. 假定 $\mathcal{E}_c(0) = \int_{\mathbb{R}^n}|x|^2\mathbf{T}(\partial_t, \partial_t)\, dx$ 有限, 我们断定：
\[
\begin{aligned}
\int_{\mathbb{R}^n}\mathbf{T}_{++}(t, \cdot)\, dx &\lesssim \frac{\mathcal{E}_c(0)}{t^2}\\
\int_{\mathbb{R}^n}\mathbf{T}_{+-}(t, \cdot)\, dx &\lesssim \frac{\mathcal{E}_c(0)}{t^2}.
\end{aligned}
\]
\eqref{eq:c1-77} 中的剩余项含有因子 $(t-r)^2$, 它沿外向零方向 $r = t + c$ 为常数. 因此对任意 $0 < \epsilon < 1$：
\[
\begin{aligned}
\int_{|x|>(1+\epsilon)t}\mathbf{T}_{--} &= O(t^{-2})\\
\int_{|x|<(1-\epsilon)t}\mathbf{T}_{--} &= O(t^{-2}).
\end{aligned}
\]
我们断定, 共形场的大部分能量由 $\mathbf{T}_{--}$ 分量携带, 并在光锥附近传播. 

\mtsubsec{$\Box\phi=0$ 的情形}
波动方程 $\Box\phi = 0$ 只在维数 $n = 1$ 时共形不变. 然而, 在任意维数下我们仍能推导出对应于共形 Killing 向量场的有用守恒律. 

\begin{lemma}\label{lem:c1-3.7}
设 $\mathbf{T}_{\alpha\beta} = \mathbf{T}_{\alpha\beta}[\phi]$ 是 $\Box\phi = 0$ 的解的相应能量-动量张量. 设 $X$ 是共形 Killing 向量场, 即 $\pi = {}^{(X)}\pi = \mathcal{L}_X m = \Omega m$, 且 $tr\, \pi = m^{\alpha\beta}\pi_{\alpha\beta}$. 容易验证 $\Box\Omega = 0$；事实上, 在 $X = K_0$ 的特殊情形, $\Omega = 4(n+1)t$. 设
\[
\bar{P}_\alpha = \mathbf{T}_{\alpha\beta}X^\beta + \frac{n-1}{4(n+1)}tr\, \pi\, \phi\, \partial_\alpha\phi - \frac{n-1}{8(n+1)}\partial_\alpha(tr\, \pi)\, \phi^2.
\]
我们有
\[
\partial^\alpha\bar{P}_\alpha = 0.
\]
\end{lemma}

现在考虑零对 $L = \partial_t + \partial_r$、$\underline{L} = \partial_t - \partial_r$, 如同上一节. 我们容易验证：
\begin{equation}
Q(L, L) = L(\phi)^2\label{eq:c1-78}
\end{equation}
\begin{equation}
Q(L, \underline{L}) = |\nabla\phi|^2\label{eq:c1-79}
\end{equation}
\begin{equation}
Q(\underline{L}, \underline{L}) = \underline{L}(\phi)^2\label{eq:c1-80}
\end{equation}
其中 $|\nabla\phi|^2 = \sum_A|e_A(\phi)|^2$, $(e_A)_{A=1, \dots, n-1}$ 是张成 $L$ 与 $\underline{L}$ 的正交补的标准正交标架. 

% ==== part02 第 2 章:一般方程 (General Equations) ====
% 源: MinerU_latex_lecture nots.tex 2577-3674 行(pdf_text.txt 1963-2788,PDF 第 39-54 页)
% 术语(新增,键=中文=英文):
%   quasi-linear=拟线性, semi-linear=半线性, fully-nonlinear=完全非线性,
%   characteristic-curve=特征曲线, bicharacteristic=双特征曲线, characteristic-hypersurface=特征超曲面,
%   non-characteristic=非特征的, overdetermined=超定的, underdetermined=欠定的,
%   well-posed=适定的, blow-up=爆破, hamilton-jacobi=Hamilton-Jacobi 方程,
%   hamiltonian-system=Hamiltonian 方程组, eikonal=eikonal 方程, harmonic-oscillator=谐振子,
%   geometric-optics=几何光学, isometric-embedding=等距嵌入, signature=符号差,
%   picard-iteration=Picard 迭代, cauchy-kowalewski=Cauchy-Kowalewski 定理
% REFS(跨章/待解析项):
%   - 公式 (1)-(5)、(3)、(6)、(20)、(24) 为导言(第 0 章)公式,已按 part00 标签写 \eqref{eq:c0-*}(可解析,协调员复核).
%   - "Eikonal equation (149)"(源 2711、3429 行)与脚注 9 中 "(148)": 属第 3 章公式,标签未知 → 正文按原文写 (149)/(148) 纯文本,待解析.
%   - 源 3442 行 "As particular solutions of (429)": 原书(pdf_text 同)即作 (429),疑为 (120) 之笔误 → 按原文保留 (429) 纯文本.【已核对: 新版 PaddleOCR 确认原书即作 (429), 照录无误】
%   - 源 3450 行 "ω ∈ R^3"(例 2 光线): 两源一致作 R^3,数学上应为 R^d → 照录 R^3.【已核对: 新版 PaddleOCR 确认原书作 R^3, 照录无误; 数学上严格应为 R^d, 忠实原文保留】
%   - 源 3512 行 "the Dirichlet condition for (??)": 原书引用即 (??) 未解析 → 按原文保留 "(??)",待人工确认.
%   - 源 3338 行 (113) 求和上限: 源与 pdf 均作 ∑_{i=1}^{n},按数学含义改为 ∑_{i=1}^{d}.
%   - 脚注 1-14 全部按 pdf_text.txt 页尾(第 39-54 页)找回并翻译为 \footnote.
% 编号说明: 公式 (81)-(126) 共 46 个,按原书顺序用 equation 环境自动编号并加 \label{eq:c2-编号};
%   其余行间公式(原书未编号)用 \[ \]. "Remark 1./2."(引言)、"例 1./2."(第 3 节)及
%   "注 1.-注 4."(第 3 节末)为原书手工编号(无节前缀),按规范写为 \noindent\textbf{...};
%   "Propagation of singularities^10." 为加粗段首词+脚注 10;Theorem 4.1(Holmgren 唯一性)为第 4 节
%   唯一编号对象,自动显示 4.1,加 \label{thm:c2-4.1};未编号定义/定理/注/例用带 * 环境;
%   第 2 节定理[常微分方程基本定理]与第 4 节定理[Cauchy-Kowalevsky (CK)]原书均未编号 → theorem*;
%   源 3148 行 "...... To fill in the proof....." 为作者占位符,照译保留.
% ====

\chapter{一般方程}

\section*{引言}

把 PDE 定义为一门研究所有偏微分方程的学科, 是很有诱惑力的, 正如, 比如说, 代数几何处理所有多项式方程一样.按照这种观点, 这门学科的目标, 是找到关于所有 PDE、或非常一般的 PDE 类的普遍理论.尽管这种观点已经相当过时, 它仍然具有重要的价值, 我希望在下面加以说明.为了看到一般理论的完整威力, 我们至少需要写下一般方程, 同时将确保在简化情形下解释主要思想.我们考虑 $\mathbb{R}^{d}$ 中的方程或方程组, 变量为 $x = (x^{1}, x^{2}, \ldots, x^{d})$.同前, 我们用
$\partial_{i} = \frac{\partial}{\partial x^{i}}$
表示关于坐标 $x^{i}$ 的偏导数, 并用
$\partial^{\alpha} = \partial_{1}^{\alpha_{1}} \partial_{2}^{\alpha_{2}} \cdots \partial_{d}^{\alpha_{d}}$
表示对应于多重指标 $\alpha = (\alpha_{1}, \alpha_{2}, \ldots, \alpha_{d}) \in \mathbb{N}^{d}$ 的混合偏导数.我们用 $\partial^{k}$ 表示所有满足 $|\alpha| = \alpha_{1} + \cdots + \alpha_{d} = k$ 的偏导数 $\partial^{\alpha}$ 组成的向量.最后我们用
$\Lambda^{k} u = (u, \partial u, \ldots, \partial^{k} u)$
表示所有阶数小于或等于 $k$ 的偏导数组成的集合.在大多数有趣的例子中, k 是 1 或 2.

\begin{example*}
为了更清楚地说明这些记号, 考虑 $\mathbb{R}^{2}$ 的情形和坐标 $x^{1}, x^{2}$.对于多重指标 $\alpha = (2, 0)$, 我们有
$\partial^{\alpha} u = \frac{\partial}{\partial x^{1}} \frac{\partial}{\partial x^{1}} u = \partial_{1}^{2} u$,
而对于 $\alpha = (1, 1)$ 我们有
$\partial^{\alpha} u = \frac{\partial}{\partial x^{1}} \frac{\partial}{\partial x^{2}} u = \partial_{1} \partial_{2} u$.又有
\[
\partial^{2} u = \left( \frac{\partial}{\partial x^{1}} \frac{\partial}{\partial x^{1}} u, \frac{\partial}{\partial x^{1}} \frac{\partial}{\partial x^{2}} u, \frac{\partial}{\partial x^{2}} \frac{\partial}{\partial x^{2}} u \right) = (\partial_{1}^{2} u, \partial_{1} \partial_{2} u, \partial_{2}^{2} u),
\]
以及
$\Lambda^{2} u = (u, \partial_{1} u, \partial_{2} u, \partial_{1}^{2} u, \partial_{1} \partial_{2} u, \partial_{2}^{2} u)$.
\end{example*}

利用这些记号, $\mathbb{R}^{d}$ 中的 Laplace 算子具有形式
$\Delta = \partial_{1}^{2} + \partial_{2}^{2} + \cdots + \partial_{d}^{2}$,
而 Minkowski 空间 $\mathbb{R}^{d+1}$ 中的 D'Alembert 算子(D'Alembertian)具有形式
$\Box = -\partial_{t}^{2} + \partial_{1}^{2} + \cdots + \partial_{d}^{2}$.要使一个出现直至 $k$ 阶偏导数的方程有意义, 我们需要处理在每个点 $k$ 次可微的函数.采用函数类 $C^{k}$——即连续、且所有阶数 $| \alpha | \leq k$ 的偏导数 $\partial^{\alpha} u$ 均连续的函数——是方便的.

\begin{definition*}
$\mathbb{R}^{d}$ 中 $k$ 阶一般偏微分方程具有如下形式,
\begin{equation}
F\bigl(x, \Lambda^{k} u(x)\bigr) = 0, \label{eq:c2-81}
\end{equation}
其中 $F$ 是指定的函数.我们还考虑 $N \times N$ 方程组\footnote{即 $N$ 个方程、$N$ 个未知数的确定方程组.}, 此时 $F$ 和 $u$ 是列 $N$ 维向量.若 $C^{k}$ 类函数 $u$ 在 $\mathbb{R}^{d}$ 的某个指定区域的所有点 $x$ 处满足方程 \eqref{eq:c2-81}, 则称它为经典解\footnote{我们称之为经典解, 以区别于将在后面各节中讨论的广义解.}.
\end{definition*}

首先考虑一维情形 $d = 1$, 此时 \eqref{eq:c2-81} 成为常微分方程 (ODE), 或 ODE 方程组.为进一步简化, 取 $k = 1$ 和 $N = 1$, 即一阶($k = 1$)常微分方程的情形.于是 \eqref{eq:c2-81} 就是 $F(x, u(x), \partial_{x} u(x)) = 0$, 其中 $F$ 是三个变量 $x$、$u$ 和 $p = \partial_{x} u$ 的给定函数, 例如 $F(x, u, p) = x \cdot p + u^{3} - \sin x$.此时求解方程 \eqref{eq:c2-81} 就是找一个 $C^{1}$ 函数 u(x), 使得
\begin{equation}
x \cdot \partial_{x} u(x) + u^{3} = \sin x. \label{eq:c2-82}
\end{equation}

现在考虑二阶 ODE 的情形, 即 $d = N = 1$ 且 $k = 2$.于是 \eqref{eq:c2-81} 成为 $F(x, u(x), \partial_{x} u(x), \partial_{x}^{2} u(x)) = 0$, 其中 $F$ 现在依赖于四个变量 $x, u, p = \partial_{x} u, q = \partial_{x}^{2} u$.作为例子取 $F = q^{2} + V'(u)$, 其中 $V = V(u)$ 是给定的函数, 此时 \eqref{eq:c2-81} 成为非线性\newterm{harmonic-oscillator}{谐振子}{harmonic oscillator}方程,
\begin{equation}
\partial_{x}^{2} u(x) + V'(u(x)) = 0. \label{eq:c2-83}
\end{equation}

转到 ODE 方程组, 取 $d = 1, k = 1$ 且 $N = 2$, 我们需要向量函数 $F = (F_{1}, F_{2})$, 其中 $F_{1}$ 和 $F_{2}$ 都依赖于五个变量 $x, u_{1}, u_{2}, p_{1} = \partial_{x} u_{1}, p_{2} = \partial_{x} u_{2}$.于是 \eqref{eq:c2-81} 成为,
\[
\begin{aligned}
F_{1}\bigl(x, u_{1}(x), u_{2}(x), \partial_{x} u_{1}(x), \partial_{x} u_{2}(x)\bigr) &= 0 \\
F_{2}\bigl(x, u_{1}(x), u_{2}(x), \partial_{x} u_{1}(x), \partial_{x} u_{2}(x)\bigr) &= 0.
\end{aligned}
\]

PDE 的情形要稍微复杂一些, 因为定义 $F$ 时涉及的变量数目很大.于是, 对于两个空间维数($d = 2$)下的一阶($k = 1$)标量方程($N = 1$), 我们需要依赖于两个空间变量 $x^{1}, x^{2}$ 以及 $u, p_{1} = \partial_{1} u$ 和 $p_{2} = \partial_{2} u$ 的函数 F.对于给定的五变量函数 $F = F(x, u, p)$, 两个空间维数中的一般一阶 PDE 具有如下形式,
\begin{equation}
F(x, u(x), \partial_{1} u(x), \partial_{2} u(x)) = 0. \label{eq:c2-84}
\end{equation}
作为特例取 $F = p_{1}^{2} + p_{2}^{2} - 1$.相应的方程是,
\begin{equation}
(\partial_{1} u(x))^{2} + (\partial_{2} u(x))^{2} = 1, \label{eq:c2-85}
\end{equation}
它在\newterm{geometric-optics}{几何光学}{geometric optics}中起重要作用.方程的经典解是 $C^{1}$ 函数 $u = u(x^{1}, x^{2})$, 它在区域 $D \subset \mathbb{R}^{2}$ 的所有点处满足 \eqref{eq:c2-85}.一个类似的例子是 eikonal \eqref{eq:c4-149} 给出的例子.

\noindent\textbf{注 1.}\quad 尽管超定的(即方程数目多于未知量数目)或欠定的(即方程数目少于未知量数目)方程组对几何和物理有显而易见的兴趣, 我们仍把它们排除在我们的定义之外.例如, Einstein 真空方程 \eqref{eq:c0-24} 乍一看是欠定的.一旦我们固定一个特定的坐标条件, 如引言中提到的波坐标条件, 它们就变成确定的.规范场论, 如 Yang-Mills, 具有类似的结构.另一方面, 由 $\mathbb{R}^{n}$ 的开集上的 De Rham 复形定义的方程构成一个超定方程组.例如, 给定一个一次形式 $\omega = \omega_{i} dx^{i}$, 方程组 $df = \omega$ 是超定的, 并且只有当 $\omega$ 的外微分消失, 即 $d\omega = 0$ 时才能局部求解.

\noindent\textbf{注 2.}\quad 事实上, 所有高阶标量方程或方程组都可以重新表述为一阶方程组, 即 $k = 1$, 方法很简单:把 $u$ 的所有高阶导数连同偏导数之间显然的相容关系一起作为未知量引入.作为例子, 考虑方程 \eqref{eq:c2-83} 并令 $v = \partial_{x} u$.于是我们可以把这个方程改写为 $N = 2$ 的一阶方程组, 即 $\partial_{x} v + V'(u) = 0$, $\partial_{x} u - v = 0$.

一个方程或方程组, 若关于最高阶导数是线性的, 则称为\newterm{quasi-linear}{拟线性}{quasi-linear}的.$\mathbb{R}^{d}$ 中一阶($k = 1$)拟线性方程组可以写成如下形式,
\begin{equation}
\sum_{i = 1}^{d} A^{i}(x, u(x)) \partial_{i} u = F(x, u(x)). \label{eq:c2-86}
\end{equation}
这里 $u$ 和 $F$ 是列 $N$ 维向量, 系数 $A^{1}, A^{2}, \ldots, A^{d}$ 是 $N \times N$ 矩阵值函数.

极小曲面方程是两个空间维数中二阶($k = 2$)拟线性标量方程($N = 1$)的一个例子.确实, 利用坐标 $x^{1}, x^{2}$ 而不是 x, y, 我们可以在 Leibniz 公式的帮助下处理 \eqref{eq:c0-6} 并将其改写成如下形式,
\begin{equation}
\sum_{i, j = 1, 2} h^{ij}(\partial u) \partial_{i} \partial_{j} u = 0, \label{eq:c2-87}
\end{equation}
其中
$h^{11}(\partial u) = 1 + (\partial_{2} u)^{2}$, $h^{22}(\partial u) = 1 + (\partial_{1} u)^{2}$, $h^{12}(\partial u) = h^{21}(\partial u) = -\partial_{1} u \cdot \partial_{2} u$,
这显然是二阶拟线性方程.

在拟线性方程的最高阶系数, 即对应于最高阶导数的那些系数, 只依赖于空间变量 $x \in \mathbb{R}^{d}$ 的特殊情形下, 该方程或方程组称为\newterm{semi-linear}{半线性}{semi-linear}的.例如, 在单值化定理的背景下导出的方程 \eqref{eq:c0-20} 就是半线性的.

k 阶线性方程或方程组可以写成如下形式,
\begin{equation}
\sum_{|\alpha| \leq k} A_{\alpha}(x) \partial^{\alpha} u(x) = F(x). \label{eq:c2-88}
\end{equation}
注意左边的微分算子在我们引言中所讨论的意义下确实是线性的.此外, 若系数 $A_{\alpha}$ 不依赖于 x, 则方程组称为常系数线性方程组.引言中讨论的五个基本方程 \eqref{eq:c0-1}–\eqref{eq:c0-5} 都是常系数线性的.典型地, 这些是唯一可以显式求解的方程.

这样我们就有了第一个有用的、确实是非常有用的 PDE 分类:完全非线性、拟线性、半线性和线性.\newterm{fully-nonlinear}{完全非线性}{fully nonlinear}方程关于最高阶导数是非线性的.典型的例子是 Monge-Ampère 方程.为简单起见, 考虑 $\mathbb{R}^{2}$ 中两个变量的函数 $u(x^{1}, x^{2})$, 其 Hessian 为 $\partial^{2} u = (\partial_{i} \partial_{j} u)_{i, j = 1, 2}$.显然, 行列式
$\det(\partial^{2} u) = (\partial_{1}^{2} u) \cdot (\partial_{2}^{2} u) - (\partial_{1} \partial_{2} u)^{2}$
关于 $u$ 的二阶导数是二次的.因此 Monge-Ampère 方程,
\begin{equation}
\det(\partial^{2} u) = f(x, u, \partial u), \label{eq:c2-89}
\end{equation}
其中 $f$ 是定义在 $\mathbb{R}^{2} \times \mathbb{R} \times \mathbb{R}^{2}$ 上的给定函数, 是完全非线性的.这个方程在几何中起重要作用, 它与\newterm{isometric-embedding}{等距嵌入}{isometric embedding}问题以及在 $\mathbb{R}^{2}$ 中寻找具有给定 Gauss 曲率的曲面的问题有关.Monge-Ampère 方程的一个针对复值函数的变体, 在复几何中与 Calabi-Yau 流形相关的部分起着核心作用.另一方面, Calabi-Yau 流形是弦论中核心的数学对象.

\begin{remark*}
物理学中的大多数基本方程, 如 Einstein 方程组, 都是拟线性的.然而, 完全非线性方程出现在与线性 PDE 的特征理论(我们将在下面详细讨论)或几何相关的背景中.
\end{remark*}

\section{一阶标量方程}

对这个经典课题给出系统的处理是没有意义的, 因为有许多 PDE 著作在这方面做得很好, 如 \cite{e} 或 \cite{j}.下面我只打算给出该理论背后的主要思想.事实证明, d 个空间维数中的标量($N = 1$)一阶($k = 1$)PDE 可以化归为一阶 ODE 方程组.

作为这一重要事实的简单说明, 考虑两个空间维数中的如下方程,
\begin{equation}
a^{1}(x^{1}, x^{2}) \partial_{1} u(x^{1}, x^{2}) + a^{2}(x^{1}, x^{2}) \partial_{2} u(x^{1}, x^{2}) = f(x^{1}, x^{2}), \label{eq:c2-90}
\end{equation}
其中 $a^{1}, a^{2}, f$ 是变量 $x = (x^{1}, x^{2}) \in \mathbb{R}^{2}$ 的给定实函数.我们给 \eqref{eq:c2-90} 配上如下一阶 $2 \times 2$ 方程组,
\begin{equation}
\frac{d x^{1}}{d s}(s) = a^{1}(x^{1}(s), x^{2}(s)), \quad \frac{d x^{2}}{d s} = a^{2}(x^{1}(s), x^{2}(s)). \label{eq:c2-91}
\end{equation}
为简化起见我们假设 $f = 0$.注意 \eqref{eq:c2-90}(其中 $f = 0$)的任意解 $u = u(x^{1}, x^{2})$ 沿 \eqref{eq:c2-91} 的任意解 $x(s) = (x^{1}(s), x^{2}(s))$ 是常数, 即
\[
\frac{d}{d s} u\bigl(x^{1}(s), x^{2}(s)\bigr) = 0.
\]
于是, 原则上, 知道 \eqref{eq:c2-91} 的解——称为 \eqref{eq:c2-90} 的\newterm{characteristic-curve}{特征曲线}{characteristic curve}——就能使我们找到 \eqref{eq:c2-90} 的所有解.我说"原则上", 是因为一般而言非线性方程组 \eqref{eq:c2-91} 并不那么容易求解.然而 ODE 更容易处理, 而且 ODE 基本定理——我们将在本节后面讨论——允许我们求解 \eqref{eq:c2-91}, 至少在 $s$ 的一个小区间内局部求解.u 沿特征曲线的常值性使我们能够获得重要的定性信息, 即使当我们找不到显式解时也是如此.例如, 假设系数 $a^{1}, a^{2}$ 是光滑的(或实解析的), 且初始数据在 $H$ 上处处光滑(或实解析), 只在某点 $x_{0} \in \mathcal{H}$ 处不连续.那么, 显然, 只要 \eqref{eq:c2-91} 的轨线定义良好且互不相同, 解 $u$ 就在所有点处保持光滑(或实解析), 除了沿从 $x_{0}$ 出发的特征曲线 $\Gamma$, 即沿 \eqref{eq:c2-91} 的满足初始条件 $x(0) = x_{0}$ 的解.$x_{0}$ 处的不连续性恰好沿 $\Gamma$ 传播.我们在这里看到一条一般原理最简单的表现, 我们将在后面陈述它: PDE 解的奇异性沿特征传播.

可以将方程 \eqref{eq:c2-90} 推广, 允许系数 $a^{1}, a^{2}$ 和 $f$ 不仅依赖于 $x = (x^{1}, x^{2})$, 还依赖于 u,
\begin{equation}
a^{1}(x, u(x)) \partial_{1} u(x) + a^{2}(x, u(x)) \partial_{2} u(x) = f(x, u(x)). \label{eq:c2-92}
\end{equation}
相应的特征方程组变为,
\begin{equation}
\frac{d x^{1}}{d s}(s) = a^{1}\bigl(x(s), u(s, x(s))\bigr), \quad \frac{d x^{2}}{d s} = a^{2}\bigl(x(s), u(s, x(s))\bigr). \label{eq:c2-93}
\end{equation}
于是, 与前面的情形一样,
\begin{equation}
\frac{d}{d s} u(x^{1}(s), x^{2}(s)) = f\bigl(x^{1}(s), x^{2}(s), u(x^{1}(s), x^{2}(s))\bigr). \label{eq:c2-94}
\end{equation}
然而与前面的情形不同, \eqref{eq:c2-93} 是不确定的;我们现在需要考虑扩大的 ODE 方程组 \eqref{eq:c2-93}–\eqref{eq:c2-94}, 其中未知量是 $x^{1}(s), x^{2}(s), u(s) = u(x^{1}(s), x^{2}(s))$.作为 \eqref{eq:c2-92} 的一个特殊例子, 考虑两个空间维数中的标量方程,
\begin{equation}
\partial_{t} u + u \partial_{x} u = 0, \qquad u(0, x) = u_{0}(x), \label{eq:c2-95}
\end{equation}
称为 Burger 方程.由于 $a^{1} = 1, a^{2} = u$, 我们可以在 \eqref{eq:c2-93} 中令 $x^{1}(s) = s$, $x^{2}(s) = x(s)$, 从而把其特征方程推导为如下形式,
\begin{equation}
\frac{d x}{d s}(s) = u(s, x(s)). \label{eq:c2-96}
\end{equation}
注意, 对于 \eqref{eq:c2-95} 的任意给定解 $u$ 和任意特征曲线 $(s, x(s))$, 我们有 $\frac{d}{d s} u(s, x(s)) = 0$.于是, 原则上, 知道 \eqref{eq:c2-96} 的解将允许我们确定 \eqref{eq:c2-95} 的解.然而, 这似乎是循环的, 因为 $u$ 本身出现在 \eqref{eq:c2-96} 中.为了看清如何绕开这个困难, 考虑 \eqref{eq:c2-95} 的初值问题, 即寻找满足 $u(0, x) = u_{0}(x)$ 的解 u.考虑一条相关的特征曲线 $x(s)$, 使得初始时 $x(0) = x_{0}$.那么, 由于 $u$ 沿该曲线为常数, 我们必有 $u(s, x(s)) = u_{0}(x_{0})$.因此, 回到 \eqref{eq:c2-96}, 我们推得 $\frac{d x}{d s} = u_{0}(x_{0})$, 从而 $x(s) = x_{0} + s u_{0}(x_{0})$.我们于是推得,
\begin{equation}
u(s, x_{0} + s u_{0}(x_{0})) = u_{0}(x_{0}), \label{eq:c2-97}
\end{equation}
这隐式地给出了解 $u$ 的形式.我们再一次从 \eqref{eq:c2-97} 看到, 如果初始数据在直线 $t = 0$ 上除某点 $x_{0}$ 外处处光滑(或实解析), 那么相应的解也在 $x_{0}$ 的某个小邻域 $V$ 内处处光滑(或实解析), 除了沿从 $x_{0}$ 出发的特征曲线.这里 $V$ 的小性是有必要的, 因为在整体上会形成新的奇异性.确实注意, u 必须沿直线 $x + s u_{0}(x)$ 为常数, 而这些直线的斜率依赖于 $u_{0}(x)$.在这些直线相交的点处, 我们会得到 $u$ 的不同值, 除非 $u$ 在该点变为奇异, 否则这是不可能的.事实上, 可以证明一阶导数 $u_{x}$ 在第一个奇异点, 即 $|t|$ 取最小值的那个奇异点处变为无穷.这种\newterm{blow-up}{爆破}{blow-up}现象对任何光滑、非恒定的初始数据 $u_{0}$ 都会发生.

\begin{remark*}
线性方程 \eqref{eq:c2-90} 与拟线性方程 \eqref{eq:c2-92} 之间有一个重要差别.前者的特征只依赖于系数 $a^{1}(x), a^{2}(x)$, 而后者的特征显式地依赖于方程的一个特定解 u.在这两种情形中, 奇异性都只能沿方程的特征曲线传播.然而, 对于非线性方程, 新的奇异性可以在整体上形成, 与数据的光滑性无关.
\end{remark*}

上述过程可以推广到 $\mathbb{R}^{d}$ 中如下形式的完全非线性标量方程,
\begin{equation}
\partial_{t} u + H(x, \partial u) = 0, \quad u(0, x) = u_{0}(x), \label{eq:c2-98}
\end{equation}
其中 $H = H(x, p)$ 是变量 $x = (x^{1}, x^{2}, \ldots, x^{d})$ 和 $p = (p_{1}, p_{2}, \ldots, p_{d})$ 的给定函数, 称为系统的 Hamiltonian, 而 $\partial u = (\partial_{1} u, \partial_{2} u, \ldots, \partial_{d} u)$.我们给 \eqref{eq:c2-98} 配上如下 ODE 方程组, 其中 $i = 1, 2, \ldots, d$,
\begin{equation}
\frac{d x^{i}}{d t} = \frac{\partial}{\partial p_{i}} H(x(t), p(t)), \quad \frac{d p_{i}}{d t} = -\frac{\partial}{\partial x^{i}} H(x(t), p(t)). \label{eq:c2-99}
\end{equation}
方程 \eqref{eq:c2-98} 称为\newterm{hamilton-jacobi}{Hamilton-Jacobi 方程}{Hamilton-Jacobi equation}, 而 \eqref{eq:c2-99} 称为\newterm{hamiltonian-system}{Hamiltonian 方程组}{Hamiltonian system}.它们之间的关系比前面讨论过的情形稍微复杂一些.为简化下面的计算, 我们假设 $d = 1$, 于是 $H = H(x, p)$ 只是两个变量的函数.设 $u$ 是 \eqref{eq:c2-98} 的解.对 \eqref{eq:c2-98} 关于 $x$ 求导并应用链式法则, 我们推得,
\begin{equation}
\partial_{t} \partial_{x} u + \partial_{p} H(x, \partial_{x} u) \partial_{x}^{2} u = -\partial_{x} H(x, \partial_{x} u). \label{eq:c2-100}
\end{equation}
现在取 $x(t)$ 为方程 $\frac{d x}{d t} = \partial_{p} H(x(t), \partial_{x} u(x(t)))$ 的解, 并令 $p(t) := \partial_{x} u(t, x(t))$.那么, 先使用链式法则, 再使用方程 \eqref{eq:c2-100}, 我们推得,
\[
\begin{aligned}
\frac{d p}{d t} &= \partial_{x} \partial_{t} u(t, x(t)) + \partial_{x}^{2} u(t, x(t)) \partial_{p} H(x(t), p(t)) \\
&= -\partial_{x} H(x(t), \partial_{x} u(t, x(t))) = -\partial_{x} H(x(t), p(t)).
\end{aligned}
\]
因此 $x(t), p(t)$ 满足 Hamilton 方程
\[
\frac{d x}{d t} = \partial_{p} H(x(t), p(t)), \quad \frac{d p}{d t} = -\partial_{x} H(x(t), p(t)).
\]
另一方面,
$\frac{d}{d t} u(t, x(t)) = \partial_{t} u(t, x(t)) + \partial_{x} u(t, x(t)) \partial_{p} H(x(t), p(t))$,
并且, 利用方程 \eqref{eq:c2-98}, $\partial_{t} u(t, x(t)) = -H(x(t), \partial_{x} u(t, x(t))) = -H(x(t), p(t))$.于是,
\[
\frac{d}{d t} u(t, x(t)) = -H(x(t), p(t)) + p(t) \partial_{p} H(x(t), p(t)),
\]
由此我们看到, 原则上, 如何仅基于解 $x(t), p(t)$——称为非线性 PDE 的\newterm{bicharacteristic}{双特征曲线}{bicharacteristic curve}——的知识来构造 u.再一次, 奇异性只能沿双特征传播.与 Burger 方程的情形一样, 奇异性本质上对一切光滑数据都会发生;因此经典解, 即连续可微的解, 只能在时间上局部构造.Hamilton-Jacobi 方程和 Hamiltonian 方程组在经典力学以及线性 PDE 中奇异性的传播理论中都起着基本作用.Hamiltonian 方程组与一阶 Hamilton-Jacobi 方程之间的深刻联系, 在量子力学中 Schrödinger 方程的引入中起过重要作用.

\section{ODE 的初值问题}

要进一步展开我们的一般论述, 需要讨论初值问题.为简单起见, 我们从一个一阶 ODE 开始,
\begin{equation}
\partial_{x} u(x) = f(x, u(x)), \label{eq:c2-101}
\end{equation}
并附带初始条件
\begin{equation}
u(x_{0}) = u_{0}. \label{eq:c2-102}
\end{equation}
读者不妨假设, 为简单起见, \eqref{eq:c2-101} 是标量方程, 且 $f$ 是 $x$ 和 $u$ 的性质良好的函数, 如 $f(x, u) = u^{3} - u + 1 + \sin x$.注意, 知道初始数据 $u_{0}$ 就能使我们确定 $\partial_{x} u(x_{0})$.对方程 \eqref{eq:c2-101} 关于 $x$ 求导并应用链式法则, 我们推得,
\[
\partial_{x}^{2} u(x) = \partial_{x} f(x, u(x)) + \partial_{u} f(x, u(x)) \partial_{x} u(x) = \cos x + 3 u^{2}(x) \partial_{x} u(x) - \partial_{x} u(x).
\]
因此, $\partial_{x}^{2} u(x_{0}) = \partial_{x} f(x_{0}, u_{0}) + \partial_{u} f(x_{0}, u_{0}) \partial_{x} u_{0}$, 而由于 $\partial_{x} u(x_{0})$ 已经确定, 我们推得 $\partial_{x}^{2} u(x_{0})$ 可以由初始数据 $u_{0}$ 显式算出.计算还涉及函数 $f$ 及其一阶偏导数.对方程 \eqref{eq:c2-101} 取更高阶导数, 我们可以递归地确定 $\partial_{x}^{3} u(x_{0})$ 以及 $u$ 在 $x_{0}$ 处的所有其他高阶导数.于是, 原则上, 可以借助 Taylor 级数
$u(x) = \sum_{k \geq 0} \frac{1}{k!} \partial_{x}^{k} u(x_{0}) (x - x_{0})^{k} = u(x_{0}) + \partial_{x} u(x_{0}) (x - x_{0}) + \frac{1}{2!} \partial_{x}^{2} u(x_{0}) (x - x_{0})^{2} + \cdots$
来确定 $u(x)$.我们说"原则上", 是因为不能保证级数收敛.然而有一个非常重要的定理, 称为\newterm{cauchy-kowalewski}{Cauchy-Kowalewski 定理}{Cauchy-Kowalewski theorem}, 它断言:若函数 $f$ 是实解析的——对我们的 $f(x, u) = u^{3} - u + 1 + \sin x$ 当然如此——那么存在 $x_{0}$ 的一个邻域 J, 在其中 Taylor 级数收敛到方程的一个实解析解 u.然后可以容易地证明, 这样得到的解是满足初始条件 \eqref{eq:c2-102} 的方程 \eqref{eq:c2-101} 的唯一解.

如果我们考虑更一般形式的方程, 同样的结论不一定成立,
\begin{equation}
a(x, u(x)) \partial_{x} u = f(x, u(x)), \qquad u(x_{0}) = u_{0}. \label{eq:c2-103}
\end{equation}
事实上, 上述递归论证在标量方程 $(x - x_{0}) \partial_{x} u = f(x, u)$ 的情形中失效, 原因很简单:我们甚至无法从初始条件 $u(x_{0}) = u_{0}$ 确定 $\partial_{x} u(x_{0})$.方程 $(u - u_{0}) \partial_{x} u = f(x, u)$ 也出现类似的问题.允许我们把前面的递归论证推广到 \eqref{eq:c2-103} 的一个显然条件是 $a(x_{0}, u_{0}) \neq 0$.否则我们称初值问题 \eqref{eq:c2-103} 是特征的.若 $a$ 和 $f$ 也都是实解析的, 则 Cauchy-Kowalewski 定理适用, 我们得到 \eqref{eq:c2-103} 在 $x_{0}$ 的某个小邻域内的唯一实解析解.对于 $N \times N$ 方程组,
\begin{equation}
A(x, u(x)) \partial_{x} u = F(x, u(x)), \qquad u(x_{0}) = u_{0}, \label{eq:c2-104}
\end{equation}
其中 $A = A(x, u)$ 是 $N \times N$ 矩阵, 非特征条件变为
\begin{equation}
\det A(x_{0}, u_{0}) \neq 0. \label{eq:c2-105}
\end{equation}
事实证明——这一点极为重要——尽管非退化条件 \eqref{eq:c2-105} 对获得方程的唯一解是本质的, 但在 ODE 的情形中, 解析性条件一点也不重要.它可以被 $A$ 和 $F$ 的简单局部 Lipschitz 条件所取代, 即只需假设, 例如, 它们的一阶偏导数存在并且只是局部有界的.如果 $A$、$F$ 的一阶导数连续, 这一点总成立.

下面的局部存在唯一性 (LEU) 定理称为 ODE 基本定理.

\begin{theorem*}[常微分方程基本定理]
若矩阵 $A(x_{0}, u_{0})$ 可逆, 且 $A$、$F$ 连续并具有局部有界的一阶导数, 那么存在一个时间区间 $x_{0} \in J \subset \mathbb{R}$ 和在 $J$ 上定义的唯一解\footnote{由于我们不对 $A$、$F$ 假定解析性, 解可能不解析, 但它具有连续的一阶导数.} u, 满足初始条件 $u(x_{0}) = u_{0}$.
\end{theorem*}

\begin{proof}
定理的证明基于\newterm{picard-iteration}{Picard 迭代}{Picard iteration}法.其思想是构造一列逼近解 $u_{(n)}(x)$, 它们收敛到所求的解.不失一般性, 我们可以假设 $A$ 是单位矩阵\footnote{由于 $A$ 可逆, 我们可以用逆矩阵 $A^{-1}$ 乘方程两边.}.首先令 $u_{(0)}(x) = u_{0}$, 然后递归地定义,
\begin{equation}
\partial_{x} u_{(n)}(x) = F(x, u_{(n-1)}(x)), \qquad u_{(n-1)}(x_{0}) = u_{0}. \label{eq:c2-106}
\end{equation}
注意在每一步我们只需要求解一个非常简单的线性问题, 这使得 Picard 迭代在数值上易于实现.我们将在下面看到, 该方法的变体也用于求解非线性 PDE.

\ldots\ldots{} 填补证明\ldots\ldots
\end{proof}

\begin{remark*}
局部存在性定理一般而言是精确的.事实上我们已经看到, $A(x_{0}, u_{0})$ 的可逆性条件是必要的.而且, 一般而言, 存在区间 $J$ 不能延拓到整条实直线.作为例子, 考虑非线性方程 $\partial_{x} u = u^{2}$, 初始数据为 $x = 0$ 处 $u = u_{0}$, 其解 $u = \frac{u_{0}}{1 - x u_{0}}$ 在有限时间内变为无穷, 即发生爆破.
\end{remark*}

一旦建立了 LEU 结果, 就可以把 ODE 数学理论的主要目标规定为:
\begin{enumerate}
\item 寻找全局存在的判据.在爆破的情形, 描述极限行为.
\item 在全局存在的情形, 描述解和解族的渐近行为.
\end{enumerate}
虽然要发展一个同时回答这两个目标的一般理论是不可能的(在实践中, 人们不得不局限于由应用所启发的特殊方程类), 上述一般 LEU 定理给出了一个有力的统一主题.如果对一般 PDE 也有类似的情形成立, 那将是极有帮助的, 确实非常美妙.

\section{PDE 的初值问题}

与一维情形类比, 自然要考虑的不是点, 而是用于规定 $u$ 的初始条件的超曲面 $\mathcal{H} \subset \mathbb{R}^{d}$.对于 $k$ 阶一般方程, 即涉及 $k$ 个导数的方程, 我们需要规定 $u$ 及其对 $H$ 的前 $k - 1$ 个法向导数\footnote{即沿 $H$ 的法线方向的导数.}的值.例如, 对于二阶波动方程 \eqref{eq:c0-3}, 我们需要在超曲面 $t = 0$ 上规定 $u$ 和 $\partial_{t} u$ 的初始数据.在这一点上不深入细节, 我们可以给出如下一般定义.

\begin{definition*}
我们说, 一个 $k$ 阶拟线性方程组的初值问题——其中我们把解 $u$ 沿 $H$ 的前 $k - 1$ 个法向导数作为数据加以规定——在 $H$ 的一点 $x_{0}$ 处是\newterm{non-characteristic}{非特征的}{non-characteristic}, 如果我们可以用这些数据形式地、唯一地确定 $u$ 在 $x_{0}$ 处的所有其他高阶偏导数.
\end{definition*}

为了理解这个定义(在这一点上它可能显得过于一般), 考虑简单得多的情形 $k = 1, N = 1$.此时我们只需要规定 $u$ 在 $H$ 上的限制 $u|_{\mathcal{H}} = u_{0}$.我们的初值问题具有如下形式,
\begin{equation}
\sum_{i = 1}^{d} a^{i}(x, u(x)) \partial_{i} u(x) = f(x, u(x)), \quad u|_{\mathcal{H}} = u_{0}, \label{eq:c2-107}
\end{equation}
其中 $a^{i}, f$ 是 $x \in \mathbb{R}^{d}$ 和 $u \in \mathbb{R}$ 的实值函数.为进一步简化, 取 $d = 2$, 即我们有关于 $x = (x^{1}, x^{2})$ 的方程,
\begin{equation}
a^{1}(x, u(x)) \partial_{1} u(x) + a^{2}(x, u(x)) \partial_{2} u(x) = f(x, u(x)), \label{eq:c2-108}
\end{equation}
即我们早先在 \eqref{eq:c2-92} 中遇到的方程.考虑 $\mathbb{R}^{2}$ 中的一条曲线 H, 参数化为 $x^{1} = x^{1}(s), x^{2} = x^{2}(s)$, 其切向量 $V(s) = (\frac{d x^{1}}{d s}, \frac{d x^{2}}{d s})$ 非退化, 即 $|V(s)| = (|\frac{d x^{1}}{d s}|^{2} + |\frac{d x^{2}}{d s}|^{2})^{1/2} \neq 0$.它有定义良好的单位法向量 $N(s) = (n_{1}(s), n_{2}(s))$, 满足条件,
\[
N(s) \cdot V(s) = 0, \qquad N(s) \cdot N(s) = 1.
\]
注意, \eqref{eq:c2-108} 中的系数 $a^{1}, a^{2}$ 可以沿着 $H$ 由初始条件 $u_{0} = u_{0}(s)$ 的知识完全确定.考虑沿 $H$ 取值的一阶导数 $(\partial_{1} u, \partial_{2} u)$, 即 $U(s) = \bigl(\partial_{1} u(x(s)), \partial_{2} u(x(s))\bigr)$.在 $H$ 的每一点, 我们的方程写作,
\begin{equation}
A(s) \cdot U(s) = f(s), \label{eq:c2-109}
\end{equation}
其中 $A(s) = \bigl(a^{1}(x(s), u_{0}(s)), a^{2}(x(s), u_{0}(s))\bigr)$ 和 $f(s) = f\bigl(x(s), u_{0}(s)\bigr)$ 完全由数据 $u_{0}(s)$ 确定.对 $u(x(s)) = u_{0}(s)$ 关于 $s$ 求导, 我们推得,
\[
U(s) \cdot V(s) = U_{0}(s), \qquad U_{0}(s) = \frac{d}{d s} u_{0}(s).
\]
要完全确定 $U(s)$, 剩下的问题是确定它在法向量 $N(s)$ 上的投影, 即 $U(s) \cdot N(s)$.事实上, 由于 $V(x)$ 和 $N(x)$ 张成 $\mathbb{R}^{2}$, 在我们曲线的所有点 $x = (x^{1}(s), x^{2}(s))$ 处, 我们有
\begin{equation}
U(s) = (U \cdot V)(s) \frac{V(s)}{|V(s)|^{2}} + (U \cdot N)(s) N(s). \label{eq:c2-110}
\end{equation}
因此, 由方程 \eqref{eq:c2-109},
\[
f(s) = A(s) \cdot U(s) = (U(s) \cdot V(s)) \frac{A(s) \cdot V(s)}{|V(s)|^{2}} + (U(s) \cdot N(s)) A(s) \cdot N(s),
\]
由此我们可以确定 $U(s) \cdot N(s)$, 只要
\begin{equation}
A(s) \cdot N(s) \neq 0. \label{eq:c2-111}
\end{equation}
另一方面, 若 $A(s) \cdot N(s) = 0$, 那么由于 $V(s) \cdot N(s) = 0$, 我们推得向量 $A(s)$ 和 $V(s) = \frac{d x}{d s}$ 必须成比例, 即 $\frac{d x}{d s} = \lambda(s) A(s)$.于是可以重新参数化曲线 $\mathcal{H}$, 即引入另一个参数 $s' = s'(s)$, 满足 $\frac{d s'}{d s} = \lambda(s)$, 使得相对于新参数有 $\lambda = 1$.这导致方程,
\[
\frac{d x^{1}}{d s} = a^{1}\bigl(x(s), u(x(s))\bigr), \quad \frac{d x^{2}}{d s} = a^{2}\bigl(x(s), u(x(s))\bigr),
\]
这正是特征方程组 \eqref{eq:c2-93}.于是,

沿特征曲线, 方程 \eqref{eq:c2-108} 是退化的, 即我们不能用数据 $u_{0}$ 唯一确定 $u$ 的一阶导数.另一方面, 非退化条件,
\begin{equation}
A(s_{0}) \cdot N(s_{0}) \neq 0, \quad \text{即} \quad a^{1}(x_{0}, u(x_{0})) n_{1}(x_{0}) + a^{2}(x_{0}, u(x_{0})) n_{2}(x_{0}) \neq 0, \label{eq:c2-112}
\end{equation}
在某个点 $x_{0} = x(s_{0}) \in \mathcal{H}$ 处, 允许我们用数据 $u_{0}$ 唯一确定 $u$ 在 $x_{0}$ 处的所有高阶导数.

事实上, 若条件 $A(s_{0}) \cdot N(s_{0}) \neq 0$ 在 $x_{0} = x(s_{0}) \in \mathcal{H}$ 处满足, 我们已经看到如何确定该点处的一阶导数 $\partial_{1} u, \partial_{2} u$.一旦有了这些, 确定 $u$ 的所有高阶导数并不困难.例如, 注意, 对方程 \eqref{eq:c2-108} 关于 $x^{1}$ 求导, 函数 $v = \partial_{1} u$ 满足如下形式的方程,
\[
a^{1}(x, u(x)) \partial_{1} v(x) + a^{2}(x, u(x)) \partial_{2} v = g(x, u(x), v(x)),
\]
其中函数 $g$ 可以容易地从系数 $a$ 和 $f$ 确定.我们现在可以像前面一样进行, 确定 $v$ 的一阶导数, 即 $\partial_{1}^{2} u, \partial_{2} \partial_{1} u$.于是, 递归地, 我们可以确定 $u$ 的任意阶的所有偏导数.

我们可以容易地把上面的讨论推广到高维情形 \eqref{eq:c2-107}.给定 $\mathbb{R}^{d}$ 中的超曲面 H, 单位法向量为 $N = (n_{1}, n_{2}, \ldots, n_{d})$, 我们发现, H 在 $x_{0}$ 处对初值问题 \eqref{eq:c2-107} 是非特征的, 若
\begin{equation}
\sum_{i = 1}^{d} a^{i}(x_{0}, u_{0}(x_{0})) n_{i}(x_{0}) \neq 0. \label{eq:c2-113}
\end{equation}

再多做一些工作, 我们可以把讨论推广到一般高阶拟线性方程或方程组, 得到 Cauchy 问题为非特征的简单充分条件.对我们特别重要的是二阶($k = 2$)标量方程($N = 1$).为简单起见, 考虑 $\mathbb{R}^{d}$ 中一个二阶半线性方程,
\begin{equation}
\sum_{i, j = 1}^{d} a^{ij}(x) \partial_{i} \partial_{j} u = f(x, u(x), \partial u(x)), \label{eq:c2-114}
\end{equation}
以及由方程 $\psi(x) = 0$ 定义、梯度 $\partial \psi$ 不消失的超曲面 H.定义点 $x_{0} \in \mathcal{H}$ 处的单位法向量为 $N = \frac{\partial \psi}{|\partial \psi|}$, 或按分量 $n_{i} = \frac{\partial_{i} \psi}{|\partial \psi|}$.作为 \eqref{eq:c2-114} 的初始条件, 我们在 $H$ 上规定 $u$ 及其法向导数 $N u(x) \doteq n_{1}(x) \partial_{1} u(x) + n_{2}(x) \partial_{2} u(x) + \cdots + n_{d}(x) \partial_{d} u(x)$,
\begin{equation}
u(x) = u_{0}(x), \quad N u(x) = u_{1}(x), \qquad x \in \mathcal{H}. \label{eq:c2-115}
\end{equation}
我们需要找到 $H$ 上的一个条件, 使得我们可以从初始数据 $u_{0}, u_{1}$ 确定解 $u$ 在 $x_{0} \in \mathcal{H}$ 处的所有高阶导数.我们可以完全像前面一样进行, 发现只要
\begin{equation}
\sum_{i, j = 1}^{d} a^{ij}(x_{0}) n_{i}(x_{0}) n_{j}(x_{0}) \neq 0, \label{eq:c2-116}
\end{equation}
u 的所有二阶导数都可以在点 $x_{0} \in \mathcal{H}$ 处确定.确实容易看到, u 的唯一不能自动由 $u_{0}, u_{1}$ 确定的二阶导数是形如 $N^{2} u(x_{0}) = N(N(u))(x_{0})$ 的那一个.后者可以在 \eqref{eq:c2-116} 得到验证的前提下由方程 \eqref{eq:c2-114} 确定.做法是把 $u$ 的所有偏导数分解为切向分量和法向分量, 正如我们在 \eqref{eq:c2-110} 中所做的那样.然后可以递归地证明, u 的所有高阶导数也都可以确定.于是, \eqref{eq:c2-116} 正是我们所要寻找的非特征条件.

另一方面, 若在所有点处 $\sum_{i, j = 1}^{d} a^{ij}(x) n_{i}(x) n_{j}(x) = 0$, 我们称 $H$ 是方程 \eqref{eq:c2-114} 的\newterm{characteristic-hypersurface}{特征超曲面}{characteristic hypersurface}.由于 $n_{i} = \frac{\partial_{i} \psi}{|\partial \psi|}$, 我们发现 $H$ 是特征的当且仅当,
\begin{equation}
\sum_{i, j = 1}^{d} a^{ij}(x) \partial_{i} \psi(x) \partial_{j} \psi(x) = 0. \label{eq:c2-117}
\end{equation}

\begin{remark*}
注意, 在确定方程的特征曲面时, 只有 \eqref{eq:c2-114} 的左边\footnote{即含二阶导数的部分}, 即所谓主部\footnote{主部}, 才是相关的.
\end{remark*}

\noindent\textbf{例 1.}\quad 假设 \eqref{eq:c2-114} 的系数 $a$ 满足,
\begin{equation}
\sum_{i, j = 1}^{d} a^{ij}(x) \xi_{i} \xi_{j} > 0, \quad \forall \xi \in \mathbb{R}^{d}, \quad \forall x \in \mathbb{R}^{d}. \label{eq:c2-118}
\end{equation}
那么 $\mathbb{R}^{d}$ 中没有曲面可以是特征的.特别是方程 $\Delta u = f$ 就是这种情况.再考虑写成形式 \eqref{eq:c2-87} 的极小曲面方程.容易验证, 与对称矩阵 $h^{ij}(\partial u)$ 相关的二次型是正定的, 与 $\partial u$ 无关.确实,
\[
h^{ij}(\partial u) \xi_{i} \xi_{j} = (1 + |\partial u|^{2})^{-1/2} \bigl(|\xi|^{2} - (1 + |\partial u|^{2})^{-1} (\xi \cdot \partial u)^{2}\bigr) > 0.
\]
于是, 即使 \eqref{eq:c2-87} 不是线性的, 我们也看到 $\mathbb{R}^{2}$ 中的所有曲面都是非特征的.

\noindent\textbf{例 2.}\quad 考虑 $\mathbb{R}^{1+d}$ 中的波动方程 $\Box u = f$.所有使得
\begin{equation}
(\partial_{t} \psi)^{2} = \sum_{i = 1}^{d} (\partial_{i} \psi)^{2}, \label{eq:c2-119}
\end{equation}
的超曲面 $\psi(t, x) = 0$ 都是特征的.这正是前面在 (149) 中出现过的同一个\newterm{eikonal}{eikonal 方程}{eikonal equation}.注意它分裂成两个 Hamilton-Jacobi 方程, 见 \eqref{eq:c2-98},
\begin{equation}
\partial_{t} \psi = \pm \Bigl(\sum_{i = 1}^{d} (\partial_{i} \psi)^{2}\Bigr)^{1/2}. \label{eq:c2-120}
\end{equation}
相应 Hamiltonian 的双特征曲线称为波动方程的双特征曲线.作为 (429) 的特解, 我们找到 $\psi_{+}(t, x) = (t - t_{0}) + |x - x_{0}|$ 和 $\psi_{-}(t, x) = (t - t_{0}) - |x - x_{0}|$, 它们的水平曲面 $\psi_{\pm} = 0$ 对应于顶点在 $p = (t_{0}, x_{0})$ 的前向和后向光锥.在物理上, 这些表示从 $p$ 处的点源发出的所有光线的并集.光线由方程 $(t - t_{0}) \omega = (x - x_{0})$ 给出, 其中 $\omega \in \mathbb{R}^{3}$ 且 $|\omega| = 1$, 它们恰好是 Hamilton-Jacobi 方程 \eqref{eq:c2-120} 的双特征曲线的 (t, x) 分量.

更一般地, 考虑线性波动方程,
\begin{equation}
g^{\alpha \beta} \partial_{\alpha} \partial_{\beta} \phi = 0, \label{eq:c2-121}
\end{equation}
其中 $g^{\alpha \beta}$ 是一般 Lorentz 度量 $g_{\alpha \beta}$\footnote{即 $g_{\alpha \beta} = g_{\beta \alpha}$, 且相应的二次型 $g_{\alpha \beta}(p) X^{\alpha} X^{\beta}$ 在每一点 $p$ 非退化, 并有\newterm{signature}{符号差}{signature} $(-1, 1, \ldots, 1)$.} 的逆.\eqref{eq:c2-121} 的特征曲面由 eikonal 方程给出,
\begin{equation}
g^{\alpha \beta} \partial_{\alpha} \psi \partial_{\beta} \psi = 0. \label{eq:c2-122}
\end{equation}
它们也称为度量 $g$ 的零超曲面.

\begin{remark*}
在 $g_{01} = \cdots = g_{0n} = 0$, $g_{00} < 0$ 且 $g_{ij}$ 正定的特殊情形下, \eqref{eq:c2-121} 具有形式,
\begin{equation}
-a^{00}(t, x) \partial_{t}^{2} u + \sum_{i, j} a^{ij}(t, x) \partial_{i} \partial_{j} u = 0, \label{eq:c2-123}
\end{equation}
其中 $a_{00} = -g_{00}^{-1}$, $a^{ij} g_{jk} = \delta_{k}^{i}$.于是特征方程具有形式 $-a^{00}(t, x) (\partial_{t} \psi)^{2} + a^{ij}(x) \partial_{i} \psi \partial_{j} \psi = 0$, 或
\begin{equation}
\partial_{t} \psi = \pm \Bigl((a^{00})^{-1} \sum_{i, j} a^{ij}(x) \partial_{i} \psi \partial_{j} \psi\Bigr)^{1/2}. \label{eq:c2-124}
\end{equation}
于是, 通过 $\mathbb{R}^{1+n}$ 中的任何点 $p$ 都有两个不同的特征曲面经过.一般情形也是如此.
\end{remark*}

相应 Hamiltonian 方程组的双特征称为 \eqref{eq:c2-123} 的双特征曲线.

\begin{remark*}
对于一阶标量方程 \eqref{eq:c2-90}, 我们已经看到如何利用特征的知识隐式地找到一般解.我们特别证明了奇异性只沿特征传播.对于二阶方程, 情形更复杂.特征通常\footnote{然而, 特征会以显式形式出现在标准波方程基本解的表达式里.这在从假设 (148) 出发的推导中尤为明显.它们在构造变系数波方程(如 \eqref{eq:c2-123})的近似解方面也起着重要作用.}不足以求解方程, 但它们继续提供重要的信息, 如奇异性的传播.例如, 对于波动方程 $\Box u = 0$, 初始数据 $u_{0}$、$u_{1}$ 除一点 $p = (t_{0}, x_{0})$ 外处处光滑, 解 $u$ 在顶点为 $p$ 的光锥 $-(t - t_{0})^{2} + |x - x_{0}|^{2} = 0$ 的所有点处都有奇异性.这一事实的一个更精细的版本表明, 奇异性沿双特征传播.这里的一般原理是, 奇异性沿 PDE 的特征超曲面传播.由于这是非常重要的原理, 值得给出一个更精确的表述, 它推广到一般的边界条件, 如 (??) 的 Dirichlet 条件.
\end{remark*}

\noindent\textbf{奇异性传播}\footnote{该原理的一个更精确的版本把奇异性的传播与双特征曲线联系起来.}\quad 若具有光滑(或实解析)系数的线性 PDE 的边界条件或系数在某个点 $p$ 处是奇异的, 而在 $p$ 之外某个小邻域 $V$ 内光滑(或实解析), 那么方程的解在 $V$ 中只能沿经过 $p$ 的特征超曲面奇异.若不存在这样的特征超曲面, 则方程的任何解在 $V \setminus \{p\}$ 中必为光滑(或实解析)的.

\noindent\textbf{注 1.}\quad 如上陈述的原理过于一般, 只有在算子的象征上作出具体假设时才能证明.然而, 它应被视为对合理方程可以预期的东西.

\noindent\textbf{注 2.}\quad 在对解作出特定的最小正则性假设下, 该原理可以推广到非线性情形.然而它在整体上是无效的.事实上, 正如我们在 Burger 方程的情形中所示, 非线性发展方程的解可以发展出与初始条件的光滑性无关的新奇异性.对于线性方程, 可以基于方程的双特征陈述该原理的全局版本, 见下文注 3.

\noindent\textbf{注 3.}\quad 根据该原理, 方程 $\Delta u = f$ 的任何满足边界条件 $u|_{\partial D} = u_{0}$ 的解——其中边界值 $u_{0}$ 只是连续的——只要 $f$ 本身在 $D$ 内部光滑, 就必然在 $D$ 内部处处光滑.而且, 若 $f$ 是实解析的, 则解是实解析的.

\noindent\textbf{注 4.}\quad 这一在一般理论中起基本作用的原理, 对线性方程可以给出更精确的版本.例如, 对于一般波动方程 \eqref{eq:c2-123}, 可以证明奇异性沿双特征传播.这些是与 Hamilton-Jacobi 方程 \eqref{eq:c2-124} 相关的双特征曲线.

\section{Cauchy-Kowalevsky 定理}

在 ODE 的情形中, 我们已经看到非特征初值问题总是有局部(在时间上)解.这一事实有没有高维的类比?答案是肯定的, 只要我们把自己限制在 Cauchy-Kowalewski 定理的一个推广上.更精确地说, 可以考虑具有实解析系数、实解析超曲面 $\mathcal{H}$ 以及 $H$ 上实解析初始数据的一般拟线性方程或方程组.

\begin{theorem*}[Cauchy-Kowalevsky 定理]
若上面作出的所有实解析性条件都满足, 且 $H$ 在 $x_{0}$\footnote{对于 \eqref{eq:c2-114} 型二阶方程, 这恰是条件 \eqref{eq:c2-116}.} 处是非特征的, 那么在 $x_{0}$ 的一个邻域内局部存在唯一的实解析解 u(x), 满足该方程组及相应的初始条件.
\end{theorem*}

CK 定理证实了通过形式展开
$u(x) = \sum_{\alpha} C_{\alpha} (x - x_{0})^{\alpha}$
(其中常数 $C_{\alpha}$ 可以由方程和 $\mathcal{H}$ 上的初始条件, 仅利用非特征条件和解析性假设, 通过简单的代数公式递归确定)寻找解的最直接的尝试.事实上, 该定理保证了以这种方式得到的朴素展开在 $x_{0} \in \mathcal{H}$ 的某个小邻域内收敛.

\begin{proof}
见 \cite{e} 或 \cite{j}.
\end{proof}

在特殊情形——线性方程 \eqref{eq:c2-88}——中, 一个重要的配套定理, 归功于 Holmgren, 断言 CK 定理给出的解析解在所有光滑解和光滑非特征超曲面 $H$ 组成的类中是唯一的.

\begin{theorem}[Holmgren 唯一性定理]\label{thm:c2-4.1}
考虑具有解析系数的 \eqref{eq:c2-88} 型线性方程的初值问题.若超曲面 $H$ 也是解析的, 且在 $x_{0} \in \mathcal{H}$ 处非特征, 那么相应的 Cauchy 问题在 $x_{0}$ 的某个小邻域内的光滑解类中是唯一的.
\end{theorem}

\begin{proof}
见 \cite{j}.
\end{proof}

\begin{remark*}
Holmgren 定理的引人注目之处在于, 即使在不能保证解存在的情形下, 它也证明唯一性.于是, 我们将在下面看到, 在超平面 $x^{1} = 0$ 上给定数据的波动方程 Cauchy 问题一般而言没有解, 然而 Holmgren 定理断言, 若解存在, 则它必唯一.
\end{remark*}

乍一看, CK 定理似乎是 ODE 基本定理的完美类比.然而事实证明, CK 定理所要求的解析性条件过于严格, 因此该结果表面的普遍性是误导性的.第一个局限性在我们考虑波动方程 $\Box u = 0$ 时立刻变得明显, 其有限传播速度的基本特征\footnote{粗略地说, 这意味着若解 $u$ 在某个 $t$ 值处具有紧支集, 则它在所有更晚的时刻也都具有紧支集.解析函数若不全恒为零, 则不可能具有紧支集.}在实解析解类中无法理解.一个相关的问题, 最早由 Hadamard 指出, 涉及在许多重要情形中, 对任意光滑、非解析的数据求解 Cauchy 问题的不可行性.例如, 考虑 $\mathbb{R}^{d}$ 中的 Laplace 方程 $\Delta u = 0$.正如我们上面所确立的, 任何超曲面 $H$ 都是非特征的, 然而对任意光滑初始条件 $u_{0}, u_{1}$, Cauchy 问题 $u|_{\mathcal{H}} = u_{0}$, $N(u)|_{\mathcal{H}} = u_{1}$ 在 $H$ 的任何一点的邻域内都可能没有局部解.事实上, 取 $H$ 为超平面 $x_{1} = 0$, 并假设对于一个给定的、非解析的光滑数据, Cauchy 问题在一个包含以原点为中心的闭球 $B$ 的区域中可以求解.相应的解也可以解释为 $B$ 中的 Dirichlet 问题的解, 其中 $u$ 的值在边界 $\partial B$ 上给定.但根据我们的启发式原理\footnote{在此情形下这一点可以很容易地严格化.}, 后者在 $B$ 的内部处处实解析, 这与我们的初始数据假设矛盾.

另一方面, 波动方程 $\Box u = 0$ 在 $\mathbb{R}^{d+1}$ 中的 Cauchy 问题, 对在类空超曲面上给定的任何光滑初始数据 $u_{0}, u_{1}$ 都有唯一解, 这里类空超曲面是指这样的超曲面 $\psi(t, x) = 0$:它在每一点 $p = (t_{0}, x_{0})$ 处的法向量都指向经过该点的前向或后向光锥的内部.形式上这意味着,
\begin{equation}
|\partial_{t} \psi(p)| > \left(\sum_{i = 1}^{d} |\partial_{i} \psi(p)|^{2}\right)^{1/2}. \label{eq:c2-125}
\end{equation}
$t = t_{0}$ 的超曲面显然满足该条件, 但任何与之接近的其他超曲面也是类空的.另一方面, 对于类时超曲面, 即满足
\begin{equation}
|\partial_{t} \psi(p)| < \left(\sum_{i = 1}^{d} |\partial_{i} \psi(p)|^{2}\right)^{1/2} \label{eq:c2-126}
\end{equation}
的超曲面, IVP 是病态的, 即不适定的.在这种情况下, 对一般的非实解析初始条件, 我们无法找到 IVP 的解.类时超曲面的一个例子由超平面 $x^{1} = 0$ 给出.

\begin{definition*}
若对属于某个指定的、包含光滑函数类\footnote{这里我们必然有些含糊.在每个具体情形下可以指定一个精确的空间.}的大函数空间的任意数据, 都能建立解的存在性和唯一性, 则称 PDE 的一个给定问题是\newterm{well-posed}{适定的}{well posed}的.此外, 解必须连续依赖于数据.
\end{definition*}

数据的连续依赖性非常重要.事实上, 如果初始条件的极小改变会立即导致相应解的极大改变, 那么 PDE 的 IVP 的解就没有什么用处.只有在光滑解类中, PDE 理论才变得真正有趣、相关和富有挑战性.这意味着我们必须放弃包罗万象的结果的希望, 转而寻找具有共同特征的方程的特殊类, 或者干脆专注于特殊的重要方程.正是在这个意义上, CK 定理的普遍性确实是一种幻觉.偏微分方程的真正研究只有在我们放弃解析性时才开始.在下一章中, 我们将详细分析最简单方程——如 Cauchy-Riemann、Laplace、热方程和波动方程——利用它们相应的基本解的主要解析性质.为此, 我们首先需要回顾分布理论.
% ==== part03 第 3 章 分布理论 (Distribution Theory) ====
% 源: MinerU_latex_lecture nots.tex 第 3675--5210 行;pdf_text.txt 第 2789--4042 行(PDF 第 55--76 页)
%
% 术语(新增,键=中文=英文):
%   distribution=分布、indicator-function=示性函数、fundamental-solution=基本解、
%   test-function=检验函数、mollification=磨光、multi-index=多重指标、support=支集、
%   compactly-supported=紧支撑、locally-integrable=局部可积、convolution=卷积、
%   bump-function=隆起函数、cutoff-function=截断函数、principal-value=主值、
%   analytic-continuation=解析延拓、partition-of-unity=单位分解、pull-back=回拉、
%   push-forward=前推、homogeneous=齐次的、characteristic-function=特征函数、
%   restriction-theorem=限制定理、trace-theorem=迹定理、wavefront-set=波前集、
%   cotangent-bundle=余切丛、tensor-product=张量积
% 注: 傅里叶变换、波动方程、Laplace 算子、Dirac delta 函数、Heaviside 函数、Gamma/Beta 函数、
%   Borel 测度、Lebesgue 积分等经典术语按普通中文书写(可能由其他部分首次引入,待协调).
%
% REFS(需跨章/待解析项):
%   - 公式编号: 原书第 3 章显式编号的公式为 (127)--(136) 共 10 个,均加了 \label{eq:c3-12x} 等;
%     本部分按指南把全部行间公式写成 \[ ... \] 自动编号(原书其余行间公式不编号),共 81 个,
%     故显示编号与原书 (127)--(136) 不对应;任务描述"~176--~230"与实际 78 个不符,请协调员核对生产环境编号规则.
%   - 第 4407 行 "the Lebesgue diferentiation theorem of section (2.6)" -> 正文写"第 2.6 节的 Lebesgue 微分定理",所在章未知 -> 目标: 待解析
%   - 第 4468 行 "Hörmander, [?] section 3.4" 文献编号 [?] 原书破损 -> 正文写"Hörmander 著作\cite{ho}第 3.4 节"
%   - 第 5076 行 "Hörmander's book [?]" 同上
%   - 第 5000 行 "the Appendix (??)" 与第 4964 行 "given in the Appendix" -> 正文写"附录",页码/引用待解析(附录 A 部分)
%   - 第 4705、5188 行 "(∗ ∗ ∗)" 为原书占位引用,正文按原文保留 "(∗∗∗)"
%   - 脚注: 本部分含两处原书脚注,均已译入正文 \footnote:
%     ① 第 4066 行 "non-metrizable1" 处;② 第 4866 行 "follow2" 处.
%     任务提及的脚注 6(Laplace-Beltrami 算子,pdf_text.txt 第 4460 行)的标记位于第 4 章(part04 范围),不在本部分.
%
% 编号说明:
%   - 原书 Proposition 3.8 为 "iff"(pdf_text 第 3492 行),源 OCR 误为 "if",按 pdf 译"当且仅当".
%   - 原书 "a fundamental solution solution for Δ"(pdf_text 第 2902 行)重复拼写,按"基本解"处理.
%   - 引言的未编号 Definition 在源文件中位于 \section{1} 之内(第 3787 行),故置于第 1 节;
%     任务结构表将其列在"引言"下,此处按源结构处理,用 definition* 无编号环境.
%   - 第 2.5 节 "Operations with distributions"、第 3.5 节 "Support of a distribution"、
%     第 4.2 节 "Other topics we have not discussed" 均为原书带编号子标题,用 \mtsubsec 推进共享计数器;
%     第 4 节内 "Proof of the Pullback Formula for m=1"、"Applications" 为未编号子标题,用 \mtsubsecn{}{}.
%   - 习题按原书分节编号: 第 1 节 1--5;第 2 节 2 个未编号"习题.";第 3 节 1--6 加 1 个未编号;第 4 节 1--6.

\chapter{分布理论}

\section*{引言}

以下是对\newterm{distribution}{分布}{distribution}理论的非常简短的概述;如需更深入地了解这一主题, 我建议参阅 F.G.~Friedlander 与 M.~Joshi 的优秀著作 \emph{Introduction to the Theory of Distributions} \cite{fr-io}.Hörmander 在 Springer 出版的 \emph{The Analysis of Linear Partial Differential Operators} 第一卷 \cite{ho} 也会很有用.

\textbf{记号.} 在本讲义中, 我们用记号 $A\lesssim B$ 表示 $A\leq cB$, 其中 $c$ 是一个与 $A$、$B$ 无关的数值常数.当 $\Omega\subseteq\mathbb{R}^n$ 是一个集合时, 我们可以用 $(x\in\Omega)$ 表示集合 $\Omega$ 的\newterm{indicator-function}{示性函数}{indicator function}.例如, $(5\leq|x|<7)$ 是一个在 $5\leq x<7$ 时等于 1、否则等于 0 的函数.

\section{分布理论导论}

对分布理论的简短描述包含一个不可避免的矛盾说法:它是一种"使能"理论, 使我们能够对完全不可微的函数求导, 并像根本不存在任何问题一样对它们进行操作.它的主要应用是在偏微分方程理论中.

我们首先回顾偏微分方程中 \newterm{fundamental-solution}{基本解}{fundamental solution} 的概念是如何通过经典电磁学诞生的.当电荷不运动时, 任何电荷分布 $\rho : \mathbb{R}^3\to\mathbb{R}$ 都会产生一个电场, 该电场(不计通常的符号与物理常数)是"势函数" $V : \mathbb{R}^3\to\mathbb{R}$ 的梯度.将 $V$ 与电荷密度 $\rho$ 联系起来的经典物理定律是 Poisson 方程

\begin{equation}
\Delta V = \rho\label{eq:c3-127}
\end{equation}

其中 $\Delta = \sum_{i=1}^3\partial_i^2$ 表示 Laplace 算子, 并且基于物理理由, 我们可以要求 $V$(或至少其导数)在无穷远处消失, 从而使远距离相互作用变弱.

与其他任何场论一样, 除非方程 \eqref{eq:c3-127} 对(合理的资料 $\rho$)存在唯一解, 且该解在某种意义下连续依赖于资料, 否则该物理理论不可能有效而完备.除了解决这些问题之外, 我们还至少希望对解的行为有定性的理解.就目前的情形而言, 由于存在大量的对称性, 我们甚至能够导出显式公式;但对于所涉及的启发式分析而言, 重要的只是算子 $\Delta = \sum_{i=1}^3\partial_i^2$ 是线性的且与平移交换(即它是常系数的线性微分算子).事实上, Laplace 算子同时关于平移和旋转不变, 其意义是:当 $T$ 是 Euclid 空间的刚体运动时, 对所有光滑函数都有 $\Delta(f\circ T) = (\Delta f)\circ T$.

想法是先求解 \eqref{eq:c3-127} 的特殊情形, 即电荷密度 $\rho(x)$ 是集中在点 $\boldsymbol{y}\in\mathbb{R}^3$ 处的单位电荷(我们形式地写 $\rho(x) = \delta_y(x)$, 其中 Dirac delta 函数对应于在 $y$ 处的单位点质量的密度函数).我们将在后面讨论 Dirac delta 函数的含义, 但此刻让我们接受它的形式定义:一个作用在连续函数 $f$ 上、在点 $y$ 处产生 $f$ 的值的算子, 即 $\delta_y$ 是测度,

\[
\delta_y(f) = f(y).
\]

因此, 我们寻找方程 $\Delta V_y(x) = \delta_y(x)$ 的解 $V_y(x)$(该方程目前无意义, 因为 $\delta_y$ 不是函数).由 $\Delta$ 的线性, 我们可以把通解写成点贡献之解的叠加:

\begin{equation}
V(x) = \int V_y(x)\rho(y)\, dy\label{eq:c3-128}
\end{equation}

形式地, 我们甚至可以对任意固定的 $y\in\mathbb{R}^3$ 求解方程 $\Delta V_y(x) = \delta_y(x)$.鉴于 $\Delta$ 的平移不变性, 我们可以假设 $y=0$.由于 $\Delta$ 旋转不变(见习题 1), $\delta_0$ 亦然, 那么若解唯一, 任何解 $V_0(x) = V_0(|x|)$ 也应旋转不变.我们称 $V_0(x)$ 为 Laplace 算子的基本解.于是, 假设 $V_0(x)$ 存在且具有球对称性, 我们得到(利用散度定理)

\[
\begin{aligned}
1 &= \int_{|x|\leq R}\delta_0(x)\, dx \\
&= \int_{|x|\leq R}\Delta V_0(x)\, dx \\
&= \int_{|x|=R}\frac{dV}{dr}(|x|)\, d\sigma(x) \\
&= 4\pi R^2\frac{dV}{dr}(R)
\end{aligned}
\]

我们选择在无穷远处衰减的唯一基本解, 即 $V_0(x) = \frac{-1}{4\pi}\frac{1}{R}$.因此, 平移回 $\delta_y$, 我们得到 $V_y(x) = \frac{-1}{4\pi}\frac{1}{|x-y|}$.通过直接计算可以看出, 在 $y$ 之外 $\Delta V_y(x) = 0$, 甚至可以证明, 对(譬如说)光滑、\newterm{compactly-supported}{紧支撑}{compactly supported}的密度 $\rho$, \eqref{eq:c3-128} 确实求解了 \eqref{eq:c3-127}.此外, 对 \eqref{eq:c3-128} 取梯度, 可以得到一个可被实验证伪的结论:远离电荷源处, 电场按

\[
\frac{\text{总电荷}}{(\text{距离})^2}
\]

渐近衰减.

\begin{definition*}
我们称 $V_y(x) = \frac{-1}{4\pi}\frac{1}{|x-y|}$ 为 $\Delta$ 在 $\mathbb{R}^3$ 中的基本解.更一般地, 我们定义线性算子 $L$ 在 $\mathbb{R}^n$ 中(即作用于 $\mathbb{R}^n$ 上的函数)的基本解为"广义"函数 $V_y(x)$, 使得 $L(V_y) = \delta_y$.
\end{definition*}

给定 $L$ 的一个基本解, 我们可以对任何光滑、紧支撑的 $f$ 求出方程 $Lu = f$ 的解:形式地(暂时不管积分可能毫无意义)令

\[
u(x) = \int V_y(x)f(y)\, dy
\]

\textbf{习题 1.} 非正式地证明:若 $L$ 与平移交换, 即对所有平移 $x\mapsto x+y$ 有 $(Lf)(\cdot+y) = L(f(\cdot+y))$, 则基本解也与平移交换, 其意义是 $V_y(x) = V(x-y)$, 其中 $V$ 满足 $L(V) = \delta_0$.

一旦找到算子 $L$ 的基本解 $V_y$, 我们既需要把它理解为广义函数, 也需要理解上述形式积分的意义.这正是分布理论所完成的.分布理论使我们能够把启发式计算严格化, 更重要的是, 它使我们能够像处理正则函数一样处理奇异对象.当然, 这种新自由是有界限的, 一个好的理论应当把这些界限讲清楚.

\textbf{习题 2.} 不难证明, 对 $\rho\in\mathcal{C}_0^{\infty}(\mathbb{R}^3)$, 势函数 $V(x) = \frac{-1}{4\pi}\int_{\mathbb{R}^3}\frac{1}{|x-y|}\rho(y)\, dy$ 在无穷远处、在 $\rho$ 的支集之外, 行为如同 $\frac{-1}{4\pi}\frac{\int_{\mathbb{R}^3}\rho(y)\, dy}{|x|} + o(|x|^{-2})$.证明这一渐近性并理解误差的一种方法是对 $\frac{1}{|x-sy|} = \frac{1}{|x|} + \int_0^1\frac{d}{ds}\frac{1}{|x-sy|}\, ds$ 作 Taylor 展开(其想法是参数 $y$ 相对较小).

当电荷分布以原点为中心(即向量值积分 $\int_{\mathbb{R}^3}y\rho(y)\, dy = 0$)时, 证明更精确的结果(带显式余项):当 $|x|\to\infty$ 时,

\[
V(x) = \frac{-1}{4\pi}\frac{\int_{\mathbb{R}^3}\rho(y)\, dy}{|x|} + O(|x|^{-3})
\]

为使计算简单, 可以把精确的一阶 Taylor 展开 $\phi(1) = \phi(0) + \phi'(0) + \int_0^1(1-s)\phi''(s)\, ds$ 应用于辅助函数 $\phi(s) = \frac{1}{|x-sy|}$.另外, 对绝对值函数求导的一个方便途径是观察到 $|x-sy|^2 = \langle x-sy, x-sy\rangle$, 其中 $\langle\, ,\, \rangle$ 表示 Euclid 内积.

\begin{remark*}
若总电荷 $\int\rho(y)\, dy$ 不为 0, 则可以找到"电荷中心" $y_c = \int y\rho\, dy/\int\rho\, dy$, 使得 $\int(y-y_c)\rho(y)\, dy = 0$.在这种情况下, 我们可以围绕 $y = y_c$ 作 Taylor 展开, 看到相应的势函数渐近行为如同以 $y_c$ 为中心:

\[
V(x) \approx \frac{-1}{4\pi}\frac{\int_{\mathbb{R}^3}\rho(y)\, dy}{|x-y_c|}.
\]

不过要注意, 当电荷以 $\int\rho(y)\, dy = 0$ 的意义"抵消"时, 相应的势函数 $V$ 在无穷远处衰减得更快, 如同 $\frac{C}{|x|^2}$.这种对局部化、振荡数据增强衰减的现象, 不仅在物理解释电荷抵消时电作用力在远距离上为何微弱方面很重要, 而且在分析中也十分重要——在许多其他自然出现的情形中也会产生类似的抵消.我们将在后面的讲义中看到这种抵消被以关键方式使用.
\end{remark*}

\textbf{习题 3.} 前一节中的推理可以推广到"求解"有界区域内部的势函数, 该区域的边界接地.即考虑问题:对边界上 $V=0$ 的有界区域 $\Omega$ 中的 $x$, 有 $\Delta V(x) = \rho(x)$.原则上, 你如何构造形如 $V(x) = \int K(x, y)\rho(y)\, dy$ 的通解?线性在哪里起作用?

\textbf{习题 4.} 假设在 $\mathbb{R}^3$ 中半径为 $R_1$ 的球面上均匀分布着一个单位负电荷, 在半径为 $R_2$ 的球面上均匀分布着一个单位正电荷.求与此电荷构型 $\rho$ 相关的静电势函数 $V$.

\textbf{习题 5.} a.~用引言中的非正式论证, 对每个 $n\geq 2$ 求出 $\Delta$ 在 $\mathbb{R}^n$ 中的基本解 $K_n(x)$, 即求解 $\Delta K_n(x) = \delta_0(x)$, 并给出 $K_n$ 的显式公式.

b.~讨论相应的解

\[
V(x) = \int K_n(x-y)\rho(y)\, dy
\]

当 $|x|\to\infty$ 时的行为, 其中 $\rho$ 紧支撑.即当 $|x|\to\infty$ 时, 主项是什么, 误差有多大?

\section{检验函数与分布}

我们从一些标准记号开始.我们用 $x = (x_1, \ldots, x_n)$ 表示 $\mathbb{R}^n$ 中的向量, 并令 $\lambda x = (\lambda x_1, \ldots, \lambda x_n)$, $x+y = (x_1+y_1, \ldots, x_n+y_n)$.我们用 $x\cdot y$ 表示标准标量积, 用 $|x| = (x\cdot x)^{\frac{1}{2}}$ 表示 $x$ 的 Euclid 长度.给定函数 $f : \Omega\to\mathbb{C}$, 我们用 $\mathrm{supp}(f)$ 表示集合 $\{x\in\Omega : f(x)\neq 0\}$ 在 $\Omega$ 中的闭包, 并称之为 $f$ 的\newterm{support}{支集}{support}.我们用 $\mathcal{C}^k(\Omega)$ 表示 $\Omega$ 上 $k$ 次连续可微的复值函数集合, 用 $\mathcal{C}_0^k(\Omega)$ 表示其中紧支撑的子集.我们还用 $\mathcal{C}^{\infty}(\Omega) = \cap_{k\in\mathbb{N}}\mathcal{C}^k(\Omega)$ 表示无穷可微函数空间, 用 $\mathcal{C}_0^{\infty}(\Omega)$ 表示其中紧支撑的子集.后者在分布理论中起着特别重要的作用;它被称为 $\Omega$ 上的\newterm{test-function}{检验函数}{test function}空间.

设 $\Omega\subset\mathbb{R}^n$, $f\in\mathcal{C}^{\infty}(\Omega)$.我们用 $\partial_i f$ 表示偏导数 $\frac{\partial f}{\partial x_i}$, $i = 1, \ldots, n$.对高阶导数我们使用标准的\newterm{multi-index}{多重指标}{multi-index}记号.一个多重指标 $\alpha$ 是非负整数的 $n$ 元组 $\alpha = (\alpha_1, \ldots, \alpha_n)$, 其长度为 $|\alpha| = \alpha_1+\cdots+\alpha_n$.令 $\alpha+\beta = (\alpha_1+\beta_1, \ldots, \alpha_n+\beta_n)$.我们用 $\alpha!$ 表示阶乘之积 $\alpha_1!\cdots\alpha_n!$.再令 $\partial^{\alpha}f = \partial_1^{\alpha_1}\cdots\partial_n^{\alpha_n}f$.显然 $\partial^{\alpha+\beta}f = \partial^{\alpha}\partial^{\beta}f$.给定两个光滑函数 $u$、$v$, 我们有 Leibniz 公式:

\[
\partial^{\alpha}(u\cdot v) = \sum_{\beta+\gamma=\alpha}\frac{\alpha!}{\beta!\gamma!}\partial^{\beta}u\, \partial^{\gamma}v.
\]

光滑函数 $f : \mathbb{R}^n\to\mathbb{C}$ 在原点附近的 Taylor 公式可以写成如下形式:

\[
f(x) = \sum_{|\alpha|\leq k}\frac{1}{\alpha!}\partial^{\alpha}f(0)\, x^{\alpha} + O(|x|^{k+1}) \quad \text{当 } x\to 0.
\]

这里 $x^{\alpha}$ 表示单项式 $x^{\alpha} = x_1^{\alpha_1}\cdots x_n^{\alpha_n}$.

我们首先解释一种一般方法(通常称为\newterm{mollification}{磨光}{mollification}), 它可以用来用光滑函数逼近较粗糙的函数.本质上, 取要逼近的函数 $f$, 按照某个支集很小的光滑概率测度 $\rho$ 随机平移 $f$ 后, 用其平均代替 $f$.直观上很清楚且容易证明, 随机扰动后的函数是原函数的光滑逼近(想象一场小地震后棱角分明的沙堆);为了得到更好的逼近, 把 $\rho$ 的支集收缩到 0.这一方法的技术实现如下.

\begin{proposition}\label{prop:c3-2.1}
设 $f\in\mathcal{C}_0^k(\mathbb{R}^n)$, $0\leq k<\infty$.设 $\rho\in\mathcal{C}_0^{\infty}(\mathbb{R}^n)$ 是支集 $\mathrm{supp}(\rho)$ 含于单位球 $B(0, 1) = \{|x|\leq 1\}$ 的光滑函数, 且 $\int\rho(x)\, dx = 1$.我们令 $\rho_{\varepsilon}(x) = \varepsilon^{-n}\rho(x/\varepsilon)$, 并设

\[
f_{\varepsilon}(x) = f*\rho_{\varepsilon}(x) = \varepsilon^{-n}\int f(y)\rho\left(\frac{x-y}{\varepsilon}\right)dy = \int f(x-\varepsilon z)\rho(z)\, dz.
\]

则我们有:

\begin{enumerate}
\item 函数 $f_{\varepsilon}$ 属于 $\mathcal{C}_0^{\infty}(\mathbb{R}^n)$, 且 $\mathrm{supp}(f_{\varepsilon})\subset\mathrm{supp}(f) + B(0, \varepsilon)$;
\item 当 $\varepsilon\to 0$ 时, $\partial^{\alpha}f_{\varepsilon}\longrightarrow\partial^{\alpha}f$ 一致收敛.
\end{enumerate}
\end{proposition}

\begin{proof}
命题的第一部分由定义立即得出, 因为关于支集的论断是显然的, 并且通过分部积分, 我们可以把 $f_{\varepsilon}$ 的所有导数转移到被积函数的光滑部分 $\rho_{\varepsilon}$ 上.为证明第二个论断, 我们只需写出

\[
\partial^{\alpha}f_{\varepsilon}(x) - \partial^{\alpha}f(x) = \int\left(\partial^{\alpha}f(x-\varepsilon z) - \partial^{\alpha}f(x)\right)\rho(z)\, dz.
\]

因此, 对 $|\alpha|\leq k$,

\[
\begin{aligned}
|\partial^{\alpha}f_{\varepsilon}(x) - \partial^{\alpha}f(x)| &\leq \int|\partial^{\alpha}f(x-\varepsilon z) - \partial^{\alpha}f(x)|\, |\rho(z)|\, dz \\
&\leq \int|\rho(z)|\, dz\sup_{|z|\leq 1}|\partial^{\alpha}f(x-\varepsilon z) - \partial^{\alpha}f(x)| \\
&\lesssim \sup_{|z|\leq 1}|\partial^{\alpha}f(x-\varepsilon z) - \partial^{\alpha}f(x)|
\end{aligned}
\]

考虑到函数 $\partial^{\alpha}f$ 的一致连续性, 证明现在可以容易地完成.
\end{proof}

作为该命题的推论, 可以容易地验证:检验函数空间 $\mathcal{C}_0^{\infty}(\Omega)$ 在空间 $\mathcal{C}^k(\Omega)$ 以及 $L^p(\Omega)$, $1\leq p<\infty$ 中稠密.当然, 首先必须至少构造一个满足 $\int\rho\, dx = 1$ 的 $\rho\in C_0^{\infty}(\mathbb{R}^n)$.隆起函数 $\rho(x) = e^{\frac{-1}{1-|x|^2}}\cdot(|x|<1)$ 的某个倍数即可.构造例子的另一种方法是从任意 $\mathcal{C}^1$ 隆起函数出发, 利用随机平移的光滑化效应(如上述命题), 同时控制支集, 通过迭代过程的极限得到光滑隆起函数.

\begin{definition}\label{def:c3-2.2}
分布 $u\in\mathcal{D}'(\Omega)$ 是线性泛函 $u : \mathcal{C}_0^{\infty}(\Omega)\to\mathbb{C}$, 它验证如下性质:

对任意紧集 $K\subset\Omega$, 存在整数 $N$ 与常数 $C = C_{K, N}$, 使得对所有支集满足 $\mathrm{supp}(\phi)\subset K$ 的 $\phi\in\mathcal{C}_0^{\infty}(\Omega)$, 有

\[
|\langle u, \phi\rangle| \leq C\sum_{|\alpha|\leq N}\sup|\partial^{\alpha}\phi|.
\]

如果在上面的定义中对每个 $K$ 都可以使用同一个整数 $N$, 则称最小的这样的 $N$ 为分布的阶.例如, Riesz 表示定理(刻画了紧 Hausdorff 空间上 $\mathcal{C}(X)$ 的对偶)保证 0 阶分布就是 Borel 测度.
\end{definition}

等价地, 分布 $u$ 是关于定义在 $\mathcal{C}_0^{\infty}(\Omega)$ 上的某个拓扑连续的线性泛函 $u : \mathcal{C}_0^{\infty}(\Omega)\to\mathbb{C}$.这个拓扑是一种相当非正统的拓扑(不可度量\footnote{该拓扑可以构造为 Fréchet 空间 $C_K$ 的归纳极限拓扑, 其中 $K\subseteq\Omega$ 紧致, $C_K$ 是支集含于 $K$ 的所有光滑函数构成的空间, 由半范数 $\phi\mapsto\sup_K|\partial^{\alpha}\phi|$(对所有多重指标 $\alpha$)赋予 Fréchet 空间结构.不过我们并不需要它的精确定义.}, 局部凸), 但不必介意这些;我们完全可以不担心其精确定义而走得很远.我们只需要知道, 在这个拓扑中, 当 $\phi_j$ 的所有支集都包含在 $\Omega$ 的一个紧子集中, 并且对每个多重指标 $\alpha$, $\partial^{\alpha}\phi_j$ 在一致范数下趋于 0 时, 序列 $\phi_j$ 在 $\mathcal{C}_0^{\infty}(\Omega)$ 中收敛到 0.记住这个定义, 我们有以下非常有用的分布刻画:

\begin{proposition}\label{prop:c3-2.3}
线性形式 $u : \mathcal{C}_0^{\infty}(\Omega)\longrightarrow\mathbb{C}$ 是 $\mathcal{D}'(\Omega)$ 中的分布, 当且仅当对每个在 $\mathcal{C}_0^{\infty}(\Omega)$ 中当 $j\to\infty$ 时收敛到 0 的检验函数序列 $\phi_j$, 都有 $\lim_{j\to\infty}u(\phi_j) = 0$.
\end{proposition}

\begin{proof}
此证明可见 Friedlander 著作第 1.3 节, 定理 1.3.2.
\end{proof}

\begin{example*}
任何\newterm{locally-integrable}{局部可积}{locally integrable}函数 $f\in L^1_{\mathrm{loc}}(\Omega)$ 定义了一个分布:

\[
\langle f, \phi\rangle = \int f\phi, \quad \forall\phi\in\mathcal{C}_0^{\infty}(\Omega).
\]

于是我们可以把 $L^1_{\mathrm{loc}}(\Omega)$ 视为 $\mathcal{D}'(\Omega)$ 的子空间.特别地, 空间 $\mathcal{C}^{\infty}(\Omega)\subset L^1_{\mathrm{loc}}(\Omega)$ 也是如此.

即使 $u\in\mathcal{D}'(\Omega)$ 不是局部可积函数, 甚至当 $\phi$ 严格说来不是检验函数时, 人们也常常使用形式记号 $\langle u, \phi\rangle = \int_{\Omega}u(x)\phi(x)\, dx$.这种记号在概念上可能更简单, 但请记住, 这绝不是真正的 Lebesgue 积分.不过, 通常可以把这种形式积分解释为经典积分的极限.
\end{example*}

\begin{example*}
支集在 $x_0\in\mathbb{R}^n$ 处、质量为 1 的 Dirac 测度定义为

\[
\langle\delta_{x_0}, \phi\rangle = \phi(x_0).
\]
\end{example*}

\begin{remark*}
我们也常常用 $u(\phi)$ 而不是 $\langle u, \phi\rangle$ 来表示分布 $u$ 作用在检验函数上的结果.于是 $\delta_{x_0}(\phi) = \phi(x_0)$.
\end{remark*}

\begin{definition}\label{def:c3-2.4}
称分布序列 $u_j\in\mathcal{D}'(\Omega)$ 弱收敛到分布 $u\in\mathcal{D}'(\Omega)$, 若对所有 $\phi\in\mathcal{C}_0^{\infty}(\Omega)$, 有 $u_j(\phi)\to u(\phi)$.
\end{definition}

例如, 当 $m\to\infty$ 时, 序列 $u_m = e^{imx}$ 在 $\mathcal{D}'(\mathbb{R})$ 中弱收敛到 0.另外, 若 $f\in L^1(\mathbb{R}^n)$, 且 $\int_{\mathbb{R}^n}f(x)\, dx = 1$, 则当 $\lambda\to\infty$ 时, 函数族 $f_{\lambda}(x) = \lambda^n f(\lambda x)$ 弱收敛到 $\delta_0$.

我们将能够证明, 任何分布都是序列 $u_n\in\mathcal{C}_0^{\infty}(\Omega)$ 的弱极限.由于这一事实, 许多最初对函数定义的操作都可以以唯一、自然的方式通过连续性延拓到 $\mathcal{D}'(\Omega)$.

\textbf{习题.} 给定紧集 $K\subseteq\mathbb{R}^n$ 与正距离 $\delta>0$, 构造一个光滑函数 $\eta : \mathbb{R}^n\to\mathbb{R}$, 使其支集包含在集合 $\{x \mid |y-x|\leq\delta \text{ 对某个 } y\in K\}$ 中, 即 $K$ 的 $\delta$-邻域内, 并且使得在 $K$ 的某个邻域中 $\eta(y) = 1$.提示:从一个粗糙的\newterm{cutoff-function}{截断函数}{cutoff function}出发, 留出一些"回旋余地".

\mtsubsec{分布上的运算}

使用分布空间的好处在于, 虽然这个空间远大于光滑函数空间, 但检验函数上最重要的操作都可以搬到分布上来.

\textbf{1. 与光滑函数相乘:} 给定 $u\in\mathcal{D}'(\Omega)$ 与 $f\in\mathcal{C}^{\infty}(\Omega)$, 我们定义

\[
\langle fu, \phi\rangle = \langle u, f\phi\rangle, \quad \forall\phi\in\mathcal{C}_0^{\infty}(\Omega)
\]

以便与 $u$ 是函数时的恒等式一致:

\[
\int(fu)\phi\, dx = \int u(f\phi)\, dx.
\]

容易验证, 乘以光滑函数是分布空间的连续自同态.

\textbf{2. 与检验函数的\newterm{convolution}{卷积}{convolution}:} 考虑 $u\in\mathcal{D}'(\mathbb{R}^n)$, $\phi\in\mathcal{C}_0^{\infty}(\mathbb{R}^n)$.以自然的方式推广两个函数的卷积, 我们定义

\[
u*\phi(x) = \langle u_y, \phi(x-y)\rangle,
\]

下标表示 $u$ 理解为作用于变量 $y$ 的函数.注意, 若 $u$ 是局部可积函数, $u\in L^1_{\mathrm{loc}}(\mathbb{R}^n)$, 则该定义与通常的定义一致:

\[
u*\phi(x) = \int u(y)\phi(x-y)\, dy.
\]

\begin{remark*}
分布与检验函数的卷积不仅仅是另一个分布.相反, 注意对每个分布 $u\in\mathcal{D}'(\mathbb{R}^n)$ 与 $\phi\in\mathcal{C}_0^{\infty}(\mathbb{R}^n)$, 我们有 $u*\phi\in\mathcal{C}^{\infty}(\mathbb{R}^n)$, 它实际上是一个光滑函数.例如, 若 $e_k$ 表示标准单位向量, 则我们可以沿 $e_k$ 方向求导如下:

\[
\begin{aligned}
\frac{u*\phi(x+he_k) - u*\phi(x)}{h} &= h^{-1}\langle u_y, \phi(x+he_k-y) - \phi(x-y)\rangle \\
&= \left\langle u_y, \int_0^1\partial_k\phi(x+the_k-y)\, dt\right\rangle.
\end{aligned}
\]

现在, 由于 $x\in K$ 被限制在某个紧集 $K\subset\mathbb{R}^n$ 中, 故对每个序列 $h_i\to 0$, 相应的函数序列 $y\mapsto\int_0^1\partial_k\phi(x+th_ie_k-y)\, dt$ 连同其所有导数一致收敛到 $\partial_k\phi(x-y)$ 及其相应的导数.而且它们都是紧支撑的, 支集包含在某个紧集 $K'$ 中.因此,

\[
\lim_{h\to 0}\frac{u*\phi(x+he_k) - u*\phi(x)}{h} = u*\partial_k\phi(x).
\]

从而 $u*\phi$ 具有连续偏导数.我们可以继续这样下去, 并得出结论:事实上 $u*\phi\in\mathcal{C}^{\infty}(\mathbb{R}^n)$.
\end{remark*}

\textbf{3. 分布的微分:} 对每个分布 $u\in\mathcal{D}'(\Omega)$, 我们定义

\[
\langle\partial^{\alpha}u, \phi\rangle = (-1)^{|\alpha|}\langle u, \partial^{\alpha}\phi\rangle.
\]

我们作此定义是为了与函数的分部积分公式一致:

\[
\int\partial_i u(x)\phi(x)\, dx = \int u(x)(-\partial_i\phi(x))\, dx, \quad \phi\in\mathcal{C}_0^{\infty}(\mathbb{R}^n),
\]

该公式可以(例如)通过考虑差商来证明.

同样, 容易验证我们由此定义了分布空间的连续自同态.当然, 上述操作是光滑函数上通常操作的唯一可能延拓.负号可以对偶地看作对测度 $\mu = \phi(x)\, dx$ 的微分.若我们暂时用 $\tau_v$ 表示"按 $v$ 平移"的运算, 则

\[
\partial_1\mu = \lim_{h\to 0}\frac{\tau_{he_1}\mu - \mu}{h}
\]

具有密度函数 $-\partial_1\phi(x)$, 且极限可以在 $\mathcal{C}_0^{\infty}$ 的拓扑中取.在这种对偶意义下, 我们有 $\partial_1u = \lim_{h\to 0}\frac{u(x+he_1)-u(x)}{h}$, 这里的极限是在弱拓扑中取的;只要我们以分布理论的意义解释所有积分, 这常常使我们能够在"积分号下求导".

我们现在可以定义一般的线性偏微分算子对分布的作用.事实上, 设

\[
P(x, \partial) = \sum_{|\alpha|\leq m}a_{\alpha}\partial^{\alpha}, \quad a_{\alpha}\in\mathcal{C}^{\infty}(\Omega),
\]

是这样一个算子.则

\[
\langle P(x, \partial)u, \phi\rangle = \langle u, P(x, \partial)^{\dagger}\phi\rangle,
\]

其中 $P(x, \partial)^{\dagger}$ 是形式共轭算子:

\[
P(x, \partial)^{\dagger}v = \sum_{|\alpha|\leq m}(-1)^{|\alpha|}\partial^{\alpha}(a_{\alpha}v).
\]

注意, 若 $u_j\in\mathcal{D}'(\Omega)$ 弱收敛到 $u\in\mathcal{D}'(\Omega)$, 则 $P(x, \partial)u_j$ 弱收敛到 $P(x, \partial)u$.

\textbf{习题.} 证明对所有 $u\in\mathcal{D}'(\Omega)$, 存在序列 $u_j\in\mathcal{C}_0^{\infty}(\Omega)$, 使得当 $j\to\infty$ 时 $u_j\to u$(按分布的意义, 即弱收敛).从而 $\mathcal{C}_0^{\infty}(\Omega)$ 在 $\mathcal{D}'(\Omega)$ 中稠密(关于后者的弱拓扑).

\section{实直线上的分布举例}

\textbf{1.} 最简单的非平凡分布是 Dirac delta 函数 $\delta_0 = \delta_0(x)$, 定义为 $\langle\delta_0(x), \phi\rangle = \phi(0)$.我们有时不带下标地写 $\delta(x)$, 表示 $\mathbb{R}$ 上原点处的点质量.

\textbf{2.} 另一个简单的例子是 Heaviside 函数 $H(x)$, 当 $x\geq 0$ 时等于 1, 当 $x<0$ 时等于 0.或者, 利用局部可积函数与分布之间的标准等同:

\[
\langle H(x), \phi\rangle = \int_0^{\infty}\phi(x)\, dx.
\]

注意 $H'(x) = \delta(x)$, 且 $H(x) = \int_{-\infty}^x\delta(t)\, dt$ 是 $\delta_0$ 的累积分布函数.

\textbf{3.} 一个更精细的例子是 $\mathrm{pv}(\frac{1}{x})$, 或简写为 $\frac{1}{x}$, 称为\newterm{principal-value}{主值}{principal value}分布:

\[
\left\langle\frac{1}{x}, \phi\right\rangle = \lim_{\varepsilon\to 0}\left(\int_{-\infty}^{-\varepsilon}\frac{1}{x}\phi(x)\, dx + \int_{\varepsilon}^{\infty}\frac{1}{x}\phi(x)\, dx\right).
\]

注意, $\log|x|$ 是局部可积的, 从而按照标准等同是一个分布.容易证明 $\frac{d}{dx}\log|x| = \mathrm{pv}(\frac{1}{x})$.注意, $\mathrm{pv}(\frac{1}{x})$ 是奇分布(它与偶检验函数正交), 并且是一阶的, 尽管在原点之外它是 0 阶的.事实上, 把 $\phi = \phi_{\mathrm{ev}} + \phi_{\mathrm{odd}}$ 分解为偶部与奇部, 我们有

\[
\left\langle\mathrm{pv}\left(\frac{1}{x}\right), \phi\right\rangle = \int\frac{\phi_{\mathrm{odd}}}{x}\, dx = \int\int_{-1}^1\phi'(tx)\, dt\, dx
\]

我们还要注意, 相比之下, 函数 $\frac{1}{|x|}\cdot(x\neq 0)$ 不能作为分布延拓到整条直线上.

\textbf{习题 1.} 证明直线上的分布 $t\frac{d}{dt}\delta(t)$ 等于 $-\delta(t)$, 这是一个非零分布.这看起来可能违反直觉, 因为 $t$ 或 $\delta'(t)$ 似乎在每个点都消失.

\textbf{习题 2.} 对满足 $0<\arg(z)<\pi$ 的 $z\in\mathbb{C}$, 设 $\log z = \log|z| + i\arg(z)$.我们可以把 $x\mapsto\log z = \log(x+iy)$ 视为一族依赖于 $y\in\mathbb{R}^+$ 的分布.对 $x\neq 0$, 我们有 $\lim_{y\to 0^+}\log z = \log|x| + i\pi(1-H(x))$.证明当 $y\to 0$, $y\in\mathbb{R}^+$, $\partial_x\log z$ 弱收敛到分布 $\frac{1}{x+i0}$, 并且

\begin{equation}
\frac{1}{x+i0} = x^{-1} - i\pi\delta_0(x).\label{eq:c3-129}
\end{equation}

\textbf{习题 3.} 若 $\Omega$ 是开连通集, $u\in\mathcal{D}'(\Omega)$, 且所有导数 $\partial_i u = 0$(按分布的意义), 则 $u$ 是常数.

\textbf{习题 4.} 任何非负分布(即当 $\phi\geq 0$ 时 $\langle u, \phi\rangle\geq 0$)事实上都是 Borel 测度.根据测度的 Riesz 表示定理, 只需证明:对每个紧集 $K\subseteq\mathbb{R}^n$, 存在常数 $C = C_K$, 使得对所有支集在 $K$ 中的 $\phi$,

\[
|u(\phi)|\leq C_K\max|\phi|
\]

成立.

\textbf{习题 5.} (利用前面的习题)若 $u : (-\infty, b)\to\mathbb{R}$ 是非降函数, 并且(为简单起见)在 $-\infty$ 处消失, 则存在唯一的 Borel 测度 $\mu\geq 0$, 使得对 Lebesgue 几乎处处 $x\in(-\infty, b]$, 有 $u(x) = \int_{-\infty}^x 1\, d\mu(t)$.用 $\mu$ 的语言来说, $u$ 何时连续?何时绝对连续?

\begin{remark*}
单调函数几乎处处可微这一经典结果可以由上述习题结合一些基本的测度论与第 2.6 节的 Lebesgue 微分定理导出.
\end{remark*}

\textbf{习题 6.} 刻画凸函数 $f : (a, b)\to\mathbb{R}$.即证明以下论断等价:

\begin{enumerate}
\item 对所有 $\alpha\in[0, 1]$ 与 $x, y\in(a, b)$, 有 $f(\alpha x+(1-\alpha)y)\leq\alpha f(x)+(1-\alpha)f(y)$;
\item $f$ 连续;$f'$(按分布理论的意义)是非降函数, 从而局部有界;
\item $f$ 连续;$f''$ 是非负分布, 从而限制在任何有界集上时是有限测度.
\end{enumerate}

提示:凸函数类在随机平移运算下保持不变, 因此磨光可能有帮助.

\begin{center}$\square$\end{center}

我们现在通过\newterm{analytic-continuation}{解析延拓}{analytic continuation}定义一个重要的分布族 $\chi_+^z$, 其中 $z\in\mathbb{C}$.我们将在后面研究波动方程的基本解时再次见到这一族, 也将在研究 傅里叶变换的\newterm{restriction-theorem}{限制定理}{restriction theorem}时再次见到它.

首先回顾 Gamma 函数的定义:

\begin{definition}\label{def:c3-3.1}
对 $\mathrm{Re}(z)>0$, 我们定义

\begin{equation}
\Gamma(z) = \int_0^{\infty}e^{-t}t^{z-1}\, dt\label{eq:c3-130}
\end{equation}

以及 Beta 函数:

\begin{equation}
B(a, b) = \int_0^1s^{a-1}(1-s)^{b-1}\, ds\label{eq:c3-131}
\end{equation}
\end{definition}

显然 $\Gamma(a+1) = a\Gamma(a)$, 且 $\Gamma(1) = 1$.从而 $\Gamma(n+1) = n!$.回忆以下恒等式成立:

\begin{equation}
B(a, b) = \frac{\Gamma(a)\cdot\Gamma(b)}{\Gamma(a+b)}\label{eq:c3-132}
\end{equation}

我们还为将来的应用记下:

\begin{equation}
\Gamma(a)\Gamma(1-a) = B(a, 1-a) = \frac{\pi}{\sin(\pi a)}\label{eq:c3-133}
\end{equation}

特别地, $\Gamma(1/2) = \pi^{1/2}$.

\textbf{习题.} 证明公式 \eqref{eq:c3-132} 与 \eqref{eq:c3-133}.可参见 Hörmander 著作[?]第 3.4 节.

\begin{definition}\label{def:c3-3.2}
对 $\mathrm{Re}(a)>0$, 我们用 $j_a(\lambda)$ 表示这样的局部可积函数:对 $\lambda<0$ 恒等于 0, 且

\begin{equation}
j_a(\lambda) = \frac{1}{\Gamma(a)}\lambda^{a-1}, \quad \lambda>0.\label{eq:c3-134}
\end{equation}
\end{definition}

下面的命题是众所周知的:

\begin{proposition}\label{prop:c3-3.3}
对所有满足 $a, b, \mathrm{Re}(a), \mathrm{Re}(b)>0$ 的 $a, b$,

\[
j_a*j_b = j_{a+b}
\]
\end{proposition}

\begin{proof}
我们有:

\[
\begin{aligned}
j_a*j_b(\lambda) &= \frac{1}{\Gamma(a)}\frac{1}{\Gamma(b)}\int_0^{\lambda}\mu^{a-1}(\lambda-\mu)^{b-1}\, d\mu \\
&= \frac{1}{\Gamma(a)}\frac{1}{\Gamma(b)}\lambda^{a+b-1}\int_0^1s^{a-1}(1-s)^{b-1}\, ds \\
&= \frac{B(a, b)}{\Gamma(a)\cdot\Gamma(b)}\lambda^{a+b-1} = \frac{1}{\Gamma(a+b)}\lambda^{a+b-1} = j_{a+b}(\lambda)
\end{aligned}
\]
\end{proof}

\begin{proposition}\label{prop:c3-3.4}
存在一族对所有 $a\in\mathbb{C}$ 定义的分布 $j_a$, 它们与 $\mathrm{Re}(a)>0$ 时的函数 $j_a$ 一致, 并且满足 $j_a*j_b = j_{a+b}$, $\frac{d}{d\lambda}j_a(\lambda) = j_{a-1}(\lambda)$, 以及 $j_0 = \delta_0$, 即原点处的 Dirac delta 函数.此外, 对所有正整数 $m$, $j_{-m}(x) = \partial_x^m\delta_0(x)$.
\end{proposition}

\begin{proof}
证明基于观察 $\frac{d}{d\lambda}j_a(\lambda) = j_{a-1}(\lambda)$.于是, 对检验函数 $\phi$,

\[
\int_{\mathbb{R}}j_{a-1}(\lambda)\phi(\lambda)\, d\lambda = -\int_{\mathbb{R}}j_a(\lambda)\phi'(\lambda)\, d\lambda
\]

基于这一观察, 我们对每个满足 $\mathrm{Re}(a)+m>0$ 的 $a\in\mathbb{C}$, 把分布 $j_a$ 定义为

\[
\langle j_a, \phi\rangle = (-1)^m\int_0^{\infty}j_{a+m}(\lambda)\phi^{(m)}(\lambda)\, d\lambda
\]

特别地,

\[
\langle j_0, \phi\rangle = -\int_0^{\infty}j_1(\lambda)\phi'(\lambda)\, d\lambda = -\int_0^{\infty}\phi'(\lambda)\, d\lambda = \phi(0)
\]

从而 $j_0 = \delta_0$.也容易看出, 对所有 $a, b\in\mathbb{C}$, 有 $j_a*j_b = j_{a+b}$.
\end{proof}

\begin{remark*}
在应用中人们常常看到分布族 $\chi_+^a = j_{a+1}$.显然 $\chi_+^a*\chi_+^b = \chi_+^{a+b+1}$, $\frac{d}{d\lambda}\chi_+^b(\lambda) = \chi_+^{b-1}(\lambda)$, 且 $\chi_+^{-1} = \delta_0$.还要注意, 对正整数 $k$, 有 $\chi_+^k(\lambda) = \frac{\lambda^k}{k!}\cdot(\lambda>0)$;更一般地, $\chi_+^a$ 是 $a$ 次齐次的, 即对任何正常数 $t$, $\chi_+^a(t\lambda) = t^a\chi_+^a(\lambda)$.当 $\mathrm{Re}(a)>-1$ 时(此时 $\chi_+^a$ 是函数), 这种齐次性显然是有意义的.你能对所有的 $a\in\mathbb{C}$ 也赋予它意义吗?
\end{remark*}

\mtsubsec{分布的支集}

分布的支集可以如下容易地导出:

\begin{definition}\label{def:c3-3.6}
对 $u\in\mathcal{D}'(\Omega)$, 我们定义 $u$ 的支集的补:

\[
\Omega\setminus\mathrm{supp}(u) = \{x\in\Omega \mid \exists\, V_x\ni x\text{ 开集, 使得 } \langle u, \phi\rangle = 0, \ \forall\phi\in\mathcal{C}_0^{\infty}(V_x)\}
\]
\end{definition}

\begin{lemma}\label{lem:c3-3.7}
若 $u\in\mathcal{D}'(\Omega)$, 且 $\phi$ 是支集满足 $\mathrm{supp}(\phi)\subset\Omega\setminus\mathrm{supp}(u)$ 的检验函数, 则 $\langle u, \phi\rangle = u(\phi) = 0$.
\end{lemma}

\begin{proof}
这可以通过\newterm{partition-of-unity}{单位分解}{partition of unity}论证容易地推出.该论证可见 Friedlander 著作第 1.4 节.
\end{proof}

上述引理可以用来证明, 任何紧支撑的分布 $u\in\mathcal{D}'(\Omega)$ 都可以延拓到 $\mathcal{C}^{\infty}(\mathbb{R}^n)$ 中的检验函数:取任意在 $u$ 的支集上等于 1 的截断函数 $\psi\in\mathcal{C}_0^{\infty}(\mathbb{R}^n)$, 并定义 $u(\phi) = u(\psi\phi)$.事实上, 下面的命题表明, 我们可以把紧支撑分布视为某个 $N$ 的 $\mathcal{C}^N(\mathbb{R}^n)$ 的对偶的元素:

\begin{proposition}\label{prop:c3-3.8}
分布 $u\in\mathcal{D}'(\mathbb{R}^n)$ 具有紧支集 $K\subset\mathbb{R}^n$, 当且仅当存在 $N\in\mathbb{N}$, 使得对所有 $\phi\in\mathcal{C}_0^{\infty}(\mathbb{R}^n)$, 有

\[
|u(\phi)|\leq C_U\sup_{x\in U}\sum_{|\alpha|\leq N}|\partial^{\alpha}\phi(x)|,
\]

其中 $U$ 是 $K$ 的任意开邻域.
\end{proposition}

\begin{proof}
这可以通过使用一个在分布的支集上恒等于 1 的截断函数来看出.
\end{proof}

\begin{remark*}
如果我们给 $\mathcal{C}^{\infty}(\mathbb{R}^n)$ 赋以由半范数族 $\phi\to\sup_{K_i}|\partial^{\alpha}\phi|$ 诱导的 Fréchet 拓扑, 其中 $\alpha\in\mathbb{N}^n$, $K_i$ 跑遍穷竭 $\mathbb{R}^n$ 的可数紧集族, 则紧支撑分布空间可以与 $\mathcal{C}^{\infty}(\mathbb{R}^n)^{*}$ 等同, 即与 $\mathcal{C}^{\infty}(\mathbb{R}^n)$ 的对偶空间等同.
\end{remark*}

我们有以下关于支集为单点的分布结构的有用事实(本质上与 Taylor 展开对偶).我们将在讲义的各个部分发现这一结果很有用, 尽管它的应用基本上总可以用重复其证明的某种变形来代替.我们不会给出全部细节, 但主要思想都在.

\begin{proposition}\label{prop:c3-3.9}
设 $u\in\mathcal{D}'(\mathbb{R}^n)$, 并假设 $\mathrm{supp}(u)\subset\{0\}$.则存在某个整数 $N$、复数 $a_{\alpha}$ 与 $\delta_0$(在 $\mathbb{R}^n$ 中支集在 0 处的 Dirac 测度), 使得

\[
u = \sum_{|\alpha|\leq N}a_{\alpha}\partial^{\alpha}(\delta_0).
\]
\end{proposition}

\begin{proof}
设 $u$ 是支集在原点的分布, $\phi\in\mathcal{C}_0^{\infty}(\mathbb{R}^n)$ 是检验函数.一般说来, $\langle u, \phi\rangle$ 并不只依赖于 $\phi$ 在 0 处的值, 但 $\langle u, \phi\rangle$ 确实只依赖于 $\phi$ 在 0 的任意小邻域上的限制.所以设 $\eta\in\mathcal{C}_0^{\infty}(\mathbb{R}^n)$ 是在原点的一个邻域上等于 1 的截断函数, 并设 $\eta_{\delta}(x) = \eta(\frac{x}{\delta})$($\delta>0$)是支集更小的截断函数.则 $\langle u, \phi\rangle = \langle u, \eta_{\delta}\phi\rangle$.我们希望证明 $\langle u, \phi\rangle$ 只依赖于 $\phi$ 的前 $N$ 阶导数, 其中

\[
|\langle u, \psi\rangle|\leq C\|\psi\|_{\mathcal{C}^N} = C\sum_{|\alpha|\leq N}\max|\partial^{\alpha}\psi|, \quad \psi\in\mathcal{C}_0^{\infty}(\mathbb{R}^n)
\]

通过 Taylor 展开, 归结为所有导数 $\partial^{\alpha}\phi = 0$(对 $|\alpha|\leq N$)的情形.更精确地说, 我们可以写

\[
\begin{aligned}
\phi(x) &= \phi(0) + \int_0^1\frac{d}{dt}\phi(tx)\, dt \\
&= \phi(0) + \sum_i\int_0^1x^i\partial^i\phi(tx)\, dt \\
&= \phi(0) + \sum_i\partial^i\phi(0)x^i + \int_0^1(1-t)\frac{d^2}{dt^2}\phi(tx)\, dt
\end{aligned}
\]

并继续分部积分, 直到把 $\phi$ 写成系数对应于 $\phi$ 在 0 处导数的多项式, 加上形如 $x^{\tau}\phi_{\tau}$ 的项之和, 其中 $\tau$ 是满足 $|\tau|>N$ 的多重指标, 函数 $\phi_{\tau}$ 光滑(但显然并非都有紧支集).

按这种方式展开 $\phi$, 我们需要证明 $\langle u, \eta_{\delta}x^{\tau}\phi_{\tau}\rangle = 0$.这里我们将使用估计 $|\langle u, \psi\rangle|\leq C\|\psi\|_{\mathcal{C}^N}$, 因此我们必须估计阶数 $|\alpha|\leq N$ 的导数 $\partial^{\alpha}(\eta_{\delta}x^{\tau}\phi_{\tau})$.注意

\[
\partial^{\alpha}\left(\eta\left(\frac{x}{\delta}\right)\right) = \delta^{-|\alpha|}(\partial^{\alpha}\eta)\left(\frac{x}{\delta}\right)
\]

这种导数与缩放的关系与量纲分析一致, 如果我们把 $\delta$ 与 $x$ 视为具有相同的"单位", 并把每个微分 $\partial^i$ 视为具有倒数单位.这也是分析中极其常见的一个计算, 我们后面将会看到.我们利用这一缩放, 把单项式因子吸收进截断函数:$\eta(\frac{x}{\delta})x^{\tau}\phi_{\tau} = \delta^{\tau}\tilde{\eta}(\frac{x}{\delta})\phi_{\tau}$.于是我们得到

\[
\begin{aligned}
|\langle u, \eta_{\delta}x^{\tau}\phi_{\tau}\rangle| &= \delta^{\tau}|\langle u, \tilde{\eta}_{\delta}\phi_{\tau}\rangle| \\
&\leq C\sum_{|\alpha|\leq N}\delta^{\tau}\max|\partial^{\alpha}(\tilde{\eta}_{\delta}\phi_{\tau})|
\end{aligned}
\]

人们可以用 Leibniz 公式展开这些导数:

\[
\partial^{\alpha}(\tilde{\eta}_{\delta}\cdot\phi_{\tau}) = \sum_{\beta+\gamma=\alpha}\frac{\alpha!}{\beta!\gamma!}\partial^{\beta}\tilde{\eta}_{\delta}\partial^{\gamma}\phi_{\tau}
\]

并产生数量庞大的项, 但求出这些乘积的精确公式可能并不有用, 尽管至少完整地过一遍细节或许是值得的.然而在实践中, 人们通过理解会出现哪些类型的项来避免那些并非问题核心的细节(例如常数的精确值), 特别是要挑出"最坏"的项.在本例中, 最坏的项出现在某个导数落到截断函数 $\tilde{\eta}_{\delta} = \tilde{\eta}(\frac{x}{\delta})$ 上时, 该截断函数随着 $\delta\to 0$ 变得越来越尖锐.对这样的项, 每个导数都产生一个 $\delta^{-1}$ 因子.但是, 至多有 $N$ 个导数可以落到这个截断函数上, 因此对某个与 $\delta$ 无关(尽管可能依赖于 $\eta$ 与 $\phi$)的常数 $C'$, 有

\[
|\langle u, \eta_{\delta}x^{\tau}\phi_{\tau}\rangle|\leq C'\delta^{\tau-N}
\]

由于 $\tau>N$, 当 $\delta\to 0$ 时它趋于 0.
\end{proof}

现在我们已经引入了支集的概念, 重要的是注意:两个分布的卷积一般不能定义, 只有当关于分布支集的某些条件满足时才能定义.我们特别注意到, 若 $u_1, u_2\in\mathcal{D}'(\mathbb{R}^n)$ 其中之一是紧支撑的, 则卷积 $u_1*u_2$ 可以定义.事实上, 假设 $u_2$ 紧支撑, 我们简单定义(∗∗∗)

\[
(u_1*u_2)*\phi = u_1*(u_2*\phi), \quad \forall\phi\in\mathcal{C}_0^{\infty}(\mathbb{R}^n).
\]

这里, $\mathrm{supp}(u_2*\phi)\subset\{x+y : x\in\mathrm{supp}(u_2), y\in\mathrm{supp}(\phi)\}$, 这是一个紧集.这一定义延拓了函数的经典卷积.

\section{分布的回拉}

设 $\Omega\subseteq\mathbb{R}^n$ 与 $\Omega'\subseteq\mathbb{R}^m$ 是开集, $u\in\mathcal{D}'(\Omega')$ 是 $\Omega'$ 上的分布, $f : \Omega\to\Omega'$.我们有时可以定义\newterm{pull-back}{回拉}{pull-back} $f^*u\in\mathcal{D}'(\Omega)$.可能存在障碍;例如, 若 $f$ 是常值函数, 则 $u\circ f$ 只对连续函数 $u$ 有意义;更一般地, 若 $f$ 把某个正测度集映到零测度集, 则对只在忽略零测度集的意义下良定义的 $u\in L^p$, $u\circ f$ 甚至没有意义.为确保上面提到的障碍都不出现, 我们假设 $f$ 光滑, 且其导数矩阵 $Df$ 在每一点都满秩, 从而至少任何开集都映到另一个开集.

为在可能时一致地定义 $u$ 经 $f$ 的回拉, 我们使用对偶性, 把函数的回拉视为测度\newterm{push-forward}{前推}{push-forward}的伴随运算, 并令

\[
\langle f^*u, \phi\rangle = \langle u, f_{\#}\phi\rangle
\]

其中 $f_{\#}\phi$ 是(有限复)测度 $\phi\, dx$ 在映射 $f$ 下的前推(的密度函数).

我们将在后面证明 $f_{\#}\phi$ 是光滑函数, 从而回拉是良定义的.显然, $f^*u$ 关于分布 $u$ 对弱极限连续.这一连续性的一个直接后果是, 光滑函数 $u$ 与 $f$ 的链式法则可以推广到 $u$ 是分布的情形.例如,

\[
\nabla u(f(x)) = u'(f(x))\nabla f(x)
\]

按分布的意义成立.

上述定义当然与测度论的形式体系一致.但是, 前推测度是否具有一个可以作为合法检验函数的密度函数并不是显而易见的, 如何用这一定义进行计算也不清楚, 所以让我们先明确讨论几个具体例子.

\begin{example*}
当 $f : \Omega\to\Omega'$ 是具有逆映射 $g$ 的 $\mathcal{C}^{\infty}$ 微分同胚时, 我们对 $y = f(x)$ 应用熟悉的换元公式:

\[
\begin{aligned}
\int u(f(x))\phi(x)\, dx &= \int u(f(x))\phi(g(f(x)))\, dx \\
&= \int u(y)\phi(g(y))\cdot|\det Dg(y)|\, dy
\end{aligned}
\]

其中 $Dg$ 是 $g$ 的导数矩阵.这一计算启发我们对这样的微分同胚定义回拉:

\[
\langle f^*u, \phi\rangle = \langle u, |\det Dg(y)|\phi\circ g\rangle, \quad \phi\in\mathcal{C}_0^{\infty}(\Omega)
\]

在这种情况下我们看到, 换元公式等价于回拉的定义.
\end{example*}

\begin{example*}
若 $f : \Omega\to\mathbb{R}$ 具有非零梯度, 则我们可以显式地得到 delta 函数 $\delta_t$ 的回拉, 即 $f^*(\delta_t) = \frac{1}{|\nabla f|}d\sigma$.这里, $d\sigma$ 表示嵌入超曲面 $f^{-1}(t) = \{f(x) = t\}\subset\mathbb{R}^n$ 上的标准曲面测度, $\nabla f$ 表示 $f$ 的梯度.

换言之, 我们可以计算前推测度的密度函数在 $t$ 处的值:

\[
f_{\#}\phi(t) = \int_{f^{-1}(t)}\phi(x)\frac{d\sigma(x)}{|\nabla f|},
\]

从而不仅可以对 delta 函数, 而且可以对任意分布 $u\in\mathcal{D}'(\mathbb{R})$ 计算 $\langle f^*u, \phi\rangle = \langle u, f_{\#}\phi\rangle$.这样, 回拉公式可以非正式地写成某种分解:

\[
u(f(x)) = \int u(t)\delta(f(x)-t)\, dt
\]

它可以由恒等式 $u(y) = \int u(t)\delta(y-t)\, dt$ 形式地导出.

作为这一公式的一个应用示例, 可以看到半径为 $R$ 的球的体积的导数就是半径为 $R$ 的球面的面积, 这得益于 $\nabla|x| = \frac{x}{|x|}$ 的梯度具有范数 1 以及如下微分:

\[
\frac{d}{dr}\int H(r-|x|)\, dx = \int\delta(r-|x|)\, dx.
\]

这个公式在几何上是清楚的:当我们比较半径为 $r$ 的球的体积与稍大的半径为 $r+\varepsilon$ 的球的体积时, 体积的变化本质上等于 $\varepsilon$ 乘以表面积.
\end{example*}

由于 delta 函数的回拉对我们非常重要, 让我们给出这个公式的证明.一旦我们证明了这一公式, 我们建立的理论就足以完整地完成前面的计算.我们将对分布 $H(r-\cdot)$ 关于变量 $r$ 取差商, 这些差商本质上支集在一个"加厚的球面"上.然后我们使用这个平凡的观察:分布的回拉 $\langle f^*u, \phi\rangle = \langle u, f_{\#}\phi\rangle$ 关于 $u$ 对弱极限连续.

现在让我们证明 delta 函数的回拉公式.证明中的几何图像基本上是前面考虑的特殊情形 $f(x) = |x|$ 的推广.

\mtsubsecn{}{回拉公式的证明($m=1$)}

如有必要, 通过取单位分解分解 $\phi$, 我们可以假设 $f$ 可以被补全为 $\phi$ 的支集上的坐标系, 这由于 $|\nabla f|$ 不消失而在 $\phi$ 的支集中任何一点的小邻域内总是可能的, 并且有限多个这样的邻域就足够了.我们考虑测度 $\mu = \phi(y)\, dy$, 并设 $\Psi(t) : \mathbb{R}\to\mathbb{C}$ 是由下式定义的分布函数:

\[
\Psi(t) = f_{\#}\mu(-\infty, t] = \int\phi(y)\cdot(f(y)\leq t)\, dy.
\]

我们现在希望证明

\[
\Psi'(t) = \lim_{\varepsilon\to 0}\frac{1}{\varepsilon}\int(t<f(y)\leq t+\varepsilon)\cdot\phi(y)\, dy
\]

在每一点 $t\in\mathbb{R}$ 都存在, 由此将推出\footnote{当 $\Psi$ 连续且在每一点 $t$ 处经典可微时, 微积分基本定理适用.} $\Psi$ 绝对连续, 并且由该公式给出的 $\Psi'(t) = f_{\#}\phi(t)$ 确实是正确的密度函数.为记号简单, 让我们假设 $\varepsilon>0$.

我们现在通过换坐标验证:在点 $y_0\in f^{-1}(t)$ 的非常近旁, 加厚超曲面 $\{y \mid t<f(y)\leq t+\varepsilon\}$ 可以被参数化为具有"高度" $\frac{\varepsilon}{|\nabla f|(y_0)} + o(\varepsilon)$ 和"宽度" $\sim d\sigma_{f^{-1}(t)}(y_0)$, 这至少从一般情形的图像来看是直观上可信的(例如, 在前面的例子 $f(x) = |x|$ 中).

我们可以(不失一般性地)假设 $\frac{\partial f}{\partial y^1}(y_0)\neq 0$, 并考虑满足

\begin{equation}
f(h(x), x') = x^1\label{eq:c3-135}
\end{equation}

的光滑函数 $h(x^1, \ldots, x^n)$, 对包含 $\phi$ 的支集的 $(t = f(y^0), y_0^1, \ldots, y_0^n)$ 的邻域中的所有 $\boldsymbol{x} = (x^1, x^2, \ldots, x^n) = (x^1, x')$.然后我们本质上用 $f$ 代替 $y^1$ 作为坐标, 作坐标变换 $(y^1, y^2, \ldots, y^n) = (h(x), x^2, \ldots, x^n)$, 得到:

\[
\begin{aligned}
&\lim_{\varepsilon\to 0}\int\left[\frac{1}{\varepsilon}(t<f(y)\leq t+\varepsilon)\cdot\phi(y)\right]|dy^1\wedge dy^2\wedge\ldots\wedge dy^n| \\
&= \lim_{\varepsilon\to 0}\int\left[\frac{1}{\varepsilon}(t<x^1\leq t+\varepsilon)\cdot\phi(h(x), x')\right]|dh\wedge dx^2\wedge\ldots\wedge dx^n| \\
&= \lim_{\varepsilon\to 0}\int\left[\frac{1}{\varepsilon}(t<x^1\leq t+\varepsilon)\cdot\phi(h(x), x')\right]\left|\frac{\partial h}{\partial x^1}(x)\right||dx^1\wedge dx^2\wedge\ldots\wedge dx^n| \\
&= \lim_{\varepsilon\to 0}\int\left[\frac{1}{\varepsilon}(t<x^1\leq t+\varepsilon)\cdot\left(\phi/\left|\frac{\partial f}{\partial y^1}\right|\right)\circ(h(x^1, x'), x')\right]dx^1\ldots dx^n \\
&= \int\left(\phi/\left|\frac{\partial f}{\partial y^1}\right|\right)\circ(h(t, x'), x')\, dx'
\end{aligned}
\]

为计算第二行中变换的 Jacobi 行列式, 我们使用了微分形式的简写, 它快速概括了这一事实:$n$ 维平行六面体的体积在一个顶点沿其他顶点张成的空间平移时保持不变, 即通过恒等式

\[
\begin{aligned}
dh\wedge dx^2\wedge\ldots\wedge dx^n &= \left(\frac{\partial h}{\partial x^1}dx^1\right)\wedge dx^2\wedge\ldots\wedge dx^n \\
&\quad + \left(\sum_{2}^n\frac{\partial h}{\partial x^i}dx^i\right)\wedge dx^2\wedge\ldots\wedge dx^n \\
&= \frac{\partial h}{\partial x^1}dx^1\wedge\ldots\wedge dx^n
\end{aligned}
\]

而 $\frac{\partial h}{\partial x^1}$ 通过对定义 $h$ 的隐式方程 \eqref{eq:c3-135} 隐式求导来计算.现在显然 $\Psi' = f_{\#}\phi$ 是 $t$ 的光滑函数.

方程 \eqref{eq:c3-135} 还表明, 对固定的 $t$, 函数 $\psi_t(x') = h(t, x')$ 把超曲面 $f^{-1}(t)$ 参数化为函数 $x^1 = \psi_t(x')$ 的图, 其中 $x'$ 变化.我们现在希望用曲面测度来解释 $f^{-1}(t)$ 上的积分:

\[
d\sigma_{f^{-1}(t)}(x') = \left(1+\sum_{2}^n\left(\frac{\partial\psi_t}{\partial x^j}\right)^2\right)^{1/2}dx'
\]

于是我们通过对 $f(\psi_t(x'), x') = t$ 隐式求导来计算曲面密度, 得到

\[
\frac{\partial f}{\partial y^1}\cdot\frac{\partial\psi_t}{\partial x^j} + \frac{\partial f}{\partial y^j} = 0 \quad j = 2, \ldots, n
\]

从而我们看到,

\[
|\nabla f| = \sqrt{\sum\left(\frac{\partial f}{\partial y^i}\right)^2} = \left|\frac{\partial f}{\partial y^1}\right|\cdot\left(1+\sum_{2}^n\left(\frac{\partial\psi_t}{\partial x^j}\right)^2\right)^{1/2}
\]

代入 $\int(\phi/|\frac{\partial f}{\partial y^1}|)\circ\psi_t(x')\, dx'$ 就给出上面例 2 中的公式.

上述证明本可以通过使用相同的变量替换 $(y^1, \ldots, y^n) = (h(x), x^2, \ldots, x^n)$ 直接证明

\[
\int u(f(y))\phi(y)\, dy = \int u(t)f_{\#}\phi(t)\, dt
\]

对所有 $u\in\mathcal{C}^{\infty}(\mathbb{R})$ 成立来简化, 或者等价地用光滑的近似 delta 函数代替尖锐的近似 delta 函数 $\frac{1}{\varepsilon}(0<t\leq\varepsilon)$.我们之所以选择上面这个较长的证明, 是因为它具有直观的几何吸引力, 也为了演示如何使用分布函数计算 $f_{\#}\phi$.更简短、更一般的证明连同一些计算工具一起在附录中给出.

\textbf{习题 1.} 设 $S_{\lambda} : \mathbb{R}^n\to\mathbb{R}^n$ 是伸缩映射 $S_{\lambda}(x) = \lambda x$.我们说分布 $u\in\mathcal{D}'(\mathbb{R}^n)$ 是 $a$ 次\newterm{homogeneous}{齐次的}{homogeneous}, 若 $S_{\lambda}^*u = \lambda^a u$.证明当 $u$ 是函数时, 该定义与通常的定义一致.证明在 $\mathbb{R}^n$ 中, $\delta_0$ 是 $-n$ 次齐次的.

\textbf{习题 2.} 证明 $\mathbb{R}^n$ 中任何既 $-n$ 次齐次又支集在原点的分布都是 $\delta_0$ 的常数倍.

上面的例子是一个更一般公式的特殊情形.当 $f : \Omega\to\mathbb{R}^m$ 是其导数处处满射的光滑映射时, 我们可以用下面的显式公式计算 $f_{\#}\phi$:

\begin{equation}
(f_{\#}\phi)(y) = \int_{f^{-1}(y)}\phi(x)\frac{d\sigma(x)}{\|f^*\omega\|(x)}\label{eq:c3-136}
\end{equation}

这里, $d\sigma$ 表示余维数 $m$ 的子流形 $f^{-1}(y)\subseteq\Omega$ 上的诱导测度, $f^*\omega = df^1\wedge\ldots\wedge df^m$ 是 $\mathbb{R}^m$ 上体积形式 $\omega = dy^1\wedge\ldots\wedge dy^m$ 的回拉, $\|\cdot\|$ 表示由 $m$ 形式上的逐点内积诱导的范数.测度 $d\sigma$ 也可以写成 $\frac{|*f^*\omega|}{\|f^*\omega\|}$, 其中 $*f^*\omega$ 是 $f^*\omega$ 的 Hodge 对偶.这些概念都将在附录中进一步回顾, 其中一般公式 \eqref{eq:c3-136} 以不同于上面证明的 $m=1$ 情形的方式得到证明.证明由分解体积形式

\[
dx^1\wedge\ldots\wedge dx^n = \frac{1}{\|f^*\omega\|^2}(f^*\omega\wedge*f^*\omega)
\]

然后先在 $f$ 的水平集上积分组成.

存在更一般的条件, 使得当分布 $u$ 的奇性被更精确地理解时, 回拉运算 $f^*u$ 可以进行.

\mtsubsecn{}{应用}

这些运算的第一个应用将是证明 Gauss 散度定理, 它把函数 $\phi$ 的某个导数在集合 $\Omega$ 内部的积分表示为 $\phi$ 的边界积分.可以把 $\Omega$ 上任何导数的积分表示为 $\phi$ 的差商积分的极限;但对偶地, 可以固定函数 $\phi$, 对 $\Omega$ 的\newterm{characteristic-function}{特征函数}{characteristic function}取伴随差商.于是, 散度定理等价于开集特征函数的微分.

我们现在计算具有光滑边界的区域 $\Omega$ 的特征函数 $\chi_{\Omega}$ 的梯度, 但让我们先用语言来描述.可以把 $\chi_{\Omega}$ 的图像想象成一张 $\Omega$ 形的桌子.如果我们略微扰动 $\chi_{\Omega}$ 以得到光滑逼近 $\tilde{\chi}_{\Omega}$, 显然 $\tilde{\chi}_{\Omega}$ 的梯度指向 $\Omega$ 内部且垂直于边界 $\partial\Omega$(沿最大增大的方向), 而且梯度本质上仍然支集在 $\Omega$ 的边界上.因此, 毫不奇怪, 我们得到

\begin{proposition}[Gauss 散度定理]\label{prop:c3-4.1}
对任何检验函数 $\phi\in\mathcal{C}_0^{\infty}(\mathbb{R}^n)$,

\[
-\int_{\Omega}\nabla\phi(x)\, dx = \int\nabla\chi_{\Omega}(x)\phi(x)\, dx = \int_{\partial\Omega}\phi(x)\vec{n}\, d\sigma_{\partial\Omega}(x),
\]

其中 $\vec{n}$ 表示内向单位法向量.
\end{proposition}

\begin{proof}
通过取单位分解, 我们可以分解 $\phi = \sum_{\alpha}\phi_{\alpha}$, 其中每个 $\phi_{\alpha}$ 是支集在小区 $V_{\alpha}$ 内的光滑函数, 使得 $V_{\alpha}\cap\Omega = \{x\in V_{\alpha} : f_{\alpha}(x)>0\}$ 是某个定义函数 $f = f_{\alpha}$ 的上水平集, 且该函数在边界上具有非零梯度.在 $V_{\alpha}$ 中, 我们(按分布的意义)有等式 $\chi_{\Omega} = H\circ f_{\alpha}$, 即

\[
\int\chi_{\Omega}(x)\phi(x)\, dx = \int H(f_{\alpha}(x))\phi(x)\, dx
\]

对任何支集在 $V_{\alpha}$ 中的检验函数 $\phi$ 成立.因此, 由链式法则,

\[
\nabla\chi_{\Omega}(x) = \delta(f_{\alpha}(x))\nabla f_{\alpha}(x) = \frac{\nabla f_{\alpha}}{|\nabla f_{\alpha}|}(x)d\sigma_{\partial\Omega}(x) = \vec{n}\, d\sigma_{\partial\Omega}(x)
\]

这证明了命题.
\end{proof}

我们可以略微修改证明, 以允许对 $f$ 宽松得多的正则性条件:首先逼近 $\Omega$ 的特征函数并取极限(例如让一族前 Heaviside 函数 $H_{\varepsilon}$ 收敛到 Heaviside 函数).例如, 当边界可以局部表示为 Lipschitz 函数的图时, 法向量几乎处处良定义, 我们以此方式对更多类型的集合(如多边形、立方体等)得到同样的公式.细节可在 Hörmander 的著作\cite{ho}中找到.

推广这些想法, 我们可以概述更一般的 Stokes 定理的一个快速证明.证明本质上如下:对定向 $k$ 维流形 $Y$(它可能嵌入到更高维流形 $X$ 中)、$X$ 上的 $k-1$ 形式 $\omega$, 以及检验函数 $\phi\in\mathcal{C}_0^{\infty}(X)$, 我们有

\[
0 = \int_Y d(\phi\omega) = \int_Y\phi\, d\omega + \int_Y d\phi\wedge\omega
\]

当 $\phi$ 的支集在一个坐标片内时这是显然的.若 $M\subseteq Y$ 是 $Y$ 的具有光滑边界 $\partial M$ 的开子集, 我们取 $\phi$ 逼近 $M$ 的特征函数(从而 $d\phi$ 的支集在 $\partial M$ 上), 在极限中得到 Stokes 定理:

\[
\int_M d\omega = \int_{\partial M}\omega
\]

分布理论向流形情形的推广大体上是直接的, 并在 Hörmander 的著作中有所概述.

讲义末尾关于子流形上积分的附录可能对下面的一些习题有帮助.

\textbf{习题 3.} 证明:若 $f$、$g$ 是 $\mathbb{R}^n$ 上处处具有非零微分的两个光滑函数, 则对所有 $a, b\in\mathbb{R}^n\times\mathbb{R}^n$:

\[
\int\delta_0(f(a)-x)\delta_0(g(b)-x)\, dx = \delta_0(f(a)-g(b)).
\]

提示:两边都应理解为 $\mathbb{R}^n\times\mathbb{R}^n$ 上的分布.可以把回拉的定义改写成 $u(g(b)) := \int u(x)\delta_0(g(b)-x)\, dx$ 的形式.用近似 delta 函数逼近, 我们可以延拓到 $u(x) = \delta_0(f(a)-x)$ 的情形.或者, 利用 $f(a) = a$ 与 $g(b) = b$ 均为恒等映射的明显特殊情形, 并对一般的 $f$ 与 $g$ 取回拉.

\textbf{习题 4.} 从去心正方形

\[
S = \{(x, y) : \max(|x|, |y|)<1\}\setminus\{(0, 0)\}
\]

中随机抽取一点.随机变量 $x\cdot y$ 的概率密度是什么?

\textbf{习题 5.} 设 $f : \mathbb{R}^3\to\mathbb{R}^2$ 是映射 $(\tau, \beta) = f(x, y, z) = (x^2+y^2+z^2, z)$, 它在直线 $x=y=0$ 之外非奇异.对 $\phi\in\mathcal{C}_0^{\infty}(\mathbb{R}^3\setminus\{x=y=0\})$, 直接证明(例如通过计算分布函数)前推测度 $f_{\#}\phi$ 由

\[
\left(\int_0^{2\pi}\phi(\tau, \beta, \theta)\, d\theta\right)\cdot(\tau>\beta^2)\frac{|d\tau\wedge d\beta|}{2}
\]

给出, 这里 $(\tau, \beta, \theta)$ 是 $\mathbb{R}^3$ 上的坐标, 其中 $\theta$ 是 $x, y$ 平面中的极角.验证公式 \eqref{eq:c3-136} 给出同样的结果.

\textbf{习题 6.} 证明:若 $\delta_0$ 是 $\mathbb{R}$ 上的 Dirac delta 函数, 则把它视为 $(t, x)\in\mathbb{R}\times\mathbb{R}^3\setminus\{(0, 0)\}$ 上的分布时, 有

\[
\langle\delta_0(t^2-|x|^2), \phi\rangle = \int_{\mathbb{R}^3}(\phi(|x|, x)+\phi(-|x|, x))\frac{dx}{2|x|}
\]

为什么这个公式作为 $\mathbb{R}^{3+1}$ 全体上的分布有意义, 即使导数 $d(t^2-|x|^2) = 2t\, dt - \sum_{i=1}^3 2x_i\, dx^i$ 在原点消失?为什么这是把 $\delta_0(t^2-|x|^2)$ 延拓到整个 $\mathbb{R}^{3+1}$ 并保持 $-2$ 次齐次性的唯一可能定义?

\mtsubsec{其他未讨论的主题}

我们没有包括关于分布相乘的一节, 因为一般说来, 无法定义分布的、有意义且结合性的乘积, 使其连续地延拓函数的通常乘法.(试构造直线上三个分布的例子, 使它们的乘积(若可定义)不能结合, 或者两个分布的例子, 使它们的乘积不能交换.)

相乘奇性不相交的分布没有困难——只需用光滑单位分解在空间中局部化即可.当奇性发生在相同点时, 问题就更微妙了.当奇性"不碰撞"时(由\newterm{wavefront-set}{波前集}{wavefront set}的概念精确化, 该概念在 $\mathbb{R}^n$ 的\newterm{cotangent-bundle}{余切丛}{cotangent bundle}中度量分布奇性的位置与方向), 相乘仍然可能——在这种情况下, 必须尝试同时在空间与"频率"上局部化以形式化乘法.例如, 可以证明 $\delta_0(x_1, x_2) = \delta(x_1)\cdot\delta(x_2)\in\mathcal{D}'(\mathbb{R}^2)$.我们将在讲义的多处提到波前集的概念, 但对这个集合究竟是什么保持模糊;关于波前集与奇性传播的严格讨论, 见(∗∗∗).对我们的目的而言, 记住奇性既有位置也有"方向"就够了.

为了给出最一般情形的更明确提示:给定两个分布 $u(x)$ 与 $v(x)\in\mathcal{D}'(X)$, 以显然的方式定义\newterm{tensor-product}{张量积}{tensor product} $u(x_1)v(x_2)\in\mathcal{D}'(X\times X)$, 而乘法(当可能时)是张量积沿对角嵌入 $X\hookrightarrow X\times X$ 的回拉 $u(x)v(x)$.不过要注意, 对角嵌入不满足我们最初定义分布回拉时假设的条件, 因为它的导数不是满射.即使在这种情形下, 有时也可以定义回拉, 但同样必须满足某些涉及该映射与被回拉分布的波前集的条件, 因为很清楚, 例如, 并非所有分布都能限制到更低维的子集上.我们将在后面涉及基本解的一些计算中, 以及研究\newterm{trace-theorem}{迹定理}{trace theorem}(它使我们能够在处理 PDE 中的某些广义函数时理解"边界值"的意义)和 傅里叶变换的限制定理时, 再次面对这个问题.

% ==== part04 第 4 章:基本解与基本线性偏微分方程 ====
% 源: MinerU_latex_lecture nots.tex 5211-6715 行;校对 pdf_text.txt 4043-5551 行(PDF 第 77-104 页)
% 术语(本部分新增,中文=英文):
%   调和函数=harmonic function(键 harmonic-function);次调和=subharmonic(键 subharmonic);
%   超调和=superharmonic(键 superharmonic);Green 函数=Green's function(键 green-function);
%   椭圆算子=elliptic operator(键 elliptic-operator);随机游走=random walk(键 random-walk);
%   前向基本解=forward fundamental solution(键 forward-fundamental-solution);
%   共面积公式=co-area formula(键 coarea-formula);弱解=weak solution(键 weak-solution)
%   其余沿用已有术语: 基本解/分布/检验函数/卷积/示性函数/Heaviside 函数/能量-动量张量/
%   波算子/程函方程/后向光锥等(已在 part00/01/03 引入,本部分直接使用不重复 \newterm)
% REFS(需协调员统一解析):
%   - 源 6154 行 "the homogeneous family of distributions j_a defined by proposition ??": 原书
%     "proposition ??" 引用未解析(疑指引言或第 2/3 章中定义齐次分布族 j_a 的命题),正文按原文
%     写 "命题 ??",待补链.
%   - 源 6182 行 "See Exercise 3 of section (4) for the n = 3 case": 原书指 "第 4 节的习题 3"
%     (n=3 情形),本部分第 4 节习题 3 为热算子最大值原理,与原意不符;疑指其他章的节,待解析.
%   - 源 6124 行 "We will return to the study of the Laplace equation ... in Chapter 3.":
%     指第 3 章(该章为分布/Fourier 相关),正文写 "第 3 章".
%   - 源 5975 行 Remark 2.6 "the measure in Exercise 5 of the Introduction": 指引言(part00)习题 5.
%   - 源 5835 行 "introduced formally in the exercises in the Introduction": 指引言习题中引入的 Green 函数.
%   - 源 6088 行 "is proven in the Appendix": 指附录 A(co-area 公式),正文写 "附录".
%   - 源 6397 行 "as discussed in subsection 4, Example 2": 指分布回拉部分(第 3 章?)的小节 4 例 2,待解析.
%   - 源 6325 行 "see [J]": 文献 [J](疑为 John 的著作),保留 [J].
%   - 源 6027 行 "Evans or Gilbarg and Trudinger": 文献名,保留.
%   - 公式 (142) 与验证式(6372 行)符号不一致(原书如此): (142) 为正号 E_+^{(1+3)}=
%     1/(2π)H(t)δ(t²-|x|²),验证式写 -1/2 π^{-1}H(t)δ_0(t²-|x|²)=-1/(4π)r^{-1}δ(t-r),照原书保留.
%   - 源 6254 行 "deduce (148) from (142)": 原文如此((148) 为 E=Aδ(u)),按数学意义疑应为
%     "(143) 从 (142)",照原文保留 \eqref 并在此注明.
%   - 源 6245 行 公式 (143) 右端以悬空 "=" 结束(两处 OCR 一致),按原书保留悬空等号.
%   - 源 6420 行 "2∂_α A ∂_β u + □u = 0": OCR 索引混乱,按数学含义与脚注 12(输运方程)还原为
%     2∂_α A ∂^α u + □u = 0.
%   - 源 6172 行 "E_+^{(2k+1+1)}" 保留原书写法(即 n=2k+1 维空间情形的 n+1=2k+2 维).
%   - FOOTNOTE: 脚注 1-14 全部已从 pdf_text.txt 对应页找回并翻译.
% 编号说明:
%   - 引言无编号(\section*{引言}),第一个编号对象为 Proposition 0.3,故先 \setcounter{mtthm}{2};
%   - 第 1 节: 命题 1.1、定理 1.2、推论 1.3、定理 1.4/1.5/1.6(习题为加粗段落,不推进计数器);
%     未编号 Remark 用 remark*;
%   - 第 2 节: 命题 2.1、2.2、定理 2.3、命题 2.4、定理 2.5、注 2.6、命题 2.7、定理 2.8、命题 2.9;
%     未编号 Remark(ODE 猜测)、Note、Example 用 remark*/example*;习题 1-8 为加粗段落;
%   - 第 3 节: 定理 3.1、3.2、命题 3.3、定理 3.4;习题 1-10 为加粗段落;
%   - 第 4 节: 无编号定理对象;习题 1-5 为加粗段落;4.1 用 \mtsubsecn{4.1}{...}(不推进计数器),
%     其下习题 1、2 亦为加粗段落;
%   - 编号公式共 18 个: (137)-(154),均用 \begin{equation}\label{eq:c4-NNN} 按原书编号;
%     其余行间公式用 \[ \] 不编号.
% ====

\chapter{基本解与基本线性偏微分方程}

\section*{引言}

在引言中， 我们以 $\mathbb{R}^{3}$ 中 Laplace 算子的特殊情形介绍了基本解的基本概念， 并利用该解推导了具有紧支集数据的 Poisson 方程的若干基本事实.现在我们开始从同一视角、但更深入地研究几个基本的线性偏微分方程.

给定具有常系数的线性偏微分算子

\[
P(\partial) = \sum_{|\alpha| \leq k} a_{\alpha}\partial^{\alpha}
\]

其中 $a_{\alpha} \in \mathbb{C}$.我们说一个分布 $E$ 是基本解， 如果它满足 $P(\partial)E = \delta_{0}$.若如此， 则对形如 $P(\partial)u = f$ 的方程（其中 $f \in \mathcal{D}'(\mathbb{R}^{n})$ 是紧支集分布）， 我们总可以通过令 $u = E*f$ 来找到解.这容易由如下事实推出：对任意 $u \in \mathcal{D}'(\mathbb{R}^{n})$ 有 $\delta_{0}*u = u$；再结合下面的命题（它最终源于 $\mathbb{R}^{n}$ 中所有平移可交换这一事实）.

\setcounter{mtthm}{2}
\begin{proposition}\label{prop:c4-0.3}
假设 $u, v \in \mathcal{D}'(\mathbb{R}^{n})$， 其中至少一个具有紧支集.则

\[
P(\partial)(u*v) = P(\partial)u*v = u*P(\partial)v.
\]
\end{proposition}

可以证明如下一般结果.

\begin{theorem}[Ehrenpreis--Malgrange]\label{thm:c4-0.4}
$\mathbb{R}^{n}$ 上任意具有常系数的线性偏微分算子 $P(\partial)$ 都有基本解.
\end{theorem}

其证明涉及初等的 Fourier 与泛函分析， 实际上与本讲义的关系较为疏远（不过有一道期中习题表明， 在更一般的背景下类似的定理可能存在障碍）.归根结底， 基本解必须足够显式， 才能用于导出所讨论方程的有趣性质.

我们转而处理一些重要的、平移不变的线性微分算子的具体例子， 它们具有特殊而有用的不变性质.凭借这些算子的简单性与特殊对称性， 我们能够导出基本解的显式公式.这使我们能够导出相应方程的重要定性性质（解的存在性与正则性、对数据的连续依赖性等）.然而， 这些定性性质往往与各种对称性无关， 并且显然在算子被轻微甚至大幅度扰动后仍然成立.因此请记住：尽管显式公式非常有用， 但为了理解 PDE 的性质， 我们最终需要发展更稳健的技巧.

\section{Cauchy-Riemann 方程组}

算子\footnote{记住这一定义的简便方法， 是把 $f$ 写成 $f(x, y) = f\!\left(\frac{z+\bar{z}}{2}, \frac{z-\bar{z}}{2i}\right)$.还要注意 $\frac{\partial f}{\partial z} = \frac{1}{2}\!\left(\frac{\partial}{\partial x} - i\frac{\partial}{\partial y}\right)$， 且 $df = \frac{\partial f}{\partial z}dz + \frac{\partial f}{\partial \bar{z}}d\bar{z}$.}
$\frac{\partial}{\partial \bar{z}} = \frac{1}{2}\!\left(\frac{\partial}{\partial x} + i\frac{\partial}{\partial y}\right)$ 是复分析的基础；复分析研究线性偏微分方程
$\frac{\partial f}{\partial \bar{z}} = 0$ 的解 $f : \mathbb{C} \to \mathbb{C}$.这样的函数称为全纯函数， 它们合在一起构成 $\mathbb{C}$ 上的一个代数， 此外还具有许多其他引人入胜的性质.联系 $f$ 的实部与虚部的这对偏微分方程称为 Cauchy-Riemann 方程组.我们假定读者已经熟悉复分析这一学科， 并进而通过应用分布理论来建立该学科的基本事实.在本节中， 我们将使用微分形式的表述， 特别地， 用 $\frac{\mathrm{d}\bar{z} \wedge \mathrm{d}z}{2i} = dx \wedge dy$ 表示 $\mathbb{C}$ 上的体积形式.给定函数 $f(z, \bar{z})$， 我们用复记号写成

\[
df = \frac{\partial f}{\partial z}dz + \frac{\partial f}{\partial \bar{z}}d\bar{z}.
\]

根据我们的一般定义， $\frac{\partial}{\partial \bar{z}}$ 的基本解是 $\mathbb{R}^{2}$ 中的一个分布 $K$， 满足 $\frac{\partial}{\partial \bar{z}}K = \delta_{0}$.与 Laplace 算子不同， $\frac{\partial}{\partial \bar{z}}$ 作为算子不与旋转 $u(z) \mapsto \tilde{u} = u(e^{it}z)$ 交换.事实上， 由计算

\[
\begin{aligned}
d\bigl[u(e^{it}z)\bigr] &= \frac{\partial u}{\partial z}(e^{it}z)\, d(e^{it}z) + \frac{\partial u}{\partial \bar{z}}(e^{it}z)\, d\bigl(\overline{e^{it}z}\bigr) \\
&= e^{it}\frac{\partial u}{\partial z}(e^{it}z)\, dz + e^{-it}\frac{\partial u}{\partial \bar{z}}(e^{it}z)\, d\bar{z}
\end{aligned}
\]

可得 $\frac{\partial}{\partial \bar{z}}\tilde{u} = e^{-it}\frac{\partial u}{\partial \bar{z}}(e^{it}z)$.

然而， 如果已经有一个基本解 $u$（满足 $\frac{\partial u}{\partial \bar{z}} = \delta_{0}$）， 那么由上面的计算可以构造出另一个基本解 $e^{it}u(e^{it}z)$， 因为 $\delta$-函数是旋转不变的.通过在旋转群上取平均， 我们可以不失一般性地假设这两个基本解相同， 即 $K(e^{it}z) = e^{-it}K(z)$， 这促使我们寻找如下形式的基本解

\[
K(re^{i\theta}) = g(r)e^{-i\theta}.
\]

由于 $\delta_{0}$ 是 $-2$ 次齐次的， 而 $\frac{\partial}{\partial \bar{z}}$ 使齐次度降低 1， 我们猜想 $g(r) = \frac{c}{r}$（其中 $c \in \mathbb{C}$ 为某常数）， 从而基本解将是 $-1$ 次齐次的.于是我们猜测 $K(z) = \frac{c}{z}$（$c$ 为某常数）.事实上， 由于 $\frac{1}{z}$ 是局部可积的， 它在 $\mathbb{R}^{2}$ 中处处定义一个分布， 且 $\frac{\partial}{\partial \bar{z}}\frac{1}{z}$ 的支集在原点.此外， 由于 $\frac{\partial}{\partial \bar{z}}\frac{1}{z}$ 是 $-2$ 次齐次的， 由支集为单点的分布的刻画我们推断它必为 $\delta_{0}$ 的常数倍；即 $\frac{\partial}{\partial \bar{z}}\frac{1}{z} = C\delta_{0}$， 其中 $C \in \mathbb{C}$ 为某常数（可能为 0）.

我们可以把该分布作用到任意我们想要的检验函数上来确定这个常数， 我们将选取单位圆盘 $D = \{|z| \leq 1\}$ 的示性函数作为检验函数.严格说来， 这样做超出了我们所讲述的分布理论的范畴， 但我们可以毫无困难地证明其合理性：要点在于， 当一个因子略有奇性时， 另一个因子至少是连续的， 这使我们能够从光滑逼近过渡.用实变量记号， 可以通过分部积分来计算

\[
\int H(1-|z|) \cdot \frac{\partial}{\partial \bar{z}}\frac{1}{z}\, dx\, dy = -\int \frac{1}{x+iy}\frac{\partial}{\partial \bar{z}}H(1-|x+iy|)\, dx\, dy.
\]

对 $\frac{\partial}{\partial \bar{z}}|z|^{2} = \frac{\partial}{\partial \bar{z}}(z\bar{z})$ 应用乘积法则， 可得 $\frac{\partial}{\partial \bar{z}}|z| = \frac{z}{2|z|}$， 从而

\[
-\int \frac{1}{x+iy}\frac{\partial}{\partial \bar{z}}H(1-|x+iy|)\, dx\, dy = \int \frac{1}{2|z|}\delta(1-|z|)\, dx\, dy = \pi.
\]

等价地， 我们可以应用 Stokes 定理\footnote{必须注意参数化 $\theta \mapsto e^{i\theta}$ 给出圆的正确定向.若天真地用顺时针参数化 $\theta \mapsto e^{-i\theta}$ 复现下面的计算， 将导致符号错误.}来计算

\[
\begin{aligned}
&\int_{D}\frac{\partial}{\partial \bar{z}}\left(\frac{1}{z}\right)\frac{\mathrm{d}\bar{z} \wedge \mathrm{d}z}{2i} = \int_{D}\mathrm{d}\left(\frac{1}{z}\frac{dz}{2i}\right) \\
&\qquad = \int_{\partial D}\frac{1}{z}\frac{dz}{2i} \\
&\qquad = \int_{0}^{2\pi}e^{-i\theta}\frac{ie^{i\theta}d\theta}{2i} \\
&\qquad = \pi
\end{aligned}
\]

由此得到命题：

\begin{proposition}\label{prop:c4-1.1}
设 $K(z) = \frac{1}{\pi}\frac{1}{z}$， 则 $\frac{\partial K}{\partial \bar{z}} = \delta_{0}$.
\end{proposition}

找到基本解后， 我们将立即得到全纯函数的表示公式.我们首先注意到， 上述计算的一个变体使我们能够计算 $\frac{\partial \chi_{\Omega}}{\partial \bar{z}}$：当 $\Omega$ 是带 Lipschitz 边界\footnote{我们证明散度定理时假设边界是光滑的， 但该证明的一个非常微小的变体对 Lipschitz 边界也成立.例如参见 Hörmander.}的开集， 且把 $\frac{1}{z}$ 换成任意检验函数时， 它表明

\[
\int \frac{\partial \chi_{\Omega}}{\partial \bar{z}} \cdot \phi(z)\frac{\mathrm{d}\bar{z} \wedge \mathrm{d}z}{2i} = -\int_{\partial\Omega}\phi(z)\frac{dz}{2i}, \qquad \phi \in \mathcal{C}_{0}^{\infty}(\mathbb{C}).
\]

例如， 这可由对恒等式

\[
d(\chi_{\Omega}\phi\, dz) = \left(\frac{\partial}{\partial \bar{z}}\chi_{\Omega}\phi + \chi_{\Omega}\frac{\partial}{\partial \bar{z}}\phi\right)d\bar{z} \wedge dz
\]

积分得到.验证这类涉及分布的积分恒等式成立的最简单方法之一， 是让奇异分布 $\chi_{\Omega}$ 被光滑函数 $\chi_{\Omega}^{\varepsilon}$ 逼近.在这种情况下， 由于 $\partial\Omega$ 是 Lipschitz 的而 $\phi$ 是连续函数， $\frac{\partial}{\partial \bar{z}}\chi_{\Omega}$ 是一个测度， 因此已经可以看出该计算是有效的.

在特殊情形 $\frac{\partial f}{\partial \bar{z}} = 0$（在 $\Omega$ 内）且 $f$ 在 $\overline{\Omega}$ 的某个邻域内是（比如说）$\mathcal{C}^{1}$ 的时， 我们可以从光滑逼近 $\phi_{n} \to f$ 过渡， 从而断定 $\int_{\partial\Omega}f(z)dz = 0$\footnote{更直接地， 可以计算 $d(f\, dz) = df \wedge dz = \frac{\partial f}{\partial z}dz \wedge dz + \frac{\partial f}{\partial \bar{z}}d\bar{z} \wedge dz = 0$， 然后在 $\Omega$ 上积分.}.

现在， 把 $\frac{\partial}{\partial \bar{z}}$ 作用到乘积

\[
\frac{\partial}{\partial \bar{z}}\left[\chi_{\Omega} \cdot \frac{1}{\pi(z-z_{0})}\right] = \delta_{z_{0}} + \frac{1}{\pi(z-z_{0})}\frac{\partial \chi_{\Omega}}{\partial \bar{z}}
\]

上， 得到一个支集仅含于 $z_{0}$ 与 $\Omega$ 边界的紧支集分布.

我们通过用 $\frac{\mathrm{d}\bar{z} \wedge dz}{2i}$ 作积分来理解这个等式， 得到恒等式：

\begin{theorem}\label{thm:c4-1.2}
设 $\Omega \subseteq \mathbb{C}$ 是带 Lipschitz 边界 $\partial\Omega$ 的有界区域， $z_{0} \in \Omega$， $\phi \in \mathcal{C}^{\infty}(\mathbb{C})$， 则

\[
\frac{1}{2\pi i}\int_{\partial\Omega}\frac{\phi(z)}{z-z_{0}}dz = \phi(z_{0}) + \frac{1}{2\pi i}\int_{\Omega}\frac{\partial \phi}{\partial \bar{z}}\frac{1}{z-z_{0}}d\bar{z} \wedge dz.
\]
\end{theorem}

我们只对光滑函数陈述了该定理， 但由逼近， 定理在远为一般的条件下成立.特别地， 当 $f$ 在 $\Omega$ 内全纯且在 $\overline{\Omega}$ 的某个邻域内是 $\mathcal{C}^{1}$ 时， 我们可以从光滑的 $\phi$ 过渡到 $f$， 由此得到如下推论：

\begin{corollary}[Cauchy 积分公式]\label{cor:c4-1.3}
设 $f$、$\Omega$ 与 $z_{0} \in \Omega$ 如上， 则

\[
\frac{1}{2\pi i}\int_{\partial\Omega}\frac{f(z)}{z-z_{0}}dz = f(z_{0}).
\]
\end{corollary}

\begin{remark*}
应用 Cauchy 积分公式并计算边界积分时必须多加小心.首先， 假设 $f$ 在边界上保持良好的性态， 这对于从光滑逼近过渡是至关重要的， 挖去原点的单位圆盘上 $\frac{1}{z}$ 的例子就说明了这一点.在这种情况下， Cauchy 积分公式不能应用于（比如说）非常接近 0 的点 $z_{0}$——边界积分 $\int_{\partial D}\frac{1}{z} \cdot \frac{1}{z-z_{0}}dz$ 的大小显然不会比弧长 $\int_{|z|=1}1\, |dz| = 2\pi$ 大多少.我们使用 Stokes 定理所涵盖的另一个重要问题是， 必须考虑边界的定向.
\end{remark*}

通过分析 Cauchy 积分公式可以证明， 全纯函数（在上述条件下）在 $\Omega$ 的任何内点处都具有收敛的幂级数展开， 因而特别是光滑的.更有用的是， 可以给出这种光滑性的定量刻画：对包含于 $\Omega$ 的紧集 $K$， 导出形如 $\|\partial^{\alpha}f\|_{L^{\infty}(K)} \lesssim \|f\|_{L^{\infty}(\partial\Omega)}$ 的估计.同样的估计还表明 $\frac{\partial f}{\partial \bar{z}} = 0$ 的解如何随其边界值连续变化（当解存在时）.这种解析性是更一般现象的一个例子：基本解在原点之外的正则性对应于 PDE 解的正则性.然而， 一般而言， 若能不依赖显式公式而得到这样的正则性结果（也许不必这么强）， 那就好了.我们很快就会重新讨论全纯函数.

\noindent\textbf{习题 1.}\quad 我们说 $u \in \mathcal{D}'(\Omega)$ 是 Cauchy-Riemann 方程组的\newterm{weak-solution}{弱解}{weak solution}， 如果 $\frac{\partial u}{\partial \bar{z}} = 0$ 在分布意义下成立.证明：一个连续函数若是弱解， 则它实际上是经典的全纯函数（因而是解析的）.（提示：全纯函数类在平移与线性组合下封闭， 因此考虑 $u$ 的一个磨光可能有用.然后利用先验估计.）

\noindent\textbf{习题 2.}\quad 由习题 1 推出：

\begin{theorem}[Schwartz 反射原理]\label{thm:c4-1.4}
设 $\Omega$ 是 $\mathbb{C}$ 的一个开子区域， 与实轴相交于一个区间 $I$.若 $f$ 在 $\Omega$ 上连续且在 $\Omega\backslash I$ 上全纯， 则 $f$ 在 $\Omega$ 上全纯.
\end{theorem}

提示：$f(x+iy) = \lim_{\varepsilon\to 0}\mathbf{1}_{\{|y|>\varepsilon\}} \cdot f(x+iy)$ 在弱意义下成立.

\noindent\textbf{习题 3.}\quad 再证明：

\begin{theorem}[Morera 定理]\label{thm:c4-1.5}
设 $f$ 在开圆盘 $D$ 中连续.则 $f$ 在 $D$ 中全纯， 当且仅当对任何边界 $\partial T$ 含于 $D$ 的直角三角形内部 $T$， 有

\[
\int \chi_{T}\frac{\partial f}{\partial \bar{z}}d\bar{z} \wedge dz \equiv \int_{\partial T}f(z)dz = 0.
\]
\end{theorem}

注意， 最左边的积分在经典意义下没有意义， 因为没有对 $\frac{\partial f}{\partial \bar{z}}$ 作任何正则性假设.（提示：这既可以利用关于 $\chi_{T}$ 的线性过渡到一般检验函数来证明， 也可以利用假设中的平移不变性来证明.）

\noindent\textbf{习题 4.}\quad Cauchy 积分公式立即推出如下形式的估计：$|f(z_{0})| \leq C\max_{z\in\partial\Omega}|f(z)|$， 其中正常数 $C = C(z_{0})$ 与 $f$ 无关.然而， 本质上由于 $\delta_{z_{0}} : f \to f(z_{0})$ 也是环同态， 我们能够选取与 $z_{0}$ 无关的 $C = 1$， 从而证明

\begin{theorem}[最大模原理]\label{thm:c4-1.6}
若 $\Omega$ 是闭包 $\overline{\Omega}$ 紧的开集， 且 $f$ 在 $\overline{\Omega}$ 的某个邻域内全纯， 则

\[
\max_{z\in\Omega}|f(z)| = \max_{z\in\partial\Omega}|f(z)|.
\]
\end{theorem}

这一思路的一种论证归功于 Landau.

另一方面， 仅由估计 $|f(z_{0})| \leq C\max_{z\in\partial\Omega}|f(z)|$ 即可推出（由 Hahn-Banach 定理）：最初对全纯函数的连续边界值定义的线性泛函 $\delta_{z_{0}} : f \mapsto f(z_{0})$ 可延拓为 $\mathcal{C}(\partial\Omega)$ 上的连续线性泛函， 因此由一个测度表示（即对在 $\Omega$ 内全纯且连续到边界的 $f$， 有 $f(z_{0}) = \int_{\partial\Omega}f(z)\, d\mu_{z_{0}}(z)$， 其中 $\mu_{z_{0}}$ 是边界上的某个有限测度， 先验上不一定唯一）.事实上， 我们已经通过证明 Cauchy 积分公式算出了这个测度.另一方面， 最大模原理可以在不使用 Cauchy 积分公式的情况下证明（习题部分请读者在下一节关于 Laplace 算子的习题 8 中给出这样的证明）.因此我们看到， 先验估计与积分表示公式的存在性是相辅相成的， 它们共同表达了解的唯一性以及对数据的连续依赖性.然而请注意：

\noindent\textbf{习题 5.}\quad 如果任意边界数据 $f \in \mathcal{C}(\partial\Omega)$ 都能被 $\Omega$ 中的全纯函数 $u$（$\frac{\partial u}{\partial \bar{z}} = 0$）实现， 那么根据上面的最大模原理， 连续泛函 $\delta_{z_{0}} : f \mapsto u(z_{0})$（$z_{0} \in \Omega$）将是定义在代数 $\mathcal{C}(\partial\Omega)$ 上的连续环同态.然而， $\mathcal{C}(\partial\Omega)$ 上任何如上延拓 $\delta_{z_{0}}$ 的连续线性泛函都不可能是环同态\footnote{事实上， 代数 $\mathcal{C}(\partial\Omega)$ 上所有这样的连续环同态都是边界本身上的点质量.}.因此， 存在不能实现为全纯函数边界值的连续函数.

因此， 我们不能仅由先验估计（如最大模原理）或表示公式就断定解存在.理解解的存在性的障碍是 PDE 中的一个有趣问题， 通常需要理解所讨论方程背后的几何或拓扑.例如， 就目前的情形而言， 全纯函数不能把一条闭曲线映为另一条定向相反的曲线.
\section{Laplace 算子 \texorpdfstring{$\Delta$}{Δ}}

正如我们在引言中所见， $\mathbb{R}^{n}$ 上的 Laplace 算子（或 Laplacian）$\Delta = \sum_{i=1}^{n}\partial_{i}^{2}$ 是最简单也是最重要的线性微分算子之一.$\Delta u = 0$ 的解称为\newterm{harmonic-function}{调和函数}{harmonic function}.在二维情形， $\Delta$ 与全纯函数的研究有关（例如， 由恒等式 $\Delta = 4\frac{\partial}{\partial \bar{z}}\frac{\partial}{\partial z}$ 以及我们前面的正则性结果可见， 全纯函数的实部与虚部都是调和的）.该算子也常常记作 $\nabla \cdot \nabla = (\partial_{1}, \partial_{2}, \ldots, \partial_{n}) \cdot (\partial_{1}, \partial_{2}, \ldots, \partial_{n})$；这种记号清楚地表明了算子的旋转对称性， 也清楚地表明了分部积分恒等式

\[
-\int \Delta u\, v\, dx = \int \nabla u \cdot \nabla v\, dx, \qquad u, v \in \mathcal{C}_{0}^{\infty}(\mathbb{R}^{n})
\]

后者可视为该算子的另一种定义\footnote{后一种定义可推广到 Riemann 流形上的 Laplace-Beltrami 算子， 此时梯度、点积与体积形式都必须关于度量来取.这一几何观点为我们提供了看待标准 Laplacian 旋转不变性的另一种方式.}.

在引言的习题中， 我们利用球对称求出了 Laplace 方程的基本解.现在我们可以严格地验证：我们形式导出的基本解是真正的基本解.

\begin{proposition}\label{prop:c4-2.1}
对所有 $n \geq 3$， 定义 $K_{n}(x) = \bigl((2-n)\omega_{n}\bigr)^{-1}|x|^{2-n}$；而对 $n = 2$， 定义 $K_{2}(x) = (2\pi)^{-1}\log|x|$.这里 $\omega_{n}$ 表示单位球面 $\mathbb{S}^{n-1}$ 的面积.则对所有 $n \geq 2$ 有

\[
\Delta K_{n} = \delta_{0}.
\]
\end{proposition}

\begin{proof}
由直接计算， $\Delta K_{n} = (\partial_{r}^{2} + \frac{n-1}{r}\partial_{r})K_{n}$ 在原点之外消失， 因此可表示为 $\delta_{0}$ 的若干导数的和.于是 $\Delta K_{n}$ 是支集在 $\mathbb{R}^{n}$ 原点处的分布， 且是 $-n$ 次齐次的， 这意味着它是 $\delta_{0}$ 的常数倍.为确定该常数， 我们可以使用任意检验函数， 并且（与 Cauchy-Riemann 方程组中的考虑相同）我们选取单位球的示性函数 $H(1-|x|)$.略加滥用记号， 记 $K_{n}(x) = K_{n}(|x|)$， 则

\[
\begin{aligned}
\int \Delta K_{n}(x)H(1-|x|)dx &\equiv -\int \nabla K_{n}(x) \cdot \nabla H(1-|x|)dx \\
&= \int \delta(1-|x|)\nabla K_{n}(|x|) \cdot \frac{x}{|x|}dx \\
&\equiv \int \delta(1-|x|)\frac{\mathrm{d}K_{n}}{\mathrm{d}r}(|x|)\frac{x}{|x|} \cdot \frac{x}{|x|}dx \\
&= \int_{|x|=1}\frac{\mathrm{d}K_{n}}{\mathrm{d}r}(1)d\sigma \\
&= 1
\end{aligned}
\]
\end{proof}

有了基本解 $K_{n}(x)$， 我们就可以求解非齐次方程

\[
\Delta V = \rho, \qquad \rho \in \mathcal{C}_{0}^{\infty}(\mathbb{R}^{n})
\]

其解由公式 $V(x) = K_{n}*\rho(x) = \int K_{n}(x-y)\rho(y)dy$ 给出.该解也常常记作 $\Delta^{-1}\rho$.

凭借一些微分几何的基本知识， 我们可以给出另一个证明.在极坐标 $x = r\omega$（$r > 0$， $|\omega| = 1$）下， $\Delta$ 取如下形式：

\[
\Delta = \partial_{r}^{2} + \frac{n-1}{r}\partial_{r} + r^{-2}\Delta_{\mathbb{S}^{n-1}},
\]

其中 $\Delta_{\mathbb{S}^{n-1}}$ 是单位球面 $\mathbb{S}^{n-1}$ 上的 Laplace-Beltrami 算子.

\noindent\textbf{习题 1.}\quad 证明：具有度量 $g$\footnote{其中 $\nabla_{g}u$ 由 $du(X) = g(\nabla_{g}u, X)$ 隐式定义， 而 $\Delta_{g}$ 由恒等式 $-\int_{M}\Delta_{g}u\, v\, dS_{g} = \int_{M}g(\nabla_{g}u, \nabla_{g}v)\, dS_{g}$（$u, v \in \mathcal{C}_{0}^{\infty}(M)$）定义.}的 Riemann 流形上的 Laplace-Beltrami 算子， 在局部坐标 $x^{i}$ 下由下式给出：

\[
\Delta_{g}\phi = \frac{1}{\sqrt{|g|}}\partial_{i}\bigl(g^{ij}\sqrt{|g|}\partial_{j}\phi\bigr).
\]

这里 $g^{ij}$ 是逆度量 $g^{-1}$ 相对于坐标 $x^{i}$ 的分量.$M$ 上的体积元 $dS_{g}$ 在局部坐标下由 $dS_{g} = \sqrt{|g|}\, dx^{1}dx^{2}\cdots dx^{n}$ 给出.注意， 在紧流形 $M$ 上

\[
\int_{M}\Delta_{g}u\, v\, dS_{g} = \int_{M}u\, \Delta_{g}v\, dS_{g}.
\]

\noindent\textbf{习题 2.}\quad 计算单位球面 $\mathbb{S}^{n-1}$ 的 Laplace-Beltrami 算子， 并验证 $\Delta$ 的极分解公式.对特殊情形 $n = 3$， 相对于坐标 $x^{1} = r\cos\theta^{1}$， $x^{2} = r\sin\theta^{1}\cos\theta^{2}$， $x^{3} = r\sin\theta^{1}\sin\theta^{2}$（$\theta^{1} \in [0, \pi)$， $\theta^{2} \in [0, 2\pi)$）证明

\[
\Delta_{\mathbb{S}^{2}} = \partial_{\theta^{1}}^{2} + \cot\theta^{1}\partial_{\theta^{1}} + \frac{1}{\sin^{2}\theta^{1}}\partial_{\theta^{2}}^{2}.
\]

此外， 面积元 $dS_{\omega}$ 取如下形式：$dS_{\omega} = r^{2}\sin\theta^{1}\, d\theta^{1}d\theta^{2}$.

\begin{proof}[几何推导]
对光滑函数 $\phi(\boldsymbol{x}) = \phi(r\omega)$， 在极坐标 $r = |x|$、$\omega \in \mathbb{S}^{n-1}$（$\mathbb{R}^{n}$ 中的单位球面）下， 有

\[
\begin{aligned}
\Delta\phi &= \left(\partial_{r}^{2} + \frac{n-1}{r}\partial_{r} + r^{-2}\Delta_{\mathbb{S}^{n-1}}\right)\phi \\
&= r^{-(n-1)}\partial_{r}\bigl(r^{n-1}\partial_{r}\phi\bigr) + r^{-2}\Delta_{\mathbb{S}^{n-1}}\phi
\end{aligned}
\]

现在我们过渡到极坐标 $x = r\omega$， 使体积元可以写成 $dx = r^{n-1}drdS_{\omega}$.在 Riemann 流形 $\mathbb{S}^{n-1}$ 上分部积分， 我们计算得

\[
\begin{aligned}
\langle \Delta K_{n}, \phi\rangle &= \langle K_{n}, \Delta\phi\rangle \\
&= \int_{|\omega|=1}\int_{0}^{\infty}K_{n}(r)\partial_{r}\bigl(r^{n-1}\partial_{r}\phi\bigr)drdS_{\omega} + \int_{|\omega|=1}\int_{0}^{\infty}K_{n}(r)\Delta_{\mathbb{S}^{n-1}}\phi\, drdS_{\omega} \\
&= \bigl((2-n)\omega_{n}\bigr)^{-1}\int_{|\omega|=1}\int_{0}^{\infty}r^{-n+2}\partial_{r}\bigl(r^{n-1}\partial_{r}\phi\bigr)drdS_{\omega} + 0 \\
&= -\int_{0}^{\infty}r^{-n+1}\bigl(r^{n-1}\partial_{r}\tilde{\phi}\bigr)dr = -\int_{0}^{\infty}\partial_{r}\tilde{\phi}(r)dr = \phi(0)
\end{aligned}
\]

其中在上述计算中， 我们定义 $\tilde{\phi}(r)$ 为 $\phi$ 在 $|x| = r$ 上的平均.我们推断， 对 $n \geq 3$ 有 $\Delta K_{n} = \delta_{0}$， 正如所愿.情形 $n = 2$ 可以同样处理.
\end{proof}

\begin{remark*}
注意， 除常数外， $K_{n}(x)$ 的表达式也可以通过寻找球对称解 $K = K(|x|)$ 轻易猜出.事实上， 方程 $\Delta K = 0$ 归结为常微分方程 $K''(r) + \frac{n-1}{r}K'(r) = 0$.
\end{remark*}

找到基本解后， 我们可以像前面一样立即推出一个表示公式.设 $\phi \in \mathcal{C}_{0}^{\infty}$ 是光滑检验函数， $\Omega$ 是带 Lipschitz 边界的有界开集， $x \in \Omega$.我们有

\[
\begin{aligned}
\phi(x) &= \int_{\Omega}\phi(y)\delta(x-y)dy \\
&= \int \chi_{\Omega}\phi\, \Delta K_{n}(x-y)dy
\end{aligned}
\]

我们的策略是分部积分， 允许至多一个导数作用到示性函数上.回忆我们关于分布回拉的讨论：$\nabla\chi_{\Omega} = \vec{n}\, d\sigma_{\partial\Omega}$， 其中 $\vec{n}$ 是内单位法向， $d\sigma$ 是边界上的曲面测度.与经典分部积分不同， 这里的处理仿佛没有边界项， 因为乘积 $\chi_{\Omega}\phi\, \Delta K_{n}(x-y)$ 具有紧支集.对在边界 $\partial\Omega$ 处具有连续一阶导数的函数 $f$， 我们用 $\frac{\partial f}{\partial\nu}$ 表示外单位法向导数.

\[
\begin{aligned}
&\phi(x) = -\int \nabla(\chi_{\Omega}\phi) \cdot \nabla K_{n}(x-y)dy \\
&\quad = -\int \phi\nabla\chi_{\Omega} \cdot \nabla K_{n}(x-y)dy - \int \chi_{\Omega}\nabla\phi \cdot \nabla K_{n}(x-y)dy \\
&\quad = \int_{\partial\Omega}\phi\frac{\partial}{\partial\nu}K_{n}(x-y)d\sigma(y) - \int \nabla\phi \cdot \bigl(\nabla(\chi_{\Omega}K_{n}) - K_{n}\nabla\chi_{\Omega}\bigr) \\
&\quad = \int_{\partial\Omega}\phi\frac{\partial}{\partial\nu}K_{n}(x-y)d\sigma(y) - \int_{\partial\Omega}K_{n}(x-y)\frac{\partial\phi}{\partial\nu}d\sigma(y) + \int_{\Omega}K_{n}(x-y)\Delta\phi\, dy
\end{aligned}
\]

于是我们得到表示公式

\begin{equation}
\phi(x) = \int_{\Omega}K_{n}(x-y)\Delta\phi\, dy + \int_{\partial\Omega}\phi\frac{\partial}{\partial\nu}K_{n}(x-y)d\sigma(y) - \int_{\partial\Omega}K_{n}(x-y)\frac{\partial\phi}{\partial\nu}d\sigma(y)\label{eq:c4-137}
\end{equation}

特别地（由逼近）， 如果 $u$ 在 $\Omega$ 内调和（并且在 $\Omega$ 的某个邻域内是 $C^{2}$ 的）：

\begin{proposition}\label{prop:c4-2.2}
\begin{equation}
u(y) = \int_{\partial\Omega}\left(u(x)\frac{\partial K_{n}}{\partial\nu}(x-y) - \frac{\partial u}{\partial\nu}(x)K_{n}(x-y)\right)d\sigma(x)\label{eq:c4-138}
\end{equation}
\end{proposition}

有了表示公式， 我们可以重复前面针对 Cauchy-Riemann 方程组所做的大部分分析.我们发现， 由于基本解是实解析的， 上述调和函数实际上是实解析的， 并且有以 $u$ 与 $\frac{\partial u}{\partial\nu}$ 的边界值表示的、关于导数的定量先验估计.我们还可以用这些估计证明：满足 $\Delta u = 0$ 的连续函数实际上是经典解.但用这个公式继续分析调和函数可能具有误导性， 因为调和函数的内点值由其边界值唯一确定， 因此法向导数不能任意给定.

事实上， 我们现在陈述的调和函数最大值原理表明：区域内部的调和函数仅由其边界值确定.

\begin{theorem}[最大值原理]\label{thm:c4-2.3}
若 $u : \Omega \to \mathbb{R}$ 在连通开集 $\Omega$ 上是 $\mathcal{C}^{2}$ 的， 且在 $\Omega$ 内 $\Delta u \geq 0$， 则 $u$ 不能在内部取得最大值， 除非 $u$ 是常数.特别地， 当 $\Omega$ 有界时，

\begin{equation}
\sup_{\overline{\Omega}}u(x) = \sup_{\partial\Omega}u(x)\label{eq:c4-139}
\end{equation}

并且作为推论， 当 $u$ 调和时， $\sup_{\overline{\Omega}}|u(x)| = \sup_{\partial\Omega}|u(x)|$.
\end{theorem}

\begin{proof}
第一个论断称为 $\mathcal{C}^{2}$ 次调和函数的强最大值原理（次调和函数指在经典意义下满足 $\Delta u \geq 0$ 的函数）；该定理表明， 区域 $\Omega$ 中的次调和函数在内部严格保持在边界值处处更大的调和函数之下， 故得名``次调和''.强最大值原理将是后面要证明的一个定理（次调和函数的平均值不等式）的简单推论， 不过我们把证明留给读者.但眼下我们至少可以证明``弱最大值原理''（即不等式 \eqref{eq:c4-139}）.

事实上， 当 $\Delta u > 0$ 严格正时， （强）最大值原理是显然的， 因为函数 $u$ 在任意点至少在一个方向上是严格凸的.通过逼近（例如把 $u$ 换成 $u + \varepsilon x_{1}^{2}$）， 我们可以在只假设 $\Delta u \geq 0$ 于有界区域 $\Omega$ 内成立、且 $u$ 连续延拓到 $\partial\Omega$ 的情况下得到弱最大值原理.

显然， 把 $u$ 换成 $-u$ 后， 对超调和函数成立``强最小值原理''\footnote{超调和指 $\Delta u \leq 0$.}， 它特别蕴含了定理中陈述的最后一个等式.
\end{proof}

证明中描述的超调和与次调和函数的概念， 即使对调和函数本身的分析也是有用的：它们更容易显式构造（用指数函数、多项式等）， 并且可以按照上面的最大值原理来界定调和函数.
按照全纯函数最大模原理讨论之后的同样推理， 调和函数必然有形如 $u(y) = \int_{\partial\Omega}u(x)d\mu_{y}(x)$ 的表示公式， 其中 $\mu_{y}$ 是依赖于 $y \in \Omega$ 的有限测度.我们可以如下得到这样的表示公式：如果能够找到在 $\Omega$ 边界上与基本解 $K_{n}(x-y)$ 重合的调和函数 $\psi_{y}(x)$， 则函数 $G(x, y) = K_{n}(x-y) - \phi_{y}(x)$ 满足

\[
\Delta_{x}G(x, y) = \delta_{y}(x) \text{ 于 } \Omega, \qquad G(x, y) = 0 \text{ 于 } \partial\Omega
\]

由最大值原理， 这样的函数至多有一个.上面的函数 $G(x, y)$ 称为 $\Omega$ 的\newterm{green-function}{Green 函数}{Green's function}， 它在引言的习题中被正式引入.按照我们前面的表示公式那样计算 $u(y) = \int\chi_{\Omega}(x)u(x)\Delta G(x, y)dx$（这次 $G$ 的边界条件消去了一个边界项：$\chi_{\Omega}\nabla G(x, y) = \nabla(\chi_{\Omega}G(x, y))$）， 我们得到所需的表示公式：

\begin{proposition}\label{prop:c4-2.4}
若 $u$ 在 $\Omega$ 内调和， 在 $\Omega$ 的某个邻域内是 $C^{2}$ 的， 且 $G(x, y)$ 如上， 则

\begin{equation}
u(y) = \int_{\partial\Omega}\left[u(x)\frac{\partial G}{\partial\nu}(x, y)\right]d\sigma(x)\label{eq:c4-140}
\end{equation}
\end{proposition}

\begin{remark*}
注意：我们并未证明公式右端定义的函数对任意区域都有定义， 也未证明它定义了一个调和函数， 甚至没有证明当 $y$ 趋向边界时它实现被积函数中的边界值.当边界足够好（比如 Lipschitz）时， 这些都可以证明， 并且任意连续边界值都可以由调和函数实现（与 Cauchy-Riemann 方程组的情形相反）.
\end{remark*}

出现在 \eqref{eq:c4-140} 中的概率测度 $\frac{\partial G}{\partial\nu}(x, y)d\sigma(x)$ 描述了从点 $y$ 出发的\newterm{random-walk}{随机游走}{random walk}首次接触边界的概率分布.于是， 调和函数在点 $y$ 处的值可以看作从 $y$ 出发的随机游走在首次接触点处边界数据所取值的期望.由这一解释， 调和函数的一些特征（例如下面的最大值原理与平均值性质）是显然的， 但我们不在此处探讨这一解释.

\begin{example*}
对于 $\mathbb{R}^{n}$ 中的半空间 $x^{n} > 0$， 可以通过在点 $y^{*} = (y^{1}, \ldots, y^{n-1}, -y^{n})$ 放置一个负点源， 并定义 $\mathcal{G}(x, y) = K(x-y) - K(x-y^{*})$， 来得到 Green 函数（$\Delta G(x, y) = \delta_{y}$， $G(x^{1}, \ldots, x^{n-1}, 0) = 0$）的显式公式.则 $G(x, y)$ 在 $x^{n} = 0$ 上为零， 因为这样的点到 $y$ 与 $y^{*}$ 等距， 而基本解只依赖于 Euclidean 距离.同样的方法也可以用来构造球 $|y| \leq 1$ 的 Green 函数.此时使用共形反射 $y \mapsto y^{*} = \frac{y}{|y|^{2}}$， 它固定球面 $|y| = 1$.Green 函数于是取如下形式：$G(x, y) = K(x-y) - K(|x|(x-y^{*}))$.在二维情形可以用全纯函数得到更多例子.
\end{example*}

我们到目前为止的分析的主要结果（最大值原理、可由 Green 公式推出的各种先验估计， 以及 $\mathbb{R}^{n}$ 上对紧支集数据解的存在性）中， 有许多对与 Laplacian 密切类似的其他算子也成立.例如， 通过换元可以看到， 当 $u$ 是调和函数且 $v(\Psi(x)) = u(x)$（$\Psi$ 是 $\mathbb{R}^{n}$ 的微分同胚）时， $v$ 将满足如下形式的方程：$L[v] = \sum_{i, j=1}^{n}a^{ij}(y)\partial_{i}\partial_{j}v + b^{i}(y)\partial_{i}v = 0$， 其中光滑函数 $a^{ij}(y)$ 是对称正定矩阵 $a^{ij}(\Psi(x)) = D\Psi(x)(D\Psi)^{t}(x)$ 的系数， 而一阶项依赖于 $\Psi$ 的二阶导数.

更一般的形如 $L = \sum_{i, j=1}^{n}a^{ij}(x)\partial_{i}\partial_{j} + \sum_{i=1}^{n}b^{i}\partial_{i}$ 的算子， 其中矩阵 $(a^{ij}(x))$ 对称正定， 称为\newterm{elliptic-operator}{椭圆算子}{elliptic operator}.它们与 Laplacian 有许多共同性质毫不奇怪， 但它们通常不一定具有与 Laplace 算子相同数量的对称性\footnote{当人们提到一个偏微分算子的``对称性''时， 通常想到的是一族与算子可交换的向量场， 或使算子不变的它们的流.$\Delta$ 在 $\mathbb{R}^{n}$ 上的对称性是底层 Euclidean 几何的对称性：无穷小平移 $\partial_{x^{1}}, \ldots, \partial_{x^{n}}$ 连同无穷小旋转.相应的流生成 Euclidean 空间的刚体运动群.在这个微分同胚不变的意义下， 如上对 Laplacian 的坐标变换具有相同数量的对称性.我们后面会看到， 这种对称性在分析微分算子时可能非常有用.无论如何， 显然仅凭椭圆性条件并不能推出这种保持算子的流的存在.}， 因此需要更稳健的方法才能成功地分析它们.不过， 也有把 Laplace 算子的结果与估计延拓、输运到更一般椭圆算子的方法.

下面的定理体现了 Laplace 算子的旋转与平移对称性， 事实上它既刻画了调和函数， 也刻画了 Laplace 算子本身.因此， 它可以用来证明用其他方法无法触及的、关于 Laplace 算子与调和函数的结果， 因此它的应用也仅限于这些目的.该定理表明 Laplacian 如何控制不同半径的球平均的变化.

\begin{theorem}[平均值性质]\label{thm:c4-2.5}
当 $u$ 在中心为 $x$、半径为 $R^{*} > R$ 的球内调和时， $u(x)$ 等于它在以 $x$ 为中心、半径为 $R$ 的球面上的平均值

\[
u(x) = \frac{1}{n\omega_{n}}\frac{1}{R^{n-1}}\int_{|y-x|=R}u(y)d\sigma(y)
\]

当 $\Delta u \geq 0$ 时， ``$=$''换成``$\leq$''.事实上， 对所有 $u \in \mathcal{C}^{\infty}(\mathbb{R}^{n})$、$0 < R_{1} < R_{2} < R^{*}$， 有

\[
\begin{aligned}
&\frac{1}{R_{2}^{n-1}}\int_{|x|=R_{2}}u(x)d\sigma(x) - \frac{1}{R_{1}^{n-1}}\int_{|x|=R_{1}}u(x)d\sigma(x) \\
&= \int_{R_{1}}^{R_{2}}\left[\int_{|y|\leq\tau}\Delta u(y)dy\right]\tau^{-(n-1)}d\tau
\end{aligned}
\]
\end{theorem}

\begin{proof}
我们证明最后一个公式， 因为定理的第一个恒等式是它的直接推论（令内半径趋于 0）.事实上， 后一个公式表明当 $\Delta u \geq 0$ 时球平均随半径增大而增大.

虽然公式以积分形式陈述， 但我们用辅助函数 $\psi(\tau) = \frac{1}{\tau^{n-1}}\int_{|x|=\tau}u(x)d\sigma(x) = \frac{1}{\tau^{n-1}}\int\delta(\tau-|x|)u(x)dx$ 证明其微分版本.记 $r = |x|$：

\[
\begin{aligned}
\frac{\mathrm{d}\psi}{\mathrm{d}\tau} &= \int\frac{1}{\tau^{n-1}}\left(\frac{\partial}{\partial\tau} - \frac{n-1}{\tau}\right)\delta(\tau-r)u(x)dx \\
&= \frac{1}{\tau^{n-1}}\int\left(-\frac{\partial}{\partial r} - \frac{n-1}{r}\right)\delta(\tau-r)u(x)dx \\
&= \frac{1}{\tau^{n-1}}\int\left(-\frac{\partial}{\partial r} - \frac{n-1}{r}\right)\left(-\frac{\partial}{\partial r}\right)H(\tau-r)u(x)dx \\
&= \frac{1}{\tau^{n-1}}\int\Delta H(\tau-r)u(x)dx \\
&= \frac{1}{\tau^{n-1}}\int_{|x|\leq\tau}\Delta u(x)dx
\end{aligned}
\]

在第二个等式中， 我们使用了如下事实：$\delta(\tau-r)$ 是 0 阶分布， $\frac{n-1}{r}$ 与 $\frac{n-1}{\tau}$ 在它的支集上到 0 阶一致；还使用了如下事实：任何由映射 $(\tau, x) \mapsto \tau-r$ 回拉的分布， 在 $\mathrm{d}\tau - \mathrm{d}r$ 的零空间中的任何向量场作用下保持不变（因此 $(\partial_{\tau} + \partial_{r})\delta(\tau-r) = 0$）.在第四行中， 我们认识到算子 $\frac{\partial^{2}}{\partial r^{2}} + \frac{n-1}{r}\frac{\partial}{\partial r}$ 作用于球对称函数时与 Laplacian 一致.

对 $\tau$ 从 $R_{1}$ 到 $R_{2}$ 积分即得所需公式.
\end{proof}

\begin{remark}\label{rem:c4-2.6}
我们实质上计算得到

\[
\int u(x)d\mu(x) = \int \Delta u\, \Delta^{-1}\mu\, dx
\]

其中 $\mu$ 是引言习题 5 中的测度.
\end{remark}

上述公式的一个特殊情形在复分析中有重要应用.当 $f$ 是圆盘 $\mathcal{D}$ 中的非零全纯函数时， 我们有恒等式

\[
\frac{1}{2\pi}\Delta\log|f(z)| = \sum_{\rho_{k}}\delta_{\rho_{k}}(z)
\]

其中 $\rho_{k}$ 遍历 $f$ 的有限个零点（按重数计）——右端的测度在代数几何中称为 $f$ 的零因子.在零点 $\rho_{j}$ 附近局部地看出这个恒等式的方法是：把 $f$ 写成 $f(z) = e^{g(z)}(z-\rho_{j})^{n}$， 其中 $g$ 在 $\rho_{j}$ 的某个邻域内全纯；并回忆全纯函数的实部是调和的， 且二维情形 Laplacian 的基本解由 $\frac{1}{2\pi}\log|z|$ 给出.这个计算特别蕴含 $\log|f(z)|$ 是次调和函数（这一事实可以与最大值原理结合使用， 给出关于全纯函数的强估计）.把\thmref{thm:c4-2.5} 中的一般公式应用于 $u = \frac{1}{2\pi}\log|f(z)|$， 我们得到

\begin{proposition}[Jensen 公式]\label{prop:c4-2.7}
设 $f$ 在中心为零、半径为 $R$ 的闭球的某个邻域内全纯， 且其（离散的）零点集合 $\{\rho_{j}\}$ 满足 $0 < |\rho_{j}| \neq R$.则

\[
\frac{1}{2\pi R}\int_{0}^{2\pi}\log|f(Re^{i\theta})|d\theta - \log|f(0)| = \sum_{|\rho_{j}|<R}\log\left|\frac{R}{\rho_{j}}\right|
\]
\end{proposition}

\begin{proof}
从\thmref{thm:c4-2.5} 中的一般公式出发， 我们通过标准的磨光构造让 $u \to \frac{1}{2\pi}\log|f(z)|$， 并注意到我们对 $f(z)$ 的假设足以保证定理中的积分公式对 $R_{2} = R$ 与 $0 < R_{1} < \min|\rho_{j}|$ 仍然有效.令 $R_{1}$ 趋于 0 即得 Jensen 公式的左端.

现在用 Fubini 定理（或分部积分）显式计算一般公式的右端：

\[
\begin{aligned}
\int_{0}^{R}\left[\int_{|y|\leq\tau}\frac{1}{2\pi}\Delta\log|f(z)|dy\right]\tau^{-1}d\tau &= \int_{0}^{R}\left[\int_{0}^{\tau}\sum_{\rho_{j}}\delta_{|\rho_{j}|}(t)dt\right]\tau^{-1}d\tau \\
&= \int_{0}^{R}\sum_{\rho_{j}}\delta_{|\rho_{j}|}(t)\left[\int_{t}^{R}\tau^{-1}d\tau\right]dt \\
&= \sum_{|\rho_{j}|<R}\log\left|\frac{R}{\rho_{j}}\right|
\end{aligned}
\]
\end{proof}

习题 1--5 中的定理是研究 Laplace 方程（连同它们到其他椭圆 PDE 的推广）的基本定理， 如果读者想查阅参考书， 其证明可以在 Evans 或 Gilbarg 与 Trudinger 的书中以某种形式找到.

\noindent\textbf{习题 1.}\quad [Hopf 引理] 断言：有界开集上的调和函数 $u : \Omega \to \mathbb{R}$ 在边界光滑的边界点 $x_{0}$ 处（其中对 $x \in \Omega\backslash\{x_{0}\}$ 有 $u(x_{0}) > u(x)$）必须满足 $\frac{\partial u}{\partial\nu}(x_{0}) > 0$.证明这一事实.一种方法是设计 $u$ 在 $x_{0}$ 附近的适当超调和扰动， 并用弱最大值原理把 $u$ 界定在超调和扰动之下.

\noindent\textbf{习题 2.}\quad 证明次调和函数的强最大值原理.（到现在为止， 你也许已经能看出不止一种证法.）尝试把这一结果也推广到椭圆方程的 $\mathcal{C}^{2}$``次解''——换言之， 假设

\[
Lu = \sum_{i, j=1}^{n}a^{ij}(x)\partial_{i}\partial_{j}u + \sum_{i=1}^{n}b^{i}(x)\partial_{i}u \geq 0
\]

其中矩阵满足 $a^{ij}(x)\xi_{i}\xi_{j} \geq \lambda|\xi|^{2}$， 即一致正定.

\noindent\textbf{习题 3.}\quad 证明边界良好的开集 $\Omega$ 的 Green 函数具有如下性质：$G(x, y) = \int\delta_{x}(z)G(z, y)dz = \int\delta_{y}(z)G(z, x)dz = G(y, x)$；对 $x \in \Omega\backslash\{y\}$ 有 $G(x, y) < 0$；且在任意边界点 $x_{0}$ 处有 $\frac{\partial G}{\partial\nu}(x_{0}) > 0$；以及 $\int_{\partial\Omega}\frac{\partial G}{\partial\nu}(x, y)d\sigma(x) = 1$.你如何从物理上或概率上解释这些事实中的任何一个？

\noindent\textbf{习题 4.}\quad [Harnack 不等式] 证明：若 $K \subseteq \Omega$ 紧， 则存在依赖于 $K$ 的常数 $C$， 使得对 $\Omega$ 中所有非负调和函数 $u$ 有

\[
\sup_{K}u \leq C\inf_{K}u
\]

\begin{theorem}[Liouville 定理]\label{thm:c4-2.8}
$\mathbb{R}^{n}$ 上处处有界（整）调和函数是常数.
\end{theorem}

\noindent\textbf{习题 5.}\quad 证明 Liouville 定理.

\noindent\textbf{习题 6.}\quad 共面积公式， 用分布理论（稍带含糊地）写成 $h(f(x)) = \int h(t)\delta_{0}(f(x)-t)dt$， 它允许我们把一般密度 $h(f(x))$ 分解为 $f$ 的水平集上的曲面测度， 然后可以逐个分析.例如， 当 $f(x) = |x|$ 时， 共面积公式归结为极坐标积分.该公式等价于分布回拉的定义， 并在附录中证明；请用它证明如下恒等式：

设 $u \in \mathcal{C}_{0}^{\infty}(\mathbb{R}^{n})$， $\rho(x) = \tilde{\rho}(|x|)$ 是球对称密度.则若 $\rho_{\varepsilon} \equiv \varepsilon^{-n}\rho(\frac{x}{\varepsilon})$， 我们有公式

\[
\frac{\partial}{\partial\varepsilon}u*\rho_{\varepsilon}(x) = \varepsilon^{-(n-1)}\int_{0}^{\infty}\left[\int_{|y|\leq\varepsilon t}\Delta u(x-y)dy\right]t\tilde{\rho}(t)dt
\]

\noindent\textbf{习题 7.}\quad 证明（至少 $n \geq 3$ 的情形）如下

\begin{proposition}\label{prop:c4-2.9}
对任意 $f \in \mathcal{C}_{0}^{\infty}(\mathbb{R}^{n})$（$n \geq 3$）， 方程 $\Delta u = f$ 有唯一的在无穷远处消失的光滑解， 即在 $|x| \to \infty$ 时趋于 0.该解由 $\int_{\mathbb{R}^{n}}K_{n}(x-y)f(y)dy$ 表示.对 $n = 2$， 同样的方程有光滑解 $u(x)$， 满足 $\lim_{|x|\to\infty}\frac{|u(x)|}{|x|} = 0$， 且当 $|x| \to \infty$ 时 $|\partial u(x)| \to 0$.该解由 $\int_{\mathbb{R}^{2}}\frac{1}{2\pi}\log|x-y|f(y)dy$ 表示， 并且（在此类中）在相差一个加法常数的意义下唯一.
\end{proposition}

\noindent\textbf{习题 8.}\quad 证明全纯函数的最大模原理.

我们将回到 Laplace 方程（及其某些推广）的研究， 见第 3 章.
\section{D'Alembert 算子}

回忆 D'Alembert 算子 $\Box = -\partial_{t}^{2} + \Delta$ 是 $\mathbb{R}^{1+n}$ 中在平移与 Lorentz 变换（即 Poincaré 群）下不变的最简单微分算子.最容易的验证方法是写成 $\Box = m^{\alpha\beta}\partial_{\alpha}\partial_{\beta}$， 其中 $m$ 是 Minkowski 度量.由于 $m$ 在 Poincaré 群下不变， $\Box$ 亦然.因此， 寻找如下形式的基本解是有意义的：$\phi(t, x) = f(\rho)$， 其中\footnote{回忆我们记 $t = x^{0}$， 并对指标 $\alpha, \beta = 0, 1, \ldots, n$ 使用求和约定.} $\rho = t^{2} - |x|^{2} = -m_{\alpha\beta}x^{\alpha}x^{\beta}$ 在 Lorentz 变换下不变.此外， 由于 $\mathbb{R}^{n+1}$ 上的分布 $\delta_{0}$ 是 $-(n+1)$ 次齐次的， 而应用 $\Box$ 使齐次度降低 2， 我们断定 $f$ 必须是 $-\frac{n-1}{2}$ 次齐次的.因此， 基本解的一个好的候选者必然具有形式 $E = c_{n}(t^{2}-|x|^{2})^{-\frac{n-1}{2}}$（$c_{n}$ 为某常数）， 于区域 $t > |x|$ 中.于是我们转而寻找一个 $-n+1$ 次齐次的分布 $E_{+}$， 它在区域 $t > |x|$ 中与 $c_{n}(t^{2}-|x|^{2})^{-\frac{n-1}{2}}$ 一致.乍一看这可能很困难， 因为 $(t^{2}-|x|^{2})^{-\frac{n-1}{2}}$ 沿 $|x| = |t|$ 的奇性阶数很高；直到我们意识到可以利用命题 ?? 定义的齐次分布族 $j_{a}$.事实上我们需要选取 $a = -\frac{n-1}{2} + 1$， 并取 $E_{+}$ 与 $j_{-\frac{n-1}{2}+1}(t^{2}-x^{2})$ 成正比， 理解为回拉 $f^{*}(j_{-\frac{n-1}{2}+1})$， 其中 $f = t^{2}-|x|^{2}$.在当前的语境下， 稍微改变记号更方便， 写成

\[
\chi_{+}^{a} := j_{a+1}
\]

于是

\[
E = \chi_{+}^{-\frac{n-1}{2}}(t^{2}-|x|^{2}).
\]

注意， 表达式 $\chi_{+}^{-\frac{n-1}{2}}(t^{2}-|x|^{2})$ 并不完全严格， 因为 $t^{2}-|x|^{2}$ 的梯度在原点处消失， 因此 $\chi_{+}^{-\frac{n-1}{2}}(t^{2}-|x|^{2})$ 只在 $\mathbb{R}^{n+1}-\{0\}$ 上定义分布.严格的表述需要更多小心， 但该分布特殊的齐次度基本上允许唯一的延拓到整个空间， $n = 3$ 的情形参见第 4 节的习题 3.现在我们得到的分布具有正确的性质， 只是它的支集覆盖整个区域 $|x| \leq |t|$.出于确定的物理原因， 我们更倾向于支集只在未来区域 $|x| \leq t$ 中的分布.这就定义了我们的\newterm{forward-fundamental-solution}{前向基本解}{forward fundamental solution}候选：

\[
E_{+}^{(n+1)}(t, x) = c_{n}H(t)\chi_{+}^{-\frac{n-1}{2}}(t^{2}-|x|^{2})
\]

其中 $H(t)$ 是支集在 $t \geq 0$ 上的 Heaviside 函数， $c_{n}$ 是待定的归一化常数.利用链式法则容易证明， $\Box E_{+}^{(n+1)} = m^{\alpha\beta}\partial_{\alpha}\partial_{\beta}\bigl(\chi_{+}^{-\frac{n-1}{2}}(f)\bigr)$ 必须在原点之外消失.由通常的齐次性考虑我们推断 $\Box E_{+}^{(n+1)}$ 与原点处的 $\delta$ 函数成正比.因此只剩下确定归一化因子 $c_{n}$.我们有如下结果：

\begin{theorem}\label{thm:c4-3.1}
由下式定义的分布 $E_{+}^{(n+1)}$

\begin{equation}
E_{+}^{(n+1)}(t, x) = -\frac{1}{2}\pi^{\frac{1-n}{2}}H(t)\chi_{+}^{-\frac{n-1}{2}}(t^{2}-|x|^{2})\label{eq:c4-141}
\end{equation}

是波动方程支集于前向区域 $|x| \leq t$ 的唯一基本解.
\end{theorem}

我们将在后面证明这个定理， 现在先作几点说明.首先， 注意 $n > 1$ 为奇数与 $n$ 为偶数的情形之间有根本的区别.事实上， 在前一种情形， $E_{+}^{(n+1)}$ 的支集只在区域 $|x| < t$ 的边界上， 即未来光锥 $|x| = t$ 上；而在后一种情形， $E_{+}^{(n+1)}$ 的支集覆盖整个前向区域 $|x| \leq t$.更精确地说， 利用链式法则以及 $\frac{\mathrm{d}}{\mathrm{d}\lambda}\chi_{+}^{s}(\lambda) = \chi_{+}^{s-1}(\lambda)$， 我们可以把基本解在远离原点处写得稍显式一些：在维数 $n = 3+2k$ 时， 基本解看起来像支集于前向光锥上的一个测度的导数

\[
c_{n}H(t)\left(\frac{-1}{2r}\partial_{r}\right)^{k}\delta(t^{2}-|x|^{2}) = c_{n}H(t)\left(\frac{1}{2t}\partial_{t}\right)^{k}\delta(t^{2}-|x|^{2})
\]

而在 $n = 2+2k$ 维时， 它具有如下形式

\[
c_{n}H(t)\left(\frac{-1}{2r}\partial_{r}\right)^{k}\frac{1}{\sqrt{t^{2}-|x|^{2}}}\cdot\mathbf{1}_{\{|x|\leq t\}} = c_{n}H(t)\left(\frac{1}{2t}\partial_{t}\right)^{k}\frac{1}{\sqrt{t^{2}-|x|^{2}}}\cdot\mathbf{1}_{\{|x|\leq t\}}
\]

（上面的分布相等， 因为 $\frac{1}{t}\partial_{t} + \frac{1}{r}\partial_{r}$ 位于 $\mathrm{d}(t^{2}-r^{2})$ 的零空间中）.在最重要的特殊情形 $n = 3$ 时， 我们有

\begin{equation}
E_{+}^{(1+3)} = \frac{1}{2\pi}H(t)\delta(t^{2}-|x|^{2}) = \frac{1}{4\pi}\delta(t-|x|)\label{eq:c4-142}
\end{equation}

另外， 对 $n = 2$ 有

\begin{equation}
E_{+}^{(1+2)} = \frac{1}{2\pi^{1/2}}H(t)(t^{2}-|x|^{2})^{1/2} =\label{eq:c4-143}
\end{equation}

重要的是要观察到， 知道奇数维的基本解就可以确定偶数维的基本解.这称为下降法.只需把 $E_{+}^{(2k+1+1)}$ 作用到与某个空间变量无关的检验函数上即可.不必担心该检验函数没有紧支集， 它会奏效， 因为对任意固定的 $t$， 基本解在 $x$ 上具有紧支集.作为例子， 读者不妨由 \eqref{eq:c4-142} 推导 \eqref{eq:c4-148}.另一个简单的说明：注意对每个固定的 $t$， $E_{+}$ 在 $x$ 上具有紧支集， 因此 $E_{+}$ 可以作用到任意在 $t$ 上具有紧支集的光滑函数上.

\noindent\textbf{习题 1.}\quad 推导维数 $n = 1$ 的基本解.事实上\footnote{在 $1+1$ 维中， 波算子可以分解为 $\Box = (\partial_{t}-\partial_{x})\cdot(\partial_{t}+\partial_{x})$.}证明 Cauchy 问题的一般解具有形式 $\psi_{1}(t+x) + \psi_{2}(t-x)$.

基本解使我们能够求解一般的 Cauchy 问题

\begin{equation}
\Box\phi = f, \qquad \phi(0, x) = f(x), \quad \partial_{t}\phi(0, x) = g(x)\label{eq:c4-144}
\end{equation}

为了明白如何做到这一点， 考虑点 $p = (t_{0}, x_{0})$（$t_{0} > 0$）， 并注意到对任何检验函数 $\phi$， 在上半空间 $\mathcal{D}_{+} = \{(t, x): t \geq 0\}$ 中有

\[
\phi(t_{0}, x_{0}) = \int_{\mathcal{D}_{+}}\phi(t, x)\delta_{p}(t, x)dtdx = \int_{\mathbb{R}^{1+n}}\chi_{+}\phi\left(m^{\alpha\beta}\partial_{\alpha}\partial_{\beta}\right)E_{p}(t, x)dtdx
\]

其中 $\chi_{+}$ 是 $\mathcal{D}_{+}$ 的示性函数， $E_{p}(t, x) = E_{+}(t-t_{0}, x-x_{0})$.因此， 分部积分得

\[
\begin{aligned}
\phi(t_{0}, x_{0}) &= -\int_{\mathbb{R}^{1+n}}m^{\alpha\beta}\partial_{\alpha}\chi\, \phi\, \partial_{\beta}E_{p} - \int_{\mathbb{R}^{1+n}}m^{\alpha\beta}\chi\, \partial_{\alpha}\phi\, \partial_{\beta}E_{p} \\
&= -\int_{\mathbb{R}^{1+n}}m^{\alpha\beta}\partial_{\alpha}\chi\, \phi\, \partial_{\beta}E_{p} + \int_{\mathbb{R}^{1+n}}\chi\, m^{\alpha\beta}\partial_{\alpha}\partial_{\beta}\phi\, E_{p} \\
&\quad + \int_{\mathbb{R}^{1+n}}m^{\alpha\beta}\partial_{\beta}\chi\, \partial_{\alpha}\phi\, E_{p}
\end{aligned}
\]

即

\[
\begin{aligned}
\phi(t_{0}, x_{0}) &= \int_{\mathcal{D}_{+}}\Box\phi\, E_{p} - \int_{\mathbb{R}^{1+n}}\partial^{\alpha}\chi\, \phi\, \partial_{\alpha}E_{p} + \int_{\mathbb{R}^{1+n}}\partial^{\alpha}\chi\, \partial_{\alpha}\phi\, E_{p} \\
&= \int_{\mathcal{D}_{+}}\Box\phi + \int_{\mathbb{R}^{1+n}}\partial_{t}\chi\, \phi\, \partial_{t}E_{p} - \int_{\mathbb{R}^{1+n}}\partial_{t}\chi\, \partial_{t}\phi\, E_{p} \\
&= \int_{\mathcal{D}_{+}}\Box\phi + \int_{\mathbb{R}^{1+n}}\delta(t)\, \phi\, \partial_{t}E_{p} - \int_{\mathbb{R}^{1+n}}\delta(t)\partial_{t}\phi\, E_{p} \\
&= \int_{\mathcal{D}_{+}}\Box\phi\, E_{p} - \int_{\mathbb{R}^{1+n}}\delta(t)\, \phi\, \partial_{t_{0}}E_{p} - \int_{\mathbb{R}^{1+n}}\delta(t)\partial_{t}\phi\, E_{p} \\
&= \int_{\mathcal{D}_{+}}\Box\phi\, E_{p} - \partial_{t_{0}}\left(\int_{\mathbb{R}^{n}}f(x)E_{p}(0, x)dx\right) - \int_{\mathbb{R}^{n}}g(x)E_{p}(0, x)dx.
\end{aligned}
\]

最终公式取如下形式：

\begin{equation}
\begin{aligned}
\phi(t_{0}, x_{0}) &= \int_{0}^{t_{0}}\int_{\mathbb{R}^{n}}E_{+}(t_{0}-t, x_{0}-x)\Box\phi(t, x)dtdx \\
&\quad - \partial_{t_{0}}\left(\int_{\mathbb{R}^{n}}E_{+}(t_{0}, x-x_{0})f(x)dx\right) - \int_{\mathbb{R}^{n}}E_{+}(t_{0}, x-x_{0})g(x)dx
\end{aligned}\label{eq:c4-145}
\end{equation}

把唯一性问题（我们稍后将单独处理）放在一边， 我们推出如下结果.

\begin{theorem}[Kirchhoff--Hadamard]\label{thm:c4-3.2}
初值问题 $\Box\phi = F$， $\phi(0, x) = f(x)$， $\partial_{t}\phi(0, x) = g(x)$ 对任意光滑函数 $f, g, F$ 都有唯一解， 由公式 \eqref{eq:c4-145} 给出.
\end{theorem}

\noindent\textbf{习题 2.}\quad 把公式 \eqref{eq:c4-145} 与 Laplacian 的 \eqref{eq:c4-137} 作比较.解释如果我们试图对 Laplace 方程证明与上面的\thmref{thm:c4-3.2} 类似的结果， 可能会出什么问题.

\noindent\textbf{习题 3.}\quad 证明：在特殊情形维数 $1+3$ 时， 公式 \eqref{eq:c4-145} 取更常见的 Kirchhoff 公式形式：

\begin{equation}
\begin{aligned}
\phi(t, x) &= \partial_{t}\left((4\pi t)^{-1}\int_{|x-y|=t}f(y)da(y)\right) + (4\pi t)^{-1}\int_{|x-y|=t}g(y)da(y) \\
&\quad + \int_{0}^{t}ds\frac{1}{t-s}\int_{|x-y|=t-s}\Box\phi(s, y)da(y)
\end{aligned}\label{eq:c4-146}
\end{equation}

推导 Kirchhoff 公式 \eqref{eq:c4-146} 的传统方法是先证明齐次情形， 即 $\Box\phi = 0$.事实上， 只需用优美的球平均方法证明 $f = 0$ 且 $g$ 任意的情形即可， 清晰的推导见 \cite{j}.齐次情形处理完之后， 就可以用 Duhamel 原理导出一般公式.具体如下：设 $W(t)g$ 表示齐次问题在数据 $f = 0$、$g$ 任意时的解 $\phi(t, \cdot)$.把它看作一族由 $t$ 参数化的算子， 把 $\mathbb{R}^{n}$ 上的光滑函数映到 $\mathbb{R}^{n}$ 上的光滑函数.于是我们需要验证方程 $\Box\phi = F$ 的解由公式

\begin{equation}
\phi(t, x) = \int_{0}^{t}W(t-s)F(s, \cdot)ds\label{eq:c4-147}
\end{equation}

给出.

\noindent\textbf{习题 4.}\quad 证明这一论断.

\noindent\textbf{习题 5.}\quad 如果在 Cauchy 问题的表述中把超曲面 $t = 0$ 换成由 $t = h(x)$ 给出的更一般的超曲面 $\Sigma_{0}$， 会发生什么？证明：如果超曲面是类空的， 即 $|\nabla h(x)| < 1$， 那么可以推出与 \eqref{eq:c4-145} 类似的公式.如果曲面在某个点 $(t_{0}, x_{0}) \in \Sigma_{0}$ 处变为类时的， 即 $|\nabla h(x_{0})| > 1$， 会发生什么？证明：在光锥 $t = |x|$ 上给定初值与法向导数的 Cauchy 问题， 一般地在锥的空间内部没有解.当你试图用 $H(t-|x|)$ 替换 $H(t)$ 来推导光锥上数据的表示公式时， 会发生什么？

\noindent\textbf{习题 6.}\quad 假设 $\rho$ 有紧支集， 且在紧集 $K$ 之外光滑.设 $u_{+}$ 是 $\Box u = \rho$ 在 $\mathbb{R}^{n+1}$ 中的解 $E_{+}*\rho$.证明 $u$ 在由 $K$ 出发的光锥集 $\{k + (|x|, x) : k \in K, x \in \mathbb{R}^{n}\}$ 之外光滑.

\begin{proof}[\thmref{thm:c4-3.1} 在 $\mathbb{R}^{1+3}$ 中的证明]
首先注意， 我们可以用球坐标把 $\Box$ 写成如下形式：

\[
\Box = -\partial_{t}^{2} + \Delta = -\partial_{t}^{2} + \partial_{r}^{2} + \frac{2}{r}\partial_{r} + r^{-2}\Delta_{\mathbb{S}^{n-1}}
\]

我们必须验证

\[
E_{+}^{(1+3)}(t, x) = -\frac{1}{2}\pi^{-1}H(t)\delta_{0}(t^{2}-|x|^{2}) = -\frac{1}{4\pi}r^{-1}\delta(t-r)
\]

其中 $r = |x|$.于是， 由于 $\Box\phi = -r^{-1}(\partial_{t}+\partial_{r})(\partial_{t}-\partial_{r})(r\phi) + \Delta_{\mathbb{S}^{2}}\phi$， 取 $\psi(t, r\omega) = (\partial_{t}-\partial_{r})\bigl(r\phi(t, r\omega)\bigr)$， 我们有

\[
\begin{aligned}
\langle E_{+}, \Box\phi\rangle &= \frac{1}{4\pi}\int_{\mathbb{S}^{2}}\int_{0}^{\infty}\int_{0}^{\infty}\delta(t-r)(\partial_{t}+\partial_{r})\psi\, dtdrdS_{\omega} \\
&= \frac{1}{4\pi}\int_{\mathbb{S}^{2}}dS_{\omega}\left(\int_{0}^{\infty}\frac{d}{dr}\psi(r, r)dr\right) \\
&= -\psi(0, 0) = \phi(0)
\end{aligned}
\]

于是 $\Box E_{+} = \delta_{0}$， 正如所愿.
\end{proof}

下面我们给出波动方程在特殊情形 $\mathbb{R}^{1+3}$ 中基本解的另一种推导.这就是所谓的几何光学推导.我们寻找 $\Box\phi = 0$ 的如下形式的解：

\begin{equation}
E = A\delta(u)\label{eq:c4-148}
\end{equation}

其中实函数 $A$ 与 $u$ 待定.这里 $\delta(u)$ 就是 $\delta_{0}$ 被 $u$ 的回拉， 正如第 4 小节的例 2 所讨论的.简单的计算得到

\[
m^{\alpha\beta}\partial_{\alpha}\partial_{\beta}\bigl(A\delta(u)\bigr) = m^{\alpha\beta}\left(\partial_{\alpha}\partial_{\beta}A\, \delta(u) + \left(2\partial_{\alpha}A\partial_{\beta}u + \Box u\right)\delta'(u) + \partial_{\alpha}u\partial_{\beta}u\, \delta''(u)\right)
\]

为消去 $\delta''(u)$ 的系数， 我们需要选取 $u$ 使得

\begin{equation}
m^{\alpha\beta}\partial_{\alpha}u\partial_{\beta}u = 0.\label{eq:c4-149}
\end{equation}

这就是 Minkowski 空间中著名的程函方程.一个简单的解族由 $u(t, x) = t-t_{0}-|x-x_{0}|$ 给出（$(t_{0}, x_{0})$ 为给定点）， 其水平超曲面就是以 $(t_{0}, x_{0})$ 为顶点的后向光锥.为了我们的目的， 我们选取 $u = t-|x|$.其次， 为消去 $\delta'(u)$ 的系数， 我们需要选取 $A$ 使得\footnote{事实证明， 下面的方程可以解释为沿后向零锥 $t-|x| = 0$ 的生成元的输运方程.}

\[
2\partial_{\alpha}A\, \partial^{\alpha}u + \Box u = 0.
\]

容易验证选取 $A = |x|^{-1}$ 即可.最后只剩下计算含 $\delta(u)$ 的项， 即

\[
(\Box A)\delta(u) = (-\Delta|x|^{-1})\delta(u) = -4\pi\delta_{0}(x)\delta(u) = -\delta_{0}(t, x)
\]

其中第一个 $\delta_{0}(x)$ 是 $\mathbb{R}^{3}$ 中的 $\delta$ 函数， 而最后的 $\delta_{0}$ 是 $\mathbb{R}^{1+3}$ 中所需的 $\delta$ 函数.于是 $E_{+}^{1+3} = -\frac{1}{4\pi}\frac{1}{|x|}\delta(t-|x|)$， 正如所愿.

\noindent\textbf{习题 7.}\quad 证明涉及分布乘积的最后一步.

基本解 $E_{+}$ 的唯一性.只需证明\thmref{thm:c4-3.2} 中一般 Cauchy 问题解的唯一性即可.

\noindent\textbf{习题 8.}\quad 验证上面的论断.

我们从一个涉及能量-动量张量的简单计算开始.为了高效计算， 记住我们对时空指标 $\alpha, \beta = 0, 1, \ldots, n$ 使用求和约定会有所帮助.我们还将使用标准的几何约定：相对于度量上升或下降指标.于是， 若 $U^{\alpha}$ 是向量（即所谓的逆变向量）， 我们定义协变向量 $U_{\beta} = m_{\beta\alpha}U^{\alpha}$.类似地， 若 $V_{\alpha}$ 是协变向量， 我们定义 $V^{\beta} = m^{\beta\alpha}V_{\alpha}$.

\begin{proposition}\label{prop:c4-3.3}
设

\begin{equation}
Q_{\alpha\beta}[\phi] = \partial_{\alpha}\phi\partial_{\beta}\phi - \frac{1}{2}m_{\alpha\beta}\left(m^{\gamma\delta}\phi_{\gamma}\phi_{\delta}\right) = \partial_{\alpha}\phi\partial_{\beta}\phi - \frac{1}{2}m_{\alpha\beta}\phi^{\gamma}\phi_{\gamma}\label{eq:c4-150}
\end{equation}

为 $\Box$ 的所谓的能量-动量张量.则\footnote{根据上升与下降指标的约定， $\partial^{\beta}Q_{\alpha\beta} = m^{\beta\gamma}\partial_{\gamma}Q_{\alpha\beta}$.这没有问题， 因为 $\partial_{\gamma}m_{\alpha\beta} = 0$.}

\begin{equation}
\partial^{\beta}Q_{\alpha\beta} = \Box\phi\, \partial_{\alpha}\phi.\label{eq:c4-151}
\end{equation}
\end{proposition}

特别地， 若 $\Box\phi = 0$， 则 $\partial^{\beta}Q_{0\beta} = 0$.现在考虑点 $p(t_{0}, x_{0}) \in \mathbb{R}^{1+n}$ 以及上面引入的方程 \eqref{eq:c4-149} 的解 $u(t, x) = t-t_{0}-|x-x_{0}|$.再设 $t_{1} < t_{2} < t_{0}$， 考虑分布 $H(u)H(t-t_{1})H(t_{2}-t)$（$H(t)$ 为 Heaviside 函数）， 并像我们多次做过的那样分部积分， 以导出恒等式：

\[
\begin{aligned}
0 &= \int_{\mathbb{R}^{1+n}}H(u)H(t-t_{1})H(t_{2}-t)\partial^{\beta}Q_{0\beta} \\
&= -\int_{\mathbb{R}^{1+n}}H'(u)H(t-t_{1})H(t_{2}-t)\partial^{\beta}u\, Q_{0\beta} \\
&\quad + \int_{\mathbb{R}^{1+n}}H(u)H(t-t_{1})H'(t_{2}-t)Q_{00} \\
&\quad - \int_{\mathbb{R}^{1+n}}H(u)H'(t-t_{1})H(t_{2}-t)Q_{00}
\end{aligned}
\]

这个恒等式可以改写成如下形式：

\begin{equation}
\int_{\mathcal{D}(t_{2})}Q_{00} + \int_{\mathcal{N}(t_{1}, t_{2})}Q_{0\beta}L^{\beta} = \int_{\mathcal{D}(t_{1})}Q_{00}\label{eq:c4-152}
\end{equation}

其中 $\mathcal{D}(t_{1})$、$\mathcal{D}(t_{2})$ 是穿过实心光锥 $|x-x_{0}| \leq t-t_{0}$ 的 $t$-截面， $\mathcal{N}(t_{1}, t_{2})$ 表示光锥 $|x-x_{0}| = t-t_{0}$ 在截面 $t = t_{1}$ 与 $t = t_{2}$ 之间的部分， 并且

\[
L^{\beta} = -\partial^{\beta}u = m^{\beta\gamma}\partial_{\gamma}u.
\]

容易验证

\[
Q_{00} = \frac{1}{2}\left(|\partial_{t}\phi|^{2} + \sum_{i=1}^{n}|\partial_{i}\phi|^{2}\right) = \frac{1}{2}\left(|\partial_{t}\phi|^{2} + |\nabla\phi|^{2}\right).
\]

并且我们有

\noindent\textbf{习题 9.}\quad 证明

\[
Q_{0\beta}L^{\beta} \geq 0.
\]

于是我们推出

\begin{theorem}[能量不等式]\label{thm:c4-3.4}
对由曲面 $t-t_{0}-|x-x_{0}| = 0$、$t = t_{1}$ 与 $t = t_{2}$ 围成的实心区域的一个邻域中的波动方程 $\Box\phi = 0$ 的每个解， 都有

\[
\int_{D(t_{2})}\frac{1}{2}\left(|\partial_{t}\phi|^{2} + |\nabla\phi|^{2}\right) \leq \int_{D(t_{1})}\frac{1}{2}\left(|\partial_{t}\phi|^{2} + |\nabla\phi|^{2}\right)
\]
\end{theorem}

特别地， 任何在 $D(t_{1})$ 上消失的光滑解也必在 $D(t_{2})$ 上消失.

\noindent\textbf{习题 10.}\quad 像我们先前从显式解推出那样， 从能量不等式推出 Cauchy 问题的有限传播速度.注意， 先前讨论的奇数空间维数 $n = 3+2k$ 中的强 Huygens 原理是一个非常特殊的现象， 它与波算子 $\Box$ 的精确形式有关.然而波动方程所展现的有限传播速度现象则更稳健， 它是许多相关的、如上能量不等式成立的方程所共有的特征.

最后， 我们在所有维数下验证前向基本解的有效性.

\begin{proof}[\thmref{thm:c4-3.1} 的证明(所有 $n$)]
我们证明公式（略去常数 $c_{n}$ 的绝对值）.

\[
\begin{aligned}
\int\!\!\int \Box\chi_{+}^{-\frac{n-1}{2}}(t^{2}-|x|^{2}) \cdot H(1-|t|)dxdt &= \int\!\!\int\partial_{t}^{2}\chi_{+}^{-\frac{n-1}{2}}(t^{2}-|x|^{2}) \cdot H(1-|t|)dxdt \\
&= \int\!\!\int\partial_{t}\chi_{+}^{-\frac{n-1}{2}}(t^{2}-r^{2})\delta(1-t)dxdt \\
&\quad - \int\!\!\int\partial_{t}\chi_{+}^{-\frac{n-1}{2}}(t^{2}-r^{2})\delta(1+t)dxdt \\
&= 2\int\!\!\int\partial_{t}\chi_{+}^{-\frac{n-1}{2}}(t^{2}-r^{2})\delta(1-t)dxdt \\
&= 2\int\!\!\int\left(\frac{-t}{r}\right)\partial_{r}\chi_{+}^{-\frac{n-1}{2}}(t^{2}-r^{2})\delta(1-t)dxdt
\end{aligned}
\]

取 $\delta(1-t)$ 的定义， 上式变为超曲面 $t = 1$ 上的积分， 我们把它写成极坐标（体积形式写为 $r^{n-1}drd\Omega_{n-1}$， 其中 $d\Omega_{n-1}$ 是 $\mathbb{R}^{n}$ 中单位球面上的曲面测度）.我们在此指出， 正如我们后面将要看到的， 一个分布并不必须是连续函数才能在低维子流形上有有意义的限制.

去掉不重要的因子 2， 我们继续：

\[
\begin{aligned}
&\int\!\!\int\left(\frac{-t}{r}\right)\partial_{r}\chi_{+}^{-\frac{n-1}{2}}(t^{2}-r^{2})\delta(1-t)dxdt \\
&= \int\!\!\int\frac{-1}{r}\partial_{r}\chi_{+}^{-\frac{n-1}{2}}(1-r^{2})r^{n-1}drd\Omega_{n-1} \\
&= |\mathbb{S}^{n-1}|\int_{0}^{\infty}-r^{n-2}\partial_{r}\chi_{+}^{-\frac{n-1}{2}}(1-r^{2})dr
\end{aligned}
\]

于是我们的命题归结为证明数

\[
-\int_{0}^{\infty}r^{n-2}\partial_{r}\chi_{+}^{-\frac{n-1}{2}}(1-r^{2})dr
\]

是正的.尽管我们还没有定义这个数\footnote{注意， 然而（由乘积法则）$r^{n-2} \cdot H(r)$ 具有相当好的正则性， 它的唯一奇点位于远离 $r = 1$ 处.}， 我们还是通过归纳证明它的正性， 并按 $n$ 的奇偶性分情形.对 $n = 1+2k$， 我们分部积分并注意到边界项消失， 得到

\[
\begin{aligned}
\int_{0}^{\infty}r^{n-2}\partial_{r}\chi_{+}^{-\frac{n-1}{2}}(1-r^{2})dr &= \int_{0}^{\infty}r^{2k}\left(\frac{-\partial_{r}}{2r}\right)^{k+1}H(1-r^{2})dr \\
&\equiv I_{k} \\
&= 2^{-(k+1)}\int\partial_{r}\left(r^{2k-1}\cdot H(r)\right)\left(\frac{-\partial_{r}}{2r}\right)^{k}H(1-r^{2})dr \\
&= 2^{-(k+1)}(2k-1)\int_{0}^{\infty}r^{2(k-1)}\left(\frac{-\partial_{r}}{2r}\right)^{k}H(1-r^{2})dr \\
&= 2^{-(k+1)}(2k-1)I_{k-1}
\end{aligned}
\]

只要我们能计算

\[
\begin{aligned}
I_{1} &= \int_{0}^{\infty}r^{2}\left(\frac{-\partial_{r}}{2r}\right)H(1-r^{2})dr = \int_{0}^{\infty}\delta(1-r^{2})dr \\
&= \int_{0}^{\infty}\frac{1}{1+r}\delta(1-r)dr = 1/2 > 0
\end{aligned}
\]

最后一个积分就应该由对 $k$ 的归纳是正的.一旦适当地选取 $c_{n} > 0$， 这就证明了奇数空间维数情形下的命题.情形 $n = 2+2k$ 类似.
\end{proof}
\section{热算子 \texorpdfstring{$\mathcal{H}$}{H}}

我们考虑作用在定义于 $\mathbb{R} \times \mathbb{R}^{n} = \mathbb{R}^{n+1}$ 上的函数的热算子 $\mathcal{H} = \partial_{t} - \Delta$.寻找 $\mathcal{H}u = 0$ 的球对称解是有意义的：也就是说， 考虑函数 $u(t, x) = u(t, |x|) = u(t, r)$.用这种方式可以找到一族局部可积解 $E_{c}(t, x) = cH(t)\, t^{-\frac{n}{2}}e^{-|x|^{2}/4t}$， 其中 $H(t)$ 是 Heaviside 函数（尽管通过 傅里叶变换处理更容易）.

事实上， 对所有 $(t, x) \neq (0, 0)$ 都有 $\mathcal{H}(E_{c}) = 0$.下面我们证明， 在整个空间上 $\mathcal{H}(E_{c})$ 与 $\delta_{0}$ 成正比， 并且可以确定常数 $c = c_{n} = 2^{-n}\pi^{-\frac{n}{2}}$， 使得相应的 $E = E_{c}$ 是 $\mathcal{H}$ 的基本解， 即 $\mathcal{H}(E) = \delta_{0}$.

我们可以非常容易地通过考虑抛物型缩放 $(t, x) \mapsto (\alpha^{2}t, \alpha x)$ 来推断 $\mathcal{H}E_{c} = C\delta_{0}$（$C$ 为某常数， 可能为 0）.为确定该常数， 我们可以使用任意检验函数， 取 $H(1-t)$ 会很方便.我们的做法大体上正是如此.

设 $\phi \in \mathcal{C}_{0}^{\infty}(\mathbb{R}^{n+1})$，

\[
\begin{aligned}
\langle \mathcal{H}(E), \phi\rangle &= \langle E, \mathcal{H}^{t}\phi\rangle = -\int E(t, x)(\partial_{t}+\Delta)\phi(t, x)dxdt \\
&= -\lim_{\varepsilon\to 0^{+}}\int_{\varepsilon}^{\infty}\int_{\mathbb{R}^{n}}E(t, x)(\partial_{t}+\Delta)\phi(t, x)dxdt \\
&= \lim_{\varepsilon\to 0^{+}}\int_{\varepsilon}^{\infty}\int_{\mathbb{R}^{n}}(\partial_{t}-\Delta)E(t, x)\phi(t, x)dxdt + \lim_{\varepsilon\to 0^{+}}\int_{\mathbb{R}^{n}}E(\varepsilon, x)\phi(\varepsilon, x)dx \\
&= \lim_{\varepsilon\to 0^{+}}\int_{\mathbb{R}^{n}}E(\varepsilon, x)\phi(\varepsilon, x)dx = c_{n}\lim_{\varepsilon\to 0^{+}}\varepsilon^{-n/2}\int_{\mathbb{R}^{n}}e^{-|x|^{2}/4\varepsilon}\phi(\varepsilon, x)dx
\end{aligned}
\]

现在我们作变量替换 $x = 2\varepsilon^{1/2}y$：

\[
\begin{aligned}
\langle \mathcal{H}(E), \phi\rangle &= 2^{n}c_{n}\lim_{\varepsilon\to 0^{+}}\int_{\mathbb{R}^{n}}\phi(\varepsilon, 2\varepsilon^{1/2}y)e^{-|y|^{2}}dy = 2^{n}c_{n}\phi(0, 0)\int_{\mathbb{R}^{n}}e^{-|y|^{2}}dy \\
&= \phi(0, 0)
\end{aligned}
\]

撇开事实 $\int_{\mathbb{R}^{n}}e^{-|y|^{2}}dy = \pi^{n/2}$（后面会证明）不谈， 我们已经证明

\begin{equation}
E(t, x) = (4\pi t)^{-n/2}H(t)e^{-|x|^{2}/4t}\label{eq:c4-153}
\end{equation}

是 $\mathcal{H}$ 的基本解.注意， 对任意固定的 $t > 0$， $E(t, x)$ 的支集遍布整个 $\mathbb{R}^{n}$， 这意味着热方程（与波动方程相反）表现出``无限传播速度''.这一现象与热算子的抛物型缩放有关， 后者与 $-\partial_{t}^{2} + \sum_{i}\partial_{i}^{2}$ 相反， 赋予时间与空间不同的``单位''.还要注意， $E(t, x)$ 在 $t > 0$ 时是光滑的；这一事实将导致初值问题 $\mathcal{H}\psi = 0$（在 $(0, \infty) \times \mathbb{R}^{n}$ 上）， $\psi(0, x) = \psi_{0}(x)$ 的瞬时磨光.

\noindent\textbf{习题 1.}\quad 推导初值问题的表示公式.为什么不可能对任意初值数据把热方程沿时间向后求解？

\noindent\textbf{习题 2.}\quad 证明：上面 Cauchy 问题的表示公式对足够光滑的数据确实给出经典解.验证得到了正确的边界值.（这与 Cauchy-Riemann 方程组的情形形成对照.）

\noindent\textbf{习题 3.}\quad 写出 $\mathcal{H}\psi = 0$ 的 $\mathcal{C}^{2}$ 解在 $(0, T] \times \overline{\Omega}$ 内部的最大值原理， 其中 $\Omega$ 是开有界集.

\noindent\textbf{习题 4.}\quad 设 $\Psi : [0, T] \times \mathbb{R}^{n}$ 是热方程的解， 对固定时间在空间无穷远处快速衰减.证明``能量''

\[
e(t) = \int_{\mathbb{R}^{n}}|\Psi(t, x)|^{2}dx
\]

随时间递减.由此推出热方程的唯一性定理.

\noindent\textbf{习题 5.}\quad 人们常用 $e^{t\Delta}\psi_{0}$ 表示 $\mathcal{H}\psi = 0$（在 $(0, \infty) \times \mathbb{R}^{n}$ 上）， $\psi(0, x) = \psi_{0}(x)$ 的（典范）解在固定时间 $t > 0$ 超曲面上的限制.证明 $e^{(t+s)\Delta} = e^{t\Delta}e^{s\Delta}$ 作为快速衰减光滑函数上的时间演化算子成立.

\mtsubsecn{4.1}{\texorpdfstring{Schrödinger 算子 $S$}{Schrödinger 算子 S}}

Schrödinger 算子 $\mathcal{S} = i\partial_{t} + \Delta$ 有一个基本解， 表面上看与热算子的完全一样：

\begin{equation}
E(t, x) = (4\pi it)^{-n/2}H(t)e^{i|x|^{2}/4t}\label{eq:c4-154}
\end{equation}

然而， 当然， 指数因子 $e^{i|x|^{2}/4t}$ 中 $i$ 的存在造成了天壤之别.

\noindent\textbf{习题 1.}\quad 证明（对适当选取的 $\log$ 的枝切）局部可积函数 $E$ 确实是 $\mathcal{S}$ 的基本解.

\noindent\textbf{习题 2.}\quad 与热方程的习题 5 类似， 人们把 Schrödinger 方程的时间演化算子记为 $e^{it\Delta}$.证明 $e^{it\Delta}$ 是酉算子， 其意义是量

\[
\int_{\mathbb{R}^{n}}|\Psi(t, x)|^{2}dx
\]

随时间保持不变.

\part{分析的基本工具}

% ==== part05 第 5 章 傅里叶变换 (Fourier transform) ====
% 源: MinerU_latex_lecture nots.tex 第 6716--7788 行;pdf_text.txt 第 5552--6455 行(PDF 第 105--120 页)
%
% 术语(新增,键=中文=英文):
%   fourier-transform=傅里叶变换、physical-space=物理空间、frequency-space=频率空间、
%   plane-wave=平面波、inversion-formula=反演公式、rapidly-decreasing=快速递减、
%   schwartz-space=Schwartz 空间、tempered-distribution=缓增分布、
%   plancherel-identity=Plancherel 恒等式、parseval-identity=Parseval 恒等式、
%   modulation=调制、uncertainty-principle=不确定性原理、
%   fourier-multiplier-operator=Fourier 乘子算子、bernstein-inequality=Bernstein 不等式、
%   dimensional-analysis=量纲分析、annihilator-subgroup=零化子群、
%   domain-of-influence=影响域、fourier-integral-operator=Fourier 积分算子
% 注: 分布、检验函数、卷积、支集、主值、张量积、解析延拓、截断函数、特征函数等已由 part03 引入,
%   此处直接使用中文;参数层(partgf1)、等距(part01)等亦同.
%
% REFS(需跨章/待解析项):
%   - 公式 (129)(1/(x+i0)): 第 3 章公式,part03 已有 \label{eq:c3-129},正文用 \eqref{eq:c3-129}.
%   - 公式 (153)、(154)(物理空间公式)与 Kirchhoff 公式 (146): 均在第 4 章(part04 未写),
%     正文用纯文本 "公式 (153) 与 (154)"、"Kirchhoff 公式 (146)",待 part04 完成后统一改 \eqref.
%   - "前一章"(第 4 章)基于基本解显式计算导出的公式: 文本引用,无需标签.
%   - 脚注 7 中 "[?] 第一卷": 原书文献编号破损,保留 [?].
%   - OCR-BROKEN: 第 7144 行公式 e^{i\xi} 缺 x,按数学正确性补为 e^{i\xi x};
%     第 7149--7150 行 "\mathrm{\bar{\ s t}{b a s i s}}^{\bar{\ y}}" 实为 "basis"(正交归一基);
%     第 6757 行孤立上标 "2"(pdf 页 106)无对应脚注文本,视为 OCR 噪音删除;第 6793、6812、
%     6934、7029 行孤立 "2" 同理删除;第 7421 行孤立 "厂" 删除.
%   - 脚注 3 中 "i d/dx_i f" 按上下文理解为 i\partial/\partial x_i.
%
% 编号说明:
%   - 原书第 5 章显式编号公式为 (155)--(178) 共 24 个,按指南全部行间公式自动编号,
%     显示编号与原书可能不对应(与 part03 相同问题,由协调员统一核对).
%   - 公式 (157)、(158) 原书(及 OCR)缺 d\xi/dy 因子,按原文保留.
%   - 第 4 节编号: Proposition 4.1 → \mtsubsec(4.2 物理空间的局部化) → \mtsubsec(4.3 频率空间的局部化)
%     → Lemma 4.4 → Lemma 4.5 → Proposition 4.6,共享计数器自然推进,无跳号.
%   - 缓增分布之 Definition 与 Example 原书无编号,用 definition* / example*;
%     Proposition 4.6 后的 Remark 无编号,用 remark*.
%   - 习题: 第 2 节 Exercise 2.5/2.6 带节号,按原文写 "习题 2.5."/"习题 2.6.";
%     第 3、5 节习题按节重起 "习题 1."/"习题 2."/"习题 3.";第 4 节一个未编号 "习题.".
%   - 原文 "Schwarz space" 与 "Schwartz space" 混用,统一为 Schwartz 空间(术语表);
%     原文 "Kirchof/Kirchoff" 拼写不一,统一为 Kirchhoff;原文 "Berenstein estimate" 拼写有误,
%     统一为 Bernstein 估计;原文 "Plancharel" 拼写有误,统一为 Plancherel;原文 "Fourer" 统一为 Fourier.

\chapter{傅里叶变换}

\section{基本性质}

回忆:若 $f \in L^1(\mathbb{R}^n)$, 则 \newterm{fourier-transform}{傅里叶变换}{Fourier transform} $\mathcal{F}(f) = \hat{f}$ 定义为连续函数

\begin{equation}
\hat{f}(\xi) = \int f(x)e^{-ix\xi}\, dx\label{eq:c5-155}
\end{equation}

若 $\hat{f} \in L^1(\widehat{\mathbb{R}}^n)$, 我们有 \newterm{inversion-formula}{反演公式}{inversion formula}

\begin{equation}
f(x) = (2\pi)^{-n}\int\hat{f}(\xi)e^{ix\xi}\, d\xi, \label{eq:c5-156}
\end{equation}

其证明我们将在后面给出.为了在概念上区分两者, 我们把 $f$ 所在的 $\mathbb{R}^n$ 称为 \newterm{physical-space}{物理空间}{physical space}, 把 $\hat{f}$ 所在的点集 $\xi$ 称为 \newterm{frequency-space}{频率空间}{frequency space}.我们用 $\widehat{\mathbb{R}}^n$ 表示频率空间, 并用归一化测度 $\frac{d\xi}{(2\pi)^n}$ 赋予它.

反演公式为我们提供了关于 傅里叶变换作用的宝贵启发式理解.我们看到, $f(x)$ 可以写成某种 \newterm{plane-wave}{平面波}{plane wave} $(x\mapsto e^{ix\cdot\xi})$ 的线性组合, 而测度 $\hat{f}(\xi)\frac{d\xi}{(2\pi)^n}$ 描述了 $f$ 在频率空间上的分布.若我们把平面波 $e^{ix\cdot\xi}$ 看作 $\mathbb{R}^n$ 上平移算子的特征向量, 则可以把 傅里叶变换视为同时对角化平移的一种尝试.类似地, 若把平面波 $e^{ix\cdot\xi}$ 看作算子族 $\{i\frac{\partial}{\partial x_j}: j = 1, \ldots, n\}$ 的特征向量——这些算子关于 $L^2$ 内积是自伴的(限制在适当的定义域上)——那么我们就会看到, 微分也被 傅里叶变换对角化了.

带着这些启发, 我们开始看到 傅里叶变换如何可能对分析有用.例如, 若 $\hat{f}$ 集中在某个频率 $\xi_0 \in \widehat{\mathbb{R}}^n$ 附近, 我们预期 $f$ 在某种程度上表现得像 $e^{ix\cdot\xi_0}$ 附近的平面波.例如, $f$ 可能允许有界、复解析地延拓到 $\mathbb{C}^n$ 的某一部分.我们还预期 $\partial_{x_j}f(x)\sim i\xi_j^0 f(x)$, 从而微分成为更容易研究的运算.事实上, 当我们稍后遇到 Littlewood-Paley 理论时, 主要思想将是把一般函数分解为频率局部化的分量, 分别分析这些分量, 然后再重新组装.

关于 傅里叶变换的另一个重要原理是光滑性与衰减在物理空间和频率空间之间的对偶性.直觉上, 一个在物理空间中其图形具有突然跳跃或尖峰的函数, 必然"由"任意大的频率"构成", 而当 $\hat{f}$ 具有紧支集时, $f$ 必是全局"驯良"的.类似地, 当 $\hat{f}$ 非常光滑时, 由于反演公式中邻近平面波之间的干涉(抵消), $f$ 可以在无穷远处衰减.这些基本原理在 Fourier 分析中处处有不止一种而是许多种正式的表现形式, 因此让我们在继续发展理论的过程中牢记它们.

反演公式在 Gauss 函数 $G(x) = e^{-|x|^2/2}$ 的情形采取了特别具体的形式.

\begin{lemma}\label{lem:c5-1.1}
如下计算对一元函数以及 $a, b \in \mathbb{R}$, $b > 0$ 成立:

\begin{equation}
\int_{-\infty}^{\infty} e^{iax}e^{-bx^2} = \left(\frac{\pi}{b}\right)^{1/2}e^{-a^2/4b}\label{eq:c5-157}
\end{equation}

于是, 在 $\mathbb{R}^n$ 中, 对 $t > 0$,

\begin{equation}
\int_{\mathbb{R}^n} e^{ix\cdot y}e^{-ty^2} = \left(\frac{\pi}{t}\right)^{n/2}e^{-|x|^2/4t}\label{eq:c5-158}
\end{equation}

特别地,

\begin{equation}
\mathcal{F}(G)(\xi) = (2\pi)^{n/2}G(\xi)\label{eq:c5-159}
\end{equation}
\end{lemma}

\begin{proof}
在复域中作变量替换 $z = b^{1/2}x - \frac{a}{2b^{1/2}}i$, 并记 $\Gamma$ 为围道 $\mathrm{Im}(z) = -\frac{a}{2b^{1/2}}$:

\[
\int_{-\infty}^{\infty} e^{iax}e^{-bx^2}\, dx = \frac{e^{-a^2/4b}}{b^{1/2}}\int_\Gamma e^{-z^2}\, dz = \frac{e^{-a^2/4b}}{b^{1/2}}\int_{-\infty}^{\infty} e^{-x^2}\, dx
\]

(利用标准围道变形论证).现在计算积分 $J = \int_{-\infty}^{\infty} e^{-x^2}\, dx = \pi^{1/2}$:我们观察到, 通过极坐标, $J^2 = \int_{\mathbb{R}^2} e^{-|x|^2}\, dx = \pi$, 由此即得 \eqref{eq:c5-157}.公式 \eqref{eq:c5-158} 随之立即成立.
\end{proof}

在回顾了 傅里叶变换的一些基本性质之后, 我们可以给出上述恒等式的另一个证明.

\begin{proposition}\label{prop:c5-1.2}
傅里叶变换是线性的, 并验证如下简单性质.

\begin{itemize}
\item 傅里叶变换把物理空间中的平移 $T_{x_0}f(x) = f(x-x_0)$ 变为频率空间中的调制 \newterm{modulation}{调制}{modulation}:$\mathcal{F}(T_{x_0}f)(\xi) = e^{-i\xi\cdot x_0}\hat{f}(\xi)$.
\item 傅里叶变换把物理空间中的调制 $M_{\xi_0}f(x) = e^{ix\cdot\xi_0}f(x)$ 变为频率空间中的平移:$\mathcal{F}(M_{\xi_0}f)(\xi) = \hat{f}(\xi-\xi_0)$.
\item 傅里叶变换把物理空间中的共轭变为频率空间中的共轭与反射, 即 $\mathcal{F}(\bar{f})(\xi) = \hat{f}(-\xi)$.
\item 傅里叶变换把物理空间中的卷积变为频率空间中的乘法:$\widehat{f*g} = \hat{f}\, \hat{g}$.
\item 傅里叶变换把物理空间中的偏导数变为频率空间中的乘法:$\mathcal{F}(\partial_{x_j}f)(\xi) = i\xi_j\hat{f}(\xi)$.
\item 傅里叶变换把物理空间中乘以 $x_j$ 变为频率空间中的偏导数 $\partial_{\xi_j}$:$\mathcal{F}(x_jf)(\xi) = i\partial_{\xi_j}\hat{f}(\xi)$.
\item 我们还有如下简单的自对偶关系:

\[
\int f(x)\hat{g}(x)\, dx = \int\hat{f}(\xi)g(\xi)\, d\xi
\]

\item 傅里叶变换把物理空间中的伸缩 $S_\lambda f(x) = f(\lambda x)$ 变为 Fourier 空间中的对偶伸缩:$\mathcal{F}(S_\lambda f)(\xi) = \lambda^{-n}\hat{f}(\xi/\lambda)$.注意, $S_\lambda(f)$ 保持大小, 即 $\|S_\lambda f\|_{L^\infty} = \|f\|_{L^\infty}$, 而对偶伸缩 $S_\lambda^* f = \lambda^{-n}f(x/\lambda)$ 保持质量, 即 $\|S_\lambda^* f\|_{L^1} = \|f\|_{L^1}$.
\end{itemize}
\end{proposition}

\begin{proof}
上述性质几乎全部归结为关于指数的简单恒等式:只要把 $f$ 与 $g$ 特殊化为物理空间或频率空间中的点质量\footnote{有限 Borel 测度的 傅里叶变换是连续函数 $\hat{\mu}(\xi) = \int_{\mathbb{R}^n} e^{-ix\cdot\xi}\, d\mu(x)$.}(从而在偶空间中就是平面波)即可——这些恒等式本身甚至可以被看作这些特殊情形的连续(双)线性延拓.
\end{proof}

利用这些性质, 我们可以给出 \eqref{eq:c5-157} 的另一个证明.由伸缩恒等式, 只需考虑 $b = \frac{1}{2}$ 并计算 $G(x) = e^{-x^2/2}$ 的 傅里叶变换.对恒等式 $\frac{dG}{dx} = -xG(x)$ 取 傅里叶变换并应用上述性质, 我们看到 $\hat{G}$ 在 $\xi$ 中满足同样的微分方程, 因此具有形式 $\hat{G}(0)e^{-\xi^2/2}$.由于我们已经证明了 $\hat{G}(0) = \int e^{-x^2/2} = \pi^{1/2}$, 这就完成了 \eqref{eq:c5-157} 的第二个证明.

设 $G_{\lambda, x_0, \xi_0}(x) = e^{ix\cdot\xi_0}G((x-x_0)/\sqrt{\lambda})$ 是平移、调制、伸缩后的 Gauss 函数.于是,

\[
\begin{aligned}
\mathcal{F}(G_{\lambda, x_0, \xi_0})(\xi) &= \lambda^{n/2}e^{-i(\xi-\xi_0)\cdot x_0}\int e^{-i\sqrt{\lambda}\, y\cdot(\xi-\xi_0)}G(y)\, dy \\
&= (\pi\lambda)^{n/2}G\big(\sqrt{\lambda}(\xi-\xi_0)\big)
\end{aligned}
\]

我们可以把这个结果解释为:$G_{\lambda, x_0, \xi_0}$ 局部化在空间位置 $x_0$ 附近, 空间展布 $\Delta x\approx\sqrt{\lambda}$;并局部化在频率位置 $\xi_0$ 附近, 频率展布 $\Delta\xi = 1/\sqrt{\lambda}$.注意 $\Delta x\cdot\Delta\xi\approx 1$, 因此我们这样同时在物理空间和频率空间中局部化的能力似乎是有限的.令人惊讶的是, 这种构造在某种意义上是我们能做到的最好情形, 它也是我们与 Fourier 分析的\newterm{uncertainty-principle}{不确定性原理}{uncertainty principle}的第一次相遇;该原理以各种形式断言, 同时在这两个频率空间和物理空间中局部化的精度存在一个界限.

我们现在来证明光滑性与衰减之间对偶性的第一个重要表现形式.\footnote{我们指出, 该命题对没有 $L^1$ 密度函数的测度不成立——例如, 点质量的 傅里叶变换是平面波, 它完全没有衰减.因此, 一定量的光滑性是必要的.}

\begin{proposition}[Riemann--Lebesgue]\label{prop:c5-1.3}
对任意的 $f \in L^1(\mathbb{R}^n)$, 我们有 $\|\hat{f}\|_{L^\infty}\lesssim\|f\|_{L^1}$.此外, 当 $|\xi|\to\infty$ 时, $\hat{f}(\xi)\to 0$.
\end{proposition}

\begin{proof}
只有最后一个陈述需要论证.注意, 若 $f\in\mathcal{C}_0^\infty(\mathbb{R}^n)$, 则我们可以用分部积分得出结论:$\hat{f}$ 快速衰减.事实上, 对任何多重指标 $\alpha$, $|\alpha| = k\in\mathbb{N}$,

\[
\begin{aligned}
\xi^\alpha\hat{f}(\xi) &= i^k\int\partial_x^\alpha e^{-ix\xi}f(x)\, dx = (-i)^k\int e^{-ix\xi}\partial_x^\alpha f(x)\, dx \\
|\xi^\alpha\hat{f}(\xi)| &\lesssim \int|\partial_x^\alpha f(x)\, dx| \leq C_\alpha
\end{aligned}
\]

对某个常数 $C_\alpha$ 成立.因此, $|\hat{f}(\xi)|\lesssim(1+|\xi|)^{-k}$, 这证明了该情形下的陈述.对一般的 $f\in L^1(\mathbb{R}^n)$, 给定 $\epsilon>0$, 我们可以选取 $g\in\mathcal{C}_0^\infty$, 使得 $\|f-g\|_{L^1}\leq\frac{\epsilon}{2}$.由前述结果, 当 $|\xi|>M = M_\epsilon$ 充分大时, $|\hat{g}(\xi)|\leq\frac{\epsilon}{2}$, 因此

\[
\sup_{|\xi|>M}|\hat{f}(\xi)| \leq \|f-g\|_{L^1(\mathbb{R}^n)} + \sup_{|\xi|>M}|\hat{g}(\xi)| \leq \epsilon
\]
\end{proof}


\section{Schwartz 空间与反演公式}

光滑函数上的许多运算自然地延拓到分布上(通过与 $\mathcal{C}_0^\infty$ 的对偶), 我们希望看到 傅里叶变换如何延拓到分布上.唯一可能的延拓必须与公式

\[
\langle\hat{u}, \phi\rangle = \langle u, \hat{\phi}\rangle
\]

相容, 但对分布 $u\in\mathcal{D}'(\mathbb{R}^n)$ 与检验函数 $\phi\in\mathcal{C}_0^\infty(\widehat{\mathbb{R}}^n)$, 这个公式没有意义, 因为不能保证 $\hat{\phi}$ 也具有紧支集.事实上, 作为不确定性原理的一个表现, $\phi$ 与 $\hat{\phi}$ 不可能同时具有紧支集, 除非 $\phi = 0$.因此, 如果我们想要一个推广 傅里叶变换的对称理论, 我们就要考虑一族关于 Fourier 对偶性表现良好的新的检验函数(以及相应的分布).

\begin{definition}\label{def:c5-2.1}
函数 $\phi\in\mathcal{C}^\infty(\mathbb{R}^n)$ 称为\newterm{rapidly-decreasing}{快速递减}{rapidly decreasing}的, 如果对一切多重指标 $\alpha, \beta$ 有

\[
\sup_{x\in\mathbb{R}^n}|x^\alpha\partial^\beta\phi(x)| < \infty.
\]
\end{definition}

这个所谓的 \newterm{schwartz-space}{Schwartz 空间}{Schwartz space} $\mathcal{S}(\mathbb{R}^n)$, 即快速递减函数的空间, 以通常的方式配备一个自然的 Frechet 拓扑.在此拓扑中, 函数序列 $\phi_j$ 收敛于零当且仅当对一切多重指标 $\alpha, \beta$, $x^\alpha\partial^\beta\phi_j$ 一致收敛于零.注意, $\mathcal{S}(\mathbb{R}^n)$ 包含紧支撑函数 $\mathcal{C}_0^\infty(\mathbb{R}^n)$.由于对 $1\leq p<\infty$, $\mathcal{C}_0^\infty(\mathbb{R}^n)$ 在 $L^p(\mathbb{R}^n)$ 空间中稠密, $\mathcal{S}(\mathbb{R}^n)$ 在 $L^p$ 空间中也稠密.还容易验证 $\mathcal{C}_0^\infty(\mathbb{R}^n)$ 在 $\mathcal{S}(\mathbb{R}^n)$ 中稠密.

我们有以下重要事实, 这就是我们在当前情形下考虑 Schwartz 空间的原因:

\begin{proposition}\label{prop:c5-2.2}
傅里叶变换是从 $\mathcal{S}(\mathbb{R}^n)$ 到 $\mathcal{S}(\widehat{\mathbb{R}}^n)$ 的同构, 其逆由反演公式 \eqref{eq:c5-156} 给出.此外, 对一切 $f, g\in\mathcal{S}(\mathbb{R}^n)$, 我们有 \newterm{plancherel-identity}{Plancherel 恒等式}{Plancherel identity}

\begin{equation}
(f, g)_{L^2} = \int_{\mathbb{R}^n} f(x)\bar{g}(x)\, dx = (2\pi)^{-n}\int\hat{f}\, \overline{\hat{g}}\, d\xi = (\hat{f}, \hat{g})_{L^2(\widehat{\mathbb{R}}^n)}\label{eq:c5-160}
\end{equation}

特别地, 我们有 \newterm{parseval-identity}{Parseval 恒等式}{Parseval identity} $\|f\|_{L^2(\mathbb{R}^n)} = \|\mathcal{F}(f)\|_{L^2(\widehat{\mathbb{R}}^n)}$.
\end{proposition}

\begin{proof}
注意 $|\xi^\alpha\partial^\beta\hat{\phi}(\xi)| = |\widehat{x^\beta\partial^\alpha\phi}|$, 且 $\partial^\alpha\phi(x)$ 的衰减比 $|x|^{-|\beta|-n-1}$ 更快.于是我们容易推出 $\mathcal{F}$ 把 $\mathcal{S}(\mathbb{R}^n)$ 映入 $\mathcal{S}(\widehat{\mathbb{R}}^n)$.设 $Rf(x) = f(-x)$, 并定义 $T = R\mathcal{F}^2$.注意, $T$ 与偏导数 $\partial_j$ 及乘以 $x_j$ 都交换.事实上, 对一切 $j = 1, \ldots, n$,

\begin{equation}
T(\partial_j f) = \partial_j(Tf), \quad T(x_jf) = x_j(Tf)\label{eq:c5-161}
\end{equation}

反演公式由如下引理得出.
\end{proof}

\begin{lemma}\label{lem:c5-2.3}
满足 \eqref{eq:c5-161} 的线性算子 $T: \mathcal{S}(\mathbb{R}^n)\longrightarrow\mathcal{S}(\mathbb{R}^n)$ 必具有形式 $T\phi = c\phi$, 其中 $c$ 是某个常数.
\end{lemma}

\begin{proof}
由交换性质 \eqref{eq:c5-161}, 我们看到 $T$ 在多项式函数的代数上是线性的.作为这一线性的推论, 我们可以证明值 $T\phi(x_0)$ 只依赖于 $\phi(x_0)$ 在点 $x_0$ 处的值.例如在一维情形, 若 $\phi$ 在点 $x_0$ 处为零, 则我们可以写成

\[
\phi(x) = (x-x_0)\int_0^1\phi'(x_0+t(x-x_0))\, dt = (x-x_0)\tilde{\phi}
\]

其中 $\tilde{\phi}\in\mathcal{S}(\mathbb{R})$.把 $T$ 作用于这个恒等式, 我们看到 $T\phi(x_0) = 0$ 也成立, 因此我们可以写成 $T\phi = f_\phi\phi$, 其中函数 $f_\phi$ 可能依赖于 $\phi$.

但 $f_\phi = f$ 不依赖于 $\phi$.若 $\psi$ 是任何另一个 Schwartz 函数, 线性组合 $\psi(x_0)\phi - \phi(x_0)\psi$ 在点 $x_0$ 处为零, 对 $T$ 应用同样的性质, 我们得出结论:在 $\psi$ 与 $\phi$ 同时非零的任何点 $x_0$ 处, $f_\phi(x_0) = f_\psi(x_0)$.显然, 为了使 $T$ 把 $\mathcal{S}(\mathbb{R}^n)$ 映入自身, 函数 $f$ 必须是光滑的;而为了使 $T$ 与微分交换, $f$ 必须是常数.

为确定常数, 我们只需指出:根据\lemref{lem:c5-1.1}, 我们有 $T(G) = \bigl((2\pi)^{n/2}\bigr)^2G = (2\pi)^nG$.因此常数 $c = (2\pi)^n$, 这就完成了对 Schwartz 函数的反演公式与命题的证明.

Plancherel 与 Parseval 恒等式是反演公式的直接推论.
\end{proof}

\begin{corollary}\label{cor:c5-2.4}
如下性质对 $\mathcal{S}(\mathbb{R}^n)$ 中的一切函数成立:

\[
\begin{aligned}
\int\hat{\phi}\psi\, dx &= \int\phi\hat{\psi}\, d\xi \\
\int\phi\bar{\psi}\, dx &= (2\pi)^{-n}\int\hat{\phi}\, \overline{\hat{\psi}}\, dx \\
\widehat{\phi*\psi} &= \hat{\phi}\, \hat{\psi} \\
\widehat{\phi\psi} &= (2\pi)^{-n}\int\hat{\phi}(\xi-\eta)\hat{\psi}(\eta)\, d\eta = \hat{\phi}*\hat{\psi}
\end{aligned}
\]

最后一个卷积是关于 $\widehat{\mathbb{R}}^n$ 上的测度取的.
\end{corollary}

我们只在维数 1 的情形完整写出了反演公式的证明, 尽管同样的证明只需对 Taylor 展开作微小的推广即可对一般 $n$ 成立.一般情形也可以从 $n = 1$ 的情形如下推出:反演公式对张量积 $f_1(x_1)\cdots f_n(x_n)$ 及其线性组合成立, delta 函数是张量积 $\delta_0(x) = \delta(x_1)\cdots\delta(x_n)$, 而任意函数都可以写成 delta 函数的线性组合 $f(x) = \int f(t)\delta(x-t)\, dt$.

\textbf{习题 2.5.} 利用近似 delta 函数, 把上述论证变成 $\mathbb{R}^n$ 上反演公式的一个严格的、自足的证明.

探讨上述经由\lemref{lem:c5-2.3} 的反演公式证明与该公式的其他证明之间的关系是值得的.正如有限维向量空间之间的线性算子可以通过矩阵表示来研究一样, 我们可以借助其核 $K$——即 $\mathbb{R}^n\times\mathbb{R}^n$ 上的分布, 使得

\[
T\phi(x) = \int\phi(x')K(x, x')\, dx'
\]

——来研究算子 $T$.在断言 $T\phi(x_0)$ 只依赖于 $\phi(x_0)$ 时, 我们已经证明了作用 $T$ 与乘以某个函数相同;用核的语言来说, 我们已经确立了

\[
T\phi(x) = \int\phi(x')f(x)\delta(x-x')\, dx'
\]

为了使 $T$ 与微分交换——这与和平移交换并无太大区别——我们得出结论:$f(x)$ 是常数.

但我们还可以直接看到, 反演公式的一个等价表述是如下分布论恒等式

\begin{equation}
\int_{\widehat{\mathbb{R}}^n} e^{i(x-x')\cdot\xi}\frac{d\xi}{(2\pi)^n} = \delta(x-x')\label{eq:c5-162}
\end{equation}

它实际上就是反演公式对 $\delta$ 函数的特例.把左边的积分看作内积, 上述恒等式可以理解为:平面波 $\xi\mapsto e^{i\xi x}$ 在 $x$ 变化的意义下是某种"正交归一"的.因此, 在写出

\[
\int f(t)\delta(x-t)\, dt = f(x) = \int\hat{f}(\xi)e^{i\xi x}\frac{d\xi}{2\pi}
\]

时, 我们可以把 傅里叶变换看作类似于从物理空间中的 delta 函数到平面波的正交归一基的更换, 从而 Plancherel 与 Parseval 恒等式应当立即成立.例如, 在一维情形, 不同的平面波是自伴算子 $i\frac{d}{d\xi}$ 的具有不同特征值的特征函数, 因此按一般原理应当正交——这一论证可以严格化, 以证明上述分布在 $x-x'\neq 0$ 处消失, 正如前一个证明中通过 Taylor 展开与在多项式环上的线性所做的那样.让我们提一下确立这个恒等式从而证明反演公式的其他几种途径.\footnote{如前所述, 平面波可视为交换的自伴算子族 $f\mapsto i\frac{\partial}{\partial x_i}f$ 的特征函数.对偶地, delta 函数——它们按 $\int\delta_0(y-x)\bar{\delta}_0(y'-x)\, dx = \delta(y-y')$ 的意义同样"正交归一"——可视为交换的算子族 $f\mapsto x_if$ 的特征函数.}

只需证明

\[
\frac{1}{(2\pi)^n}\int e^{i\xi\cdot x}\, d\xi = \delta_0(x)
\]

作为变量 $x$ 上的分布、在 $\mathbb{R}^n$ 上成立即可——这种平移不变性对应于前一个证明中 $T$ 与 $\frac{d}{dx}$ 交换.把上述分布看作张量积, 只需考虑 $n = 1$ 的情形, 但我们不妨不这样做.回忆:每个支集在原点的分布都是 $\delta(x)$ 的导数的有限线性组合, 因此 $\delta$ 函数本身(不计常数)是唯一支集在 $0$ 处、$-n$ 次齐次的分布——这些事实由 Taylor 展开容易建立.由于左边的积分显然关于 $x$ 是 $-n$ 次齐次的, 只要我们能够证明

\begin{equation}
\int e^{i\xi\cdot x}\, d\xi\label{eq:c5-163}
\end{equation}

的支集在原点, 就在不计常数的意义下证明了该恒等式——在这个精确的意义上, 平面波 $\xi\mapsto e^{i\xi\cdot x}$ "平均地"为零.

启发式地, 让我们概述完成这一计算的几种途径.假装积分 \eqref{eq:c5-163} 是经典积分, 且 $x\neq 0$ 固定.若把平面波看作微分算子的特征函数, 我们可以利用恒等式

\[
\int e^{i\xi\cdot x}\, d\xi = \int\frac{1}{|x|^2}\Delta_\xi e^{i\xi\cdot x}\, d\xi,
\]

在 $\xi$ 中分部积分;或者我们也可以旋转到 $x = |x|(1, 0, \ldots, 0)$ 的情形, 利用恒等式 $\frac{1}{i|x|}\frac{\partial e^{i\xi\cdot x}}{\partial\xi_1}$ 分部积分.如果我们更愿意把平面波看作平移算子的特征函数, 则可以通过在 $\xi$ 变量中平移来证明该积分在 $x\neq 0$ 时为零:

\[
\int e^{i\xi\cdot x}\, d\xi = \int e^{i(\xi-\xi')\cdot x}\, d\xi = e^{-i\xi'\cdot x}\int e^{i\xi\cdot x}\, d\xi
\]

其中 $\xi'$ 是依赖于 $x$ 的某个适当的频率.当然, 这些都是把积分当作经典积分来处理的启发式思路, 我们不能把 $x$ 当作固定点, 但我们可以通过固定一个局部化在 $x$ 附近的检验函数, 给出一个与下面类似的完整论证, 从而使这些论证严格化.

\begin{proof}
设 $\psi\in\mathcal{C}_0^\infty(\mathbb{R}^n)$ 的支集远离原点.设 $\phi$ 是光滑、快速递减的函数, 满足 $\phi(0) = 1$.

\[
\begin{aligned}
\int\int\psi(x)e^{i\xi\cdot x}\, dx\, d\xi &= \lim_{\delta\to 0^+}\int\Big(\int\psi(x)e^{i\xi\cdot x}\, dx\Big)\phi(\delta\xi)\, d\xi \\
&= \lim_{\delta\to 0^+}\int\psi(\delta x)\Big(\int e^{i\xi\cdot x}\phi(\xi)\, d\xi\Big)\, dx
\end{aligned}
\]

由控制收敛定理, 该极限为零($d\xi$ 积分是 $x$ 的快速递减函数, 且 $\psi(\delta x)\to 0$).
\end{proof}

若不对 $\psi$ 的支集作任何假设, 则只要选取一个其 傅里叶变换已知的 $\phi$(例如 Gauss 函数), 上述证明就直接建立了带常数的反演公式.事实上, 如果我们知道 Gauss 函数的反演公式, 则反演公式对 Gauss 函数的伸缩与平移也成立.作为极限情形, 反演公式对任何 $\delta$ 函数成立, 从而由分解 $f(x) = \int f(t)\delta(x-t)\, dt$, 反演公式对任意函数成立.

\textbf{习题 2.6.} 从 Gauss 函数的情形出发, 给出反演公式的一个自足的、直接的证明.

在 $n = 1$ 的情形, 还有一种更偏向复分析的方法来求分布论积分 $\int_{-\infty}^{\infty} e^{i\xi x}\, d\xi$, 它不仅确定了常数\footnote{顺便指出, 这一论证也间接地确定了 Gauss 积分的值.}, 而且直接把反演公式中的 $2\pi$ 与 Cauchy 积分公式中的 $2\pi$(圆的周长)联系起来.基本做法是:利用平面波的复解析延拓以及 $\frac{1}{x+i0}$ 的公式 \eqref{eq:c3-129}, 把

\[
\begin{aligned}
\int_{-\infty}^{\infty} e^{i\xi x}\, d\xi &= \int_{-\infty}^{0} e^{i\xi x}\, d\xi + \int_{0}^{\infty} e^{i\xi x}\, d\xi \\
&= \lim_{y\to 0^+}\int_{-\infty}^{0} e^{i\xi(x-iy)}\, d\xi + \int_{0}^{\infty} e^{i\xi(x+iy)}\, d\xi \\
&= i\left(\frac{-1}{x-i0} + \frac{1}{x+i0}\right) \\
&= 2\pi\delta(x)
\end{aligned}
\]

分解.


\section{傅里叶变换的推广}

作为 Parseval 与 Plancherel 公式的推论, 我们可以通过一个简单的稠密性论证把 傅里叶变换的定义延拓到 $L^2(\mathbb{R}^n)$ 函数.事实上, 对任何 $u\in L^2$, 我们可以选取 $\mathcal{S}(\mathbb{R}^n)\subset L^1$ 中的函数序列 $u_j$ 在 $L^2$ 范数下收敛于 $u$.由 Plancherel, $\|\mathcal{F}(u_j)-\mathcal{F}(u_k)\|_{L^2}\lesssim\|u_j-u_k\|_{L^2}$.因此序列 $\mathcal{F}(u_j)$ 在 $L^2$ 中构成 Cauchy 列, 从而收敛于一个极限, 我们可以称之为 $\hat{u}$.显然, 这个定义不依赖于特定的序列.此外, 容易验证 Parseval 恒等式可以延拓到一切 $L^2$ 函数.由于 $\mathcal{F}$ 因此是到其像上的等距, 其像必是闭的;但该像也包含 Schwartz 函数, 所以它必然是整个 $L^2$.于是我们证明了:

\begin{theorem}\label{thm:c5-3.1}
傅里叶变换是 Hilbert 空间 $L^2(\mathbb{R}^n)$ 与 $L^2(\widehat{\mathbb{R}}^n)$ 之间的等距.
\end{theorem}

我们还可以把 傅里叶变换进一步延拓到定义在 $\mathbb{R}^n$ 上的一类特殊分布.

\begin{definition*}
我们称 Schwartz 空间的对偶空间中的元素为 \newterm{tempered-distribution}{缓增分布}{tempered distribution}.
\end{definition*}

\begin{example*}
虽然 $e^x$ 不是 $\mathbb{R}$ 上的缓增分布, 因为它增长得太快, 但函数 $e^x\cos(e^x) = \frac{d}{dx}\sin(e^x)$ 是.这里我们作函数与分布的通常等同.
\end{example*}

注意, 缓增分布连续嵌入到前面定义的普通分布空间中.与普通分布的性质类比, 对每个缓增分布 $u$, 存在一个自然数 $N$ 与常数 $C = C_{\alpha, \beta}$, 使得

\[
|\langle u, \phi\rangle| \leq C\sum_{|\alpha|, |\beta|\leq N}\sup|x^\alpha\partial^\beta\phi|, \quad \phi\in\mathcal{S}(\mathbb{R}^n)
\]

我们现在可以定义缓增分布的 傅里叶变换, 即

\[
\langle\hat{u}, \phi\rangle = \langle u, \hat{\phi}\rangle.
\]

容易验证, 对每个缓增分布 $u$, 这定义了一个缓增分布 $\hat{u}$.此外, 所有已在 $\mathcal{S}(\mathbb{R}^n)$ 中对 Schwartz 函数验证过的 傅里叶变换性质都可以容易地延拓到一切缓增分布.特别地, 由于一切 $L^p$ 空间都包含在 $\mathcal{S}'(\mathbb{R}^n)$ 中, 我们对所有这些空间都有了 傅里叶变换的定义.注意, 在 $L^1$ 的情形, 这个定义与 \eqref{eq:c5-155} 中给出的定义一致.

Dirac 测度 $\delta_0$ 的 傅里叶变换的如下简单而非常有用的公式现在有了意义:

\begin{equation}
\mathcal{F}(\delta_0) = 1, \qquad \mathcal{F}(1) = (2\pi)^n\delta_0\label{eq:c5-164}
\end{equation}

另外注意, 若我们用 $\mathrm{sign}(x)$ 表示由局部可积函数 $\frac{x}{|x|}$ 给出的一维缓增分布, 则我们有

\begin{equation}
\widehat{\mathrm{sign}}(\xi) = -2i\, \mathrm{pv}\Big(\frac{1}{\xi}\Big)\label{eq:c5-165}
\end{equation}

事实上, $\mathrm{sign}'(x) = 2\delta_0$.因此, $i\xi\, \widehat{\mathrm{sign}}(\xi) = 2$.于是, 对任何快速递减的 $\phi$, 我们有

\[
i\int\mathrm{sign}(x)\, \widehat{x\phi}(x)\, dx = 2\hat{\phi}(0) = 2\int\phi(x)\, dx
\]

又注意, $\widehat{\mathrm{sign}}(x)$ 是奇分布, 因此每当 $\phi(x) = \phi(-x)$ 是偶检验函数时, $\langle\widehat{\mathrm{sign}}, \phi\rangle = 0$.现在给定一般检验函数 $\phi$, 写成 $\phi = \frac{1}{2}(\phi(x)+\phi(-x)) + \frac{1}{2}(\phi(x)-\phi(-x)) = \phi_{\mathrm{ev}} + \phi_{\mathrm{odd}}$.于是, 由前述, 我们推出

\[
\langle\widehat{\mathrm{sign}}, \phi\rangle = \Big\langle\widehat{\mathrm{sign}}, \xi\Big(\frac{1}{\xi}\phi_{\mathrm{odd}}\Big)\Big\rangle = -2i\Big\langle\mathrm{pv}\Big(\frac{1}{\xi}\Big), \phi\Big\rangle
\]

如所期望.

这一事实也可以通过求如下分布论积分

\[
\int\mathrm{sign}(x)e^{-i\xi x}\, dx
\]

更直接地观察到, 所用的思路与上一节概述的 Fourier 反演公式的复解析证明相同.

\textbf{习题 1.} 证明:是缓增分布的调和函数只有多项式.

\textbf{习题 2.} 设 $f(x) = e^{-|x|}\in L^1(\mathbb{R})$.利用 $f$ 满足一个简单的二阶微分方程这一事实, 计算 $\hat{f}(\xi)$(从而 $\hat{f}(0) = \int f(x)\, dx = 2$).讨论 $f$ 与 $\hat{f}$ 的精确的正则性与衰减量, 以及如何从物理空间表示中预见它们.注意, $\hat{f}$ 可亚纯延拓到复平面——通过考虑与复平面波 $x\to e^{izx}$, $z\in\mathbb{C}$ 的相关性, 你可以从 $f$ 在物理空间中的形式预见极点的位置.

\textbf{习题 3.} 假设 $u$ 是一个缓增分布, 在 $\mathbb{R}^n$ 的子群 $S$ 的平移下不变——例如 $u$ 可以是周期的, 或是少于 $n$ 个变量的函数.为什么我们可以假设 $S$ 是闭的?证明 傅里叶变换 $\hat{u}$ 的支集包含在 $S$ 下不变的平面波的 \newterm{annihilator-subgroup}{零化子群}{annihilator subgroup} $S^\perp$ 中:

\[
S^\perp = \{\xi \mid e^{i\xi\cdot x} = 1 \text{ 对一切 } x\in S\}
\]


\section{不确定性原理与局部化}

在实直线上设算子 $X, D$ 定义为

\[
Xf(t) = tf(t), \qquad Df(t) = -if'(t)
\]

注意,

\[
[D, X]f = DXf - XDf = -if
\]

这种不可交换性是如下结论的原因:

\begin{proposition}[Heisenberg 不确定性原理]\label{prop:c5-4.1}
如下不等式成立:

\[
\|Xf\|_{L^2}\cdot\|Df\|_{L^2}\geq\frac{1}{2}\|f\|_{L^2}^2
\]
\end{proposition}

\begin{proof}
注意, 利用上述对易关系,

\[
0\leq\|(aX+ibD)f\|_{L^2}^2 = a^2\|Xf\|_{L^2}^2 + b^2\|Df\|_{L^2}^2 - ab\|f\|_{L^2}^2
\]

现在, 选取 $a = \|Df\|_{L^2}$, $b = \|Xf\|_{L^2}$, 使右边最小化.
\end{proof}

不确定性原理非正式地描述为 $\Delta x\cdot\Delta\xi\geq 1/2$\footnote{这一记号来自量子力学:在那里归一化 $\|f\|_{L^2} = 1$, 并把 $|f(x)|^2$ 与 $\frac{|\hat{f}(\xi)|^2}{2\pi}$ 分别解释为位置空间与动量空间中的概率密度.在这种设定下, Heisenberg 不确定性原理给出位置与动量的标准差之积的一个下界.}, 它对我们在空间和频率中同时局部化一个函数(或任何其他相关对象)的精度设置了一个界限.让我们更详细地研究这些局部化.

\mtsubsec{物理空间的局部化}

如果我们想把函数 $f$ 局部化到区域 $D\subset\mathbb{R}^n$, 可以简单地用特征函数 $\chi_D$ 乘 $f$.这种局部化的问题是, 所得函数 $\chi_Df$ 即使 $f$ 光滑也不是光滑的.为纠正这一点, 我们选取 $\phi_D\in\mathcal{C}_0^\infty(D)$, 使得 $\phi_D$ 与 $\chi_D$ 相差不大.在 $D$ 是以 $x_0$ 为中心的球 $B(x_0, R)$ 的特殊情形, 我们可以选取 $\phi_D$ 在球 $B(x_0, R)$ 上等于 1, 在球 $B(x_0, 2R)$ 外等于 0.这导致如下对 $\phi_D$ 的导数的界

\[
|\partial^\alpha\phi_D|\lesssim R^{-|\alpha|}.
\]

一般地, 给定一个可以关联一个长度尺度 $R$(例如其直径, 或到其内部一个固定点的距离)的区域 $D$, 我们可以找到函数 $\phi_D\in\mathcal{C}_0^\infty(D)$, 使得对一切多重指标 $\alpha\in\mathbb{N}^n$,

\begin{equation}
|\partial^\alpha\phi_D|\lesssim R^{-|\alpha|}, \label{eq:c5-166}
\end{equation}

一个一般性的注记:形如 \eqref{eq:c5-166} 的导数估计在分析中非常常见, 而且几乎总是出现于满足该估计的函数来自另一个导数有界的函数的伸缩版本时.这就是为什么出现的特定指数与\newterm{dimensional-analysis}{量纲分析}{dimensional analysis}一致.

\mtsubsec{频率空间的局部化}

和前面一样, 我们可以通过 $\mathcal{F}^{-1}(\chi_D\hat{f})$ 把函数 $f$ 在频率空间中局部化到区域 $D\subset\mathbb{R}^n$.同样, 使用更光滑的截断版本常常是有益的, 因此我们设

\[
\widehat{P_Df}(\xi) = \phi_D\hat{f}(\xi).
\]

$P_D$ 是 \newterm{fourier-multiplier-operator}{Fourier 乘子算子}{Fourier multiplier operator}的一个例子, 即如下类型的算子:

\begin{equation}
\widehat{T_mf}(\xi) = m(\xi)\hat{f}(\xi).\label{eq:c5-167}
\end{equation}

其中 $m = m(\xi)$ 是给定的函数, 称为该算子的符号.显然,

\begin{equation}
T_mf(x) = f*K(x) = \int f(x-y)K(y)\, dy\label{eq:c5-168}
\end{equation}

其中 $K$, 即 $T$ 的核, 是 $m$ 的逆 傅里叶变换:

\[
K(x) = (2\pi)^{-n}\int e^{ix\cdot\xi}m(\xi)\, d\xi
\]

显然, 任何线性微分算子 $P(\partial)$ 都是以 $P(i\xi)$ 为符号的乘子.

为了比较粗糙与光滑截断算子在物理空间中的作用, 只需考察相应的核 $K$.设 $I = [-1, 1]\subset\mathbb{R}$, $\chi_I$ 为粗糙截断(忽略 $2\pi$ 常数).相应的核

\[
K(x) = \int_{-1}^{1} e^{ix\cdot\xi}\, d\xi = 2\frac{\sin x}{x}
\]

当 $|x|\to\infty$ 时衰减得非常慢.正因如此, 算子

\[
\mathcal{F}^{-1}(\chi_I\hat{f})(x) = 2\int\frac{\sin(x-y)}{(x-y)}f(y)\, dy
\]

具有非常差的局部化性质.事实上, 该算子把任何支集在某个集合 $J\subset\mathbb{R}$ 中的函数扩散到整个 $\mathbb{R}$.这种情况对应于频率空间中的完美局部化与物理空间中的极差局部化.当我们对物理空间中的粗糙截断局部化 $\chi_If$ 时, 恰好出现相反的情形.另一方面, 当我们在频率空间中使用光滑截断 $\phi_I$ 时, 频率截断算子 $P_If = \mathcal{F}^{-1}(\phi_I\hat{f})$ 具有形式 $f\to K*f$, 其中核

\[
K(x) = \int_{\mathbb{R}} e^{ix\cdot\xi}\phi_I(\xi)\, d\xi
\]

是快速递减的.在这种情况下, 我们可以证明:

\begin{lemma}\label{lem:c5-4.4}
设 $I = [-1, 1]$, $\phi_I$ 是 $I$ 上的光滑截断, $P_If = \mathcal{F}^{-1}(\phi_I\hat{f})$.则, 若 $f$ 是任何支集在集合 $D\subset\mathbb{R}$ 上的 $L^2$ 函数,

\[
|P_I(f)(x)|\lesssim C_j\|f\|_{L^2}\bigl(1+\mathrm{dist}(x, D)\bigr)^{-j}
\]

对一切 $j\in\mathbb{N}$ 成立.
\end{lemma}

因此, $P_I$ 把任何函数 $f$ 的支集扩散一个 $O(1)$ 的距离, 再加上一个快速递减的尾部.

\textbf{习题.} 证明:不存在非平凡函数 $\phi$, 使得 $\phi$ 与 $\mathcal{F}(\phi)$ 都具有紧支集.

上述讨论可以容易地推广到高维情形.特别地, 我们可以得到 $\mathbb{R}^n$ 中 Fourier 支集限制在中心在原点的球 $B_R = B(0, R)$ 内的函数的定性描述.设 $\phi_R$ 是 $B_R$ 的光滑截断.更精确地说, 我们取它形如

\[
\phi_R(\xi) = \phi(\xi/R)
\]

其中 $\phi$ 是 $B_1$ 的光滑截断, 即 $\phi$ 光滑, 在 $B_1$ 上恒等于 1, 支集在(比如说)$B_2$ 中.容易验证, 对任何多重指标 $\alpha$,

\[
\sup_\xi|\partial_\xi^\alpha\phi_R(\xi)|\leq c_\phi R^{-|\alpha|},
\]

其中常数 $c_\phi$ 只依赖于固定的 $\phi$ 及其导数.

若 $f$ 的 Fourier 支集限制在 $B_R$ 内, 则 $\hat{f} = \phi_R\hat{f}$.因此,

\begin{equation}
f(x) = \int_{\mathbb{R}^n} f(y)K_R(x-y)\, dy\label{eq:c5-169}
\end{equation}

其中 $K_R(x) = \mathcal{F}^{-1}(\phi_R)$.

\begin{lemma}\label{lem:c5-4.5}
核 $K_R(x)$ 满足如下估计:

\begin{equation}
|\partial_x^\alpha K_R(x)| \leq C_{N, \alpha}R^{|\alpha|}R^n(1+|x|R)^{-N}, \label{eq:c5-170}
\end{equation}

对一切 $R>0$、任何 $N\in\mathbb{N}$ 与多重指标 $\alpha\in\mathbb{N}^n$ 成立, 其中常数 $C_{N, \alpha}$ 只依赖于 $N$、$\alpha$、维数 $n$ 以及固定检验函数 $\phi$ 的选取.
\end{lemma}

\begin{proof}
事实上, 分部积分,

\[
\begin{aligned}
K_R(x) &= \int_{\mathbb{R}^n} e^{ix\cdot\xi}\phi_R(\xi)\, d\xi = \int_{\mathbb{R}^n}\Big(\frac{-1}{ix}\Big)^\alpha\partial_\xi^\alpha(e^{ix\cdot\xi})\phi_R(\xi)\, d\xi \\
&= \int_{\mathbb{R}^n}\Big(\frac{1}{ix}\Big)^\alpha e^{ix\cdot\xi}\partial_\xi^\alpha\phi_R(\xi)\, d\xi
\end{aligned}
\]

因此, 对任何 $\alpha$, $|\alpha| = N$, 记 $|B_R| = c_nR^n$ 为 $B_R$ 的体积,

\[
|x^\alpha K_R(x)| \leq \int_{\mathbb{R}^n}|\partial_\xi^\alpha\phi_R(\xi)| \leq c_\phi R^{-N}|B_R| \leq c_nc_\phi R^{-N+n}
\]

于是, 对某个依赖于 $N$、$n$ 与固定 $\phi$ 的常数 $C_N$, $|K_R(x)|\leq C_NR^n(|x|R)^{-N}$.另一方面, 当 $|x|\leq R^{-1}$ 时, $|K_R(x)|\lesssim R^n$.因此, 对每个 $N\in\mathbb{N}$,

\[
|K_R(x)|\lesssim C_NR^n(1+|x|R)^{-N}.
\]

还容易验证, $K_R$ 的每个导数都花费一个 $R$ 的因子, 从而证明了 \eqref{eq:c5-170}.
\end{proof}

现在回到 \eqref{eq:c5-169}, 我们有

\[
\begin{aligned}
|\partial^\alpha f(x)| = \Big|\int_{\mathbb{R}^n} f(y)\partial^\alpha K_R(x-y)\, dy\Big| &\lesssim R^{|\alpha|+n}\int_{\mathbb{R}^n}|f(y)|(1+R|x-y|)^{-N}\, dy \\
&\lesssim R^{|\alpha|+n}\|f\|_{L^1}
\end{aligned}
\]

又, 由 Hölder 不等式, 其中 $\frac{1}{p}+\frac{1}{p'} = 1$,

\[
\begin{aligned}
|\partial^\alpha f(x)| &\lesssim \|f\|_{L^p}\|\partial^\alpha K_R\|_{L^{p'}}\lesssim R^{|\alpha|}R^nR^{-n/p'}\|f\|_{L^p} \\
&\lesssim R^{|\alpha|+n/p}\|f\|_{L^p}
\end{aligned}
\]

我们刚刚证明了非常重要的 \newterm{bernstein-inequality}{Bernstein 不等式}{Bernstein inequality}的如下版本($L^p-L^\infty$ 版本):

\begin{proposition}\label{prop:c5-4.6}
假设 $f$ 是 $L^p$ 函数, 其 傅里叶变换的支集在球 $B_R = B(0, R)$ 内.则 $f$ 具有无穷多个在 $L^\infty$ 中有界的导数, 并且有

\[
\|\partial^\alpha f\|_{L^\infty(\mathbb{R}^n)}\lesssim R^{n/p+|\alpha|}\|f\|_{L^p}.
\]
\end{proposition}

\begin{remark*}
注意, 该命题可以通过化归到 $R = 1$ 的特殊情形来证明.更精确地说, 假设结论对 $R = 1$ 成立, 并考虑 傅里叶变换支集在 $B_R$ 内的函数 $f$.设 $g(x) = R^{-n}f(R^{-1}x)$, 注意, $\operatorname{supp}\hat{g}(\xi) = \operatorname{supp}\hat{f}(R\xi)\subset B_1$, 因此我们有 $\|\partial^\alpha g\|_{L^\infty(\mathbb{R}^n)}\lesssim\|g\|_{L^1} = R^{-n}R^{n/p}\|f\|_{L^p}$.于是, $\|\partial^\alpha f\|_{L^\infty(\mathbb{R}^n)}\lesssim R^{n/p+|\alpha|}\|f\|_{L^p}$.
\end{remark*}

正如我们在这些讲义中将非常频繁地看到的那样, 量纲分析可以用来迅速弄清上面出现的指数.例如, 如果我们把空间变量看作具有尺度 $L$, 从而体积元 $dx$ 具有尺度 $L^n$, 那么频率变量具有单位 $R\sim L^{-1}$.我们看到 $\|\partial^\alpha f\|_{L^\infty}$ 具有尺度 $L^{-|\alpha|}$, 而 $\|f\|_{L^p} = \big(\int|f(x)|^p\, dx\big)^{1/p}$ 具有尺度 $L^{n/p}$——因此出现的 $R$ 的幂使得该估计与量纲分析一致.

定性地说, Bernstein 估计体现了关于 $L^p$ 范数与 傅里叶变换的一些基本直觉.在某种意义上, 函数的较高 $L^p$ 范数(如 $L^\infty$ 范数)控制函数在空间的局部区域中爆破的能力, 而较低的 $L^p$ 范数控制其在无穷远处的增长.Bernstein 估计说的是:如果一个频率局部化的函数在无穷远处不增长(即具有有界的 $L^p$ 范数, $1\leq p<\infty$), 那么由于没有高频分量, 该函数也被禁止具有局部奇性(即 $L^\infty$ 范数有界).

我们对这类估计感兴趣的一个主要原因, 并不是我们经常遇到 傅里叶变换紧支撑的函数, 而是我们经常把更一般的函数分解为频率局部化的各部分之和.我们将在 Littlewood-Paley 理论的研究中回到这个想法.


\section{在 PDE 中的应用}

考虑我们的基本 PDE 在 $\mathbb{R}\times\mathbb{R}^n$ 中的初值问题, 写成如下形式:

\begin{equation}
\partial_t\phi = \Delta\phi, \qquad \phi(0, x) = f(x)\label{eq:c5-171}
\end{equation}

\begin{equation}
\partial_t\phi = i\Delta\phi, \quad \phi(0, x) = f(x)\label{eq:c5-172}
\end{equation}

\begin{equation}
\partial_t^2\phi = \Delta\phi, \qquad \phi(0, x) = f(x), \quad \partial_t\phi(0, x) = g(x)\label{eq:c5-173}
\end{equation}

\begin{equation}
\partial_t^2\phi = -\Delta\phi, \quad \phi(0, x) = f(x), \quad \partial_t\phi(0, x) = g(x)\label{eq:c5-174}
\end{equation}

在每一种情形下, 我们都可以用 傅里叶变换方法写出解.更精确地说, 我们可以对每个方程取 傅里叶变换, 设

\[
\hat{\phi}(t, \xi) = \int e^{-ix\cdot\xi}\phi(t, x)\, dx
\]

并求解关于 $t$ 的所得微分方程.完成这一步之后, 我们只需用逆 傅里叶变换即可得到我们的解, 即

\[
\phi(t, x) = \int e^{ix\cdot\xi}\hat{\phi}(t, \xi)\frac{d\xi}{(2\pi)^n}
\]

在热方程 \eqref{eq:c5-171} 的情形, 我们导出

\begin{equation}
\phi(t, x) = (2\pi)^{-1}\int_{-\infty}^{+\infty} e^{ix\xi}e^{-t|\xi|^2}\hat{u}_0(\xi)\frac{d\xi}{(2\pi)^n}\label{eq:c5-175}
\end{equation}

而在 Schrödinger 方程的情形,

\begin{equation}
\phi(t, x) = (2\pi)^{-1}\int_{-\infty}^{+\infty} e^{ix\xi}e^{-it|\xi|^2}\hat{u}_0(\xi)\frac{d\xi}{(2\pi)^n}\label{eq:c5-176}
\end{equation}

\textbf{习题 1.} 说明如何把公式 \eqref{eq:c5-175} 与 \eqref{eq:c5-176} 与物理空间\eqref{eq:c4-153} 和 (154) 联系起来.

在波动方程 \eqref{eq:c5-173} 的特殊情形, 我们导出

\begin{equation}
\phi(t, x) = \int_{\mathbb{R}^n} e^{ix\cdot\xi}\Big(\cos t|\xi|\hat{f}(\xi) + \frac{\sin t|\xi|}{|\xi|}\hat{g}(\xi)\Big)\frac{d\xi}{(2\pi)^n}\label{eq:c5-177}
\end{equation}

\textbf{习题 2.} 对 Laplace 方程 \eqref{eq:c5-174} 导出与 \eqref{eq:c5-177} 类似的公式.利用这些公式证明:\eqref{eq:c5-173} 对一切 $f, g\in\mathcal{S}(\mathbb{R}^n)$ 都有解, 而 \eqref{eq:c5-174} 没有.不过要证明:若我们只指定 $\phi(0, x) = f$(这就是 $\mathbb{R}^{n+1}$ 中 Laplace 算子 $\partial_t^2+\Delta$ 的 Dirichlet 问题), 则该问题对一切函数 $f\in\mathcal{S}(\mathbb{R}^n)$ 都有唯一解 $\phi$, 且当 $|t|+|x|\to\infty$ 时它衰减为零.

\textbf{习题 3.} 证明:在维数 $1+3$ 的特殊情形, 如何从公式 \eqref{eq:c5-177} 过渡到 Kirchhoff \eqref{eq:c4-146}:

\begin{equation}
\phi(t, x) = \partial_t\Big((4\pi t)^{-1}\int_{|x-y|=t} f(y)\, da(y)\Big) + (4\pi t)^{-1}\int_{|x-y|=t} g(y)\, da(y)\label{eq:c5-178}
\end{equation}

它与前一章基于基本解的显式计算导出的公式一致.

比较基于 Fourier 的公式 \eqref{eq:c5-177} 与 Kirchhoff 公式 \eqref{eq:c5-178} 是很有趣的.注意, 利用 Parseval, 从 \eqref{eq:c5-177} 导出全局能量恒等式是相当容易的:

\[
\int_{\mathbb{R}^n}\big(|\partial_t\phi|^2 + |\nabla\phi|^2\big) = \int_{\mathbb{R}^n}\big(|\nabla f|^2 + |g|^2\big)\, dx
\]

而从 \eqref{eq:c5-178} 得到这样的恒等式似乎一点也不显然, 事实上相当不可信.另一方面, \eqref{eq:c5-178} 非常适合给出影响域信息.事实上, 我们立即从公式中读出:若资料 $f, g$ 支集在球 $B_a = \{|x-x_0|\leq a\}$ 内, 则对任何时刻 $t$, $\phi(t, x)$ 的支集在球 $B_{a+|t|}$ 内.另一方面, 这一事实在基于 Fourier 的公式\footnote{支集信息可由 Paley-Wiener 型结果导出.}\eqref{eq:c5-177}\footnote{然而, 支集信息可以借助 Paley-Wiener 方法从 \eqref{eq:c5-177} 中提取出来, 参见 [?] 第一卷.}中一点也不透明.解的不同表示具有不同甚至相反的优点与缺点, 这一事实对构造更复杂的、线性变系数或非线性波动方程的参数层(即近似解)具有重要意义.有两种可能的构造类型:物理空间中的构造, 模仿物理空间公式 \eqref{eq:c5-178};或 Fourier 空间中的构造, 模仿公式 \eqref{eq:c5-177}.前者称为 Kirchhoff-Sobolev 参数层或 Hadamard 参数层, 后者称为 Lax 参数层, 或更一般地, \newterm{fourier-integral-operator}{Fourier 积分算子}{Fourier integral operator}.

% ==== part06a1 第 6 章 基本泛函不等式(第 1--2 节) ====
% 源: MinerU_latex_lecture nots.tex 第 7789--9184 行;pdf_text.txt 第 6456--7557 行(PDF 第 121--140 页)
%
% 术语(新增,键=中文=英文):
%   interpolation=插值、multiplier=乘子、distribution-function=分布函数、simple-function=简单函数、
%   analytic-family-of-operators=解析算子族、sublinear-operator=次线性算子、weak-lp=弱 L^p、
%   maximal-function=极大函数、covering-lemma=覆盖引理、regularity=正则性、
%   dyadic-decomposition=二进分解、sobolev-space=Sobolev 空间、embedding-theorem=嵌入定理、
%   holder-space=Hölder 空间、trace-operator=迹算子、trace-theorem=迹定理、
%   characteristic-function=特征函数(键复用 part03)、extension=延拓
% 注: 傅里叶变换、Plancherel/Parseval 恒等式、Hölder 不等式、Chebyschev 不等式、Minkowski 不等式、
%   Banach/Hilbert 空间、Lebesgue 测度、Borel 测度、卷积、特征函数、截断函数、局部可积、检验函数、
%   Lipschitz 函数、缓增分布等经典术语按普通中文书写(可能由其他部分首次引入,待协调).
%
% REFS(需跨章/待解析项):
%   - 源 7822 行 "the integral representation (168)" → \eqref{eq:c5-168},公式 (168) 位于第 6 章引言
%     (本部分之前的部分),标签需由前一部分定义.
%   - 源 8623 行 "For this we refer to (***)" → 按原文保留 "(***)"(原书占位引用).
%   - 源 8745 行 "Following Nirenberg, [?]" → 正文写"参照 Nirenberg 的著作[?]"([?] 为原书破损文献号).
%   - 源 8867、8959 行 "Evans ([E], ch. 5)"、"see [E], section 5.6.2" → 文献 [E] = Evans 著作,照原样保留.
%   - 源 8704 行末尾的 "27" 为 OCR 噪音(pdf_text 对应行无此标记,原书该处无脚注),已删除.
%   - 脚注: 本部分四处原书脚注均按 pdf_text.txt 译入 \footnote:
%     ① 源 7859 行 "measure μ^1"(PDF 页 122);② 源 8124 行 "form^2"(页 125);
%     ③ 源 8375 行 "proved below^3"(页 128,Vitali 覆盖引理);④ 源 8836 行 "weak^4"(页 135).
%   - 源 7811 行 "Young-Hausdorf inequality" 按 pdf_text 译 "Young-Hausdorff 不等式"(原书拼写).
%
% 编号说明:
%   - 第 1 节: 1.1、1.2、1.3、1.5、1.9、1.11 为带编号子标题(\mtsubsec),自然推进共享计数器,
%     故 Lemma 1.4、Definition 1.6、Remark 1.7、Theorem 1.8、Theorem 1.10、Definition 1.12、
%     Theorem 1.13 编号与原书一致;三处未编号 Remark 用 remark* 环境,未编号习题用 \textbf{习题.}.
%   - 第 2 节: 2.1、2.6、2.8、2.10、2.12、2.15、2.17--2.19(原书 "Exercise 2.17/2.18/2.19",
%     处于共享序列中,以 \mtsubsec{习题} 推进计数器)、2.21、2.23、2.26 为子标题,其余编号对象
%     (2.2--2.5、2.7、2.9、2.11、2.13、2.14、2.16、2.20、2.22、2.24、2.25、2.27)自动对齐原书.
%   - 公式 (179)--(213) 共 35 个,按原书顺序以 \begin{equation} 自动编号,正文引用用 \eqref{eq:c6-NNN}.
% ====

\chapter{基本泛函不等式}

\section{基本插值理论}

\mtsubsec{引言}

把 傅里叶变换视为线性算子 $\mathcal{F}:\mathcal{S}(\mathbb{R}^n)\to\mathcal{S}(\widehat{\mathbb{R}}^n)$.根据 Plancherel 恒等式, 有 $\|\mathcal{F}(f)\|_{L^2}\le\|f\|_{L^2}$.另一方面, 有 $\|\mathcal{F}(f)\|_{L^\infty}\le\|f\|_{L^1}$.我们能否得到其他类型的界 $\|\mathcal{F}(f)\|_{L^q}\lesssim\|f\|_{L^p}$?事实证明, 这类估计可以通过在上述两个估计之间插值而轻易建立.复\newterm{interpolation}{插值}{interpolation}使我们能够对满足 $p^{-1}+q^{-1}=1$ 且 $q\ge 2$ 的任意 $p$、$q$ 值推出 $L^p$ 到 $L^q$ 的估计.这就是所谓的 Young-Hausdorff 不等式.插值理论对形如

\[
\widehat{T_m f}(\xi) = m(\xi)\hat{f}(\xi)
\]

的线性\newterm{multiplier}{乘子}{multiplier}算子特别有用, 其中 $m$ 是有界乘子.根据 Parseval 恒等式, 很容易验证 $L^2$-$L^2$ 估计 $\|T_m f\|_{L^2}\lesssim\|f\|_{L^2}$.为获得更多的估计, 我们通常使用积分表示 \eqref{eq:c5-168}:$T_m f(x) = f*K(x) = \int f(x-y)K(y)\, dy$, 其中 $K$ 是 $m$ 的逆 傅里叶变换.例如, 若能建立 $K\in L^1$, 则容易推出 $\|T_m f\|_{L^1}\lesssim\|f\|_{L^1}$, 因为 $\|f*K\|_{L^1}\le\|f\|_{L^1}\cdot\|K\|_{L^1}$.于是我们对 $T_m$ 同时有 $L^1$-$L^1$ 与 $L^2$-$L^2$ 估计, 因而有理由猜测对一切 $1\le p\le 2$ 都可能存在 $L^p$-$L^p$ 估计.这样的估计确实成立, 它可由插值得到.另一方面, 若能建立 $K\in L^\infty$, 则 $\|f*K\|_{L^\infty}\lesssim\|f\|_{L^1}$, 从而我们可以通过插值证明与 Hausdorff-Young 不等式相同的 $L^p$-$L^q$ 估计.

\mtsubsec{$L^p$ 空间回顾}

给定可测子集 $\Omega\subset\mathbb{R}^n$, 空间 $L^p(\Omega)$, $1\le p<\infty$, 由所有 $L^p$ 范数有限的可测函数 $f:\Omega\to\mathbb{C}$ 组成,

\[
\|f\|_{L^p} = \left(\int_\Omega |f(x)|^p\, dx\right)^{1/p} < \infty.
\]

空间 $L^\infty(\Omega)$ 由所有几乎处处有界的可测函数组成, 即

\[
\|f\|_{L^\infty} = \inf\left\{\alpha : \int_\Omega (|f(x)| > \alpha)\, dx = 0\right\} = \operatorname{ess\, sup}|f| < \infty.
\]

对一切 $1\le p\le\infty$, 空间 $L^p(\Omega)$ 都是 Banach 空间.当我们把 Lebesgue 测度 $dx$ 换成一般的正测度 $\mu$\footnote{不过, 当整个空间不能写成有限 $\mu$-测度子集的可数并时, 会出现一些复杂情况.这在本讲义中与我们无关.}时, $L^p$ 空间的理论依然成立.下面的不等式称为 Hölder 不等式

\begin{equation}
\|fg\|_{L^p} \le \|f\|_{L^q}\|g\|_{L^r}, \label{eq:c6-179}
\end{equation}

只要 $1/p = 1/q + 1/r$.指数之间的这种关系是必要的, 它保证两边关于测度都是 $\frac{1}{p}$ 次齐次的.特别地, 对 $p=1$,

\[
\|fg\|_{L^1} \le \|f\|_{L^q}\|g\|_{L^{q'}},
\]

其中满足 $\frac{1}{q'} = 1 - \frac{1}{q}$ 的 $q'$ 是与 $q$ 对偶的指数.该不等式意味着我们可以把每个元素 $g\in L^{q'}$ 等同于 Banach 空间 $L^q$ 上的有界线性泛函 $f\mapsto\int f(x)g(x)\, dx$.对一切 $1\le q<\infty$, 空间 $L^{q'}(\Omega)$ 是 $L^q(\Omega)$ 的对偶空间, 其含义是上述等同是等距的(特别地, $L^q$ 上每个有界线性泛函都以这种方式由唯一的 $g$ 产生), 而 $L^\infty(\Omega)$ 的对偶包含 $L^1(\Omega)$, 但远大于后者.常常扮演 $L^\infty$ 角色的是在无穷远处消失的连续函数空间 $\mathcal{C}_0(\mathbb{R}^n)$(因为它们在 $L^\infty$ 范数下构成 $\mathcal{C}_0^\infty$ 的闭包), 其对偶空间是 $\mathbb{R}^n$ 上有限 Borel 测度的集合.

不同的 $L^p$ 范数度量函数大小的不同方面.较高 $L^p$ 范数(如 $\|f\|_{L^\infty}\lesssim 1$)的估计保证 $|f|$ 在局部不会过大, 而较低 $L^p$ 范数(如 $\|f\|_{L^1}\lesssim 1$)的估计则控制 $f$ 在无穷远处的行为.空间 $L^2(\Omega)$ 尤为重要, 因为它自对偶, 并且具有由内积

\[
\langle f, g\rangle = \int_\Omega f\, \bar{g}\, dx
\]

给出的 Hilbert 空间结构.

\textbf{习题.} 证明对一切 $1\le p<\infty$, $\mathcal{C}_0^\infty(\Omega)$ 在 $L^p(\Omega)$ 中稠密.

给定一个可测函数 $f$ 和一个正数 $\alpha$, 记 $\Lambda(f, \alpha)$ 为 $f$ 的\newterm{distribution-function}{分布函数}{distribution function}, 定义为

\[
\Lambda(f, \alpha) = |\{x\in\Omega : |f(x)| > \alpha\}|.
\]

对 $1\le p<\infty$, 我们有 Chebyschev 不等式

\begin{equation}
\Lambda(f, \alpha) \le \alpha^{-p}\|f\|_{L^p}^p, \label{eq:c6-180}
\end{equation}

它定量地表达了如下事实:$L^p$ 函数的上水平集具有有限测度.注意到(至少作为记忆手段)两边具有相同的量纲是有帮助的, 因为 $f$ 与 $\alpha$ 具有相同的量纲.

\begin{proof}
\[
\Lambda(f, \alpha) = \int(|f(x)| > \alpha)\, d\mu(x) \le \int\left(\frac{|f(x)|^p}{\alpha^p}\right)\cdot(|f(x)| > \alpha)\, d\mu(x) \le \alpha^{-p}\|f\|_{L^p}^p.
\]
\end{proof}

我们可以用分布函数来写出 $f$ 的 $L^p$ 范数.事实上, 积分 $\int|f|^p$ 就是由图形 $\{( \beta, x) : 0 < \beta < |f(x)|^p \}$ 围成的区域的测度, 因此

\begin{equation}
\int|f(x)|^p\, dx = \int\int_0^\infty (|f(x)|^p > \beta)\, d\beta\, dx = p\int_0^\infty \alpha^{p-1}\Lambda(f, \alpha)\, d\alpha, \label{eq:c6-181}
\end{equation}

其中最后一个积分由代换 $\beta = \alpha^p$ 得到.

如果可测函数 $f:\Omega\to\mathbb{C}$ 的值域只包含 $\mathbb{C}$ 中的有限个点, 则称其为\newterm{simple-function}{简单函数}{simple function}, 即 $f = \sum_{i=1}^N a_i\chi_{A_i}$, 其中 $a_i\in\mathbb{C}$, $A_i\subset\Omega$ 可测.在本节中, 我们用 $S(\Omega)$ 表示 $\Omega$ 中所有简单函数组成的集合.回忆:对一切 $1\le p\le\infty$, $S(\Omega)$ 在 $L^p(\Omega)$ 中稠密.证明通常涉及用特征函数 $(f(x)\in E_\alpha)$ 的线性组合逼近给定的 $f(x)$, 并让集合族 $\{E_\alpha\}$ 趋于 $\mathbb{C}$ 的一个精细而完备的分割.

\textbf{习题.} 设 $f(x, y)$ 是 $\Omega_1\times\Omega_2\subset\mathbb{R}^{n_1}\times\mathbb{R}^{n_2}$ 上的可测函数.证明下列形式的 Minkowski 不等式,

\[
\left\|\int_{\Omega_2} f(x, y)\, dy\right\|_{L_x^p(\Omega_1)} \le \int_{\Omega_2}\|f(x, y)\|_{L_x^p(\Omega_1)}\, dy
\]

对 $1\le p\le\infty$ 成立, 既用对偶方法, 也直接证明.

\mtsubsec{三线引理}

对作用在 $L^p$ 空间上的线性算子而言, 解析插值方法基于最大模定理在条形区域上的一个变体, 称为三线引理.考虑条形区域

\[
D = \{z\in\mathbb{C} : 0 < \operatorname{Re}(z) < 1\}.
\]

我们用 $\mathcal{A}_{BC}$ 表示 $D$ 的闭包上的有界连续函数中在 $D$ 上解析的函数所成的集合.

\begin{lemma}[三线引理]\label{lem:c6-1.4}
设 $f\in\mathcal{A}_{BC}$ 满足

\[
|f(0+ib)| \le M_0, \qquad |f(1+ib)| \le M_1,
\]

对一切 $b\in\mathbb{R}$ 成立.则对一切 $0<a<1$ 与 $b\in\mathbb{R}$, 有

\[
|f(a+ib)| \le M_0^{1-a}M_1^a.
\]
\end{lemma}

\begin{remark*}
回忆:当 $f$ 全纯(且非平凡)时, $\log|f(z)|$ 是次调和函数.两条带上的界保证 $\log|f(z)|$ 在 $D$ 的边界上取得小于调和函数 $\phi(a+bi) = (1-a)\log M_0 + a\log M_1$ 的值, 而三线引理的结论断言在区域 $D$ 内部有 $\log|f(a+bi)| \le \phi(a+bi)$;然而我们不能简单地应用弱最大原理, 因为它对无界区域并非完全成立(如 $D$ 上的例子 $z\mapsto e^{-ie^{i\pi z}}$ 所示).
\end{remark*}

\begin{proof}
首先, 如有必要, 把 $f$ 换成 $\frac{f(z)}{M_0^{1-z}M_1^z}$, 我们可以假设 $M_0 = M_1 = 1$, 只需证明在整个区域上 $|f(z)|\le 1$.当 $f$ 在 $|\Im z|\to\infty$ 时衰减到 0, 则可以简单地对 $D$ 的足够大的子集应用通常的最大模原理, 得到在整个区域上 $|f(z)|\le 1$.若情况并非如此, 则(因为我们已假设 $|f(z)|$ 在 $|\Im z|\to\infty$ 时不会显著增长)可以对逼近函数 $F_\varepsilon(z) = e^{\varepsilon(1-z)z}f(z)$(它在 $|\Im z|$ 大时确实衰减)应用同样的论证, 得到

\[
|f(z)| = \lim_\varepsilon |F_\varepsilon(z)| \le 1
\]

在整个 $D$ 上成立.
\end{proof}

\mtsubsec{Stein-Riesz-Thorin 插值}

\begin{definition}\label{def:c6-1.6}
我们称一族由 $z\in D$ 指标化的线性算子 $T_z$ 为\newterm{analytic-family-of-operators}{解析算子族}{analytic family of operators}, 如果

\begin{enumerate}
\item $T_z$ 把简单函数映为可测函数;
\item 对任意一对简单函数 $f, g\in S(\Omega)$, 映射 $z\mapsto\int g(x)T_z f(x)\, dx$ 属于 $\mathcal{A}_{BC}$.
\end{enumerate}
\end{definition}

\begin{remark}\label{rem:c6-1.7}
在前一定义中选择简单函数作为检验函数的原因是:它们易于操作, 并且在每个 $p\in[1, \infty)$ 的 $L^p$ 中构成稠密集合.
\end{remark}

\begin{theorem}\label{thm:c6-1.8}
设 $T_z$ 是解析算子族, 并假设存在正常数 $M_0$、$M_1$, 使得对每个 $b\in\mathbb{R}$, 有

\[
\|T_{ib}f\|_{L^{q_0}(d\mu)} \le M_0\|f\|_{L^{p_0}(d\nu)}, \qquad \|T_{1+ib}f\|_{L^{q_1}(d\mu)} \le M_1\|f\|_{L^{p_1}(d\nu)},
\]

其中 $1\le q_0, p_0, q_1, p_1\le\infty$.则对 $z = a+ib\in D$, $T_z$ 延拓为从 $L^p(d\nu)$ 到 $L^q(d\mu)$ 的有界算子, 并且

\[
\|T_z f\|_{L^q(d\mu)} \le M_0^{1-a}M_1^a\|f\|_{L^p(d\nu)},
\]

其中

\[
\frac{1}{p} = \frac{1-a}{p_0} + \frac{a}{p_1}, \qquad \frac{1}{q} = \frac{1-a}{q_0} + \frac{a}{q_1}.
\]
\end{theorem}

\begin{proof}
通过把测度 $\mu$ 与 $\nu$ 本身乘以标量, 可以假设 $M_0 = M_1 = 1$.采用双线性表述, 我们需要证明

\begin{equation}
\left|\int g(x)T_z f(x)\, dx\right| \le 1\label{eq:c6-182}
\end{equation}

对每对满足 $\|f\|_{L^p} = \|g\|_{L^{q'}} = 1$ 的简单函数 $f, g$ 成立.固定这样一对 $f, g$, 考虑相关的(解析)简单函数族

\[
f_z(x) = |f(x)|^{\frac{p}{p(z)}-1}f(x), \qquad g_z(x) = |g(x)|^{\frac{q'}{q'(z)}-1}g(x),
\]

其中指数为

\[
\frac{1}{p(z)} = \frac{1-z}{p_0} + \frac{z}{p_1}, \qquad \frac{1}{q'(z)} = \frac{1-z}{q_0'} + \frac{z}{q_1'}.
\]

我们容易验证

\[
|f_{ib}| \le |f|^{p/p_0}, \quad |f_{1+ib}| \le |f|^{p/p_1}, \quad |g_{ib}| \le |g|^{q'/q_0'}, \quad |g_{1+ib}| \le |g|^{q'/q_1'}.
\]

这里我们使用约定 $1/\infty = 0$, 特别地, 若 $p_0 = p_1 = \infty$, 则 $p = p(z) = \infty$ 且 $f_z\equiv f$;类似地, 若 $q_0' = q_1' = \infty$, 则 $q' = q'(z) = \infty$ 且 $g_z\equiv g$.不难验证 $\|f_z\|_{L^{\operatorname{Re}(p(z))}} = \|f\|_{L^p} = 1$, 且 $\|g_z\|_{L^{\operatorname{Re}(q'(z))}} = \|g\|_{L^{q'}} = 1$.

现在考虑定义在 $D$ 上的映射

\[
h(z) = \int g_z(x)T_z f_z(x)\, dx.
\]

根据我们的构造以及 $T_z$ 的线性与解析性质, 不难看出 $h\in\mathcal{A}_{BC}$.由假设(以及 Hölder 不等式), 对每个 $b\in\mathbb{R}$ 有 $|h(ib)|\le 1$ 与 $|h(1+ib)|\le 1$.由三线引理推出 $|h(z)|\le 1$, 特别地 \eqref{eq:c6-182}.
\end{proof}

\mtsubsec{Young 不等式}

我们常常需要估计形如\footnote{事实上, Schwartz 核定理表明, 从 $\mathcal{C}_0^\infty(\Omega_1)$ 到 $\mathcal{D}'(\Omega_2)$ 的每个连续线性映射都具有 (183) 的形式, 其中 $k(x, y)\in\mathcal{D}'(\Omega_1\times\Omega_2)$ 是某个分布.}

\begin{equation}
Tf(x) = \int k(x, y)f(y)\, dy, \label{eq:c6-183}
\end{equation}

的积分算子.这类结果中最简单的是下面的 Young 定理.

\begin{theorem}[Young 不等式]\label{thm:c6-1.10}
设 $k(x, y)$ 是可测函数, 并假设对某个 $1\le r\le\infty$, 有

\[
\sup_x\|k(x, \cdot)\|_{L^r} \le 1, \qquad \sup_y\|k(\cdot, y)\|_{L^r} \le 1.
\]

则对 $1\le p\le r'$ 以及

\begin{equation}
1 + \frac{1}{q} = \frac{1}{r} + \frac{1}{p}, \label{eq:c6-184}
\end{equation}

有

\begin{equation}
\|Tf\|_{L^q} \le \|f\|_{L^p}.\label{eq:c6-185}
\end{equation}
\end{theorem}

\begin{proof}
由 Hölder 不等式,

\begin{equation}
\|Tf\|_{L^\infty} \le \|f\|_{L^{r'}}.\label{eq:c6-186}
\end{equation}

另一方面, 对偶算子 $T^*$ 与 $T$ 具有相同的形式,

\[
T^*g(y) = \int \overline{k(x, y)}g(x)\, dx,
\]

因此

\[
\|T^*g\|_{L^\infty} \le \|g\|_{L^{r'}},
\]

由对偶性给出另一个端点

\begin{equation}
\|Tf\|_{L^r} \le \|f\|_{L^1}.\label{eq:c6-187}
\end{equation}

现在, 我们可以用\thmref{thm:c6-1.8}(取 $T_z\equiv T$)在 \eqref{eq:c6-186} 与 \eqref{eq:c6-187} 之间插值, 得到 \eqref{eq:c6-185}.
\end{proof}

作为直接推论, 当 $k$ 平移不变, 即 $k(x, y) = k(x-y)$ 时, 我们得到著名的卷积估计:

\begin{equation}
\|k*f\|_{L^q} \le \|k\|_{L^r}\|f\|_{L^p}, \label{eq:c6-188}
\end{equation}

只要指数 $1\le p, q, r\le\infty$ 满足 \eqref{eq:c6-184}.注意, 指数之间的这种关系是必要的, 它保证两边关于测度具有相同的齐次次数.

\textbf{习题.} 更一般地, 当 $\|k\|_{L_x^\infty L_y^r}\le 1$ 不成立, 且对 $\|k\|_{L_y^\infty L_x^r}$ 也类似时, 可以通过改变 $x$ 与 $y$ 变量中的测度, 归结为\thmref{thm:c6-1.10} 的假设.这样做, 你能得到怎样的"更一般"的 Young 不等式?

\textbf{习题.} 用复插值证明 傅里叶变换 $\mathcal{F}$ 的 Hausdorff-Young 不等式,

\[
\|\mathcal{F}(f)\|_{L^q} \lesssim \|f\|_{L^p}, \qquad \text{对一切 } q\ge 2, \quad 1/q + 1/p = 1.
\]

\mtsubsec{Marcinkiewicz 插值}

比函数 $f$ 的 $L^p$ 可积性稍弱一点的条件是所谓的弱-$L^p$ 性质.

\begin{definition}\label{def:c6-1.12}
对 $1\le p<\infty$, 如果对每个 $\alpha>0$ 都有 $\Lambda(f, \alpha)\lesssim\alpha^{-p}$, 则称 $f$ 属于\newterm{weak-lp}{弱 $L^p$}{weak-$L^p$}.若 $p = \infty$, 我们令弱-$L^\infty$ 与 $L^\infty$ 重合.
\end{definition}

由 Chebyschev 不等式 \eqref{eq:c6-180}, $L^p$ 中的任何函数也在弱-$L^p$ 中.下面是实插值最简单的例子.它适用于\newterm{sublinear-operator}{次线性算子}{sublinear operator}, 即满足

\[
|T(f+g)(x)| \lesssim |Tf(x)| + |Tg(x)|
\]

的算子.

\begin{theorem}\label{thm:c6-1.13}
考虑把 $X$ 上的可测函数映为 $Y$ 上的可测函数的次线性算子 $T$.假设 $T$ 把 $L^{p_i}(X)$ 映到弱-$L^{p_i}(Y)$, 带界

\[
\Lambda(Tf, \alpha) \lesssim \alpha^{-p_i}\|f\|_{L^{p_i}}^{p_i},
\]

其中 $i = 1, 2$, $1\le p_1 < p_2\le\infty$.则对任何满足 $p_1 < p < p_2$ 的 $p$, $T$ 把 $L^p(X)$ 映到 $L^p(Y)$, 带界

\[
\|Tf\|_{L^p} \lesssim \|f\|_{L^p}.
\]
\end{theorem}

\begin{proof}
给定 $f\in L^p(X)$ 与 $\alpha>0$, 我们写 $f = f^\alpha + f_\alpha$, 其中 $f^\alpha(x) = f(x)\cdot(|f(x)| > \alpha)$ 和 $f_\alpha(x) = f(x)\cdot(|f(x)| \le \alpha)$ 是 $f$ 的截断.特别地, 由 Hölder 不等式, $f^\alpha\in L^{p_1}$ 且 $f_\alpha\in L^{p_2}$.

先考虑 $p_2<\infty$ 的情形.由对 $T$ 的假设, 有

\begin{equation}
\Lambda(Tf, 2\alpha) \lesssim \Lambda(Tf^\alpha, \alpha) + \Lambda(Tf_\alpha, \alpha) \lesssim \alpha^{-p_1}\|f^\alpha\|_{L^{p_1}}^{p_1} + \alpha^{-p_2}\|f_\alpha\|_{L^{p_2}}^{p_2}.\label{eq:c6-189}
\end{equation}

利用 $\|Tf\|_{L^p}$ 的分布刻画与 Fubini 定理, 我们推出

\[
\int|Tf(x)|^p\, dx \lesssim \int_{0<\alpha<|f(x)|}|f(x)|^{p_1}\alpha^{p-p_1-1}\, d\alpha\, dx + \int_{|f(x)|\le\alpha}|f(x)|^{p_2}\alpha^{p-p_2-1}\, d\alpha\, dx.
\]

但 $\int_0^{|f(x)|}\alpha^{p-p_1-1}\, d\alpha \simeq |f(x)|^{p-p_1}$, 因为 $p-p_1-1>-1$, 且 $\int_{|f(x)|}^\infty\alpha^{p-p_2-1}\, d\alpha \simeq |f(x)|^{p-p_2}$, 因为 $p-p_2-1<-1$, 结论随之而来.

在 $p_2 = \infty$ 的情形, 证明实际上更简单.我们只需注意到 $|Tf(x)|\gg\alpha$ 蕴含 $|Tf^\alpha(x)|\gg\alpha$, 因为 $|Tf_\alpha(x)|\lesssim\|f_\alpha\|_{L^\infty}\le\alpha$.因此可以把 \eqref{eq:c6-189} 换成

\[
\Lambda(Tf, C\alpha) \lesssim \Lambda(Tf^\alpha, \alpha) \lesssim \alpha^{-p_1}\|f^\alpha\|_{L^{p_1}}^{p_1},
\]

其中 $C$ 是某个正常数, 证明如前进行.
\end{proof}
\section{极大函数、分数次积分及其应用}

\mtsubsec{极大函数}

对某个 $1\le p\le\infty$ 属于 $L^p(\mathbb{R}^n)$ 的函数 $f$, 可能具有非常差的\newterm{regularity}{正则性}{regularity}性质.给定 $\alpha>0$, $|f(x)|>\alpha$ 的点 $x$ 的集合可以仅仅是任意可测集(若 $p<\infty$ 则具有有限测度).人们常常希望用一个与 $f$ 具有(几乎)相同可积性性质、但具有更好局部正则性的正函数来代替 $f$.这可以通过考虑 $f$ 的极大平均来实现.

\begin{definition}\label{def:c6-2.2}
给定 $\mathbb{R}^n$ 上的可测函数, 我们定义其\newterm{maximal-function}{极大函数}{maximal function}为

\[
\mathcal{M}f(x) = \sup_{x\in B}\frac{1}{|B|}\int_B |f(y)|\, dy.
\]

这里的上确界取遍所有包含 $x$ 的 Euclid 球 $B$(不仅限于以 $x$ 为球心的球).
\end{definition}

\begin{remark}\label{rem:c6-2.3}
由定义立即看出 $\mathcal{M}f$ 是下半连续的.事实上, 对每个 $\alpha\ge 0$, 集合 $E_\alpha = \{x\in\mathbb{R}^n : \mathcal{M}f(x) > \alpha\}$ 总是开集:若 $x\in E_\alpha$, 则存在包含 $x$ 的球 $B$, 使得

\begin{equation}
\frac{1}{|B|}\int_B |f(y)|\, dy > \alpha, \label{eq:c6-190}
\end{equation}

这也意味着对每个 $y\in B$ 有 $\mathcal{M}f(y) > \alpha$, 因此 $B\subset E_\alpha$.
\end{remark}

由三角不等式我们还看到, $f\mapsto\mathcal{M}f$ 是次加性算子,

\begin{equation}
\mathcal{M}(f+g)(x) \le \mathcal{M}f(x) + \mathcal{M}g(x).\label{eq:c6-191}
\end{equation}

平均过程可能改善局部正则性, 但由于上确界的缘故, $\mathcal{M}f$ 是否保持 $f$ 的可积性性质并不清楚.若 $f$ 本质有界, 则 $\mathcal{M}f$ 有界, 且

\begin{equation}
\|\mathcal{M}f\|_{L^\infty} \le \|f\|_{L^\infty}.\label{eq:c6-192}
\end{equation}

但是, 若 $f$ 是可积函数, 并不推出 $\mathcal{M}f$ 可积.例如取 $f = \chi_B\in L^1$, 即球的特征函数, 则 $\mathcal{M}f(x)\gtrsim(1+|x|)^{-n}$, 它勉强不属于 $L^1$.幸运的是, 极大函数仍然保留关于 $f$ 的可积性性质的大部分信息.

\begin{theorem}\label{thm:c6-2.4}
若 $f\in L^1$, 则 $\mathcal{M}f$ 弱属于 $L^1$, 其含义是:对 $\alpha>0$, 有

\begin{equation}
|E_\alpha| = \Lambda(\mathcal{M}f(x), \alpha) \lesssim \frac{1}{\alpha}\|f\|_{L^1}.\label{eq:c6-193}
\end{equation}

若 $1<p\le\infty$ 时 $f\in L^p$, 则 $\mathcal{M}f\in L^p$, 且

\begin{equation}
\|\mathcal{M}f\|_{L^p} \lesssim \|f\|_{L^p}.\label{eq:c6-194}
\end{equation}
\end{theorem}

\begin{proof}
命题的第二部分由第一部分以及极大算子的 $L^\infty$ 有界性, 通过 Marcinkiewicz 插值(\thmref{thm:c6-1.13})得出.因此, 我们只需证明 \eqref{eq:c6-193}.

设 $f\in L^1$ 并固定 $\alpha>0$.由\remref{rem:c6-2.3} 中的讨论, 我们可以找到一族球 $\mathcal{B} = \{B\}$, 使得 $E_\alpha = \cup_{B\in\mathcal{B}}B$, 且每个球 $B$ 都满足 \eqref{eq:c6-190}.若这些球两两不相交, 则容易得出结论, 因为此时

\[
|E_\alpha| \le \sum_{B\in\mathcal{B}}|B| < \frac{1}{\alpha}\sum_B\int_B |f(y)|\, dy \le \frac{1}{\alpha}\int_{\mathbb{R}^n}|f(y)|\, dy.
\]

一般情形下这些球并不不相交, 我们必须更加小心.

设 $K$ 是 $E_\alpha$ 的紧子集, 则可以在 $\mathcal{B}$ 中选择覆盖 $K$ 的有限子族 $\mathcal{B}'$.利用下面证明的\newterm{covering-lemma}{覆盖引理}{covering lemma}——\lemref{lem:c6-2.5}\footnote{这有时被称为 Vitali 覆盖引理.}, 我们可以从 $\mathcal{B}'$ 中的球里选出另一个由不相交球组成的有限子族 $\mathcal{B}''$(它可能不再覆盖 $K$), 使得

\[
|\cup_{B'\in\mathcal{B}'}B'| \lesssim \sum_{B''\in\mathcal{B}''}|B''|.
\]

然后如前进行, 我们得到

\[
|K| \lesssim \frac{1}{\alpha}\|f\|_{L^1},
\]

对所有可能的紧集 $K$ 取上确界, 最终得到 \eqref{eq:c6-193}.
\end{proof}

\begin{lemma}\label{lem:c6-2.5}
设 $B_1, \ldots, B_N$ 是 $\mathbb{R}^n$ 中球的有限族, 则可以选出由不相交球组成的子族 $B_{j_1}, \ldots, B_{j_M}$, $M\le N$, 使得

\[
\left|\cup_{j=1}^N B_j\right| \lesssim \sum_{k=1}^M |B_{j_k}|.
\]
\end{lemma}

\begin{proof}
我们可以假设球 $B_j = B(x_j, r_j)$ 按半径非增的顺序标记, $r_1\ge r_2\ge\cdots\ge r_N$.

取 $j_1 = 1$, 则 $B_{j_1}$ 是半径最大的球.然后归纳地定义 $j_{k+1}$ 为与前已选出的球 $B_{j_1}, \ldots, B_{j_k}$ 不相交的球 $B_j$ 中指标的最小者;若不存在这样的球, 则在第 $k$ 步停止.

由这个构造, 每个球 $B_j$ 都与某个选出的球 $B_{j_k}$ 相交, 且 $r_j\le r_{j_k}$, 因此 $B_j\subset B(x_{j_k}, 3r_{j_k})$.这意味着

\[
\left|\cup_{j=1}^N B_j\right| \le \left|\cup_{k=1}^M B(x_{j_k}, 3r_{j_k})\right| \le 3^n\sum_{k=1}^M |B_{j_k}|.
\]
\end{proof}

\mtsubsec{Lebesgue 微分定理}

若函数 $f$ 连续, 则显然

\begin{equation}
\lim_{r\to 0}\frac{1}{|B(x, r)|}\int_{B(x, r)} f(y)\, dy = f(x).\label{eq:c6-195}
\end{equation}

作为\thmref{thm:c6-2.4} 的一个应用, 我们可以证明这个性质对局部可积函数仍然成立.

\begin{corollary}[Lebesgue 微分定理]\label{cor:c6-2.7}
若 $f\in L^1_{\mathrm{loc}}(\mathbb{R}^n)$, 则 \eqref{eq:c6-195} 对几乎每个 $x$ 成立.
\end{corollary}

\begin{proof}
由于该论断是局部的, 我们可以假设 $f\in L^1$.

设 $A_r$ 是由 $A_r f(x) = |B(x, r)|^{-1}\int_{B(x, r)} f(y)\, dy$ 定义的平均算子.证明分两步.首先证明当 $r\to 0$ 时在 $L^1$ 中 $A_r f\to f$, 然后只需证明 $\lim_{r\to 0}A_r f(x)$ 几乎处处存在.

第一步:给定 $\epsilon>0$, 利用 $\mathcal{C}_0$ 在 $L^1$ 中的稠密性, 我们总能找到紧支撑的连续函数 $g$, 它在 $L^1$ 中逼近 $f$, 并且对 $r$ 一致地有 $\|A_r f - A_r g\|_{L^1}\le\|f-g\|_{L^1}<\epsilon$.然后由 $g$ 的一致连续性, 当 $r\to 0$ 时在 $L^1$ 中 $A_r g\to g$, 因此存在 $r_\epsilon$, 使得

\[
\|A_r f - f\|_{L^1} \le \|A_r f - A_r g\|_{L^1} + \|A_r g - g\|_{L^1} + \|f - g\|_{L^1} \le 3\epsilon,
\]

对 $r < r_\epsilon$ 成立.

第二步:我们定义 $L^1$ 函数 $f$ 的振荡为

\[
\Omega f(x) = \limsup_{r\to 0}A_r f(x) - \liminf_{r\to 0}A_r f(x).
\]

振荡是次加性算子, $\Omega(f+g)\le\Omega f + \Omega g$, 并且被极大函数算子控制, $\Omega f\le 2\mathcal{M}f$;此外, 连续函数的振荡为零.若 $g$ 是在 $L^1$ 中逼近 $f$ 的连续函数, 则

\[
\Omega f \le \Omega(f-g) + \Omega g = \Omega(f-g) \le 2\mathcal{M}(f-g).
\]

现在我们可以应用极大函数的弱-$L^1$ 性质, 对任何正的 $\alpha$, 得到

\[
|\{x : \Omega f(x) > \alpha\}| \le |\{x : \mathcal{M}(f-g)(x) > \alpha/2\}| \lesssim \frac{1}{\alpha}\|f-g\|_{L^1}.
\]

由于 $\|f-g\|_{L^1}$ 可以任意小, 我们推出 $f$ 的振荡为正的点的集合是零测集.
\end{proof}

\mtsubsec{分数次积分}

设 $T$ 是作用在 $\mathbb{R}^n$ 上定义的函数上的积分算子, 其核 $k$ 如 \eqref{eq:c6-183} 所示.若我们对 $k(x, y)$ 所知的唯一信息是形如

\[
|k(x, y)| \lesssim |x-y|^{-\gamma}
\]

的衰减估计(对某个 $\gamma>0$), 那么 Young 不等式(\thmref{thm:c6-1.10})不允许我们对 $Tf$ 获得好的控制, 因为函数 $|x|^{-\gamma}$ 勉强不属于 $L^{n/\gamma}$.然而, 卷积具有某种光滑化性质, 它蕴含一些肯定性结果, 包含在下面这个重要定理中, 该定理最初由 Hardy 与 Littlewood 对 $n=1$ 证明, 后来由 Sobolev 推广到 $n>1$.

\begin{theorem}[Hardy--Littlewood--Sobolev 不等式]\label{thm:c6-2.9}
设 $0<\gamma<n$, $1<p<q<\infty$, 使得

\begin{equation}
1 - \frac{\gamma}{n} = \frac{1}{p} - \frac{1}{q}, \label{eq:c6-196}
\end{equation}

则

\begin{equation}
\left\||\cdot|^{-\gamma}*f\right\|_{L^q(\mathbb{R}^n)} \lesssim \|f\|_{L^p(\mathbb{R}^n)}.\label{eq:c6-197}
\end{equation}
\end{theorem}

\begin{proof}
我们可以把与奇异核的卷积分成两部分:

\[
I_\gamma f(x) = |\cdot|^{-\gamma}*f(x) = \int_{|y|\ge R}\frac{f(x-y)}{|y|^\gamma}\, dy + \int_{|y|<R}\frac{f(x-y)}{|y|^\gamma}\, dy,
\]

其中半径 $R$ 是稍后要选取的正常数.我们只用 Hölder 不等式估计第一项,

\[
\left|\int_{|y|\ge R}\frac{f(x-y)}{|y|^\gamma}\, dy\right| \le \|f\|_{L^p}\left(\int_{|y|\ge R}|y|^{-\gamma p'}\, dy\right)^{1/p'} \lesssim R^{\frac{n}{p'}-\gamma}\|f\|_{L^p},
\]

这里需要可积性条件 $\gamma p' > n$, 由 \eqref{eq:c6-196} 它等价于 $q<\infty$.

对第二部分, 我们在奇点周围作\newterm{dyadic-decomposition}{二进分解}{dyadic decomposition}, 得到用极大函数表达的估计:

\[
\begin{aligned}
\left|\int_{|y|<R}\frac{f(x-y)}{|y|^\gamma}\, dy\right| &\le \sum_{k=0}^\infty\int_{2^{-k-1}\le\frac{|y|}{R}\le 2^{-k}}\frac{|f(x-y)|}{|y|^\gamma}\, dy \\
&\lesssim \sum_{k=0}^\infty\frac{1}{(2^{-k}R)^\gamma}\int_{|y|\le 2^{-k}R}|f(x-y)|\, dy \\
&\lesssim \sum_{k=0}^\infty(2^{-k}R)^{n-\gamma}\mathcal{M}f(x) \simeq R^{n-\gamma}\mathcal{M}f(x),
\end{aligned}
\]

这里我们需要 $\gamma<n$ 以保证最后一个几何级数收敛.

此时我们发现, 对每个 $x\in\mathbb{R}^n$ 与每个 $R>0$, 有

\[
\left||\cdot|^{-\gamma}*f(x)\right| \lesssim R^{\frac{n}{p'}-\gamma}\|f\|_{L^p} + R^{n-\gamma}\mathcal{M}f(x),
\]

常数与 $R$ 和 $x$ 无关.我们优化这个不等式:对每个 $x$, 选择半径 $R = R(x)$ 使右端两项相等,

\[
R^{\frac{n}{p'}-\gamma}\|f\|_{L^p} = R^{n-\gamma}\mathcal{M}f(x),
\]

即

\[
R(x) = \left(\frac{\|f\|_{L^p}}{\mathcal{M}f(x)}\right)^{p/n},
\]

又因为 $(n-\gamma)p/n = 1-p/q$, 我们有

\[
|I_\gamma f(x)| \lesssim \|f\|_{L^p}^{1-\frac{p}{q}}\mathcal{M}f(x)^{\frac{p}{q}}.
\]

然后两边取 $L^q$ 范数,

\[
\|I_\gamma f\|_{L^q} \lesssim \|f\|_{L^p}^{1-\frac{p}{q}}\|\mathcal{M}f\|_{L^p}^{\frac{p}{q}}.
\]

若 $p>1$, 我们可以利用极大函数的估计 \eqref{eq:c6-194} 得出结论.
\end{proof}

\begin{remark*}
Hardy-Littlewood-Sobolev 不等式有等价的双线性表述:

\[
\iint\frac{f(x)g(y)}{|x-y|^\gamma}\, dx\, dy \lesssim \|f\|_{L^{p_1}}\|g\|_{L^{p_2}},
\]

其中 $0<\gamma<n$, $1<p_1, p_2<\infty$, 且

\[
\frac{1}{p_1'} + \frac{1}{p_2'} = \frac{\gamma}{n}.
\]

重要的是要理解, 指数之间的关系可以很快由尺度论证推出.若给变量 $x$ 赋予长度尺度 $L$, 则表达式

\[
\left\||\cdot|^{-\gamma}*f\right\|_{L^q(\mathbb{R}^n)} = \left\|\int|x-y|^{-\gamma}f(y)\, dy\right\|_{L_x^q(\mathbb{R}^n)}
\]

具有单位 $L^{-\gamma}\cdot L^n\cdot L^{\frac{n}{q}}$, 而 $\|f\|_{L^p}$ 具有单位 $L^{\frac{n}{p}}$.指数 $\gamma, q, p$ 必须以使两个量的指数匹配的方式相联系.事实上, 若非如此, 则可以取任意非平凡的 $f$ 并对其作尺度变换导出矛盾, 从而推出估计 \eqref{eq:c6-197} 不成立.
\end{remark*}

\begin{remark*}
在我们对 \eqref{eq:c6-197} 的证明中, 并没有充分利用极大函数的力量(例如, 只考虑了以 $x$ 为球心的球).事实上, 把核 $|x-y|^{-\gamma}$ 换成任何具有相同分布函数的核 $k(x, y)$, 同样的估计仍然成立.沿着这条思路的证明需要建立 Lorentz 空间的理论以及更一般形式的 Marcinkiewicz 插值定理.对此我们参考 (***).
\end{remark*}

利用 Hardy-Littlewood-Sobolev 不等式, 我们现在说明可以给出 Sobolev 不等式的一个非常简短的证明,

\[
\|f\|_{L^q} \lesssim \|\partial f\|_{L^p},
\]

其中 $n/q = n/p-1$, 在非临界区域 $p>1$.与 Hardy-Littlewood-Sobolev 不等式一样, 指数可以通过考虑两边的长度尺度轻松导出.Sobolev 不等式定量地表达了如下直观事实:一个函数只有在它的导数以更坏的方式爆破时, 才可能在某个地方爆破.

\begin{proof}
假设 $f\in\mathcal{C}_0^\infty(\mathbb{R}^n)$.对每个单位向量 $\omega$, 有

\[
f(x) = -\int_0^\infty\frac{d}{dr}f(x+\omega r)\, dr.
\]

为了考虑 $f$ 可能增长的所有方向, 我们对整个球面求积分, 并回忆 $\mathbb{R}^n$ 中极坐标下的体积元为 $dy = r^{n-1}\, dr\, d\sigma_\omega$, 我们发现

\[
|f(x)| \lesssim \int\frac{|\partial f(y)|}{|x-y|^{n-1}}\, dy = (|\cdot|^{1-n}*|\partial f|)(x).
\]

我们取 $L^q$ 范数并利用 \eqref{eq:c6-197}, 得到

\[
\|f\|_{L^q} \lesssim \||\cdot|^{1-n}*|\partial f|\|_{L^q} \lesssim \|\partial f\|_{L^p},
\]

只要 $p>1$ 且

\[
1 - \frac{n-1}{n} = \frac{1}{p} - \frac{1}{q}.
\]
\end{proof}
\mtsubsec{Sobolev 不等式}

在上一小节中, 我们已经看到如何用 $f$ 的梯度的 $L^p$ 范数($1 - \frac{n-1}{n} = \frac{1}{p} - \frac{1}{q}$, $p>1$)来估计函数在 $L^q(\mathbb{R}^n)$ 中的范数.我们现在证明它的一个更强的版本.

\begin{theorem}[Gagliardo--Nirenberg--Sobolev 不等式]\label{thm:c6-2.11}
不等式

\begin{equation}
\|f\|_{L^q(\mathbb{R}^n)} \lesssim \|\partial^m f\|_{L^p(\mathbb{R}^n)}, \qquad f\in\mathcal{C}_0^\infty(\mathbb{R}^n), \label{eq:c6-198}
\end{equation}

在

\begin{equation}
\frac{1}{q} = \frac{1}{p} - \frac{m}{n} > 0, \qquad m\in\mathbb{N}, \qquad (1\le p < q < \infty)\label{eq:c6-199}
\end{equation}

时成立.而当 $q=\infty$ 时, 有

\begin{equation}
\|f\|_{L^\infty(\mathbb{R}^n)} \lesssim \sum_{k=0}^m\|\partial^k f\|_{L^p(\mathbb{R}^n)}, \qquad f\in\mathcal{C}_0^\infty(\mathbb{R}^n), \label{eq:c6-200}
\end{equation}

当 $\frac{n}{p} - m < 0$ 时成立.
\end{theorem}

\begin{remark*}
我们不需要记住精确的条件 \eqref{eq:c6-199};它可以通过简单的量纲分析推出.由于估计是齐次的, 它必须在伸缩下不变, 而 \eqref{eq:c6-199} 只是说 \eqref{eq:c6-198} 两边具有相同的尺度.条件 $\frac{n}{p}-m<0$ 也是对 \eqref{eq:c6-200} 两边尺度的比较, 它排除了一个非常局部化、尖锐的反例.
\end{remark*}

\begin{remark*}
估计 \eqref{eq:c6-198} 的下列非临界版本也对一切 $1\le p<q<\infty$ 且 $1/p - m/n < 1/q$ 成立:

\begin{equation}
\|f\|_{L^q(\mathbb{R}^n)} \lesssim \sum_{|\alpha|\le m}\|\partial^\alpha f\|_{L^p}, \qquad f\in\mathcal{C}_0^\infty(\mathbb{R}^n), \label{eq:c6-201}
\end{equation}
\end{remark*}

\textbf{习题.} 举例说明不等式 \eqref{eq:c6-200} 在 $p = n/m$ 时可以不成立.利用\thmref{thm:c6-2.11} 的结果证明 $m=1$ 时的 \eqref{eq:c6-201}.

\textbf{习题.} 用尺度论证说明:若不等式 \eqref{eq:c6-201} 在 $1/p = 1/q + m/n < 0$ 时成立, 则齐次不等式 \eqref{eq:c6-198} 也成立.

\begin{proof}[证明 \eqref{eq:c6-198}]
$m>1$ 的情形通过对 $m=1$ 的情形反复迭代得到.因此我们可以假设 $m=1$, 由 \eqref{eq:c6-199},

\[
1\le p < n, \qquad \frac{n}{n-1}\le q = \frac{np}{n-p}<\infty.
\]

一旦我们有了 $p=1$、$q = n/(n-1)$ 的估计, 则 $p>1$、$q > n/(n-1)$ 的情形只需应用 Hölder 不等式即可得到.事实上, 设 $q = \lambda n/(n-1)$, 其中 $\lambda>1$, 则

\[
\|f\|_{L^q}^\lambda = \||f|^\lambda\|_{L^{\frac{n}{n-1}}} \lesssim \||f|^{\lambda-1}\partial f\|_{L^1} \le \||f|^{\lambda-1}\|_{L^{p'}}\|\partial f\|_{L^p},
\]

我们只需验证

\[
(\lambda-1)p' = \frac{\frac{n-1}{n}q - 1}{1 - \frac{1}{n} - \frac{1}{q}} = q.
\]

但这本质上无需验证——由不等式的尺度性质, 指数必然匹配.

现在只剩下证明特殊情形 $m=1$, $p=1$, $q = n/(n-1)$.参照 Nirenberg 的著作[?], 可以证明更强的结果:对 $f\in\mathcal{C}_0^\infty(\mathbb{R}^n)$, 有

\begin{equation}
\|f\|_{L^{\frac{n}{n-1}}(\mathbb{R}^n)} \lesssim \prod_{j=1}^n\|\partial_j f\|_{L^1(\mathbb{R}^n)}^{1/n}.\label{eq:c6-202}
\end{equation}

当 $n=1$ 时, 这容易由写出

\[
f(x) = \int_{-\infty}^x f'(y)\, dy
\]

得到.当 $n=2$ 时, 我们对每个变量做同样的事情, 然后相乘并积分:

\[
\begin{aligned}
\iint|f(x_1, x_2)|^2\, dx_1\, dx_2 &\le \iiint|\partial_1 f(y_1, x_2)|\, dy_1\int|\partial_2 f(x_1, y_2)|\, dy_2\, dx_1\, dx_2 \\
&= \|\partial_1 f\|_{L^1}\|\partial_2 f\|_{L^1}.
\end{aligned}
\]

当 $n\ge 3$ 时事情变得更加棘手, 为了分离变量, 我们必须反复使用 Hölder 不等式.只看 $n=3$ 的情形.为简化记号, 令 $f_j = \partial_j f$, 并令 $\int\phi(x)\, dx_j = \int_j\phi(\hat{x}_j)$.我们从

\[
|f(x)|^{\frac{3}{2}} \le \left(\int_1|f_1(\cdot, x_2, x_3)|\right)^{\frac{1}{2}}\left(\int_2|f_2(x_1, \cdot, x_3)|\right)^{\frac{1}{2}}\left(\int_3|f_3(x_1, x_2, \cdot)|\right)^{\frac{1}{2}}
\]

开始.然后关于 $x_1$ 积分.右端第一个因子不依赖于 $x_1$, 而我们用 Hölder 把第二个与第三个因子分开:

\[
\int_1|f(\cdot, x_2, x_3)|^{\frac{3}{2}} \le \left(\int_1|f_1(\cdot, x_2, x_3)|\right)^{\frac{1}{2}}\left(\int_{1, 2}|f_2(\cdot, \cdot, x_3)|\right)^{\frac{1}{2}}\left(\int_{1, 3}|f_3(\cdot, x_2, \cdot)|\right)^{\frac{1}{2}}.
\]

对关于 $x_2$ 的积分同样进行:

\[
\int_{1, 2}|f(\cdot, \cdot, x_3)|^{\frac{3}{2}} \le \left(\int_{1, 2}|f_1(\cdot, \cdot, x_3)|\right)^{\frac{1}{2}}\left(\int_{1, 2}|f_2(\cdot, \cdot, x_3)|\right)^{\frac{1}{2}}\left(\int_{1, 2, 3}|f_3(\cdot)|\right)^{\frac{1}{2}},
\]

最后对 $x_3$ 也这样做:

\[
\int_{1, 2, 3}|f(\cdot)|^{\frac{3}{2}} \le \left(\int_{1, 2, 3}|f_1(\cdot)|\right)^{\frac{1}{2}}\left(\int_{1, 2, 3}|f_2(\cdot)|\right)^{\frac{1}{2}}\left(\int_{1, 2, 3}|f_3(\cdot)|\right)^{\frac{1}{2}}.
\]

当 $n>3$ 时, 过程完全相同.
\end{proof}

\begin{proof}[证明 \eqref{eq:c6-200}]
显然只需考虑 $m=1$ 的情形, 因为 $m>1$ 的情形应用 \eqref{eq:c6-198} 即可由它推出.因此假设 $m=1$ 且 $p>n$, 我们要证明

\[
|f(0)| \lesssim \|f\|_{L^p} + \|Df\|_{L^p}.
\]

首先假设 $f$ 的支集包含在单位球 $B = \{|x|<1\}$ 中, 则

\begin{equation}
f(0) = -\int_0^1\frac{d}{dr}f(r\omega)\, dr, \qquad \omega\in\mathbb{S}^{n-1}.\label{eq:c6-203}
\end{equation}

关于 $\omega$ 积分, 然后应用 Hölder 不等式,

\begin{equation}
|f(0)| \lesssim \int_B\frac{|\partial f(x)|}{|x|^{n-1}}\, dx \lesssim \|\partial f\|_{L^p}\left(\int_B\frac{dx}{|x|^{(n-1)p'}}\right)^{1/p'} \lesssim \|\partial f\|_{L^p}, \label{eq:c6-204}
\end{equation}

这里需要的可积性条件是 $(n-1)p'<n$, 它正是 $p>n$.

一般情形, 固定支集在 $B$ 中且 $\phi(0)=1$ 的截断函数 $\phi\in\mathcal{C}_0^\infty$, 则由上述讨论, $|f(0)| = |\phi(0)f(0)| \lesssim \|\partial(\phi f)\|_{L^p} \lesssim \|f\|_{L^p} + \|\partial f\|_{L^p}$.
\end{proof}

\mtsubsec{经典 Sobolev 空间}

\thmref{thm:c6-2.11} 的 Sobolev 不等式引导我们引入\newterm{sobolev-space}{Sobolev 空间}{Sobolev space}.

\begin{definition}\label{def:c6-2.13}
设 $\Omega$ 是 $\mathbb{R}^n$ 的开子集.固定 $1\le p\le\infty$, 并设 $s\in\mathbb{N}$ 是非负整数.空间 $W^{s, p}(\Omega)$ 由 $\Omega$ 上所有局部可积的实值(或复值)函数 $u$ 组成, 使得对满足 $|\alpha|\le s$ 的所有多重指标 $\alpha$, 弱\footnote{即按分布意义理解的导数.}导数 $\partial^\alpha u$ 都属于 $L^p(\Omega)$.这些空间配备范数:

\[
\begin{aligned}
\|u\|_{W^{s, p}(\Omega)} &= \left(\sum_{|\alpha|\le s}\|\partial^\alpha u\|_{L^p(\Omega)}^p\right)^{1/p}, \qquad \text{对 } 1\le p<\infty, \\
\|u\|_{W^{s, \infty}(\Omega)} &= \sum_{|\alpha|\le s}\|\partial^\alpha u\|_{L^\infty(\Omega)}.
\end{aligned}
\]
\end{definition}

我们还用 $W_0^{k, p}(\Omega)$ 表示 $\mathcal{C}_0^\infty(\Omega)$ 在 $W^{k, p}(\Omega)$ 中的闭包.

特别地, 当 $p=2$ 时, 我们记 $H^s(\Omega) = W^{s, 2}(\Omega)$.显然 $H^0(\Omega) = L^2(\Omega)$.我们也记 $H_0^s(\Omega) = W_0^{s, 2}(\Omega)$.这些空间因其 Hilbert 空间结构而尤为重要.

特别地, 当 $p=\infty$ 时, 我们使用更小的空间 $\mathcal{C}^s(\overline{\Omega})\subset W^{s, \infty}(\Omega)$, 即 $s$ 次连续可微且具有有界 $\|\cdot\|_{W^{s, \infty}}$ 范数的函数所成的集合.

\textbf{习题.} 证明对每个 $s\in\mathbb{N}$ 与 $1\le p\le\infty$, 空间 $W^{s, p}(\Omega)$ 都是 Banach 空间.

关于区域 $\Omega\subset\mathbb{R}^n$ 上的 Sobolev 空间, 还有很多可说的.例如, 我们请读者参阅 Evans 的著作 (\cite{e}, 第 5 章).现在我们把注意力限制到 $\Omega = \mathbb{R}^n$ 的情形.

\textbf{习题.} 证明空间 $W^{k, p}(\mathbb{R}^n)$ 与 $W_0^{k, p}(\mathbb{R}^n)$ 重合.这意味着 $\mathcal{C}_0^\infty$ 在 $W^{k, p}(\mathbb{R}^n)$ 中稠密.

上一小节证明的 Sobolev 不等式可以解释为\newterm{embedding-theorem}{嵌入定理}{embedding theorem}.事实上, \eqref{eq:c6-198} 与 \eqref{eq:c6-201} 可以解释为:只要 $\frac{1}{p} - \frac{m}{n}\le\frac{1}{q}$, Sobolev 空间 $W^{m, p}(\mathbb{R}^n)$ 就包含在 Lebesgue 空间 $L^q(\mathbb{R}^n)$ 中.

\begin{proposition}\label{prop:c6-2.14}
下列包含关系是连续的:

\[
W^{m, p}(\mathbb{R}^n)\subset L^q(\mathbb{R}^n), \qquad \text{若 } \frac{1}{p} - \frac{m}{n}\le\frac{1}{q}.
\]

此外, 对 $q=\infty$, 只要 $m > n/p$, $W^{m, p}(\mathbb{R}^n)$ 就嵌入 $\mathbb{R}^n$ 上有界连续函数的空间.
\end{proposition}

\begin{proof}
由\thmref{thm:c6-2.11} 以及 $\mathcal{C}_0^\infty(\mathbb{R}^n)$ 在 $W^{m, p}(\mathbb{R}^n)$ 中的稠密性得出.
\end{proof}

\mtsubsec{Hölder 空间}

与 Sobolev 空间一起, \newterm{holder-space}{Hölder 空间}{Hölder space}在分析中起着非常重要的作用, 特别是在与椭圆方程的联系中.在引入这些空间之前, 我们先回顾 $m$ 次连续可微函数 $u:\Omega\to\mathbb{R}$ 的空间 $\mathcal{C}^m(\overline{\Omega})$ 的定义, 其中 $\Omega$ 是开区域, 且 $W^{m, \infty}$ 范数有界:

\[
\|u\|_{\mathcal{C}^m(\overline{\Omega})} = \sum_{|\alpha|\le m}\|\partial^\alpha u(x)\|_{L^\infty(\Omega)} < \infty.
\]

\begin{definition}\label{def:c6-2.16}
设 $\Omega$ 是 $\mathbb{R}^n$ 中的开区域.如果

\begin{equation}
[u]_{\mathcal{C}^{0, \gamma}(\overline{\Omega})} = \sup_{x\neq y\in\Omega}\frac{|u(x)-u(y)|}{|x-y|^\gamma} < \infty, \label{eq:c6-205}
\end{equation}

则称函数 $u:\Omega\to\mathbb{R}$ 是指数为 $0<\gamma\le 1$ 的 Hölder 连续函数.
\end{definition}

Hölder 空间 $\mathcal{C}^{k, \gamma}(\overline{\Omega})$ 由所有满足范数

\begin{equation}
\|u\|_{\mathcal{C}^{k, \gamma}(\overline{\Omega})} = \|u\|_{\mathcal{C}^k(\overline{\Omega})} + \sum_{|\alpha|=k}[\partial^\alpha u]_{\mathcal{C}^{0, \gamma}(\overline{\Omega})}\label{eq:c6-206}
\end{equation}

有限的函数 $u\in\mathcal{C}^k(\overline{\Omega})$ 组成.

\mtsubsec{习题}

空间 $\mathcal{C}^{k, \gamma}(\overline{\Omega})$ 是 Banach 空间.

\mtsubsec{习题}

证明 $\mathcal{C}^{0, 1}((a, b))$(区间上的 Lipschitz 函数空间)恰好由那些导数属于 $L^\infty$ 的分布组成.

\mtsubsec{习题}

设 $f(x) = (a\le x\le b)$ 是区间的特征函数.证明七重卷积 $f*\cdots*f$ 属于 Hölder 类 $\mathcal{C}^{5, 1}$.

下面的 $L^\infty$ 中 Sobolev 嵌入的更强版本在椭圆理论中很重要.和往常一样, 所涉指数之间的关系可以由量纲分析推出.

\begin{theorem}[Morrey 不等式]\label{thm:c6-2.20}
假设 $n<p\le\infty$.则对一切 $u\in\mathcal{C}_0^\infty(\mathbb{R}^n)$, 有

\begin{equation}
\|u\|_{\mathcal{C}^{0, \gamma}(\mathbb{R}^n)} \lesssim \|u\|_{W^{1, p}(\mathbb{R}^n)}, \label{eq:c6-207}
\end{equation}

只要 $\gamma = 1 - n/p$.
\end{theorem}

\begin{proof}
在一维情形, 这是微积分基本定理与 Hölder 不等式的简单应用.一般情形见 \cite{e} 第 5.6.2 节.
\end{proof}
\mtsubsec{分数次 $H^s$-Sobolev 空间}

考虑 Sobolev 空间

\[
H^s(\mathbb{R}^n) = \{u\in L^2 : \partial^\alpha u\in L^2, \quad \forall|\alpha|\le s\}.
\]

\begin{proposition}\label{prop:c6-2.22}
Sobolev 空间 $H^s(\mathbb{R}^n)$ 与所有满足下列条件的缓增分布 $u\in\mathcal{S}'(\mathbb{R}^N)$ 所成的集合重合:$\hat{u}$ 局部可积, 且

\begin{equation}
\|u\|_{H^s}^2 = \int_{\mathbb{R}^N}(1+|\xi|^2)^s|\hat{u}(\xi)|^2 < \infty.\label{eq:c6-208}
\end{equation}
\end{proposition}

\begin{proof}
由 Parseval 恒等式、$\mathcal{C}_0^\infty$ 在每个空间中的稠密性以及 $\widehat{\partial_j f}(\xi) = i\xi_j\hat{f}(\xi)$ 这一事实容易得出.
\end{proof}

注意, \propref{prop:c6-2.22} 的等价定义不仅对正整数有意义, 而且对一切实数 $s$ 都有意义.因此我们可以谈论所有实数值 $s$ 的 Sobolev 空间 $H^s$.我们还将用到下列齐次 Sobolev 范数, 对一切 $s\ge 0$,

\begin{equation}
\|u\|_{\dot{H}^s}^2 = \int_{\mathbb{R}^N}|\xi|^{2s}|\hat{u}(\xi)|^2 < \infty.\label{eq:c6-209}
\end{equation}

\textbf{问题.} 为什么若把物理空间变量的单位视为 $L$, 则 $\|u\|_{\dot{H}^s}$ 具有单位 $L^{\frac{N}{2}-s}$?

\textbf{习题.} 对 $s\in(0, 1)$, 空间 $H^s(\mathbb{R}^n)$ 与满足

\begin{equation}
\left(\int_{\mathbb{R}^n}\int_{\mathbb{R}^n}\frac{|u(x)-u(x+y)|^2}{|y|^{n+2s}}\, dx\, dy + \|u\|_{L^2(\mathbb{R}^n)}^2\right)^{1/2} < \infty\label{eq:c6-210}
\end{equation}

的局部可积函数所成的空间重合.

\textbf{习题.} 证明对 $s > n/2$, Sobolev 空间 $H^s(\mathbb{R}^n)$ 嵌入有界连续函数的空间.

\mtsubsec{一个迹定理}

为了给偏微分方程中广义函数的边界值赋予意义, 重要的是证明限制运算(它对连续函数显然有意义)在函数不连续时仍然有意义.这样的定理称为迹定理.为简单起见, 考虑 $\mathbb{R}^n$ 中超平面 $x_n = 0$ 的情形, 定义\newterm{trace-operator}{迹算子}{trace operator},

\begin{equation}
Tf(x_1, \ldots, x_n) = f(x_1, \ldots, x_{n-1}, 0).\label{eq:c6-211}
\end{equation}

显然, 该算子对 $\mathbb{R}^n$ 中任何连续函数 $f$ 都有意义, 特别地对任何检验函数都有意义.

\begin{theorem}\label{thm:c6-2.24}
下列估计对任何检验函数 $f\in\mathcal{C}_0^\infty(\mathbb{R}^n)$, $n\ge 2$, 以及任何 $s>1/2$, 一致地成立:

\begin{equation}
\|Tf\|_{H^{s-\frac{1}{2}}(\mathbb{R}^{n-1})} \lesssim \|f\|_{H^s(\mathbb{R}^n)}.\label{eq:c6-212}
\end{equation}

因此 $T$ 延拓为线性映射 $T : H^s(\mathbb{R}^n)\to H^{s-1/2}(\mathbb{R}^{n-1})$.
\end{theorem}

\begin{remark*}
注意这个结果在量纲上是尖锐的, 如果把它与 Sobolev 空间 $H^s(\mathbb{R})$ 到 $L^\infty(\mathbb{R})$ 的通常嵌入相比较, 这多少有些令人惊讶, 因为在那种情形我们知道尖锐情形 $s = 1/2$ 是不成立的.事实上, 上述迹定理在 $s = 1/2$ 的情形也不成立.
\end{remark*}

\begin{proof}
取 $f$ 光滑, 令 $g(x') = f(x', 0)$.设 $\tilde{f}$ 是 $f$ 仅关于 $x_n$ 的 傅里叶变换, 而 $\hat{f}$、$\hat{g}$ 分别是 $f$ 与 $g$ 在 $\mathbb{R}^n$ 和 $\mathbb{R}^{n-1}$ 中的 傅里叶变换.

即

\[
\tilde{f}(x', \xi_n) = \int_{-\infty}^{\infty}f(x', x_n)e^{-ix_n\xi_n}\, dx_n.
\]

应用(取 $x_n = 0$ 的)Fourier 反演, 然后作 傅里叶变换, 得到

\[
\begin{aligned}
g(x') &= f(x', 0) = \frac{1}{2\pi}\int_{-\infty}^{\infty}\tilde{f}(x', \xi_n)\, d\xi_n, \\
\hat{g}(\xi') &= \frac{1}{2\pi}\int_{-\infty}^{\infty}\hat{f}(\xi', \xi_n)\, d\xi_n.
\end{aligned}
\]

然后, 利用我们对分数次 $H^s$ 空间的了解以及 Cauchy-Schwartz 不等式, 可以看到:

\[
\begin{aligned}
\|g\|_{H^{s-1/2}(\mathbb{R}^{n-1})} &\lesssim \int_{\mathbb{R}^{n-1}}|\hat{g}(\xi')|^2(1+|\xi'|^2)^{s-1/2}\, d\xi' \\
&\lesssim \int_{\mathbb{R}^{n-1}}\left|\int_{-\infty}^{\infty}\hat{f}(\xi)\, d\xi_n\right|^2(1+|\xi'|^2)^{s-1/2}\, d\xi' \\
&\lesssim \int_{\mathbb{R}^{n-1}}(1+|\xi'|^2)^{s-1/2}\left(\int_{-\infty}^{\infty}|\hat{f}(\xi)|^2(1+|\xi|^2)^s\, d\xi_n\right)J(\xi')\, d\xi',
\end{aligned}
\]

其中

\[
J(\xi') = \int_{-\infty}^{\infty}(1+|\xi|^2)^{-s}\, d\xi_n.
\]

又因为 $s>1/2$, 我们有

\[
\begin{aligned}
J(\xi') &= \int_{-\infty}^{\infty}(1+|\xi|^2)^{-s}\, d\xi_n = \int_{-\infty}^{\infty}(1+|\xi'|^2+|\xi_n|^2)^{-s}\, d\xi_n \\
&= (1+|\xi'|^2)^{-s+1/2}\int_{-\infty}^{\infty}(1+y^2)^{-s}\, dy.
\end{aligned}
\]

把它代入上面关于 $\|g\|_{H^{s-1/2}}$ 的估计, 就证明了结果.
\end{proof}

对于更高余维超曲面上的迹, 也有类似的结果.下面就是这样一个结果, 它可以用初等方法证明.

\begin{proposition}\label{prop:c6-2.25}
考虑 $\mathbb{R}^3$ 中的迹算子 $T$, 它把连续函数 $f(t, x^1, x^2)$ 映为 $Tf(t) = f(t, 0, 0)$.对任何检验函数 $f$, 有

\begin{equation}
\|\partial_t(Tf)\|_{L^2(\mathbb{R})} \lesssim \|\partial^2 f\|_{L^2(\mathbb{R}^3)}.\label{eq:c6-213}
\end{equation}

因此 $T$ 延拓为以 $H^2(\mathbb{R}^3)$ 为定义域、取值于 $H^1(\mathbb{R})$ 的有界线性算子.
\end{proposition}

\begin{proof}

\[
\begin{aligned}
\int_{\mathbb{R}}|\partial_t f(t, 0, 0)|^2\, dt &= \int_0^\infty dx^1\int_0^\infty dx^2\left(\partial_1\partial_2\int_{\mathbb{R}}\partial_t f(t, x)\partial_t f(t, x)\, dt\right) \\
&= 2\int_0^\infty dx^1\int_0^\infty dx^2\left(\int_{\mathbb{R}}\partial_1\partial_2\partial_t f(t, x)\partial_t f(t, x)\, dt\right) \\
&\quad + 2\int_0^\infty dx^1\int_0^\infty dx^2\left(\int_{\mathbb{R}}\partial_1\partial_t f(t, x)\partial_2\partial_t f(t, x)\, dt\right).
\end{aligned}
\]

显然,

\[
\left|\int_0^\infty dx^1\int_0^\infty dx^2\left(\int_{\mathbb{R}}\partial_1\partial_t f(t, x)\partial_2\partial_t f(t, x)\, dt\right)\right| \lesssim \|\partial^2 f\|_{L^2(\mathbb{R}^3)}^2.
\]

另一方面, 关于 $t$ 分部积分,

\[
\int_{\mathbb{R}}\partial_1\partial_2\partial_t f(t, x)\partial_t f(t, x)\, dt = -\int_{\mathbb{R}}\partial_1\partial_2 f(t, x)\partial_t^2 f(t, x)\, dt.
\]

因此,

\[
\left|\int_0^\infty dx^1\int_0^\infty dx^2\left(\int_{\mathbb{R}}\partial_1\partial_2\partial_t f(t, x)\partial_t f(t, x)\, dt\right)\right| \lesssim \|\partial^2 f\|_{L^2(\mathbb{R}^3)}.
\]

如所期望.
\end{proof}

\textbf{习题.} 用 傅里叶变换证明同样的结果, 并将其推广到所有维数与一般的 $H^s$ 空间.

\textbf{习题.} 把该结果推广到有界区间上的 $t$.

\mtsubsec{延拓}

为了把对 $\mathbb{R}^n$ 中函数成立的结论推广到 $\mathbb{R}^n$ 中的区域, 我们需要以受控的方式\newterm{extension}{延拓}{extension}这些函数.我把讨论限制在半空间 $\mathbb{R}_+^n = \{x\in\mathbb{R}^n / x^n\ge 0\}$ 的情形.考虑 Sobolev 空间 $W^{1, p}(\mathbb{R}_+^n)$.我们要证明如下结论.

\begin{proposition}\label{prop:c6-2.27}
存在有界线性算子 $E : W^{1, p}(\mathbb{R}_+^n)\to W^{1, p}(\mathbb{R}^n)$, 使得对任何连续的 $f$, 有

\[
Ef|_{\mathbb{R}_+^n} = f
\]

并且

\[
\|Ef\|_{W^{1, p}(\mathbb{R}^n)} \lesssim \|f\|_{W^{1, p}(\mathbb{R}_+^n)}.
\]
\end{proposition}

\begin{proof}
只需对函数 $f\in\mathcal{C}^1(\mathbb{R}_+^n)$ 证明该结果.给定这样的函数, 我们利用其高阶反射定义其延拓 $\bar{f}$, 它在 $\mathbb{R}_+^n$ 中与 $f$ 重合, 并且对所有 $x^n<0$,

\[
\bar{f}(x', x^n) = -3f(x', -x^n) + 4f\left(x', -\frac{1}{2}x^n\right).
\]

首先注意到 $\bar{f}$ 也是 $\mathcal{C}^1$ 的.事实上, $\bar{f}$ 在 $x^n = 0$ 处连续, 其关于变量 $x' = (x^1, \ldots, x^{n-1})$ 的导数也连续.另一方面, 对 $x^n<0$,

\[
\partial_n\bar{f}(x', x^n) = 3\partial_n f(x', -x^n) - 2\partial_n f\left(x', -\frac{1}{2}x^n\right).
\]

因此, 令 $x^n$ 从 $x^n<0$ 趋于零,

\[
(\partial_n\bar{f})^-(x', 0) = \partial_n f(x', 0).
\]

利用这些计算, 我们立刻得到所需的估计.
\end{proof}

\textbf{习题.} 把该结果推广到 $s\in\mathbb{N}$ 的 $W^{s, p}$ 空间.分数次 $H^s$ 空间又如何?

% ==== part06a2 第 6 章 基本函数不等式(第 3--4 节)====
% 源: MinerU_latex_lecture nots.tex 第 9185--9902 行(第 5 节自第 9904 行起归 part06b);
%     pdf_text.txt 第 7558--8150 行(PDF 第 140--150 页)
%
% 术语(新增,键=中文=英文):
%   littlewood-paley-theory=Littlewood-Paley 理论、dyadic-annulus=二进环带、
%   scaling-identity=伸缩恒等式、cutoff-operator=截断算子、almost-orthogonality=几乎正交性、
%   finite-band-property=有限带宽性质、commutator=换位子、square-function=平方函数、
%   maximal-function=极大函数、multiplier=乘子、besov-space=Besov 空间、
%   triebel-space=Triebel 空间、trichotomy-formula=三分公式、paradifferential-calculus=仿微分演算、
%   product-estimate=乘积估计、wedge-product=楔积、hodge-duality=Hodge 对偶、
%   canonical-solution=典范解、wentes-inequality=Wente 不等式
% 注: 隆起函数(bump function)、截断函数(cutoff function)已在第 3 章(part03)定义,不再重复 \newterm.
%
% REFS(需跨章/待解析项):
%   - 同章公式引用: 卷积不等式 (188)、Sobolev 不等式 (201) 在第 6 章第 1--2 节(part06a1)中,
%     本部分以 \eqref{eq:c6-188}、\eqref{eq:c6-201} 引用,依赖 part06a1 定义相应标签.
%   - 脚注(2 处,已译入 \footnote): ① 第 9288 行 "cut-off operator^5"(pdf_text 第 7651--7653 行);
%     ② 第 9754 行 "The proof^6"(pdf_text 第 8018 行).
%   - OCR-BROKEN: 原书 Theorem 4.1 在源 OCR 中缺失(第 9821--9830 行之间),已按 pdf_text
%     第 8079--8080 行还原为 \label{thm:c6-4.1}.
%   - 源第 9199 行 "ξ≠0 2"、第 9586 行 "1/q . 8"、第 9747 行 "s>n/2 19" 中孤立的数字 2/8/19
%     经与 pdf_text 核对均为 OCR 噪音(原书无对应脚注),已删除.
%
% 编号说明:
%   - 第 3 节编号对象: Theorem 3.1 → 3.2--3.5 为带编号子标题(\mtsubsec)→ Lemma 3.6 → Theorem 3.7.
%   - 第 4 节编号对象: Theorem 4.1(OCR 缺失已还原)→ Lemma 4.2.
%   - 本部分显式编号公式为原书 (214)--(245),按项目约定全部行间公式 \[...\] 自动编号,
%     显式编号公式加 \label{eq:c6-NNN};正文引用用 \eqref.
%   - 原书无编号的 Definition(乘子算子相似性、Besov 空间)与全部 Remark 用带 * 环境.

\section{Littlewood-Paley 理论}

在其最简单的形式中, \newterm{littlewood-paley-theory}{Littlewood-Paley 理论}{Littlewood-Paley theory}是一种系统且非常有用的方法:把定义在 $\mathbb{R}^n$ 上的函数 $f$ 分解为无穷二进和 $f = \sum_{k\in\mathbb{Z}} f_k$, 其中频率局部化的分量 $f_k$ 满足:对\newterm{dyadic-annulus}{二进环带}{dyadic annulus} $2^{k-1} \leq |\xi| \leq 2^{k+1}$ 之外的所有 $\xi$, 都有 $\hat{f}_k(\xi) = 0$.这样的分解很容易实现:在 Fourier 空间中选取一个支集在 $\frac{1}{2} \leq |\xi| \leq 2$ 中的检验函数 $\chi(\xi)$, 使得对所有 $\xi \neq 0$,

\begin{equation}
\sum_{k\in\mathbb{Z}} \chi(2^{-k}\xi) = 1.\label{eq:c6-214}
\end{equation}

事实上, 取 $\phi(\xi)$ 为支集在 $|\xi| \leq 2$ 中的实值径向隆起函数, 它在球 $|\xi| \leq 1$ 上等于 1.则函数 $\chi(\xi) = \phi(\xi) - \phi(2\xi)$ 具有所要求的性质.

现在我们定义

\begin{equation}
\widehat{P_k f}(\xi) = \chi(\xi/2^k)\hat{f}(\xi), \label{eq:c6-215}
\end{equation}

或者, 在物理空间中,

\begin{equation}
P_k f = f_k = m_k * f, \label{eq:c6-216}
\end{equation}

其中 $m_k(x) = 2^{nk} m(2^k x)$, 而 $m(x)$ 是 $\chi$ 的逆 傅里叶变换.显然, 由 \eqref{eq:c6-214},

\begin{equation}
f = \sum_{k\in\mathbb{Z}} P_k f, \label{eq:c6-217}
\end{equation}

正是所要求的.把"乘以 $2^{-k}$"的映射记作 $2^{-k}$, 我们还有重要的\newterm{scaling-identity}{伸缩恒等式}{scaling identity}

\begin{equation}
(P_k f) \circ 2^{-k} = P_0(f \circ 2^{-k}).\label{eq:c6-218}
\end{equation}

注意到 $P_k f$ 的 傅里叶变换的支集在二进区间 $2^{k-1} \leq |\xi| \leq 2^{k+1}$ 中, 因此

\[
P_{k'} P_k f = 0, \quad \forall\, k, k'\in\mathbb{Z}, \quad |k-k'| > 2.
\]

于是

\[
P_k f = \sum_{k'\in\mathbb{Z}} P_{k'}(P_k f) = \sum_{|k-k'|\leq 1} P_{k'} P_k f.
\]

从而, 由于 $P_{k-1}, P_k, P_{k+1}$ 彼此相差不大, 我们可以写

\[
P_k = \sum_{|k-k'|\leq 1} P_{k'} P_k \approx P_k^2.
\]

正因如此, \newterm{cutoff-operator}{截断算子}{cut-off operator} $P_k$ 才被(不严格地)称为 LP 投影.

对所有区间 $J \subset \mathbb{Z}$, 记 $P_J = \sum_{k\in J} P_k$.特别地, 我们写 $P_{\leq k} = P_{(-\infty, k]}$, $P_{<k} = P_{\leq k-1}$.显然, $P_k = P_{\leq k} - P_{<k}$.

下列 LP 投影的性质处于经典 LP 理论的核心地位:

\begin{theorem}\label{thm:c6-3.1}
LP 投影具有下列性质:

LP 1. \newterm{almost-orthogonality}{几乎正交性}{almost orthogonality}.算子 $P_k$ 是自伴的, 并且对一切满足 $|k_1 - k_2| \geq 2$ 的整数对, 有 $P_{k_1} P_{k_2} = 0$.特别地,

\begin{equation}
\|f\|_{L^2}^2 \approx \sum_k \|P_k f\|_{L^2}^2.\label{eq:c6-219}
\end{equation}

LP 2. $L^p$ 有界性:对任意 $1 \leq p \leq \infty$ 与任意区间 $J \subset \mathbb{Z}$,

\begin{equation}
\|P_J f\|_{L^p} \lesssim \|f\|_{L^p}.\label{eq:c6-220}
\end{equation}

LP 3. \newterm{finite-band-property}{有限带宽性质}{finite band property}.任何偏导数 $\partial P_k f$ 都可以写成 $\partial P_k f = 2^k \tilde{P}_k f$ 的形式, 而 $\tilde{P}_k$ 的象征是一个满足性质 LP2 的截断算子\footnote{它与一个稍不同的检验函数 $\tilde{\chi}$ 相关联, 该函数仍保持支集在 $\frac{1}{2} \leq |\xi| \leq 2$ 中, 但可能不满足 \eqref{eq:c6-214}.}.特别地, 对任意 $1 \leq p \leq \infty$,

\begin{equation}
\|\partial P_k f\|_{L^p} \lesssim 2^k \|f\|_{L^p}, \label{eq:c6-221}
\end{equation}

\begin{equation}
2^k \|P_k f\|_{L^p} \lesssim \|\partial f\|_{L^p}.\label{eq:c6-222}
\end{equation}

LP 4. Bernstein 不等式.对任意 $1 \leq p \leq q \leq \infty$, 我们有 Bernstein 不等式

\begin{equation}
\|P_k f\|_{L^q} \lesssim 2^{kn(1/p-1/q)} \|f\|_{L^p}, \quad \forall\, k\in\mathbb{Z}, \label{eq:c6-223}
\end{equation}

\begin{equation}
\|P_{\leq 0} f\|_{L^q} \lesssim \|f\|_{L^p}.\label{eq:c6-224}
\end{equation}

特别地,

\[
\|P_k f\|_{L^\infty} \lesssim 2^{kn/p} \|f\|_{L^p}.
\]

LP5. \newterm{commutator}{换位子}{commutator}估计.考虑换位子

\[
[P_k, f] \cdot g = P_k(f \cdot g) - f \cdot P_k g,
\]

其中 $f, g \in \mathcal{C}_0^\infty(\mathbb{R}^n)$.我们有

\[
\|[P_k, f] \cdot g\|_{L^p} \lesssim 2^{-k} \|\nabla f\|_{L^\infty} \|g\|_{L^p}.
\]

LP6. \newterm{square-function}{平方函数}{square function}不等式.设 $Sf$ 为向量值函数 $\mathbf{S}f = (P_k f)_{k\in\mathbb{Z}}$.量

\begin{equation}
Sf(x) = |\mathbf{S}f(x)| = \big(\sum_{k\in\mathbb{Z}} |P_k f(x)|^2\big)^{1/2}\label{eq:c6-225}
\end{equation}

称为 Littlewood-Paley 平方函数.对每个 $1 < p < \infty$, 存在依赖于 $p$ 的常数, 使得对所有 $f \in \mathcal{C}_0^\infty$,

\begin{equation}
\|f\|_{L^p} \lesssim \|Sf\|_{L^p} \lesssim \|f\|_{L^p}.\label{eq:c6-226}
\end{equation}
\end{theorem}
\begin{proof}
只有 LP6 的证明不是直截了当的, 我们把它推迟到下一节. LP1 的证明是直接的.事实上, 我们只需验证 \eqref{eq:c6-219}.显然,

\[
\begin{aligned}
\|f\|_{L^2}^2 &= \|\sum_k P_k f\|_{L^2}^2 = \sum_{|k-k'|\leq 1} \langle P_k f, P_{k'} f\rangle_{L^2}\\
&\leq \sum_{|k-k'|\leq 1} \|P_k f\|_{L^2} \|P_{k'} f\|_{L^2}\\
&\lesssim \sum_k \|P_k f\|_{L^2}^2
\end{aligned}
\]

要证明 $\sum_k \|P_k f\|_{L^2}^2 \lesssim \|f\|_{L^2}^2$, 我们只需利用 Parseval 恒等式以及投影 $P_k$ 的定义.

证明 LP2 只需对形如 $J = (-\infty, k] \subset \mathbb{Z}$ 的区间进行, 即证明 $P_{\leq k}$ 的 $L^p$ 有界性.若 $\chi(\xi) = \phi(\xi) - \phi(2\xi)$, 则 $\widehat{P_{\leq k} f} = \phi(\xi/2^k)\hat{f}(\xi)$.因此

\[
P_{\leq k} f = \bar{m}_k * f,
\]

其中 $\bar{m}_k(x) = 2^{nk} \bar{m}(2^k x)$, 而 $\bar{m}(x)$ 是 $\phi$ 的逆 傅里叶变换.注意到 $\|\bar{m}_k\|_{L^1} = \|\bar{m}\|_{L^1} \lesssim 1$.于是, 利用卷积不等式 \eqref{eq:c6-188},

\[
\|P_{\leq k} f\|_{L^p} \leq \|\bar{m}_k\|_{L^1} \|f\|_{L^p} \lesssim \|f\|_{L^p},
\]

正是所要求的.

为证明 LP3, 我们写 $\partial_i(P_k f) = 2^k (\partial_i m)_k * f$, 其中 $(\partial_i m)_k(x) = 2^{nk} \partial_i m(2^k x)$.显然 $\|(\partial_i m)_k\|_{L^1} = \|\partial_i m\|_{L^1} \lesssim 1$.因此,

\[
\|\partial_i(P_k f)\|_{L^p} \lesssim 2^k \|f\|_{L^p},
\]

这就建立了 \eqref{eq:c6-221}.为证明 \eqref{eq:c6-222}, 我们写 $\hat{f}(\xi) = \sum_{j=1}^n \frac{\xi_j}{i|\xi|^2} \widehat{\partial_{x_j} f}(\xi)$.于是,

\[
2^k \widehat{P_k f}(\xi) = \sum_{j=1}^n 2^k \frac{\xi_j}{i|\xi|^2} \chi(\xi/2^k) \widehat{\partial_{x_j} f}(\xi) = \sum_{j=1}^n 2^k \psi_j(\xi/2^k) \widehat{\partial_{x_j} f}(\xi),
\]

其中 $\psi_j(\xi) = \frac{\xi_j}{i|\xi|^2} \chi(\xi)$.于是, 在物理空间中,

\[
2^k P_k f = \sum_{j=1}^n ({}^j\bar{m})_k * \partial_j f,
\]

其中 $({}^j\bar{m})_k(x) = 2^{nk} \cdot {}^j\bar{m}(2^k x)$, 而 ${}^j\bar{m}$ 是 $\psi_j$ 的逆 傅里叶变换.于是, 如前,

\[
2^k \|P_k f\|_{L^p} \lesssim \sum_{j=1}^n \|\partial_j f\|_{L^p} = \|\partial f\|_{L^p},
\]

正是所要求的.

性质 LP4 是物理空间表示 \eqref{eq:c6-216} 与卷积不等式 \eqref{eq:c6-188} 的直接推论:

\[
\|P_k f\|_{L^q} = \|m_k * f\|_{L^q} \lesssim \|m_k\|_{L^r} \|f\|_{L^p},
\]

其中 $1 + q^{-1} = r^{-1} + p^{-1}$.现在,

\[
\|m_k\|_{L^r} = 2^{nk} \big(\int_{\mathbb{R}^n} |m(2^k x)|^r\, dx\big)^{1/r} = 2^{nk} 2^{-nk/r} \|m\|_{L^r} \lesssim 2^{nk(1-1/r)} \lesssim 2^{nk(1/p-1/q)}.
\]

只剩下证明 LP5.由 \eqref{eq:c6-216}, 我们可以写

\[
P_k(fg)(x) - f(x) P_k g(x) = \int_{\mathbb{R}^n} m_k(x-y)\big(f(y) - f(x)\big) g(y)\, dy.
\]

另一方面,

\[
\begin{aligned}
|f(y) - f(x)| &\lesssim \big|\int_0^1 \frac{d}{ds} f(x + s(y-x))\, ds\big|\\
&\lesssim |x-y| \|\partial f\|_{L^\infty}
\end{aligned}
\]

因此,

\[
|P_k(fg)(x) - f(x) P_k g(x)| \lesssim 2^{-k} \|\partial f\|_{L^\infty} \int_{\mathbb{R}^n} |\bar{m}_k(x-y)|\, |g(y)|\, dy,
\]

其中 $\bar{m}_k(x) = 2^{nk} \bar{m}(2^k x)$, $\bar{m}(x) = |x| m(x)$.于是,

\[
\|P_k(fg) - f P_k g\|_{L^p} \lesssim 2^{-k} \|\partial f\|_{L^\infty} \|g\|_{L^p}.
\]

我们把性质 LP6 的证明留到下一节.
\end{proof}

\begin{remark*}
在前面的证明中, 如果只证明性质 LP2--4 在 $k = 0$ 的情形, 再通过伸缩恒等式 \eqref{eq:c6-218} 推出更一般的估计, 事情会更简单.特别地, 注意 Bernstein 不等式无非是说:当 $f$ 在频率空间局部化时(与相反的情形——即 $f$ 在物理空间局部化时——相对), 较低的 $L^p$ 范数控制较高的 $L^p$ 范数.这与我们对 $L^p$ 范数的直觉一致:虽然频率局部化的函数在物理空间的无穷远处可能太大而不可能可积, 但我们不必担心函数在局部爆炸的突然跳跃或尖峰, 因此只需控制前一种现象.
\end{remark*}

\begin{definition*}
我们说一个乘子算子 $\tilde{P}_k$ 与标准的 LP 投影 $P_k$ 相似, 如果它的象征 $\tilde{\chi}_k$ 是适应二进区域 $|\xi| \sim 2^k$ 的隆起函数.更精确地说, 对某个支集在区域 $c^{-1} 2^k \lesssim |\xi| \leq c 2^k$ 中的隆起函数 $\tilde{\chi}$, 其中 $c > 0$ 固定, 我们可以写 $\tilde{\chi}_k(\xi) = \tilde{\chi}\big(\frac{\xi}{2^k}\big)$.
\end{definition*}

\begin{remark*}
注意到不等式 $\|P_k f\|_{L^p} \lesssim \|f\|_{L^p}$ 对每个与 $P_k$ 相似的算子 $\tilde{P}_k$ 都成立.性质 LP3、LP4 与 LP5 亦然.
\end{remark*}

\begin{remark*}
关于算子 $\tilde{P}_k$ 与极大函数, 我们有如下的逐点关系:

\begin{equation}
|\tilde{P}_{\leq k} f| \lesssim \mathcal{M} f(x).\label{eq:c6-227}
\end{equation}

事实上, 如前,

\[
\tilde{P}_{\leq k} f = \tilde{m}_k * f,
\]

其中 $\tilde{m}_k(x) = 2^{nk} \tilde{m}(2^k x)$, $\tilde{m}(x) \in \mathcal{S}(\mathbb{R}^n)$.因此,

\[
\begin{aligned}
|\tilde{P}_{\leq k} f| &\lesssim 2^{nk} \int |f(y)| \tilde{m}\big(2^k(x-y)\big)\, dy \lesssim 2^{nk} \int |f(y)| (1 + 2^k |x-y|)^{-n-1}\, dy\\
&\lesssim 2^{nk} \int_{B(x, 2^{-k})} |f(y)| (1 + 2^k |x-y|)^{-n-1}\, dy\\
&\quad + 2^{nk} \sum_{j=0}^\infty \int_{2^j \leq 2^k |x-y| \leq 2^{j+1}} |f(y)| (1 + 2^k |x-y|)^{-n-1}\, dy\\
&\lesssim 2^{nk} \big(\int_{B(x, 2^{-k})} |f(y)|\, dy + \sum_{j \geq 0} 2^{-(n+1)j} \int_{|x-y| \leq 2^{j+1-k}} |f(y)|\, dy\big)\\
&\lesssim \mathcal{M} f(x) + \sum_{j > 0} 2^{-(n+1)j} 2^{nk} 2^{n(j+1-k)} \frac{1}{|B(x, 2^{-k+j+1})|} \int_{B(x, 2^{-k+j+1})} |f(y)|\, dy\\
&\lesssim \mathcal{M} f(x) + 2^n \sum_{j > 0} 2^{-j} \mathcal{M} f(x) \lesssim \mathcal{M} f(x)
\end{aligned}
\]

正是所要求的.
\end{remark*}

性质 LP3--LP4 在很大程度上解释了为什么 LP 理论是偏微分方程如此有用的工具.有限带宽性质使我们可以把二进分量 $f_k$ 的导数换成与 $2^k$ 相乘.$L^p \to L^\infty$ Bernstein 不等式是 Sobolev 空间 $W^{\frac{n}{p}, p}(\mathbb{R}^n)$ 到 $L^\infty(\mathbb{R}^n)$ 的嵌入失效的一种二进补救.事实上, 鉴于有限带宽性质, Bernstein 不等式确实对每个 LP 分量 $f_k$ 蕴含所要求的 Sobolev 不等式;对 $f$ 而言 Sobolev 不等式失效的原因在于求和 $f = \sum_k f_k$.

下面我们给出 LP 演算的几个应用.

\mtsubsec{插值不等式}

下列不等式对 $\mathcal{C}_0^\infty(\mathbb{R}^n)$ 中的任意函数和任意整数 $0 \leq i \leq m$ 成立:

\begin{equation}
\|\partial^i f\|_{L^p} \lesssim \|f\|_{L^p}^{1-i/m} \|\partial^m f\|_{L^p}^{i/m}.\label{eq:c6-228}
\end{equation}

为证明它, 我们分解 $f = P_{\leq k} f + P_{> k} f = f_{\leq k} + f_{> k}$.现在, 利用 LP2--LP4, 对任意固定的 $k \in \mathbb{Z}$,

\[
\begin{aligned}
\|\partial^i f\|_{L^p} &\leq \|\partial^i f_{\leq k}\|_{L^p} + \|\partial^i f_{> k}\|_{L^p}\\
&\leq 2^{ki} \|f\|_{L^p} + 2^{k(i-m)} \|\partial^m f\|_{L^p}
\end{aligned}
\]

因此,

\[
\|\partial^i f\|_{L^p} \leq \lambda^i \|f\|_{L^p} + \lambda^{i-m} \|\partial^m f\|_{L^p}
\]

对任意 $\lambda \in 2^{\mathbb{Z}}$.为了完成证明, 我们希望选取 $\lambda$, 使得右端两项相等, 即

\[
\lambda_0 = \left(\frac{\|\partial^m f\|_{L^p}}{\|f\|_{L^p}}\right)^{1/m}.
\]

由于我们限于 $\lambda \in 2^{\mathbb{Z}}$, 我们选取二进数 $\lambda \in 2^{\mathbb{Z}}$, 使得 $\lambda \leq \lambda_0 \leq 2\lambda$.于是,

\[
\|\partial^i f\|_{L^p} \leq \lambda_0^i \|f\|_{L^p} + \left(\frac{2}{\lambda_0}\right)^{m-i} \|\partial^m f\|_{L^p} \lesssim \|f\|_{L^p}^{1-i/m} \|\partial^m f\|_{L^p}^{i/m}.
\]

一般地, 当关于 $\mathbb{R}^n$ 上函数的某个估计在量纲上不一致时(即两边的伸缩不相同), 这样的估计可以被"放大"为一个看似更强的估计(或者被证明是错误的).

\textbf{习题.} 假定不等式

\[
\|\partial^i f\|_{L^p} \lesssim \|f\|_{L^p} + \|\partial^m f\|_{L^p},
\]

通过考虑伸缩 $f \mapsto f \circ \lambda$ 推导估计 \eqref{eq:c6-228}.

\mtsubsec{非临界 Sobolev 不等式}

我们将证明不等式 \eqref{eq:c6-201} 的一个稍加改进的版本:对函数 $f \in C_0^\infty(\mathbb{R}^n)$ 与指数 $1 \leq p < q < \infty$, $1/p - m/n < 1/q$,

\[
\|f\|_{L^q(\mathbb{R}^n)} \lesssim \|f\|_{L^p} + \|\partial^m f\|_{L^p}.
\]

我们分解 $f = P_{\leq 0} f + \sum_{k\in\mathbb{N}} P_k f = f_{< 0} + \sum_{k > 0} f_k$.于是, 先用 LP4, 再用 LP3,

\[
\begin{aligned}
\|f\|_{L^q} &\leq \|f_{< 0}\|_{L^q} + \sum_{k > 0} \|f_k\|_{L^q} \lesssim \|f\|_{L^p} + \sum_{k > 0} 2^{kn(1/p-1/q)} \|f\|_{L^p}\\
&\lesssim \|f\|_{L^p} + \sum_{k > 0} 2^{kn(m/n-\varepsilon)} \|f\|_{L^p} \lesssim \|f\|_{L^p} + \sum_{k > 0} 2^{-kn\varepsilon} \|\partial^m f\|_{L^p}\\
&\lesssim \|f\|_{L^p} + \|\partial^m f\|_{L^p}
\end{aligned}
\]
\mtsubsec{函数空间}

Littlewood-Paley 理论既可以用来给出 Sobolev 空间的等价描述, 也可以引入新的、更精细的函数空间.我们首先注意到, 鉴于几乎正交性性质 LP1,

\[
\begin{aligned}
\|f\|_{L^2}^2 &= \|\sum_{k\in\mathbb{Z}} P_k f\|_{L^2}^2 \lesssim \sum_{k\in\mathbb{Z}} \|P_k f\|_{L^2}^2\\
\sum_{k\in\mathbb{Z}} \|P_k f\|_{L^2}^2 &\lesssim \|f\|_{L^2}
\end{aligned}
\]

于是我们可以给出齐次 Sobolev 范数 $\|\cdot\|_{\dot{H}^s(\mathbb{R}^n)}$ 的一个 LP 描述:

\begin{equation}
\|f\|_{\dot{H}^s}^2 \approx \sum_{k\in\mathbb{Z}} 2^{2ks} \|P_k f\|_{L^2}^2.\label{eq:c6-229}
\end{equation}

对 $k \in \mathbb{Z}^+$, 定义算子 $\Delta_k = P_k$(若 $k > 0$), $\Delta_0 = P_{\leq 0}$.对 $H^s$ 范数, 还有

\begin{equation}
\|f\|_{H^s}^2 \approx \sum_{k=0}^\infty 2^{2ks} \|\Delta_k f\|_{L^2}^2.\label{eq:c6-230}
\end{equation}

Littlewood-Paley 分解可以用来定义新的函数空间, 如\newterm{besov-space}{Besov 空间}{Besov space}.

\begin{definition*}
Besov 空间 $B_{p, q}^s(\mathbb{R}^n)$ 是 $\mathcal{C}_0^\infty(\mathbb{R}^n)$ 关于范数

\begin{equation}
\|f\|_{B_{p, q}^s} = \big(\sum_{k=0}^\infty 2^{ksq} \|\Delta_k f\|_{L^p}^q\big)^{1/q}\label{eq:c6-231}
\end{equation}

的闭包.
\end{definition*}

相应的齐次 Besov 范数定义为

\begin{equation}
\|f\|_{\dot{B}_{p, q}^s} = \big(\sum_{k\in\mathbb{Z}} 2^{sqk} \|P_k f\|_{L^p}^q\big)^{1/q}, \label{eq:c6-232}
\end{equation}

类似地, 把 \eqref{eq:c6-231} 中的 $L^p$ 范数与 $l^q$ 范数对调, 可以定义\newterm{triebel-space}{Triebel 空间}{Triebel space} $F_{p, q}^s$.于是, 例如, $H^s$ 范数与 Besov 范数 $B_{2, 2}^s$ 等价.注意到 $H^s \subset B_{2, 1}^s$.更大的空间 $B_{2, 1}^s$ 之所以有用, 原因之一是下列事实:

\begin{equation}
\|f\|_{L^\infty} \lesssim \|f\|_{\dot{B}_{2, 1}^{n/2}}, \label{eq:c6-233}
\end{equation}

它由 Bernstein 不等式 LP4 推出.\eqref{eq:c6-233} 将在下一节中起关键作用.使用 Besov 范数 $B_{2, 1}^s$ 的另一个原因将在下一节讨论乘积估计时变得清楚.

\mtsubsec{乘积估计}

LP 演算对非线性估计特别有用.设 $f, g$ 是 $\mathbb{R}^n$ 上的两个函数.考虑

\begin{equation}
P_k(fg) = \sum_{k', k''\in\mathbb{Z}} P_k(P_{k'} f\, P_{k''} g).\label{eq:c6-234}
\end{equation}

现在, 由于 $P_{k'} f$ 的 Fourier 支集在集合 $D' = 2^{k'-1} \leq |\xi| \leq 2^{k'+1}$ 中, $P_{k''} f$ 的 Fourier 支集在 $D'' = 2^{k''-1} \leq |\xi| \leq 2^{k''+1}$ 中, 可知 $P_{k'} f\, P_{k''} g$ 的 Fourier 支集在 $D' + D''$ 中.只有当 $D' + D''$ 与 $2^{k-1} \leq |\xi| \leq 2^{k+1}$ 相交时, 和式 \eqref{eq:c6-234} 中才得到非零贡献.因此, 记 $f_k = P_k f$, $f_{< k} = P_{< k} f$, 并对任意区间 $J \subset \mathbb{Z}$ 记 $f_J = P_J f$, 我们得到

\begin{lemma}\label{lem:c6-3.6}
给定函数 $f, g$, 我们有如下分解:

\begin{equation}
\begin{aligned}
P_k(f \cdot g) &= HH_k(f, g) + LL_k(f, g) + LH_k(f, g) + HL_k(f, g)\\
HH_k(f, g) &= \sum_{k', k'' > k+5, \, |k'-k''| \leq 3} P_k\big(f_{k'} \cdot P_{k''} g\big)\\
LL_k(f, g) &= P_k\big(f_{[k-5, k+5]} \cdot g_{[k-5, k+5]}\big)\\
LH_k(f, g) &= P_k\big(f_{\leq k-5} \cdot g_{[k-3, k+3]}\big)\\
HL_k(f, g) &= P_k\big(f_{[k-3, k+3]} \cdot g_{\leq k-5}\big)
\end{aligned}\label{eq:c6-235}
\end{equation}
\end{lemma}

项 $HH_k(f, g)$ 对应高-高相互作用.更精确地说, 定义 $HH_k(f, g)$ 的和式中每一项的频率为 $\sim 2^m$, 其中某个 $2^m \gg 2^k$.我们将示意性地写

\begin{equation}
HH_k(f, g) = P_k\big(\sum_{m > k} f_m \cdot g_m\big).\label{eq:c6-236}
\end{equation}

项 $LL_k(f, g)$ 由有限多项组成, 通常可以忽略.事实上, 在任何估计中, 它们既可以当作有限多个 HH 项, 也可以当作有限多个 LH 与 HL 项来处理.我们示意性地写

\begin{equation}
LL_k(f, g) = 0.\label{eq:c6-237}
\end{equation}

最后, $LH_k$ 与 $HL_k$ 项分别由低-高与高-低相互作用组成.我们示意性地写

\begin{equation}
LH_k(f, g) = P_k\big(f_{< k} \cdot g_k\big), \label{eq:c6-238}
\end{equation}

\begin{equation}
HL_k(f, g) = P_k\big(f_k \cdot g_{< k}\big).\label{eq:c6-239}
\end{equation}

\begin{remark*}
在 \eqref{eq:c6-235} 给出的 $LH_k$ 的精确表达式中, 形如 $f_{\leq k-5} \cdot g_{k''}$, $k'' \in [k-3, k+3]$ 的项, 其 Fourier 支集在二进区域 $\sim 2^k$ 中.因此 $P_k$ 可以安全地忽略, 我们可以写

\[
LH_k(f, g) \sim f_{< k} \cdot g_k.
\]

我们从而建立了著名的\newterm{trichotomy-formula}{三分公式}{trichotomy formula}

\begin{equation}
P_k(f \cdot g) = LH_k(f, g) + HL_k(f, g) + HH_k(f, g), \label{eq:c6-240}
\end{equation}

它是\newterm{paradifferential-calculus}{仿微分演算}{paradifferential calculus}的基础.在实际应用中, 每当我们应用公式 \eqref{eq:c6-240} 时, 都必须记住公式 \eqref{eq:c6-237}--\eqref{eq:c6-239} 只是近似的;精确的定义由 \eqref{eq:c6-235} 给出.然而在任何估计中, 我们都可以安全地忽略额外的项, 因为它们的估计方式与我们保留的项完全一样.
\end{remark*}

我们现在利用三分公式来证明一个\newterm{product-estimate}{乘积估计}{product estimate}.

\begin{theorem}\label{thm:c6-3.7}
下列估计对所有 $s > 0$ 成立:

\begin{equation}
\|fg\|_{H^s} \lesssim \|f\|_{L^\infty} \|g\|_{H^s} + \|g\|_{L^\infty} \|f\|_{H^s}.\label{eq:c6-241}
\end{equation}

因此对所有 $s > n/2$,

\begin{equation}
\|fg\|_{H^s} \lesssim \|f\|_{H^s} \|g\|_{H^s}.\label{eq:c6-242}
\end{equation}
\end{theorem}

下面我们给出\thmref{thm:c6-3.7} 的一个相当简单且很有启发性的证明.这个证明\footnote{我感谢 Igor Rodnianski 为我指出这一论证.}表明, 有时最好不要依赖完整的分解 \eqref{eq:c6-235}, 而是在需要时节省地使用分解.事实上, 我们写

\[
\|fg\|_{\dot{H}^s}^2 \lesssim \sum_k 2^{2ks} \|P_k(fg)\|_{L^2}^2 \lesssim \sum_k 2^{2ks} \|P_k(f_{< k} g)\|_{L^2}^2 + \sum_k 2^{2ks} \|P_k(f_{\geq k} g)\|_{L^2}^2.
\]

现在,

\[
\begin{aligned}
\sum_k 2^{2ks} \|P_k(f_{\geq k} g)\|_{L^2}^2 &\lesssim \|g\|_{L^\infty}^2 \sum_k 2^{2ks} \|f_{\geq k}\|_{L^2}^2\\
&\lesssim \|g\|_{L^\infty}^2 \sum_k \sum_{k' \geq k} 2^{2(k-k')s} \|2^{k's} f_{k'}\|_{L^2}^2\\
&= \|g\|_{L^\infty}^2 \sum_{k'} \big(\sum_{k \leq k'} 2^{2(k-k')s}\big) \|2^{k's} f_{k'}\|_{L^2}^2\\
&\lesssim \|g\|_{L^\infty}^2 \|f\|_{H^s}^2
\end{aligned}
\]

为了估计 $\sum_k 2^{2ks} \|P_k(f_{\leq k} g)\|_{L^2}^2$, 我们将进一步分解, 仿照分解 \eqref{eq:c6-235} 进行.但首先注意到, 项 $\sum_k 2^{2ks} \|P_k(f_{[k-3, k]} g)\|_{L^2}^2$ 可以像 $\sum_k 2^{2ks} \|P_k(f_{> k} g)\|_{L^2}^2$ 一样精确地处理.事实上, 我们还不如直接估计 $\sum_k 2^{2ks} \|P_k(f_{> k-3} g)\|_{L^2}^2$.现在,

\[
\begin{aligned}
P_k(f_{\leq k-3} g) &= \sum_{k'} P_k(f_{\leq k-3} g_{k'}) = \sum_{k' < k-2} P_k(f_{\leq k-3} g_{k'}) + \sum_{k-2 \leq k' \leq k+2} P_k(f_{\leq k-3} g_{k'})\\
&\quad + \sum_{k' > k+2} P_k(f_{\leq k-3} g_{k'})
\end{aligned}
\]

注意到第一项与最后一项为零, 因此,

\[
P_k(f_{\leq k-3} g) = \sum_{k-2 \leq k' \leq k+2} P_k(f_{\leq k-3} g_{k'}) \approx P_k(f_{\leq k-3} g_k).
\]

为简单起见, 我们常常直接写

\begin{equation}
P_k(f_{< k} g) \approx f_{< k} \cdot g_k.\label{eq:c6-243}
\end{equation}

当然, 这个公式并不完全正确, 但在精神上是正确的.现在,

\[
\begin{aligned}
\sum_k 2^{2ks} \|P_k(f_{< k} g)\|_{L^2}^2 &= \sum_k 2^{2ks} \|f_{< k} g_k\|_{L^2}^2\\
&\lesssim \|f\|_{L^\infty}^2 \sum_k 2^{2ks} \|g_k\|_{L^2}^2 = \|f\|_{L^\infty}^2 \|g\|_{\dot{H}^s}^2
\end{aligned}
\]

正是所要求的.

\begin{remark*}
鉴于 \eqref{eq:c6-243}, 我们有如下的部分分解公式

\begin{equation}
P_k(fg) = f_{< k} g_k + P_k(f_{\geq k} g) = LH_k(f, g) + P_k(f_{\geq k} g), \label{eq:c6-244}
\end{equation}

把它与完整的三分分解 \eqref{eq:c6-240} 对比.
\end{remark*}

类似的估计在 Besov 空间中成立, 而且更容易证明.事实上, 对每个 $s > 0$, 我们有

\begin{equation}
\|fg\|_{B_{2, 1}^s} \lesssim \|f\|_{L^\infty} \|g\|_{B_{2, 1}^s} + \|g\|_{L^\infty} \|f\|_{B_{2, 1}^s}.\label{eq:c6-245}
\end{equation}

\textbf{习题.} 证明估计 \eqref{eq:c6-245}.

\section{Wente 不等式}

在本节中, 我们作为 Littlewood-Paley 理论的一个应用来证明\newterm{wentes-inequality}{Wente 不等式}{Wente's inequality}.下面给定 $\mathbb{R}^2$ 中的两个函数 $f, g$, 我们考虑双线性表达式 $(df \wedge dg)^* = \partial_x f\, \partial_y g - \partial_y f\, \partial_x g$, 其中 $*$ 表示 $\mathbb{R}^2$ 中平凡的\newterm{hodge-duality}{Hodge 对偶}{Hodge duality}.为了语言上的方便, 下面我们略去对偶记号, 直接写 $df \wedge dg$, 其中 $\wedge$ 表示\newterm{wedge-product}{楔积}{wedge product}.

\begin{theorem}\label{thm:c6-4.1}
在 $\mathbb{R}^2$ 上, 设 $f, g \in H^1(\mathbb{R}^2)$, $\Delta u = (df \wedge dg)$.则 $u \in L^\infty$, 事实上是连续的.
\end{theorem}

\begin{remark*}
事实上, 若把 $\wedge$ 换成普通乘法, 我们至多能得到 $df \cdot dg \in L^1$.这显然不足以得到 $u \in L^\infty$.然而, 事实证明 $df \wedge dg$ 具有特殊的结构, 使我们能够导出所要求的估计.在上述定理中, 我们指的是通过\newterm{canonical-solution}{典范解}{canonical solution}得到的典范解 $u = \Delta^{-1}(df \wedge dg)$.
\end{remark*}

\begin{proof}
由有限带宽性质容易看出, $\Delta$ 是从 $\dot{B}_{p, 1}^s$ 到 $\dot{B}_{p, 1}^{s-2}$ 的等距算子.事实上, 我们将取 $p = 2$.鉴于 Sobolev 不等式 \eqref{eq:c6-233}, 只需证明 $df \wedge dg \in \dot{B}_{2, 1}^{-1}(\mathbb{R}^2)$.利用三分公式以及 LP 投影 $P_k$ 与 $d$ 交换这一事实, 我们写

\[
\begin{aligned}
I &= df \wedge dg = LH_k + HL_k + HH_k\\
LH_k &= dP_{< k} f \wedge dP_k g\\
HL_k &= dP_k \wedge dP_{< k} g\\
HH_k &= P_k\big(\sum_{m \geq k} (dP_m f \wedge dP_m g)\big)
\end{aligned}
\]

由对称性, 我们只需处理 LH 与 HH.$LH$ 项不用楔积的特殊结构就能平凡地估计.利用 Bernstein 不等式, 我们写

\[
\begin{aligned}
2^{-k} \|LH_k\|_{L^2} &\lesssim 2^{-k} \sum_{l < k} \|dP_l f\|_{L^\infty} \|dP_k(g)\|_{L^2}\\
&\lesssim \sum_{l < k} 2^{l-k} \|DP_l f\|_{L^2} \|DP_k f\|_{L^2}
\end{aligned}
\]

证明接下来用 Young 不等式的如下离散版本.

\begin{lemma}\label{lem:c6-4.2}
设 $f(k) \in l^1(\mathbb{Z})$, $g(k), h(k) \in l^2(\mathbb{Z})$.则

\[
\sum_{k, l} f(k-l)\, g(l)\, h(k) \leq \|f\|_{l^1} \|g\|_{L^2} \|h\|_{l^2}.
\]
\end{lemma}

利用该引理, 我们导出

\[
\begin{aligned}
\sum_k 2^{-k} \|LH_k\|_{L^2} &\lesssim \big(\sum_l \|DP_l f\|_{L^2}^2\big)^{1/2} \big(\sum_k \|DP_k f\|_{L^2}^2\big)^{1/2}\\
&\lesssim \|Df\|_{L^2} \|Dg\|_{L^2}
\end{aligned}
\]

现在我们考虑 $HH_k$.正是在这里我们需要利用楔积的特殊结构.事实上, 我们只需利用恒等式 $df \wedge dg = d(f \wedge dg)$.于是,

\[
\begin{aligned}
HH_k &= \sum_{m \geq k} P_k(dP_m f \wedge dP_m g)\\
&= \sum_{m \geq k} dP_k(P_m f \wedge dP_m g)
\end{aligned}
\]

于是, 利用有限带宽性质与 Bernstein 不等式,

\[
\begin{aligned}
\|HH_k\|_{L^2} &\lesssim 2^{2k} \|P_m f \wedge dP_m g\|_{L^1}\\
&\lesssim 2^{2k} \|P_m f\|_{L^2} \|DP_m g\|_{L^2}\\
&\lesssim 2^{2k-m} \|DP_m f\|_{L^2} \|DP_m g\|_{L^2}
\end{aligned}
\]

因此,

\[
2^{-k} \|HH_k\|_{L^2} \lesssim 2^{k-m} \|DP_m f\|_{L^2} \|DP_m g\|_{L^2}.
\]

于是, 再次利用上面引理中的离散 Young 不等式,

\[
\sum_k 2^{-k} \|LH_k\|_{L^2} \lesssim \|Df\|_{L^2} \|Dg\|_{L^2},
\]

正是所要求的.
\end{proof}

% ==== part06b 第 6 章:第 5-8 节 ====
% 源: MinerU_latex_lecture nots.tex 9904-12106(第 5 节起;pdf_text.txt 8153-9833,PDF 第 150-179 页)
% 术语: 新术语(中文=英文):
%   二进立方体=dyadic cube; 父立方体=parent cube; 子立方体=child cube;
%   消没条件=cancellation condition; 平方函数=square function;
%   Rademacher 函数=Rademacher functions; 反向 Hölder 不等式=reverse Hölder inequality;
%   反向 Minkowski 不等式=reverse Minkowski inequality; 主曲率=principal curvature;
%   Stein 算子=Stein operator; Knapp 反例=Knapp counterexample;
%   演化算子=evolution operator; 双线性形式=bilinear form; 三线性形式=trilinear form;
%   混合范数=mixed norm; 相空间=phase space; 投影不相交=projectionally disjoint;
%   Zygmund 类=Zygmund class; 解析插值=analytic interpolation; 帽=cap;
%   延拓算子=extension operator; Riesz 位势=Riesz potentials; 主值分布=principal value distribution
% REFS: 跨章引用(正文暂用纯文本,待统一解析):
%   - 第 6 章第 5 节 "Corollary 2.7"(Lebesgue 微分定理)→ 目标: 第 2 章推论 2.7
%   - 第 6 章第 5 节 "Theorem 2.4"(极大函数弱 L1 性质)→ 目标: 第 2 章定理 2.4
%   - 第 6 章第 8 节 "Stein interpolation theorem, Thm. 1.8" → 目标: 第 1 章定理 1.8【已核对: 新 OCR 确认原书引 Theorem 1.8(见 md 3211/4839 行), 但译本第 1 章(part01)无此编号, 待人工定位标签】
%   - 第 6 章第 8 节 "the mixed norm notation defined in section 7"(疑为"第 5 节"之误)→ 目标: 第 6 章第 5 节
%   - 第 6 章第 5 节 "the Trichotomy Formula"(三分公式,第 6 章第 4 节 (240) 附近)→ 目标: 第 6 章第 4 节
%   - "性质 LP1-LP6"(Littlewood-Paley 投影性质,第 5 章/第 6 章前部定义)→ 目标: part06a 相关部分
%   - 例 2 中 "see (??)" 为原书破损引用(Laplace 方程基本解公式)→ 目标: 第 1 章对应公式
%   - 文献引用保持原样: [Kl-Rodn3], [Ge-S], [K-T], [?]
% FOOTNOTE: 脚注 7(例 2)、8 与 9(定理 6.11 证明)、10(问题 3)、11(问题 6)已按 pdf_text.txt 对应页补全并翻译.
% 编号说明:
%   - 公式 (246)-(306) 顺次自动编号.(254a)-(254e) 为原书同一编号下的五个子式,无法用自动编号
%     表达,采用五个 \[ \] 显示内 \tag{254a}...\tag{254e}(\[ \] 不推进计数器;指南"禁止
%     \tag"在本特例下无法满足,否则后续编号整体偏移,同 partgf1 先例);五式之后用
%     \stepcounter{equation} 手动推进一次(占去编号 254),使下一式从 (255) 开始.
%   - 第 7 节问题 1-8 原书显示为 "Problem 1."..."Problem 8."(无节前缀,与共享计数器显示的
%     "7.x" 不同),故用加粗段落 \textbf{问题 N.} 手工编号,不使用 problem 环境.
%   - 第 8 节 "8.13. First proof: analytic interpolation" 正文 "According to Remark 8.12 and
%     Remark 8.4" 照原样保留(8.12 实为小节,原书如此写法).
%   - 引理 6.7 证明中 dist(Q_x, δΩ)、dist(Q*, δΩ) 的 δΩ 按 ∂Ω 处理(与引理陈述一致,原书笔误).
%   - n=2 情形 "||Sf||³_{L⁶(R²)} = ... ≃ A^{1/2}||f||³_{L²(S²)}" 中 S² 疑为 S¹ 之误,照原书保留.
%   - 命题 6.18 源文件/原书均有(任务清单未列),照译.
% ====

\section{一个尖锐的迹定理}

本节我们给出 LP 理论的另一个应用:Besov 空间中的迹定理的一个更强版本, 见 \cite{kl-rodn3}.

为简单起见, 设 $I = [0, 1]$, 考虑 $I \times \mathbb{R}^2$.我们将使用\newterm{mixed-norm}{混合范数}{mixed norm}记号:
\begin{align*}
\|f\|_{L^q_t L^p_x} &= \left(\int_0^1 \|f(t, \cdot)\|^q_{L^p_x(\mathbb{R}^2)} \mathrm{d}t\right)^{\frac{1}{q}} \\
\|f\|_{L^p_x L^q_t} &= \left(\int_{\mathbb{R}^2} \|f(\cdot, x)\|^p_{L^q_t(I)} \mathrm{d}x\right)^{\frac{1}{p}}
\end{align*}
当 $p = \infty$ 或 $q = \infty$ 时作显然的修改.

我们将得到下述迹型估计:
\begin{equation}
\left\| \int_I |\partial_t f|^2 \mathrm{d}t \right\|_{B_{2, 1}^{1}} \lesssim \|f\|^2_{H^2(I \times \mathbb{R}^2)}\label{eq:c6-246}
\end{equation}

我们注意到
\[
\|g\|_{B_{2, 1}^{1}} \lesssim \|\nabla g\|_{B_{2, 1}^{0}} + \|g\|_{L^2}
\]

因此, \eqref{eq:c6-246} 由下面的"尖锐双线性迹"定理推出.

\begin{theorem}\label{thm:c6-5.1}
对 $I \times \mathbb{R}^2$ 上任意光滑标量函数 $g, h$, 有
\begin{equation}
\left\| \int_I \partial_t g \cdot h\, \mathrm{d}t \right\|_{B_{2, 1}^{0}} \lesssim \|g\|_{H^1(I \times \mathbb{R}^2)} \cdot \|h\|_{H^1(I \times \mathbb{R}^2)}\label{eq:c6-247}
\end{equation}
\end{theorem}

\begin{proof}
立刻可见:
\[
\begin{array}{rcl}
\left\| \int_I \partial_t g \cdot h\, \mathrm{d}t \right\|_{B_{2, 1}^{0}} &=& \sum_{k \geq 0} \left\| P_k \int_0^1 \partial_t g \cdot h\, \mathrm{d}t \right\|_{L_x^2} + \left\| P_{<0} \int_0^1 \partial_t g \cdot h\, \mathrm{d}t \right\|_{L_x^2} \\
&\lesssim& \sum_{k \geq 0} \left\| P_k \int_0^1 \partial_t g \cdot h\, \mathrm{d}t \right\|_{L_x^2}
\end{array}
\]

然后我们将 $g$ 与 $h$ 按 $x$ 分解:$g = \sum_k P_k g = \sum_k g_k$, $h = \sum_k P_k h = \sum_k h_k$.于是可以把 $P_k \int_0^1 (\partial_t g \cdot h)$ 分解为 $A_k + B_k + C_k + D_k$, 其中
\begin{align*}
A_k &= P_k \int_0^1 (\partial_t g)_{<k} \cdot h_{\geq k} \\
B_k &= P_k \int_0^1 (\partial_t g)_{\geq k} \cdot h_{<k} \\
C_k &= P_k \int_0^1 (\partial_t g)_{<k} \cdot h_{<k} \\
D_k &= P_k \int_0^1 (\partial_t g)_{\geq k} \cdot h_{\geq k}
\end{align*}

如同三分公式中那样, $C_k$ 本质上为零(有限个例外项可以归入 $A_k$、$B_k$ 或 $D_k$).

我们现在简要勾勒如何估计 $A_k$、$B_k$、$D_k$, 细节留待补全.注意 $P_k$ 平凡地与积分 $\int_0^1 \mathrm{d}t$ 以及任何偏导数 $\partial_t$ 交换.

为估计 $A_k$, 注意到我们可以写出(利用 LP2):
\[
\|A_k\|_{L_x^2} \lesssim \sum_{k' < k \leq k''} \int_0^1 \|(\partial_t g)_{k'} \cdot h_{k''}\|_{L_x^2} \mathrm{d}t
\]

然后我们可以用 Bernstein 不等式 LP4 与作用在 $h$ 上的性质 LP3 提取出幂 $2^{k'-k''}$.把 $2^{k'-k''}$ 写成 $2^{(k'-k)/2 + (k-k'')/2}$, 利用 LP1, 并对 $k$ 求和, 可以得到:
\[
\sum_{k \geq 0} \|A_k\|_{L_t^\infty L_x^2} \lesssim \|\partial_t g\|_{L_t^\infty L_x^2} \cdot \|\nabla h\|_{L_t^\infty L_x^2}
\]

为估计 $D_k = P_k \int_0^1 (\partial_t g)_{\geq k} \cdot h_{\geq k}$, 写
\[
D_k = D_k^1 + D_k^2 = \sum_{k \leq k' \leq k''} P_k \int_0^1 (\partial_t g)_{k'} \cdot h_{k''} + \sum_{k \leq k' \leq k''} P_k \int_0^1 (\partial_t g)_{k''} \cdot h_{k'}
\]

$D_k^1$ 可以直接估计, 无需分部积分.用 LP4 与 LP3 写出
\[
\|D_k^1\|_{L_x^2} \lesssim 2^{k-k'} \|\partial_t g\|_{L_t^2 L_x^2} \cdot \|\nabla h\|_{L_t^2 L_x^2}
\]
然后对 $k$ 求和并利用 LP1 得到:
\[
\sum_{k \geq 0} \|D_k^1\|_{L_x^2} \lesssim \|\partial_t g\|_{L_t^2 L_x^2} \cdot \|\nabla h\|_{L_t^2 L_x^2}
\]

为估计 $D_k^2$, 我们用分部积分把 $\partial_t$ 从高频的 $g_{k''}$ 转移到低频的 $h_{k'}$.分部积分之后, 把所得结果与 $D_k^1$ 完全同样地处理.因此, 我们只需估计边界项:$\|I_k(1) - I_k(0)\|_{L_x^2} \lesssim \|I_k\|_{L_t^\infty L_x^2}$, 其中
\[
I_k = \sum_{k \leq k' < k''} P_k (g_{k''} \cdot h_{k'})
\]

我们用下面的引理来完成:

\begin{lemma}\label{lem:c6-5.2}
对任意 $k, k', k''$, 有
\[
\|P_k (g_{k'} \cdot h_{k''})\| \lesssim 2^{-\frac{1}{4}(|k'-k| + |k''-k|)} \|g_{k'}\| \|h_{k''}\|
\]
\end{lemma}

利用此引理, 我们分部积分并像 $D_k^1$ 那样估计 $D_k^2$, 再加上边界项, 最终得到:
\[
\sum_k \|D_k^2\|_{L_x^2} \lesssim \|g\|_{H^1} \cdot \|h\|_{H^1}
\]

现在我们类似地分解来估计 $B_k$:$B_k = \sum_{k' < k \leq k''} P_k \int_0^1 (\partial_t g)_{k''} \cdot h_{k'}$.同上, 我们分部积分并用引理估计边界项 $J_k = \sum_{k' < k < k''} P_k (g_{k''}) \cdot h_{k'})$.然后不难经过整理并对 $k$ 求和得到
\[
\sum_k \|B_k\|_{L_x^2} \lesssim \|g\|_{H^1} \cdot \|h\|_{H^1}
\]

综合对 $A_k$、$B_k$ 与 $D_k$ 的所有估计, 即完成定理的证明.
\end{proof}

剩下只需证明上面帮助我们估计边界项的引理.不深入全部细节, 这是通过考虑三种情形完成的:
\[
k' \geq k'' \geq k, \qquad k' \geq k > k'', \qquad k > k' \geq k''
\]
我们注意到第三种("低-低")情形不可能出现.另外两种情形用 LP3 与下面这个(简单的)微积分不等式来界定:
\begin{equation}
\|f\|_{L_t^\infty L_x^2} \lesssim \|\partial_t f\|_{L_t^2 L_x^2}^{\frac{1}{2}} \cdot \|f\|_{L_t^2 L_x^2}^{\frac{1}{2}} + \|f\|_{L_t^2 L_x^2}\label{eq:c6-248}
\end{equation}
用 \eqref{eq:c6-248} 与 LP3 估计 $\|P_k(g_{k'} \cdot h_{k''})\|_{L_t^\infty L_x^2}$ 即得引理中的估计.

\textbf{习题.} 补全上面定理证明中缺失的步骤.

\section{Calderon-Zygmund 理论}

下面的 $L^2$ 恒等式
\[
\sum_{i, j=1}^n \|\partial_i \partial_j u\|^2_{L^2} = \|\Delta u\|^2_{L^2}.
\]
对任意 $u \in \mathcal{C}_0^\infty(\mathbb{R}^n)$ 都可通过分部积分容易地建立, 见下文 \eqref{eq:c6-252}.于是,
\begin{equation}
\|\partial^2 u\|_{L^2} \lesssim \|\Delta u\|_{L^2}\label{eq:c6-249}
\end{equation}

自然会问, 这样的估计对其他的 $L^p$ 范数是否仍然成立.事实证明, 该问题可以归结为研究一类非常重要的线性算子——称为 Calderon-Zygmund 算子——的 $L^p$ 有界性.

\begin{definition}[Calderon-Zygmund 算子]\label{def:c6-6.1}
作用在 $L^2(\mathbb{R}^n)$ 上的线性算子 $T$ 称为 Calderon-Zygmund 算子, 如果:
\begin{enumerate}
\item $T$ 从 $L^2$ 到 $L^2$ 有界.
\item 存在一个可测核 $k$, 使得对每个具有紧支集的 $f \in L^2$ 以及 $x \notin \operatorname{supp} f$, 有
\[
Tf(x) = \int_{\mathbb{R}^n} k(x-y) f(y) \mathrm{d}y,
\]
其中积分在 $\operatorname{supp} f$ 的补集中对一切 $x$ 绝对收敛.
\item 存在常数 $C > 1$ 与 $A > 0$, 使得
\begin{equation}
\int_{|x| \geq C|y|} |k(x-y) - k(x)| \mathrm{d}x \leq A, \label{eq:c6-250}
\end{equation}
对 $y$ 一致成立.为简单起见, 可以取 $C = 2$.
\end{enumerate}
\end{definition}

\begin{proposition}\label{prop:c6-6.2}
假设核 $k(x)$ 对一切 $x \neq 0$ 满足
\begin{equation}
|k(x)| \lesssim |x|^{-n}, \qquad |\partial k(x)| \lesssim |x|^{-n-1}\label{eq:c6-251}
\end{equation}
则 $k$ 满足消没条件 \eqref{eq:c6-250}.
\end{proposition}

\textbf{习题.} 证明该命题.

\textbf{例 1.} Hilbert 变换 $Hf(x) = \int e^{ix \cdot \xi} \operatorname{sign} \xi\, \hat{f}(\xi) \mathrm{d}\xi$.由 Plancherel 定理容易验证 $H$ 是 $L^2$ 上的有界线性算子.另一方面, 我们知道 $\operatorname{sign} \xi$ 的逆 傅里叶变换与主值分布 $\operatorname{pv}(1/x)$ 成正比.因此, 若 $x \notin \operatorname{supp} f$, 则
\[
Hf(x) = c \int_{-\infty}^{+\infty} \frac{1}{x-y} f(y) \mathrm{d}y.
\]
容易验证核 $k(x) = \frac{1}{x}$ 满足上面的条件 3.

\textbf{例 2.} 考虑 $\mathbb{R}^n$($n \geq 3$)中的方程 $\Delta u = f$, 其中 $f$ 光滑且具紧支集.回忆(见 (??)), 任何在 $\infty$ 处消失\footnote{在 $n = 2$ 的情形, 指任何一阶导数在 $\infty$ 处消失的解.}的解 $u$ 都可以表示为 $u = K_n * f$, 其中 $K_n(x) = c_n |x|^{2-n}$.于是, 若 $x \notin \operatorname{supp} f$, 就可以在积分号下求导, 得到
\[
\partial_i \partial_j u = \partial_i \partial_j K_n * f = \int_{\mathbb{R}^n} \partial_i \partial_j K_n (x-y) f(y) \mathrm{d}y.
\]
容易验证核 $k(x) = \partial_i \partial_j K_n(x)$ 满足条件 3.为证明算子
\[
R_{ij} f(x) = \int_{\mathbb{R}^n} \partial_i \partial_j K_n (x-y) f(y) \mathrm{d}y
\]
是 Calderon-Zygmund 算子, 只剩下 $L^2$ 有界性需要验证.这容易由方程 $\Delta u = f$ 推出.事实上, $u = K_n * f$ 是该方程在 $\infty$ 处消失的唯一解.此外, $|u(x)| \lesssim |x|^{2-n}$, $|\partial u(x)| \lesssim |x|^{1-n}$, 且 $R_{ij} f = \partial_i \partial_j u(x)$.于是我们可以在下面的表达式中分部积分,
\begin{equation}
\begin{array}{rcl}
\int_{\mathbb{R}^n} |f(x)|^2 \mathrm{d}x &=& \int_{\mathbb{R}^n} \Delta u(x) \Delta u(x) \mathrm{d}x = \sum_{i, j=1}^n \int_{\mathbb{R}^n} |\partial_i \partial_j u(x)|^2 \mathrm{d}x \\
&=& \sum_{i, j=1}^n \int_{\mathbb{R}^n} |R_{ij} f(x)|^2 \mathrm{d}x
\end{array}\label{eq:c6-252}
\end{equation}
因此对每一对 $1 \leq i, j \leq n$,
\[
\|R_{ij} f\|_{L^2} \leq \|f\|_{L^2}.
\]
于是算子 $R_{ij}$ 是 Calderon-Zygmund 的.我们将示意地写 $R_{ij} = \partial_i \partial_j (-\Delta)^{-1}$.

\begin{theorem}\label{thm:c6-6.3}
Calderon-Zygmund 算子从 $L^1$ 有界地映入弱 $L^1$.
\end{theorem}

作为推论我们得到,

\begin{corollary}\label{cor:c6-6.4}
Calderon-Zygmund 算子对任意 $1 < p < \infty$ 从 $L^p$ 有界地映入 $L^p$.一般地, 它们对 $p = 1$ 与 $p = \infty$ 不有界.
\end{corollary}

\begin{proof}
对 $1 < p < 2$, $L^p$ 上的有界性由弱 $L^1$ 与 $L^2$ 有界性经 Marcinkiewicz 插值得到.$p > 2$ 的情形由对偶性从如下事实得出:核为 $k(x)$ 的 Calderon-Zygmund 算子的对偶, 是核为 $k(-x)$ 的 Calderon-Zygmund 算子.更精确地说, 若 $f, g$ 具有不相交的支集, 则
\[
\int_{\mathbb{R}^n} Tf(x) g(x) \mathrm{d}x = \int_{\mathbb{R}^n} \int_{\mathbb{R}^n} k(x-y) f(y) g(x) \mathrm{d}x = \int_{\mathbb{R}^n} f(y) T^* g(y) \mathrm{d}y
\]
其中
\[
T^* g(y) = \int_{\mathbb{R}^n} k(-y+x) g(x) \mathrm{d}x, \qquad \forall y \notin \operatorname{supp} g.
\]
另一方面, $\|T^* f\|_{L^2} = \|Tf\|_{L^2} \lesssim \|f\|_{L^2}$.因此 $T^*$ 确实是 CZ 算子.现在, 利用 $L^p$ 与 $L^{p'}$ 之间的对偶性($1/p + 1/p' = 1$)以及 $T^*$ 对 $p' \leq 2$ 的 $L^{p'}$ 有界性,
\[
\begin{array}{ll}
\|Tf\|_{L^p} &= \sup_{\|g\|_{L^{p'}} \leq 1} \left|\int_{\mathbb{R}^n} Tf(x) g(x) \mathrm{d}x\right| = \sup_{\|g\|_{L^{p'}} \leq 1} \left|\int_{\mathbb{R}^n} f(x) T^* g(x) \mathrm{d}x\right| \\
&= \sup_{\|g\|_{L^{p'}} \leq 1} \|f\|_{L^p} \cdot \|T^* g\|_{L^{p'}} \lesssim \|f\|_{L^p}.
\end{array}
\]
\end{proof}

我们将在下面两个小节中证明主\thmref{thm:c6-6.3}.

\mtsubsec{Calderon-Zygmund 分解}

在导言中研究 Laplace 算子时, 我们发现, 对 $f \in \mathcal{C}_0^\infty$, 若 $\int f(x) \mathrm{d}x = 0$, 则 $\Delta^{-1}(f)$ 在 $f$ 的支集之外迅速衰减.这一事实在物理上很重要:它解释了为什么我们必须把手"接触"一个物体才能移动它, 尽管众所周知, 当存在正负电荷的聚集时, 同样的电磁力可以在远得多的距离上移动物体.我们还发现, 与振荡相关的一个特殊行为对 CZO($\Delta^{-1}$ 本身不是其例, 但它的近亲 Riesz 位势是)的分析相当重要.因此, 我们接下来用一节来讨论把一般函数分解为一部分有界、其余部分振荡且在物理上局部化的方法, 这个分解将使我们能够证明\thmref{thm:c6-6.3}.

\begin{definition}[二进立方体]\label{def:c6-6.6}
我们把 $\mathbb{R}^n$ 中的\newterm{dyadic-cube}{二进立方体}{dyadic cube}定义为形如
\[
Q = [2^k a_1, 2^k(a_1+1)] \times \cdots \times [2^k a_n, 2^k(a_n+1)],
\]
的立方体 $Q$, 其中 $k, a_1, \ldots, a_n \in \mathbb{Z}$.此时我们说 $\operatorname{size}(Q) = 2^k$.若 $Q$ 是一个二进立方体, 则它的父立方体是唯一满足 $Q \subset Q^*$ 且 $\operatorname{size}(Q^*) = 2 \operatorname{size}(Q)$ 的二进立方体 $Q^*$, 并称 $Q$ 是 $Q^*$ 的子立方体.
\end{definition}

\begin{lemma}[Whitney 分解]\label{lem:c6-6.7}
$\mathbb{R}^n$ 中任何真开集 $\Omega$ 都可以被一族互不相交的二进立方体 $\mathcal{Q} = \{Q\}$ 覆盖:
\[
\Omega = \cup_{Q \in \mathcal{Q}} Q,
\]
其中每个立方体 $Q \in \mathcal{Q}$ 满足性质
\begin{equation}
\operatorname{size}(Q) \approx \operatorname{dist}(Q, \partial \Omega).\label{eq:c6-253}
\end{equation}
\end{lemma}

\begin{proof}
对每个 $x \in \Omega$, 记 $Q_x$ 为包含 $x$ 的、满足性质 $\operatorname{dist}(Q_x, \partial \Omega) > \operatorname{size}(Q_x)$ 的最大二进立方体.若 $Q^*$ 表示 $Q_x$ 的父立方体, 则 $\operatorname{dist}(Q^*, \partial \Omega) \leq \operatorname{size}(Q^*)$.由三角不等式得
\[
\operatorname{dist}(Q_x, \partial \Omega) \leq \sqrt{n}\, \operatorname{size}(Q_x) + \operatorname{dist}(Q^*, \partial \Omega) \leq (\sqrt{n} + 2) \operatorname{size}(Q_x).
\]
因此, $Q_x$ 满足 \eqref{eq:c6-253}.若 $y \in Q_x$, 则由 $Q_x$ 与 $Q_y$ 的极大性, 必有 $Q_y = Q_x$.因此, 族 $\mathcal{Q} = \{Q_x\}_{x \in \Omega}$ 由互不相交的立方体组成并覆盖 $\Omega$.
\end{proof}

\begin{proposition}[Calderon--Zygmund 分解]\label{prop:c6-6.8}
设 $f \in L^1(\mathbb{R}^n)$, $\alpha > 0$.则可以找到可数族互不相交的二进立方体 $\mathcal{Q} = \{Q\}$ 以及一个分解 $f = g + \sum_{Q \in \mathcal{Q}} b_Q$, 使得:
\begin{equation}
\|g\|_{L^\infty} \lesssim \alpha, \tag{254a}\label{eq:c6-254a}
\end{equation}
\begin{equation}
\operatorname{supp} b_Q \subseteq Q, \tag{254b}\label{eq:c6-254b}
\end{equation}
\begin{equation}
\int b_Q (x) \mathrm{d}x = 0, \tag{254c}\label{eq:c6-254c}
\end{equation}
\begin{equation}
\|b_Q\|_{L^1} \lesssim \alpha |Q|, \tag{254d}\label{eq:c6-254d}
\end{equation}
\begin{equation}
\sum_Q |Q| \lesssim \frac{1}{\alpha} \|f\|_{L^1}.\tag{254e}\label{eq:c6-254e}
\end{equation}
\stepcounter{equation}% (254a)-(254e) 占原书编号 254,手动推进使下一式从 (255) 开始
\end{proposition}

\begin{remark*}
注意在上面的 $\alpha, g, b_Q$ 与 $f$ 都具有相同的量纲, 因此这些关于 $g$ 与 $b_Q$ 的大小和支集的估计, 是唯一可能的、量纲上仍然正确的估计.
\end{remark*}

\begin{proof}
设 $\mathcal{Q}$ 是开集 $\Omega = \{\mathcal{M} f(x) > \alpha\}$ 的 Whitney 分解, 如\lemref{lem:c6-6.7} 所述.对每个 $Q$, 定义 $f_Q = |Q|^{-1} \int_Q f(x) \mathrm{d}x$.令
\[
g(x) = \left\{ \begin{array}{ll} f(x), & \text{若 } x \notin \Omega, \\ f_Q, & \text{若 } x \in Q, \end{array} \right.
\]
并令 $b_Q(x) = \chi_Q(x) (f(x) - f_Q)$, 其中 $\chi_Q$ 是立方体 $Q$ 的特征函数.当然我们有 $f = g + \sum_Q b_Q$.由 \eqref{eq:c6-253} 推出的重要性质是:每个立方体 $Q$ 都包含在一个球 $B$ 内, 该球并不完全包含在 $\Omega$ 中, 且 $|Q| \approx |B|$.设 $x \in B \setminus \Omega$, 有
\begin{equation}
|f_Q| \leq \frac{1}{|Q|} \int_Q |f(y)| \mathrm{d}y \lesssim \frac{1}{|B|} \int_B |f(y)| \mathrm{d}y \leq \mathcal{M} f(x) \leq \alpha.\label{eq:c6-255}
\end{equation}

我们现在来验证这个分解具有所要求的性质.对 $\Omega$ 外几乎处处 $x$, 由 Lebesgue 微分定理(推论 2.7), 有 $|g(x)| \leq \mathcal{M} f(x) \leq \alpha$.当 $x \in \Omega$ 时, 由 \eqref{eq:c6-255} 得 $g(x) \lesssim \alpha$.因此 \eqref{eq:c6-254a} 成立.性质 \eqref{eq:c6-254b} 与 \eqref{eq:c6-254c} 是 $b_Q$ 定义的直接推论.性质 \eqref{eq:c6-254d} 由 \eqref{eq:c6-255} 推出.最后, \eqref{eq:c6-254e} 正是定理 2.4 中证明的 $\mathcal{M} f$ 的弱 $L^1$ 性质.
\end{proof}
\mtsubsec{\thmref{thm:c6-6.3} 的证明}

考虑 $f \in L^1$ 与 $\alpha > 0$.设 $f = g + \sum_Q b_Q = g + b$ 是按\propref{prop:c6-6.8} 给出的 $f$ 的 Calderon-Zygmund 分解.由于
\[
\{|Tf(x)| > \alpha\} \subseteq \{|Tg(x)| > \alpha/2\} \cup (\{|Tb(x)| > \alpha/2\})
\]
并且根据 \eqref{eq:c6-254e}, 只需分别证明
\begin{equation}
|\{|Tg(x)| > \alpha/2\}| \lesssim \frac{1}{\alpha} \|f\|_{L^1}, \label{eq:c6-256}
\end{equation}
\begin{equation}
|\{|Tb(x)| > \alpha/2\}| \lesssim \frac{1}{\alpha} \|f\|_{L^1}\label{eq:c6-257}
\end{equation}

估计 \eqref{eq:c6-256} 由 Chebyshev 不等式、$T$ 在 $L^2$ 上的有界性以及 $g$ 的一致界得出:
\[
\begin{array}{l}
|\{|Tg(x)| > \alpha/2\}| \lesssim \frac{1}{\alpha^2} \|Tg\|^2_{L^2} \lesssim \frac{1}{\alpha^2} \|g\|^2_{L^2} \lesssim \frac{1}{\alpha} \|g\|_{L^1} \leq \\
\qquad \leq \frac{1}{\alpha} \left(\|f\|_{L^1} + \sum_Q \|b_Q\|_{L^1}\right) \lesssim \frac{1}{\alpha} \|f\|_{L^1} + \sum_Q |Q| \lesssim \frac{1}{\alpha} \|f\|_{L^1}.
\end{array}
\]

剩下要推出 \eqref{eq:c6-257}.由于族 $\mathcal{Q}$ 是可数的, 我们记它的成员为 $Q_j$, $j \in \mathbb{N}$.对每个 $Q_j$, 设 $y_{(j)}$ 为其中心, 取 $\hat{Q}_j$ 为具有相同中心、但各边按 $2n^{1/2}$ 扩张的立方体, 使得对 $\hat{Q}_j$ 补集中的一切 $x$,
\[
|x - y_{(j)}| \geq 2 \max_{y \in Q_j} |y - y_{(j)}|
\]

设 $\Omega = \cup_j \hat{Q}_j$, $F$ 为其补集.记 $b_j = b_{Q_j}$.由于 $\int b_j \mathrm{d}y = 0$, 对 $x \in F$ 写
\[
T(b_j)(x) = \int_{Q_j} \left(k(x-y) - k(x - y_{(j)})\right) b_j(y) \mathrm{d}y,
\]
或者, 由于立方体 $Q_j$ 互不相交,
\[
T(b_j)(x) = \int_{Q_j} \left(k(x-y) - k(x - y_{(j)})\right) b(y) \mathrm{d}y,
\]

于是, 根据 \eqref{eq:c6-250},
\[
\begin{array}{rcl}
\int_F |T(b)(x)| \mathrm{d}x &\leq& \sum_j \int_F |T(b)(x)| \mathrm{d}x \lesssim \sum_j \int_{x \in \mathbb{R}^n \setminus \hat{Q}_j} \int_{y \in Q_j} |k(x-y) - k(x - y_{(j)})|\, |b(y)|, \\
&=& \sum_j \int_{y \in Q_j} |b_j(y)| \int_{x \in \mathbb{R}^n \setminus \hat{Q}_j} |k(x-y) - k(x - y_{(j)})| \\
&\leq& \sum_j \int_{y \in Q_j} |b(y)| \int_{x \in \mathbb{R}^n \setminus \{\hat{Q}_j - y_{(j)}\}} |k(x - (y - y_j)) - k(x)| \\
&\lesssim& \sum_j \int_{y \in Q_j} |b(y)| \int_{|x| \geq 2|(y - y_j)|} |k(x - (y - y_j)) - k(x)| \\
&\lesssim& A \sum_j \int_{y \in Q_j} |b(y)| \lesssim \|f\|_{L^1}
\end{array}
\]

因此,
\[
|\{x \in F: |Tb(x)| > \alpha/2\}| \lesssim \alpha^{-1} \|f\|_{L^1}
\]

另一方面, $F$ 的补集(即 $\Omega = \cup \hat{Q}_j$)的测度也可控:
\[
|\Omega| \leq \sum_j |\hat{Q}_j| \lesssim \sum_j |Q_j| \lesssim \alpha^{-1} \|f\|_{L^1}.
\]

因此,
\[
\left|\left\{x \in \mathbb{R}^n: |Tb(x)| > \alpha/2\right\}\right| \lesssim \alpha^{-1} \|f\|_{L^1}
\]
如所愿.

\mtsubsec{Michlin-Hörmander 定理}

一类重要的 CZ 算子可以通过 Fourier 乘子算子来定义.回忆这些算子由 傅里叶变换定义:
\begin{equation}
\widehat{Tf}(\xi) = m(\xi) \hat{f}(\xi), \label{eq:c6-258}
\end{equation}
其中 $m$ 是一个有界函数, 称为乘子.我们可以把这些算子看作卷积算子 $Tf = k * f$, 其中 $\hat{k} = m$.自然会问, Fourier 乘子算子何时给出一个 CZ 算子.既然我们知道 CZO 会赋予一个均值零的局部化函数额外的衰减, 我们应当预期乘子 $m$ 应该离原点相当远.这正是下面定理的内容.

\begin{theorem}\label{thm:c6-6.11}
设 $l > n/2$.假设 $m$ 是 $\widehat{\mathbb{R}}^n \setminus 0$ 上 $C^l$ 类的 Fourier 乘子, 满足
\[
|\partial_\xi^\alpha m(\xi)| \lesssim |\xi|^{-|\alpha|}, \qquad \forall \xi \in \widehat{\mathbb{R}}^n \setminus 0
\]
对每个满足 $|\alpha| \leq l$ 的多重指标 $\alpha$ 成立.则由 \eqref{eq:c6-258} 定义的算子是一个 Calderon-Zygmund 算子.
\end{theorem}

\begin{proof}
考虑与 LP 投影中所用相同的二进单位分解:
\[
1 = \sum_{\lambda \in 2^{\mathbb{Z}}} \chi_\lambda(\xi) \qquad \text{对} \quad \xi \in \mathbb{R}^n \setminus 0,
\]
由 $\chi \in C_0^\infty$ 生成, 其中 $\operatorname{supp} \chi \subseteq \{1/2 \leq |\xi| \leq 2\}$, 且 $\chi_\lambda(\xi) = \chi(\xi/\lambda)$.

把 $m$ 分解为二进片, $m = \sum_\lambda m_\lambda$, 其中 $m_\lambda = \chi_\lambda m$.由于 $|\partial^\gamma m(\xi)| \lesssim |\xi|^{-|\gamma|}$ 且 $\chi(\xi)$ 的所有导数有界,
\begin{equation}
|\partial_\xi^\alpha m_\lambda(\xi)| \leq \sum_{|\beta| + |\gamma| \leq |\alpha|} |\partial^\beta \chi_\lambda\, |\xi|^{-\gamma}| \lesssim \sum_{|\beta| + |\gamma| \leq |\alpha|} \lambda^{-|\beta|} \lambda^{-|\gamma|} \approx \lambda^{-|\alpha|}.\label{eq:c6-259}
\end{equation}

设 $k_\lambda$ 是 $m_\lambda$ 的逆 傅里叶变换.由于 $m_\lambda$ 具有紧支集, $k_\lambda$ 是光滑函数.此外, 对任何整数 $N$, 有\footnote{回忆:由分部积分, 有 $|F_\xi^{-1} f(x)| \leq |x|^{-N} \|\partial^N f\|_{L^1}$.}
\[
|k_\lambda(x)| \lesssim |x|^{-N} \|\partial^N m_\lambda\|_{L^1} \lesssim |x|^{-N} \lambda^{n-N}.
\]

现在取 $N > n$ 并对 $\lambda \in 2^{\mathbb{Z}}$ 求和.注意到 $\sum_\lambda k_\lambda$ 收敛到 $\mathbb{R}^n \setminus 0$ 上一个良定义的可测函数 $k$, 容易看出 $k$ 满足\defref{def:c6-6.1} 的性质 2.

$T$ 在 $L^2$ 上的有界性直接由 $m$ 在 $\widehat{\mathbb{R}}^n$ 上的有界性得出.

对 $0 \leq j \leq l$, 由 Plancherel 定理与 \eqref{eq:c6-259} 得
\[
\int |x|^{2j} |k_\lambda(x)|^2 \mathrm{d}x \simeq \sum_{|\alpha| = j} \int |\partial_\xi^\alpha m_\lambda(\xi)|^2 \mathrm{d}\xi \lesssim \lambda^{n-2j}.
\]

设 $R > 0$, 用 $j = 0$ 的情形我们发现
\begin{equation}
\int_{|x| \leq R} |k_\lambda(x)| \mathrm{d}x \lesssim \left(\int |k_\lambda(x)|^2 \mathrm{d}x\right)^{1/2} R^{n/2} \lesssim (\lambda R)^{n/2}, \label{eq:c6-260}
\end{equation}
而用 $j = l$ 的情形我们发现
\begin{equation}
\int_{|x| \geq R} |k_\lambda(x)| \mathrm{d}x \lesssim \left(\int |x|^{2l} |k_\lambda(x)|^2 \mathrm{d}x\right)^{1/2} \left(\int_{|x| > R} \frac{\mathrm{d}x}{|x|^{2l}}\right)^{1/2} \lesssim (\lambda R)^{n/2 - l}.\label{eq:c6-261}
\end{equation}

若取 $R = 1/\lambda$, 把 \eqref{eq:c6-260} 与 \eqref{eq:c6-261} 相加, 得 $\|k_\lambda\|_{L^1} \lesssim 1$ 对 $\lambda$ 一致成立.我们可以对 $\partial k_\lambda$ 应用同样的过程, 其符号为 $\xi m_\lambda \approx \lambda m_\lambda$, 以证明 $\|\partial k_\lambda\|_{L^1} \lesssim \lambda$.于是,
\begin{equation}
\int_{|x| \gg |y|} |k_\lambda(x-y) - k_\lambda(x)| \mathrm{d}x \leq \int \int_0^{|y|} |\partial k_\lambda (x - t y/|y|)| \mathrm{d}t \mathrm{d}x\label{eq:c6-262}
\end{equation}
\begin{equation}
= |y| \cdot \|\partial k_\lambda\|_{L^1} \lesssim \lambda |y|, \label{eq:c6-263}
\end{equation}
而且, 由 \eqref{eq:c6-261},
\begin{equation}
\int_{|x| \gg |y|} |k_\lambda(x-y) - k_\lambda(x)| \mathrm{d}x \leq 2 \int_{|x| \geq |y|} |k_\lambda(x)| \mathrm{d}x \lesssim (\lambda |y|)^{n/2 - l}.\label{eq:c6-264}
\end{equation}

当 $\lambda |y| \leq 1$ 时用 \eqref{eq:c6-262}, 当 $\lambda |y| > 1$ 时用 \eqref{eq:c6-264}, 对 $\lambda$ 求和, 得\footnote{这里我们使用了二进记号下几何级数的如下求和性质:对 $\alpha > 0$, 有 $\sum_{\lambda \leq L} \lambda^\alpha \simeq L^\alpha$ 与 $\sum_{\lambda \geq L} \lambda^{-\alpha} \simeq L^{-\alpha}$.}
\[
\int_{|x| \gg |y|} |k(x-y) - k(x)| \mathrm{d}x \lesssim |y| \sum_{\lambda \leq |y|^{-1}} \lambda + |y|^{n/2 - l} \sum_{\lambda > |y|^{-1}} \lambda^{n/2 - l} \lesssim 1.
\]
如所愿.
\end{proof}

\textbf{习题.} 设 $\phi \in \mathcal{C}_0^\infty(\mathbb{C})$, $f$ 是不齐次 Cauchy-Riemann 方程 $\frac{\partial f}{\partial \bar{z}} = \phi$ 的、在无穷远处衰减的解.证明对 $1 < p < \infty$ 有估计
\[
\|\partial f\|_{L^p} \lesssim \|\phi\|_{L^p}
\]

\mtsubsec{平方函数估计}

我们回忆平方函数的性质 LP6:
\[
Sf = \left(\sum_k |P_k f|^2\right)^{1/2},
\]

\begin{theorem}[Littlewood--Paley]\label{thm:c6-6.13}
我们有
\begin{equation}
\|f\|_{L^p} \lesssim \|Sf\|_{L^p} \lesssim \|f\|_{L^p}\label{eq:c6-265}
\end{equation}
对所有 $1 < p < \infty$ 成立.
\end{theorem}

我们给出这个估计的两个证明.

\begin{proof}[第一个证明]
首先我们用对偶论证证明 \eqref{eq:c6-265} 中的第一个不等式由第二个推出.事实上, 利用 Plancherel 定理、事实 $P_k P_{k'} = 0$(除非 $k \sim k'$)以及 Cauchy-Schwarz 不等式, 我们得到
\[
\begin{array}{lll}
\int f(x) g(x) \mathrm{d}x &\simeq& \int \sum_{k \approx k'} P_k f(x) P_{k'} g(x) \mathrm{d}x \\
&\lesssim& \int \left(\sum_k |P_k f(x)|^2\right)^{1/2} \left(\sum_{k'} |P_{k'} g(x)|^2\right)^{1/2} \mathrm{d}x \leq \\
&\lesssim& \|Sf\|_{L^p} \|Sg\|_{L^{p'}} \lesssim \|Sf\|_{L^p} \|g\|_{L^{p'}}.
\end{array}
\]
\eqref{eq:c6-265} 中的左端不等式现在通过对所有满足 $\|g\|_{L^{p'}} = 1$ 的 $g$ 取上确界得出.

为证明 \eqref{eq:c6-265} 中的右端不等式, 我们需要引入 Rademacher 函数 $r_k(t)$, 它们定义在 $\mathbb{R}$ 上:对每个 $k \geq 0$, $k \in \mathbb{Z}$, 以及 $t \in \mathbb{R}$, 令 $r_k(t) = r_0(2^k t)$, 其中 $r_0(t)$ 是周期函数, $r_0(t+1) = r_0(t)$, 满足 $0 \leq t < 1/2$ 时 $r_0(t) = 1$, $1/2 \leq t < 1$ 时 $r_0(t) = -1$.这些 Rademacher 函数在 $L^2[0, 1]$ 中构成正交序列, 并且构成一列独立同分布的随机变量.我们所需的基本性质是:Rademacher 函数的线性组合的 $L^p$ 范数等价于其系数的 $l^2$ 范数.

\begin{lemma}\label{lem:c6-6.14}
给定满足 $\sum_{k=0}^{\infty} a_k^2 < \infty$ 的实数序列 $\{a_k\}$, 定义
\[
F(t) = \sum_{k=0}^{\infty} a_k r_k(t).
\]
则 $F \in L^2([0, 1])$, 且 $\|F\|_{L^2} = (\sum_{k=0}^{\infty} a_k^2)^{1/2}$.此外, 对 $1 < p < \infty$, 有 $F \in L^p([0, 1])$, 并且存在常数 $A_p$ 使得
\[
A_p^{-1} \|F\|_{L^p} \leq \|F\|_{L^2} \leq A_p \|F\|_{L^p}.
\]
\end{lemma}

此引理的证明见 Stein, [?, 附录 D].

定义算子 $T_t$ 如下:
\[
T_t f = \sum_{k=0}^{\infty} r_k(t) P_k f
\]

显然 $T_t$ 是以 $m_t(\xi) = \sum_k r_k(t) \chi(2^{-k}\xi)$ 为符号的 Fourier 乘子算子, 其中 $\chi$ 是用于定义 LP 投影的光滑截断函数.对 $\xi \neq 0$, 定义 $m_t(\xi)$ 的和式中至多只有三项非零.于是我们可以容易地验证 $m_t$ 满足\thmref{thm:c6-6.11} 的条件.即
\[
|\partial_\xi^\alpha m_t(\xi)| \leq C_\alpha |\xi|^{-|\alpha|},
\]
其中常数 $C_\alpha$ 与 $t$ 无关.因此, 由 Calderon-Zygmund 理论(特别是\corref{cor:c6-6.4}), 有:
\[
\|T_t f\|_{L^p} \lesssim \|f\|_{L^p}
\]
于是,
\[
\left(\int_0^1 \|T_t f\|^p_{L^p} \mathrm{d}t\right)^{1/p} \lesssim \|f\|_{L^p}
\]

此外, 我们可以用\lemref{lem:c6-6.14} 看到:
\[
\begin{array}{ll}
\int_0^1 \|T_t f\|^p_{L^p} \mathrm{d}t &= \int_0^1 \int_{\mathbb{R}} \left|\sum_k r_k(t) (P_k f)(x)\right|^p \mathrm{d}x \mathrm{d}t \\
&\gtrsim \int_{\mathbb{R}} \left(\sum_k |(P_k f)(x)|^2\right)^{p/2} \mathrm{d}x
\end{array}
\]

于是综合我们的结果得到:
\[
\|Sf\|_{L^p} \lesssim \|f\|_{L^p}
\]

(注意这个论证只在单维情形 $n = 1$ 证明了定理.不过, 它可以像 Stein 的《Singular Integrals》第 IV 章第 5 节那样推广到 $\mathbb{R}^n$.)
\end{proof}

\begin{proof}[第二个证明]
我们回忆向量值函数的定义:
\[
\mathbf{S} f(x) = (P_k f(x))_{k \in \mathbb{Z}}.
\]
显然, 若 $f \in \mathcal{S}(\mathbb{R}^n)$, 则对每个 $x \in \mathbb{R}^n$, $\mathbf{S} f(x) \in l^2$, 且 $Sf(x) = |\mathbf{S} f(x)|$ 表示 $\mathbf{S} f(x)$ 的 $l^2$ 范数.我们断言
\[
\mathbf{S} f(x) = \int \mathbf{K}(x-y) f(y) \mathrm{d}y
\]
是一个 $l^2$ 值的 Calderon-Zygmund 算子, 其 $l^2$ 值核定义为
\[
\mathbf{K}(x) = (K_k(x))_{k \in \mathbb{Z}}, \qquad K_k(x) = 2^{nk} \hat{\chi}(2^k x)
\]
记
\[
|\mathbf{K}(x)| = \left(\sum_k |K_k(x)|^2\right)^{1/2}, \qquad |\partial \mathbf{K}(x)| = \left(\sum_k |\partial K_k(x)|^2\right)^{1/2}
\]
我们容易验证条件 \eqref{eq:c6-251} 的 $l^2$ 值版本成立,
\begin{equation}
|\mathbf{K}(x)| \lesssim |x|^{-n}, \qquad |\partial \mathbf{K}(x)| \lesssim |x|^{-(n+1)}, \qquad \text{对} \quad x \neq 0.\label{eq:c6-266}
\end{equation}
另一方面,
\[
\|\mathbf{S} f\|_{L^2} := \|Sf\|_{L^2} \lesssim \|f\|_{L^2}.
\]
因此 $\mathbf{S}$ 确实是一个 $l^2$ 值的 C-Z 算子, 从而, 根据\thmref{thm:c6-6.3} 及其推论的直接推广, 我们推得
\[
\|\mathbf{S} f\|_{L^p} := \||\mathbf{S} f|\|_{L^p} = \|Sf\|_{L^p} \lesssim \|f\|_{L^p}
\]
根据我们定理第一个证明的开头, 我们还推得
\[
\|f\|_{L^p} \lesssim \|Sf\|_{L^p}.
\]
\end{proof}

注意, 根据\thmref{thm:c6-6.13}, $|\sum_k P_k f| \approx (\sum_k |P_k f|^2)^{1/2}$.一个更一般的原理断言:若一列函数 $f_1, f_2, \ldots, f_k, \ldots$ 以不同的速率振荡, 即任意两个相位不同, 则
\[
\left|\sum_k f_k\right| \approx \left(\sum_k |f_k|^2\right)^{1/2}
\]

性质 LP6 与\thmref{thm:c6-6.13} 的如下版本对与 $P_k$ 相似的 LP 投影 $\tilde{P}_k \sim P_k$ 也成立.更精确地说,
\begin{equation}
\left\|\left(\sum_k |\tilde{P}_k f|^2\right)^{1/2}\right\|_{L^p} \lesssim \|f\|_{L^p}, \qquad 1 < p < \infty.\label{eq:c6-267}
\end{equation}
这可以用与 $\|Sf\|_{L^p} \lesssim \|f\|_{L^p}$ 相同的方式证明:引入 $l^2$ 值算子 $\tilde{\mathbf{S}} f = (\tilde{P}_k f)_{k \in \mathbb{Z}}$, 并完全按照\thmref{thm:c6-6.13} 第二个证明的过程进行.给定一个 $l^2$ 值向量函数 $\mathbf{g} = (g_k)_{k \in \mathbb{Z}}$, 注意到
\[
\langle \tilde{\mathbf{S}} f, \mathbf{g} \rangle = \int_{\mathbb{R}^n} \tilde{\mathbf{S}} f(x) \cdot \bar{\mathbf{g}}(x) \mathrm{d}x = \int_{\mathbb{R}^n} \sum_k \tilde{P}_k f(x) \overline{g_k}(x) \mathrm{d}x = \int_{\mathbb{R}^n} f(x) \overline{\sum_k \tilde{P}_k g_k (x)} \mathrm{d}x
\]
于是,
\begin{equation}
\tilde{\mathbf{S}}^* \mathbf{g} = \sum_k \tilde{P}_k g_k\label{eq:c6-268}
\end{equation}
因此, 与 \eqref{eq:c6-267} 对偶的估计具有形式 $\|\tilde{\mathbf{S}}^* \mathbf{g}\|_{L^{p'}} \lesssim \|\mathbf{g}\|_{L^{p'}}$, 其中 $1/p + 1/p' = 1$.换言之,
\begin{equation}
\left\|\sum_k \tilde{P}_k g_k\right\|_{L^p} \lesssim \left\|\left(\sum_k |g_k|^2\right)^{1/2}\right\|_{L^p}, \qquad 1 < p < \infty.\label{eq:c6-269}
\end{equation}

下面是\thmref{thm:c6-6.13} 的一个容易的推论.

\begin{corollary}\label{cor:c6-6.15}
对 $2 \leq p < \infty$, 有
\begin{equation}
\|f\|^2_{L^p} \lesssim \sum_{k \in \mathbb{Z}} \|P_k f\|^2_{L^p}.\label{eq:c6-270}
\end{equation}
对 $1 < p \leq 2$, 有
\begin{equation}
\sum_{k \in \mathbb{Z}} \|P_k f\|^2_{L^p} \lesssim \|f\|^2_{L^p}.\label{eq:c6-271}
\end{equation}
\end{corollary}

\begin{proof}
回忆 $Sf(x)^2 := \sum_{k \in \mathbb{Z}} |P_k f|^2$.若 $p/2 \geq 1$, 根据 LP6 与 Minkowski 不等式, 有
\[
\|f\|^2_{L^p} \lesssim \|Sf\|^2_{L^p} = \left\|\sum_k |P_k f|^2\right\|_{L^{p/2}} \leq \sum_k \left\||P_k f|^2\right\|_{L^{p/2}} = \sum_k \|P_k f\|^2_{L^p}.
\]
若 $p/2 \leq 1$, 我们改而使用反向 Minkowski 不等式,
\[
\|f\|^2_{L^p} \gtrsim \left\|\sum_k |P_k f|^2\right\|_{L^{p/2}} \geq \sum_k \left\||P_k f|^2\right\|_{L^{p/2}} = \sum_k \|P_k f\|^2_{L^p}.
\]
\end{proof}

我们在这里使用的反向 Minkowski 不等式断言:对 $0 < q \leq 1$ 和一列正函数 $(f_k)_{k \in \mathbb{Z}}$,
\begin{equation}
\left\|\sum_k |f_k|\right\|_{L^q} \geq \sum_k \|f_k\|_{L^q}.\label{eq:c6-272}
\end{equation}

我们简要勾画 \eqref{eq:c6-272} 的证明;它可以在许多书中找到(例如 Garling 的《Inequalities》或 DiBenedetto 的《Real Analysis》, 我们这里采用的正是这个特定证明).

一种方法是先证明一个反向 Hölder 不等式:对 $0 < p < 1$, $q < 0$, $\frac{1}{p} + \frac{1}{q} = 1$, $f \in L^p$, $g \in L^q$, 有
\[
\int |fg| \geq \|f\|_{L^p} \|g\|_{L^q}.
\]
这可以容易地通过写出 $\|f\|_{L^p} = \left(\int \frac{|fg|^p}{|g|^p}\right)^{1/p}$, 并对指数 $\tilde{p} = 1/p > 1$ 与 $\tilde{q} = 1/(1-p) > 1$ 应用通常的 Hölder 不等式来证明.

有了这个, 两项的反向 Minkowski 不等式($\||f| + |g|\|_{L^q} \geq \|f\|_{L^q} + \|g\|_{L^q}$, $0 < q \leq 1$)如下推出(写 $\frac{1}{q'} = 1 - \frac{1}{q}$):
\[
\begin{array}{rcl}
\||f| + |g|\|^q_{L^q} &=& \int (|f| + |g|)^{q-1} (|f| + |g|) \\
&\geq& \left(\int (|f| + |g|)^{(q-1)q'}\right)^{1/q'} (\|f\|_{L^q} + \|g\|_{L^q}) \\
&\geq& \||f| + |g|\|^{q-1}_{L^q} (\|f\|_{L^q} + \|g\|_{L^q})
\end{array}
\]

\mtsubsec{$W^{s, p}$ Sobolev 空间}

我们回忆已经定义的函数的 $W^{s, p}$ 范数:
\[
\|f\|_{W^{s, p}} = \sum_{j=0}^s \|\partial^j f\|_{L^p}
\]
我们断言下面的事实:

\begin{lemma}\label{lem:c6-6.17}
对任何 $j \geq 0$, $1 < p < \infty$, 有
\[
\|\partial^j f\|_{L^p} \approx \left\|\left(\sum_k |2^{jk} P_k f|^2\right)^{1/2}\right\|_{L^p}
\]
\end{lemma}

\begin{proof}
我们先写
\[
\|\partial^j f\|_{L^p} \lesssim \left\|\sum_k \partial^j P_k f\right\|_{L^p}
\]
如同性质 LP5 的证明中那样, 我们可以把 $\nabla^j P_k f$ 表示为 $\nabla^j P_k f = 2^{jk} \tilde{P}_k P_k f$, 其中 $\tilde{P}_k$ 与 $P_k$ 相似.于是, 利用估计 \eqref{eq:c6-269},
\[
\|\partial^j f\|_{L^p} \lesssim \left\|\sum_k 2^{jk} \tilde{P}_k P_k f\right\|_{L^p} \lesssim \left\|\left(\sum_k |2^{jk} P_k f|^2\right)^{1/2}\right\|_{L^p}.
\]
另一方面, 我们也可以对某个其他的相似 LP 投影写 $2^{jk} P_k f = \tilde{P}_k \partial^j f$.于是, 根据 \eqref{eq:c6-267},
\[
\left\|\left(\sum_k |2^{jk} P_k f|^2\right)^{1/2}\right\|_{L^p} \lesssim \left\|\left(\sum_k |\tilde{P}_k \partial^j f|^2\right)^{1/2}\right\|_{L^p} \lesssim \|\partial^j f\|_{L^p}
\]
\end{proof}

利用该引理, 我们现在可以找到使用 LP 投影的等价定义:

\begin{proposition}\label{prop:c6-6.18}
对任何 $1 < p < \infty$ 与任何 $s \in \mathbb{N}$, 有
\begin{equation}
\|f\|_{W^{s, p}} \approx \left\|\sum_k (1 + 2^k)^s P_k f\right\|_{L^p}.\label{eq:c6-273}
\end{equation}
此外, 对齐次 $W^{s, p}$ 范数 $\|f\|_{\dot{W}^{s, p}} = \|\partial^s f\|_{L^p}$,
\begin{equation}
\|f\|_{\dot{W}^{s, p}} \approx \left\|\sum_k 2^{ks} P_k f\right\|_{L^p}.\label{eq:c6-274}
\end{equation}
\end{proposition}

注意 \eqref{eq:c6-273} 与 \eqref{eq:c6-274} 右端的表达式对每个 $s \in \mathbb{R}$ 都有意义.因此我们可以把 $W^{s, p}$ 与 $\dot{W}^{s, p}$ 空间的定义推广到一切实数 $s$.

齐次 Sobolev 范数 $\|\cdot\|_{\dot{W}^{s, p}}$ 的另外一些刻画可以用下面的事实给出:

\begin{proposition}\label{prop:c6-6.19}
对 $2 \leq p < \infty$ 与任何 $s$, 有
\begin{equation}
\left(\sum_k 2^{kps} \|P_k f\|^p_{L^p}\right)^{1/p} \lesssim \|f\|_{\dot{W}^{s, p}} \lesssim \left(\sum_k 2^{2ks} \|P_k f\|^2_{L^p}\right)^{1/2}.\label{eq:c6-275}
\end{equation}
对 $1 < p \leq 2$ 与 $s \in \mathbb{R}$, 有
\begin{equation}
\left(\sum_k 2^{2ks} \|P_k f\|^2_{L^p}\right)^{1/2} \lesssim \|f\|_{\dot{W}^{s, p}} \lesssim \left(\sum_k 2^{kps} \|P_k f\|^p_{L^p}\right)^{1/p}.\label{eq:c6-276}
\end{equation}
\end{proposition}

\begin{proof}
若 $p/2 \geq 1$, 由\thmref{thm:c6-6.13} 与 Minkowski 不等式, 有
\[
\|f\|^2_{L^p} \lesssim \left\|\sum_k |P_k f|^2\right\|_{L^{p/2}} \leq \sum_k \left\||P_k f|^2\right\|_{L^{p/2}} = \sum_k \|P_k f\|^2_{L^p}.
\]
若 $p/2 \leq 1$, 由\thmref{thm:c6-6.13} 与反向 Minkowski 不等式, 有
\[
\|f\|^2_{L^p} \gtrsim \left\|\sum_k |P_k f|^2\right\|_{L^{p/2}} \geq \sum_k \left\||P_k f|^2\right\|_{L^{p/2}} = \sum_k \|P_k f\|^2_{L^p}.
\]
其余细节显然可以补全.
\end{proof}
\section{问题}

\noindent\textbf{问题 1.}~[$\mathbb{R}$ 中的分布]

设 $f(z)$ 是区域 $D_+ = \{z \in \mathbb{C}: 0 < \Im(z) < \epsilon\}$ 中的解析函数, 满足对一切 $z \in D$ 有 $|f(z)| \lesssim |\Im(z)|^{-N}$.证明存在一个分布 $f_+ = f(\cdot + i0)$, 使得对每个 $\phi \in \mathcal{C}_0^\infty(\mathbb{R}^n)$,
\[
\lim_{y \to 0, y > 0} \int_{\mathbb{R}} f(x + iy) \phi(x) \mathrm{d}x = \langle f_+, \phi \rangle,
\]
类似地, 对定义在 $D_- = \{z \in \mathbb{C}: -\epsilon < \Im(z) < 0\}$ 上的解析函数, 我们可以定义分布 $f_- = f(\cdot - i0)$,
\[
\lim_{y \to 0, y < 0} \int_{\mathbb{R}} f(x + iy) \phi(x) \mathrm{d}x = \langle f_-, \phi \rangle
\]
这特别地在 $f = \frac{1}{z} = \frac{1}{x + iy}$ 时定义了分布 $(x + i0)^{-1}$ 与 $(x - i0)^{-1}$.证明公式
\[
(x + i0)^{-1} - (x - i0)^{-1} = -2\pi i \delta_0(x).
\]
再证明
\[
(x + i0)^{-1} = x^{-1} - i\pi \delta_0(x)
\]
其中 $\frac{1}{x}$ 是正文中定义的主值分布.

\noindent\textbf{问题 2.}~[基本解] 考虑 $\mathbb{R}^3$ 中的算子 $Lu = \Delta u + u$.求 $Lu = 0$ 的所有球对称解.证明
\[
K(x) = -\frac{\cos |x|}{4\pi |x|}
\]
是 $L$ 的基本解.

\noindent\textbf{问题 3.}~[初值问题] 考虑 $\mathbb{R} \times \mathbb{R}^n$ 中下列四个演化方程的初值问题:
\begin{equation}
\partial_t u = \Delta u, \qquad u(0, x) = f(x)\label{eq:c6-277}
\end{equation}
\begin{equation}
\partial_t u = i\Delta u, \qquad u(0, x) = f(x)\label{eq:c6-278}
\end{equation}
\begin{equation}
\partial_t^2 u = \Delta u, \qquad u(0, x) = f(x), \quad \partial_t u(0, x) = g(x)\label{eq:c6-279}
\end{equation}
\begin{equation}
\partial_t^2 u = -\Delta u, \qquad u(0, x) = f(x), \quad \partial_t u(0, x) = g(x)\label{eq:c6-280}
\end{equation}
在每种情形下, 用 傅里叶变换方法写出解.换言之, 对每个方程取 傅里叶变换, 令
\[
\hat{u}(t, \xi) = \int e^{-ix \cdot \xi} u(t, x) \mathrm{d}x
\]
并解所得的关于 $t$ 的微分方程.比较最后两个方程的结果.证明 \eqref{eq:c6-279} 对一切 $f, g \in \mathcal{S}(\mathbb{R}^n)$ 都有解, 而 \eqref{eq:c6-280} 没有.不过要证明:如果我们只给定 $u(0, x) = f$(这是 $\mathbb{R}^{n+1}$ 中 Laplace 算子 $\partial_t^2 + \Delta$ 的 Dirichlet 问题), 则对一切函数 $f \in \mathcal{S}(\mathbb{R}^n)$, 该问题都有唯一解 $u$, 且当 $|t| + |x| \to \infty$ 时 $u$ 衰减到零.在所有情形下, 把所得的解表示为作用于初值(在物理空间中)的积分算子\footnote{你将需要执行逆 傅里叶变换, $u(t, x) = \mathcal{F}^{-1} \hat{u}(t, \xi)$.对波动方程, 一般来说这更困难, 但你可以在维数 $n = 3$ 时做到.}.

\noindent\textbf{问题 4.}~[延拓算子] 设 $H$ 是 $\mathbb{R}^n$ 中的半空间 $x_n > 0$, $1 \leq p \leq \infty$.证明存在一个延拓算子, 即有界线性算子 $E: W^{1, p}(H) \to W^{1, p}(\mathbb{R}^n)$, 使得对所有 $u \in W^{1, p}(H)$, 在 $H$ 中几乎处处有 $Eu = u$, 且
\[
\|Eu\|_{W^{1, p}(\mathbb{R}^n)} \lesssim \|u\|_{W^{1, p}(H)}.
\]
把结果推广到任何 $s \in \mathbb{N}$.你能把结果推广到任意区域 $U \subset \mathbb{R}^n$ 吗?具有光滑边界的区域呢?

\noindent\textbf{问题 5.}~[圆周上的分布与 Fourier 分析] 圆周 $\mathbb{R}/\mathbb{Z}$ 上的光滑函数是 $\mathbb{R}$ 上满足 1-周期性的光滑函数
\[
f(x + k) = f(x), \quad k \in \mathbb{Z}
\]
圆周有一个离散的频率空间 $m \in \widehat{(\mathbb{R}/\mathbb{Z})} = \mathbb{Z}$, 对应于函数 $x \mapsto e^{2\pi i m x}$.频率空间的离散性与圆周的紧性密切相关.圆周上的 Schwartz 函数就是光滑函数;$\mathbb{Z}$ 上的 Schwartz 函数是在无穷远处比任何多项式都更快衰减的函数.

\begin{enumerate}
\renewcommand{\labelenumi}{\alph{enumi}.}
\item 我们定义周期函数的 傅里叶变换 $\hat{f}(m) = \int_0^1 f(x) e^{-2\pi i m x}$.证明光滑函数在圆周上的 Fourier 反演公式
\[
f(x) = \sum_{m \in \mathbb{Z}} \hat{f}(m) e^{2\pi i m x}
\]
并推出 Plancherel 公式 $\langle f, g \rangle = \langle \hat{f}, \hat{g} \rangle$.
\item 我们把圆周上的分布 $u$ 定义为某个 $k$ 的 $\mathcal{C}^k(\mathbb{R}/\mathbb{Z})$ 的对偶中的元素, 即对某个 $k, C$ 以及所有 $\phi \in \mathcal{C}^\infty(\mathbb{R}/\mathbb{Z})$, 有 $|\langle u, \phi \rangle| \leq C \|\phi\|_{C^k}$.圆周具有光滑结构, 因此可以表述微分算子的基本解的概念(圆周上的群结构也使得卷积有意义)——但并非总能找到这样的解.证明算子 $\frac{d}{dx}$ 在圆周上没有基本解 $u$.换言之, 不存在满足 $\frac{du}{dx} = \delta(x)$ 的分布 $u$, 其意义是
\[
\langle \frac{du}{dx}, \phi \rangle \equiv -\langle u, \frac{d\phi}{dx} \rangle = \phi(0), \quad \phi \in \mathcal{C}^\infty(\mathbb{R}/\mathbb{Z})
\]
证明这一点有很多方法.你能同时在物理空间和频率空间中看到这一点吗?如果我们把向量场 $\frac{d}{dx}$ 换成另一个不消失的向量场 $\tilde{D} = \psi \frac{d}{dx}$, 其中 $\psi \in \mathcal{C}^\infty(\mathbb{R}/\mathbb{Z})$ 是不消失的光滑函数, 会怎么样?
\end{enumerate}

\noindent\textbf{问题 6.}~[迹定理] 设 $\mathbb{R}^{n-1}$ 是 $\mathbb{R}^n$ 中的超平面, 例如 $x_n = 0$.对任何 $f \in \mathcal{S}(\mathbb{R}^n)$, 记 $Rf$ 为 $f$ 到 $\mathbb{R}^{n-1}$ 的限制.

\begin{enumerate}
\renewcommand{\labelenumi}{\roman{enumi}.}
\item 证明:对任何 $s > \frac{1}{2}$,
\begin{equation}
\|Rf\|_{L^2(\mathbb{R}^{n-1})} \lesssim \|f\|_{H^s(\mathbb{R}^n)}\label{eq:c6-281}
\end{equation}
\item 证明该结果对 $s \leq 1/2$ 不成立.不过要证明, 对一切 $s > 0$, 下面的尖锐迹定理成立:
\begin{equation}
\|Rf\|_{H^s(\mathbb{R}^{n-1})} \lesssim \|f\|_{H^{s+1/2}(\mathbb{R}^n)}\label{eq:c6-282}
\end{equation}
\item 证明:若 $f$ 的 Fourier 支集在球 $|\xi| \lesssim 2^k$ 中(对某个整数 $k$), 则对所有 $1 \leq p \leq \infty$ 与 $s > 1/p$,
\[
\|f\|_{L^p(\mathbb{R}^{n-1})} \lesssim 2^{k/p} \|f\|_{W^{s, p}(\mathbb{R}^n)}
\]
你能由此推出一个用 $L^p$ 范数表述的、推广 \eqref{eq:c6-281} 的迹结果吗?\eqref{eq:c6-282} 呢?
\item 设 $H$ 是半空间 $x_n > 0$.根据上面的考虑, 我们可以谈论 $W^{1, p}(H)$ 中函数到超平面 $x_n = 0$ 的迹(证明这一点!).证明:函数 $f \in W^{1, p}(H)$ 属于 $W_0^{1, p}(H)$\footnote{回忆 $W_0^{1, p}(H)$ 是 $\mathcal{C}_0^\infty(H)$ 在 $W^{1, p}(H)$ 中的闭包.}当且仅当它到 $x_n = 0$ 的迹为零.
\end{enumerate}

\noindent\textbf{问题 7.}~[Littlewood-Paley] 考虑空间 $\Lambda_\gamma = C^{0, \gamma}(\mathbb{R}^n)$, 其范数为
\[
\|f\|_{\Lambda_\gamma} = \|f\|_{L^\infty(\mathbb{R}^n)} + \sup_{x \neq y \in \mathbb{R}^n} \frac{|u(x) - u(y)|}{|x - y|^\gamma}
\]
\begin{enumerate}
\renewcommand{\labelenumi}{\roman{enumi}.}
\item 证明, 用 Littlewood-Paley 投影 $P_k$, 有
\[
\|f\|_{\Lambda^\gamma} \approx \|P_{\leq 0} f\|_{L^\infty} + \sup_{k > 0} 2^{k\gamma} \|P_k f\|_{L^p}
\]
\item 定义 Zygmund 类 $\Lambda_*$, 其范数为
\[
\|f\|_{\Lambda^*} = \|f\|_{L^\infty} + \sup_{x \in \mathbb{R}^n, 0 \leq h \leq 1} \frac{|f(x+h) + f(x-h) - 2f(x)|}{h}
\]
证明
\[
\|f\|_{\Lambda_*} \approx \|P_{\leq 0} f\|_{L^\infty} + \sup_{k > 0} 2^k \|P_k f\|_{L^p}.
\]
\item 证明 Besov 空间 $B^s = H^{s, 1}$, $s > 0$ 中的乘积估计
\[
\|fg\|_{B^s} \lesssim \|f\|_{L^\infty} \|g\|_{B^s} + \|g\|_{L^\infty} \|f\|_{B^s}
\]
\end{enumerate}

\noindent\textbf{问题 8.}~自己阅读关于 Calderon-Zygmund 算子的那一节.说明该理论如何推广到取值于给定 Hilbert 空间(如 $l^2$)的算子.
\section{限制定理}

众所周知, 当 $f \in L^1(\mathbb{R}^n)$ 时, 它的 傅里叶变换 $\hat{f}$ 是有界连续函数, 因此 $\hat{f}$ 到任何超曲面上的限制都是完全良定义的.另一方面, 若 $f \in L^2(\mathbb{R}^n)$, 则 $\hat{f}$ 可以是 $L^2$ 中的任何函数, 因此它只在几乎处处有定义, 而在像超曲面这样的零测集上完全任意.

当 $f$ 属于某个 $1 < p < 2$ 的 $L^p$ 时, 能否给 $\hat{f}$ 到光滑超曲面 $S$ 的限制赋予意义?这是现代 Fourier 分析中的一个基本问题, 我们将看到, 它最终与波动方程解的正则性性质密切相关.

若取 $S$ 为超平面, 我们立刻看到答案是否定的.事实上, 设 $f(x_1, x') = u(x_1) v(x')$, $\hat{f}(\xi_1, \xi') = \hat{u}(\xi_1) \hat{v}(\xi')$, 其中 $x_1, \xi_1 \in \mathbb{R}$, $x', \xi' \in \mathbb{R}^{n-1}$.$\hat{f}$ 到超平面 $\xi_1 = 0$ 上的限制仅在 $\hat{u}(0) = \int u(x) \mathrm{d}x$ 良定义时才有定义.对任何 $p > 1$, 总能找到 $u \in L^p(\mathbb{R})$ 使得 $\int u \mathrm{d}x$ 没有意义.我们推断, 当 $p > 1$ 时, 傅里叶变换在超平面上的限制无法定义.

如果我们考虑具有非零曲率的超曲面, 答案就不同了.为简单起见, 我们考虑球面的模型情形.

\mtsubsec{Stein-Tomas 定理}

下面这类结果最初由 Stein 证明, 然后由 Tomas 推广, 并再次由 Stein 给出最终形式.

\begin{theorem}[Stein--Tomas]\label{thm:c6-8.2}
设 $\mathbb{S} = \mathbb{S}^{n-1}$ 是 $\mathbb{R}^n$ 中的标准单位球面, $\mathrm{d}\sigma$ 是其标准体积元.设 $f \in L^p(\mathbb{R}^n)$, 其中
\[
1 \leq p \leq \frac{2(n+1)}{n+3}.
\]
则 $\mathcal{R} f = \hat{f}\big|_{\mathbb{S}} \in L^2(\mathbb{S})$, 且
\[
\|\mathcal{R} f\|_{L^2(\mathbb{S})} \lesssim \|f\|_{L^p(\mathbb{R}^n)}.
\]
\end{theorem}

这个定理有一个等价的对偶表述.定义\newterm{stein-operator}{Stein 算子}{Stein operator}为 Fourier 限制算子 $\mathcal{R} f = \hat{f}\big|_{\mathbb{S}}$ 的对偶,
\[
\mathcal{S} g(x) = \mathcal{R}^* g(x) = \int_{\mathbb{S}} e^{ix \cdot \xi} g(\xi) \mathrm{d}\sigma_\xi \simeq (g \mathrm{d}\sigma)^\vee (x),
\]
其中 $g$ 现在是定义在球面上的函数.

\begin{theorem}\label{thm:c6-8.3}
设 $f \in L^2(\mathbb{S})$, 且
\[
\frac{2(n+1)}{n-1} \leq p \leq \infty.
\]
则 $\mathcal{S} f \in L^p(\mathbb{R}^n)$, 且
\begin{equation}
\|\mathcal{S} f\|_{L^p(\mathbb{R}^n)} \lesssim \|f\|_{L^2(\mathbb{S})}.\label{eq:c6-283}
\end{equation}
\end{theorem}

\begin{remark}\label{rem:c6-8.4}
只需对 $p = p_* = 2(n+1)/(n-1)$ 证明\thmref{thm:c6-8.3}.事实上, 对 $p > p_*$, 由 Sobolev 不等式有
\[
\|\mathcal{S} f\|_{L^p} \lesssim \|D^s \mathcal{S} f\|_{L^{p_*}}
\]
其中 $s = n(1/p_* - 1/p) > 0$, $(D^s u)^\wedge (\xi) = |\xi|^s \hat{u}(\xi)$.但这里
\[
D^s \mathcal{S} f = \mathcal{S}(|\cdot|^s f) = \mathcal{S} f.
\]
因此, 如果我们能在 $p = p_*$ 时证明定理, 则
\[
\|\mathcal{S} f\|_{L^p} \lesssim \|\mathcal{S} f\|_{L^{p_*}} \lesssim \|f\|_{L^2(\mathbb{S})}
\]
\end{remark}

\begin{remark}\label{rem:c6-8.5}
如果用 $\mathrm{d}\mu = \psi \mathrm{d}\sigma$($\psi \in C_0^\infty(\mathbb{R}^n)$)替换 $\mathrm{d}\sigma$, 结果仍然成立, 因为定理推出
\[
\|(f \mathrm{d}\mu)^\vee\|_{L^p} \lesssim \|f \psi\|_{L^2(\mathbb{S})} \lesssim \|f\|_{L^2(\mathbb{S})}.
\]
此外, 利用单位分解, 只需对 $\boldsymbol{S} f = (f \mathrm{d}\mu)^\vee$(其中 $\mathrm{d}\mu = \psi \mathrm{d}\sigma$, $\psi \in C_0^\infty(\mathbb{R}^n)$ 的支集在球面某点的小邻域内)证明\thmref{thm:c6-8.3}.虽然这是显然的, 但我们可以局部化限制估计这一事实非常重要, 正如我们将来会看到的那样.
\end{remark}

\mtsubsec{Knapp 反例}

根据下面的反例([?]), \thmref{thm:c6-8.3} 的结果对任何 $p < p_*$ 都不成立.

对某个小的 $\delta > 0$, 定义\newterm{phase-space}{相空间}{phase space}中的区域
\[
D = \left\{\xi \in \mathbb{R}^n: |\xi_1 - 1| < \delta^2, |\xi'| < \delta\right\}.
\]
现在令 $f = \chi_{\mathbb{S} \cap D}$ 为帽 $\mathbb{S} \cap D$ 的特征函数, 则
\[
\|f\|_{L^2(\mathbb{S})} = |\mathbb{S} \cap D|^{1/2} \sim \delta^{(n-1)/2}.
\]
我们可以写
\[
\mathcal{S} f(x) = e^{ix_1} \int_{\mathbb{S} \cap D} e^{i\phi(x, \xi)} \mathrm{d}\sigma_\xi,
\]
其中相位 $\phi(x, \xi) = x_1(\xi_1 - 1) + x' \cdot \xi'$.于是可以固定物理空间中的一个区域,
\[
R = \left\{x \in \mathbb{R}^n: |x_1| < \frac{\pi}{6} \delta^{-2}, |x'| < \frac{\pi}{6} \delta^{-1}\right\},
\]
使得对 $x \in R$ 与 $\xi \in D$, 有 $|\phi(x, \xi)| \le \pi/3$, 因此, 当 $x \in R$ 时,
\[
|\mathcal{S} f(x)| \geq \operatorname{Re}\left(e^{-ix_1} \mathcal{S} f(x)\right) = \int_{\mathbb{S} \cap D} \cos(\phi(x, \xi)) \mathrm{d}\sigma_\xi \geq \frac{1}{2} |\mathbb{S} \cap D|.
\]
这推出
\[
\frac{\|\mathcal{S} f\|_{L^p}}{\|f\|_{L^2}} \gtrsim |\mathbb{S} \cap D|^{1/2} |R|^{1/p} \sim \delta^{\frac{n-1}{2} - \frac{n+1}{p}}.
\]
对小的 $\delta$, 形如 \eqref{eq:c6-283} 的估计必然要求
\[
\frac{n-1}{2} - \frac{n+1}{p} \geq 0
\]
这只有在 $p \geq p_* = 2(n+1)/(n-1)$ 时才可能成立.

这个例子提示, 算子 $\mathcal{S}$ 的结构中存在着某种抛物型标度性质, 它来自球面的非零曲率.

\mtsubsec{曲率的重要性}

如果我们把标准球面 $\mathbb{S}^{n-1}$ 换成具有非零 Gauss 曲率的紧超曲面 $H \subset \mathbb{R}^n$, 限制定理及其对偶仍然成立.非零 Gauss 曲率的重要性由下面的结果说明.

\begin{lemma}\label{lem:c6-8.8}
设 $H \subset \mathbb{R}^n$ 是带非零 Gauss 曲率(即所有主曲率都不为零)与体积元 $\mathrm{d}\sigma$ 的紧超曲面.则对任何光滑函数 $\psi$, 有
\begin{equation}
|(\psi \mathrm{d}\sigma)^\vee (x)| \lesssim (1 + |x|)^{-\frac{n-1}{2}}\label{eq:c6-284}
\end{equation}
若恰好有一个主曲率为零, 则我们有
\[
|(\psi \mathrm{d}\sigma)^\vee (x)| \lesssim (1 + |x|)^{-\frac{n-2}{2}}
\]
\end{lemma}

\begin{proof}
一般的证明基于稳定相方法, 见 Stein 的《Harmonic Analysis》一书.对标准球面 $H = \mathbb{S}^{n-1}$ 与奇数 $n$ 的特殊情形, 可以通过极坐标直接计算来证明.
\end{proof}

\textbf{习题.} 对 $\mathbb{S}^2 \subset \mathbb{R}^3$ 证明该引理.

\begin{remark}\label{rem:c6-8.9}
另一个有趣的观察把这些限制定理与偏微分方程联系起来.事实上, 若 $u = \mathrm{d}\sigma^\vee * f$, 则 $u$ 是线性椭圆方程
\[
\Delta u + u = 0,
\]
的解, 这只要取 傅里叶变换就可以容易地看到:
\[
\mathcal{F}(u + \Delta u)(\xi) \simeq (1 - |\xi|^2) \delta(1 - |\xi|) \hat{f}(\xi) = 0.
\]
其中 $\delta$ 是 Dirac 分布.
\end{remark}

\mtsubsec{$TT^*$ 原理}

下面这个简单的泛函分析结果在限制与 Strichartz 型估计中起重要作用.设 $B$ 是 Banach 空间, 记 $B'$ 为其对偶.设 $H$ 是带内积 $\langle \cdot, \cdot \rangle$ 的 Hilbert 空间.考虑线性算子 $T: H \to B'$.由于我们可以把 $H$ 与其对偶等同, 我们可以把 $T$ 看作算子 $T^*: B \to H$ 的伴随, 后者由下式定义:
\[
\langle h, T^*(x) \rangle = Th(x).
\]
实际上, 当 $B$ 自反时, $T^*$ 是 $T$ 的伴随, 但为了我们的目的, 我们将继续称 $T^*$ 为 $T$ 的伴随.

$TT^*$ 原理断言 $T$ 的有界性与 $TT^*$ 的有界性等价.更精确地说, 我们有:

\begin{proposition}\label{prop:c6-8.11}
下面的陈述是等价的:
\begin{enumerate}
\renewcommand{\labelenumi}{(\roman{enumi})}
\item $T: H \to B'$ 有界且 $\|T\| = M$;
\item $T^*: B \to H$ 有界且 $\|T^*\| = M$;
\item $TT^*: B \to B'$ 有界且 $\|TT^*\| = M^2$;
\item 双线性形式 $(x, y) \mapsto \langle T^* x, T^* y \rangle$ 在 $B \times B$ 上有界, 范数为 $M^2$.
\end{enumerate}
\end{proposition}

证明是泛函分析中的一个标准练习.


\mtsubsec{限制定理的 $TT^*$ 表述}

Stein 算子的 $TT^*$ 表述对应于与球面上测度的(逆)傅里叶变换做卷积.形式上, 我们有
\[
SS^* f(x) = S\mathcal{R} f(x) = \int_{\mathbb{S}} e^{ix \cdot \xi} \hat{f}(\xi) \mathrm{d}\sigma_\xi = \int_{\mathbb{R}^n} \int_{\mathbb{S}} e^{i(x-y) \cdot \xi} \mathrm{d}\sigma_\xi f(y) \mathrm{d}y = \mathrm{d}\sigma^\vee * f(x).
\]
于是我们被引向限制定理的如下等价形式,
\begin{equation}
\|\mathrm{d}\sigma^\vee * f\|_{L^p(\mathbb{R}^n)} \lesssim \|f\|_{L^{p'}(\mathbb{R}^n)}, \label{eq:c6-285}
\end{equation}
对 $p \geq p^*$ 成立.

对\thmref{thm:c6-8.3} 可以给出三个不同的证明.我们将勾勒基于解析插值的第一个证明.这本质上是 Stein 与 Tomas 的原始证明.第二个证明基于引入一个时间参数并把 $\boldsymbol{S} f$ 看作演化算子, 它使我们能把限制定理视为一个更一般框架的一部分, 该框架包含各种带常系数的线性偏微分方程的 Strichartz 估计.最后, 第三个方法只适用于特定的指数, 它将使我们能够与双线性估计联系起来.

\mtsubsec{第一个证明:解析插值}

根据注 8.12 与\remref{rem:c6-8.4}, 只需证明 $\boldsymbol{U} f = \mathrm{d}\sigma^\vee * f$ 满足
\begin{equation}
\|Uf\|_{L^{p_*}(\mathbb{R}^n)} \lesssim \|f\|_{L^{p_*'}(\mathbb{R}^n)}, \label{eq:c6-286}
\end{equation}
其中 $p_* = 2(n+1)/(n-1)$, $p_*' = 2(n+1)/(n+3)$.

一般来说, 直接得到 $L^{p'} - L^p$ 估计通常非常复杂, 除了 $p$ 是像 $p = 4, 6$ 这样的好指数(这只发生在 $n = 2$ 或 $n = 3$)的情形外, 我们不知道任何直接证明.我们更习惯于 $L^2 - L^2$ 型估计(Plancherel 定理在其中是强大的工具), 或 $L^1 - L^\infty$ 型估计, 因为振荡积分的逐点衰减估计可以由稳定相方法得到.这提示对 $L^p$ 空间使用某种插值理论.但是, 对算子 $U$ 的 $L^2 - L^2$ 估计被 Knapp 反例排除, 而 $L^\infty - L^1$ 估计太琐碎, 不能回答我们的问题.正是在这里, Stein 插值定理(定理 1.8)显示出它的威力, 因为它允许我们从其他(合理的)不同于 $U$ 的算子的 $L^2 - L^2$ 与 $L^\infty - L^1$ 估计得到 $U$ 的 $L^{p'} - L^p$ 估计.

我们将通过构造一族卷积算子 $U_z f = \mu_z^\vee * f$ 来实现这一点, 其中 $\mu_z$ 是依赖于 $z$ 的解析分布.参数 $z$ 本质上将反映分布 $\mu_z$ 的齐次度.为此, 自然地把我们的目标放在 $z = -1$, 要求 $U_{-1} = U$ 或 $\mu_{-1} = \mathrm{d}\sigma$, 因为 $\mathrm{d}\sigma$ 可以写成球面上一个 delta 分布(它是 $-1$ 次齐次的)的回拉:
\[
\mathrm{d}\sigma \simeq \delta(1 - |\xi|) \mathrm{d}\xi.
\]

若 $\mu_z$ 与一个有界函数一致, 则 $U_z$ 的 $L^2 - L^2$ 估计将成立;事实上, 由 Plancherel 定理, 有
\begin{equation}
\|U_z f\|_{L^2} \simeq \|(U_z f)^\wedge\|_{L^2} \simeq \|\mu_z \cdot \hat{f}\|_{L^2} \lesssim \|\mu_z\|_{L^\infty} \|f\|_{L^2}.\label{eq:c6-287}
\end{equation}
要使 $\mu_z(\xi)$ 有界, 我们必须要求 $\mu_z(\xi)$ 本质上是 0 次齐次的, 即当 $z$ 位于直线 $\mathrm{Re}(z) = 0$ 上时.

相反, 当 $\mu_z^\vee$ 与一个有界函数一致时, $U_z$ 的 $L^1 - L^\infty$ 估计将成立, 因为我们直接有
\begin{equation}
\|U_z f\|_{L^\infty} \lesssim \|\mu_z^\vee\|_{L^\infty} \|f\|_{L^1}.\label{eq:c6-288}
\end{equation}
为了从 \eqref{eq:c6-287} 与 \eqref{eq:c6-288} 的解析插值得到 \eqref{eq:c6-286}, 我们希望后者发生在直线 $\mathrm{Re}(z) = a$ 上, 其中 $a$ 的选择使得
\[
-1 = \theta a + (1-\theta) 0, \qquad \frac{1}{p_*} = \frac{\theta}{\infty} + \frac{1-\theta}{2}, \qquad \frac{1}{p_*'} = \frac{\theta}{1} + \frac{1-\theta}{2},
\]
而这恰好发生在 $\mathrm{Re}(z) = a = -(n+1)/2$ 时.

这个论证引向我们将要使用的 Stein 解析插值定理的精确版本.

\begin{proposition}\label{prop:c6-8.14}
设 $U_z$ 是满足下列条件的线性算子解析族:
\begin{enumerate}
\renewcommand{\labelenumi}{(\roman{enumi})}
\item $U_{-1} = U$;
\item $\|U_z f\|_{L^2} \lesssim \|f\|_{L^2}$, 在直线 $\mathrm{Re}(z) = 0$ 上一致;
\item $\|U_z f\|_{L^\infty} \lesssim \|f\|_{L^1}$, 在直线 $\mathrm{Re}(z) = -(n+1)/2$ 上一致.
\end{enumerate}
则由此推出
\[
\|Uf\|_{L^{p_*}} \lesssim \|f\|_{L^{p_*'}}.
\]
\end{proposition}


上面的讨论表明, 当我们把 $U_z f$ 写成卷积 $\mu_z^\vee * f$ 时, 只要 $\mu_z$ 是满足如下条件的解析分布族, 命题的假设就得到满足:

(i') $\mu_{-1} = \mathrm{d}\sigma$;

(ii') $\mu_z(\xi)$ 与一个有界函数一致, 在直线 $\mathrm{Re}(z) = 0$ 上一致有界;

(iii') $\mu_z^\vee(x)$ 与一个有界函数一致, 在直线 $\mathrm{Re}(z) = -(n+1)/2$ 上一致有界.

于是剩下要定义分布 $\mu_z$ 并验证这些性质.

受恒等式 $\delta = \chi_+^{-1}$ 与 $\mathrm{d}\sigma_\xi \simeq \delta(1 - |\xi|)$ 的启发, 我们定义分布族
\begin{equation}
\mu_z(\xi) = e^{z^2} \chi_+^z (1 - |\xi|) \psi(|\xi|), \label{eq:c6-289}
\end{equation}
其中 $\psi \in C_0^\infty(\mathbb{R})$ 是支集在 1 的小邻域(比如 $[1/2, 3/2]$)内的截断函数, 且 $\psi(1) = 1$.

我们回忆齐次分布 $\chi_+^z$, 当 $\mathrm{Re}(z) > -1$ 时, 它们与函数一致:
\[
\chi_+^z(t) = \left\{ \begin{array}{ll} t^z / \Gamma(z+1) & \text{若 } t \geq 0, \\ 0 & \text{若 } t < 0, \end{array} \right.
\]
其中 Gamma 函数定义为 $\Gamma(z+1) = \int_0^\infty t^z e^{-t} \mathrm{d}t$.由恒等式 $\Gamma(z+1) = z\Gamma(z)$, 推出
\begin{equation}
\frac{\mathrm{d}}{\mathrm{d}t} \chi_+^z(t) = \chi_+^{z-1}(t).\label{eq:c6-290}
\end{equation}
利用这个公式, 通过反复分部积分, $\chi_+^z$ 可以对所有 $z \in \mathbb{C}$ 解析延拓.为此我们首先注意到, 对 $\mathrm{Re}(z) > -1$ 与 $\phi \in C_0^\infty$, 有
\[
\int \chi_+^z(t) \phi(t) \mathrm{d}t = -\int \chi_+^{z+1}(t) \phi'(t) \mathrm{d}t = \cdots = (-1)^m \int \chi_+^{z+m}(t) \phi^{(m)}(t) \mathrm{d}t.
\]
因此, 分部积分足够多次后, 当 $\mathrm{Re}(z) > -1 - m$ 时(对任何 $m$, 从而对所有 $z$), 我们可以使 $\int \chi_+^z \phi \mathrm{d}t$ 有意义.要看到 $\chi_+^{-1} = \delta$, 只需一次分部积分, 事实上
\[
\int \chi_+^{-1} \phi \mathrm{d}t = -\frac{1}{\Gamma(1)} \int_0^\infty \phi'(t) \mathrm{d}t = \phi(0).
\]
关于 $\chi_+^z$ 与分布理论的更多信息, 可以查阅 Gel'fand 与 Shilov 的著作 \cite{ge-s} 或 Hörmander \cite{ho}.

$\mu_z$ 定义中的因子 $e^{z^2}$ 是为了保证我们的算子对大的 $\Im(z)$ 一致有界;事实上, 当 $\Im(z) \to \infty$ 时, $e^{z^2}$ 在带形区域 $-(n+1)/2 \leq \mathrm{Re}(z) \leq 0$ 上一致地指数衰减.这使得下面各个不等式中的常数可以按 $b = \Im(z)$ 多项式增长.

显然 $\mu_{-1} \simeq \delta(1 - |\xi|) \psi(|\xi|) \simeq \mathrm{d}\sigma$.这验证了 (i').

条件 (ii') 立即可验证, 因为当 $\mathrm{Re}(z) = 0$ 时, $\chi_+^{-z}$ 总是一个有界函数.条件 (iii') 将由稳定相论证推出, 更一般地我们有:

\begin{proposition}\label{prop:c6-8.15}
\begin{equation}
|\mu_z^\vee(x)| \lesssim (1 + |x|)^{-\mathrm{Re}(z) - 1 - \frac{n-1}{2}}.\label{eq:c6-291}
\end{equation}
\end{proposition}

\mtsubsec{第二个证明:演化算子方法}

在本节中, 我们对 $f$ 作如下假设:
\begin{equation}
f \in C^\infty(\mathbb{S}), \qquad \operatorname{supp} f \subset \{\xi_1 > 1/2\}.\label{eq:c6-292}
\end{equation}

在这个假设下, 我们可以把 $x_1 = t$ 重新标记为时间参数, 并把 $\boldsymbol{S} f$ 重写为
\[
\begin{array}{c}
\mathcal{S} f(t, x') = \int_{|\xi'| < \sqrt{3}/2} e^{it\sqrt{1 - |\xi'|^2}} e^{ix' \cdot \xi'} f(\sqrt{1 - |\xi'|^2}, \xi') \frac{\mathrm{d}\xi'}{\sqrt{1 - |\xi'|^2}} \\
= \int e^{it\sqrt{1 - |\xi'|^2}} e^{ix' \cdot \xi'} \beta(|\xi'|) g(\xi') \mathrm{d}\xi'.
\end{array}
\]
其中 $\beta \in C_0^\infty$ 的支集在 $|\xi'| < 1$ 内, 且 $g(\xi') = f(\sqrt{1 - |\xi'|^2}, \xi')/\sqrt{1 - |\xi'|^2}$.注意到
\[
\int |g(\xi')|^2 \mathrm{d}\xi' = \int_{\mathbb{S}} \frac{|f(\xi)|^2}{|\xi_1|^2} \mathrm{d}\sigma_\xi \simeq \|f\|^2_{L^2(\mathbb{S})}
\]
由关于 $f$ 支集的假设得出.

\begin{theorem}\label{thm:c6-8.17}
设 $\beta \in C_0^\infty(\mathbf{R}^{n-1})$ 的支集在单位球 $\{\xi \in \mathbf{R}^{n-1}: |\xi| < 1\}$ 内, 考虑算子
\[
Tg(t, x) = \int_{\mathbf{R}^{n-1}} e^{it\sqrt{1 - |\xi|^2}} e^{ix \cdot \xi} \beta(\xi) g(\xi) \mathrm{d}\xi, \qquad t \in \mathbf{R}, x \in \mathbf{R}^{n-1}.
\]
设 $q, r$ 是满足下列条件的 Lebesgue 指数:
\begin{equation}
0 \leq \frac{2}{q} \leq \min\{1, \gamma(r)\}, \label{eq:c6-293}
\end{equation}
\begin{equation}
\left(\frac{2}{q}, \gamma(r)\right) \neq (1, 1), \label{eq:c6-294}
\end{equation}
其中 $\gamma(r) = (n-1)(1/2 - 1/r)$.则对一切 $g \in C_0^\infty(\mathbf{R}^{n-1})$, 下面的估计成立:
\begin{equation}
\|Tg\|_{L^q_t L^r_x(\mathbf{R} \times \mathbf{R}^{n-1})} \lesssim \|g\|_{L^2(\mathbf{R}^{n-1})}.\label{eq:c6-295}
\end{equation}
其中我们使用第 7 节定义的混合范数记号.

由\remref{rem:c6-8.5}, \thmref{thm:c6-8.3} 由特殊情形 $q = r = 2\frac{n+1}{n-1}$ 推出.
\end{theorem}

\begin{remark}\label{rem:c6-8.18}
当 $q \geq 2$ 时, 我们可以再次运行 Knapp 例子来证明条件 \eqref{eq:c6-293} 的必要性.事实上, 设 $D \subset \mathbb{R}^{n-1}$ 是由 $|\xi| \leq \delta$ 定义的圆盘, $\delta > 0$ 足够小, 取 $g = \chi_D$ 为 $D$ 的特征函数.我们写
\[
Tg(t, x) = e^{it} \int_D e^{it(\sqrt{1 - |\xi|^2} - 1)} e^{ix \cdot \xi} \beta(\xi) \mathrm{d}\xi
\]
并注意到, 对 $|t| \leq \delta^{-2}$ 与 $|x| \leq \delta^{-1}$, 存在固定的常数 $c > 0$ 使得 $|Tg(t, x)| \geq c$.事实上, 这容易由 $|\xi| \le \delta$ 与 $|\sqrt{1 - |\xi|^2} - 1| \lesssim \delta^2$ 推出.因此, 若 \eqref{eq:c6-295} 成立, 则对一切足够小的 $\delta > 0$, 必有
\[
c \delta^{-\frac{2}{q}} \delta^{-\frac{n-1}{r}} \lesssim \|Tg\|_{L^q_t L^r_x} \lesssim \|\chi_D\|_{L^2} \lesssim \delta^{-\frac{n-1}{2}}
\]
由此推出 \eqref{eq:c6-293}, $q \geq 2$.
\end{remark}

\begin{remark}\label{rem:c6-8.19}
当 $n \neq 3$ 时, 由于 Keel 与 Tao 的一个著名结果 \cite{k-t}(《Endpoint Strichartz Inequalities》), 端点限制 \eqref{eq:c6-294} 可以被去掉.\eqref{eq:c6-293} 中隐含的另一个限制 $q \geq 2$ 将在下一章讨论.
\end{remark}

我们从计算 $T^*$ 与 $TT^*$ 开始:
\[
\begin{array}{l}
\langle T^* F, g \rangle = \langle F, Tg \rangle = \iint F\, \overline{Tg}\, \mathrm{d}t \mathrm{d}x = \\
\qquad = \iint F(t, x) \int e^{-it\sqrt{1 - |\xi|^2}} e^{-ix \cdot \xi} \overline{\beta(\xi) g(\xi)} \mathrm{d}\xi \mathrm{d}t \mathrm{d}x = \\
\qquad = \int \overline{g(\xi) \beta(\xi)} \left(\iint e^{-it\sqrt{1 - |\xi|^2}} e^{-ix \cdot \xi} F(t, x) \mathrm{d}t \mathrm{d}x\right) \mathrm{d}\xi.
\end{array}
\]
因此
\[
T^* F(\xi) = \overline{\beta(\xi)} \iint e^{-it\sqrt{1 - |\xi|^2}} e^{-ix \cdot \xi} F(t, x) \mathrm{d}t \mathrm{d}x,
\]
且
\[
\begin{array}{rl}
& TT^* F(t, x) = \int e^{it\sqrt{1 - |\xi|^2}} e^{ix \cdot \xi} \beta(\xi) T^* F(\xi) \mathrm{d}\xi \\
& \qquad = \iint e^{i(t-s)\sqrt{1 - |\xi|^2}} e^{ix \cdot \xi} |\beta(\xi)|^2 \hat{F}(s, \xi) \mathrm{d}\xi \mathrm{d}s,
\end{array}
\]
其中 $\hat{F}(s, \xi) = \int e^{-ix \cdot \xi} F(s, x) \mathrm{d}x$.如果我们引入算子族
\[
U(t) f(x) = \int e^{it\sqrt{1 - |\xi|^2}} e^{ix \cdot \xi} |\beta(\xi)|^2 \hat{f}(\xi) \mathrm{d}\xi,
\]
我们可以把 $TT^*$ 写为卷积算子,
\begin{equation}
TT^* F(t, \cdot) = \int U(t-s) F(s, \cdot) \mathrm{d}s.\label{eq:c6-296}
\end{equation}

由\propref{prop:c6-8.11}, 要证明 $T$ 是从 $L^q_t L^r_x(\mathbf{R}^n)$ 到 $L^2(\mathbf{R}^{n-1})$ 的有界算子, 只需证明 $TT^*$ 是从 $L^{q'}_t L^{r'}_x(\mathbf{R}^n)$ 到 $L^q_t L^r_x(\mathbf{R}^n)$ 的有界算子.

我们将首先证明关于 $U(t)$ 的一个估计:

\begin{proposition}\label{prop:c6-8.20}
设 $2 \leq r \leq \infty$, $\gamma(r) = (n-1)(1/2 - 1/r)$.则 $U(t)$ 满足估计
\begin{equation}
\|U(t) f\|_{L^r(\mathbf{R}^{n-1})} \lesssim (1 + |t|)^{-\gamma(r)} \|f\|_{L^{r'}(\mathbf{R}^{n-1})}.\label{eq:c6-297}
\end{equation}
\end{proposition}

\begin{proof}
一旦我们证明了两个极端情形 $r = 2$ 与 $r = \infty$,
\begin{equation}
\|U(t) f\|_{L^2(\mathbf{R}^{n-1})} \lesssim \|f\|_{L^2(\mathbf{R}^{n-1})}\label{eq:c6-298}
\end{equation}
\begin{equation}
\|U(t) f\|_{L^\infty(\mathbf{R}^{n-1})} \lesssim (1 + |t|)^{-(n-1)/2} \|f\|_{L^1(\mathbf{R}^{n-1})}\label{eq:c6-299}
\end{equation}
该估计就由标准的 Riesz 插值定理得出.

我们利用 Plancherel 公式立即得到 \eqref{eq:c6-298}, 因为
\[
(U(t) f)^\wedge (\xi) \simeq e^{it\sqrt{1 - |\xi|^2}} |\beta(\xi)|^2 \hat{f}(\xi).
\]

为证明 \eqref{eq:c6-299}, 我们写
\[
U(t) f(x) = \int K_t (x-y) f(y) \mathrm{d}y,
\]
其中
\[
\begin{array}{rl}
& K_t(x) = \int e^{ix \cdot \xi} e^{it\sqrt{1 - |\xi|^2}} |\beta(\xi)|^2 \mathrm{d}\xi \\
& \quad \simeq \iint e^{ix \cdot \xi} e^{it\tau} \delta(1 - \tau^2 - |\xi|^2) \sqrt{1 - |\xi|^2}\, |\beta(\xi)|^2 \mathrm{d}\tau \mathrm{d}\xi \\
& \quad \simeq \iint e^{i(t, x) \cdot (\tau, \xi)} \delta(1 - |(\tau, \xi)|) \beta_1(\tau, \xi) \mathrm{d}\tau \mathrm{d}\xi, \qquad (\beta_1(\tau, \xi) = \tau |\beta(\xi)|^2), \\
& \quad = (\beta_1 \mathrm{d}\sigma_{n-1})^\vee (t, x).
\end{array}
\]
因此 $K_t$ 正是支集在球面 $\mathbb{S}^{n-1}$ 上的一个测度的 傅里叶变换, 对它有衰减估计
\[
|K_t(x)| \lesssim (1 + |t| + |x|)^{-(n-1)/2},
\]
这推出 \eqref{eq:c6-299}.
\end{proof}

接下来我们把\propref{prop:c6-8.20} 应用于 \eqref{eq:c6-296},
\begin{equation}
\|TT^* F(t, \cdot)\|_{L^r_x} \lesssim \int \frac{1}{(1 + |t-s|)^{\gamma(r)}} \|F(s, \cdot)\|_{L^{r'}_x} \mathrm{d}s.\label{eq:c6-300}
\end{equation}

最后, 我们能够应用 Hardy-Littlewood-Sobolev 不等式, 若 $0 < \gamma(r) < 1$, 则得到
\[
\|TT^* F\|_{L^q_t L^r_x} \lesssim \|F\|_{L^{q'}_t L^{r'}_x},
\]
当 $-\gamma(r) + 1 + 1/q = 1/q'$, 从而 $\gamma(r) = 2/q$ 时成立.因此, 我们在 $0 < \gamma(r) = 2/q < 1$ 的情形证明了\thmref{thm:c6-8.17}.

另一方面, 若 $q = 2$ 且 $\gamma(r) > 1$, 由 \eqref{eq:c6-300},
\[
\|TT^* F\|_{L^2_t L^r_x} \lesssim \|F\|_{L^2_t L^{r'}_x},
\]
应用标准的 Hausdorff-Young 不等式即可.

最后, 若 $2/q < 1$ 且 $\gamma(r) > 2/q$, 结果由 $\gamma(r) = 2/q$ 的情形利用 Sobolev 不等式得出.

\mtsubsec{第三个证明:双线性形式($n = 2$ 与 $n = 3$)}

我们现在给出证明球面限制定理的另一个方法, 它适用于特殊情形 $n = 2$, $p = 6$ 或 $n = 3$, $p = 4$.其思想是:当 $p$ 是偶整数时, 限制定理可以看作一个多重线性形式的 $L^2$ 估计, 而该形式通过 傅里叶变换具有卷积结构, 提供某种光滑化效应.下面给出的证明是所谓双线性/三线性估计的根源, 后者在现代非线性波动方程与色散方程理论中起根本作用.

让我们先看 $n = 3$ 的情形.我们考虑 Stein 算子 $\boldsymbol{S} f = (f \mathrm{d}\sigma)^\vee$, 并利用事实 $(Sf \cdot Sf)^\wedge \simeq (f \mathrm{d}\sigma) * (f \mathrm{d}\sigma)$.设 $B(f, g) = Sf \cdot Sg$, 则 $\boldsymbol{S} f$ 的 $L^4$ 估计对应于 $B(f, f)$ 的 $L^2$ 估计.我们有
\[
\hat{B}(f, g)(\xi) \simeq (f \mathrm{d}\sigma) * (g \mathrm{d}\sigma)(\xi) = \int_{\mathbf{R}^3} \delta(1 - |\xi - \eta|) \delta(1 - |\eta|) f(\xi - \eta) g(\eta) \mathrm{d}\eta,
\]
并对测度 $\delta(1 - |\xi - \eta|) \delta(1 - |\eta|) \mathrm{d}\eta$ 应用 Cauchy-Schwarz, 我们得到
\[
|\hat{B}(f, g)(\xi)|^2 \leq \hat{B}(1, 1)(\xi) \hat{B}(|f|^2, |g|^2)(\xi).
\]
对 $\xi$ 积分, 我们得到
\begin{equation}
\|B(f, g)\|^2_{L^2(\mathbf{R}^3)} \lesssim A \|f\|^2_{L^2(\mathbb{S}^2)} \|g\|^2_{L^2(\mathbb{S}^2)}, \label{eq:c6-301}
\end{equation}
其中
\begin{equation}
A = \sup_\xi |\hat{B}(1, 1)(\xi)| = \sup_\xi \int \delta(1 - |\xi - \eta|) \delta(1 - |\eta|) \mathrm{d}\eta.\label{eq:c6-302}
\end{equation}

因此, 在这种情形下证明定理只需验证 $A$ 有限.对 $A(\xi) = \hat{B}(1, 1)(\xi)$ 作显式计算是有用的.对任何维数 $n \geq 2$, 有:

\begin{lemma}\label{lem:c6-8.22}
\begin{equation}
A(\xi) = \int_{\mathbf{R}^n} \delta(1 - |\xi - \eta|) \delta(1 - |\eta|) \mathrm{d}\eta \simeq \frac{1}{|\xi|} \left(4 - |\xi|^2\right)_+^{\frac{n-3}{2}}.\label{eq:c6-303}
\end{equation}
\end{lemma}

\begin{proof}
\[
\begin{array}{c}
A(\xi) = \int \delta(1 - |\xi - \eta|) \delta(1 - |\eta|) \mathrm{d}\eta \simeq \int_{|\eta| = 1} \delta(1 - |\xi - \eta|^2) \mathrm{d}\sigma_\eta = \\
= \int_{|\eta| = 1} \delta(|\xi|^2 - 2\xi \cdot \eta) \mathrm{d}\sigma_\eta \simeq \frac{1}{|\xi|} \int_{|\eta| = 1} \delta\left(\frac{|\xi|}{2} - \frac{\xi}{|\xi|} \cdot \eta\right) \mathrm{d}\sigma_\eta.
\end{array}
\]
由旋转对称性, 我们可以假设 $\boldsymbol{\xi} = (|\xi|, 0, \ldots, 0)$, 于是
\[
\begin{array}{rcl}
A(\xi) &\simeq& \frac{1}{|\xi|} \int_0^\pi \delta\left(\frac{|\xi|}{2} - \cos\theta\right) (\sin\theta)^{n-2} \mathrm{d}\theta = \\
&=& \frac{1}{|\xi|} \int_{-1}^1 \delta\left(\frac{|\xi|}{2} - u\right) (1 - u^2)^{\frac{n-3}{2}} \mathrm{d}u = \frac{1}{|\xi|} \left(1 - \frac{|\xi|^2}{4}\right)^{\frac{n-3}{2}},
\end{array}
\]
当 $|\xi|/2 \in [-1, 1]$ 时成立.
\end{proof}

当 $n = 3$ 时, $A(\xi) \simeq 1/|\xi|$ 只在 $\xi = 0$ 处奇异, 但我们可以通过假设 $f$ 与 $g$ 的支集在 $\mathbb{S}^2$ 中某点的小邻域内来避免这个困难(回忆不失一般性, 我们可以把估计局部化到球面上的一个小帽上).于是 \eqref{eq:c6-302} 中的上确界只需对一切 $\xi \in \operatorname{supp}(f) + \operatorname{supp}(g)$ 取, 而这是一个远离 0 的有界集合.因此我们可以在 \eqref{eq:c6-302} 中限制 $|\xi| \ge C > 0$, 奇异性消失, 留下 $A < \infty$.

由双线性形式 $B(f, g)$ 的 $L^2$ 估计 \eqref{eq:c6-301}, 推出 Stein 算子 $\mathcal{S} f$ 的 $L^4$ 估计:
\[
\|\mathcal{S} f\|^2_{L^4(\mathbf{R}^3)} = \|B(f, f)\|_{L^2} \simeq A^{1/2} \|f\|^2_{L^2(\mathbb{S}^2)},
\]
其中假设 $f$ 的支集在球面上的一个小帽内.

在 $n = 2$ 的情形, 我们要的是 $\boldsymbol{S} f$ 的 $L^6$ 估计.由于 $6 = 3 \times 2$, 我们可以尝试用三线性形式重复同样的计算, $T(f, g, h) = Sf \cdot Sg \cdot Sh$, 并利用事实 $\|Sf\|^3_{L^6} = \|T(f, f, f)\|_{L^2}$.我们有
\[
\begin{array}{l}
\hat{T}(f, g, h)(\xi) \simeq (f \mathrm{d}\sigma) * (g \mathrm{d}\sigma) * (h \mathrm{d}\sigma)(\xi) = \\
\qquad = \iint_{\mathbb{R}^2 \times \mathbb{R}^2} \delta(1 - |\xi - \eta - \zeta|) \delta(1 - |\eta|) \delta(1 - |\zeta|) f(\xi - \eta) g(\eta) h(\zeta) \mathrm{d}\eta \mathrm{d}\zeta,
\end{array}
\]
并对测度 $\delta(1 - |\xi - \eta|) \delta(1 - |\eta|) \delta(1 - |\zeta|) \mathrm{d}\eta \mathrm{d}\zeta$ 应用 Cauchy-Schwarz, 我们得到
\[
|\hat{T}(f, g, h)(\xi)|^2 \leq \hat{T}(1, 1, 1)(\xi) \hat{T}(|f|^2, |g|^2, |h|^2)(\xi).
\]
对 $\xi$ 积分, 我们得到
\begin{equation}
\|T(f, g, h)\|^2_{L^2(\mathbb{R}^2)} \lesssim A \|f\|^2_{L^2(\mathbb{S}^1)} \|g\|^2_{L^2(\mathbb{S}^1)} \|h\|^2_{L^2(\mathbb{S}^1)}, \label{eq:c6-304}
\end{equation}
其中
\begin{equation}
A = \sup_\xi |\hat{T}(1, 1, 1)(\xi)| = \sup_\xi \iint \delta(1 - |\xi - \eta|) \delta(1 - |\eta|) \delta(1 - |\zeta|) \mathrm{d}\eta \mathrm{d}\zeta.\label{eq:c6-305}
\end{equation}

卷积结构允许我们把 $\xi$ 限制在集合 $\operatorname{supp} f + \operatorname{supp} g + \operatorname{supp} h$ 上, 如果我们假设 $f, g, h$ 的支集在球面的一个小帽内, 就可以假设 $1 \leq |\xi| \leq 3$.利用\lemref{lem:c6-8.22}, 我们可以计算 $T(1, 1, 1)$ 并证明 $A$ 有界,
\[
\begin{array}{rl}
& T(1, 1, 1)(\xi) = \int B(1, 1)(\xi - \zeta) \delta(1 - |\zeta|) \mathrm{d}\zeta \sim \\
& \qquad \sim \int_{|\xi - \zeta| < 2} \frac{\delta(1 - |\zeta|)}{(4 - |\xi - \zeta|^2)^{1/2}} \mathrm{d}\zeta = \int_{|\xi - \zeta| < 2}^{\zeta \in \mathbb{S}^1} \frac{\mathrm{d}\sigma_\zeta}{(3 - 2\xi \cdot \zeta + |\xi|^2)^{1/2}} \simeq \\
& \qquad \simeq \int_{a(\xi)}^1 \frac{\mathrm{d}a}{(3 - |\xi|^2 + 2|\xi| a)^{1/2} (1 - a^2)^{1/2}} \sim \int_{a(\xi)}^1 \frac{\mathrm{d}a}{(a - a(\xi))^{1/2} (1 - a)^{1/2}} \simeq 1,
\end{array}
\]
其中 $a(\xi) = -\frac{3 - |\xi|^2}{2|\xi|}$.由三线性形式 $T(f, g, h)$ 的 $L^2$ 估计 \eqref{eq:c6-304}, 推出 Stein 算子 $\mathcal{S} f$ 的 $L^6$ 估计:
\[
\|\mathcal{S} f\|^3_{L^6(\mathbb{R}^2)} = \|T(f, f, f)\|_{L^2} \simeq A^{1/2} \|f\|^3_{L^2(\mathbb{S}^2)}.
\]

我们也可以尝试对 $n = 2$ 重复双线性论证.如前, 对 $B(f, g) = Sf \cdot Sg$, 有
\[
|\hat{B}(f, g)(\xi)|^2 \leq \hat{B}(1, 1)(\xi) \hat{B}(|f|^2, |g|^2)(\xi).
\]
对 $\xi$ 积分, 并利用\lemref{lem:c6-8.22} 计算 $\hat{B}(1, 1)$,
\[
\|B(f, g)\|^2_{L^2(\mathbb{R}^2)} \lesssim \iint \frac{\delta(1 - |\xi - \eta|) \delta(1 - |\eta|)}{|\xi| (4 - |\xi|^2)^{1/2}} |f(\xi - \eta)|^2 |g(\eta)|^2 \mathrm{d}\eta \mathrm{d}\xi.
\]
作变量替换 $\xi \to \zeta = \xi - \eta$, 注意到当 $|\eta| = |\zeta| = 1$ 时, 有
\[
\begin{array}{c}
|\xi| = |\eta + \zeta| \simeq (1 + \eta \cdot \zeta)^{1/2}, \\
(4 - |\xi|^2)^{1/2} = (4 - |\eta + \zeta|^2)^{1/2} \simeq (1 - \eta \cdot \zeta)^{1/2},
\end{array}
\]
因此
\begin{equation}
\|B(f, g)\|^2_{L^2(\mathbb{R}^2)} \lesssim \iint_{\mathbb{S}^1 \times \mathbb{S}^1} \frac{|f(\zeta)|^2 |g(\eta)|^2}{(1 - (\eta \cdot \zeta)^2)^{1/2}} \mathrm{d}\sigma_\eta \mathrm{d}\sigma_\zeta.\label{eq:c6-306}
\end{equation}

这是一个有趣的公式.注意, 如果 $f$ 与 $g$ 在 $\mathbb{S}^1$ 上的支集投影不相交, 即不含同一方向上的点, 则量 $1 - (\eta \cdot \zeta)^2$ 有正的下界, 此时我们得到双线性限制估计
\[
\|B(f, g)\|_{L^2(\mathbb{R}^2)} \lesssim \|f\|_{L^2(\mathbb{S}^1)} \|g\|_{L^2(\mathbb{S}^1)}.
\]

我们还可以考虑其他类型的双线性形式, 它们具有特殊的结构, 可以消去分母中的奇异性.例如取 $Q(f, g) = \partial_1 Sf \partial_2 Sg - \partial_2 Sf \partial_1 Sg$, 则取 傅里叶变换并如前进行, 我们看到
\[
\begin{array}{rcl}
\|Q(f, g)\|^2_{L^2(\mathbb{R}^2)} &\lesssim& \iint_{\mathbb{S}^1 \times \mathbb{S}^1} \frac{|\eta_1 \zeta_2 - \eta_2 \zeta_1|^2}{(1 - (\eta \cdot \zeta)^2)^{1/2}} |f(\zeta)|^2 |g(\eta)|^2 \mathrm{d}\sigma_\eta \mathrm{d}\sigma_\zeta \\
&\lesssim& \|f\|^2_{L^2(\mathbb{S}^1)} \|g\|^2_{L^2(\mathbb{S}^1)},
\end{array}
\]
因为我们有恒等式 $|\eta_1 \zeta_2 - \eta_2 \zeta_1|^2 = 1 - (\eta \cdot \zeta)^2 \leq 1$.


\part{Minkowski 空间中波动方程的分析}

% ==== part07 第 7 章 衰减估计 (Decay estimates) ====
% 源: MinerU_latex_lecture nots.tex 第 12107--12484 行;pdf_text.txt 第 9839--10174 行(PDF 第 181--186 页)
%
% 术语(新增,键=中文=英文):
%   pseudodifferential-operator=伪微分算子、commuting-vectorfields=交换向量场、
%   scaling-vector-field=伸缩向量场、inverted-translations=倒置平移、boost=加速算子、
%   spherical-means=球平均、wave-zone=波区
%   (注: 所附论文部分 partgf1 已用同键;倒置平移与附录 A 部分所用"逆平移"译名不一致,待协调)
%
% REFS(需跨章/待解析项):
%   - "见附录 4.2"(Killing 与共形 Killing 向量场族) -> 目标: 附录 A(partapx)相应小节,标签待解析
%   - "局部存在定理"(Remark 0.23) -> 原书所在章未知,正文纯文本
%   - Remark 0.27 与 Prop 0.30 证明中的 "Proposition 0.26" 系原书对 Theorem 0.26 的误称,
%     译文按原文写"命题 0.26",数字用 \ref{thm:c7-0.26} 可点击跳转
%   - 原书 (312) 后 "slightly larger than W^{(n+1)/2}, 1." 中孤立 ", 1." 疑为排版残留,已省略
%   - 原书 (318)/(320) 分母为 t-|x|,无外绝对值号(与 (316)/(317) 的 (1+|t-|x||) 不同),照原书
%   - 原书 "[?]" 文献占位符缺失,保留 "[?]"
%   - 原书 "and and"、"repeted" 等笔误按正常语义翻译;源 OCR "Proposition 0.28. Let ⋆φ = 0"
%     按 PDF 还原为 □φ = 0;(311) 中 f'(ξ)、g'(ξ) 按 PDF 还原为 f̂(ξ)、ĝ(ξ)
%   - 原书 (316) 后 "satisfy 315" 无括号,按上下文作 \eqref{eq:c7-315}
%
% 编号说明:
%   - 本章无节,编号对象显示 "0.x";原书第一个编号对象为 Remark 0.23(0.1--0.22 在原书中不存在),
%     故在 \section 后加 \setcounter{mtthm}{22} 使显示与原书一致(实测本模板 \section 会把
%     subsection 与 theorem 计数器清零,若不设置则首对象会显示 0.1)
%   - 原书第 8 章编号接续本章(第 8 章首个编号为 0.30.1/0.31),part08 需在 \section 后续接计数器
%   - 公式 (307)--(325) 用 \begin{equation} 自动编号(本模板 [ \] 不编号,与 partapx 同法);
%     其余行间公式用 [ \] 不占号
% ====

\chapter{衰减估计}

\setcounter{mtthm}{22}% 原书本章从 Remark 0.23 开始(0.1--0.22 不存在)

考虑 Minkowski 空间 $\mathbb{R}^{n+1}$ 中的标准波动方程

\begin{equation}
\Box\phi = 0.\label{eq:c7-307}
\end{equation}

$\mathbb{R}^{n+1}$ 中的标准惯性坐标记为 $x^{\mu}$, $\mu = 0, 1, \ldots, n$, 相对于这些坐标, Minkowski 度量取对角形式 $\mathbf{m}_{\mu\nu} = \mathrm{diag}(-1, 1, \ldots, 1)$.我们有 $x^0 = t$, 而 $x = (x^1, \ldots, x^n)$ 表示空间坐标.我们采用关于重复指标的标准求和约定, 以及关于向量和张量的指标升降的相应约定.特别地, 若 $x_{\mu} = m_{\mu\nu}x^{\nu}$, 则 $x_0 = -t$, $x_i = x^i$, $i = 1, \ldots, n$.我们用 $\Sigma_{t_0}$ 表示类空超平面 $t = t_0$.波算子定义为
$\Box = \mathbf{m}^{\alpha\beta}\partial_{\alpha\beta} = -\partial_t^2 + \sum_i\partial_i^2$.
我们研究如下初值问题,

\begin{equation}
\phi(0, x) = f(x), \quad\partial_t\phi(0, x) = g(x)\label{eq:c7-308}
\end{equation}

为方便起见, 记 $\phi[0] = (f, D^{-1}g)$, 其中 $D^{-1}$ 是以 $|\xi|^{-1}$ 为符号的\newterm{pseudodifferential-operator}{伪微分算子}{pseudodifferential operator}.设

\begin{equation}
E[\phi](t) = \int_{\Sigma_t}\left(|\partial_t\phi|^2 + \sum_i|\partial_i\phi|^2\right)dx\label{eq:c7-309}
\end{equation}

为 $\phi$ 在时刻 $t$ 的总能量.能量的守恒律为

\begin{equation}
E[\phi](t) = E[\phi](0)\label{eq:c7-310}
\end{equation}

作为推论, 对所有 $s\geq 0$ 有能量不等式

\[
\left\|\partial\phi(t)\right\|_{H^s(\mathbb{R}^n)} \leq \left\|\partial\phi(0)\right\|_{H^s(\mathbb{R}^n)}.
\]

能量恒等式既可用前面各节讨论过的几何技巧证明, 其中只涉及分部积分, 也可用 Fourier 方法证明, 后者使用 Plancherel 公式以及如下 Fourier 表示公式,

\begin{equation}
\phi(t, x) = (2\pi)^{-n}\int e^{ix\cdot\xi}\Big(\cos t|\xi|\, \hat{f}(\xi) + \frac{\sin t|\xi|}{|\xi|}\hat{g}(\xi)\Big)d\xi\label{eq:c7-311}
\end{equation}

\begin{remark}\label{rem:c7-0.23}
标准 Sobolev 嵌入 $H^s(\mathbb{R}^n)\subset L^{\infty}(\mathbb{R}^n)$($s > \frac{n}{2}$)使我们无需使用显式表示即可得到 \eqref{eq:c7-307} 的解的 $L^{\infty}$ 界.这一方法可推广到非线性方程, 并在局部存在定理的证明中起重要作用.
\end{remark}

\begin{proposition}[色散不等式]\label{prop:c7-0.24}
\eqref{eq:c7-307}、\eqref{eq:c7-308} 的解满足,

\begin{equation}
\|\phi(t)\|_{L^{\infty}} \leq c|t|^{-\frac{n-1}{2}}\|\phi[0]\|_{B_{1, 1}^{\frac{n+1}{2}}}\label{eq:c7-312}
\end{equation}

其中 $B_{1, 1}^{\frac{n+1}{2}}$ 是比 $W^{\frac{n+1}{2}}$ 稍大的 Besov 空间.更确切地说,

\[
\|f\|_{B_{1, 1}^{\frac{n+1}{2}}} \approx \sum_{k\in\mathbb{Z}}2^{k\frac{n+1}{2}}\|P_k f\|_{L^1(\mathbb{R}^n)}
\]
\end{proposition}

\noindent\textbf{习题.}\quad 证明不等式 \eqref{eq:c7-312} 可由其频率局部化版本推出.换言之, 证明只需验证如下不等式,

\begin{equation}
\|\phi(t)\|_{L^{\infty}} \leq c|t|^{-\frac{n-1}{2}}\|\phi[0]\|_{L^1}\label{eq:c7-313}
\end{equation}

其中初值 $f, g$ 的 傅里叶变换局部化于 $\frac{1}{2}\leq|\xi|\leq 2$.

\begin{proof}
\eqref{eq:c7-312} 的标准证明基于对表示 \eqref{eq:c7-311} 应用稳定相方法.在奇数维中, 可以利用解的\newterm{spherical-means}{球平均}{spherical means}表示证明色散估计的一个相关形式.这在 $n = 3$ 时特别容易.我们将在后面讨论 \eqref{eq:c7-312} 的一个不依赖任何表示公式的推导.
\end{proof}

\begin{remark}\label{rem:c7-0.25}
色散不等式提供两类信息.第一类涉及 $\|\phi(t)\|_{L^{\infty}}$ 当 $t\to\infty$ 时的精确衰减率, 第二类提供 $\|\phi(t)\|_{L^{\infty}}$ 对 $t>0$ 的正则性性质的信息.就改善正则性而言, 与 Sobolev 嵌入 $L^{\infty}(\mathbb{R}^n)\subset W^{1, n}(\mathbb{R}^n)$ 相比, 估计 \eqref{eq:c7-312} 对 $t>0$ 增益了 $\frac{n-1}{2}$ 阶导数.
\end{remark}

在许多应用中, 尤其是对非线性方程的应用中, \eqref{eq:c7-312} 并不十分有用.推导波动方程解的渐近性质的更有效方法, 是基于广义能量估计, 它由\newterm{commuting-vectorfields}{交换向量场}{commuting vectorfields}方法结合整体 Sobolev 不等式得到.下面我们回顾用于推导上述衰减率估计的交换向量场方法.其思想是将能量恒等式 \eqref{eq:c7-310} 与和波算子交换的向量场以及经典 Sobolev 不等式的整体形式结合使用.我们请读者参阅 [?] 与 [?] 以了解细节.

Minkowski 时空 $\mathbb{R}^{n+1}$ 配备一族 Killing 与共形 Killing 向量场(见附录 4.2):平移 $\mathbf{T}_{\mu} = \partial_{\mu}$、Lorentz 旋转 $\mathbf{L}_{\mu\nu} = x_{\mu}\partial_{\nu} - x_{\nu}\partial_{\mu}$、\newterm{scaling-vector-field}{伸缩向量场}{scaling vector field} $\mathbf{S} = t\partial_t + x^i\partial_i$ 以及\newterm{inverted-translations}{倒置平移}{inverted translations} $\mathbf{K}_{\mu} = -2x_{\mu}\mathbf{S} + \langle x, x\rangle\partial_{\mu}$.回顾 $x^{\mu}$ 表示标准变量 $x^0 = t$, $x^1, \ldots, x^n$, 而 $x_{\mu} = m_{\mu\nu}x^{\nu}$.Killing 向量场 $\mathbf{T}_{\mu}$ 与 $\mathbf{L}_{\mu\nu}$ 与 $\Box$ 交换, 而 $\mathbf{S}$ 保持解空间, 其含义是:$\Box\phi = 0$ 蕴含 $\Box S\phi = 0$, 因为 $[\Box, S] = 2\Box$.我们可以把算子 $\mathbf{L}_{\mu\nu}$ 分解为角旋转算子 $^{(ij)}\mathbf{O} = x_i\partial_j - x_j\partial_i$ 与\newterm{boost}{加速算子}{boost} $^{(i)}\mathbf{L} = x_i\partial_t + t\partial_i$, 其中 $i, j, k = 1, \ldots, n$.回顾 \eqref{eq:c7-309} 中的能量表达式.基于上述交换性质, 我们定义如下``广义能量'':

\begin{equation}
E_k[\phi] = \sum_{X_{i_1}, \ldots, X_{i_j}}E[X_{i_1}X_{i_2}\cdots X_{i_j}\phi]\label{eq:c7-314}
\end{equation}

其中求和取遍 $0\leq j\leq k$ 以及所有 Killing 向量场 $\mathbf{T}$、$\mathbf{L}_{\mu\nu}$ 以及伸缩向量场 $\mathbf{S}$.交换向量场方法的关键在于, 量 $E_k$, $k\geq 1$, 被 \eqref{eq:c7-307} 的解保持.因此, 若

\begin{equation}
\sum_{0\leq k\leq s}\int(1+|x|)^{2k}\left(|\nabla^{k+1}f(x)|^2 + |\nabla^k g(x)|^2\right)dx \leq C_s < \infty\label{eq:c7-315}
\end{equation}

则对所有 $t$ 有 $E_s[\phi](t)\leq C_s$.\eqref{eq:c7-307} 的解的期望衰减估计现在可以从如下整体形式的 Sobolev 不等式(见 [?]、[?])导出:

\begin{theorem}[全局 Sobolev 估计]\label{thm:c7-0.26}
设 $\phi$ 是 $\mathbb{R}^{n+1}$ 中的任意函数, 使得对某个整数 $s > \frac{n}{2}$, $E_s[\phi]$ 有限.则

\begin{equation}
|\partial\phi(t, x)| \lesssim (1+t+|x|)^{-\frac{n-1}{2}}(1+|t-|x||)^{-\frac{1}{2}}\sup_{0\leq t'\leq t}E_s[\phi](t')\label{eq:c7-316}
\end{equation}

对所有 $t>0$ 成立.因此, 若 \eqref{eq:c7-307} 中的数据 $f, g$ 满足 \eqref{eq:c7-315}, 且 $s > \frac{n}{2}$, 则对所有 $t\geq 0$,

\begin{equation}
|\partial\phi(t, x)| \lesssim \frac{1}{(1+t+|x|)^{\frac{n-1}{2}}(1+|t-|x||)^{\frac{1}{2}}}\label{eq:c7-317}
\end{equation}
\end{theorem}

\begin{remark}\label{rem:c7-0.27}
显然, 该估计(其证明纯粹是几何的\footnote{特别是, 它不需要任何解的显式表示.})蕴含色散不等式 \eqref{eq:c7-312} 给出的衰减性质.事实上, 它在\newterm{wave-zone}{波区}{wave zone} $|x|\sim t$ 之外提供了更多信息, 这与线性方程 $\Box\phi = 0$ 的预期传播性质非常吻合.另一方面, 由于 \eqref{eq:c7-316} 确实是 Sobolev 不等式的整体形式, 命题 0.26 的估计似乎与 \eqref{eq:c7-312} 的改进正则性特征无关.然而事实并非完全如此, 我们将在后面看到.
\end{remark}
\begin{proof}
下面我们只概述证明的主要思想.考虑标准零对 $L_{\pm} = \partial_t\pm\partial_r$, 相应的零标架 $e_1, \ldots, e_{n-1}$, $e_n = L_-$, $e_{n+1} = L_+$, 以及角向量场 $A_i = \partial_i - \frac{x_i}{r}\partial_r$.显然,

\[
\sum_i|A_i\phi| \lesssim |\nabla\!\!\!/\phi| \lesssim \sum_i|A_i\phi|.
\]

其中 $|\nabla\!\!\!/\phi|^2 = \sum_{i=1}^{n-1}|e_i(\phi)|^2$.此外,

\[
|\partial_r\phi| + \sum_i|A_i\phi| \lesssim |\nabla\phi| \lesssim |\partial_r\phi| + \sum_i|A_i\phi|.
\]

我们还可以容易地验证如下简单的代数恒等式,

\[
\begin{aligned}
\frac{1}{2}\big((t+r)L_+ + (t-r)L_-\big) &= \mathbf{S} \\
\frac{1}{2}\big((t+r)L_+ - (t-r)L_-\big) &= \sum_i\frac{x_i}{|x|}\mathbf{L}_i \\
tA_i &= L_i - \frac{x_i}{|x|}\sum_j\frac{x_j}{|x|}L_j \\
t\mathbf{O}_{ij} &= x_i\mathbf{L}_j - x_j\mathbf{L}_i
\end{aligned}
\]

由前两个恒等式我们容易推出,

\begin{equation}
\begin{aligned}
|L_+\phi(t, x)| &\lesssim \frac{1}{t}|\Gamma\phi(t, x)| \\
|L_-\phi(t, x)| &\lesssim \frac{1}{t-|x|}|\Gamma\phi(t, x)|
\end{aligned}\label{eq:c7-318}
\end{equation}

其中 $|\Gamma\phi| = |\mathbf{S}\phi| + |\mathbf{L}\phi|$.

\begin{equation}
|\nabla\!\!\!/\phi(t, x)| \lesssim \frac{1}{t}|\Gamma\phi(t, x)|.\label{eq:c7-319}
\end{equation}

显然我们还有,

\[
|\partial\phi(t, x)| \lesssim \frac{1}{t-|x|}|\Gamma\phi(t, x)|
\]

或更一般地,

\begin{equation}
|\partial^N\phi(t, x)| \lesssim \frac{1}{(t-|x|)^N}|\Gamma^N\phi(t, x)|\label{eq:c7-320}
\end{equation}

其中 $|\Gamma^N\phi| = \sum|\Gamma_1\cdots\Gamma_N\phi|$, $\Gamma_1, \ldots, \Gamma_N$ 是向量场 $S, L_1, \ldots, L_n$ 中的任意一个.

将上述不等式与我们范数的定义结合, 我们推出

\[
\begin{aligned}
t\|E_+\phi(t)\|_{L^2} &\lesssim \|\Gamma\phi(t)\|_{L^2(\mathbb{R}^n)} \\
t\|\nabla\!\!\!/\phi(t)\|_{L^2} &\lesssim \|\Gamma\phi(t)\|_{L^2(\mathbb{R}^n)} \\
\|uE_-\phi(t)\|_{L^2} &\lesssim \|\Gamma\phi(t)\|_{L^2(\mathbb{R}^n)}
\end{aligned}
\]

其中 $u = |t-|x||$.剩下的任务是从上述 $L^2$ 估计推出 sup 范数估计.
\end{proof}

\begin{proposition}\label{prop:c7-0.28}
设 $\Box\phi = 0$, 其初值满足上述假设.则对所有 $t\geq 0$, $s > \frac{n}{2}$,

\begin{equation}
\|\phi(t)\|_{L^{\infty}} \lesssim \left(\frac{1}{t}\right)^{\frac{n-1}{2}}\|\Gamma^s\phi(t)\|_{L^2(\mathbb{R}^n)}\label{eq:c7-321}
\end{equation}

\begin{equation}
\|(1+|u|)^k\partial^k\phi(t)\|_{L^{\infty}} \lesssim \left(\frac{1}{t}\right)^{\frac{n-1}{2}}\|\Gamma^s\phi(t)\|_{L^2(\mathbb{R}^n)}\label{eq:c7-322}
\end{equation}

此外,

\[
\begin{aligned}
\|E_+\phi(t)\|_{L^{\infty}} &\lesssim \left(\frac{1}{t}\right)^{\frac{n+1}{2}}\|\Gamma^{s+1}\phi(t)\|_{L^2(\mathbb{R}^n)} \\
\|\nabla\!\!\!/\phi(t)\|_{L^{\infty}} &\lesssim \left(\frac{1}{t}\right)^{\frac{n+1}{2}}\|\Gamma^{s+1}\phi(t)\|_{L^2(\mathbb{R}^n)} \\
\|(1+|u|)E_-\phi(t)\|_{L^{\infty}} &\lesssim \left(\frac{1}{t}\right)^{\frac{n-1}{2}}\|\Gamma^{s+1}\phi(t)\|_{L^2(\mathbb{R}^n)}
\end{aligned}
\]
\end{proposition}

证明基于如下引理.

\begin{lemma}\label{lem:c7-0.29}
设 $u(x)$ 是 $\mathbb{R}^n$($n\geq 2$)上的光滑、紧支撑函数.我们有,

\begin{equation}
|u(x)| \leq C\frac{1}{|x|^{n-1}}\left(\|\partial_r u\|_{L^1} + \|(r\nabla\!\!\!/)^{n-1}\partial_r u\|_{L^1}\right)\label{eq:c7-323}
\end{equation}
\end{lemma}

\begin{proposition}[交换向量场方法]\label{prop:c7-0.30}
交换向量场方法蕴含色散不等式 \eqref{eq:c7-312}(见 [?]).
\end{proposition}

\begin{proof}
不失一般性, 可假设 $\partial_t\phi = g = 0$, 且 $f = \phi(0)$ 的 傅里叶变换支集位于壳层 $\frac{\lambda}{2}\leq|\xi|\leq 2\lambda$ 内, 其中 $\lambda\in 2^{\mathbf{N}}$.通过简单的伸缩论证, 我们事实上可以假设 $\lambda = 1$.对 Fourier 支集限制在 $1/2\leq|\xi|\leq 2$ 内的这类初值条件, 只需证明

\[
\|\phi(t)\|_{L^{\infty}} \lesssim (1+|t|)^{-\frac{n-1}{2}}\|f\|_{L^1(\mathbb{R}^n)}.
\]

由于 $\hat{\phi}$——$\phi$ 关于空间变量 $x$ 的 傅里叶变换——的支集也位于同一壳层内, 只需证明对 $\nabla\phi$ 的估计即可, 而非对 $\phi$ 本身.

接下来我们用中心位于整数坐标点 $I\in\mathbb{Z}^n$ 的圆盘 $D_I$ 覆盖 $\mathbb{R}^n$, 使得每个 $D_I$ 至多与有限个 $c_n$ 个圆盘 $D_J$ 相交, 其中 $c_n$ 只依赖于维数 $n$.考虑一个光滑的单位分解 $(\chi_I)_{I\in\mathbb{Z}^n}$, 满足 $\operatorname{supp}\chi_I\subset D_I$, 且每个 $\chi_I$ 非负.显然我们可以安排, 对所有 $k$,

\begin{equation}
\sum_{I\in\mathbb{Z}^n}|\nabla^k\chi_I(x)| \leq C_{k, n}\label{eq:c7-324}
\end{equation}

关于 $x\in\mathbb{R}^n$ 一致成立, 其中常数 $C_{n, k}$ 只依赖于 $n$ 与 $k$.现在设 $f_I = \chi_I\cdot f$, 并令 $\phi_I$ 为 \eqref{eq:c7-307} 的以 $\phi_I(0) = f_I$, $\partial_t\phi_I(0) = 0$ 为数据的相应解.显然 $f = \sum_I f_I$, $\phi = \sum_I\phi_I$.只需证明对所有 $I$,

\begin{equation}
\|\nabla^k\phi_I(t)\|_{L^{\infty}} \lesssim (1+t)^{-\frac{n-1}{2}}\sum_{j=0}^{n+k+1}\|D^j f_I\|_{L^1}\label{eq:c7-325}
\end{equation}

事实上, 若 \eqref{eq:c7-325} 成立, 我们容易推出

\[
\|\nabla^k\phi(t)\|_{L^{\infty}} \leq \sum_I\|\nabla^k\phi_I(t)\|_{L^{\infty}} \lesssim (1+t)^{-\frac{n-1}{2}}\sum_{j=0}^{n+k+1}\sum_I\|D^j f_I\|_{L^1}.
\]

根据 \eqref{eq:c7-324} 我们有,

\[
\begin{aligned}
\sum_I\|D^j f_I\|_{L^1} &= \sum_I\|D^j(\chi_I f)\|_{L^1} \lesssim \sum_I\sum_{0\leq i\leq j}\int_{\mathbb{R}^n}|D^i\chi_I(x)||D^{j-i}f(x)|dx \\
&= \sum_{0\leq i\leq j}\int_{\mathbb{R}^n}\Big(\sum_I|D^i\chi_I(x)|\Big)|D^{j-i}f(x)|dx \\
&\lesssim \sum_{0\leq i\leq j}c_{i, n}\|D^{j-i}f\|_{L^1(\mathbb{R}^n)} \lesssim \|f\|_{L^1(\mathbb{R}^n)}.
\end{aligned}
\]

因此,

\[
\|\nabla^k\phi(t)\|_{L^{\infty}} \lesssim (1+t)^{-\frac{n-1}{2}}\|f\|_{L^1(\mathbb{R}^n)},
\]

如所期望.

于是只需验证 \eqref{eq:c7-325}.不失一般性, 通过空间平移, 可假设 $I = 0$.将命题~\ref{thm:c7-0.26} 应用于 $\psi = \nabla\phi_0$, 对 $s_*$——严格大于 $\frac{n}{2}$ 的第一个整数——推得

\[
\begin{aligned}
\|\psi(t)\|_{L^{\infty}} &\leq c(1+t)^{-\frac{n-1}{2}}E_{s_*}[\phi_0](t) \\
&\leq c(1+t)^{-\frac{n-1}{2}}E_{s_*}[\phi_0](0).
\end{aligned}
\]

由于 $\phi_0$ 的支集包含在以原点为中心、半径为 1 的球内, 我们有

\[
E_{s_*}[\phi_0](0) \leq C_n\sum_{j=0}^{s_*+1}\|D^j f_0\|_{L^2}.
\]

最后, 根据 $\mathbb{R}^n$ 中的标准 Sobolev 不等式 $\|f\|_{L^2} \leq c\|\nabla^{\frac{n}{2}}f\|_{L^1}$, 我们得出结论:

\[
\|\psi(t)\|_{L^{\infty}} \leq c(1+t)^{-\frac{n-1}{2}}\sum_{j=0}^{n+2}\|D^j f_0\|_{L^1},
\]

如所期望.
\end{proof}


% ==== part08 第 8 章:Strichartz 不等式 ====
% 源: MinerU_latex_lecture nots.tex 12485-13555 + MinerU_latex_A-非.tex 78-301
%     (pdf_text.txt 10175-11161, PDF 第 187-204 页)
% 术语(新增, 键=中文=英文):
%   strichartz-inequality=Strichartz 不等式=Strichartz inequality;
%   admissible-wave-pair=容许波对=admissible wave pair;
%   dispersive-inequality=色散不等式=dispersive inequality;
%   truncated-cone-operator=截断锥算子=truncated cone operator;
%   retarded-operator=推迟算子=retarded operator;
%   endpoint-case=端点情形=endpoint case. 
%   其余按普通中文书写(可能由其他部分首次引入, 待协调): 限制定理、基本解、能量估计、
%   二进分解、LP(Littlewood-Paley)投影、光锥/零锥、球对称、自由波、演化算子、
%   双线性型、Whitney 分解、Hilbert 变换、截断函数、Duhamel 原理、Knapp 反例、
%   Sobolev 嵌入、Plancherel 公式、Hausdorff--Young 不等式、Hardy--Littlewood--Sobolev 不等式. 
% REFS(跨章/待解析项):
%   - 源 12509 "eqrefeq:genwave"(原书破损引用)→ 正文写 "上述方程"(即 (326)-(327)), 待核实;
%   - 源 12675 "see section on the fundamental solution of □ in R^{3+1}" → 正文写 "见 R^{3+1} 中 □ 的基本解那一节", 所在章待解析;
%   - 源 12944 "corollary 6.15" → 正文写 "推论 6.15", 目标: 第 6 章(源 10796 行), 待解析;
%   - 源 13098 "see lemma 8.8" → 正文写 "引理 8.8", 目标: 限制定理部分;【已核对: 译本限制定理部分无编号 8.8 的引理(仅 lemma 0.48/0.56), 原书 Lemma 8.8 在译本中编号待人工确认, 暂保留纯文本】
%   - 源 13118 "the argument following formula (300)" → 正文写 "公式 (300) 之后的论证", 目标: 限制定理部分, 待解析;
%   - 源 13054 "We only need to check (??)" → 正文写 "只需验证 (357)", 待核实;
%   - 源 13300 "due to Hörmander,  [?]" → 正文保留 "[?]", 文献编号破损;
%   - A-非 97 行 "(see [] and [])"(Christ-Kiselev 文献)→ 正文保留 "[]", 待补文献编号;
%   - 图 1: 源 images/41a61c93816f45bd38d116772bd400980527b33ed71fd417f933227a865de530.jpg
%     缺失, 暂以文字占位(REFS: 待补图). 
% OCR/数学修正(对照 pdf_text.txt, 均按数学含义处理;原书/OCR 疑误处照排并注明):
%   - 源 12508-12509 波算子: "\Xi = -\partial_t^2 u + \Delta" → "\Box = -\partial_t^2 + \Delta"(去掉多余 u, □ 补回);
%   - 源 12580 定理 0.32(2) 维数条件 "2/q" → "1/q"(与 (333) 及脚注 2 "γ=1/2" 一致);
%   - 源 12966 截断函数支集 "1.2 ≤ |ξ| ≤ 2" → "1/2 ≤ |ξ| ≤ 2"(与命题 0.40 及二进环一致);
%   - 源 13115 γ(r) 定义负号: "γ(r) = -(n-1)/2(1-2/r)" → "γ(r) = (n-1)/2(1-2/r)"
%     (否则与下文 "0<γ(r)<1" 及 HLS 应用矛盾, 原书笔误);
%   - 源 13178 定理 0.42 估计 (363) 右端两项均写 ||f||(原书如此, 疑第二项应为 ||g||, 照排未改);
%   - 源 12994 CC*F: 已核对(新版 PaddleOCR 确认原书用 χ, 译文正文亦用 χ(见 eq:c8-349/353), 一致; 旧源误作 β, 无需改动).
%   - 源 13082: 已核对(新版 PaddleOCR 确认原书作 χ_μ, 译文正文亦作 χ_μ(见 eq:c8-357 附近), 一致; 旧源误作 β, 无需改动).
%   - 源 12866 条带 D_T 定义后多余 "2" 为 OCR 噪音, 已删. 
% 编号说明:
%   - 第 8 章无节, 所有编号对象显示 "0.x", 共享 theorem 计数器;
%   - 第一个编号子标题 "0.30.1 齐次波动方程" 用 \mtsubsecn{0.30.1}{...}(不推进计数器), 
%     其后 \setcounter{mtthm}{30} 使 Definition 0.31 编号正确;
%   - 带编号子标题 0.36/0.37/0.38/0.44/0.46/0.49/0.54 用 \mtsubsec 自然推进计数器;
%   - "0.56.1 具有受限核的积分算子" 用 \mtsubsecn{0.56.1}{...}, 其后 \setcounter{mtthm}{56}
%     保证 Proposition 0.57 编号正确;
%   - 公式 (326)-(379) 在 (a) 段、(380)-(383) 在 (b) 段, 全部用 equation 环境自动编号, 
%     与上一部分(part07)结束于 (325) 相接;未编号行间公式用不编号显示公式. 
%   - 脚注 1-7 全部按 pdf_text.txt 对应页补全并翻译. 
% ====

\chapter{Strichartz 不等式}

\newterm{strichartz-inequality}{Strichartz 不等式}{Strichartz inequality}是研究线性与非线性波动方程的重要工具. 它们与限制定理密切相关. 本章我们只考虑标准线性波动方程的情形. 然而, 类似的等式对线性色散方程也成立, 如 Schrödinger 方程、线性 KdV 方程等. 

\mtsubsecn{0.30.1}{齐次波动方程}

考虑方程
\begin{equation}
\Box u=F, \label{eq:c8-326}
\end{equation}
\begin{equation}
u(0, x)=f(x), \qquad \partial_t u(0, x)=g(x), \label{eq:c8-327}
\end{equation}
的解 $u=u(t, x)$, 其中 $t\in\mathbb{R}$, $x\in\mathbb{R}^n$, $\Box$ 为波算子 $\Box=-\partial_t^2+\Delta$. 显然, 上述方程的解可以写成齐次波动方程
\begin{equation}
\Box u=0, \label{eq:c8-328}
\end{equation}
的解(它在 $t=0$ 时刻满足初始条件 \eqref{eq:c8-327})与纯非齐次波动方程
\begin{equation}
\Box u=F, \label{eq:c8-329}
\end{equation}
的解(具有零初始数据
\[
u(0, x)=0, \qquad \partial_t u(x, 0)=0
\]
)的叠加. 

我们用 $W(t)h$ 表示齐次问题 \eqref{eq:c8-328} 的基本解, 即 $u(t, x)=\bigl(W(t)h\bigr)(x)$ 是 \eqref{eq:c8-328} 的满足初始条件
\[
u(0, x)=0, \qquad \partial_t u(0, x)=h(x)
\]
的唯一解. 由 Duhamel 原理, 非齐次方程的任意解本身也可以写成齐次方程解的叠加, 按如下公式:
\begin{equation}
u(t)=\int_0^t W(t-t')F(t')\, \mathrm{d}t', \label{eq:c8-330}
\end{equation}

在陈述本节的主要结果之前, 我们给出如下定义. 

\setcounter{mtthm}{30}

\begin{definition}\label{def:c8-0.31}
我们称实数对 $(q, r)$ 为\newterm{admissible-wave-pair}{容许波对}{admissible wave pair}, 如果它们满足条件
\[
\begin{array}{rcl}
q & \geq & 2, \\[2pt]
\dfrac{2}{q} & \leq & (n-1)\left(\dfrac{1}{2}-\dfrac{1}{r}\right), \\[4pt]
(q, r, n) & \neq & (2, \infty, 3).
\end{array}
\]
\end{definition}

我们现在准备陈述如下结果. 

\begin{theorem}\label{thm:c8-0.32}
设 $n\geq2$, 且 $(q, r)$ 是 $r<\infty$ 的波容许对\footnote{当 $r=\infty$ 时, 只要把下面估计左端的空间改为适当的 Besov 空间, 该情形也可以包含进来. }. 
\begin{enumerate}
\def\labelenumi{(\arabic{enumi})}
\item 假设维数条件
\[
\frac{1}{q}+\frac{n}{r}=\frac{n}{2}-\gamma .
\]
则若 $u$ 以初始条件 \eqref{eq:c8-327} 满足齐次方程 \eqref{eq:c8-328}, 
\begin{equation}
\|u\|_{L_t^q L_x^r}+\|u\|_{L_t^\infty\dot{H}_x^\gamma}+\|\partial_t u\|_{L_t^\infty\dot{H}_x^{\gamma-1}}
\lesssim\|f\|_{\dot{H}^\gamma}+\|g\|_{\dot{H}^{\gamma-1}}\label{eq:c8-331}
\end{equation}
\item 假设维数条件
\[
\frac{1}{q}+\frac{n}{r}=\frac{n}{2}-\gamma=\frac{1}{q'}+\frac{n}{r'}-2, 
\]
其中 $q'$ 是 $q$ 的对偶指数, $r'$ 是 $r$ 的对偶指数\footnote{因此事实上 $\gamma=1/2$. }. 则若 $u$ 以零初始条件满足纯非齐次问题 0.30.1, 那么在有限时间区间 $[0, T]$ 上有:
\begin{equation}
\|u\|_{L^q([0, T];L^r)}+\|u\|_{C([0, T];\dot{H}^\gamma)}+\|\partial_t u\|_{C([0, T];\dot{H}^{\gamma-1})}
\lesssim\|F\|_{L^{q'}([0, T];L^{r'})}\label{eq:c8-332}
\end{equation}
\item 对满足维数条件的容许对 $(q_1, r_1)$、$(q_2, r_2)$(其中 $r_1, r_2<\infty$), 我们还有 \eqref{eq:c8-332} 的如下更一般版本:
\[
\frac{1}{q_1}+\frac{n}{r_1}=\frac{n}{2}-\gamma=\frac{1}{q_2'}+\frac{n}{r_2'}-2 .
\]
则
\begin{equation}
\|u\|_{L^{q_1}([0, T];L^{r_1})}+\|u\|_{C([0, T];\dot{H}^\gamma)}+\|\partial_t u\|_{C([0, T];\dot{H}^{\gamma-1})}
\lesssim\|F\|_{L^{q_2'}([0, T];L^{r_2'})}\label{eq:c8-333}
\end{equation}
\end{enumerate}
\end{theorem}

\begin{remark}\label{rem:c8-0.33}
对 $n\geq4$, 容许指数的区域对应于平面 $(1/q, 1/r)$ 中以
$\mathrm{O}=(1/\infty, 1/\infty)$、$\mathrm{E}=(1/\infty, 1/2)$、$\mathrm{P}=\bigl(1/2, \frac{n-3}{2(n-1)}\bigr)$ 与 $\mathrm{Q}=(1/2, 1/\infty)$ 为顶点的四边形 $\mathrm{OEPQ}$. 当 $n=3$ 时, 点 $P$ 与 $Q$ 重合, 区域退化为三角形 $\mathrm{OEQ}$. 当 $n=2$ 时, 我们得到更小的三角形 $\mathrm{OEQ}_2$, 其中 $\mathrm{Q}_2=(1/4, 1/\infty)$. 

对 $n=3$, 三角形区域的边界除端点 $P$ 外全部允许. 对 $n\geq4$, 四边形区域的边界完全允许, 下面我们会指出这一点. 

有趣的情形是线段 $EP$ 上的点以及 $PQ$ 上靠近 $P$ 的点, 因为其余情形都可以利用 Sobolev 嵌入由这些情形推出. 点 $E$ 对应于能量估计. 当 $n=3$ 时, 存在排除点 $P$ 的反例, 而高维情形下 $P$ 的包含性最近由 Keel 与 Tao 获得 \cite{k-t}. 
\end{remark}

\begin{figure}[htbp]
\centering
% 原书图 1:Admissible exponents for n ≥ 4. 源文件引用的图片
% images/41a61c93816f45bd38d116772bd400980527b33ed71fd417f933227a865de530.jpg 缺失, 
% 暂以文字占位;若图片可用, 请把下面的 \fbox 替换为
% \includegraphics[width=0.6\linewidth]{images/41a61c93816f45bd38d116772bd400980527b33ed71fd417f933227a865de530.jpg}
\fbox{\parbox{0.6\linewidth}{\centering (原书此处为图:$n\geq4$ 时容许指数的四边形区域 $\mathrm{OEPQ}$)}}%
\caption{$n\geq4$ 时的容许指数}
\end{figure}

标准 Strichartz 估计\footnote{即 Strichartz 本人实际证明的那个估计. }对应于点 $S=\Bigl(\frac{n-1}{2(n+1)}, \frac{n-1}{2(n+1)}\Bigr)$. 

\begin{remark}\label{rem:c8-0.34}
我们指出, 尽管\newterm{endpoint-case}{端点情形}{endpoint case} $n=3, q=\infty, r=2$ 被禁止, 估计在球对称情形下仍然成立. 事实上, 设 $\phi$ 是 $\mathbb{R}^{3+1}$ 中齐次波动方程 $\Box\phi=0$ 的、满足初始条件
\[
\phi(0, x)=0, \qquad \partial_t\phi(0, x)=f(x)
\]
的解, 并假设 $f$ 是球对称的, 即 $f(x)=f(|x|)$. 则
\begin{equation}
\int_0^\infty\|\phi(t, \cdot)\|_{L^\infty(\mathbb{R}^3)}^2\, \mathrm{d}t\leq c\|f\|_{L^2}^2.\label{eq:c8-334}
\end{equation}
证明是 Hardy--Littlewood 极大定理\footnote{在区域 $r\leq t$ 中这是显然的, 而对 $r\geq t$ 论证是初等的. }的直接推论, 鉴于对球对称的 $f$, 
\[
\phi(x, t)=\frac{c}{|x|}\int_{|\, |x|-t\, |}^{|x|+t}\lambda f(\lambda)\, \mathrm{d}\lambda .
\]
\end{remark}

\begin{remark}\label{rem:c8-0.35}
下面给出一个初等例子, 说明端点结果 $n=3, q=\infty, r=2$ 在一般情况下如何因沿零光线的可能集中而失效. 我们下面证明, 存在 $C_0^\infty(\mathbb{R}^3)$ 中的函数序列 $f_n$, 满足 $\|f_n\|_{L^2}=1$, 使得对相应的解 $\phi_n$, 
\begin{equation}
\int_0^\infty|\phi_n(t, t, 0, 0)|^2\, \mathrm{d}t\geq n.\label{eq:c8-335}
\end{equation}

反证, 假设事实上对所有满足 $\|f\|_{L^2}=1$ 的 $f\in C_0^\infty(\mathbb{R}^3)$ 及某个 $\varphi\in\mathcal{S}(\mathbb{R})$, $\varphi\not\equiv0$, 有
\[
J:=\int_0^\infty\phi(t, t, 0, 0)\, \varphi(t)\, \mathrm{d}t<C .
\]
鉴于公式(见 $\mathbb{R}^{3+1}$ 中 $\Box$ 的基本解那一节)
\[
\phi(t, x)=(4\pi)^{-1}t\int_{|\xi|=1}f(x+t\xi)\, \mathrm{d}\xi, 
\]
我们发现
\[
J=(4\pi)^{-1}\int_{\mathbb{R}^3}|y|^{-1}f_1(y_1+|y|, y_2, y_3)\, \varphi(|y|)\, \mathrm{d}y, 
\]
或者, 作变量替换 $z=y+(|y|, 0, 0)$, 
\[
J=(4\pi)^{-1}\int_{z_1>0}\frac{1}{z_1}f(z)\, \varphi\!\left(\frac{|z|^2}{2z_1}\right)\mathrm{d}z<c .
\]
由于 $f$ 是任意的 $C_0^\infty(\mathbb{R}^3)$ 函数, $\|f\|_{L^2}=1$, 我们必然有
\[
z\mapsto\frac{1}{z_1}\, \varphi\!\left(\frac{|z|^2}{2z_1}\right)
\]
属于 $L^2(\mathbb{R}_+^3)$, 而当 $\varphi\not\equiv0$ 时这是错误的. 事实上, 
\[
\int_{\mathbb{R}_+^3}\frac{1}{z_1^2}\, \varphi^2\!\left(\frac{|z|^2}{2z_1}\right)\mathrm{d}z
=\int_{\mathbb{R}^3}\frac{1}{(y_1+|y|)|y|}\, \varphi^2(|y|)\, \mathrm{d}y
=2\pi\int_0^\infty\varphi^2(\lambda)\int_0^\pi\frac{\sin\theta}{1+\cos\theta}\, \mathrm{d}\theta
\]
当 $\varphi\not\equiv0$ 时对数发散. 
\end{remark}
\mtsubsec{解的 Fourier 表示}

我们可以用 Fourier 方法求解齐次问题 \eqref{eq:c8-328}. 回顾:若我们对空间变量作 傅里叶变换, 初值问题 \eqref{eq:c8-328}、\eqref{eq:c8-327} 就变成一个常微分方程的 Cauchy 问题:
\[
\partial_t^2\widehat{u}+|\xi|^2\widehat{u}=0, \qquad
\widehat{u}(0, \xi)=\widehat{f}(\xi), \qquad
\partial_t\widehat{u}(0, \xi)=\widehat{g}(\xi), 
\]
它可以显式求解:
\begin{equation}
\widehat{u}(t, \xi)=\cos(t|\xi|)\, \widehat{f}(\xi)+\sin(t|\xi|)\, \frac{\widehat{g}(\xi)}{|\xi|}\label{eq:c8-336}
\end{equation}
于是, 上面定义的基本解 $W(t)h$ 具有形式
\begin{equation}
W(t)h(x)=\int_{\mathbb{R}^n}e^{ix\cdot\xi}\, \frac{\sin(t|\xi|)}{|\xi|}\, \widehat{h}(\xi)\, \mathrm{d}\xi .\label{eq:c8-337}
\end{equation}
由 Duhamel 原理(见 \eqref{eq:c8-330}), 非齐次方程 $\Box u=F$ 的一般解可以表示为
\begin{equation}
u(t)=\partial_t W(t)f+W(t)g+\int_0^t W(t-s)F(s)\, \mathrm{d}s.\label{eq:c8-338}
\end{equation}
设 $D=(-\Delta)^{1/2}$ 为 Fourier 空间中符号为 $|\xi|$ 的算子. 注意, 
\[
(DW(t))f(x)=(W(t)Df)(x)=\int_{\mathbb{R}^n}e^{ix\cdot\xi}\sin(t|\xi|)\, \widehat{f}(\xi)\, \mathrm{d}\xi .
\]
由于 $\sin(t|\xi|)$ 与 $\cos(t|\xi|)$ 有界, 算子 $\partial_t W(t)$ 与 $DW(t)$ 把 $H^s(\mathbb{R}^n)$ 映入自身. 特别地, \eqref{eq:c8-328}、\eqref{eq:c8-327} 的解 $u$ 保持初始数据 $f$ 与 $g$ 的 (Sobolev) 正则性. 更确切地说, 若对某个 $s\in\mathbb{R}$ 有 $f, D^{-1}g\in H^s$, 则 $u(t), D^{-1}\partial_t u(t)\in H^s$ 关于 $t\in\mathbb{R}$ 一致成立. 我们还可以写出
\[
\|u\|_{L_t^\infty\dot{H}_x^\gamma}+\|\partial_t u\|_{L_t^\infty\dot{H}_x^{\gamma-1}}\lesssim\|f\|_{\dot{H}^\gamma}+\|g\|_{\dot{H}^{\gamma-1}}, 
\]
它给出了估计 \eqref{eq:c8-331} 的容易部分\footnote{另一种基于能量恒等式的推导将在下一小节给出. }. 因此, 要证明 \eqref{eq:c8-331}, 只需证明
\begin{equation}
\|u\|_{L_t^q L_x^r}\lesssim\|f\|_{\dot{H}^\gamma}+\|g\|_{\dot{H}^{\gamma-1}}\label{eq:c8-339}
\end{equation}
对任意波容许对 $(q, r)$ 成立. 

我们还注意到
\[
\partial_t W(t)h(x)=\int_{\mathbb{R}^n}e^{ix\cdot\xi}\cos(t|\xi|)\, \widehat{h}(\xi)\, \mathrm{d}\xi
\]
以及
\[
D^{-1}W(t)h(x)=\int_{\mathbb{R}^n}e^{ix\cdot\xi}\, \frac{\cos(t|\xi|)}{|\xi|}\, \widehat{h}(\xi)\, \mathrm{d}\xi .
\]
我们可以把 \eqref{eq:c8-336} 改写为
\[
\widehat{u}(t, \xi)=e^{it|\xi|}\widehat{f}^{+}(\xi)+e^{-it|\xi|}\widehat{f}^{-}(\xi), 
\]
其中 $f^{\pm}=\frac{1}{2}\bigl(f\pm D^{-1}g\bigr)$. 由此得 $u=u^{+}+u^{-}$, 其中
\[
u^{\pm}=\int e^{i(x\cdot\xi\pm t|\xi|)}\, \widehat{f}_{\pm}(\xi)\, \mathrm{d}\xi .
\]
注意, 要证明 \eqref{eq:c8-339} 只需证明
\begin{equation}
\|u^{+}\|_{L_t^q L_x^r}\lesssim\|f^{+}\|_{\dot{H}^\gamma}\label{eq:c8-340}
\end{equation}
以及对 $f^{-}$ 的类似估计. 

\mtsubsec{能量估计}

我们将通过分部积分为 $\Box u=F$ 的一般解导出一个简单的 $L^2$ 估计. 一切都来自如下简单的代数恒等式:
\begin{equation}
-\frac{1}{2}\partial_t\bigl(|\partial_t u|^2+|\nabla u|^2\bigr)+\partial_i(\partial_t u\, \partial_i u)=\partial_t u\cdot F\label{eq:c8-341}
\end{equation}
其中 $|\nabla u|^2=\sum_{i=1}^n(\partial_i u)^2$, $\partial_i=\partial_{x^i}$. 关于 $x$ 积分, 并假设 $u$ 及其导数在无穷远处消失\footnote{这可以由波动方程解的有限传播速度性质容易地证明. }, 我们得到
\[
\partial_t\int_{\mathbb{R}^n}\bigl((\partial_t u)^2+|\nabla u|^2\bigr)\mathrm{d}x=2\int_{\mathbb{R}^n}\partial_t u\cdot F\, \mathrm{d}x .
\]
于是对 $t$ 积分, 
\[
\|\partial_t u(t)\|_{L^2}^2+\|\nabla u(t)\|_{L^2}^2
\leq\|\partial_t u(0)\|_{L^2}^2+\|\nabla u(0)\|_{L^2}^2+2\int_0^t\!\int_{\mathbb{R}^n}\partial_t u\cdot F\, \mathrm{d}x\, \mathrm{d}s, 
\]
我们把它改写为(记 $|\partial u|^2=|\partial_t u|^2+|\nabla u|^2$)
\begin{equation}
\|\partial u(t)\|_{L^2}^2=\|\partial u(0)\|_{L^2}^2+2\int_0^t\!\int_{\mathbb{R}^n}\partial_t u\cdot F\, \mathrm{d}x\, \mathrm{d}s.\label{eq:c8-342}
\end{equation}
特别地, 应用 Hölder 不等式, 
\[
\|\partial u(t)\|_{L^2}^2\leq\|\partial u(0)\|_{L^2}^2+2\int_0^t\|\partial_t u(s)\|_{L^2}\, \|F(s)\|_{L^2}\, \mathrm{d}s, 
\]
由此我们导出非齐次能量估计
\begin{equation}
\sup_{t\in[0, T]}\|\partial u(t)\|_{L^2}\lesssim\|\partial u(0)\|_{L^2}+\int_0^T\|F(s)\|_{L^2}\, \mathrm{d}s\label{eq:c8-343}
\end{equation}
现在设 $D^s$ 为算子 $D^s=(-\Delta)^{s/2}$, 其在 Fourier 空间中的符号为 $|\xi|^s$. 由于 $D^s$ 与 $\Box$ 交换, 我们容易推出
\[
\|\partial D^s u(t)\|_{L^2}^2=\|\partial D^s u(0)\|_{L^2}^2+2\int_{\mathbb{R}^n}\partial_t D^s u\cdot D^s F\, \mathrm{d}x .
\]
利用关于 $x$ 变量的 Plancherel 公式, 我们可以写出
\[
\int_{\mathbb{R}^n}\partial_t D^s u\cdot D^s F\, \mathrm{d}x=\int_{\mathbb{R}^n}\partial_t D^{2s}u\cdot F\, \mathrm{d}x, 
\]
因此, 由 Hölder 不等式, 在条带 $\mathcal{D}_T=[0, T]\times\mathbb{R}^n$ 上, 
\[
\sup_{t\in[0, T]}\|\partial D^s u(t)\|_{L^2}^2
\leq\|\partial D^s u(0)\|_{L^2}^2+2\|D^{2s}\partial_t u\|_{L_t^q L_x^r(\mathcal{D}_T)}\, \|F\|_{L_t^{q'} L_x^{r'}(\mathcal{D}_T)} .
\]
取 $s=-1$, 我们推得
\[
\sup_{t\in[0, T]}\|\partial D^{-1/2}u(t)\|_{L^2}^2
\leq\|\partial D^{-1/2}u(0)\|_{L^2}^2+2\|D^{-1}\partial_t u\|_{L_t^q L_x^r(\mathcal{D}_T)}\, \|F\|_{L_t^{q'} L_x^{r'}(\mathcal{D}_T)} .
\]
我们把这个能量估计应用于具有零初始条件的非齐次问题 0.30.1 的解. 我们还假设维数条件
\[
\frac{1}{q}+\frac{n}{r}=\frac{n}{2}-\gamma=\frac{1}{q'}+\frac{n}{r'}-2
\]
成立. 这蕴含 $\gamma=\frac{1}{2}$. 于是我们有
\[
\sup_{t\in[0, T]}\|\partial D^{-1/2}u(t)\|_{L^2}^2
\leq 2\|D^{-1}\partial_t u\|_{L_t^q L_x^r(\mathcal{D}_T)}\, \|F\|_{L_t^{q'} L_x^{r'}(\mathcal{D}_T)} .
\]
暂时假设我们能证明估计
\begin{equation}
\|D^{-1}\partial_t u\|_{L_t^q L_x^r(\mathcal{D}_T)}\lesssim\|F\|_{L_t^{q'} L_x^{r'}(\mathcal{D}_T)}\label{eq:c8-344}
\end{equation}
则
\[
\sup_{t\in[0, T]}\|\partial D^{-1/2}u(t)\|_{L^2}\lesssim\|F\|_{L_t^{q'} L_x^{r'}(\mathcal{D}_T)}, 
\]
它等价于
\[
\sup_{t\in[0, T]}\Bigl(\|u(t)\|_{\dot{H}^\gamma}+\|\partial_t u\|_{\dot{H}^{\gamma-1}}\Bigr)\lesssim\|F\|_{L^{q'}([0, T];L^{r'})}, 
\]
从而证明了估计 \eqref{eq:c8-332} 的一半. 因此非齐次估计 \eqref{eq:c8-332} 归结为证明
\begin{equation}
\|u\|_{L^q([0, T];L^r)}+\|D^{-1}\partial_t u\|_{L^q([0, T];L^r)}\lesssim\|F\|_{L^{q'}([0, T];L^{r'})}\label{eq:c8-345}
\end{equation}

\mtsubsec{齐次情形}

本节我们证明估计 \eqref{eq:c8-340}, 从而完成\thmref{thm:c8-0.32} 的齐次 Strichartz 估计的证明. 利用时空 傅里叶变换, 即对 $t$ 与 $x$ 都作 傅里叶变换, 
\begin{equation}
\widetilde{u}_{+}(\tau, \xi)=\delta(\tau-|\xi|)\, \widehat{f}_{+}(\xi), \qquad
\widetilde{u}_{-}(\tau, \xi)=-\delta(\tau+|\xi|)\, \widehat{f}_{-}(\xi), \label{eq:c8-346}
\end{equation}
它们是 $\widetilde{u}$ 的位于前向零锥 $C_{+}=\{\tau=|\xi|\}$ 与后向零锥 $C_{-}=\{\tau=-|\xi|\}$ 上的分量. 因此我们可以从半光锥 $C_{+}$ 或 $C_{-}$ 的限制定理的观点来解释 \eqref{eq:c8-340}. 下面我们证明, 只需对 $\widehat{f}_{+}$ 包含在固定二进片中的情形证明 \eqref{eq:c8-340}. 更确切地说, 略去下标 $+$, 只需证明
\begin{equation}
\|u_{k}^{+}\|_{L_t^q L_x^r}\lesssim 2^{k\gamma}\|f_{k}^{+}\|_{L^2}\label{eq:c8-347}
\end{equation}
其中 $u^{+}=\sum_{k\in2^{\mathbb{Z}}}u_{k}^{+}$, $u_{k}^{+}=P_k u^{+}$, $f_{k}^{+}=P_k f^{+}$, 而 $P_k$ 是关于空间变量 $x$ 的标准 LP 投影. 

要证明 \eqref{eq:c8-348} 蕴含 \eqref{eq:c8-340} 是高度非平凡的\footnote{若不使用推论 6.15, 我们只能得到右端用 Besov 范数 $\dot{B}_{2, 1}^{\gamma}$ 代替 $\dot{H}^{\gamma}$ 范数的较弱估计. }, 因为我们需要依赖于适用于混合范数 $L_t^q L_x^r$($q$ 与 $r$ 均大于 2)的推论 6.15. 于是
\[
\|u^{+}\|_{L_t^q L_x^r}^2\lesssim\sum_{k\in\mathbb{Z}}\|u_{k}^{+}\|_{L_t^q L_x^r}^2\lesssim\sum_{k\in\mathbb{Z}}2^{2k\gamma}\|f_{k}^{+}\|_{L^2}^2\lesssim\|f^{+}\|_{\dot{H}^\gamma}
\]
最后我们注意到, 利用一个简单的伸缩论证, \eqref{eq:c8-348} 可由下式推出:
\begin{equation}
\|u_{0}^{+}\|_{L_t^q L_x^r}\lesssim\|f_{0}^{+}\|_{L^2}\label{eq:c8-348}
\end{equation}
我们现在把\newterm{truncated-cone-operator}{截断锥算子}{truncated cone operator} $C$ 定义为算子
\begin{equation}
Cf(t, x)=\int e^{it|\xi|}e^{ix\cdot\xi}\, \chi(\xi)\, \widehat{f}(\xi)\, \mathrm{d}\xi, \label{eq:c8-349}
\end{equation}
其中 $\chi$ 是支集在 $1/2\leq|\xi|\leq2$ 中的截断函数, 例如 LP 投影定义中用到的那一种(见 (214)). 算子 $C$ 可以看作 傅里叶变换到截断锥的限制的伴随算子:
\[
\widehat{C^{*}F}(\xi)=\overline{\chi(\xi)}\, \widetilde{F}(|\xi|, \xi)
\]
估计 \eqref{eq:c8-348} 是如下定理的直接推论. 

\begin{theorem}\label{thm:c8-0.39}
设 $(q, r)$、$(q_1, r_1)$、$(q_2, r_2)$ 是容许指数对. 则我们有估计
\begin{equation}
\|Cf\|_{L_t^q L_x^r}\lesssim\|f\|_{L^2}, \label{eq:c8-350}
\end{equation}
以及
\begin{equation}
\|CC^{*}F\|_{L_t^{q_1} L_x^{r_1}}\lesssim\|F\|_{L_t^{q_2'} L_x^{r_2'}}.\label{eq:c8-351}
\end{equation}
\end{theorem}

把 $C$ 与 $C^{*}$ 复合, 我们得到
\[
CC^{*}F(t, x)\simeq\int e^{i[(t-s)|\xi|+(x-y)\cdot\xi]}\, |\beta(\xi)|^2 F(s, y)\, \mathrm{d}s\, \mathrm{d}y\, \mathrm{d}\xi, 
\]
它可以改写为卷积
\begin{equation}
CC^{*}F(t, \cdot)=\int U(t-s)F(s, \cdot)\, \mathrm{d}s, \label{eq:c8-352}
\end{equation}
其中演化算子
\begin{equation}
U(t)f(x)=\int e^{i(t|\xi|+x\cdot\xi)}\, |\chi(\xi)|^2\, \widehat{f}(\xi)\, \mathrm{d}\xi .\label{eq:c8-353}
\end{equation}
(注意 $U$ 本质上与 $C$ 是同一个算子!)由 $TT^{*}$ 原理, 我们知道估计 \eqref{eq:c8-350} 等价于关于 $CC^{*}$ 的如下估计:
\begin{equation}
\|CC^{*}F\|_{L_t^q L_x^r}\lesssim\|F\|_{L_t^{q'} L_x^{r'}}, \label{eq:c8-354}
\end{equation}
它也等价于极化形式 \eqref{eq:c8-351}. 因此, 要证明定理只需证明 \eqref{eq:c8-354}. 与上一节限制定理的第二个证明一样, 要证明 \eqref{eq:c8-354}, 我们需要证明演化算子 $U(t)$ 的如下性质. 

\begin{proposition}\label{prop:c8-0.40}
设 $\chi(\xi)$ 是支集在 $1/2\leq|\xi|\leq2$ 中的固定 $C_0^{\infty}(\mathbb{R}^n)$ 函数, 且
\begin{equation}
U(t)f(x)=\int e^{i(t|\xi|+x\cdot\xi)}\, \chi(\xi)\, \widehat{f}(\xi)\, \mathrm{d}\xi .\label{eq:c8-355}
\end{equation}
则
\begin{equation}
\|U(t)f\|_{L^2}\lesssim C\|f\|_{L^2}\label{eq:c8-356}
\end{equation}
\begin{equation}
\|U(t)f\|_{L^\infty}\lesssim(1+|t|)^{-\frac{n-1}{2}}\|f\|_{L^1}\label{eq:c8-357}
\end{equation}
由此插值, 对所有 $2\leq r\leq\infty$, 
\begin{equation}
\|U(t)f\|_{L^r}\lesssim(1+|t|)^{-\frac{n-1}{2}\left(1-\frac{2}{r}\right)}\|f\|_{L^{r'}}\label{eq:c8-358}
\end{equation}
此外, 若 $\chi=\chi_{\mu}$ 的支集在大小为 $\mu$ 的立方体中, 则 \eqref{eq:c8-357} 可以加强为
\begin{equation}
\|U(t)f\|_{L^\infty}\lesssim\mu(1+|t|)^{-\frac{n-1}{2}}\|f\|_{L^1}\label{eq:c8-359}
\end{equation}
\end{proposition}

\begin{proof}
我们直接证明更强的版本 \eqref{eq:c8-359}. 我们只需验证 \eqref{eq:c8-357}. 我们写出
\[
U(t)f=K_t*f, \qquad K_t(x)=\int e^{i(x\cdot\xi+t|\xi|)}\, \chi_{\mu}(\xi)\, \mathrm{d}\xi .
\]
只需证明
\[
|K_t(x)|\lesssim\mu\, \frac{1}{1+|t|+|x|}.
\]
在区域 $|x|<|t|/2$ 与 $|x|\geq2|t|$ 中, 我们关于算子
\[
L=-i\sum_j\frac{x_j+t\, \frac{\xi_j}{|\xi|}}{\bigl|x+t\, \frac{\xi}{|\xi|}\bigr|^2}\, \partial_{\xi_j}
\]
分部积分 $k$ 次, 使得 $L\bigl(e^{i(x\cdot\xi+t|\xi|)}\bigr)=e^{i(x\cdot\xi+t|\xi|)}$. 我们还利用简单估计 $|\partial_{\xi}^{\alpha}\chi_{\mu}(\xi)|\lesssim\mu^{-|\alpha|}$ 导出 $|K_t(x)|\lesssim(1+|t|)^{-k}\mu^{n-k}$, 或者, 取 $k=\frac{n-1}{2}$, 
\[
|K_t(x)|\lesssim(1+|t|)^{-\frac{n-1}{2}}\mu^{\frac{n+1}{2}}.
\]
另一方面, 在区域 $|t|\approx|x|$ 中, 我们写出(其中 $\beta(|\xi|)$ 在 $h_{\mu}$ 的支集上消失)
\[
K_t(x)=\int_{1-2\mu}^{1+2\mu}e^{it\lambda}\, \chi(\lambda)\int_{|\xi|=\lambda}e^{ix\cdot\xi}\, h_{\mu}(\xi)\, \mathrm{d}\sigma(\xi)
\]
我们现在需要依赖如下估计:
\begin{equation}
\sup_{1/2\leq\lambda\leq2}\Bigl|\int_{|\xi|=\lambda}e^{ix\cdot\xi}\, h(\xi)\, \mathrm{d}\sigma(\xi)\Bigr|\lesssim(1+|x|)^{-\frac{n-1}{2}}\label{eq:c8-360}
\end{equation}
它容易由上一节讨论的、支集在 $\mathbb{S}^{n-1}$ 上的测度的 傅里叶变换的衰减推出, 见引理 8.8. 因此, 对 $|t|\sim|x|$, 
\[
|K_t(x)|\lesssim\mu(1+|x|)^{-\frac{n-1}{2}}\lesssim\mu(1+|t|)^{-\frac{n-1}{2}}, 
\]
证毕. 
\end{proof}
我们现在准备按照限制定理第二个证明中的同样论证来证明 \eqref{eq:c8-354}. 事实上, 鉴于 \eqref{eq:c8-352} 与 \eqref{eq:c8-358}, 我们导出
\begin{equation}
\|CC^{*}F\|_{L_x^r}(t)\lesssim\int_{-\infty}^{+\infty}(1+|t-s|)^{-\gamma(r)}\|F(s)\|_{L_x^{r'}}\, \mathrm{d}s\label{eq:c8-361}
\end{equation}
其中 $\gamma(r)=\frac{n-1}{2}\bigl(1-\frac{2}{r}\bigr)$. 我们现在正处于与限制定理第二个证明完全相同的情形, 见\eqref{eq:c6-300} 之后的论证. 若 $0<\gamma(r)<1$, 我们可以应用 Hardy--Littlewood--Sobolev 不等式得到
\[
\|CC^{*}F\|_{L_t^q L_x^r}\lesssim\|F\|_{L_t^{q'} L_x^{r'}}, 
\]
当 $-\gamma(r)+1+1/q=1/q'$, 从而 $\gamma(r)=2/q$ 时. 这在 $0<\gamma(r)=2/q<1$ 的情形证明了 \eqref{eq:c8-348}, 从而证明了\thmref{thm:c8-0.39}. 若 $q=2$ 且 $\gamma(r)>1$, 由 \eqref{eq:c8-361} 我们有
\[
\|CC^{*}F\|_{L_t^2 L_x^r}\lesssim\|F\|_{L_t^2 L_x^{r'}}, 
\]
这是应用标准 Hausdorff--Young 不等式得到的. 

最后, 若 $2/q<1$ 且 $\gamma(r)>2/q$, 结果可以由 $\gamma(r)=2/q$ 的情形利用 Sobolev 不等式推出. 由于光锥的一个主曲率消失, 波动方程的 Strichartz 估计并不像可能的那样强. 利用改进的色散估计 \eqref{eq:c8-359}, 我们仍然可以导出更强的结论, 它在应用中非常有用. 

\begin{proposition}\label{prop:c8-0.41}
设 $0<\mu<1$. 设 $f$ 是 $L^2$ 函数, 其 傅里叶变换的支集在距离原点 1 处、大小为 $\mu$ 的立方体中. 设 $(q, r)$ 是 Strichartz 估计的容许指数对. 则
\begin{equation}
\|Cf\|_{L_t^q L_x^r}\lesssim\mu^{\left(\frac{1}{2}-\frac{1}{r}\right)}\|f\|_{L^2}.\label{eq:c8-362}
\end{equation}
\end{proposition}

证明基于改进的色散估计 \eqref{eq:c8-359}. 把它与 \eqref{eq:c8-356} 插值, 我们导出
\[
\|U(t)f\|_{L^r}\lesssim\mu^{1-\frac{2}{r}}(1+|t|)^{-\frac{n-1}{2}\left(1-\frac{2}{r}\right)}\|f\|_{L^{r'}}
\]
证明随后完全如上继续, 导出
\[
\|CC^{*}F\|_{L_t^q L_x^r}\lesssim\mu^{1-\frac{2}{r}}\|F\|_{L_t^{q'} L_x^{r'}}, 
\]
从而由 $TT^{*}$ 论证, $\|Cf\|_{L_t^q L_x^r}\lesssim\mu^{\left(\frac{1}{2}-\frac{1}{r}\right)}\|f\|_{L^2}$, 证毕. 作为命题的直接推论, 我们得到:

\begin{theorem}\label{thm:c8-0.42}
考虑 $\Box u=0$ 的一般解, 其数据 $f, g$ 在 Fourier 空间中的支集位于大小为 $\mu$ 的立方体中, 该立方体位于大小为 $\lambda$ 的二进壳中, $\lambda$ 远大于 $\mu$, 比如 $\lambda\geq8\mu$. 则
\begin{equation}
\|u\|_{L_t^q L_x^r}\lesssim\mu^{1-\frac{2}{r}}\Bigl(\|f\|_{\dot{H}^\gamma}+\|f\|_{\dot{H}^{\gamma-1}}\Bigr)\label{eq:c8-363}
\end{equation}
\end{theorem}

\begin{proof}
证明容易由上面的命题通过伸缩论证推出. 
\end{proof}

最后, 我们陈述下面另一个结果, 它容易由衰减估计 \eqref{eq:c8-357} 推出. 

\begin{theorem}\label{thm:c8-0.43}
设 $u$ 是自由波, 即齐次方程 $\Box u=0$ 的、具有初始数据 $(f, g)$ 的解. 则
\[
\begin{array}{rcl}
\|u(t)\|_{L^\infty} & \lesssim & |t|^{-\frac{n-1}{2}}\sum_{\lambda\in2^{\mathbb{Z}}}\Bigl(\lambda^{\frac{n+1}{2}}\|f_{\lambda}\|_{L^1}+\lambda^{\frac{n-1}{2}}\|g_{\lambda}\|_{L^1}\Bigr)\\[4pt]
& = & |t|^{-\frac{n-1}{2}}\Bigl(\|f\|_{\dot{B}_{1, 1}^{n+1/2}}+\|g\|_{\dot{B}_{1, 1}^{n-1/2}}\Bigr).
\end{array}
\]
\end{theorem}

大 $t$ 时的一致衰减率 $|t|^{-\frac{n-1}{2}}$ 在标准波动方程的非线性扰动研究中起着非常重要的作用. 

\mtsubsec{非齐次 Strichartz 估计}

我们已经把\thmref{thm:c8-0.32} 的非齐次 Strichartz 估计 \eqref{eq:c8-332} 归结为估计 \eqref{eq:c8-345}. 与齐次估计的情形一样, 我们现在可以把 \eqref{eq:c8-345} 归结为 $F$ 的空间 傅里叶变换支集在单位二进环 $1/2\leq|\xi|\leq2$ 中的情形. 此外, 像前面一样把 $u$ 分解为 $\pm$ 两部分, 只需分别对 $u_{+}$ 与 $u_{-}$ 证明估计. 因此我们需要证明
\begin{equation}
\|u^{+}\|_{L^q([0, T];L^r)}+\|D^{-1}\partial_t u^{+}\|_{L^q([0, T];L^r)}\lesssim\|F\|_{L^{q'}([0, T];L^{r'})}\label{eq:c8-364}
\end{equation}
我们有
\[
\begin{array}{rcl}
u_{+}(t, \cdot) & = & \displaystyle\int_0^t U(t-s)F(s, \cdot)\, \mathrm{d}s\\[6pt]
D^{-1}\partial_t u_{+}(t, \cdot) & = & \displaystyle\int_0^t \partial_t D^{-1}U(t-s)F(s, \cdot)\, \mathrm{d}s
\end{array}
\]
由于根据二进限制有 $\partial_t D^{-1}U(t)\sim U(t)$, 只需证明关于 $\|u^{+}\|_{L^q([0, T];L^r)}$ 的估计. 显然, $u^{+}$ 与 \eqref{eq:c8-352} 中的 $CC^{*}F$ 只相差积分区间限制到 $[0, t]$. 鉴于此事实, 我们写 $u_{+}=(CC^{*})_R F$. 我们于是被引导到如下定理, 由它得到 \eqref{eq:c8-364}, 从而得到 \eqref{eq:c8-332}. 

\begin{theorem}\label{thm:c8-0.45}
设 $U(t)$ 如 \eqref{eq:c8-355} 中定义, 并设
\[
(CC^{*})_R F(t, \cdot)=\int_0^t U(t-s)F(s, \cdot)\, \mathrm{d}s
\]
则对所有容许对 $(q_1, r_1)$、$(q_2, r_2)$, 
\begin{equation}
\|(CC^{*})_R F\|_{L_t^{q_1} L_x^{r_1}([0, T]\times\mathbb{R}^n)}\lesssim\|F\|_{L_t^{q_2'} L_x^{r_2'}([0, T]\times\mathbb{R}^n)}\label{eq:c8-365}
\end{equation}
\end{theorem}

\begin{proof}
在 $(q_1, r_1)=(q_2, r_2)=(q, r)$ 的情形, 证明是直接的. 事实上, 在这种情形我们只需重复估计 \eqref{eq:c8-354} 的证明, 并把积分的上下限考虑进去. 我们还处理了 $q_1=\infty$, $r_1=2$ 的情形, 见能量估计那一小节. 其余非对角情形要稍微困难一些, 将在本节后面讨论的更一般的抽象框架中处理. 不过我们给出的证明已经覆盖了估计 \eqref{eq:c8-332} 最有趣的情形. 于是我们给出了\thmref{thm:c8-0.32} 前两部分的完整证明. 
\end{proof}

\mtsubsec{容许条件的必要性}

为了理解 $q$ 与 $r$ 的最优指数范围, 我们在 \eqref{eq:c8-349} 中定义的截断锥算子 $C$ 的背景下考虑 Knapp 反例的类比. 对某个小的 $\delta>0$, 设
\[
D=\bigl\{\xi\in\mathbb{R}^n:|\xi_1-1|<1/2, \ |\xi'|<\delta\bigr\}, 
\]
并考虑 $f=\chi_D$. 我们有
\[
Cf(t, x)=e^{i(t+x_1)}\int_D e^{i[t(|\xi|-\xi_1)+(t+x_1)(\xi_1-1)+x'\cdot\xi']}\, \mathrm{d}\xi, 
\]
并注意到
\[
|\xi|-\xi_1=\frac{|\xi'|^2}{|\xi|+\xi_1}\lesssim\delta^2 .
\]
于是我们可以选择由
\[
|t|\lesssim\delta^{-2}, \qquad |t+x_1|\lesssim1, \qquad |x'|\lesssim\delta^{-1}
\]
定义的时空区域 $R$, 使得当 $(t, x)\in R$ 且 $\xi\in D$ 时, 最后一个积分内的振荡因子可以当作常数处理. 因此, 对 $(t, x)\in R$ 有 $|Cf(t, x)|\gtrsim|D|$, 并且我们有
\[
\frac{\|Cf\|_{L_t^q L_x^r}}{\|f\|_{L^2}}\gtrsim\frac{|D|\, \|\chi_R\|_{L_t^q L_x^r}}{|D|^{1/2}}\sim\delta^{\frac{n-1}{2}-\frac{2}{q}-\frac{n-1}{r}} .
\]
在极限 $\delta\to0$ 中, \eqref{eq:c8-354} 形式的估计必然蕴含 $q$ 与 $r$ 满足条件
\begin{equation}
\frac{2}{q}\leq(n-1)\left(\frac{1}{2}-\frac{1}{r}\right).\label{eq:c8-366}
\end{equation}
对 $q$ 的范围的另一限制, 即 $q\geq2$, 是算子 $CC^{*}$ 在时间平移下不变性的推论. 事实上, 对平移不变的算子, 我们有如下归功于 Hörmander 的一般结果 [?]. 

\begin{proposition}\label{prop:c8-0.47}
设 $T:L^p(\mathbb{R}^n)\to L^q(\mathbb{R}^n)$ 是一个(非平凡的)线性算子, 它与平移交换, 即 $(Tf)\circ\tau_y=T(f\circ\tau_y)$, 其中 $\tau_y(x)=x+y$, $x, y\in\mathbb{R}^n$. 若 $T$ 从 $L^p$ 到 $L^q$ 有界, 则必然有 $q\geq p$. 
\end{proposition}

证明基于如下引理. 

\begin{lemma}\label{lem:c8-0.48}
设 $f\in L^p(\mathbb{R}^n)$, 则
\[
\lim_{|y|\to\infty}\|f+f\circ\tau_y\|_{L^p}=2^{1/p}\|f\|_{L^p}.
\]
\end{lemma}

\begin{proof}
对每个 $R>0$ 考虑分解 $f=g_R+h_R$, 其中当 $|x|<R$ 时 $g_R(x)=f(x)$, 当 $|x|\geq R$ 时 $g_R(x)=0$;当 $|x|<R$ 时 $h_R(x)=0$, 当 $|x|\geq R$ 时 $h_R(x)=f(x)$. 则
\[
\lim_{R\to\infty}\|g_R\|_{L^p}=\|f\|_{L^p}, \qquad \lim_{R\to\infty}\|h_R\|_{L^p}=0.
\]
取 $R=|y|/2$, 有
\[
f+f\circ\tau_y=g_R+g_R\circ\tau_y+h_R+h_R\circ\tau_y .
\]
函数 $g_R$ 与 $g_R\circ\tau_y$ 的支集不相交, 所以
\[
\|g_R+g_R\circ\tau_y\|_{L^p}^p=\|g_R\|_{L^p}^p+\|g_R\circ\tau_y\|_{L^p}^p=2\|g_R\|_{L^p}^p, 
\]
而
\[
\lim_{|y|\to\infty}\|h_R+h_R\circ\tau_y\|_{L^p}\leq\lim_{|y|\to\infty}2\|h_R\|_{L^p}=0, 
\]
因此
\[
\lim_{|y|\to\infty}\|f+f\circ\tau_y\|_{L^p}=\lim_{|y|\to\infty}2^{1/p}\|g_R\|_{L^p}=2^{1/p}\|f\|_{L^p}.
\]
\end{proof}

\begin{proof}[命题 \ref{prop:c8-0.47} 的证明]
设 $C>0$ 是估计
\[
\|Tf\|_{L^q}\leq C\|f\|_{L^p}, \qquad \forall f\in L^p
\]
的最优常数. 则由线性与平移不变性, 
\[
\|Tf+(Tf)\circ\tau_y\|_{L^q}\leq C\|f+f\circ\tau_y\|_{L^p}.
\]
当 $|y|\to\infty$ 时, 应用引理得到
\[
2^{1/q}\|Tf\|_{L^q}\leq C\, 2^{1/p}\|f\|_{L^p}, \qquad \forall f\in L^p .
\]
$C$ 的最优性蕴含 $2^{\frac{1}{p}-\frac{1}{q}}\geq1$, 从而 $q\geq p$. 
\end{proof}

该命题容易推广到向量值 $L^p$ 空间;若我们把 $CC^{*}$ 看作从 $L^{q'}(\mathbb{R};L_x^{r'})$ 到 $L^q(\mathbb{R};L_x^r)$ 的算子, 则必须有 $q\geq q'$, 这就是条件 $q\geq2$. 
\mtsubsec{一个一般的抽象框架}

事实证明, 上面描述的证明 Strichartz 估计的方法适用于许多其他方程, 如 Schrödinger 方程、KdV 方程等. 因此建立一个适用于所有这些情形的统一框架是值得的. 

设 $(X, \mathrm{d}\mu)$ 是一个测度空间, H 是一个 Hilbert 空间. 考虑一族算子 $(U(t))_{t\in\mathbb{R}}$, $U(t):H\to L^2(X)$, 它描述某个数据在 $H$ 中的系统的演化. 我们假设该演化满足如下两个性质:

\begin{itemize}
\item 对所有 $t\in\mathbb{R}$ 与 $f\in H$, 我们有能量估计:
\begin{equation}
\|U(t)f\|_{L^2(X)}\lesssim\|f\|_H;\label{eq:c8-367}
\end{equation}
\item 对所有 $t\neq s$ 与 $g\in L^1(X)$, 我们有\newterm{dispersive-inequality}{色散不等式}{dispersive inequality}:
\begin{equation}
\bigl\|U(t)U^{*}(s)g\bigr\|_{L^\infty(X)}\lesssim|t-s|^{-\gamma_0}\|g\|_{L^1(X)}, \label{eq:c8-368}
\end{equation}
对某个 $\gamma_0>0$. 
\end{itemize}

在 \eqref{eq:c8-367} 与 \eqref{eq:c8-368} 之间插值, 我们得到估计
\begin{equation}
\bigl\|U(t)U^{*}(s)g\bigr\|_{L^r(X)}\lesssim|t-s|^{-\gamma(r)}\|g\|_{L^{r'}(X)}, \label{eq:c8-369}
\end{equation}
对 $r\geq2$, 其中
\[
\gamma(r)=\gamma_0\left(1-\frac{2}{r}\right).
\]

\begin{theorem}\label{thm:c8-0.50}
若演化算子 $U(t)$ 满足 \eqref{eq:c8-367} 与 \eqref{eq:c8-368}, 则估计
\begin{equation}
\|U(t)f\|_{L_t^q L_X^r}\lesssim\|f\|_H, \label{eq:c8-370}
\end{equation}
对所有满足
\begin{equation}
\frac{2}{q}=\gamma(r), \qquad (q, r, \gamma_0)\neq(2, \infty, 1)\label{eq:c8-371}
\end{equation}
的 $q, r\geq2$ 成立. 
\end{theorem}

\begin{remark}\label{rem:c8-0.51}
这种形式的 Strichartz 不等式适用于线性色散方程, 如 Schrödinger 方程. 
\end{remark}

\begin{proof}
若我们考虑由 $Tf(t, x)=(U(t)f)(x)$ 定义的算子 $T:H\to L_t^q L_X^r$, 则容易验证 $T$ 的对偶是算子 $T^{*}:L_t^{q'} L_X^{r'}\to H$, 由 $T^{*}F=\int U^{*}(s)F(s, \cdot)\, \mathrm{d}s$ 给出. 由 $TT^{*}$ 方法, \eqref{eq:c8-370} 等价于估计
\begin{equation}
\Bigl\|\int U(t)U^{*}(s)F(s)\, \mathrm{d}s\Bigr\|_{L_t^q L_X^r}\lesssim\|F\|_{L_t^{q'} L_X^{r'}}.\label{eq:c8-372}
\end{equation}
由对偶性与对称性考虑, 这又等价于
\begin{equation}
|B(F, G)|\lesssim\|F\|_{L_t^{q'} L_X^{r'}}\, \|G\|_{L_t^{q'} L_X^{r'}}, \label{eq:c8-373}
\end{equation}
其中 $B(F, G)$ 是双线性型
\begin{equation}
B(F, G)=\iint_{s<t}\bigl\langle U^{*}(t)F(t), U^{*}(s)G(s)\bigr\rangle\, \mathrm{d}t\, \mathrm{d}s.\label{eq:c8-374}
\end{equation}
由 \eqref{eq:c8-369} 的双线性版本, 我们有
\begin{equation}
|B(F, G)|\lesssim\iint\frac{\|F(t)\|_{L^{r'}}\|G(s)\|_{L^{r'}}}{|t-s|^{\gamma(r)}}\, \mathrm{d}s\, \mathrm{d}t.\label{eq:c8-375}
\end{equation}
若 $\gamma(r)<1$, 我们可以应用 Hardy--Littlewood--Sobolev 不等式得到 \eqref{eq:c8-373}. 这完成了 $q=2/\gamma(r)>2$ 情形的证明. 

端点情形, 即 $\gamma(r)=2/q=1$ 的情形, 在 $r<\infty$ 时是允许的. 其证明将在下一节描述. 
\end{proof}

\begin{remark}\label{rem:c8-0.52}
若我们把色散条件 \eqref{eq:c8-368} 加强为
\begin{equation}
\bigl\|U(t)U^{*}(s)g\bigr\|_{L^\infty(X)}\lesssim(1+|t-s|)^{-\gamma_0}\|g\|_{L^1(X)}, \label{eq:c8-376}
\end{equation}
则 \eqref{eq:c8-375} 可以改进为
\begin{equation}
|B(F, G)|\lesssim\iint\frac{\|F(t)\|_{L^{r'}}\|G(s)\|_{L^{r'}}}{(1+|t-s|)^{\gamma(r)}}\, \mathrm{d}s\, \mathrm{d}t.\label{eq:c8-377}
\end{equation}
现在我们可以从 Young 不等式得到 \eqref{eq:c8-373}, 当 $2/q=1/p$ 且 $(1+|t|)^{-\gamma(r)}\in L^p(\mathbb{R})$, 即 $\gamma(r)p>1$ 时. 因此, \eqref{eq:c8-376} 使我们能把\thmref{thm:c8-0.50} 中的 Strichartz 估计 \eqref{eq:c8-370} 推广到范围
\begin{equation}
\frac{2}{q}\leq\gamma(r), \qquad (q, r, \gamma_0)\neq(2, \infty, 1).\label{eq:c8-378}
\end{equation}
这种情形适用于线性波动方程. 
\end{remark}

\begin{remark}\label{rem:c8-0.53}
我们注意到, 这个抽象表述中的对象具有自然的伸缩. 更确切地说, \thmref{thm:c8-0.50} 中的估计 \eqref{eq:c8-370} 在如下定义的尺度变换下不变:
\begin{equation}
U(t)\leftarrow U(t/\lambda), \quad U^{*}(s)\leftarrow U^{*}(s/\lambda), \quad \mathrm{d}\mu\leftarrow\lambda^{\gamma_0}\, \mathrm{d}\mu, \quad \langle f, g\rangle_H\leftarrow\lambda^{\gamma_0}\langle f, g\rangle_H.\label{eq:c8-379}
\end{equation}
我们还可以考虑端点情形
\[
q=2, \qquad r=\frac{2\gamma_0}{\gamma_0-1}, \qquad \gamma_0>1.
\]
事实上, 这比之前的非端点情形更困难, 需要比插值 \eqref{eq:c8-369} 给出的单参数族更好的双参数估计. 这个证明已在前面提到的 Keel--Tao 的论文 \emph{Endpoint Strichartz Estimates} 中给出. 我们在此略去. 
\end{remark}

\mtsubsec{非齐次估计}

说一个算子 $T$ 把 Hilbert 空间 $H$ 映入 $L_t^q L_X^r$, 等价于说它的对偶 $T^{*}$ 把 $L_t^{q'} L_X^{r'}$ 映入 H, 也等价于说 $TT^{*}$ 算子把 $L_t^{q'} L_X^{r'}$ 映入 $L_t^q L_X^r$. 若对 $(q, r)$ 允许在容许指数集 $E$ 中变化, 我们可以把 $TT^{*}$ 看作与不同指数对相关联的两个算子的复合. 由此可知, $TT^{*}$ 实际上满足更大的一组映射性质, 因为它把 $L_t^{\widetilde{q}'} L_X^{\widetilde{r}'}$ 映入 $L_t^q L_X^r$, 对任意两对 $(q, r)$、$(\widetilde{q}, \widetilde{r})\in E$. 

上一小节定义的算子 $Tf(t)=U(t)f$ 可以看作某个齐次的、平移不变的线性演化方程的解. 对应的非齐次问题的解, 利用 Duhamel 原理, 将由\newterm{retarded-operator}{推迟算子}{retarded operator}表示:
\[
RF(t)=\int_{s<t}U(t)U^{*}(s)F(s)\, \mathrm{d}s .
\]
注意, 算子 $R$ 与 $TT^{*}$ 算子非常相似, 后者由
\[
TT^{*}F(t)=\int U(t)U^{*}(s)F(s)\, \mathrm{d}s
\]
给出. 

然而, R 的定义中 $s<t$ 的限制破坏了 $TT^{*}$ 的复合结构. 幸运的是, 我们上面导出的 $TT^{*}$ 的所有映射性质都可以转移到 R. 

\begin{theorem}\label{thm:c8-0.55}
算子 $R$ 把 $L_t^{\widetilde{q}'} L_X^{\widetilde{r}'}$ 映入 $L_t^q L_X^r$, 对任意两对 $(q, r)$、$(\widetilde{q}, \widetilde{r})$, 只要 Strichartz 估计 \eqref{eq:c8-370} 成立. 
\end{theorem}

\begin{proof}
首先注意, 在\thmref{thm:c8-0.50} 的证明中, 我们实际上证明了对角情形 $(q, r)=(\widetilde{q}, \widetilde{r})$. 事实上, \eqref{eq:c8-374} 中定义的双线性型可以写成 $B(F, G)=\iint R(F)\cdot G\, \mathrm{d}x\, \mathrm{d}t$, 而 \eqref{eq:c8-373} 是 $R$ 的映射性质的对偶表述. 

满足 $\frac{1}{q}+\frac{1}{\widetilde{q}}<1$ 的非对角情形由 $TT^{*}$ 的映射性质推出, 利用 Christ 与 Kiselev 关于积分算子的一般论证(见 [] 与 []), 我们在下面的\propref{prop:c8-0.57} 中总结它. 

剩下要考虑 $q=\widetilde{q}=2$ 且 $r\neq\widetilde{r}$ 的情形, 假设演化 $U(t)$ 满足更强的色散不等式 \eqref{eq:c8-376}, 其中 $\gamma_0>1$. 由于我们已经证明了 $r=\widetilde{r}$ 的情形, 由插值只需考虑极端情形:$r=r_{*}=\frac{2\gamma_0}{\gamma_0-1}$, $\widetilde{r}=\infty$, 并证明
\[
|B(F, G)|\lesssim\|F\|_{L_t^2 L_X^{r_{*}'}}\, \|G\|_{L_t^2 L_X^1}.
\]
这个估计通过把 $B(F, G)$ 分解为二进片 $B=\sum_{\lambda\in2^{\mathbb{Z}}}B_{\lambda}$ 得到, 其中
\begin{equation}
B_{\lambda}(F, G)=\iint_{\lambda/2\leq|t-s|\leq2\lambda}\bigl\langle U^{*}(t)F(t), U^{*}(s)G(s)\bigr\rangle\, \mathrm{d}t\, \mathrm{d}s, \label{eq:c8-380}
\end{equation}
所需的结论立即由下面的引理推出. 
\end{proof}

\begin{lemma}\label{lem:c8-0.56}
设 $B_{\lambda}(F, G)$ 是 \eqref{eq:c8-380} 中定义的双线性型. 则存在 $\varepsilon>0$ 使得
\[
|B_{\lambda}(F, G)|\lesssim\min\bigl\{\lambda, \lambda^{-1}\bigr\}^{\varepsilon}\, \|F\|_{L_t^2 L_X^{r_{*}'}}\, \|G\|_{L_t^2 L_X^1}.
\]
\end{lemma}

\begin{proof}
我们可以假设 $F$ 与 $G$ 支集在长度 $O(\lambda)$ 的不相交时间区间上, 彼此相距 $O(\lambda)$. 则 $B_{\lambda}(F, G)=\langle T^{*}F, T^{*}G\rangle_H$. 我们用能量估计来界定 $\|T^{*}F\|_H$, 用 $q=2$、$r=\infty$ 的 Strichartz 估计来界定 $\|T^{*}G\|_H$, 于是
\[
|B_{\lambda}(F, G)|\lesssim\|F\|_{L_t^1 L_X^2}\, \|G\|_{L_t^2 L_X^1}.
\]
然后我们应用 Hölder 不等式, 并利用关于 $F$ 支集的假设, 得到
\[
|B_{\lambda}(F, G)|\lesssim\lambda^{1/2}\|F\|_{L_t^2 L_X^2}\, \|G\|_{L_t^2 L_X^1}.
\]
我们还可以写出 $B_{\lambda}(F, G)=\iiint F(t)\cdot U(t)U^{*}(s)G(s)\, \mathrm{d}x\, \mathrm{d}s\, \mathrm{d}t$, 并利用色散不等式, 
\[
|B_{\lambda}(F, G)|\lesssim(1+\lambda)^{-\gamma_0}\|F\|_{L_t^1 L_X^1}\, \|G\|_{L_t^1 L_X^1}.
\]
再次应用 Hölder 不等式, 并利用关于 $F$ 与 $G$ 支集的假设, 得到
\[
|B_{\lambda}(F, G)|\lesssim\frac{\lambda}{(1+\lambda)^{\gamma_0}}\, \|F\|_{L_t^2 L_X^1}\, \|G\|_{L_t^2 L_X^1}.
\]
因此, $B_{\lambda}$ 在 $L_t^2 L_X^2\times L_t^2 L_X^1$ 上有界, 常数为 $\lambda^{1/2}$;在 $L_t^2 L_X^1\times L_t^2 L_X^1$ 上有界, 常数为 $\frac{\lambda}{(1+\lambda)^{\gamma_0}}$. 通过 $L^p$ 空间的标准插值, 我们得到 $B_{\lambda}$ 在 $L_t^2 L_X^{r_{*}'}\times L_t^2 L_X^1$ 上有界, 常数为 $C_{\lambda}$, 其中
\[
C_{\lambda}=\lambda^{\theta/2}\left(\frac{\lambda}{(1+\lambda)^{\gamma_0}}\right)^{1-\theta}, \qquad
\frac{1}{r_{*}'}=\frac{\theta}{2}+\frac{1-\theta}{1}, \qquad
r_{*}=\frac{2\gamma_0}{\gamma_0-1}.
\]
化简表达式, 我们发现
\[
C_{\lambda}=\frac{\lambda^{\frac{\gamma_0+1}{2\gamma_0}}}{1+\lambda}\lesssim\min\bigl\{\lambda, \lambda^{-1}\bigr\}^{\varepsilon}, 
\]
其中
\[
\varepsilon=\min\left\{\frac{\gamma_0+1}{2\gamma_0}, \ 1-\frac{\gamma_0+1}{2\gamma_0}\right\}
=\frac{\gamma_0-1}{2\gamma_0}=\frac{1}{r_{*}}>0 .
\]
\end{proof}

\mtsubsecn{0.56.1}{具有受限核的积分算子}

在本小节中, 我们给出上面提到的 Christ--Kiselev 结果的自包含阐述. 考虑具有可测核 $K(s, t)$ 的积分算子
\[
Tf(t)=\int_{\mathbb{R}}K(s, t)f(s)\, \mathrm{d}s, 
\]
以及与核 $K(s, t)\chi(s<t)$ 相关联的受限版本
\[
Rf(t)=\int_{s<t}K(s, t)f(s)\, \mathrm{d}s .
\]
若 $T$ 把 $L^p$ 映入 $L^q$, 且 $1\leq p<q\leq\infty$, 则 $R$ 也把 $L^p$ 映入 $L^q$. 这个事实的等价表述由下面的命题给出. 

\setcounter{mtthm}{56}

\begin{proposition}\label{prop:c8-0.57}
设 $K(s, t)$ 是 $\mathbb{R}\times\mathbb{R}$ 上的可测函数. 设 $B(f, g)$ 是以 $K$ 为核的双线性型, 
\[
B(f, g)=\iint K(s, t)f(s)g(t)\, \mathrm{d}s\, \mathrm{d}t, 
\]
并设 $\widetilde{B}(f, g)$ 是核限制在区域 $s<t$ 上的双线性型, 
\[
\widetilde{B}(f, g)=\iint_{s<t}K(s, t)f(s)g(t)\, \mathrm{d}s\, \mathrm{d}t .
\]
设 $p, q\geq1$, 满足条件
\begin{equation}
\frac{1}{p}+\frac{1}{q}>1.\label{eq:c8-381}
\end{equation}
若 $B$ 在 $L^p\times L^q$ 上有界, 
\[
|B(f, g)|\lesssim\|f\|_{L^p}\|g\|_{L^q}, 
\]
则 $\widetilde{B}$ 也在 $L^p\times L^q$ 上有界, 
\[
\bigl|\widetilde{B}(f, g)\bigr|\lesssim\|f\|_{L^p}\|g\|_{L^q}.
\]
\end{proposition}

\begin{remark}\label{rem:c8-0.58}
有些情形不允许条件 \eqref{eq:c8-381} 取等号. 例如考虑 Hilbert 变换的情形, 它对应于核 $K(s, t)=\frac{1}{s-t}$, 且 $p=q=2$. 
\end{remark}

\begin{proof}
设 $f\in L^p$, $g\in L^q$, 且 $\|f\|_{L^p}=\|g\|_{L^q}=1$. 

定义 $F(t)=\int_{s<t}|f(s)|^p\, \mathrm{d}s$. F 是连续非降函数, 把 $[-\infty, +\infty]$ 映到 $[0, 1]$ 上. 特别地, $I=[a, b]\subset[0, 1]$ 型区间的原像也是同型区间, $F^{-1}(I)=[A, B]$, 其中 $F(A)=a$, $F(B)=b$, 且 $\int_A^B|f(s)|^p\, \mathrm{d}s=F(B)-F(A)=b-a$. 因此
\begin{equation}
\|f\|_{L^p(F^{-1}(I))}=|I|^{1/p}.\label{eq:c8-382}
\end{equation}

现在考虑集合 $\Omega=\{(x, y)\in\mathbb{R}^2:x<y\}$ 的 Whitney 分解, 把它分解为不相交的二进正方形, 如\lemref{lem:c6-6.7}, $\Omega=\cup_Q Q$, 其中每个正方形 $Q=I\times J$ 具有性质
\begin{equation}
\operatorname{dist}(I, J)\approx|I|=|J|=\lambda, \label{eq:c8-383}
\end{equation}
对某个二进值 $\lambda$. 若我们只看覆盖三角形 $\Omega\cap[0, 1]^2$ 所需的那些正方形, 则 $\lambda\leq1/2$. 

注意 $s<t$ 蕴含或者 $F(s)<F(t)$, 或者 $f$ 在区间 $[s, t]$ 上几乎处处等于 0. 因此我们可以写出
\[
\widetilde{B}(f, g)=\iint_{F(s)<F(t)}K(s, t)f(s)g(t)\, \mathrm{d}s\, \mathrm{d}t
=\sum_Q B\bigl(\chi_{F^{-1}(I)}f, \chi_{F^{-1}(J)}g\bigr).
\]
利用 $B$ 在 $L^p\times L^q$ 上的有界性, 我们得到
\[
\bigl|\widetilde{B}(f, g)\bigr|\lesssim\sum_Q\|f\|_{L^p(F^{-1}(I))}\, \|g\|_{L^q(F^{-1}(J))}.
\]
现在我们利用 \eqref{eq:c8-382}、\eqref{eq:c8-383} 以及如下事实:对每个给定的二进区间 J, 使得 $I\times J$ 是 Ω 分解中一个正方形的区间 $I$ 的数目被一个普适常数界定. 因此
\[
\bigl|\widetilde{B}(f, g)\bigr|\lesssim\sum_{\lambda\leq1/2}\lambda^{\frac{1}{p}}\sum_{|J|=\lambda}\|g\|_{L^q(F^{-1}(J))}.
\]
接下来, 我们对长度为 $\lambda$ 的二进区间 $J$ 的求和应用 Hölder 不等式;由于 $[0, 1]$ 中有 $\lambda^{-1}$ 个这样的区间, 我们有
\[
\bigl|\widetilde{B}(f, g)\bigr|\lesssim\sum_{\lambda\leq1/2}\lambda^{\frac{1}{p}}\lambda^{-\frac{1}{q'}}\|g\|_{L^q}
=\sum_{\lambda\leq1/2}\lambda^{\frac{1}{p}+\frac{1}{q}-1}\lesssim1 .
\]
\end{proof}


% ==== part09 第 9 章:双线性估计 ====
% 源: MinerU_latex_A-非.tex 302-717 行;校对 pdf_text.txt 11162-11551 行(PDF 第 205-212 页)
% 术语(新增，键=中文=英文):
%   bilinear-operator=双线性算子=bilinear operator; bilinear-estimate=双线性估计=bilinear estimate;
%   reduced-initial-data=约化初值=reduced initial data; admissible-pair=容许对=acceptable pair
%   (原书用 acceptable，同 admissible pair); dimensional-analysis=量纲分析=dimensional analysis;
%   dimensional-condition=量纲条件=dimension condition; trichotomy-formula=三分公式=trichotomy formula;
%   low-high-term=低-高项=low-high term; high-low-term=高-低项=high-low term;
%   high-high-term=高-高项=high-high term; littlewood-paley-projection=Littlewood-Paley 投影=
%   Littlewood-Paley projection; ellipsoid-of-revolution=旋转椭球面=ellipsoid of revolution;
%   hyperboloid-of-revolution=旋转双曲面=hyperboloid of revolution; null-quadratic-form=零二次型=
%   null quadratic form; null-form=零形式=null form; symbol=象征=symbol;
%   Strichartz 估计=Strichartz estimate(专名，直接书写)
% REFS:
%   跨章引用(第 8 章，均写成 eqref/autoref，由协调员统一核对标签):
%     - 第 1 节 "The Strichartz estimate (331)" → \eqref{eq:c8-331}
%     - 第 1 节 "recalling (346)" → \eqref{eq:c8-346}
%     - 第 2 节 "apply theorem 0.42" → \thmref{thm:c8-0.42}
%   原书未解析引用(照排保留 "??"/"[?]"):
%     - 源 495 行 "theorem \ref{eq:10.16}" → 正文写 "定理 ??"，原书与 pdf_text 均显示 ??
%     - 源 564 行 "see [?]"、源 584 行 "see [?]，[?]，[?] and [?]" 文献引用缺失，照排
%     - 源 646 行 "special case of theorem ??"、源 716 行 "particular cases of theorem ??"
%     - 源 612 行 "presented in the next section": 所指第 9 章第 4 节在原书不存在(第 9 章止于第 3 节)
%   FOOTNOTE: 定理 2.1 脚注 1 按 pdf_text.txt 第 11469-11470 行补全并翻译
%   OCR/原书笔误修正(均对照 pdf_text.txt，按数学含义处理):
%     - 源 382/386 行 (386)(387): 积分变量 s 与 x 混用，统一为 x;F 的第二自变量 (τ+x|ξ|)/2 改为
%       (τ-x|ξ|)/2(原书两源一致印错，与推导 J_F^- = ∫F(λ，τ-λ)dλ、λ=(τ+x|ξ|)/2 不符)
%     - 源 402 行 δ 恒等式第三行 2λξcosθ → 2λ|ξ|cosθ;末行 2a|ξ| 补 λ → 2a|ξ|λ(与 (388) 一致)
%     - 源 428 行 J_F^- 第一行 [1-(…)] 补平方 → [1-(…)²]
%     - 源 478 行 "g_k" → "g_λ"，"H^α" → "H^a"(与 b/2=γ-a 一致)
%     - 源 490 行 "By symmetry" 后标注 LH → HL(原书笔误，数学上应为高-低项)
%     - 源 589 行 "\Box = \Box v = 0" → "\Box u = \Box v = 0"
%     - 源 625 行 2Q_0(u，v)=□(uv)-□u v-u□u 末项 → u□v(数学上必须，与后续 □u=□v=0 用法一致)
%     - 源 655 行 "|||f|||" → ‖f‖
%     - 源 660 行 "with initial." 原书句子残缺，按"带初值"译出
%     - 源 713 行 两式第二因子 "u[0]" → "v[0]"(原书笔误)
%     - "Cauchy-Schwartz" → "Cauchy-Schwarz"(原书拼写)
% 编号说明:
%   - 公式 (384)-(400) 共 17 个，用 \begin{equation} 自动编号(全书连续,part08 应止于 (383));
%     其余行间公式用 \[ \] 不编号(与原书一致)
%   - 定理类 1.1/1.2/2.1/3.1-3.7 共享计数器自动对齐;Lemma 3.7 在 OCR 源中显示为
%     \section{LEMMA 3.7...},实为引理,用 lemma 环境
%   - 第 9 章无章引言,直接从第 1 节开始;源第 717 行后为附录 A(partapx),未越界
% ====

\chapter{双线性估计}

\section{若干 Strichartz 估计的双线性证明}

考虑 $\mathbb{R}^{1+3}$ 中的齐次波动方程 $\Box u = 0$. 取 $q = r = 4$, $\gamma = 1/2$, Strichartz 估计 \eqref{eq:c8-331} 具有如下形式:
\[
\|u\|_{L^4(\mathbb{R}^{1+3})} \lesssim \|f\|_{\dot H^{1/2}} + \|g\|_{\dot H^{-1/2}}
\]
将 $u = u^{+} + u^{-}$ 分解, 只需证明
\begin{equation}
\|u^{+}\|_{L^4(\mathbb{R}^{1+3})} \lesssim \|f^{+}\|_{\dot H^{1/2}}\label{eq:c9-384}
\end{equation}
其中
\[
u^{+}(t, x) = \int e^{ix\cdot\xi + t|\xi|}\hat f(\xi)\, d\xi
\]
显然，
\[
\|u^{+}\|^2_{L^4(\mathbb{R}^{1+3})} = \|u^{+}\cdot u^{+}\|_{L^2} = \|\widetilde{u^{+}} * \widetilde{u^{+}}\|_{L^2}
\]
现在， 回忆 \eqref{eq:c8-346}， 并略去指标 +，
\[
\begin{aligned}
\tilde{u}*\tilde{u}(\tau, \xi) &= \int\int\delta(\tau - \lambda - |\xi-\eta|)\hat f(\xi-\eta)\delta(\lambda - |\eta|)\hat f(\eta)\, d\lambda\, d\eta\\
&= \int\delta(\tau - |\eta| - |\xi-\eta|)\hat f(\eta)\hat f(\xi-\eta)\, d\eta
\end{aligned}
\]
显然， \eqref{eq:c9-384} 由下面的事实推出:

\begin{theorem}\label{thm:c9-1.1}
\newterm{bilinear-operator}{双线性算子}{bilinear operator}
\[
B(F, G) = \int\delta(\tau - |\eta| - |\xi-\eta|)\frac{F(\xi-\eta)}{|\xi-\eta|^{1/2}}\frac{G(\eta)}{|\eta|^{1/2}}\, d\eta.
\]
满足估计
\begin{equation}
\|B(F, G)\|_{L^2(\mathbb{R}^{1+3})} \lesssim \|F\|_{L^2(\mathbb{R}^3)}\|G\|_{L^2(\mathbb{R}^{1+3})}\label{eq:c9-385}
\end{equation}
\end{theorem}

\begin{proof}
由 Cauchy-Schwarz 不等式，
\[
\begin{aligned}
|B(F, G)(\tau, \xi)|^2 &\lesssim J(\tau, \xi)\int\delta(\tau - |\eta| - |\xi-\eta|)|F(\xi-\eta)|^2|G(\eta)|^2\, d\eta\\
J(\tau, \xi) &= \int\delta(\tau - |\eta| - |\xi-\eta|)\frac{1}{|\xi-\eta|}\frac{1}{|\eta|}\, d\eta
\end{aligned}
\]
只需证明 $J$ 一致有界. 事实上， 若确实如此， 则
\[
\begin{aligned}
\|B(F, G)\|_{L^2(\mathbb{R}^{1+3})} &\lesssim \sup_{\tau, \xi}J(\tau, \xi)\int\int\delta(\tau - |\eta| - |\xi-\eta|)|F(\xi-\eta)|^2|G(\eta)|^2\, d\eta\, d\tau\, d\xi\\
&\lesssim \sup_{\tau, \xi}J(\tau, \xi)\|F\|_{L^2}^2\|G\|_{L^2}^2
\end{aligned}
\]
因此， \newterm{bilinear-estimate}{双线性估计}{bilinear estimate}是 $J$ 的一致有界性的直接推论. 而这由下面这个更一般的引理推出.
\end{proof}

\begin{lemma}\label{lem:c9-1.2}
设 $F$ 是两变量的任意函数， $J_F$ 为积分
\[
J_F^{\mp}(\tau, \xi) = \int_{\mathbb{R}^n}\delta(\tau - |\eta| \mp |\xi-\eta|)F(|\eta|, |\xi-\eta|)
\]
则，
\begin{equation}
J_F^{-}(\tau, \xi) = (\tau^2 - |\xi|^2)^{\frac{n-3}{2}}\int_{-1}^{1}F\big(\frac{\tau + x|\xi|}{2}, \frac{\tau - x|\xi|}{2}\big)(\tau^2 - x^2|\xi|^2)(1 - |x|^2)^{\frac{n-3}{2}}\, dx, \label{eq:c9-386}
\end{equation}
\begin{equation}
J_F^{+}(\tau, \xi) = (\tau^2 - |\xi|^2)^{\frac{n-3}{2}}\int_{1}^{\infty}F\big(\frac{\tau + x|\xi|}{2}, \frac{\tau - x|\xi|}{2}\big)(\tau^2 - x^2|\xi|^2)(1 - |x|^2)^{\frac{n-3}{2}}\, dx\label{eq:c9-387}
\end{equation}
\end{lemma}

\begin{proof}
注意到， 在 $\mp = -$ 的情形， 测度 $\delta(\tau - |\eta| - |\xi-\eta|)$ 支在焦点为 $0$ 与 $\xi$ 的\newterm{ellipsoid-of-revolution}{旋转椭球面}{ellipsoid of revolution} $\mathcal{E}(\tau, \xi) = \{\eta\in\mathbb{R}^n: |\eta| + |\xi-\eta| = \tau\}$ 上， 此时 $|\xi| \leq \tau$. 在 $\mp = +$ 的情形， 测度 $\delta(\tau - |\eta| + |\xi-\eta|)$ 支在焦点为 $0$ 与 $\xi$ 的\newterm{hyperboloid-of-revolution}{旋转双曲面}{hyperboloid of revolution} $\mathcal{H}(\tau, \xi) = \{\eta\in\mathbb{R}^n: |\eta| - |\xi-\eta| = \tau\}$ 上， 这是一个体积无穷的无界超曲面. 此时 $|\xi|^2 \leq \tau^2$. 在分布意义下， 我们有恒等式
\[
\begin{aligned}
\delta(\tau - |\eta| \mp |\xi-\eta|) &= \delta\big(\frac{(\tau - |\eta|)^2 - |\xi-\eta|^2}{2(\tau - |\eta|)}\big)\\
&= 2(\tau - |\eta|)\delta\big((\tau - |\eta|)^2 - |\xi-\eta|^2\big)\\
&= 2(\tau - |\eta|)\delta\big(\tau^2 - |\xi|^2 - 2\tau\lambda + 2\lambda|\xi|\cos\theta\big)\\
&= 2(\tau - |\eta|)\delta\big(\tau^2 - |\xi|^2 - 2\tau\lambda + 2a|\xi|\lambda\big)
\end{aligned}
\]
其中 $a$ 是 $\eta$ 与 $\xi$ 之间夹角的余弦. 于是， 对固定的 $\tau$ 与 $\xi$， 在测度的支集上必须有
\begin{equation}
a = -\frac{\tau^2 - |\xi|^2 - 2\tau\lambda}{2|\xi|\lambda}\label{eq:c9-388}
\end{equation}
注意到， 在椭球情形， $a$ 可以取区间 $[-1, 1]$ 中的任意值， 从而， 由于 $\lambda = \frac{\tau^2 - |\xi|^2}{2(\tau - a|\xi|)}$， 我们有
\[
\frac{\tau - |\xi|}{2} \leq \lambda \leq \frac{\tau + |\xi|}{2}.
\]
另一方面， 在 $|\xi|^2 > \tau^2$ 的双曲情形， 还必须有如下限制:
\[
\frac{\tau}{|\xi|} \leq a.
\]
从而 $\lambda = \frac{-\tau^2 + |\xi|^2}{2(-\tau + a|\xi|)} \geq \frac{\tau + |\xi|}{2}$.
于是， 由于 $d\eta = \lambda^{n-1}d\lambda\, dS_\omega = (1 - a^2)^{\frac{n-3}{2}}\lambda^{n-1}d\lambda\, dS_{\omega'}$，
\[
\begin{aligned}
J_F^{-} &= \frac{1}{|\xi|}\int_{\frac{\tau-|\xi|}{2}}^{\frac{\tau+|\xi|}{2}}F(\lambda, \tau-\lambda)(\tau-\lambda)\lambda^{n-2}\Big[1 - \Big(\frac{\tau^2 - |\xi|^2 - 2\tau\lambda}{2|\xi|\lambda}\Big)^2\Big]^{\frac{n-3}{2}}d\lambda\\
&= \frac{(\tau^2 - |\xi|^2)^{\frac{n-3}{2}}}{|\xi|^{n-2}}\int_{\frac{\tau-|\xi|}{2}}^{\frac{\tau+|\xi|}{2}}F(\lambda, \tau-\lambda)(\tau-\lambda)\lambda\Big[\Big(\frac{\tau+|\xi|}{2} - \lambda\Big)\Big(\lambda - \frac{\tau-|\xi|}{2}\Big)\Big]^{\frac{n-3}{2}}d\lambda
\end{aligned}
\]
最后， 我们作变量替换 $x = \frac{2\lambda - \tau}{|\xi|}$， 以导出所期望的公式 \eqref{eq:c9-386}. \eqref{eq:c9-387} 的证明与此相同.
\end{proof}

\section{改进的双线性 Strichartz 估计}

考虑两个齐次波动方程 $\Box u = \Box v = 0$ 的解. 为简单起见， 并且不失一般性， 我们假设 $u$， $v$ 在 $t = 0$ 时满足\newterm{reduced-initial-data}{约化初值}{reduced initial data}:
\[
u(0, x) = f(x), \quad v(0, x) = g(x), \quad \partial_t u(0, x) = \partial_t v(0, x) = 0.
\]
我们考虑如下形式的估计
\[
\|D^{-b}(uv)\|_{L_t^{q/2}L_x^{r/2}} \lesssim \|f\|_{\dot H^a}\|g\|_{\dot H^a}
\]
其中 $(q, r)$ 是一个\newterm{admissible-pair}{容许对}{acceptable pair}. 由量纲分析并回忆 \eqref{eq:c8-331} 中的指数 $\gamma = n\big((\frac{1}{2} - \frac{1}{r})\big) - \frac{1}{q}$， 我们必须有
\begin{equation}
2a = -b + 2\bigl(n\big(\frac{1}{2} - \frac{1}{r}\big) - \frac{1}{q}\bigr) = -b + 2\gamma\label{eq:c9-389}
\end{equation}
我们用三分公式分解乘积 $u\cdot v$:
\[
\begin{aligned}
u\cdot v &= \sum_{\mu<\lambda}u_\mu v_\lambda + \sum_{\mu<\lambda}v_\mu u_\lambda + \sum_{\mu\leq\lambda}P_\mu(u_\lambda v_\lambda)\\
&= (u\cdot v)_{LH} + (u\cdot v)_{HL} + (u\cdot v)_{HH}
\end{aligned}
\]
这里 $\mu, \lambda\in 2^{\mathbb{Z}}$， $u_\lambda = P_\lambda u$， 而 $P_\lambda$ 是通常的 \newterm{littlewood-paley-projection}{Littlewood-Paley 投影}{Littlewood-Paley projection}. 现在，
\[
\|D^{-b}(uv)_{LH}\|_{L_t^{q/2}L_x^{r/2}} \leq \sum_{\mu\leq\lambda}\lambda^{-b}\|u_\mu v_\lambda\|_{L_t^{q/2}L_x^{r/2}} \leq \sum_{\mu\leq\lambda}\lambda^{-b}\|u_\mu\|_{L_t^q L_x^r}\|v_\lambda\|_{L_t^q L_x^r}
\]
根据前一节的 Strichartz 估计
\[
\|u_\mu\|_{L_t^q L_x^r} \lesssim \mu^{(\gamma-a)}\|f_\mu\|_{\dot H^a} = \mu^{b/2}\|f_\mu\|_{\dot H^a}, \qquad
\|v_\lambda\|_{L_t^q L_x^r} \lesssim \lambda^{(\gamma-a)}\|g_\lambda\|_{\dot H^a} = \lambda^{b/2}\|g_\lambda\|_{\dot H^a}
\]
并且， 对 $b > 0$，
\[
\|D^{-b}(uv)_{LH}\|_{L_t^{q/2}L_x^{r/2}} \lesssim \sum_{\mu\leq\lambda}\big(\frac{\mu}{\lambda}\big)^b\|f_\mu\|_{\dot H^a}\|g_\lambda\|_{\dot H^a} \lesssim \|f\|_{\dot H^a}\|g\|_{\dot H^a}
\]
由对称性，
\[
\|D^{-b}(uv)_{HL}\|_{L_t^{q/2}L_x^{r/2}} \lesssim \|f\|_{\dot H^a}\|g\|_{\dot H^a}
\]
因此， 只剩下估计高-高项 $\|(u\cdot v)_{HH}\|_{L_t^{q/2}L_x^{r/2}}$. 这需要基于定理 ?? 的更为精细的论证. 我们写出
\[
\|D^{-b}(u\cdot v)_{HH}\|_{L_t^{q/2}L_x^{r/2}} \lesssim \sum_{\mu\leq\lambda}\mu^{-b}\|P_\mu(u_\lambda v_\lambda)\|_{L_t^{q/2}L_x^{r/2}}
\]
如果我们使用标准的 Strichartz 估计， 即
\begin{equation}
\|P_\mu(u_\lambda v_\lambda)\|_{L_t^{q/2}L_x^{r/2}} \lesssim \|u_\lambda\|_{L_t^q L_x^r}\|v_\lambda\|_{L_t^q L_x^r} = \lambda^{2(\gamma-a)}\|f\|_{\dot H^a}\|g\|_{\dot H^a} = \lambda^b\|f\|_{\dot H^a}\|g\|_{\dot H^a}\label{eq:c9-390}
\end{equation}
则会导出
\[
\|D^{-b}(u\cdot v)_{HH}\|_{L_t^{q/2}L_x^{r/2}} \lesssim \sum_{\mu\leq\lambda}\lambda^b\mu^{-b}\|f\|_{\dot H^a}\|g\|_{\dot H^a}
\]
这是发散的. 我们需要用更强的估计取代 \eqref{eq:c9-390}， 该估计要考虑 $P_\mu$ 位于 $u_\lambda v_\lambda$ 之前这一事实. 为此， 我们首先需要利用一些正交性性质. 我们在 Fourier 空间中把数据 $f_\lambda$， $g_\lambda$ 分解为支在大小为 $\mu$ 的立方体上的小块， $f_\lambda = \sum_Q f_Q$， $g_\lambda = \sum_Q g_Q$， 并用 $u_Q$， $v_Q$ 记相应的解. 显然， 该分解与波算子 $\Box$ 可交换. 于是， $u_\lambda \sim \sum_Q u_Q$， $v_\lambda \sim \sum_Q v_Q$， 并且
\[
P_\mu(u_\lambda\cdot v_\lambda) \sim \sum_{Q_1, Q_2}P_\mu\big(u_{Q_1}v_{Q_2}\big)
\]
注意到， 只有当 $Q_1 + Q_2$ 与 $P_\mu$ 支集的、大小为 $\mu$ 的频率区域相交时， $P_\mu(u_{Q_1}u_{Q_2}) \neq 0$. 对每个大小为 $\mu$ 的立方体 $Q_1$， 只有有限多个(其个数只依赖于 $n$)立方体 $Q_2$ 使得上述情形发生. 从道理上说， 必要时通过放大立方体， 我们可以假设 $Q_2 = -Q_1$， 于是
\[
P_\mu(u_\lambda\cdot v_\lambda) \sim \sum_Q u_Q v_{-Q}.
\]
因此，
\[
\|P_\mu(u_\lambda v_\lambda)\|_{L_t^{q/2}L_x^{r/2}} \lesssim \sum_Q\|u_Q v_{-Q}\|_{L_t^q L_x^r} \lesssim \sum_Q\|u_Q\|_{L_t^q L_x^r}\|v_{-Q}\|_{L_t^q L_x^r}.
\]
我们现在可以应用\thmref{thm:c8-0.42}. 于是，
\[
\|u_Q\|_{L_t^q L_x^r} \lesssim \big(\frac{\mu}{\lambda}\big)^{\frac{1}{2}-\frac{1}{r}}\|f_Q\|_{\dot H^\gamma}
\]
对 $v_{-Q}$ 也类似. 因此，
\[
\begin{aligned}
\|P_\mu(u_\lambda v_\lambda)\|_{L_t^{q/2}L_x^{r/2}} &\lesssim \big(\frac{\mu}{\lambda}\big)^{1-\frac{2}{r}}\sum_Q\|f_Q\|_{\dot H^\gamma}\|g_Q\|_{\dot H^\gamma}\\
&\lesssim \big(\frac{\mu}{\lambda}\big)^{1-\frac{2}{r}}\|f_\lambda\|_{\dot H^\gamma}\|g_\lambda\|_{\dot H^\gamma}\\
&\lesssim \big(\frac{\mu}{\lambda}\big)^{1-\frac{2}{r}}\lambda^{2\gamma-2a}\|f_\lambda\|_{\dot H^a}\|g_\lambda\|_{\dot H^a}\\
&\lesssim \big(\frac{\mu}{\lambda}\big)^{1-\frac{2}{r}}\lambda^b\|f_\lambda\|_{\dot H^a}\|g_\lambda\|_{\dot H^a}
\end{aligned}
\]
并且， 从而，
\[
\begin{aligned}
\|D^{-b}(u\cdot v)_{HH}\|_{L_t^{q/2}L_x^{r/2}} &\lesssim \sum_{\mu<\lambda}\big(\frac{\mu}{\lambda}\big)^{1-\frac{2}{r}-b}\|f_\lambda\|_{\dot H^a}\|g_\lambda\|_{\dot H^a}\\
&\lesssim \|f\|_{\dot H^a}\|g\|_{\dot H^a}
\end{aligned}
\]
只要 $b < 1 - \frac{2}{r}$. 我们刚刚证明了如下的双线性估计， 参见 [?].

\begin{theorem}\label{thm:c9-2.1}
下面的估计\footnote{此处 $\|u[0]\|_{\dot H^a} = \|u(0)\|_{\dot H^a} + \|\partial_t u(0)\|_{\dot H^a}$.}对 $\Box u = \Box v = 0$ 的解、任意容许对 $(q, r)$ 以及任意 $0 \leq b < 1 - \frac{2}{r}$ 成立:
\begin{equation}
\|D^{-b}(u\cdot v)\|_{L_t^{q/2}L_x^{r/2}} \lesssim \|u[0]\|_{\dot H^a}\|v[0]\|_{\dot H^a}\label{eq:c9-391}
\end{equation}
只要如下的\newterm{dimensional-condition}{量纲条件}{dimensional condition}成立:
\begin{equation}
a = -\frac{b}{2} + \gamma, \qquad \gamma = n\big(\frac{1}{2} - \frac{1}{r}\big) - \frac{1}{q}\label{eq:c9-392}
\end{equation}
\end{theorem}

\section{零形式的双线性估计}

在本小节中， 我们讨论零二次型的最简单的双线性估计， 参见 [?]， [?]， [?] 与 [?].

\begin{definition}\label{def:c9-3.1}
设 $u$， $v$ 是 $\mathbb{R}^{n+1}$ 上满足 $\Box u = \Box v = 0$ 的两个光滑解. \newterm{null-quadratic-form}{标准零二次型}{null quadratic form}为
$Q_0(u, v) = -\partial_t u\, \partial_t v + \sum_{i=1}^{n}\partial_i u\, \partial_i v$，
以及
$Q_{ij}(u, v) = \partial_i u\, \partial_j v - \partial_i v\, \partial_j u$，
和
$Q_{0i}(u, v) = \partial_i u\, \partial_t v - \partial_i v\, \partial_t u$，
其中 $i, j = 1, \ldots, n$.
\end{definition}

\begin{theorem}\label{thm:c9-3.2}
对任意\newterm{null-form}{零形式}{null form} $Q$ 以及 $\mathbb{R}^{n+1}$ 上满足 $\Box u = \Box v = 0$ 的任意解， $n \geq 2$， 我们有
\begin{equation}
\|Q(u, v)\|_{L^2(\mathbb{R}^{n+1})} \lesssim \|u[0]\|_{\dot H^1(\mathbb{R}^n)}\|v[0]\|_{H^{\frac{n+1}{2}}(\mathbb{R}^n)}\label{eq:c9-393}
\end{equation}
\end{theorem}

\begin{remark}\label{rem:c9-3.3}
不失一般性， 只需考虑约化初值问题
\begin{equation}
u(0, x) = f(x), \quad v(0, x) = g(x), \quad \partial_t u(0, x) = \partial_t v(0, x) = 0\label{eq:c9-394}
\end{equation}
下面我们说明如何从下一节给出的更一般的双线性估计形式推导出 \thmref{thm:c9-3.2} 中的估计.
\end{remark}

\begin{definition}\label{def:c9-3.4}
设 $D^{\alpha}$， $D_+^{\alpha}$ 与 $D_-^{\alpha}$ 是 $\mathbb{R}^{n+1}$ 中分别由具有如下\newterm{symbol}{象征}{symbol}的乘子所定义的算子:
\[
|\xi|^{\alpha}, \qquad (|\tau| + |\xi|)^{\alpha}, \qquad \big||\tau| - |\xi|\big|^{\alpha}.
\]
\end{definition}

注意到， 对任意光滑函数 $u$， $v$， 我们可以写出
\[
2Q_0(u, v) = \Box(uv) - \Box u\, v - u\, \Box v
\]
于是， 若 $\Box u = \Box v = 0$， 利用 Plancherel 恒等式，
\[
\begin{aligned}
\|Q_0(u, v)\|_{L^2(\mathbb{R}^{n+1})} &\leq \frac{1}{2}\|\Box(uv)\|_{L^2(\mathbb{R}^{n+1})} = \frac{1}{2}(2\pi)^{-n}\|(\tau^2 - |\xi|^2)\widetilde{uv}\|_{L^2(\mathbb{R}^{n+1})}\\
&\lesssim \|D_+ D_-(uv)\|_{L^2(\mathbb{R}^{n+1})}
\end{aligned}
\]
因此，
\begin{equation}
\|Q_0(u, v)\|_{L^2(\mathbb{R}^{n+1})} \leq \|D_+ D_-(uv)\|_{L^2(\mathbb{R}^{n+1})}\label{eq:c9-395}
\end{equation}
于是， 在零形式 $Q_0$ 的情形， \thmref{thm:c9-3.2} 归结为
\begin{equation}
\|D_+ D_-(uv)\|_{L^2(\mathbb{R}^{n+1})} \lesssim \|u[0]\|_{\dot H^1(\mathbb{R}^n)}\|v[0]\|_{H^{\frac{n+1}{2}}(\mathbb{R}^n)}\label{eq:c9-396}
\end{equation}
这是定理 ?? 的一个特殊情形.

下面我们证明， 对其他零形式 $Q_{ij}$， $Q_{0i}$， 类似的估计也成立.

\begin{remark}\label{rem:c9-3.5}
给定 $\Box u = 0$ 的一个解 $u$， 其初值为 $u(0, x) = f(x)$， $\partial_t u(0, x) = 0$， 我们用 $u'$ 记同一方程的、具有数据 $u'(0, x) = f'(x)$， $\partial_t u'(0, x) = 0$ 的解， 其中 $f' = \mathcal{F}^{-1}(|\hat f|)$. 当然， 注意到 $\|f'\|_{\dot H^a} = \|f\|_{\dot H^a}$， 因此， 从我们正在考虑的 $L^2$ 型估计的角度看， $u$ 与 $u'$ 是无法区分的.
\end{remark}

\begin{proposition}\label{prop:c9-3.6}
设 $u$， $v$ 是齐次波动方程的带初值的光滑解. 下面的估计成立:
\begin{equation}
\|Q_{ij}(u, v)\|_{L^2(\mathbb{R}^{n+1})} \lesssim \|D^{1/2}D_-^{1/2}\big(D^{1/2}u'\cdot D^{1/2}v'\big)\|_{L^2(\mathbb{R}^{n+1})}\label{eq:c9-397}
\end{equation}
\begin{equation}
\|Q_{0i}(u, v)\|_{L^2(\mathbb{R}^{n+1})} \lesssim \|D_+^{1/2}D_-^{1/2}\big(D^{1/2}u'\cdot D^{1/2}v'\big)\|_{L^2(\mathbb{R}^{n+1})}\label{eq:c9-398}
\end{equation}
\end{proposition}

\begin{proof}
我们首先像前面一样分解 $u = u^{+} + u^{-}$， $v = v^{+} + v^{-}$. 在 Fourier 变量中写出
\[
\widetilde{Q_{ij}(u^{+}, v^{\pm})}(\tau, \xi) = \int q_{ij}(\eta, \xi-\eta)\delta(\tau - |\eta| \pm |\xi-\eta|)\hat f(\eta)\hat g(\xi-\eta)\, d\eta
\]
其中 $q_{ij}(\eta, \xi-\eta) = \eta_i(\xi-\eta)_j - \eta_j(\xi-\eta)_i = (\xi\wedge\eta)_{ij}$. 现在我们依赖下面这个简单的引理.

\begin{lemma}\label{lem:c9-3.7}
下面的不等式成立:
\begin{equation}
|\xi \wedge \eta| \lesssim |\xi|^{1/2}|\eta|^{1/2}|\xi+\eta|^{1/2}(|\xi| + |\eta| - |\xi+\eta|)^{1/2}\label{eq:c9-399}
\end{equation}
\begin{equation}
|\xi \wedge \eta| \lesssim |\xi|^{1/2}|\eta|^{1/2}|\xi+\eta|^{1/2}\big(|\xi+\eta| - \big||\xi| - |\eta|\big|\big)^{1/2}\label{eq:c9-400}
\end{equation}
\end{lemma}

我们确实有
\[
\begin{aligned}
4|\xi \wedge \eta|^2 &= 4(|\xi||\eta| - \xi\cdot\eta)(|\xi||\eta| + \xi\cdot\eta)\\
&= (|\xi| + |\eta| - |\xi+\eta|)(|\xi| + |\eta| + |\xi+\eta|)\big(|\xi+\eta| - \big||\xi| - |\eta|\big|\big)\big(|\xi+\eta| + \big||\xi| - |\eta|\big|\big)
\end{aligned}
\]
由此引理立即得证.

因此， 在两种情形下， $|\widetilde{Q_{ij}(u^{+}, v^{\pm})}(\tau, \xi)|$ 都可以被下面的表达式界定:
\[
\begin{aligned}
&\int |q_{ij}(\eta, \xi-\eta)|\, \delta(\tau - |\eta| \pm |\xi-\eta|)\, |\hat f(\eta)|\, |\hat g(\xi-\eta)|\, d\eta\\
&\lesssim \big||\tau| - |\xi|\big|^{1/2}|\xi|^{1/2}\int\delta(\tau - |\eta| \pm |\xi-\eta|)\, |\eta|^{1/2}|\xi-\eta|^{1/2}|\hat f(\eta)|\, |\hat g(\xi-\eta)|\, d\eta\\
&= D^{1/2}D_-^{1/2}(D^{1/2}u'\cdot D^{1/2}v')
\end{aligned}
\]
如所愿.
\end{proof}

根据\propref{prop:c9-3.6}， 对于 $Q = Q_{ij}$ 与 $Q = Q_{0i}$， \thmref{thm:c9-3.2} 分别归结为如下命题:
\[
\begin{aligned}
\|D^{1/2}D_-^{1/2}(u\cdot v)\|_{L^2(\mathbb{R}^{n+1})} &\lesssim \|u[0]\|_{\dot H^{1/2}}\cdot\|v[0]\|_{\dot H^{n/2}}\\
\|D_+^{1/2}D_-^{1/2}(u\cdot v)\|_{L^2(\mathbb{R}^{n+1})} &\lesssim \|u[0]\|_{\dot H^{1/2}}\cdot\|v[0]\|_{\dot H^{n/2}}
\end{aligned}
\]
它们是定理 ?? 的特殊情形.

\appendix
% ==== partapx 附录 A:Riemann 与 Lorentz 几何的基本概念 + 参考文献 ====
% 源: MinerU_latex_A-非.tex 第 718-2283 行;校对 pdf_text.txt 第 11552-12658 行(PDF 第 213-234 页)
% 术语: 时空=spacetime, 符号差=signature, 可定向的=orientable, 切空间=tangent space,
%   惯性坐标=inertial coordinates, 零锥=null cone, 类光的=null, 零对=null pair,
%   外微分=exterior derivative, Lie 括号=Lie bracket, k 维分布=k-distribution,
%   对合的=involutive, 可积的=integrable, Levi-Civita 联络=Levi-Civita connection,
%   Christoffel 符号=Christoffel symbols, 平行移动=parallel transport, 仿射参数=affine parameter,
%   世界线=world line, 固有时=proper time, 指数映射=exponential map, 测地完备的=geodesically complete,
%   法邻域=normal neighbourhood, Weyl 张量=Weyl tensor, 共形的=conformal, 等距=isometry,
%   共形等距=conformal isometry, 共形 Killing 向量场=conformal Killing vector field,
%   变形张量=deformation tensor, 体积形式=volume form, 散度定理=divergence theorem,
%   透镜形区域=lens shaped domain, D'Alembert 算子=D'Alembertian, 时间函数=time function,
%   程函方程=eikonal equation, 因果未来=causal future, 因果过去=causal past,
%   保守区域=conservative domain, 零测地线生成元=null geodesic generator,
%   零基本形式=null fundamental form, 测地叶状结构=geodesic foliation,
%   膨胀系数=expansion coefficient, 零时距函数=null lapse function,
%   正能量条件=positive energy condition, 无散度的=divergenceless, 焦散=caustic,
%   逆平移=inverted translations, 线汇=congruence, Poincaré 群=Poincaré group,
%   共形群=conformal group, 共形不变的=conformally invariant
% REFS:
%   - 源 986 行 "See [?], Hawking and Ellis, section 2.4" 与源 1347-1348 行
%     "(see Hawking-Ellis, [?], chapter 2, section 2.6)": 原书 [?] 引文占位符缺失,已省略
%   - 源 1392 行 "citeHawkEll, chapter 2, section 2.6": OCR 破损引用,按 Hawking--Ellis 第 2 章第 2.6 节处理
%   - 源 2207 行 "According to eqrefeq:cons-law1": 破损引用,按上下文解析为 \eqref{eq:cA-443}(即刚推出的 (443))
%   - 脚注 1 中 "The large scale structure of space-time, [?]" 的 [?] 占位符省略
%   - OCR-BROKEN: 源 1439 行 D/dt 系统;源 1595 行分母 (k-1)!(原书笔误,应作 (n-1)!,照抄);
%     源 1779 行 Σ_2 也取 t_1(原书笔误,应作 t_2,照抄);源 1652 行 "12lodtsn" 系 OCR 乱码,
%     还原为 ε_{12...n}
% 编号说明:
%   - 公式 (401)-(446) 共 46 个,按原文顺序自动编号;原书未编号的行间公式用 [ ] 不占号
%   - 定理类编号与原书一致: 第 2 节 2.1-2.6、第 3 节 3.1-3.8、第 4 节 4.1-4.7
%   - 原书第 3 节 3.2 内的三个未编号 Definition 与一个未编号 Remark 用带 * 环境
%   - 原书第 4 节 Theorem 4.3 之后的未编号 Remark(437 式)用 remark*
%   - 附录公式从 (401) 起自动接续(主文件前 9 章公式止于 (400))
%   - 参考文献 [Kl-Rodn2] 两条显示标签相同,照原书;内部键分 kl-rodn2a / kl-rodn2b
% ====
\chapter{Riemann 与 Lorentz 几何的基本概念}

\section{引言}

在下面的内容中, 我们简要概述 Riemann 与 Lorentz 几何的基本概念. 关于更详细的综述, 我们分别参见 \emph{\cite{pet}}(Riemann 几何)以及 \emph{\cite{car}}、\emph{\cite{ha-e}}、\emph{\cite{wa}}(Lorentz 几何).

一个\newterm{pseudo-riemannian-manifold}{伪 Riemann 流形}{pseudo-Riemannian manifold}\footnote{我们假定读者已经熟悉微分几何的基本概念, 如流形、张量场、协变、Lie 与外微分. 关于这些概念的简短介绍, 参见 Hawking 与 Ellis 的《The large scale structure of space-time》第 2 章.}, 或简称为时空, 由一对 $(\mathbf{M}, \mathbf{g})$ 组成, 其中 $\mathbf{M}$ 是一个\newterm{orientable}{可定向的}{orientable} $p+q=n$ 维流形, $\mathbf{g}$ 是定义在其上的伪 Riemann 度量, 即一个光滑的、非退化的、符号差为 $(p, q)$ 的二阶协变对称张量场. 这意味着在每一点 $p\in\mathbf{M}$ 可以选取\newterm{tangent-space}{切空间}{tangent space} $T\mathbf{M}_p$ 中由 $p+q$ 个向量组成的基 $\{e_{(\alpha)}\}$, 使得

\begin{equation}
\mathbf{g}(e_{(\alpha)}, e_{(\beta)}) = \eta_{\alpha\beta}\label{eq:cA-401}
\end{equation}

对一切 $\alpha, \beta=0, 1, \ldots, n$ 成立, 其中 $\eta$ 是对角矩阵, 其前 $p$ 个对角元为 $-1$, 后 $q$ 个对角元为 $+1$. 若 $X$ 是 $p$ 处的任意向量, 按基 $\{e_{(\alpha)}\}$ 表为 $X=X^{\alpha}e_{(\alpha)}$, 则我们有

\begin{equation}
\mathbf{g}(X, X) = -(X^{1})^{2} - \dots - (X^{p})^{2} + (X^{p+1})^{2} + \dots + (X^{p+q})^{2}\label{eq:cA-402}
\end{equation}

当 $p=0$、$q=n$ 时, 对应于 $n$ 维 Riemann 流形. 我们感兴趣的另一个情形是 $p=1$、$q=n$, 它对应于 $n+1$ 维 Lorentz 流形. Riemann 流形的主要例子是 Euclidean 空间 $R^{n}$. 任何其他 Riemann 流形在局部上都与 $R^{n}$ 相似. 类似地, Lorentz 流形的主要例子是 Minkowski 时空, 即狭义相对论的时空. 它在 Lorentz 几何中的地位与 Euclidean 空间在 Riemann 几何中的地位相同. 此时流形 $M$ 与 $R^{n+1}$ 微分同胚, 并且存在整体定义的坐标系 $x^{\alpha}$, 相对于该坐标系度量取对角形式 $-1, 1, \ldots, 1$. 所有这些坐标系都通过 Lorentz 变换相联系, 并称为\newterm{inertial-coordinates}{惯性坐标}{inertial coordinates}. 我们将 $n+1$ 维 Minkowski 时空记作 $(\mathbb{R}^{n+1}, \mathbf{m})$.

相对于给定的坐标系 $x^{\mu}$, 伪 Riemann 度量的分量取如下形式

\[
g_{\mu\nu} = \mathbf{g}(\partial_{\mu}, \partial_{\nu})
\]

其中 $\partial_{\mu}=\frac{\partial}{\partial x^{\mu}}$ 是相应的坐标向量场. 我们用 $g^{\mu\nu}$ 表示逆度量 $g^{-1}$ 在同一坐标系 $x$ 下的分量, 用 $|g|$ 表示矩阵 $g_{\mu\nu}$ 的行列式. $\mathbf{M}$ 的体积元 $dv_{\mathbf{M}}$ 在局部坐标下表示为 $\sqrt{|g|}dx=\sqrt{|g|}dx^{1}\ldots dx^{n}$. 因此, 支集在坐标卡 $U\subset M$ 中的函数 $f$ 的积分 $\int_{M}fdv_{M}$ 定义为 $\int_{U}f(x)\sqrt{|g(x)|}dx$. 任意函数 $f$ 在 $M$ 上的积分通过作从属于坐标卡覆盖的单位分解来定义. 容易验证该定义与局部坐标系的选取无关.

由 \eqref{eq:cA-402} 可见, Lorentz 度量把每一点 $p$ 处切空间 $TM_{p}$ 中的向量分为\newterm{timelike-vector}{类时}{timelike}、\newterm{null-vector}{类光}{null}或\newterm{spacelike-vector}{类空}{spacelike}三类, 分别取决于二次型

\begin{equation}
(X, X) = g_{\mu\nu} X^{\mu} X^{\nu}\label{eq:cA-403}
\end{equation}

为负、为零或为正. 类光向量的集合 $N_{p}$ 构成一个双锥, 称为点 $p$ 的零锥. 类时向量的集合 $I_{p}$ 构成该锥的内部. $I_{p} \cup N_{p}$ 中的向量称为\newterm{causal-vector}{因果的}{causal}. 类空向量的集合 $S_{p}$ 是 $I_{p} \cup N_{p}$ 的补集.

满足 \eqref{eq:cA-401} 的标架 $e_{(\alpha)}$ 称为标准正交标架. 在 Lorentz 流形的情形下, 除了标准正交标架之外, 考虑零标架也是有意义的. 零标架是这样的向量场\footnote{只要不会引起混淆, 为简化记号我们写 $e_{\alpha}$ 而非 $e_{(\alpha)}$.}集合 $e_{\alpha}$, 它由两个类光向量 $e_{n+1}, e_{n}$ 与标准正交的类空向量 $(e_{a})_{a=1, \ldots, n-1}$ 组成, 满足

\[
\begin{array}{l}
\mathbf{g}(e_{n}, e_{n}) = \mathbf{g}(e_{n+1}, e_{n+1}) = 0, \quad \mathbf{g}(e_{n}, e_{n+1}) = -2 \\
\mathbf{g}(e_{n}, e_{a}) = \mathbf{g}(e_{n+1}, e_{a}) = 0, \quad \mathbf{g}(e_{a}, e_{b}) = \delta_{ab}
\end{array}
\]

一次形式 $A = A_{\alpha}dx^{\alpha}$ 是 $M$ 的余切丛的截面. 我们用 $A(X)$ 表示 $A$ 与向量场 $X$ 之间的自然配对. 我们可以通过 $A^{\alpha} = g^{\alpha\beta}A_{\beta}$ 提升 $A$ 的指标. $A' = A^{\alpha}\partial_{\alpha}$ 定义 $M$ 上的向量场, 且有 $A(X) = \mathbf{g}(A', X)$. k 阶协变张量 $A$ 是 TM 上的 $k$ 重线性形式.

\par\noindent\textbf{记号:}我们将使用以下记号约定:用粗体字母表示重要的张量, 如度量 $g$ 与 Riemann 曲率张量 R. 它们相对于任意标架的分量也用粗体字母表示. 因此, 给定标架 $\{e_{(\alpha)}\}$, 我们写 $\mathbf{g}_{\alpha\beta} = \mathbf{g}(e_{\alpha}, e_{\beta})$、$\mathbf{R}_{\alpha\beta\gamma\delta} = \mathbf{R}(e_{\alpha}, e_{\beta}, e_{\gamma}, e_{\delta})$, 并且对任意张量 T,

\[
T_{\alpha\beta\gamma\delta\dots} \equiv T(e_{\alpha}, e_{\beta}, e_{\gamma}, e_{\delta}, \dots)
\]

在 Riemann 流形的情形下, 我们常用拉丁字母 $i, j, k, l, \ldots$ 表示坐标 $x^{1}, x^{2}, \ldots, x^{n}$ 或张量的指标. 对 Lorentz 流形, 我们用希腊字母 $\alpha, \beta, \gamma, \ldots$ 表示指标 $0, 1, \ldots, n$. 对一般的符号差为 $s$ 的伪 Riemann 度量, 我们也使用希腊指标.

下面我们将回顾以下主题:

\begin{enumerate}
\renewcommand{\labelenumi}{\arabic{enumi}.)}
\item 向量场的 Lie 括号. Frobenius 定理
\item 张量场的 Lie 导数
\item 多重线性形式与外微分
\item 联络与协变导数
\item 伪 Riemann 度量. Riemann 与 Lorentz 几何.
\item 与伪 Riemann 度量相伴的 Levi-Civita 联络.
\item 平行移动、测地线、指数映射、完备性
\item 伪 Riemann 流形的曲率张量. 对称性. 第一与第二 Bianchi 恒等式.
\item 等距与共形等距. Killing 与共形 Killing 向量场.
\end{enumerate}


\section{各种微分概念}

这里我们回顾 Riemann 或 Lorentz 流形上微分几何的三个基本算子:外微分、\newterm{lie-derivative}{Lie 导数}{Lie derivative}, 以及 Levi-Civita 联络及其相应的\newterm{covariant-derivative}{协变导数}{covariant derivative}.

\mtsubsec{外微分}

给定标量函数 $f$, 其微分 $df$ 是由下式定义的一次形式

\[
df(X) = X(f)
\]

对任意向量场 X. 这个定义可以按如下方式推广到 $M$ 上的所有微分形式, 即 $\Omega^{k}(\mathbf{M})$:

\begin{enumerate}
\renewcommand{\labelenumi}{\roman{enumi})}
\item $d$ 是从 $\mathbf{M}$ 上所有 $k$ 形式的空间到 $k+1$ 形式的空间的线性算子. 因此对一切 $k$ 形式 $A$、$B$ 与实数 $\lambda, \mu$,
\end{enumerate}

\[
d(\lambda A + \mu B) = \lambda dA + \mu dB
\]

\begin{enumerate}
\renewcommand{\labelenumi}{\roman{enumi})}
\setcounter{enumi}{1}
\item 对任意 $k$ 形式 $A$ 与任意形式 B,
\end{enumerate}

\[
d(A \wedge B) = dA \wedge B + (-1)^{k} A \wedge dB
\]

\begin{enumerate}
\renewcommand{\labelenumi}{\roman{enumi})}
\setcounter{enumi}{2}
\item 对任意形式 A,
\end{enumerate}

\[
d^{2}A = 0.
\]

我们回忆, 若 $\Phi$ 是从 $M$ 到另一流形 $M'$ 的光滑映射, 则

\[
d(\Phi^{*}A) = \Phi^{*}(dA)
\]

最后, 若 $A$ 是一次形式而 $X$、$Y$ 是任意向量场, 我们还有方程

\[
dA(X, Y) = \frac{1}{2}\left(X(A(Y)) - Y(A(X)) - A([X, Y])\right)
\]

其中 $[X, Y]$ 是换位子 $X(Y)-Y(X)$. 这可以很容易地推广到任意的 $k$ 形式, 参见 \emph{\cite{sp}} 第 I 卷第 7 章定理 13.

\mtsubsec{Lie 导数}

考虑任意向量场 X. 在局部坐标 $x^{\mu}$ 下, X 的流由微分方程组

\[
\frac{dx^{\mu}}{dt} = X^{\mu}(x^{1}(t), \dots, x^{p+q}(t)).
\]

给出. 相应的曲线 $x^{\mu}(t)$ 称为 $X$ 的积分曲线. 对每一点 $p \in \mathbf{M}$, 存在一个开邻域 $\mathcal{U}$、一个小量 $\epsilon > 0$ 和一族微分同胚 $\Phi_t: \mathcal{U} \to \mathbf{M}$, $|t| \leq \epsilon$, 它把 $\mathcal{U}$ 中的每一点沿 $X$ 的积分曲线送到参数距离为 $t$ 处. 我们利用这些微分同胚来构造:对 $p$ 处的任意张量 T, $\Phi_t(p)$ 处的张量族 $(\Phi_t)_*T$.

张量场 $T$ 关于 $X$ 的 Lie 导数 $L_{X}T$ 为

\[
\mathcal{L}_{X}T|_{p} \equiv \lim_{t \to 0} \frac{1}{t}\left(T|_{p} - (\Phi_{t})_{*}T|_{p}\right).
\]

它具有以下性质:

\begin{enumerate}
\renewcommand{\labelenumi}{\roman{enumi})}
\item $L_{X}$ 线性地把 $(p, q)$ 型张量场映为同型的张量场.
\item $\mathcal{L}_X$ 与缩并可交换.
\item 对任意张量场 $S$、T,
\end{enumerate}

\[
\mathcal{L}_{X}(S \otimes T) = \mathcal{L}_{X}S \otimes T + S \otimes \mathcal{L}_{X}T.
\]

若 $X$ 是向量场, 容易验证

\[
\mathcal{L}_{X}Y = [X, Y]
\]

方法是写出 $(\mathcal{L}_XY)^i = -\frac{d}{dt}((\Phi_t)_*Y)^i\Big|_{t = 0}$, 并把 $((\Phi_t)_*Y)^i\Big|_p = \frac{\partial x^i(\Phi_t(q))}{\partial x^j(q)} Y^j\Big|_q$ 代入, 其中 $q = \Phi_{-t}(p)$. (详见 Hawking 与 Ellis 第 2.4 节.)

若 $A$ 是 $k$ 形式, 作为外微分与回拉 $\Phi^{*}$ 的交换公式的推论, 我们有

\[
d(\mathcal{L}_{X}A) = \mathcal{L}_{X}(dA).
\]

对给定的 $k$ 阶协变张量场 T, 我们有

\[
\mathcal{L}_{X}T(Y_{1}, \dots, Y_{k}) = XT(Y_{1}, \dots, Y_{k}) - \sum_{i=1}^{k} T(Y_{1}, \dots, \mathcal{L}_{X}Y_{i}, \dots, Y_{k})
\]

我们注意到两个坐标向量场的 Lie 括号为零,

\[
\left[\frac{\partial}{\partial x^{\mu}}, \frac{\partial}{\partial x^{\nu}}\right] = 0.
\]

其逆命题也成立, 即(参见 \emph{\cite{sp}} 第 I 卷第 5 章):

\begin{proposition}\label{prop:cA-2.3}
若 $X_{(0)}, \ldots, X_{(k)}$ 是点 $p$ 的邻域中的线性无关向量场, 且其中任意两个的 Lie 括号均为零, 则存在 $p$ 附近的坐标系 $x^{\mu}$, 使得对每个 $\rho = 0, \ldots, k$ 有 $X_{(\rho)} = \frac{\partial}{\partial x^{\rho}}$.
\end{proposition}

\begin{proof}
当 k = 2 时, 该命题等价于如下事实:若两个线性无关的向量场 $X$、$Y$ 可交换, 即 $[X, Y] = 0$, 则它们的流在如下意义下也可交换. 从任意点 $p$ 出发, 若先沿从 $p$ 出发的 $X$ 的积分曲线行进参数距离 t, 再沿 $Y$ 的积分曲线行进参数距离 s, 则到达的点与先沿 $Y$ 的积分曲线行进距离 $s$、再沿 $X$ 的积分曲线行进距离 $t$ 所到达的点相同. 为了证明该断言, 我们首先考虑相对于 $X = \partial_{t}$ 的坐标 t, x. Y 在同一坐标系下的系数与 $t$ 无关. 然后作另一个坐标变换 $x = x(y)$, 使得 $Y = A(y)\partial_{t} + \partial_{y^{1}}$. 最后可以再作形如 $s = a(y) + t$ 的坐标变换, 其中 $a$ 是 $A$ 关于 $y^{1}$ 的原函数. 在新坐标 s, y 下, 必须有 $X = \partial_{s}$, $Y = \partial_{y^{1}}$, 正如所愿. 一般情形可以同样证明.
\end{proof}

上述命题是证明 Frobenius 定理的主要步骤. 为了陈述该定理, 我们回忆 $M$ 中 $k$ 维\newterm{k-distribution}{分布}{k-distribution}的定义. 所谓 $k$ 维分布, 是指在 $U$ 的定义域中每一点处光滑地指定一个 $k$ 维平面 $\pi_{p}$. 若对 $U$ 上的任意向量场 $X$、Y, 只要对任意 $p \in U$ 有 $X|_{p}, Y|_{p} \in \pi_{p}$, 就有 $[X, Y]|_{p} \in \pi_{p}$, 则该分布称为\newterm{involutive}{对合的}{involutive}. 对\newterm{integrable}{可积的}{integrable}分布\footnote{回忆:$U$ 上的分布 $\pi$ 称为可积的, 如果通过每一点 $p \in U$ 恰有一条维数为 $k$ 的子流形 N, 使得 $\pi_p = TN_p$.}显然如此. 事实上, 若对一切 $p \in N$ 有 $X|_{p}, Y|_{p} \in TN_{p}$, 则 $X$、$Y$ 切于 N, 它们的换位子 $[X, Y]$ 亦然. Frobenius 定理断言其逆命题也成立\footnote{证明参见 Spivak 的《[Sp]》第 I 卷第 6 章.}, 即处于对合状态也是分布可积的充分条件:

\begin{theorem}[Frobenius 定理]\label{thm:cA-2.4}
分布 $(\pi_p)_{p\in\mathcal{U}}$ 可积的充分必要条件是它是对合的.
\end{theorem}

\begin{proof}
若 $\pi$ 是对合的, 我们可以找到线性无关的向量场 $X_{1}, \ldots, X_{k}$, 它们在每一点张成 $\pi$, 并且其中任意两个的 Lie 括号都是 $X_{1}, \ldots, X_{k}$ 的线性组合. 然后可以重新定义 $X_{1}^{\prime}$ 为其他向量场的线性组合, 使得它与 $X_{2}, \ldots, X_{k}$ 可交换. 接着可以归纳地重新定义其余向量场, 使它们两两可交换. 证明于是由前一命题得出.
\end{proof}

\mtsubsec{联络与协变导数}

\newterm{connection}{联络}{connection} D 是这样一种法则:它给每个向量场 $X$ 指派一个微分算子 $D_{X}$. 该算子把向量场 $Y$ 映为向量场 $D_{X}Y$, 且对 $\alpha, \beta \in R$ 与 $M$ 上的标量函数 $f$、g, 满足

\begin{equation}
\begin{array}{l}
a) \quad \mathbf{D}_{fX + gY}Z = f\mathbf{D}_{X}Z + g\mathbf{D}_{Y}Z \\
b) \quad \mathbf{D}_{X}(\alpha Y + \beta Z) = \alpha\mathbf{D}_{X}Y + \beta\mathbf{D}_{X}Z \\
c) \quad \mathbf{D}_{X}fY = X(f)Y + f\mathbf{D}_{X}Y
\end{array}\label{eq:cA-404}
\end{equation}

因此, 在点 $p$ 处,

\begin{equation}
\mathbf{D}Y \equiv Y^{\alpha}_{;\beta}\, \theta^{(\beta)} \otimes e_{(\alpha)}\label{eq:cA-405}
\end{equation}

其中 $\theta^{(\beta)}$ 是关于标准正交标架 $e_{(\beta)}$ 的对偶基的一次形式. 注意 $Y_{;\beta}^{\alpha} = \theta^{(\alpha)}(\mathbf{D}_{e_{(\beta)}}Y)$. 另一方面, 由 c) 得

\[
\mathbf{D}(fY) = df \otimes Y + f\mathbf{D}Y
\]

所以

\[
\mathbf{D}Y = \mathbf{D}(Y^{\alpha}e_{(\alpha)}) = dY^{\alpha} \otimes e_{(\alpha)} + Y^{\alpha}\mathbf{D}e_{(\alpha)}
\]

最后, 利用 $df(\cdot) = e_{(\alpha)}(f)\theta^{(\alpha)}(\cdot)$, 得

\begin{equation}
\mathbf{D}Y = \left(e_{(\beta)}(Y^{\alpha}) + Y^{\gamma}\theta^{(\alpha)}(\mathbf{D}_{e_{(\beta)}}e_{(\gamma)})\right)\theta^{(\beta)} \otimes e_{(\alpha)}\label{eq:cA-406}
\end{equation}

因此

\[
Y^{\alpha}_{;\beta} = e_{(\beta)}(Y^{\alpha}) + \Gamma^{\alpha}_{\beta\gamma}Y^{\gamma}
\]

从而联络由其联络系数

\begin{equation}
\Gamma^{\alpha}_{\beta\gamma} = \theta^{(\alpha)}(\mathbf{D}_{e_{(\beta)}}e_{(\gamma)})\label{eq:cA-407}
\end{equation}

决定, 这些联络系数在坐标基下就是通常的Christoffel 符号, 具有表达式

\[
\Gamma^{\mu}_{\rho\nu} = dx^{\mu}\left(\mathbf{D}_{\frac{\partial}{\partial x^{\rho}}}\frac{\partial}{\partial x^{\nu}}\right)
\]

最后

\begin{equation}
\mathbf{D}_{X}Y = \left(X(Y^{\alpha}) + \Gamma^{\alpha}_{\beta\gamma}X^{\beta}Y^{\gamma}\right)e_{(\alpha)}\label{eq:cA-408}
\end{equation}

在坐标标架的特殊情形下, 我们有

\[
\mathbf{D}_{X}Y = \left(X^{\mu}\frac{\partial Y^{\nu}}{\partial x^{\mu}} + \Gamma^{\nu}_{\rho\sigma}X^{\rho}Y^{\sigma}\right)\frac{\partial}{\partial x^{\nu}}
\]

若 $\mathbf{D}\mathbf{g} = 0$, 则联络称为\newterm{levi-civita-connection}{Levi-Civita 联络}{Levi-Civita connection}. 也就是说, 对任意三个向量场 X, Y, Z,

\begin{equation}
Z(\mathbf{g}(X, Y)) = \mathbf{g}(\mathbf{D}_{Z}X, Y) + \mathbf{g}(X, \mathbf{D}_{Z}Y)\label{eq:cA-409}
\end{equation}

微分几何中一个非常简单而基本的结果断言:对任意给定的(伪 Riemann)度量, 存在唯一与其相伴的仿射联络.

\begin{proposition}\label{prop:cA-2.6}
在 $M$ 上存在唯一的联络, 称为 Levi-Civita 联络, 它满足 Dg = 0. 该联络是无挠的, 即

\[
\mathbf{D}_{X}Y - \mathbf{D}_{Y}X = [X, Y].
\]

此外, 相对于坐标系 $x^{\mu}$, 该联络的 Christoffel 符号由标准公式

\[
\Gamma^{\mu}_{\rho\nu} = \frac{1}{2}g^{\mu\tau}\left(\partial_{\rho}g_{\nu\tau} + \partial_{\nu}g_{\tau\rho} - \partial_{\tau}g_{\nu\rho}\right).
\]

给出.
\end{proposition}

\par\noindent\textbf{习题:}自己证明上述命题, 不要看书.

到目前为止, 我们只定义了向量场的协变导数. 我们可以容易地把定义推广到一次形式 $A = A_{\alpha}dx^{\alpha}$:要求对一切向量场 $X$、$Y$ 有

\[
X(A(Y)) = \mathbf{D}_{X}A(Y) + A(D_{X}Y),
\]

给定 $k$ 阶协变张量场 T, 我们按如下法则定义其协变导数 $D_{X}T$:

\[
\mathbf{D}_{X}T(Y_{1}, \dots, Y_{k}) = XT(Y_{1}, \dots, Y_{k}) - \sum_{i=1}^{k} T(Y_{1}, \dots, \mathbf{D}_{X}Y_{i}, \dots, Y_{k})
\]

我们可以把 DT 看作秩为 $k+1$ 的协变张量, 定义为

\[
\mathbf{D}T(X, Y_{1}, \dots, Y_{k}) = \mathbf{D}_{X}T(Y_{1}, \dots, Y_{k}).
\]

给定标架 $e_{\alpha}$, 我们用 $T_{\alpha_{1}\ldots, \alpha_{k};\beta} = \mathbf{D}T(e_{\beta}, e_{\alpha_{1}}, \ldots, e_{\alpha_{k}})$ 表示 DT 相对于该标架的分量. 通过重复的协变微分, 我们可以定义 $D^{2}T, \ldots D^{m}T$. 相对于标架 $e_{\alpha}$, 我们写

\[
\mathbf{D}_{\beta_{1}}\dots\mathbf{D}_{\beta_{m}}T_{\alpha_{1}\dots\alpha_{k}} = T_{\alpha_{1}\dots\alpha_{k};\beta_{1}\dots\beta_{m}} = \mathbf{D}^{m}T(e_{\beta_{1}}\dots, e_{\beta_{m}}, e_{\alpha_{1}}, \dots, e_{\alpha_{k}}).
\]

Levi-Civita 联络无挠这一事实使我们能把协变微分与 Lie 导数联系起来. 于是, 若 $T$ 是 $k$ 阶协变张量, 在坐标基下我们有

\[
(\mathcal{L}_{X}T)_{\sigma_{1}\ldots\sigma_{k}} = X^{\mu}T_{\sigma_{1}\ldots\sigma_{k};\mu} + X^{\mu}_{;\sigma_{1}}T_{\mu\sigma_{2}\ldots\sigma_{k}} + \ldots + X^{\mu}_{;\sigma_{k}}T_{\sigma_{1}\ldots\sigma_{k-1}\mu}.
\]

协变导数还与外导数按如下简单公式相联系. 若 $A$ 是 $k$ 形式, 我们有\footnote{$[\sigma_1\ldots\sigma_k;\mu]$ 表示对所有指标取反对称化(即张量在所有指标置换上的交错和乘以 $1/k!$), 而 "$, \mu$" 表示对 $x^{\mu}$ 的普通偏导数.}$A_{[\sigma_{1}\ldots\sigma_{k};\mu]} = A_{[\sigma_{1}\ldots\sigma_{k}, \mu]}$, 并且

\[
dA = \sum A_{\sigma_{1}\dots\sigma_{k};\mu}\, dx^{\mu} \wedge dx^{\sigma_{1}} \wedge dx^{\sigma_{2}} \wedge \dots \wedge dx^{\sigma_{k}}.
\]

给定光滑曲线 $x : [0, 1] \to M$, 以 $t$ 为参数, 令 $T = \left(\frac{\partial}{\partial t}\right)_{\mathbf{x}}$ 为沿曲线的相应切向量场. 定义在曲线上的向量场 $X$ 称为沿曲线\newterm{parallel-transport}{平行移动}{parallel transport}的, 如果 $D_{T}X = 0$. 若曲线相对于坐标系有参数方程 $x^{\nu} = x^{\nu}(t)$, 则 $T^{\mu} = \frac{dx^{\mu}}{dt}$, 并且分量 $X^{\mu} = X^{\mu}(\mathbf{x}(t))$ 满足常微分方程组

\[
\frac{\mathbf{D}}{dt}X^{\mu} \equiv \frac{dX^{\mu}}{dt} + \Gamma^{\mu}_{\rho\sigma}(\mathbf{x}(t))\frac{dx^{\rho}}{dt}X^{\sigma} = 0.
\]

若在曲线的每一点处 $D_{T}T$ 切于曲线, 即 $D_{T}T = \lambda T$, 则称该曲线为\newterm{geodesic}{测地线}{geodesic}. 此时可以重新参数化曲线, 使得相对于新参数 s, 切向量 $S = \left(\frac{\partial}{\partial s}\right)_{\mathbf{x}}$ 满足 $D_{S}S = 0$. 这样的参数称为"仿射参数". 仿射参数在变换 $s = as' + b$(a、$b$ 为常数)意义下确定. 相对于仿射参数 $s$ 与任意坐标 $x^{\mu}$, 测地线满足方程

\[
\frac{d^{2}x^{\mu}}{ds^{2}} + \Gamma^{\mu}_{\rho\sigma}\frac{dx^{\rho}}{ds}\frac{dx^{\sigma}}{ds} = 0.
\]

以仿射参数参数化的测地曲线简称为测地线. 在 Lorentz 几何中, 类时测地线对应于在由联络系数表示的引力场中自由下落的粒子的\newterm{world-line}{世界线}{world line}. 此时仿射参数 $s$ 称为粒子的\newterm{proper-time}{固有时}{proper time}.

给定点 $p \in M$ 与切空间 $T_{p}M$ 中的向量 X, 令 $\mathbf{x}(t)$ 为从 $p$ 出发、具有"速度" X 的唯一测地线. 我们定义\newterm{exponential-map}{指数映射}{exponential map}:

\[
\exp_{p}: T_{p}\mathbf{M} \to \mathbf{M}.
\]

这个映射不一定对所有 $X \in T_{p}M$ 都有定义. 常微分方程组的存在唯一性定理表明指数映射在 $T_{p}M$ 中原点的一个邻域上有定义. 若对每一点 p, 指数映射对整个 $T_{p}M$ 都有定义, 则称流形 $M$ 是\newterm{geodesically-complete}{测地完备的}{geodesically complete}. 一般而言, 若联络是 $C^{r}$ 联络\footnote{所谓 $C^{r}$ 联络是指:若 $Y$ 是 $C^{r+1}$ 向量场, 则 DY 是 $C^{r}$ 向量场.}, 则存在 $T_{p}M$ 中原点的开邻域 $U_{0}$ 与 $M$ 中点 $p$ 的开邻域 $V_{p}$, 使得映射 $\exp_{p}$ 是 $U_{0}$ 到 $V_{p}$ 上的 $C^{r}$ 微分同胚. 邻域 $V_{p}$ 称为 $p$ 的\newterm{normal-neighbourhood}{法邻域}{normal neighbourhood}.


\section{Riemann 曲率张量、Ricci 张量与 Bianchi 恒等式}

在平坦时空中, 若沿任意闭曲线平行移动一个向量, 则得到出发时的那个向量. 一般而言这不再成立, 因为向量场的二阶协变导数不可交换. 这种不可交换性由\newterm{curvature-tensor}{曲率张量}{curvature tensor}度量:

\begin{equation}
\mathbf{R}(X, Y)Z = \mathbf{D}_{X}(\mathbf{D}_{Y}Z) - \mathbf{D}_{Y}(\mathbf{D}_{X}Z) - \mathbf{D}_{[X, Y]}Z\label{eq:cA-410}
\end{equation}

或者写成相对于任意标架的分量形式:

\begin{equation}
\mathbf{R}^{\alpha}_{\beta\gamma\delta} = \theta^{(\alpha)}\big((\mathbf{D}_{\gamma}\mathbf{D}_{\delta} - \mathbf{D}_{\delta}\mathbf{D}_{\gamma})e_{(\beta)}\big)\label{eq:cA-411}
\end{equation}

相对于坐标系 $x^{\mu}$, 用 $g_{\mu\nu}$ 分量写出, Riemann 分量的表达式为

\begin{equation}
R^{\mu}_{\nu\rho\sigma} = \frac{\partial\Gamma^{\mu}_{\sigma\nu}}{\partial x^{\rho}} - \frac{\partial\Gamma^{\mu}_{\rho\nu}}{\partial x^{\sigma}} + \Gamma^{\mu}_{\rho\tau}\Gamma^{\tau}_{\sigma\nu} - \Gamma^{\mu}_{\sigma\tau}\Gamma^{\tau}_{\rho\nu}\label{eq:cA-412}
\end{equation}

曲率张量的基本性质(由 Riemann 首先证明)断言:若 $\mathbf{R}$ 在点 $p$ 的邻域中恒为零, 则可以找到一族局部坐标, 使得在 $p$ 的邻域中 $g_{\mu\nu} = \eta_{\mu\nu}$\footnote{关于此事实的详细讨论与证明, 参见 Spivak 的《\cite{sp}》第 II 卷.}.

曲率张量关于度量 $g$ 的迹是一个对称张量, 称为\newterm{ricci-tensor}{Ricci 张量}{Ricci tensor},

\[
\mathbf{R}_{\alpha\beta} = \mathbf{g}^{\gamma\delta}\mathbf{R}_{\alpha\gamma\beta\delta}
\]

标量曲率是 Ricci 张量的迹

\[
\mathbf{R} = \mathbf{g}^{\alpha\beta}\mathbf{R}_{\alpha\beta}.
\]

任意时空 $(\mathbf{M}, \mathbf{g})$ 的 Riemann 曲率张量具有以下对称性质,

\begin{equation}
\begin{aligned}
& \mathbf{R}_{\alpha\beta\gamma\delta} = -\mathbf{R}_{\beta\alpha\gamma\delta} = -\mathbf{R}_{\alpha\beta\delta\gamma} = \mathbf{R}_{\gamma\delta\alpha\beta} \\
& \mathbf{R}_{\alpha\beta\gamma\delta} + \mathbf{R}_{\alpha\gamma\delta\beta} + \mathbf{R}_{\alpha\delta\beta\gamma} = 0
\end{aligned}\label{eq:cA-413}
\end{equation}

\eqref{eq:cA-413} 中的第二个恒等式称为第一 Bianchi 恒等式.

它还满足第二 Bianchi 恒等式, 我们在这里称之为 Bianchi 方程, 在一般标架下具有如下形式:

\begin{equation}
\mathbf{D}_{[\epsilon}\mathbf{R}_{\gamma\delta]\alpha\beta} = 0\label{eq:cA-414}
\end{equation}

曲率张量的无迹部分 $\mathbf{C}$ 称为\newterm{weyl-tensor}{Weyl 张量}{Weyl tensor}, 在任意标架下具有如下表达式:

\begin{equation}
\begin{aligned}
\mathbf{C}_{\alpha\beta\gamma\delta} = &\ \mathbf{R}_{\alpha\beta\gamma\delta} - \frac{1}{n-1}\left(\mathbf{g}_{\alpha\gamma}\mathbf{R}_{\beta\delta} + \mathbf{g}_{\beta\delta}\mathbf{R}_{\alpha\gamma} - \mathbf{g}_{\beta\gamma}\mathbf{R}_{\alpha\delta} - \mathbf{g}_{\alpha\delta}\mathbf{R}_{\beta\gamma}\right) \\
& + \frac{1}{n(n-1)}(\mathbf{g}_{\alpha\gamma}\mathbf{g}_{\beta\delta} - \mathbf{g}_{\alpha\delta}\mathbf{g}_{\beta\gamma})\mathbf{R}
\end{aligned}\label{eq:cA-415}
\end{equation}

注意 $C$ 满足 Riemann 张量的所有对称性质:

\begin{equation}
\begin{aligned}
& \mathbf{C}_{\alpha\beta\gamma\delta} = -\mathbf{C}_{\beta\alpha\gamma\delta} = -\mathbf{C}_{\alpha\beta\delta\gamma} = \mathbf{C}_{\gamma\delta\alpha\beta} \\
& \mathbf{C}_{\alpha\beta\gamma\delta} + \mathbf{C}_{\alpha\gamma\delta\beta} + \mathbf{C}_{\alpha\delta\beta\gamma} = 0
\end{aligned}\label{eq:cA-416}
\end{equation}

此外还有 $\mathbf{g}^{\alpha\gamma}\mathbf{C}_{\alpha\beta\gamma\delta} = 0$.

我们说两个度量 $g$ 与 $\hat{g}$ 是\newterm{conformal}{共形的}{conformal}的, 如果存在某个非零可微函数 $\lambda$ 使得 $\hat{g} = \lambda^{2}g$. 于是下面的定理成立(参见 \emph{Hawking--Ellis}, 第 2 章第 2.6 节):

\begin{theorem}\label{thm:cA-3.1}
设 $\hat{\mathbf{g}} = \lambda^2\mathbf{g}$, $\hat{\mathbf{C}}$ 是相对于 $\hat{\mathbf{g}}$ 的 Weyl 张量, $\mathbf{C}$ 是相对于 $\mathbf{g}$ 的 Weyl 张量. 则

\[
\hat{\mathbf{C}}^{\alpha}_{\beta\gamma\delta} = \mathbf{C}^{\alpha}_{\beta\gamma\delta}.
\]
\end{theorem}

因此 $\mathbf{C}$ 是\newterm{conformally-invariant}{共形不变的}{conformally invariant}.

\mtsubsec{等距与共形等距、Killing 与共形 Killing 向量场}

\begin{definition*}
微分同胚 $\Phi : \mathcal{U} \subset \mathbf{M} \to \mathbf{M}$ 称为\newterm{conformal-isometry}{共形等距}{conformal isometry}, 如果在每一点 $p$ 有 $\Phi_{*}\mathbf{g} = \Lambda^{2}\mathbf{g}$, 即

\[
(\Phi^{*}\mathbf{g})(X, Y)|_{p} = \mathbf{g}(\Phi_{*}X, \Phi_{*}Y)|_{\Phi(p)} = \Lambda^{2}\mathbf{g}(X, Y)|_{p}
\]

其中 $\Lambda\neq 0$. 若 $\Lambda = 1$, 则称 $\Phi$ 为 $M$ 的等距.
\end{definition*}

\begin{definition*}
生成单参数等距(相应地, 共形等距)群的向量场 $K$ 称为\newterm{killing-vector-field}{Killing 向量场}{Killing vector field}(相应地, \newterm{conformal-killing-vector-field}{共形 Killing 向量场}{conformal Killing vector field}).
\end{definition*}

设 $K$ 是这样的向量场, $\Phi_{t}$ 是相应的单参数群. 由于 $(\Phi_{t})_{*}$ 是共形等距, 我们推得 $L_{K}g$ 必与度量 $g$ 成比例. 此外, 若 $K$ 是 Killing 向量场, 则 $L_{K}g = 0$.

\begin{definition*}
给定任意向量场 X, 我们用 $^{(X)}\pi$ 表示 $X$ 的变形张量, 由如下公式定义:

\[
{}^{(X)}\pi_{\alpha\beta} = (\mathcal{L}_{X}g)_{\alpha\beta} = \mathbf{D}_{\alpha}X_{\beta} + \mathbf{D}_{\beta}X_{\alpha}.
\]
\end{definition*}

张量 $(X)\pi$ 在精确的意义上度量了由 $X$ 生成的微分同胚与等距或共形等距相差多少. 下面的命题成立(参见 \emph{Hawking--Ellis}, 第 2 章第 2.6 节):

\begin{proposition}\label{prop:cA-3.3}
向量场 $X$ 是 Killing 的当且仅当 $(X)\pi = 0$. 它是共形 Killing 的当且仅当 $(X)\pi$ 与 $g$ 成比例.
\end{proposition}

\begin{remark*}
可以选择局部坐标使得 $X = \frac{\partial}{\partial x^{\mu}}$. 于是立即推得, 相对于这些坐标, 度量 $g$ 与分量 $x^{\mu}$ 无关.
\end{remark*}

\begin{proposition}\label{prop:cA-3.4}
在任意维数为 $n = p + q$ 的伪 Riemann 时空 $M$ 上, 线性无关的 Killing 向量场不可能超过 $\frac{1}{2}(p + q)(p + q + 1)$ 个.
\end{proposition}

\begin{proof}
\hyperref[prop:cA-3.4]{\propref{prop:cA-3.4}} 是以下关系的简单推论, 该关系对任意向量场 $X$ 成立, 可由直接计算并利用 $R$ 的对称性得到:

\begin{equation}
\mathbf{D}_{\beta}\mathbf{D}_{\alpha}X_{\lambda} = \mathbf{R}_{\lambda\alpha\beta\delta}X^{\delta} + {}^{(X)}\Gamma_{\alpha\beta\lambda}\label{eq:cA-417}
\end{equation}

其中

\begin{equation}
{}^{(X)}\Gamma_{\alpha\beta\lambda} = \frac{1}{2}\left(\mathbf{D}_{\beta}\pi_{\alpha\lambda} + \mathbf{D}_{\alpha}\pi_{\beta\lambda} - \mathbf{D}_{\lambda}\pi_{\alpha\beta}\right)\label{eq:cA-418}
\end{equation}

而 $\pi \equiv {}^{(X)}\pi$ 是 $X$ 的变形张量.

若 $X$ 是 Killing 向量场, 方程 \eqref{eq:cA-417} 变为

\begin{equation}
\mathbf{D}_{\beta}(\mathbf{D}_{\alpha}X_{\lambda}) = \mathbf{R}_{\lambda\alpha\beta\delta}X^{\delta}\label{eq:cA-419}
\end{equation}

再结合常微分方程的存在唯一性定理, 这意味着任何 Killing 向量场完全由 $X$ 与 DX 在给定点的 $\frac{1}{2}(n+1)(n+2)$ 个值确定. 事实上, 设 $p$、$q$ 是由具有切向量 $T$ 的曲线 $x(t)$ 相连的两点. 令 $L_{\alpha\beta} \equiv D_{\alpha}X_{\beta}$, 注意沿 $x(t)$, X、$L$ 满足微分方程组

\[
\frac{\mathbf{D}}{dt}X = T \cdot L, \qquad \frac{\mathbf{D}}{dt}L = \mathbf{R}(\cdot, \cdot, X, T)
\]

因此 $X$、$L$ 沿曲线的值由它们在 $p$ 处的值唯一确定.
\end{proof}

具有最多 Killing 向量场个数的 $n$ 维 Riemann 流形是 Euclidean 空间 $R^{n}$. 类似地, Minkowski 时空 $R^{n+1}$ 是具有最多 Killing 向量场个数的 Lorentz 流形.

\mtsubsec{体积形式、Hodge 对偶与散度定理}

设 $(\mathbf{M}, \mathbf{g})$ 是维数为 $n$、符号差为 $s$ 的可定向伪 Riemann 流形 M. 设 $e_{(\alpha)}$ 是 $M$ 上任意正定向标架, $e^{(\alpha)}$ 是相应的对偶标架. 我们可以把它与 $n$ 形式

\[
e^{(1)} \wedge e^{(2)} \wedge \ldots \wedge e^{(n)}
\]

相联系. 若 $e_{(\alpha)}' = M_{\alpha}^{\beta}e_{(\beta)}$ 表示另一个基, 而 $e'^{(\alpha)} = (M^{-1})_{\beta}^{\alpha}e^{(\beta)}$ 是相应的一次形式对偶基, 则

\[
e'^{(1)} \wedge e'^{(2)} \wedge \dots \wedge e'^{(n)} = \det(M^{-1})e^{(1)} \wedge e^{(2)} \wedge \dots \wedge e^{(n)}
\]

另一方面, 若用 $\mathbf{g}$ 与 $\mathbf{g}'$ 表示矩阵 $\mathbf{g}(e_{(\alpha)}, e_{(\beta)})$、$\mathbf{g}(e'_{(\alpha)}, e'_{(\beta)})$ 的行列式, 则 $\mathbf{g}' = (\det M)^2\mathbf{g}$. 因此, n 形式

\begin{equation}
\varepsilon = |\mathbf{g}|^{1/2}e^{(1)} \wedge \dots \wedge e^{(n)}\label{eq:cA-420}
\end{equation}

与保持定向的基变换无关. 于是, 若 $\mathbf{M}$ 可定向, 我们可以定义一个唯一的 $n$ 形式, 称为 $\mathbf{M}$ 的\newterm{volume-form}{体积形式}{volume form}. 相对于正定向标准正交标架 $e_{(\alpha)}$, 当有两个指标相同时 $\varepsilon_{\alpha_1\dots\alpha_n} = 0$, 否则等于置换 $\alpha_{1}\alpha_{2}\ldots\alpha_{n}$ 的符号. $\varepsilon$ 的以下乘法性质在计算中非常有用,

\begin{equation}
\varepsilon^{\alpha_{1}\dots\alpha_{n}}\varepsilon_{\beta_{1}\dots\beta_{n}} = (-1)^{s}\det(\delta^{\alpha_{i}}_{\beta_{j}})_{i, j = 1, \dots, n}\label{eq:cA-421}
\end{equation}

其中 $s$ 是度量的符号差. 另外, 通过重复缩并,

\[
\begin{aligned}
\varepsilon^{\alpha_{1}\alpha_{2}\ldots\alpha_{n}}\varepsilon_{\alpha_{1}\beta_{2}\ldots\beta_{n}} &= (-1)^{s}\det(\delta^{\alpha_{i}}_{\beta_{j}})_{i, j = 2, \ldots, n} \\
\varepsilon^{\alpha_{1}\alpha_{2}\alpha_{3}\ldots\alpha_{n}}\varepsilon_{\alpha_{1}\alpha_{2}\beta_{3}\ldots\beta_{n}} &= 2(-1)^{s}\det(\delta^{\alpha_{i}}_{\beta_{j}})_{i, j = 3, \ldots, n} \\
\ldots &= \ldots \\
\varepsilon^{\alpha_{1}\ldots\alpha_{n}}\varepsilon_{\alpha_{1}\ldots\beta_{n}} &= n!(-1)^{s}
\end{aligned}
\]

注意 $\varepsilon$ 的协变导数消失. 事实上, 这可由对公式 $\varepsilon^{\alpha_{1}\cdots\alpha_{n}}\varepsilon_{\alpha_{1}\cdots\alpha_{n}} = n!(-1)^{s}$ 求导直接推出, 由此得 $\varepsilon^{\alpha_{1}\cdots\alpha_{n}}\varepsilon_{\alpha_{1}\cdots\alpha_{n};\mu} = 0$, 于是利用 \eqref{eq:cA-421} 有 $D_{\mu}\varepsilon_{\alpha_{1}\cdots\alpha_{n}} = \varepsilon_{\alpha_{1}\cdots\alpha_{n};\mu} = 0$.

Hodge 对偶是自然定义的从 $\Omega^{k}(\mathbf{M})$ 到 $\Omega^{n-k}(\mathbf{M})$ 的算子. 在固定的正定向标架中, 若用 $A_{\alpha_{0}\ldots\alpha_{k-1}}$ 表示 $A \in \Omega^{k}(\mathbf{M})$ 的系数, 则我们有

\[
\ast A_{\alpha_{1}\dots\alpha_{n-k}} = \frac{1}{k!}\varepsilon_{\alpha_{1}\dots\alpha_{n-k}}{}^{\beta_{1}\dots\beta_{k}}A_{\beta_{1}\dots\beta_{k}}
\]

容易验证, 对 $k$ 形式 $A$ 有\footnote{特别地, 若 $\mathbf{g}$ 是维数为 4 的 Lorentz 流形, 则 $*(*A) = (-1)^{k+1}A$.},

\begin{equation}
*(*A) = (-1)^{s}(-1)^{k(n-k)}A\label{eq:cA-422}
\end{equation}

我们按如下公式定义散度算子 $\delta : \Omega^{k}(\mathbf{M}) \to \Omega^{k-1}(\mathbf{M})$:

\begin{equation}
\delta A = *d*A\label{eq:cA-423}
\end{equation}

给定 $\mathbf{M}$ 上连续、紧支集的函数 f, 我们按如下公式定义其积分:

\[
\int_{\mathbf{M}}fdv_{\mathbf{g}} := \int f\varepsilon
\]

若 $f$ 的支集在(正定向)坐标卡内, 则

\[
\int_{\mathbf{M}}fdv_{\mathbf{g}} = \int f(x^{1}, \ldots, x^{n})|\mathbf{g}|^{1/2}dx^{1} \wedge \ldots \wedge dx^{n}.
\]

于是我们可以引入 $M$ 上函数的标量积,

\begin{equation}
\langle f_{1}, f_{2}\rangle = \int_{\mathbf{M}}f_{1}\overline{f_{2}}\, dv_{\mathbf{g}}\label{eq:cA-424}
\end{equation}

给定两个(紧支集的)形式 $A \in \Omega^k(\mathbf{M})$、$B \in \Omega^{n-k}(\mathbf{M})$, 我们定义

\[
\int_{\mathbf{M}}A \wedge B
\]

其中积分理解为 $M$ 上 $n$ 形式的通常积分. 另外, 给定 $A, B \in \Omega^{k}(\mathbf{M})$, 相对于固定的标架 $e_{1}, e_{2}$ 我们定义

\[
\mathbf{g}(A, B) = \frac{1}{k!}A^{\alpha_{1}\dots\alpha_{k}}B_{\alpha_{1}\dots\alpha_{k}}
\]

并且

\begin{equation}
\langle A, B\rangle := \int_{\mathbf{M}}\mathbf{g}(A, B)\label{eq:cA-425}
\end{equation}

利用这些定义, 我们容易验证, 对 $A \in \Omega^k(\mathbf{M})$、$B \in \Omega^{n-k}(\mathbf{M})$ 有

\[
\langle {}^{\ast}A, B\rangle = \int_{\mathbf{M}}A \wedge B
\]

还注意到散度算子 $\delta$ 在如下意义下是 $d$ 的对偶(对 $A \in \Omega^{k}(\mathbf{M})$、$B \in \Omega^{k-1}(\mathbf{M})$):

\[
\langle \delta A, B\rangle = \langle A, dB\rangle
\]

若 $X$ 是 $M$ 上的(紧支集)向量场, 我们定义其 Hodge 对偶 $*X$ 为与其相伴的一次形式的 Hodge 对偶, 即

\[
(\ast X)_{\alpha_{1}\ldots\alpha_{n-1}} = \varepsilon_{\alpha_{1}\ldots\alpha_{n-1}\mu}X^{\mu}
\]

容易验证, 相对于任意正定向坐标标架, 即 $*X = \frac{1}{(k-1)!}\varepsilon_{\alpha_{1}\ldots\alpha_{n-1}\mu}X^{\mu}dx^{\alpha_{1}} \wedge \ldots \wedge dx^{\alpha_{n-1}}$, 我们有

\[
d(*X) = (-1)^{n-1}(\mathbf{D}^{\alpha}X_{\alpha})dv_{\mathbf{g}}
\]

从而

\[
\delta X =: \mathbf{Div}X = \mathbf{D}^{\alpha}X_{\alpha}
\]

将这一事实与 Stokes 定理相结合, 我们推得如下结论.

\begin{proposition}\label{prop:cA-3.6}
设 $X$ 是 $M$ 上的向量场, $\ast X$ 是其 Hodge 对偶. 设 $U$ 是 $M$ 中具有光滑边界、且按 $M$ 诱导的正定向定向的区域. 则

\[
\int_{\mathcal{U}}\mathbf{Div}X\, dv_{\mathbf{g}} = (-1)^{n-1}\int_{\partial\mathcal{U}}{}^{\ast}X.
\]
\end{proposition}

我们将该命题应用于维数为 $n+1$ 的 Lorentz 流形. 给定具有诱导 Riemann 度量 $g$ 的类空超曲面 $\Sigma$, 我们用 $dv_{g}$ 表示 $g$ 相对于诱导正定向的体积元. 现在我们把\newterm{divergence-theorem}{散度定理}{divergence theorem}应用于 $M$ 中由两个类空超曲面 $\Sigma_{0}, \Sigma_{1}$ 围成的\newterm{lens-shaped-domain}{透镜形区域}{lens shaped domain} U.

\begin{corollary}\label{cor:cA-3.7}
设 $U$ 是由两个类空超曲面 $\Sigma_{0}, \Sigma_{1}$ 围成的透镜形区域. 设 $T$ 表示边界两侧的未来定向单位法向量, X 是 $M$ 上的任意光滑向量场. 则

\[
\int_{\Sigma_{1}}(X, T)dv_{g} - \int_{\Sigma_{0}}(X, T)dv_{g} = -\int_{\mathcal{U}}\mathbf{Div}X\, dv_{\mathbf{g}}.
\]
\end{corollary}

\begin{proof}
首先考虑边界的 $\Sigma_{0}$ 部分, 并在 $\Sigma_{0}$ 的每一点 $p$ 附近局部定义坐标系 $x^{0}, x^{1}, \ldots, x^{n}$, 使得

\begin{enumerate}
\renewcommand{\labelenumi}{(\arabic{enumi})}
\item 在 $\Sigma_0$ 上 $x^0 = 0$, 在 $\mathcal{U}$ 中 $x^0 < 0$.
\item $x^{1}, x^{2}, \ldots, x^{n}$ 在 $\Sigma_{0}$ 中是正定向的.
\item 在 $p$ 附近 $\Sigma_0$ 上 $\mathbf{g}_{00} = -1$、$\mathbf{g}_{0i} = 0$.
\end{enumerate}

相对于这些坐标, 我们有 $dv_{g} = \varepsilon_{12\ldots n}dx^{1} \wedge \ldots \wedge dx^{n}$, 其中 $\varepsilon_{12\ldots n} = \varepsilon_{01\ldots n}$. 因此在 $\Sigma_{0}$ 上,

\[
(\ast X)_{i_{1}\dots i_{n}} = \varepsilon_{i_{1}\dots i_{n}\mu}X^{\mu} = \varepsilon_{i_{1}\dots i_{n}0}X^{0} = (-1)^{n}X^{0}dv_{g}
\]
\end{proof}

\mtsubsec{Laplace--Beltrami 算子}

伪 Riemann 流形 $M$ 上的标量Laplace--Beltrami 算子定义为

\begin{equation}
\Delta_{\mathbf{M}}u(x) = g^{\mu\nu}\mathbf{D}_{\mu}\mathbf{D}_{\nu}u\label{eq:cA-426}
\end{equation}

其中 $u$ 是 $M$ 上的标量函数. 或者, 在局部坐标下

\begin{equation}
\Delta_{\mathbf{M}}u(x) = \frac{1}{\sqrt{|g(x)|}}\partial_{\mu}(g^{\mu\nu}\sqrt{|g(x)|}\partial_{\nu})u(x)\label{eq:cA-427}
\end{equation}

在 Lorentz 流形的特殊情形下, Laplace--Beltrami 算子称为D'Alembert 算子, 记作 $\Box_{\mathbf{M}}$. 在任意伪 Riemann 流形上, $\Delta_{\mathbf{M}}$ 关于下面标量函数 $u$、$v$ 的标量积是对称的:

\[
(u, v)_{\mathbf{M}} = \int u(x)v(x)dv_{\mathbf{M}}
\]

事实上, 以下恒等式通过对任意两个光滑、紧支集\footnote{若流形 $\mathbf{M}$ 是紧的, 则此条件自动满足.}的函数 $u$、$v$ 分部积分容易建立:

\begin{equation}
(-\Delta u, v)_{\mathbf{M}} = \int_{\mathbf{M}}\nabla u \cdot \nabla v\, dv_{\mathbf{M}} = (u, -\Delta v)_{\mathbf{M}}\label{eq:cA-428}
\end{equation}

其中 $\nabla u\cdot\nabla v = g^{ij}\partial_{i}u\partial_{j}v$. 在 u = v 的特殊情形下, 我们推得 $(-\Delta u, v)_{\mathbf{M}} = \int_{\mathbf{M}}|\nabla u|^{2}$, 其中 $|\nabla u|^{2} = \nabla u\cdot\nabla u$. 于是, 对函数 $u \in \mathcal{C}_{0}^{\infty}(\mathbf{M})$, $-\Delta = -\Delta_{M}$ 是对称的. 若流形 $M$ 是 Riemann 的, 则它是正定的. 对 Lorentz 流形则不然:$\Box_{M}$ 不是正定的.


\section{Minkowski 空间}

\mtsubsec{基本定义}

我们记为 $R^{n+1}$ 的 $n+1$ 维 Minkowski 空间由流形 $R^{n+1}$ 连同 Lorentz 度量 $m$ 与一个特殊的坐标系 $x^{\alpha}$, $\alpha = 0, 1, \ldots, n$(称为惯性坐标系)组成, 相对于该坐标系度量具有对角形式 $\mathbf{m}_{\alpha\beta} = \text{diag}(-1, 1, \ldots, 1)$. 两个惯性坐标系之间由平移或 Lorentz 变换相联系. 我们使用关于 $m$ 及其逆 $m^{-1} = m$ 的指标升降的标准几何约定, 以及重复指标上的通常求和约定. 坐标向量场 $\frac{\partial}{\partial x^{\alpha}}$ 记作 $\partial_{\alpha}$, 任意向量场记作 $X = X^{\alpha}\partial_{\alpha}$, 其中 $X^{\alpha} = X^{\alpha}(x^{0}, \ldots, x^{n})$. 注意, 相对于 $m$ 下降指标, 得 $X_{0} = -X^{0}$, 且对一切 $i = 1, \ldots, n$ 有 $X_{i} = X^{i}$. 我们用 $D$ 表示 $R^{n+1}$ 的平坦协变导数, 即对任意一次形式 $w = \omega_{\alpha}dx^{\alpha}$ 有 $D_{\alpha}\omega_{\beta} = \partial_{\alpha}\omega_{\beta}$. 我们还把时空坐标 $x^{\alpha}$ 分为时间分量 $x^{0} = t$ 与空间分量 $x = x^{i}, \ldots, x^{n}$. 注意 $t_{0} = -t$, 且对 $i = 1, \ldots, n$ 有 $x^{i} = x_{i}$.

向量 $X$ 称为类时、类光或类空的, 分别取决于 $\mathbf{m}(X, X)$ 是 $<0$、$=0$ 还是 $>0$. 相应地, 光滑曲线 $x^{\alpha}(s)$ 称为类时、类光或类空的, 如果其切向量 $\frac{dx^{\alpha}}{ds}$ 在其每一点处是类时、类光或类空的. 因果曲线可以是类时的或类光的. 类似地, 超曲面 $u(x^{0}, \ldots, x^{n})=0$ 称为类空、类光或类时的, 如果其法向量 $N^{\alpha} = -\mathbf{m}^{\alpha\beta}\partial_{\beta}u$ 分别是类时、类光或类空的. m 在类空超曲面上诱导的度量必然正定, 即是 Riemann 的. 函数 $\mathbf{t}(x^{0}, x^{1}, \ldots, x^{n})$ 称为\newterm{time-function}{时间函数}{time function}, 如果其水平超曲面 t = t 是类空的. 在零超曲面上, 诱导度量沿法向方向是退化的, 即 $\mathbf{m}(N, N) = 0$. 特别地, 水平面 u = u 为零(即类光)的函数 $\mathbf{u} = \mathbf{u}(x^{0}, \ldots x^{n})$ 必须满足\newterm{eikonal-equation}{程函方程}{eikonal equation}

\begin{equation}
\mathbf{m}^{\alpha\beta}\partial_{\alpha}\mathbf{u}\partial_{\beta}\mathbf{u} = 0\label{eq:cA-429}
\end{equation}

方程 \eqref{eq:cA-429} 也可以写成 $D_{N}N = 0$ 的形式. 我们称 $N$ 为 $u$ 的水平超曲面的\newterm{null-geodesic-generator}{零测地线生成元}{null geodesic generator}.

因果曲线在其每一点处可以是类时的或类光的. $R^{n+1}$ 的典范时间定向由向量场 $T_{0} = \partial_{0}$ 给出. 类时向量 $X$ 称为未来定向的, 如果 $\mathbf{m}(X, T_{0}) < 0$;称为过去定向的, 如果 $\mathbf{m}(X, T_{0}) > 0$. 集合 $S$ 的\newterm{causal-future}{因果未来}{causal future} $J^{+}(S)$ 由 $R^{n+1}$ 中所有能通过未来定向因果曲线与 $S$ 相连的点组成. \newterm{causal-past}{因果过去}{causal past} $\mathcal{J}^{-}(S)$ 以同样方式定义. 于是, 对点 $p = (t, x)$, $\mathcal{J}^{+}(p) = \{(t \geq t_{0}, x)\, /\, |x - x_{0}| \leq t - t_{0}\}$. 给定光滑区域 D, 其未来集合 $\mathcal{J}^{+}(D)$ 一般可能具有不光滑的边界, 这是由于\newterm{caustic}{焦散}{caustic}引起的.

我们考虑\newterm{conservative-domain}{保守区域}{conservative domain} $\mathcal{J}^{+}(D_{1})\cap\mathcal{J}^{-}(D_{2})$, 其中 $D_{1}\subset\Sigma_{1}$、$D_{2}\subset\Sigma_{2}$ 是类空超曲面. 若 $D_{1}, D_{2}$ 都是正则的, 且其非类空边界 $\mathcal{N}_1\subset\partial(\mathcal{J}^+(D_1))\setminus D_1$ 与 $\mathcal{N}_2\subset\partial(\mathcal{J}^{-}(D_2))\setminus D_2$ 都是光滑的, 则称该区域为正则的. 在特殊情形 $D_{1} = \Sigma_{1}$、$D = D_{2}\subset\Sigma_{2}$ 下, 我们得到 $\mathcal{J}^{+}(\Sigma_{1})\cap\mathcal{J}^{-}(D)$, 称为 $D$ 关于 $\Sigma_{1}$ 的\newterm{domain-of-dependence}{依赖域}{domain of dependence}, 它由 $D\subset\Sigma_{2}$ 的因果过去中位于 $\Sigma_{1}$ 未来的所有点组成. 类似地, $\mathcal{J}^{+}(D)\cap\mathcal{J}^{-}(\Sigma_{2})$(其中 $D\subset\Sigma_{1}$)称为 $D$ 关于 $\Sigma_{2}$ 的影响依赖域. 特别有用的例子由时间函数 $\mathbf{t}$ 给出, 取 $\Sigma_1 = \{(t, x)\, /\, \mathbf{t}(t, x) = t_1\}$、$\Sigma_{2} = \{(t, x)\, /\, \mathbf{t}(t, x) = t_{1}\}$ 为两个不相交的水平超曲面, $\Sigma_{2}$ 位于 $\Sigma_{1}$ 的未来.

一对类光向量场 $L, \underline{L}$ 构成零对, 如果 $\mathbf{m}(L, \underline{L}) = -2$. 零对 $e_n = L$、$e_{n+1} = \underline{L}$ 连同满足 $\mathbf{m}(L, e_a) = \mathbf{m}(\underline{L}, e_a) = 0$ 且对一切 $a, b = 1, \ldots, n-1$ 有 $\mathbf{m}(e_a, e_b) = \delta_{ab}$ 的向量场 $e_1, \ldots, e_{n-1}$, 称为零标架. 零对

\begin{equation}
\mathbf{L} = \partial_{t} + \partial_{r}, \qquad \underline{\mathbf{L}} = \partial_{t} - \partial_{r}\label{eq:cA-430}
\end{equation}

(其中 $r = |x|$, $\partial_{r} = x^{i}/r\, \partial_{i}$)称为典范的. 类似地, 满足 $e_{n} = L$、$e_{n+1} = \underline{L}$ 的零标架 $e_{1}, \ldots, e_{n+1}$ 称为典范零标架. 此时 $e_{1}, \ldots, e_{n-1}$ 在任意点构成经过该点、$t$ 与 $r$ 恒定的球面 $S_{t, r}$ 的标准正交基. 还注意到 $L$ 是与 u = t - r 相关联的零测地线生成元, 而 $\underline{L}$ 是 $\underline{u} = t + r$ 的零测地线.

\mtsubsec{共形 Killing 向量场}

设 $x^{\mu}$ 是 Minkowski 空间 $R^{n+1}$ 的惯性坐标系. 以下是 $R^{n+1}$ 的全部等距与共形等距.

\begin{enumerate}
\renewcommand{\labelenumi}{\arabic{enumi}.}
\item \textbf{平移:}对任意给定的向量 $a = (a^{0}, a^{1}, \ldots, a^{n}) \in \mathbb{R}^{n+1}$,
\end{enumerate}

\[
x^{\mu} \rightarrow x^{\mu} + a^{\mu}
\]

\begin{enumerate}
\renewcommand{\labelenumi}{\arabic{enumi}.}
\setcounter{enumi}{1}
\item \textbf{Lorentz 旋转:}给定任意 $\Lambda = \Lambda_{\sigma}^{\rho} \in \mathbf{O}(1, n)$,
\end{enumerate}

\[
x^{\mu} \to \Lambda_{\nu}^{\mu}x^{\nu}
\]

\begin{enumerate}
\renewcommand{\labelenumi}{\arabic{enumi}.}
\setcounter{enumi}{2}
\item \textbf{伸缩:}给定任意非零实数 $\lambda$,
\end{enumerate}

\[
x^{\mu} \rightarrow \lambda x^{\mu}
\]

\begin{enumerate}
\renewcommand{\labelenumi}{\arabic{enumi}.}
\setcounter{enumi}{3}
\item \textbf{反演:}考虑变换 $x^{\mu} \rightarrow I(x^{\mu})$, 其中
\end{enumerate}

\[
I(x^{\mu}) = \frac{x^{\mu}}{(x, x)}
\]

对一切满足 $(x, x) \neq 0$ 的点 $x \in \mathbb{R}^{n+1}$ 有定义.

前两类变换是 $R^{n+1}$ 的等距, 由它们生成的群称为 Poincaré 群. 后两类变换是共形等距. 由上述所有变换生成的群称为共形群. 事实上, Liouville 定理(其无穷小版本将在后面证明)断言它就是 $R^{n+1}$ 的全部共形等距群.

下面我们列出生成上述变换的 Killing 与共形 Killing 向量场.

\begin{enumerate}
\renewcommand{\labelenumi}{\roman{enumi}.}
\item $x^{\mu}$ 方向平移的生成元, $\mu = 0, 1, \ldots, n$:
\end{enumerate}

\[
\mathbf{T}_{\mu} = \frac{\partial}{\partial x^{\mu}}
\]

\begin{enumerate}
\renewcommand{\labelenumi}{\roman{enumi}.}
\setcounter{enumi}{1}
\item $(\mu, \nu)$ 平面内 Lorentz 旋转的生成元:
\end{enumerate}

\[
\mathbf{L}_{\mu\nu} = x_{\mu}\partial_{\nu} - x_{\nu}\partial_{\mu}
\]

\begin{enumerate}
\renewcommand{\labelenumi}{\roman{enumi}.}
\setcounter{enumi}{2}
\item 伸缩变换的生成元:
\end{enumerate}

\[
\mathbf{S} = x^{\mu}\partial_{\mu}
\]

\begin{enumerate}
\renewcommand{\labelenumi}{\roman{enumi}.}
\setcounter{enumi}{3}
\item 逆平移\footnote{注意, 向量场 $K_{\mu}$ 可通过对向量场 $T_{\mu}$ 作用 $I_{*}$ 得到.}的生成元:
\end{enumerate}

\[
\mathbf{K}_{\mu} = 2x_{\mu}x^{\rho}\frac{\partial}{\partial x^{\rho}} - (x^{\rho}x_{\rho})\frac{\partial}{\partial x^{\mu}}
\]

下面我们还列出这些向量场之间的换位关系,

\begin{equation}
\begin{aligned}
&[\mathbf{L}_{\alpha\beta}, \mathbf{L}_{\gamma\delta}] = \eta_{\alpha\gamma}\mathbf{L}_{\beta\delta} - \eta_{\beta\gamma}\mathbf{L}_{\alpha\delta} + \eta_{\beta\delta}\mathbf{L}_{\alpha\gamma} - \eta_{\alpha\delta}\mathbf{L}_{\beta\gamma} \\
&[\mathbf{L}_{\alpha\beta}, \mathbf{T}_{\gamma}] = \eta_{\alpha\gamma}\mathbf{T}_{\beta} - \eta_{\beta\gamma}\mathbf{T}_{\alpha} \\
&[\mathbf{T}_{\alpha}, \mathbf{T}_{\beta}] = 0 \\
&[\mathbf{T}_{\alpha}, \mathbf{S}] = \mathbf{T}_{\alpha} \\
&[\mathbf{T}_{\alpha}, \mathbf{K}_{\beta}] = 2(\eta_{\alpha\beta}\mathbf{S} + \mathbf{L}_{\alpha\beta}) \\
&[\mathbf{L}_{\alpha\beta}, \mathbf{S}] = [\mathbf{K}_{\alpha}, \mathbf{K}_{\beta}] = 0 \\
&[\mathbf{L}_{\alpha\beta}, \mathbf{K}_{\gamma}] = \eta_{\alpha\gamma}\mathbf{K}_{\beta} - \eta_{\beta\gamma}\mathbf{K}_{\alpha}
\end{aligned}\label{eq:cA-431}
\end{equation}

记 $\mathcal{P}(1, n)$ 为由向量场 $T_{\alpha}, L_{\beta\gamma}$ 生成的 李代数, $\underline{\mathcal{K}}(1, n)$ 为由所有向量场 $T_{\alpha}, L_{\beta\gamma}, S, K_{\delta}$ 生成的 李代数, 我们陈述 Liouville 定理的如下版本:

\begin{theorem}\label{thm:cA-4.3}
以下陈述成立.

\begin{enumerate}
\renewcommand{\labelenumi}{\arabic{enumi})}
\item $\mathcal{P}(1, n)$ 是 $\mathbb{R}^{n+1}$ 中所有 Killing 向量场的 李代数.
\item 若 $n > 1$, $\underline{\mathcal{K}}(1, n)$ 是 $\mathbb{R}^{n+1}$ 中所有共形 Killing 向量场的 李代数.
\item 若 $n = 1$, $\mathbb{R}^{1+1}$ 中所有共形 Killing 向量场的集合由如下表达式给出
\end{enumerate}

\[
f(x^{0} + x^{1})(\partial_{0} + \partial_{1}) + g(x^{0} - x^{1})(\partial_{0} - \partial_{1})
\]

其中 $f$、$g$ 是单变量的任意光滑函数.
\end{theorem}

\begin{proof}
定理第 1 部分的证明作为\hyperref[prop:cA-3.4]{\propref{prop:cA-3.4}} 的特殊情形立即得出. 由 \eqref{eq:cA-417}, 因 R = 0 且 $X$ 是 Killing 的, 我们有

\[
D_{\mu}D_{\nu}X_{\lambda} = 0.
\]

因此, 存在常数 $a_{\mu\nu}, b_{\mu}$ 使得 $X^{\mu} = a_{\mu\nu}x^{\nu} + b_{\mu}$. 由于 $X$ 是 Killing 的, $D_{\mu}X_{\nu} = -D_{\nu}X_{\mu}$, 这意味着 $a_{\mu\nu} = -a_{\nu\mu}$. 从而 $X$ 可以写成向量场 $T_{\alpha}, L_{\beta\gamma}$ 的实系数线性组合.

现在设 $X$ 是共形 Killing 向量场. 存在函数 $\Omega$ 使得

\begin{equation}
{}^{(X)}\pi_{\rho\sigma} = \Omega\eta_{\rho\sigma}\label{eq:cA-432}
\end{equation}

由 \eqref{eq:cA-417} 与 \eqref{eq:cA-418} 推得

\begin{equation}
D_{\mu}D_{\nu}X_{\lambda} = \frac{1}{2}\left(\Omega_{, \mu}\eta_{\nu\lambda} + \Omega_{, \nu}\eta_{\mu\lambda} - \Omega_{, \lambda}\eta_{\nu\mu}\right)\label{eq:cA-433}
\end{equation}

对 \eqref{eq:cA-433} 两边关于 $\mu, \nu$ 取迹, 我们推得

\begin{equation}
\begin{aligned}
\Box X_{\lambda} &= -\frac{n-1}{2}\Omega_{, \lambda} \\
D^{\mu}X_{\mu} &= \frac{n+1}{2}\Omega
\end{aligned}\label{eq:cA-434}
\end{equation}

再对第一个方程作用 $D^{\lambda}$, 对第二个方程作用 $\Box$, 相减得

\begin{equation}
\Box\Omega = 0\label{eq:cA-435}
\end{equation}

对 \eqref{eq:cA-434} 的第一个方程作用 $D_{\mu}$ 并利用 \eqref{eq:cA-435}, 得

\begin{equation}
\begin{aligned}
(n-1)D_{\mu}D_{\lambda}\Omega &= \frac{n-1}{2}(D_{\mu}D_{\lambda}\Omega + D_{\lambda}D_{\mu}\Omega) = -\Box(D_{\mu}X_{\lambda} + D_{\lambda}X_{\mu}) \\
&= -(\Box\Omega)\eta_{\mu\lambda} = 0
\end{aligned}\label{eq:cA-436}
\end{equation}

因此对 $n \neq 1$, 有 $D_{\mu}D_{\lambda}\Omega = 0$. 这意味着 $\Omega$ 必为 $x^{\mu}$ 的线性函数. 于是我们可以找到常系数的线性组合 $cS + d^{\alpha}K_{\alpha}$, 使得 $X - (cS + d^{\alpha}K_{\alpha})$ 的变形张量为零. 确实如此, 因为 $^{(S)}\pi = 2\eta$ 且 $^{(K_{\mu})}\pi = 4x_{\mu}\eta$. 因此 $X - (cS + d^{\alpha}K_{\alpha})$ 是 Killing 的, 由定理的第一部分即得结论.

第 3 部分可以通过求解 \eqref{eq:cA-432} 容易导出. 事实上, 设 $X = a\partial_0 + b\partial_1$, 我们得到 $2D_0X_0 = -\Omega$、$2D_1X_1 = \Omega$、$D_0X_1 + D_1X_0 = 0$. 因此 $a$、$b$ 满足方程组

\[
\frac{\partial a}{\partial x^{0}} = \frac{\partial b}{\partial x^{1}}, \qquad \frac{\partial b}{\partial x^{0}} = \frac{\partial a}{\partial x^{1}}.
\]

于是一次形式 $adx^0 + bdx^1$ 是恰当的, $adx^0 + bdx^1 = d\phi$, 并且 $\frac{\partial^2a}{\partial x^{0^2}} = \frac{\partial^2b}{\partial x^{1^2}}$, 即 $\Box\phi = 0$. 总之

\[
X = \frac{1}{2}\left(\frac{\partial\phi}{\partial x^{0}} + \frac{\partial\phi}{\partial x^{1}}\right)(\partial_{0} + \partial_{1}) + \frac{1}{2}\left(\frac{\partial\phi}{\partial x^{0}} - \frac{\partial\phi}{\partial x^{1}}\right)(\partial_{0} - \partial_{1})
\]

这就证明了结论.
\end{proof}


\begin{remark*}
相对于典范零对表示, 有

\begin{equation}
\mathbf{T}_{0} = 2^{-1}(\mathbf{L} + \underline{\mathbf{L}}), \qquad \mathbf{S} = 2^{-1}(\underline{u}\mathbf{L} + u\underline{\mathbf{L}}), \qquad \mathbf{K}_{0} = 2^{-1}(\underline{u}^{2}\mathbf{L} + u^{2}\underline{\mathbf{L}}).\label{eq:cA-437}
\end{equation}

$T_{0} = \partial_{t}$ 与 $\mathbf{K}_{0} = (t^{2} + |x|^{2})\partial_{t} + 2tx^{i}\partial_{i}$ 都是因果的. 这使得它们在推导能量估计时很重要. 注意 $S$ 仅在 $\mathcal{J}^{+}(0) \cup \mathcal{J}^{-}(0)$ 中是因果的.
\end{remark*}

\mtsubsec{零超曲面}

零超曲面特别重要, 因为它们对应于 Minkowski 空间中波动方程或 Maxwell 方程解的传播前沿\footnote{或更一般地在 Lorentz 时空上.}. 描述零超曲面几何的最简单方法是从余维数为一的超曲面 $S_{0} \subset \Sigma_{0}$ 出发, 其中 $\Sigma_{0}$ 是 $M^{n+1}$ 的固定类空超曲面. 在 $S_{0}$ 的每一点 $p \in S_{0}$, 恰有两个与切空间 $T_{p}(S_{0})$ 正交的类光方向. 设 $L$ 表示与 $S_{0}$ 正交的光滑类光向量场, 并考虑由 $L$ 的积分曲线生成的零测地线汇\footnote{在 Minkowski 空间中, 这些实际上是直线.}. 只要这些零测地线不相交, 该线汇就构成光滑的零超曲面 N. 我们还可以通过平行移动把 $L$ 延拓到 $N$ 的所有点. 显然 $D_{L}L = 0$、$\mathbf{m}(L, L) = 0$, 此外对一切切于 $N$ 的向量 $X$ 有 $\mathbf{m}(L, X) = 0$. 还注意到 $L$ 在乘以一个只依赖于 $S_{0}$ 的共形因子的意义下唯一确定. 对一切切于 $N$ 的向量场 $X$、Y, 定义

\begin{equation}
\gamma(X, Y) = \mathbf{m}(X, Y), \qquad \chi(X, Y) = \mathbf{m}(D_{X}L, Y)\label{eq:cA-438}
\end{equation}

它们都是对称张量, 分别称为 $N$ 的\newterm{null-fundamental-form}{零基本形式}{null fundamental form}第一与第二零基本形式. 注意 $\chi$ 在乘以与 $L$ 相联系的同一个共形因子的意义下唯一确定. 显然, 对一切切于 $N$ 的 $X$ 有 $\gamma(L, X) = \chi(L, X) = 0$, 因此它们在固定的 $p\in N$ 处只依赖于与 $L_{p}$ 横截的一个固定超平面. 定义 s(称为仿射参数), 由条件 $L(s)=1$ 及在 $S_{0}$ 上 s=0 确定. 其水平面定义了 $N$ 的\newterm{geodesic-foliation}{测地叶状结构}{geodesic foliation}. 给定 $S_{0}$ 上的坐标 $w=(\omega^{a})$, $a=1, \ldots, n-1$, 我们可以用流 $x^{\mu}(s, \omega)$ 参数化 $S_{s}$ 上的点, 该流由 $\frac{dx^{\mu}}{ds}=L^{\mu}$ 定义, 其中 $x^{\mu}(0, \omega)$ 是 $S_{0}$ 上坐标为 $w$ 的点. 令

\[
\gamma_{ab} = \gamma\left(\frac{\partial}{\partial\omega^{a}}, \frac{\partial}{\partial\omega^{b}}\right), \qquad \chi_{ab} = \chi\left(\frac{\partial}{\partial\omega^{a}}, \frac{\partial}{\partial\omega^{b}}\right)
\]

表示 $\gamma$ 与 $\chi$ 相对于这些坐标的分量. 容易验证 $\frac{d}{ds}\gamma_{ab} = 2\chi_{ab}$. $S_{s}$ 的体积元由

\[
da_{S_{s}} = \sqrt{|\gamma|}\, d\omega^{1}\dots d\omega^{n-1}
\]

给出, 其中 $\gamma$ 是度量 $\gamma$ 的行列式. 注意 $\frac{d}{ds}\log|\gamma| = \gamma^{ab}\frac{d}{ds}\gamma_{ab} = 2\operatorname{tr}\chi$, 其中 $\operatorname{tr}\chi = \gamma^{ab}\chi_{ab}$ 是零超曲面的\newterm{expansion-coefficient}{膨胀系数}{expansion coefficient}. 因此,

\[
\frac{d}{ds}\sqrt{|\gamma|} = \operatorname{tr}\chi\sqrt{|\gamma|}.
\]

总体积 $|S_{s}|$ 的变化率由如下公式给出:

\begin{equation}
\frac{d}{ds}|S_{s}| = \int_{S_{s}}\operatorname{tr}\chi\, da_{S_{s}}\label{eq:cA-439}
\end{equation}

我们还注意到 $\chi$ 满足如下 Riccati 型方程,

\begin{equation}
\frac{d}{ds}\chi + \chi^{2} = 0\label{eq:cA-440}
\end{equation}

它可以显式积分. 于是可以验证, 若在 $S_{0}$ 的某点处 $\operatorname{tr}\chi(0, \omega_{0}) < 0$, 则 $\operatorname{tr}\chi(s, \omega_{0})$ 可能在有限的 $s>0$ 值处变为 $-\infty$. 这一现象对应于焦散的形成.

N 上的任意叶状结构 $S_{v}$ 可以用 $v(s, \omega)$ 参数化, 其中 $(s, \omega)$ 是上面定义的测地坐标. 我们把 $\Omega = \frac{dv}{ds}$ 称为该叶状结构的\newterm{null-lapse-function}{零时距函数}{null lapse function}, 并用 $\gamma'$ 与 $\chi'$ 表示 $\gamma, \chi$ 在 $S_{v}$ 上的限制. 若 $X$ 是切于测地叶状结构 $S_{s}$ 的向量场, 则 $X' = X - \Omega^{-1}X(v)L$ 切于 $S_{v}$. 于是, 若 $X$、$Y$ 切于 $S_{s}$, 则 $\gamma(X, Y) = \gamma(X', Y')$, $\chi(X', Y') = \chi(X, Y)$. 相对于坐标 $(v, \omega)$, 我们有

\[
\gamma_{ab}^{\prime} = \gamma_{ab}, \qquad \chi_{ab}^{\prime} = \chi_{ab}.
\]

为了在零超曲面 $N$ 上定义体积元, 我们选择任意叶状结构 v, 其零时距函数为 $\frac{dv}{ds} = \Omega$, 诱导度量为 $\gamma$, 并令

\begin{equation}
da_{\mathcal{N}} = \Omega^{-1}da_{S_{v}}dv\label{eq:cA-441}
\end{equation}

其中 $da_{S_{v}}$ 表示由 $\gamma$ 在 $S_{v}$ 上诱导的面积元. 该定义不依赖于叶状结构的选取.

\mtsubsec{能量-动量张量}

$R^{n+1}$ 中的\newterm{energy-momentum-tensor}{能量-动量张量}{energy-momentum tensor}是满足正能量条件的对称二阶张量 $Q_{\alpha\beta}$:

\[
Q(X, Y) \geq 0,
\]

对一切因果的、未来定向的 $X$、Y. 若

\begin{equation}
D^{\beta}Q_{\alpha\beta} = 0\label{eq:cA-442}
\end{equation}

我们称 $Q$ 是\newterm{divergenceless}{无散度的}{divergenceless}. 给定任意向量场 X,

\[
D^{\alpha}(Q_{\alpha\beta}X^{\beta}) = Q^{\alpha\beta}D_{\alpha}X_{\beta} = \frac{1}{2}Q^{\alpha\beta}\, {}^{(X)}\pi_{\alpha\beta},
\]

其中 ${}^{(X)}\pi = \mathcal{L}_{X}\mathbf{m}$ 表示 $X$ 的变形张量. 回忆 ${}^{(X)}\pi_{\alpha\beta} = \partial_{\alpha}X_{\beta} + \partial_{\beta}X_{\alpha}$. 在 $X$ 是 Killing 向量场的特殊情形下, 即 ${}^{(X)}\pi = 0$, 我们推得

\begin{equation}
D^{\alpha}(Q_{\alpha\beta}X^{\beta}) = 0.\label{eq:cA-443}
\end{equation}

若 $X$ 是共形 Killing 的且 $Q$ 无迹, 即 $\mathbf{m}^{\alpha\beta}Q_{\alpha\beta} = 0$, 同样的恒等式成立.

一个典型的守恒律是这样得到的:在正则保守时空区域(见第 4.1 节)$\mathcal{J}^{+}(D_{1})\cap\mathcal{J}^{-}(D_{2})$ 上积分后面的恒等式并应用 Stokes 定理, 其中光滑类空边界 $D_{i}\subset\Sigma_{i}$, 零边界 $\mathcal{N}_i$, $i = 1, 2$. 我们用 $T_{1}, T_{2}$ 表示类空超曲面 $\Sigma_1, \Sigma_2$ 的未来单位法向量, 并选取零法向量 $L_{1}, L_{2}$, 使得沿边界 $D_{i}\subset\Sigma_{i}$, $i = 1, 2$ 有 $\mathbf{m}(L_i, T_i) = -1$. 为简单起见, 只要不会引起混淆, 我们用 $T$ 表示两个类时法向量, 用 $L$ 表示两个零法向量.

\begin{proposition}\label{prop:cA-4.6}
设 $Q_{\alpha\beta}$ 是无散度的能量-动量张量, X 是上述正则保守区域 $\mathcal{J}(D_{1}, D_{2})$ 邻域中的 Killing 向量场. 则

\begin{equation}
\int_{\mathcal{N}_{2}}Q(X, L) + \int_{D_{2}}Q(X, T) = \int_{\mathcal{N}_{1}}Q(X, L) + \int_{D_{1}}Q(X, T)\label{eq:cA-444}
\end{equation}
\end{proposition}

积分分别关于零超曲面 $N_{1}, N_{2}$ 上的面积元 $da_{N}$ 以及 $m$ 在 $\Sigma_{1}, \Sigma_{2}$ 上诱导的 Riemann 度量的面积元进行. 注意, 若 $X$ 是因果的, 则所有被积函数都是正的. 若 $X$ 是共形 Killing 的且 $Q$ 无迹, 恒等式 \eqref{eq:cA-444} 仍然成立.

\begin{proof}
设 $P_{\alpha} = Q_{\alpha\beta}X^{\beta}$. 根据 \eqref{eq:cA-443} 有 $D^{\alpha}P_{\alpha} = 0$.
\end{proof}

对于依赖域 $\mathcal{J}^{+}(\Sigma_{1})\cap\mathcal{J}^{-}(D\subset\Sigma_{2})$, 或影响域 $\mathcal{J}^{+}(D\subset\Sigma_{1})\cap\mathcal{J}^{-}(\Sigma_{2})$($\Sigma_{2}$ 位于 $\Sigma_{1}$ 的未来), 结果可以简化. 我们在 $\partial D\subset\Sigma_{2}$ 上以条件 $\mathbf{m}(L, T) = -1$ 规范化 L, 其中 $T$ 表示 $\Sigma_{1}, \Sigma_{2}$ 的单位法向量.

\begin{corollary}\label{cor:cA-4.7}
若 $Q$ 无散度, X 是 Killing 的, 且 $D \subset \Sigma_2$, 则

\begin{equation}
\int_{\mathcal{N}}Q(X, L) + \int_{D \subset \Sigma_{2}}Q(X, T) = \int_{\mathcal{J}^{-}(D)\cap\Sigma_{1}}Q(X, T)\label{eq:cA-445}
\end{equation}

类似地, 若 $D\subset\Sigma_{1}$, 则

\begin{equation}
\int_{\mathcal{N}}Q(X, L) + \int_{D \subset \Sigma_{1}}Q(X, T) = \int_{\mathcal{J}^{+}(D)\cap\Sigma_{2}}Q(X, T)\label{eq:cA-446}
\end{equation}

若 $X$ 是共形 Killing 的且 $Q$ 无迹, 恒等式仍然成立.
\end{corollary}

\section*{参考文献}
\markboth{参考文献}{参考文献}
\addcontentsline{toc}{section}{参考文献}
\begin{mybibliography}{Wa}
\bibitem[Car]{car} S. M. Carroll, \emph{Lecture Note on General Relativity}, arXiv:gr-qc/9712019v1, Dec 1997.
\bibitem[E]{e} L. C. Evans, \emph{Partial Differential Equations}, Providence: American Mathematical Society, 1998.
\bibitem[Fr-Io]{fr-io} F. G. Friedlander and M. Ioshi, \emph{An introduction to the theory of distributions}, Cambridge University Press, second edition, 1998.
\bibitem[G-T]{g-t} D. Gilbarg and N. Trudinger, \emph{Elliptic partial differential equations of second order}, Springer Verlag, second edition, 1983.
\bibitem[Ge-S]{ge-s} I. M. Gelfand and G. E. Shilov, \emph{Generalized Functions, Volume I: Properties and Operations}, Academic Press, 1964.
\bibitem[Ho]{ho} L. Hörmander, \emph{The Analysis of linear partial differential operators}, vols 1-4, Springer Verlag, 1983, 1985.
\bibitem[Ha-E]{ha-e} S. W. Hawking and G. F. R. Ellis, \emph{The Large Scale Structure of Space-Time}, Cambridge University Press, 1973.
\bibitem[J]{j} F. John, \emph{Partial Differential Equations}, Providence: American Mathematical Society, fourth edition, 1982.
\bibitem[K-T]{k-t} Markus Keel and Terrence Tao, ``Endpoint Strichartz Estimates'', \emph{Am. J. Math.} 120 (1998), 955--980.
\bibitem[Kl-Rodn2]{kl-rodn2a} S. Klainerman and I. Rodnianski, Causal geometry of Einstein-vacuum spacetimes with finite curvature flux, \emph{Invent. Math.} 159 (2005), 437-529.
\bibitem[Kl-Rodn2]{kl-rodn2b} S. Klainerman and I. Rodnianski, A geometric approach to the Littlewood-Paley theory, \emph{Geom. Funct. Anal.} 16 (2006), 126--163.
\bibitem[Kl-Rodn3]{kl-rodn3} S. Klainerman and I. Rodnianski, Sharp trace theorems for null hypersurfaces on Einstein metrics with finite curvature flux, \emph{Geom. Funct. Anal.} 16 (2006), 164--229.
\bibitem[Pet]{pet} P. Petersen, \emph{Riemannian Geometry}, GTM 171, Springer, 1998.
\bibitem[Sp]{sp} M. Spivak, \emph{A Comprehensive introduction to Differential Geometry}, vols I-V, Wilmington: Publish or Perish.
\bibitem[T]{t} M. E. Taylor, \emph{Partial Differential Equations}, Springer Verlag, Texts in Applied Mathematics, vol. 23, 1996.
\bibitem[Wa]{wa} R. M. Wald, \emph{General Relativity}, The University of Chicago Press, 1984.
\end{mybibliography}


% ================================================================
%   所附论文已单独抽取排版, 正书不再收录
%   原书末尾附论文《几何与 Fourier 方法在非线性波动方程中的应用》
%   (GEOMETRIC AND FOURIER METHODS FOR NONLINEAR WAVE EQUATIONS).
%   该论文为 article 而非 book, 已用简洁的 article 模板单独重排为中译,
%   见独立文档 ipam/main.pdf (源文件 ipam/ 目录: main.tex + lec1..6.tex + references.tex).
% ================================================================
\clearpage

% ================================================================
%                       术语索引: 全文最后
% ================================================================
\printterminology

\end{document}
